










\documentclass[acmjacm]{acmsmall}

\usepackage[shortlabels]{enumitem}
 
\usepackage{latexsym,xspace,calc,amsthm}
\usepackage{amsmath,amssymb} %
\usepackage{centernot}

\setcounter{tocdepth}{4}
\setcounter{secnumdepth}{3}


\usepackage{lmodern}
\usepackage{tgtermes}
\usepackage{fix-cm}


\usepackage{cancel}

\usepackage{savesym}
\savesymbol{subcaption}
\usepackage{subcaption}

\usepackage{graphbox} %

\makeatletter
\def\@biblabel#1{[#1]} %
\def\thebibliography#1{%
    \footnotesize
    \refsection*{{\refname}
        \@mkboth{\uppercase{\refname}}{\uppercase{\refname}}%
    }
    \list{\@biblabel{\@arabic\c@enumiv}}%
       {\settowidth\labelwidth{\@biblabel{#1}}%
        \leftmargin\labelwidth
        \advance\leftmargin\bibindent
        \itemindent-\bibindent
        \itemsep2pt
        \parsep \z@
        \usecounter{enumiv}%
        \let\p@enumiv\@empty
        \renewcommand\theenumiv{\@arabic\c@enumiv}%
    }%
    \let\newblock\@empty
    \sloppy
    \sfcode`\.=1000\relax
}
\makeatother

\let\odlthebibliography=\thebibliography
 \let\endodlthebibliography=\endthebibliography
 \renewenvironment{thebibliography}[1]{%
   \begin{odlthebibliography}{#1}%
     \setlength{\parskip}{0ex}%
     \setlength{\itemsep}{3pt}%
     \fontsize{9.5}{9.5} %
     \selectfont
}%
 {%
   \end{odlthebibliography}%
 }


\usepackage{tikz}
\usepackage{tikz-cd}
\usetikzlibrary{cd,decorations.markings}
\tikzset{
    dharrow/.style={
        <->,
        postaction={decorate,-},
        }
}
\tikzset{
    dhdashedarrow/.style={
        <->,
        dashed,
        postaction={decorate,-},
        }
    }
\tikzset{
    lrharpoonarrow/.style={
        <[harpoon]->[harpoon],
        postaction={decorate,-},
        }
}
\tikzset{
    lrharpoondashedarrow/.style={
        <[harpoon]->[harpoon],
        dashed, %
        postaction={decorate,-},
        }
}
\usetikzlibrary{arrows}
\usetikzlibrary {arrows.meta} 
\usetikzlibrary{calc}
\usetikzlibrary{positioning}
\usetikzlibrary{snakes,automata,chains}
\usetikzlibrary{graphs}

\usepackage{binarytree}

\usepackage{amssymb,stmaryrd,amsmath}
\usepackage{mdwlist} %
\usepackage{float}   %
\usepackage{centernot} %


\usepackage[colorlinks]{hyperref}
\usepackage{breakurl}  %


\usepackage{bm} %





\newtheoremstyle{jamiestyle}%
  {4pt}%
  {0pt}%
  {\it}%
  {0pt}%
  {\sc}%
  {.}%
  { }%
  {}%
\theoremstyle{jamiestyle}
\newtheorem{thrm}{Theorem}[subsection]
\newtheorem{prop}[thrm]{Proposition}
\newtheorem{lemm}[thrm]{Lemma}
\newtheorem{corr}[thrm]{Corollary}
\newtheorem{conj}[thrm]{Conjecture}
\newtheoremstyle{jamienfstyle}%
  {4pt}%
  {0pt}%
  {\normalfont}%
  {0pt}%
  {\sc}%
  {.}%
  { }%
  {}%
\theoremstyle{jamienfstyle}
\newtheorem{nttn}[thrm]{Notation}
\newtheorem{defn}[thrm]{Definition}
\newtheorem{xmpl}[thrm]{Example}
\newtheorem{rmrk}[thrm]{Remark}


\usepackage{color}
\definecolor{mygreen}{rgb}{0,0.6,0}
\definecolor{mygray}{rgb}{0.5,0.5,0.5}
\definecolor{mymauve}{rgb}{0.58,0,0.82}

\usepackage{listings}
 
\definecolor{gray}{RGB}{128, 128, 128}
\definecolor{lightgray}{RGB}{200, 200, 200}
\definecolor{cyan}{RGB}{0, 255, 255}
\definecolor{blue}{RGB}{0, 0, 255}
\definecolor{red}{RGB}{255, 0, 0}
\definecolor{pink}{RGB}{255, 128, 128}
\definecolor{green}{RGB}{0, 128, 0}
\definecolor{lightyellow}{RGB}{255, 255, 200}
\definecolor{purple}{RGB}{128, 0, 128}

\lstdefinestyle{all}
    {basicstyle=\ttfamily\scriptsize,
     keywordstyle=\color{blue}\ttfamily\scriptsize,
     commentstyle=\color{green}\ttfamily\scriptsize,
     stringstyle=\color{red}\ttfamily\scriptsize}

\lstdefinelanguage{hask}{%
    frame=none,
    xleftmargin=2pt,
    belowcaptionskip=\bigskipamount,
    captionpos=b,
    tabsize=2,
    numbers=left,
    numberstyle=\tiny\color{gray},
    emphstyle={\bf},
	morecomment=[s][\color{green}]{\{-}{-\}},
    stringstyle=\mdseries\rmfamily,
    commentstyle=\color{green},
    keywords={},
    keywords=[1]{case, of, data, if, then, else, where, let, in, do},
    keywords=[2]{Chip, Config, CurrencySymbol, TokenName, PubKeyHash, Integer, Value, State, Action, Text, Maybe, Void, TxConstraints,  Contract},
    keywords=[3]{HasNative},
    keywords=[4]{=>},
    keywords=[5]{Just, Nothing, MkChip, MkConfig, SetPrice, Buy},
    keywordstyle=[1]\mdseries\sffamily\color{red},
    keywordstyle=[2]\mdseries\sffamily\color{blue},
    keywordstyle=[3]\mdseries\sffamily\color{green},
    keywordstyle=[4]\mdseries\sffamily,
    keywordstyle=[5]\mdseries\sffamily\color{purple},
    columns=flexible,
    basicstyle=\small\sffamily,
    showstringspaces=false,
    breaklines=false,
    showspaces=false,
    escapeinside={--}{\^^M},escapebegin={\color{green}--},escapeend={},
    literate= {+}{{$+$}}1 {/}{{$/$}}1 {*}{{$*$}}1 {=}{{$=$}}1
              {>}{{$>$}}1 {<}{{$<$}}1 {\\}{{$\lambda$}}1
              {\\\\}{{\char`\\\char`\\}}1
              {->}{{$\rightarrow$}}2 {>=}{{$\geq$}}2 {<-}{{$\leftarrow$}}2
              {<=}{{$\leq$}}2 {=>}{{$\Rightarrow$}}2
              {\ .}{{$\circ$}}2 {\ .\ }{{$\circ$}}2
              {>>}{{>>}}2 {>>=}{{>>=}}2
              {|}{{$\mid$}}1
              {\_}{{\underline{\hspace{2mm}}}}2
}

\lstdefinelanguage{solidity}{%
    frame=none,
    xleftmargin=2pt,
    belowcaptionskip=\bigskipamount,
    captionpos=b,
    tabsize=2,
    numbers=left,
    numberstyle=\tiny\color{gray},
    emphstyle={\bf},
	morecomment=[s][\color{green}]{\{-}{-\}},
    stringstyle=\mdseries\rmfamily,
    commentstyle=\color{green},
    keywords={},
    keywords=[1]{pragma, solidity, contract, event, constructor, require, function, return, emit},
    keywords=[2]{address, uint, mapping},
    keywords=[3]{public, payable, external, view, returns},
    keywords=[4]{=>, +=, -=, =, <=, ==},
    keywords=[5]{msg, sender, transfer, value},
    keywordstyle=[1]\mdseries\sffamily\color{red},
    keywordstyle=[2]\mdseries\sffamily\color{blue},
    keywordstyle=[3]\mdseries\sffamily\color{green},
    keywordstyle=[4]\mdseries\sffamily,
    keywordstyle=[5]\mdseries\sffamily\color{purple},
    columns=flexible,
    basicstyle=\small\sffamily,
    showstringspaces=false,
    breaklines=false,
    showspaces=false,
    escapeinside={--}{\^^M},escapebegin={\color{green}--},escapeend={},
    literate= {+}{{$+$}}1 {/}{{$/$}}1 {*}{{$*$}}1 {=}{{$=$}}1
              {>}{{$>$}}1 {<}{{$<$}}1 {\\}{{$\lambda$}}1
              {\\\\}{{\char`\\\char`\\}}1
              {->}{{$\rightarrow$}}2 {>=}{{$\geq$}}2 {<-}{{$\leftarrow$}}2
              {<=}{{$\leq$}}2 {=>}{{$\Rightarrow$}}2
              {\ .}{{$\circ$}}2 {\ .\ }{{$\circ$}}2
              {>>}{{>>}}2 {>>=}{{>>=}}2
              {|}{{$\mid$}}1
              {\_}{{\underline{\hspace{2mm}}}}2
}




\newcommand\flanks{\ltimes}
\newcommand\setminusskip{\hskip 2pt plus 1pt minus 2pt}
\newcommand\mysetminus{\setminusskip{{\setminus}}\setminusskip}

\newcommand\declaresoundness[2]{The case of \rulefont{#1} with $#2$.}
\newcommand\declaresoundnessshort[1]{The case of \rulefont{#1}.}


\makeatletter
\newcommand\hpn[2][]{%
  \ext@arrow 9999{\hpnfill@}{#1}{#2}}
\newcommand\hpnfill@{%
  \arrowfill@\leftharpoonup\relbar\rightharpoondown}
\makeatother

\newcommand\arrowlength{0.5}
\newcommand\arrowr{{\tikz \draw [semithick,->] (0,0) -- (\arrowlength*0.8,0);}}
\newcommand\Arrowr{{\tikz \draw [semithick,->,decorate,decoration={zigzag,segment length=3,amplitude=0.5}] (0,0) -- (\arrowlength*0.8,0);}}
\newcommand\Arrowrr{{\tikz \draw [semithick,->>,decorate,decoration={zigzag,segment length=3,amplitude=0.5}] (0,0) -- (\arrowlength*0.8,0);}}
\newcommand\arrowlr{{\tikz \draw [semithick,<->] (0,0) -- (\arrowlength*0.8,0);}}
\newcommand\arrowllrr{{\tikz \draw [semithick,<<->>] (0,0) -- (\arrowlength,0);}}
\newcommand\arrowrr{{\tikz \draw [semithick,->>] (0,0) -- (\arrowlength*0.8,0);}}
\newcommand\arrowrrr{{\tikz \draw [thick,-Stealth] (0,0) -- (\arrowlength*0.8,0);}}
\newcommand\arrowll{{\tikz \draw [semithick,<<-] (0,0) -- (\arrowlength*0.8,0);}}
\newcommand\arrowerr{{\tikz \draw [semithick,-.>>] (0,0) -- (\arrowlength*0.8,0);}}
\newcommand\arrowlrr{{\tikz \draw [semithick,<->>] (0,0) -- (\arrowlength,0);}}
\newcommand\arrowllr{{\tikz \draw [semithick,<<->] (0,0) -- (\arrowlength,0);}}
\newcommand\arrowelrr{{\tikz \draw [semithick,<.->>] (0,0) -- (\arrowlength,0);}}
\newcommand\arrowelr{{\tikz \draw [semithick,<.->] (0,0) -- (\arrowlength,0);}}
\newcommand\arrower{{\tikz \draw [semithick,-.>] (0,0) -- (\arrowlength*0.8,0);}}


\newcommand\everyone{\modality{\Box}}
\newcommand\someone{\modality{\Diamond}}
\newcommand\someoneAll{{\mathring{\exists}_{\hspace{-0pt}\someone}}}
\newcommand\everyoneAll{{\mathring{\forall}_{\hspace{-1pt}\everyone}}}

\newcommand\Quorum{\modality{\mathring{\exists}}}
\newcommand\QuorumBox{\modality{\mathring{\exists}_{\Box}}} %
\newcommand\Coquorum{\modality{\mathring{\forall}}}
\newcommand\CoquorumDiamond{\modality{\mathring{\forall}_{\hspace{-1pt}\Diamond}}} %


\newcommand{\ite}[2]{[#1,\ \text{\it o\hspace{-1pt}'wise}\ #2]}
\newcommand\mruarrow{\mathbin{\rlap{\scalebox{0.73}{$\recentsimple$}}\hspace{-2pt}\rightarrow}}
\newcommand{\mru}[2]{#1\mruarrow#2}
\newcommand{\mrup}[3]{#2[#1\mruarrow #3]}
\newcommand{\mruX}[2]{{#1\mathbin{\recent}#2}}
\newcommand{\established}[2]{#1 \mathbin{\Downarrow_{\QuorumBox}} #2}
\newcommand{\lrharpoon}{\mathbin{\hpn{}}}
\newcommand{\nlrharpoon}{\mathbin{\not\hpn{}}}
\newcommand\id{\f{id}}
\newcommand\st{{:}}

\usepackage{tikz}
\newcommand*{\myhash}{%
  \begin{tikzpicture}
    \pgfmathsetlengthmacro\myWidth{.8*width("=")}%
    \pgfmathsetlengthmacro\myHeight{height("H")}%
    \pgfmathsetlengthmacro\mySepY{.3333*\myWidth}%
    \pgfmathsetlengthmacro\mySideBearing{.1*\myWidth}%
    \def\myAngle{70}%
    \pgfmathsetlengthmacro\mySepX{\mySepY/sin(\myAngle)}%
    \pgfmathsetlengthmacro\mySlantX{\myHeight/tan(\myAngle)}%
    \draw[line cap=round]
      (0, {(\myHeight - \mySepY)/2}) -- ++(\myWidth, 0)
      (0, {(\myHeight + \mySepY)/2}) -- ++(\myWidth, 0)
      ({(\myWidth - \mySepX - \mySlantX)/2}, 0)
      -- ({(\myWidth - \mySepX + \mySlantX)/2}, \myHeight)
      ({(\myWidth + \mySepX - \mySlantX)/2}, 0)
      -- ({(\myWidth + \mySepX + \mySlantX)/2}, \myHeight)
    ;%
    \useasboundingbox
      (-\mySideBearing, 0)
      (\myWidth + \mySideBearing, \myHeight)
    ;%
  \end{tikzpicture}%
}

\newcommand\THREE{{\mathbf 3}}
\newcommand\threeValid{\THREE_{\f{valid}}}
\newcommand\threeCorrect{\THREE_{\f{correct}}}
\newcommand\threeTrue{\THREE_{\f{true}}}
\newcommand\threeFalse{\THREE_{\f{false}}}
\newcommand\binaryconnectives{\ensuremath{\{\tand,\tor,\tnotor,\timpc\}}} %
\newcommand\unaryconnectives{\ensuremath{\{\tneg,\modT,\modB,\modF,\modTB,\modTF\}}}
\newcommand\modalityconnectives{\ensuremath{\{\someone,\everyone,\Quorum,\Coquorum,\sometime,\forever,\tomorrow\}}} %
\newcommand\doublemodalities{\ensuremath{\{\QuorumBox,\CoquorumDiamond,\someoneAll,\everyoneAll\}}}
\newcommand\singlemodalities{\ensuremath{\{\someone,\everyone,\Quorum,\Coquorum\}}}
\newcommand\Pnt{\ns{Point}}
\newcommand\Val{\ns{Value}}
\newcommand\Valuation{\ns{Valuation}}
\newcommand\PredSymb{\ns{PredSymb}}

\newcommand\FOUR{{\mathbf 4}}
\newcommand\tvT{{\mathbf t}}
\newcommand\tvF{{\mathbf f}}
\newcommand\tvB{{\mathbf b}}
\newcommand\tvN{{\mathbf n}}
\newcommand\tvsTB{{\mathit{tb}}}
\newcommand\tvsFB{{\mathit{fb}}}
\newcommand\tvsTT{{\mathit{tt}}}
\newcommand\tvsFF{{\mathit{ff}}}
\newcommand*\Rcirc{\ensuremath{\smash{\overset{\circ}{\smash{\mathsf{R}}\rule{0pt}{1.2ex}}\vphantom{\mathbb{R}}}}}
\newcommand\recentsimple{{\reflectbox{$\mathsf R$}}} 
\newcommand\recent{\recentsimple} %
\newcommand\urecent{{\ubar\recentsimple}} %
\newcommand\mrecent[2]{#1\recentsimple #2}

\newcommand\figunderline[1]{\text{\rlap{\underline{\emph{#1}}}}}




\usepackage{stackengine}
\newcommand\ubar[1]{\stackunder[0.6pt]{$#1$}{\rule{1ex}{.18ex}\hspace{0.5pt}}}


\newcommand\figskip{\\[1ex]}
\newcommand\from[1]{{\uparrow}_{#1}}
\newcommand\fs[1]{\mathtt{#1}}
\newcommand\ps[1]{\mathtt{#1}}
\newcommand\thea{a}
\newcommand\theb{b}
\newcommand\thec{c}
\newcommand\aval{v}
\newcommand\xor{\mathbin{\mathsf{\small xor}}}
\newcommand\modality[1]{#1\hspace{1.5pt}}
\newcommand\modT{\modality{\tf T}}
\newcommand\modTB{\modality{\tf T\hspace{-2.5pt}\tf B}}
\newcommand\modFB{\modality{\tf F\hspace{-1.5pt}\tf B}}
\newcommand\modB{\modality{\tf B}}
\newcommand\modTF{\modality{\tf T\hspace{-2.5pt}\tf F}}
\newcommand\modF{\modality{\tf F}}
\newcommand\witno{\prec}
\newcommand\atopen{T}
\newcommand\afilter{F}
\newcommand\apoint{P}
\newcommand\acontext{\f{ctx}}
\newcommand\avaluation{\varsigma} %
\newcommand\indicator[1]{\delta_{#1}} 
\newcommand\charop[1]{\f{char}(#1)}
\newcommand\resilient{R}
\newcommand\rulefont[1]{\ensuremath{{\mathsf{(#1)}}}}
\newcommand\leftopeninterval[1]{(#1]}
\newcommand\rightopeninterval[1]{[#1)}
\newcommand\openinterval[1]{(#1)}
\newcommand\closedinterval[1]{[#1]}
\newcommand\opens{\ns{Open}}
\newcommand\opensne{\opens_{\hspace{-1pt}{\neq}\varnothing}}
\newcommand\regularOpens{\tf{Open}_{\f{reg}}}
\newcommand\topens{\tf{Topen}}
\newcommand\closed{\tf{Closed}}
\newcommand\regularClosed{\tf{Closed}_{\f{reg}}}
\newcommand\thenext[0]{\f{ad}}
\newcommand\closure[1]{|#1|}
\newcommand\dotclosure[1]{\closure{#1}^{\bigcdot}}
\newcommand{\dotarrow}{%
   \mathrel{\ooalign{\hss\raise.85ex\hbox{\scalebox{1.25}{\normalfont .}}%
   \kern0.35ex\hss\cr$\rightarrow$}}}
\newcommand{\dotleftarrow}{\mathrel{\reflectbox{$\dotarrow$}}}
\newcommand\mycard{\myhash}
\newcommand\narrowmath[1]{\scalebox{0.8}[1.0]{$\mathsf{#1}$}}
\newcommand\sometime{\modality{\narrowmath{sometime}}}
\newcommand\forever{\modality{\narrowmath{forever}}}
\newcommand\tomorrow{\modality{\narrowmath{tomorrow}}}
\newcommand\yesterday{\modality{\narrowmath{yesterday}}}
\newcommand\infinitely{\modality{\narrowmath{infinitely}}}
\newcommand\final{\modality{\narrowmath{finally}}}
\newcommand\arity{\f{arity}}
\newcommand\correct[1]{\tf{correct}[#1]}
\newcommand\correctterm{correct\xspace}
\newcommand\pointwise[1]{\tf{pointwise}[#1]}
\newcommand\pointwiseterm{pointwise\xspace}
\newcommand\incorrect[1]{\tf{incorrect}[#1]}
\newcommand\uniform[1]{\tf{uniform}[#1]}
\newcommand\udfn{\tf{undef}}
\newcommand\gslt{\f{GSLT}}
\newcommand\dense{\f{dense}}
\newcommand\noi{\f{noi}}

\newcommand{\upperset}[2]{
  \underset{
            \text{\raisebox{0.9ex}{\smash{\scalebox{0.8}{$#1$}}\ \ }}%
           }
           {
            \text{\raisebox{0ex}{\smash{$#2$}}}
           }
                       }
\newcommand{\convergesto}[1]{\mathrel{\upperset{#1}{\dotarrow}}}
\newcommand\betwixt{\asymp}
\newcommand\notbetween{\mathbin{\cancel{\between}}}
\newcommand\notintertwinedwith{\mathrel{\notbetween}}
\newcommand\intersectswith{\between}
\newcommand\notintersectswith{\notbetween}
\newcommand\stronglybetween{\mathbin{\rlap{$\between$}\hspace{-1.5pt}\between}}
\newcommand\xxtwined{$2$-twined\xspace}
\newcommand\xxxtwined{$3$-twined\xspace}
\newcommand\intertwined[1]{#1_{\between}}
\newcommand\intertwinedwith{\mathrel{\between}}
\newcommand\stronglyintertwinedwith{\mathrel{\stronglybetween}}
\newcommand\leqk{\leq_{\hspace{-.7pt}\intertwinedwith}}
\newcommand\geqk{\geq_{\hspace{-.7pt}\intertwinedwith}}
\newcommand\cw{\leftrightarrow}  %
\newcommand\ctr{\rightarrow}  %
\newcommand\notctr{{\not\rightarrow}}  %
\newcommand\ctl{\leftarrow}  %
\newcommand\cti{\leq}  %
\newcommand\ctieq{{\stackrel{=}{\hookrightarrow}}}  %
\newcommand\notcti{{\not\hookrightarrow}}  %
\newcommand\nbhd[0]{\f{nbhd}}
\newcommand\interior[0]{\f{interior}}
\newcommand\kiss[0]{\f{kiss}}
\newcommand\community[0]{\f{K}}
\newcommand\framecommunity[0]{\f{k}}
\newcommand\kernel[0]{\f{ker}}
\newcommand\Kernel[0]{\f{Ker}}
\newcommand\witness[0]{\f{W}}
\newcommand\Kmod[1]{\tf K_{\hspace{-2pt}#1}}
\newcommand\cast[1]{#1^{\ast c}}
\newcommand\cclo[1]{#1^c}
\newcommand\mbot{\mathsf{e}}
\newcommand\mtop{\mathsf{f}}
\newcommand\isdef[2]{#1\bullet #2}
\newcommand\deffont[1]{{\bfseries #1}}
\newcommand\powerset{\f{pow}}
\newcommand\comppowerset{\f{pow}_{\hspace{-2pt}c}}
\newcommand\finpow{\f{fin}}
\newcommand\finto{\stackrel{\f{fin}}{\rightharpoonup}}
\newcommand\f[1]{\mathit{#1}}
\newcommand\tf[1]{\mathsf{#1}}
\DeclareMathAlphabet{\mathsb}{OT1}{cmss}{sbc}{n}
\newcommand\ns[1]{{\mathsb{#1}}}
\usepackage{slantsc}
\newcommand\theory[1]{\ensuremath{\text{\itshape\scshape #1}}\xspace}
\newcommand\ThyPax{\theory{TyPax}}
\newcommand\ThySPax{\theory{TySPax}}
\newcommand\ThyBB{\theory{TyBB}}
\newcommand\finsubseteq{\mathbin{\subseteq_{\text{\it fin}}}}
\newcommand\at{\text{@}}
\newcommand\tx{\f{tx}}
\newcommand\ty{\f{ty}}
\newcommand\ctx{\f{ctx}}
\newcommand\txs{\f{txs}}
\newcommand\utxi{\f{utxi}}
\newcommand\utxo{\f{utxo}}
\newcommand\stx{\f{stx}}
\newcommand\utxio{\f{utxio}}
\newcommand\UTxIO{\f{UTxIO}}
\newcommand\inlinehask[1]{\lstinline[language=hask]{#1}}
\newcommand\inlinesolidity[1]{\lstinline[language=solidity]{#1}}
\newcommand\valid{\f{valid}}
\newcommand\liff{\Longleftrightarrow}
\newcommand\limp{\Longrightarrow}
\newcommand\nlimp{\centernot{\Longrightarrow}}
\newcommand\aeq{\mathrel{=_{\alpha}}}
\newcommand\ssm{{{:}\text{=}}}
\DeclareMathSymbol{\shortminus}{\mathbin}{AMSa}{"39}
\newcommand\minus{{\shortminus}}
\newcommand\plus{{+}}
\newcommand\Forall[1]{\forall #1.}
\newcommand\Exists[1]{\exists #1.}
\newcommand\sm{\mapsto}
\newcommand\rewritesto{\longrightarrow}
\newcommand\cent{\vdash}
\newcommand\ncent{\not\vdash}
\newcommand\ment{\vDash}
\newcommand\nment{\not\vDash}
\newcommand\mentval{\ment_{\avaluation}}
\newcommand\nmentval{\nment_{\avaluation}}
\newcommand\supp{\f{supp}}
\newcommand\boundary{\f{boundary}}
\newcommand\lmodel{[\hspace{-0.2em}[}
\newcommand\rmodel{]\hspace{-0.2em}]}
\newcommand\model[1]{{\lmodel #1 \rmodel}}
\newcommand\modellabel[2]{{\lmodel #1 \rmodel}_{#2}}
\newcommand\synmodel[1]{{\langle #1 \rangle}}
\newcommand\mone{{\text{-}1}}
\newcommand\fv{\f{fv}}
\newcommand\Ngeqz{\mathbb N_{\scriptscriptstyle\geq 0}}
\newcommand\Time{\Ngeqz}
\newcommand\weakmodusponens{weak modus ponens (Proposition~\ref{prop.mp.for.tnotor}(\ref{item.mp.for.tnotor}))\xspace}
\newcommand\strongmodusponens{strong modus ponens (Proposition~\ref{prop.mp.for.tnotor}(\ref{item.mp.for.timpc}))\xspace}
\newcommand\weakmodusponensnoref{weak modus ponens\xspace}
\newcommand\strongmodusponensnoref{strong modus ponens\xspace}

\makeatletter
\DeclareRobustCommand{\barcent}{\mathbin{\mathpalette\barcent@@\relax}}
\newcommand{\barcent@@}[2]{%
  \vbox{\offinterlineskip
    \sbox\z@{$\m@th#1\cent$}%
    \ialign{%
      \hfil##\hfil\cr
      $\m@th#1{}_{\minus}\kern-\scriptspace$\cr
      \noalign{\kern-.3\ht\z@}
      \box\z@\cr
    }%
  }%
}
\makeatother

\makeatletter
\def\pmb@#1#2{\setbox8\hbox{$\m@th#1{#2}$}%
  \setboxz@h{$\m@th#1\mkern-.1mu$}\pmbraise@\wdz@
  \binrel@{#2}%
  \dimen@-\wd8 %
  \binrel@@{%
    \mkern-.1mu\copy8 %
    \kern\dimen@\mkern-.2mu\copy8 %
    \kern\dimen@\mkern-.3mu\copy8 %
    \kern\dimen@\mkern-.4mu\copy8 %
    \kern\dimen@\mkern.1mu\copy8 %
    \kern\dimen@\mkern.2mu\copy8 %
    \kern\dimen@\mkern.3mu\copy8 %
    \kern\dimen@\mkern.0mu\raise\pmbraise@\copy8 %
    \kern\dimen@\mkern.4mu\box8 %
           }%
}
\makeatother


\newcommand\overbracebig[2]{\overbrace{#2}^{\mbox{$#1$}}}
\newcommand\compressthislight[1]{{\hspace{.8pt}\raisebox{.5pt}{\scalebox{.85}{$#1$}}\hspace{.2pt}}}
\newcommand\compressthis[1]{\pmb{\compressthislight{#1}}}
\newcommand\tempty{\compressthis{\varnothing}}
\newcommand\tcap{{\compressthis{\cap}}}
\newcommand\tcup{{\compressthis{\cup}}}
\newcommand\tsubseteq{{\compressthis\subseteq}}
\newcommand\tneg{{\pmb\neg}}
\newcommand\tnegnf{{\pmb{\sim}}}
\newcommand\tnegneg{{\pmb{{\sim}{\sim}}}}
\newcommand\ttop{{\pmb\top}}
\newcommand\tbot{{\pmb\bot}}
\newcommand\teq{{\pmb{\text{=}}}}
\newcommand\tneq{{\pmb{\neq}}}
\newcommand\tand{\mathrel{\pmb\wedge}}
\newcommand\tor{\mathbin{\pmb\vee}}
\newcommand\tgg{{\compressthis\gg}}
\newcommand\tleq{\compressthis{\leq}}
\newcommand\tlesssim{\compressthis{\lesssim}}
\newcommand\tapprox{\compressthis{\approx}}
\newcommand\timp{{\compressthis{\Rightarrow}}}
\newcommand\ntimp{{\pmb\not\Rightarrow}}
\newcommand\tnotor{\mathrel{\arrowr}}
\newcommand\tnotore{\mathrel{\arrower}}
\newcommand\tlatticeiff{\mathrel{\arrowlr}}
\newcommand\tlatticeiffec{\mathbin{\arrowelr}}
\newcommand\timpc{\mathrel{\arrowrr}}
\newcommand\timpcc{\mathrel{\arrowrrr}}
\newcommand\timpcleft{\mathrel{\arrowll}}
\newcommand\timpec{\mathrel{\arrowerr}}
\newcommand\tlatticeiffc{{\pmb{\bar\equiv}}}
\newcommand\tiff{\compressthis{\Leftrightarrow}} %
\newcommand\tiffc{\mathbin{\arrowlrr}}
\newcommand\tiffcc{\mathbin{\arrowllrr}}
\newcommand\tiffec{\mathbin{\arrowelrr}}
\newcommand\tlam{{\pmb\lambda}}
\newcommand\tsim{\compressthis{{\approx}}}
\newcommand\tin{{\pmb{\in}}}
\newcommand\ntin{{\pmb{\not\in}}}
\newcommand\qall{{\compressthis{\mathsf Q}}}
\newcommand\tall{{\compressthis{\forall}}}
\newcommand\texi{{\compressthis{\exists}}}
\newcommand\texiunique{{\compressthis{\exists\hspace{0.5pt}!}}}
\newcommand\texiaffine{{\compressthis{\exists_{01}}}}
\newcommand\texiR{{\compressthis{\exists_{\mathsf R}}}}
\newcommand\tlot{\textbf{?}}
\newcommand\AllBut[2]{\tf{AllBut}_{#2}(#1)}
\newcommand\Nset[1]{\{0\dots #1\minus 1\}}

\newcommand\oimplies{\mathbin{\rlap{$+$}\hspace{.7pt}{\supset}}}

\newcommand\QLogic{Coalition Logic\xspace}

\makeatletter
\newcommand{\circlearrow}{}%
\DeclareRobustCommand{\circlearrow}{%
  \mathrel{\vphantom{\shortrightarrow}\mathpalette\circle@arrow\relax}%
}
\newcommand{\circle@arrow}[2]{%
  \m@th
  \ooalign{%
    \hidewidth$#1\circ\mkern1mu$\hidewidth\cr
    $#1\longrightarrow$\cr}%
}
\makeatother
\newcommand\minusoplus{{\pmb{\circlearrow}}} %

\makeatletter
\newcommand*\bigcdot{\mathpalette\bigcdot@{.5}}
\newcommand*\bigcdot@[2]{\mathbin{\vcenter{\hbox{\scalebox{#2}{$\m@th#1\bullet$}}}}}
\makeatother

\newcommand\act{{\cdot}}
\newcommand\pact{{\cdot}}
\newcommand\mact{\mathbin{\bigcdot}}

\usepackage{datetime}
\yyyymmdddate

\begin{document}
\title{A declarative approach to specifying distributed algorithms using three-valued modal logic}
\newcommand\titlerunning{\emph{Declarative distributed algorithms via modal logic}} %
\newcommand\authorrunning{\emph{Murdoch J. Gabbay and Luca Zanolini}}
\author{Murdoch J. Gabbay \affil{Heriot-Watt University, Edinburgh, UK}
Luca Zanolini \affil{Ethereum Foundation, London, UK}}

\begin{abstract}
  We present \emph{\QLogic}, a three-valued modal fixed-point logic designed for declaratively specifying and reasoning about distributed algorithms, such as the Paxos consensus algorithm.

  Our methodology represents a distributed algorithm as a logical theory, enabling correctness properties to be derived directly within the framework—or revealing logical errors in the algorithm's design when they exist.
  
  \QLogic adopts a declarative approach, specifying the overall logic of computation without prescribing control flow. Notably, message-passing is not explicitly modeled, distinguishing our framework from approaches like TLA+. This abstraction emphasises the logical essence of distributed algorithms, offering a novel perspective on their specification and reasoning.
  
  We define the syntax and semantics of \QLogic, explore its theoretical properties, and demonstrate its applicability through a detailed treatment of the Paxos consensus algorithm. By presenting Paxos as a logical theory and deriving its standard correctness properties, we showcase the framework’s capacity to handle non-trivial distributed systems.
  
  We envision \QLogic as a versatile tool for specifying and reasoning about distributed algorithms. The Paxos example highlights the framework’s ability to capture intricate details, offering a new lens through which distributed algorithms can be specified, studied, and checked.


\keywords{Modal logic, distributed consensus} 
\end{abstract}


\maketitle
\thispagestyle{empty}

\tableofcontents

\section{Introduction}
\label{sect.intro}

\subsection{Initial overview}

\subsubsection*{Introducing \QLogic}
\emph{\QLogic}is a logical framework for studying distributed algorithms in a \emph{declarative} manner.
In our framework an algorithm is presented as a logical theory (a signature and a set of axioms), and then its correctness properties are derived from those axioms.
We apply our framework to present and study an axiomatisation of Paxos, a canonical distributed consensus algorithm. 

Logic \emph{per se} is widely-applied in distributed algorithms: for example the logic TLA+ (based on Linear Temporal Logic) is an industry standard to model and verify distributed protocols~\cite{DBLP:journals/toplas/Lamport94,DBLP:books/aw/Lamport2002,DBLP:conf/charme/YuML99}. 
However, TLA+ specifications describe state changes in a sequential, step-by-step manner.
A \QLogic axiomatisation feels more abstract: we dispense with step-by-step descriptions of the algorithm and with explicit message-passing, to capture an abstract but still entirely identifiable essence of the algorithm.

Abstracting an algorithm to a logical theory \emph{per se} is also not new: \emph{declarative programming} is based on this.
However, turning a distributed algorithm (like Paxos) into an axiomatic theory, is new, and this paper shows that it can be done.
Examples follow below and in the body of the text.\footnote{Or, if the reader wants to skip the preliminaries and jump to a toy example that illustrates our method, then see Subsection~\ref{subsect.simple}.}

For reasons of length, we focus on the example of Paxos.
This is a canonical example with sufficient complexity to be interesting, but we have checked that \QLogic can be adapted for other algorithms; it is intended as a general logical tool that can be tailored to particular case studies.
In the long run, we hope this will help do all the things that declarative techniques are supposed to: help find errors, decrease development times, increase productivity, and lead to new mathematics.
 
\subsubsection*{What is declarative programming}
By `declarative' we mean that we specify the logic of a computation without specifying its control flow.

Thus our axiomatisation of Paxos (in Figure~\ref{fig.logical.paxos}) relates to a concrete Paxos implementation and/or its TLA+ specification, much as this axiomatic, algebraic specification of the factorial function\footnote{By `axiomatic' we mean that the function is specified just by asserting properties that are true of it.  By `algebraic' we mean that these assertions are purely equational, i.e. we just assert some equalities between terms.}
$$
0!=1
\qquad 
(n+1)!=(n+1)*n!
$$
relates to this concrete implementation of factorial in BASH
\begin{quote}
{\tt\small\hspace{-2em} value=1; while [ \$number -gt 1 ]; do value=\$((value * number)); number=\$((number - 1)); done}
\end{quote}
The reader may be familiar with programming languages based on declarative principles, such as Haskell, Erlang, or Prolog.
Their design philosophy is that we should assert things that we want to be true, and leave it to the compiler to figure out how best to make this happen.
Contrast with C, which is a canonical imperative language.\footnote{We are not interested here in taking sides.  With a variety of tools to choose from, we pick the one that is best for the job.  What matters is that the widest possible \emph{variety} of levels of abstraction is made available.}

This paper %
explores a methodology for specifying distributed algorithms which sits decidedly at the `Haskell/Erlang/Prolog' side of the spectrum, compared to TLA+ which (while still quite high-level) sits more at the C side of the spectrum. 

Consider a tiny concrete example of a simple voting system: we will briefly exhibit this in \QLogic, just to give a flavour of how it works.
Consider a voting system such that 
\begin{quote}
`if a quorum of voters vote for candidate $c$, then $c$ wins the election'.
\end{quote}
We write this axiomatically simply as
$$
(\QuorumBox\tf{vote}(c)) \timpc \tf{wins}(c) .
$$ 
The $\QuorumBox$ is a logical modality that is true when there exists a quorum of participants at which $\tf{vote}(v)$ holds.\footnote{Details of this modality in our framework are in Proposition~\ref{prop.pointwise.dense.char}(\ref{item.dense.char.pointwise.QuorumBox}) and Lemma~\ref{lemm.allbut.concretely}.}

Even without going into the details of what these symbols mean, it is clear that details of message-passing, of how votes are communicated, and of how they are counted up, are absent.
In fact in this case the predicate is just a direct rendering of the English statement above it, without specifying how it is concretely implemented.
Thus, in keeping with a declarative philosophy, we specify what we want the algorithm to do, but not how it should do it: it is the point of the axiom to express formally and precisely our idea --- and to add as little as possible to that, beyond what is necessary to make it logically precise.


\subsubsection*{Who should read this document}
We have in mind two communities: 
\begin{enumerate*}
\item
\emph{logicians}, who might be interested in exploring a logic with an interesting combination of features (three-valued modal fixedpoint logic) with specific practical motivations in programming distributed systems, and 
\item
\emph{researchers in distributed systems}, who might find \QLogic useful for analysing and communicating their designs.
\end{enumerate*}
We will be particularly interested in consensus algorithms as applied to blockchain, because this is an important and safety-critical use case (the second author is part of the core team developing consensus algorithms for Ethereum).

\subsubsection*{References to key points in the text}

The reader who might want to jump to where the action is and read backwards and forwards from there, is welcome to proceed as follows:
\begin{itemize*}
\item
The axioms of Declarative Paxos are in Figure~\ref{fig.logical.paxos}.
\item
The syntax and semantics of \QLogic are in Figures~\ref{fig.predicate.syntax}, \ref{fig.3.phi.f}, and~\ref{fig.3.derived}.
\item
Truth-tables for three-valued logic and its connectives are in Figure~\ref{fig.3}.
\item
An outline of the Paxos algorithm (for the reader who might like to be reminded what it is) is in Remark~\ref{rmrk.high-level.paxos}.
\item
See also toy axiomatisations for an imaginary `simple' protocol in Subsection~\ref{subsect.simple}, designed to be vaguely reminiscent of Paxos, but simplified almost to triviality for the sake of exposition.
\end{itemize*}
 

\subsection{Motivation and background (further details)}

Let us unpack what consensus in distributed systems (like blockchains) is, why it matters, and why it is hard:

\subsubsection*{Distributed systems}
A distributed computing system is is one that distributes its computation over several machines.
Distributed systems are essential to computing: they enable a network of parties to collaborate to achieve common goals~\cite{cachinbook}. 

Being distributed is attractive to software engineers because it can give us desirable properties, including some or all of: scalability, redundancy, lower latency, reliability, and resilience.
For this reason, almost any nontrivial commercial software system that the reader might care to name is implemented in a distributed manner.
Furthermore, blockchain systems (like Bitcoin and Ethereum\footnote{\url{https://bitcoin.org/en/} and \url{https://ethereum.org/en/}}) and peer-to-peer systems are inherently distributed by design; if they were centralised, then they would not even make sense.

\subsubsection*{The challenge of consensus}
Among the many challenges in distributed systems, \emph{consensus} represents a fundamental abstraction.
Consensus captures the problem of 
\begin{quote}
reaching agreement among multiple participants on a common value, in the face of possibly unreliable communication, and in the presence of possibly faulty participants.
\end{quote}
Consensus algorithms are an essential prerequisite for coordination in any distributed system, and they play a particularly critical role in blockchain networks.
In such networks, consensus is the mechanism that ensures all participants agree on the state of the ledger --- the blockchain's source of truth --- thereby solving problems like \emph{double spending} (manipulating the blockchain ledger to spend the same token twice).
Consensus protocols also enable the secure ordering of transactions, and they ensure the integrity and consistency of the blockchain, even in the face of adversarial actors.

Consensus protocols in blockchain systems are particularly challenging to design and implement correctly. 
Indeed, there are numerous examples in which incomplete or insufficiently-specified protocols have led to serious issues, such as network halts or vulnerabilities that adversaries could exploit.
For example:
\begin{enumerate*}
\item
Solana (\url{https://solana.com}) has experienced instances where its consensus mechanism contributed to temporary network halts~\cite{yakovenko2018solana,shoup2022poh}.
\item 
Ripple's (\url{https://ripple.com}) consensus protocol has been analysed for conditions in which safety and liveness may not be guaranteed~\cite{DBLP:conf/opodis/Amores-SesarCM20}. 
\item
Avalanche's (\url{https://www.avax.network}) consensus protocol has faced scrutiny regarding its assumptions and performance in adversarial scenarios~\cite{DBLP:conf/opodis/Amores-SesarCT22}.
\end{enumerate*}
Even mature blockchains like Ethereum are not immune to challenges.
Attacks such as \emph{balancing attacks}~\cite{DBLP:conf/sp/NeuTT21} and \emph{ex-ante reorganisation} (\emph{reorg}) attacks~\cite{DBLP:conf/fc/Schwarz-Schilling22} have highlighted vulnerabilities in Ethereum's protocol.

These examples underscore how subtle flaws or ambiguities in protocol design can have significant consequences, emphasising the need for rigorous analysis and well-defined specifications in this complex and safety-critical field.

Compounding the fact that consensus algorithms are foundational to the security and functionality of blockchain systems, is the fact that designing and implementing these algorithms is notoriously challenging.
Distributed algorithms --- including consensus protocols --- are difficult (at least for human brains) because of the inherent difficulty of coordinating multiple participants, which may be faulty, hostile, drifting in and out of network connectivity~\cite{sleepy}, and/or have different local clocks, states, agendas, and so forth.

In short, this is a notoriously difficult field which is full of `gotchas', and within this difficult field, consensus protocols are amongst the most intricate, demanding, and practically important objects of study. 

\subsubsection*{Introducing a logic-based declarative approach (applied to Paxos)}
In this paper, we develop a logic-based declarative approach to modeling distributed algorithms.
It stands in relation to existing techniques in roughly the same way as declarative programming languages stand to imperative ones. %

Technically speaking, \QLogic is a \emph{three-valued modal fixedpoint logic}. 
This means:
\begin{itemize*}
\item
There are three truth-values: $\tvT$ (true), $\tvB$ (no value returned, i.e. `crashed'), and $\tvF$ (false)~\cite{sep-logic-manyvalued}.
\item
Modalities let us express things like `some participant has property $\phi$' or `a quorum of participants has property $\psi$' or `$\phi$ has been true in the past' or `$\psi$ will be true in the future'.\footnote{In order, these are as follows in our logic: $\someone\phi$, $\QuorumBox\phi$, $\recent\phi$, and $\sometime\phi$.  Precise semantics are in Figures~\ref{fig.3.phi.f} and~\ref{fig.3.derived}.} 
\item
A fixedpoint operator $\mu$ imports the expressive power of iteration~\cite{wiki:Fixedpoint_logic}.
Consensus algorithms are often iterative, so it is natural to include this in our logic.
\end{itemize*} 
We will apply these constructs to give axiomatic models of the Paxos consensus protocol~\cite{DBLP:journals/tocs/Lamport98,lamport2001paxos,cachin2011paxos} and then we will state and derive standard correctness properties from the axioms.
Our example is elementary but nontrivial and strikes (we think) a good balance between being compact but still exhibiting enough richness of structure to illustrate the essential features of how our logic works. 


The first half of this paper will focus on the logic, and the second half of this paper will apply the logic to the examples.
The reader from a logic background might like to start with the first half and refer to examples for motivation; the reader from a distributed systems background might like to do the opposite and start with the examples in the second half, referring back to the logic as required to understand what is being done.
 
\subsubsection*{Related work using logic}
Logic-based approaches have long been used to model and verify distributed protocols, with TLA+~\cite{DBLP:journals/toplas/Lamport94,DBLP:books/aw/Lamport2002,DBLP:conf/charme/YuML99} being one of the most prominent frameworks in the field. 
Introduced by Leslie Lamport in 1999, TLA+ uses Linear Temporal Logic (LTL)~\cite{DBLP:conf/focs/Pnueli77} to describe state changes in a sequential, step-by-step manner, which has become an industry standard for specifying and verifying distributed algorithms.

But TLA+ does things in a very particular way, which is not declarative.
We would argue that it is good for specifying \emph{algorithms}, and so is not particularly \emph{declarative}.
Specifically, TLA+ excels at specifying and verifying system behaviours through state transitions; it is operational and state-based.
This can sometimes distance TLA+ protocol specifications from the invariants that they are intended to uphold.

This separation arises because TLA+ primarily describes how states evolve \emph{over time}, which can make it challenging to directly reason about the essential properties that protocols are designed to maintain. 
So, although TLA+ provides powerful tools for specifying and checking invariants, the process of verifying that these invariants are preserved across all possible state transitions can be complex and unintuitive, particularly when dealing with the intricacies of distributed algorithms.

\subsubsection*{Our approach}
Our novel framework leverages the mathematics of semitopology and three-valued logic. 

Semitopology is from previous work~\cite{gabbay:semdca,gabbay:semtad}.
It provides a natural way to describe quorums --- sets of participants that work together to progress the system state --- by representing them as open sets.
This perspective is particularly suited to distributed systems, where the behaviour of the system emerges from the interactions of coalitions its components, rather than from a sequence of steps by a central controller.

Three-valued logic (a good overview is in \cite[Chapter~5]{bergmann:intmvf}) helps us to reason about the (possibility for) faults that is inherent in distributed systems.
Unlike classical logic with its binary truth-values $\tvT$ and $\tvF$, three-valued logic introduces a third truth value $\tvB$ that captures the indeterminate states often encountered when dealing with faulty parties.
This turns out to give particularly natural and compact descriptions of distributed protocols, distinguishing between normal operation and faulty behaviour. 
It also provides a formal foundation that would be amenable to verification through tools like SAT solvers and theorem provers (building tooling for what we see in this paper would be future work), enhancing both the expressiveness and verifiability of protocol specifications.

Our approach is particularly relevant in the context of blockchain consensus protocols, which are distributed and decentralised systems that require robust and clear specifications to function correctly.
It is not an exaggeration to say that blockchains in particular have put a rocket underneath the development of new consensus protocols~\cite{DBLP:conf/wdag/CachinV17}.
Many of these are not fully specified or formally verified. 
Often, they are even described informally in blog posts and other non-academic sources --- leading to confusion about their properties and behaviours. 
This lack of formalisation creates a fragmented landscape, where the guarantees and limitations of the protocols remain unclear, hindering their adoption and undermining trustworthiness.

We hope that by providing a logic-based declarative language, and the ability to axiomatise consensus protocols, our framework will offer a way to formally specify, reason about, and verify the properties of the many distributed systems under development --- thus lowering costs, increasing reliability, and reducing development time. 
We would further argue that effective formalisation is essential in reducing ambiguity in the design and evaluation of the many new consensus mechanisms being developed in the blockchain space in particular.

Our approach bridges the gap between the declarative specification of invariants and the practical verification of protocol correctness. 

\subsubsection*{We will explain everything, and that's a feature}

Different readers will find some parts of this paper easy, and other parts less easy --- but which parts these are depends on the reader and their background.\footnote{Not a hypothetical: the first author (with a logic background) took a long time to get his head around how Paxos works, whereas the second author (who knows Paxos well and is familiar with logic in general terms, but is not a logician) was not immediately familiar with modal fixedpoint logic and models of axiomatic logical theories.  The community of readers who approach this paper and are already well-versed in both fields may not be the largest.  We hope this paper will help ameliorate this lamentable state of affairs.}

For example:
\begin{itemize*}
\item
three-valued modal logic and its models, and technical results like the Knaster-Taski fixedpoint theorem used in Proposition~\ref{prop.mu.fixedpoint} may be known to the reader with a background in logic (like the first author), but perhaps not to a distributed systems theorist (like the second author); conversely, 
\item
the Paxos algorithm, and standard technical devices like the discussion motivating $n$-twined semitopologies in Remark~\ref{rmrk.motivate.n-twined}, are the bread-and-butter of a distributed systems expert's work --- but this application might be new to someone from a logic background. 
\end{itemize*}
We play safe and explain everything. 
If the reader finds we are going a bit slow in places, then better that than going too fast and skipping details that might help some other reader follow the argument.
 


\subsection{Structure of the paper}

\begin{enumerate*}
\item
Section~\ref{sect.intro} is the Introduction; You are Here.
The rest of the paper splits into two halves: Sections~\ref{sect.three.truth.values} to Section~\ref{sect.n-twined} have to do with logic and semitopologies; and Sections~\ref{sect.declarative.paxos} to~\ref{sect.simpler.paxos} use \QLogic to declaratively specify and study Paxos.
\ \\
\item
Section~\ref{sect.three.truth.values} introduces the basics of three-valued logic, including notable connectives ($\tnotor$ and $\timpc$) and modalities ($\modT$ and $\modTB$).
We also introduce a notion of validity for truth-values.
\item
Section~\ref{sect.quorum.logic} defines semitopologies (our model of the quorums that are ubiquitous in consensus algorithms), then it builds the syntax and semantics of \QLogic.
\item
Section~\ref{sect.predicates} studies \QLogic predicates; notably developing some logical equivalences, and well-behavedness criteria like being \emph{pointwise}.
\item 
Section~\ref{sect.n-twined} returns to semitopologies, and shows how antiseparation properties can be represented in logical form in \QLogic.
One particular antiseparation property corresponds to the \emph{quorum intersection property} which is used in the correctness proofs for Paxos (see Remark~\ref{rmrk.ntwined.and.the.literature}), but we will put this fact in a more general logical/semitopological context.
\ \\
\item
Section~\ref{sect.declarative.paxos} introduces Declarative Paxos; this is representative for the kind of algorithm that we built \QLogic to express.
We take our time to explain how Paxos works, state the axiomatisation, and discuss it in detail.
\item
Section~\ref{sect.paxos.correctness.properties} states the standard correctness properties for Paxos, phrases them in declarative form, and then proves them from the axioms of Declarative Paxos.
What all our work has accomplished is to convert proving correctness properties into a (relatively straightforward) argument of proving facts of an abstract model from its axioms.
\item
Section~\ref{sect.simpler.paxos} revisits the axioms of Declarative Paxos from Section~\ref{sect.declarative.paxos}, in the light of the proofs in Section~\ref{sect.paxos.correctness.properties} --- and then prunes them down to a minimal theory required to prove correctness.

We call this minimal theory \emph{Simpler Declarative Paxos}: it is simpler, and because further details have been elided, it is more abstract.
The reader may prefer Declarative Paxos to Simpler Declarative Paxos because it is more identifiably related to the Paxos algorithm, but if we want to prove correctness properties, then Simpler Declarative Paxos distils the logical essence of what makes those proofs work. 
\ \\
\item
Section~\ref{sect.conclusions} concludes with conclusions and reflections on future work.
\end{enumerate*} 


\section{Three truth-values}
\label{sect.three.truth.values}

\subsection{Basic definition of $\THREE$, and some truth-tables}

\begin{defn}[Three-valued truth-values]
\label{defn.THREE}
Let $\THREE$ be the poset of \deffont{truth-values} with elements $\tvF$ (false), $\tvB$ (both), and $\tvT$ (true) ordered as 
$$
\tvF<\tvB<\tvT.
$$
\end{defn}

\begin{figure} %
$$
\begin{array}{c}
\begin{array}{c@{\qquad}c@{\qquad}c@{\qquad}c@{\qquad}c@{\qquad}c}
\begin{array}{c c c c}
p\tand q & \tvT & \tvB & \tvF
\\
\tvT &  \tvT & \tvB & \tvF
\\
\tvB &  \tvB & \tvB & \tvF 
\\
\tvF &  \tvF & \tvF & \tvF 
\end{array}
&
\begin{array}{c c c c}
p\tor q & \tvT & \tvB & \tvF
\\
\tvT &  \tvT & \tvT & \tvT
\\
\tvB &  \tvT & \tvB & \tvB
\\
\tvF &  \tvT & \tvB & \tvF
\end{array}
&
\begin{array}{c c c c}
p\tnotor q & \tvT & \tvB & \tvF
\\
\tvT &  \tvT & \tvB & \tvF
\\
\tvB &  \tvT & \tvB & \tvB
\\
\tvF &  \tvT & \tvT & \tvT
\end{array}
&
\begin{array}{c c c c}
p \timpc q & \tvT & \tvB & \tvF
\\
\tvT       & \tvT & \tvF & \tvF
\\
\tvB       & \tvT & \tvB & \tvB
\\
\tvF       & \tvT & \tvT & \tvT
\end{array}
\end{array}
\\[7ex]
\begin{array}{c@{\qquad}c@{\qquad}c@{\qquad}c@{\qquad}c@{\qquad}c@{\qquad}c}
\begin{array}{cc}
\tneg p
\\
\tvT & \tvF
\\
\tvB & \tvB
\\
\tvF & \tvT
\end{array}
&
\begin{array}{cc}
\modT p 
\\
\tvT & \tvT
\\
\tvB & \tvF
\\
\tvF & \tvF
\end{array}
&
\begin{array}{cc}
\modF p 
\\
\tvT & \tvF
\\
\tvB & \tvF
\\
\tvF & \tvT
\end{array}
&
\begin{array}{cc}
\modB p
\\
\tvT & \tvF
\\
\tvB & \tvT
\\
\tvF & \tvF
\end{array}
&
\begin{array}{cc}
\modTB p  
\\
\tvT & \tvT
\\
\tvB & \tvT
\\
\tvF & \tvF
\end{array}
&
\begin{array}{cc}
\modTF p  
\\
\tvT & \tvT
\\
\tvB & \tvF
\\
\tvF & \tvT
\end{array}
&
\begin{array}{cc}
\modFB p  
\\
\tvT & \tvF
\\
\tvB & \tvT
\\
\tvF & \tvT
\end{array}
\end{array}
\\[7ex]
\begin{array}{c@{\qquad}c@{\qquad}c@{\qquad}c@{\qquad}c@{\qquad}c}
\begin{array}{c c c c}
p\timp q & \tvT & \tvB & \tvF
\\
\tvT &  \tvT & \tvB & \tvF
\\
\tvB &  \tvT & \tvB & \tvF 
\\
\tvF &  \tvT & \tvT & \tvT 
\end{array}
&
\begin{array}{c c c c}
p\tleq q & \tvT & \tvB & \tvF
\\
\tvT &  \tvT & \tvF & \tvF
\\
\tvB &  \tvT & \tvB & \tvF 
\\
\tvF &  \tvT & \tvT & \tvT 
\end{array}
\end{array}
\end{array}
$$
\emph{Above, the vertical axis of a table indicates values for $p$; the horizontal axis (if nontrivial) denotes values for $q$.}
\caption{Truth-tables for some operations on $\THREE$ (Remark~\ref{rmrk.three.elementary})}
\label{fig.3}
\end{figure}


\begin{figure} %
$$
\begin{array}{c@{\qquad}c@{\qquad}c@{\qquad}c@{\qquad}c@{\qquad}c}
\begin{array}{c c c c}
p\tlatticeiff q & \tvT & \tvB & \tvF
\\
\tvT &  \tvT & \tvB & \tvF
\\
\tvB &  \tvB & \tvB & \tvB
\\
\tvF &  \tvF & \tvB & \tvT
\end{array}
&
\begin{array}{c c c c}
p \tiffcc q & \tvT & \tvB & \tvF
\\
\tvT       & \tvT & \tvF & \tvF
\\
\tvB       & \tvF & \tvB & \tvB
\\
\tvF       & \tvF & \tvB & \tvT
\end{array}
&
\begin{array}{c c c c}
p \equiv q & \tvT & \tvB & \tvF
\\
\tvT       & \tvT & \tvF & \tvF
\\
\tvB       & \tvF & \tvT & \tvF
\\
\tvF       & \tvF & \tvF & \tvT
\end{array}
\end{array}
$$
\caption{Truth-tables for notions of equivalence (Definition~\ref{defn.equivalent})}
\label{fig.3.iff}
\end{figure}


\begin{rmrk}
\label{rmrk.three.elementary}
We start with some elementary observations about $\THREE$.
\begin{enumerate*}
\item\label{item.three.elementary.bounded.lattice}
$\THREE$ %
has bottom element $\tvF$, top element $\tvT$, and meets $\tand$ and joins $\tor$.
This makes it a \emph{finite lattice}.
\item
$\THREE$ has a negation operation $\tneg$.
\item
$\THREE$ has several natural modalities that identify truth-values, including $\modT$, $\modF$, $\modB$, $\modTB$, and $\modTF$.

For example, the reader can examine the truth-table for $\modTF$ in Figure~\ref{fig.3} and see that it identifies elements of $\{\tvT,\tvF\}$.\footnote{One could also call these modalities \emph{characteristic functions} for various subsets of $\THREE$, by which we mean functions which (literally) return $\tvT$ on elements in the set and $\tvF$ on elements not in the set.}
\item
$\THREE$ also has multiple notions of implication and equivalence, some of which are illustrated in Figures~\ref{fig.3} and~\ref{fig.3.iff}.
We will discuss these implications in Subsection~\ref{subsect.3.imp}.
\end{enumerate*}
\end{rmrk}


\begin{rmrk}
We will use the third truth-value $\tvB$ to represent participants in a distributed algorithm \emph{crashing}.
The method is simple: the truth-value of a predicate describing the outcome of a crashed computation, is deemed to be $\tvB$.

This is similar in spirit to how in domain theory an additional value --- usually written $\bot$ --- is added to a denotation to represent nontermination; the function-value of a nonterminating computation is deemed to be $\bot$~\cite{abramsky:domt}.
The nonterminating computation does not return a value, but we represent this absence of a value in the denotation, using $\bot$. 

In our declarative axiomatisation of Paxos, we use $\tvB$ in much the same way, as a truth-value to describe the \emph{absence} of a return value (either $\tvT$ or $\tvF$) from participants that have crashed.
\end{rmrk}



\subsection{Valid and correct truth-values}



\begin{defn}
\label{defn.tv.ment}
\leavevmode
\begin{enumerate*}
\item\label{item.tv.ment.correct}
Define subsets $\threeValid,\threeCorrect,\threeTrue,\threeFalse\subseteq\THREE$ by
$$
\threeValid = \{\tvT,\tvB\}
\quad\text{and}\quad
\threeCorrect = \{\tvT,\tvF\} 
\quad\text{and}\quad
\threeTrue = \{\tvT\} 
\quad\text{and}\quad
\threeFalse = \{\tvF\} .
$$ 
\item
We refer to $\f{tv}\in\THREE$ as being \deffont{valid} / \deffont{invalid},\ \deffont{correct} / \deffont{incorrect},\ and \deffont{true} / \deffont{untrue},\ and \deffont{false} / \deffont{unfalse}\footnote{\dots `unfalse' is `$\tvT$-or-$\tvB$', not `$\tvT$'.  Similarly, `untrue' is $\tvB$-or-$\tvF$.  We cannot resist stating the following gem of three-valued clarity: not-$\tvT$, is not $\tneg\tvT$, and not-$\tvF$ is not $\tneg\tvF$, and not-$\tvB$ \emph{is} not-$\tneg\tvB$, because $\tvB$ is $\tneg\tvB$.} depending on whether it is / is not in $\threeValid$, $\threeCorrect$, $\threeTrue$, and $\threeFalse$ respectively.
\item\label{item.tv.ment}
Define $\ment\f{tv}$ \ (read `\deffont{$\f{tv}$ is valid}') by
$$
\ment\f{tv} 
\quad\text{when}\quad
\f{tv}\in\threeValid ,
\ \ \text{which just means}\ \ 
\f{tv}\in\{\tvT,\tvB\} .
$$
Thus, `$\f{tv}$ is valid' holds precisely when $\f{tv}$ is a valid truth-value.
\end{enumerate*}
\end{defn}

\begin{rmrk}
\label{rmrk.motivate.validity}
We take a moment to discuss Definition~\ref{defn.tv.ment}:
\begin{enumerate}
\item
We can think of $\tvT$ and $\tvF$ as specific `yes/no' results from the successful completion of some computation.
Therefore, we will call $\tvT$ and $\tvF$ \emph{correct}.
\item
We can think of $\tvB$ as representing failure of a computation to complete correctly --- e.g. because a participant has crashed and so can return neither $\tvT$ nor $\tvF$.
\item\label{item.tb.valid}
In Definition~\ref{defn.tv.ment} we set $\threeValid=\{\tvT,\tvB\}$ and define $\ment\f{tv}$ when $\f{tv}\in\threeValid$, and we call $\tvT$ and $\tvB$ \emph{valid}.
\end{enumerate}
It might seem strange at first to let $\tvB$ be valid; if $\tvB$ represents the (non)behaviour of a crashed participant, then surely crashing should \emph{not} be valid?

Not necessarily.
To see why, consider that many of assertions of interest that we will consider will have the form 
\begin{quote}
`if $p$ is not crashed, then $\ment\phi$ at $p$', 
\end{quote}
or equivalently 
\begin{quote}
`either $p$ is crashed, or $\ment\phi$ at $p$'.
\end{quote}
If we let $\tvB$ be valid (i.e. if we define $\ment\tvB$ to hold), then we can simplify the above to the conveniently succinct assertion 
\begin{quote}
`$\ment\phi$ at $p$'.
\end{quote}
The `if $p$ is not crashed' precondition gets swallowed up into the fact that if $p$ is crashed then $\phi$ returns $\tvB$, and $\ment\tvB$ so there is nothing to prove. 
\end{rmrk}


\begin{rmrk}
Later on in Definition~\ref{defn.validity.judgement}(\ref{item.ctx.ment.phi}) we will extend validity to predicates-in-context, just by evaluating a predicate-in-context to a truth-value using a denotation function, and checking that truth-value is valid in the sense we define here. 
\end{rmrk}


We make some elementary observations about $\THREE$ and Definition~\ref{defn.tv.ment}:
\begin{lemm}
\label{lemm.tand.tor}
Suppose $\f{tv}\in\THREE$ is a truth-value.
Then:
\begin{enumerate*}
\item
$\ment\tvT$ and $\ment\tvB$.
\item
$\nment\tvF$.
\item\label{item.tand.tor.tor}
$\ment\f{tv}\tor\f{tv}'$ if and only if ${\ment\f{tv}}\lor{\ment\f{tv}'}$.
\item
$\ment\f{tv}\tand\f{tv}'$ if and only if ${\ment\f{tv}}\land{\ment\f{tv}'}$.
\end{enumerate*}
\end{lemm}
\begin{proof}
Fact of Definition~\ref{defn.tv.ment} and the truth-tables in Figure~\ref{fig.3}.
\end{proof} 

Negation is conspicuous by its absence in Lemma~\ref{lemm.tand.tor}.
We consider this next:
\begin{lemm}
\label{lemm.para}
Suppose $\f{tv}\in\THREE$ is a truth-value.
Then:
\begin{enumerate*}
\item\label{item.para.em}
$\ment\f{tv}\tor\tneg\f{tv}$, and as a corollary $\nment\f{tv}$ implies $\ment\tneg\f{tv}$.
\item\label{item.para.para}
$\ment\f{tv}\tand\tneg\f{tv}$ is possible, and as corollaries: $\ment\f{tv}$ does not necessarily imply $\nment\tneg\f{tv}$; and it is possible that $\ment\f{tv}$ and $\ment\tneg\f{tv}$ both hold.
\end{enumerate*}
\end{lemm}
\begin{proof}
We consider each part in turn:
\begin{enumerate}
\item
Using Lemma~\ref{lemm.tand.tor}(3) it would suffice to show that $\ment\f{tv}$ or $\ment\tneg\f{tv}$.
So suppose $\nment\f{tv}$.
By Definition~\ref{defn.tv.ment}(\ref{item.tv.ment}) $\f{tv}\not\in\threeValid$, thus $\f{tv}=\tvF$ and by Figure~\ref{fig.3} $\tneg\f{tv}=\tvT\in\threeValid$ and so $\ment\tneg\f{tv}$.
\item
We just take $\f{tv}=\tvB$ and check from Figure~\ref{fig.3} that $\tneg\tvB=\tvB$.
We use Lemma~\ref{lemm.tand.tor}(1\&4).
\qedhere\end{enumerate}
\end{proof}

\begin{rmrk}
Some comments on Lemma~\ref{lemm.para}:
\begin{itemize}
\item
Lemma~\ref{lemm.para}(\ref{item.para.para}) notes a typical fact of validity in three-valued logic, that it is naturally \deffont{paraconsistent}, meaning that it is possible for a thing, and the negation of that thing, to both be valid.
This is a feature of $\tvB\in\THREE$, not a bug, and we will make good use of it later when we use $\tvB$ %
as a return value for `I return no correct truth-value'. 
\item
Lemma~\ref{lemm.para}(\ref{item.para.em}) asserts validity of the \emph{excluded middle} (`$\phi$-or-negation-of-$\phi$'), but its corollaries implicitly illustrate another fact of three-valued logic, by what they do \emph{not} say: 
looking at the proof of Lemma~\ref{lemm.para}(\ref{item.para.em}), it is clearly also true that $\ment\f{tv}\tor\modF\f{tv}$, and that $\ment\f{tv}\tor\modT\tneg\f{tv}$, and that $\nment\f{tv}$ implies $\ment\modF\f{tv}$, and also that $\nment\f{tv}$ implies $\ment\modT\tneg\f{tv}$.  

These are stronger assertions, so why did we not state them in the lemma instead?
Because we will not need the result in those forms later.  

In two-valued logic, if something is true then it is probably a useful lemma.
In contrast, in three-valued logic, there is so much more structure with the third truth-value that more things are \emph{true} than are necessarily \emph{useful}.
So, we have to pick and choose our lemmas, and it is not always the case that the strongest form of a result is the most useful.
\end{itemize}
\end{rmrk}


Lemmas~\ref{lemm.tv.ment.TF} and~\ref{lemm.TF.correct} are easy, but they spell out some useful facts about truth-values and validity judgements:
\begin{lemm}
\label{lemm.tv.ment.TF}
Suppose $\f{tv}\in\THREE$ is a truth-value.
Then:
\begin{enumerate*}
\item\label{item.tv.ment.TF.1}
$\threeCorrect\cap\threeValid=\threeTrue$.

In words: a truth-value is true if and only if it is correct and valid.
\item\label{item.tv.ment.TF.2b}
$\ment \modT\f{tv}$ if and only if $\f{tv}\in\threeTrue$. 

In words: `$\modT\f{tv}$ is valid' holds if and only if $\f{tv}$ is true.
\item\label{item.tv.ment.TF.2b.F}
$\ment \modF\f{tv}$ if and only if $\f{tv}\in\threeFalse$. 

In words: `$\modT\f{tv}$ is valid' holds if and only if $\f{tv}$ is true.
\item\label{item.tv.ment.TF.TB}
$\ment \modTB\f{tv}$ if and only if $\f{tv}\in\threeValid$ if and only if $\ment\f{tv}$. 

In words: `$\modTB\f{tv}$ is valid' if and only if $\f{tv}$ is valid.\footnote{Note that $\modTB\f{tv}$ is not necessarily \emph{literally equal} to $\f{tv}$, in the sense of $\modT\f{tv}\equiv\f{tv}$ from Definition~\ref{defn.equivalent}.  For example, $\modTB\tvB=\tvT\neq\tvB$.}
\item\label{item.tv.ment.TF.2}
$\ment \modTF\f{tv}$ if and only if $\f{tv}\in\threeCorrect$. 

In words: `$\modTF\f{tv}$ is valid' if and only if $\f{tv}$ is correct.
\end{enumerate*}
\end{lemm}
\begin{proof}
Part~\ref{item.tv.ment.TF.1} is a fact of sets.
Parts~\ref{item.tv.ment.TF.2b} and~\ref{item.tv.ment.TF.TB} and~\ref{item.tv.ment.TF.2} are just from the truth-table for $\modT$, $\modTB$, and $\modTF$ in Figure~\ref{fig.3} and Definition~\ref{defn.tv.ment}. 
\end{proof}


\begin{lemm}
\label{lemm.TF.correct}
Suppose $\f{tv}\in\THREE$ is a truth-value.
Then the following are all equivalent:
$$
\ment\modTF\f{tv}
\ \liff\ 
\modTB\f{tv}=\modT\f{tv}
\ \liff\ 
\modT\f{tv}=\f{tv} 
\ \liff\ 
\modFB\f{tv}=\modF\f{tv}
\ \liff\ 
\modF\f{tv}=\tneg\f{tv} 
$$ 
Informally --- we do not build the machinery to make this formal, but it may still be a helpful observation --- we can write that the modality $\modTF$ is equivalent to the modal equivalences $\modTB=\modT=\f{id}$ and $\modFB=\modF=\tneg$.
\end{lemm}
\begin{proof}
By Lemma~\ref{lemm.tv.ment.TF}(\ref{item.tv.ment.TF.2}) and Definition~\ref{defn.tv.ment}(\ref{item.tv.ment.correct}),
$\ment\modTF\f{tv}$ when 
$$
\f{tv}=\tvT \ \lor\ \f{tv}=\tvF
\quad\text{and when}\quad
\f{tv}\neq\tvB .
$$
We check the truth-tables in Figure~\ref{fig.3}, and see that all of the other statements above precisely characterise the same condition.
\end{proof}


\subsection{The rich implication and equivalence structure of $\THREE$}
\label{subsect.3.imp}

\subsubsection{Weak and strong implications, and others}


Three-valued logic supports sixteen different kinds of implication~\cite[note~5, page~22]{arieli:idepl}.
We discuss two of them in Definition~\ref{defn.three.imp}, and mention two more in Remark~\ref{rmrk.rich.implication}:
\begin{defn}
\label{defn.three.imp}
\leavevmode
\begin{enumerate*}
\item\label{item.three.elementary.weak.implication}
Define \deffont{weak implication} $\tnotor$, using $\tneg$ and $\tor$, as follows:
$$
\f{tv}\tnotor\f{tv}'=(\tneg\f{tv})\tor\f{tv}'.
$$
\item\label{item.three.elementary.strong.implication}
Define \deffont{strong implication} $\timpc$, using $\tnotor$ and $\modT$, as follows:
$$
\f{tv}\timpc\f{tv}'=\f{tv}\tnotor \modT\f{tv}'. 
$$
\end{enumerate*}
These and other encodings are assembled in Figure~\ref{fig.3.derived}.
\end{defn}

We motivate the weak and strong implications from Definition~\ref{defn.three.imp}(\ref{item.three.elementary.weak.implication}\&\ref{item.three.elementary.strong.implication}):
\begin{rmrk}[How to interpret $\tnotor$ and $\timpc$]
\label{rmrk.interpret.implication}
\leavevmode
For this Remark, let us interpret $\f{tv}\in\THREE$ as the result of a computation process which either succeeds and returns true (truth-value $\tvT$), or it succeeds and returns false (truth-value $\tvF$), or the process does not succeed (e.g. it crashes) and it returns an invalid value / no value (represented by the truth-value $\tvB$).
Recall also from Definition~\ref{defn.tv.ment}(\ref{item.tv.ment}) that $\ment\f{tv}$ means $\f{tv}\in\{\tvT,\tvB\}$.
Then by this reading:
\begin{enumerate*}
\item
$\ment\f{tv}$ means \emph{``the process $\f{tv}$ returns $\tvT$, if valid''}. 
\item
$\ment\f{tv}\tnotor\f{tv}'$ means \emph{``\emph{if} process~$\f{tv}$ returns true \emph{then} process~$\f{tv}'$ returns $\tvT$, if valid''}.

Implication does not imply sequentiality: we are just making a logical assertion about outcomes. 
\item
$\ment\f{tv}\timpc\f{tv}'$ means \emph{``\emph{if} process~$\f{tv}$ returns true \emph{then} process~$\f{tv}'$ is \emph{also} valid, and it returns $\tvT$''}.
\end{enumerate*}
\end{rmrk}


We can sum this up as a \emph{modus ponens} result, which relates the validity $\ment$ from Definition~\ref{defn.tv.ment}, the weak implication $\tnotor$ and strong implication $\timpc$ from Figure~\ref{fig.3}, to real-world implication:
\begin{prop}[Weak and strong modus ponens]
\label{prop.mp.for.tnotor}
Suppose $\f{tv},\f{tv}'\in\THREE$ are truth-values.
Then:
\begin{enumerate*}
\item\label{item.mp.for.tnotor}
$\ment\f{tv}\tnotor\f{tv}'$ if and only if\ \ $\ment\modT\f{tv}$ implies $\ment\f{tv}'$. 

In words: `$\f{tv}$ weakly implies $\f{tv}'$' is valid, if and only if $\f{tv}$ is true implies $\f{tv}'$ is valid (but $\f{tv}'$ need not be true). 

We may call this logical equivalence %
\deffont{weak modus ponens}.
\item\label{item.mp.for.timpc}
$\ment\f{tv}\timpc \f{tv}'$ if and only if $\ment\modT\f{tv}$ implies $\ment\modT\f{tv}'$. 

In words: `$\f{tv}$ strongly implies $\f{tv}'$' is valid if and only if $\f{tv}$ is true implies $\f{tv}'$ is true. 

We may call this logical equivalence \deffont{strong modus ponens}. 
\end{enumerate*}
\end{prop}
\begin{proof}
We consider each part in turn:
\begin{enumerate}
\item
By direct inspection of the truth-table for $\tnotor$ in Figure~\ref{fig.3}, using Definition~\ref{defn.tv.ment}(\ref{item.tv.ment}).
\item
By a similar direct inspection.
\qedhere\end{enumerate}
\end{proof}


\begin{rmrk}[Explicit and implicit implications]
\label{rmrk.rich.implication}
So far, the structure we have seen (the lattice structure, the negation, the modalities, and the two implications $\tnotor$ and $\timpc$) is reflected explicitly in the syntax of our logic in Figures~\ref{fig.predicate.syntax} and~\ref{fig.3.derived}.
For instance, Figure~\ref{fig.predicate.syntax} has a unary negation connective $\tneg$ whose denotation as per Figure~\ref{fig.3.phi.f} is precisely the negation operator $\tneg$ in Figure~\ref{fig.3}.
So far, so explicit.

However, $\THREE$ has further implication structure, which we do not mention explicitly in our logic, but which is still there, is important, and which appears implicitly in the proofs.
Notably: 
\begin{enumerate*}
\item\label{item.rich.implication.timp}
An implication $\timp$ defined by $\f{tv}\timp \f{tv}'=\tneg\f{tv}'\timpc\tneg\f{tv}=(\modTB\f{tv})\tnotor\f{tv}'$ has a truth-table spelled out in Figure~\ref{fig.3} (bottom left).
This can be read as ``if $\f{tv}$ is valid then $\f{tv}'$ is valid''.

$\timp$ is clearly implicit in results like Proposition~\ref{prop.someone.implies.someoneAll} -- see all the parts whose statement uses logical implication $\limp$ --- and in Proposition~\ref{prop.tv.ment.TF.model}.
\item
Another implication is implicit in the use of $\leq$ in results like Lemma~\ref{lemm.positive.monotonicity} and Corollary~\ref{corr.recent.everyone.commute}.
That symbol refers to the lattice order on $\THREE$, but this could easily be internalised in our logic --- e.g. we can set $\f{tv}\tleq \f{tv'}$ to be $(\f{tv}\timpc\f{tv}')\tand(\f{tv}\timp\f{tv}')$, to obtain the connective $\tleq$ also illustrated in Figure~\ref{fig.3} (bottom right).
\end{enumerate*}
So in fact, we use all four of $\timp$ and $\tleq$, and $\tnotor$ and $\timpc$.

The reason that $\tnotor$ and $\timpc$ make it into our logic as (sugared) syntax, and $\timp$ and $\leq$ do not, is that $\tnotor$ and $\timpc$ naturally arise in our axioms (see most of the axioms in Figures~\ref{fig.logical.paxos}) %
whereas we do not need $\timp$ or $\tleq$ in the same way for the axiomatisations in this paper.
Different axiomatisations of different algorithms might make different demands on our logic, and it would be straightforward to include (sugar for) connectives accordingly. 
\end{rmrk}

\subsubsection{Equivalence $\equiv$}

Our three-valued domain of truth-values $\THREE$ has rich notions of equivalence, to match its rich notions of implication.
We mention three of them:
\begin{defn}[Notions of equivalence in $\THREE$]
\label{defn.equivalent}
Define $\phi\tlatticeiff\phi'$ and $\phi\tiffcc\phi'$ and $\phi\equiv\phi'$ by:
$$
\begin{array}{r@{\ }l}
\phi\tlatticeiff\phi' =& (\phi\tnotor\phi')\tand(\phi'\tnotor\phi)
\\
\phi\tiffcc\phi' =& (\phi\timpc\phi')\tand(\phi'\timpc\phi)
\\
\phi\equiv\phi' =& (\modT\phi\tand\modT\phi')\tor(\modB\phi\tand\modB\phi')\tor(\modF\phi\tand\modF\phi')
\end{array}
$$
Truth-tables for $\tlatticeiff$, $\tiffcc$, and $\equiv$ are given in Figure~\ref{fig.3.iff}, and we unpack this further in Remark~\ref{rmrk.equiv}.
\end{defn}


\begin{rmrk}
\label{rmrk.equiv}
We can read 
\begin{itemize*}
\item
$\f{tv}\tlatticeiff \f{tv}'$ as `$\f{tv}$ is valid if and only if $\f{tv}'$ is valid', 
and 
\item
$\f{tv}\tiffcc \f{tv}'$ as `$\f{tv}$ is true if and only if $\f{tv}'$ is true', 
and 
\item
$\f{tv}\equiv \f{tv}'$ as `$\f{tv}$ is literally equal to $\f{tv}'$'.
\end{itemize*}
Of these three notions of equivalence, $\equiv$ will be most useful to us later.
We will consider it in our predicate language in Subsection~\ref{subsect.equivalence.of.predicates}, and (for instance) it will help us to express notions of extensional equivalence in Figure~\ref{fig.easy.equivalences}.

We will call $\equiv$ \deffont{literal equivalence}, and $\equiv$ appears in the sugared syntax of Figure~\ref{fig.3.derived}.
As per the discussion in Remark~\ref{rmrk.rich.implication}, there would be no barrier to including $\tlatticeiff$ or $\tiffcc$ but we do not need to for now.
\end{rmrk}








\section{\QLogic}
\label{sect.quorum.logic}

\subsection{A preliminary: semitopologies}

\begin{rmrk}
A semitopology is like a topology, but generalised to not require that intersections of open sets necessarily be opens.

The intuition of an open set is an \emph{actionable coalition}; a set of participants with the voting power (or other resources) to make updates to their local state.
The reader might be familiar with the notion of \emph{quorum} from literature like~\cite{DBLP:journals/tocs/Lamport98,DBLP:journals/dc/MalkhiR98,naor:loacaq}, or the notion of \emph{winning coalition} in social choice theory (the theory of voting)~\cite[Item~5, page~40]{riker:thepc}.
Semitopologies are similar to this, except that open sets are always closed under arbitrary (possibly empty) unions, whereas 
\begin{itemize*}
\item
quorums may be closed under arbitrary nonempty unions (in which case the correspondence is precise, except that $\varnothing$ is an open set but is not a quorum), or 
\item
quorums are not closed under arbitrary unions (in which case a quorum corresponds to a minimal nonempty open set).
\end{itemize*}
Our presentation of semitopologies will be quite brief --- we just define the parts we need to give semantics to the modalities in \QLogic --- and the reader who wants to know more is referred to previous work~\cite{gabbay:semdca}. 
\end{rmrk} 

\begin{defn}
\label{defn.semitopology}
A \deffont{semitopology} is a pair $(\Pnt,\opens)$ where:
\begin{enumerate*}
\item
$\Pnt$ is a set of \deffont{points}, which following terminology from distributed systems one might also call \deffont{participants}, \deffont{processes}, \deffont{parties}, or \deffont{nodes}.
\item
$\opens\subseteq\powerset(\Pnt)$ is a set of \deffont{open sets} such that:
\begin{enumerate*}
\item
If $\mathcal O\subseteq\powerset(\Pnt)$ then $\bigcup\mathcal O\in\Pnt$.
In words: $\opens$ is closed under arbitrary (possibly empty) unions.
\item
$\Pnt\in\opens$.
\end{enumerate*} 
\end{enumerate*} 
We may write $\opensne$ for the set of nonempty open sets, so
$$
\opensne = \{O\in\opens \mid O\neq\varnothing\} .
$$
\end{defn}


\begin{nttn}
\label{nttn.empty.semitopology}
Suppose $(\Pnt,\opens)$ is a semitopology.
\begin{enumerate*}
\item
We call $(\Pnt,\opens)$ the \deffont{empty semitopology} when its set of points is empty --- i.e. when $\Pnt=\varnothing$ --- so that necessarily $\opens=\{\varnothing\}$.
\item
We call $(\Pnt,\opens)$ a \deffont{nonempty semitopology} when its set of points is nonempty --- i.e. when $\Pnt\neq\varnothing$.
\end{enumerate*}
\end{nttn}


\begin{defn}[All-but-$f$-many semitopology]
\label{defn.allbut}
\leavevmode
\begin{enumerate*}
\item
Suppose $X$ is a finite set and $N\in\Ngeqz$.
Define $\powerset_{\geq N}(X)$ to be the set of subsets of $X$ that contain at least $N$ elements.
In symbols:
$$
\powerset_{\geq N}(X) = \{O\subseteq X \mid \mycard O\geq N\} .
$$
\item\label{item.allbut}
Suppose $N,f\in\Ngeqz$ and $N\geq f$.
Define a semitopology $\AllBut{N}{f}$ by
$$
\AllBut{N}{f} = \bigl(\Nset{N}, \{\varnothing\}\cup\powerset_{\geq N\minus f}(\Nset{N}\bigr)\bigr).
$$
Thus: points are numbers $0\leq n<N$; and open sets are sets of points of cardinality at least $N\minus f$.
\end{enumerate*}
\end{defn}

\begin{lemm}
Suppose $N,f\in\Ngeqz$ and $N\geq f$.
Then:
\begin{enumerate*}
\item
$\AllBut{N}{f}$ is indeed a semitopology.
\item
$\AllBut{N}{f}$ is a topology precisely when $f=N$ or $f=0$.
\end{enumerate*}
\end{lemm}
\begin{proof}
We consider each part in turn:
\begin{enumerate}
\item
It is routine to check that the conditions on open sets in Definition~\ref{defn.semitopology} are satisfied.
The key observation is just the elementary fact that if $\mycard O\geq N\minus f$ and $\mycard O'\geq N\minus f$ then $\mycard(O\cup O')\geq N\minus f$.
\item
For $\AllBut{N}{f}$ to be a topology, open sets need to be closed under intersections.
It is routine to check that this can happen only when $\AllBut{N}{f}$ is the full powerset (so $f=N$) or when open sets are just $\varnothing$ and $\Pnt$ (so $f=0$).
\qedhere\end{enumerate}
\end{proof}

\begin{rmrk}
The original Paxos paper by Lamport~\cite{DBLP:journals/tocs/Lamport98} used semitopologies of the form $\AllBut{N}{f}$ where $N=2n$ and $f=n\minus 1$, for $n\geq 1$ (i.e. open sets are strict majorities of points).
This invites a question: why bother with semitopologies at all?
\begin{enumerate*}
\item
Because we can. 
More specifically, our axiomatisation is parametric over the underlying semitopology.
As good mathematicians, if we can generalise at no cost, then we should.\footnote{\dots and this applies not just in mathematics papers.  Software engineers follow a similar principle and have many names for it: e.g. encapsulation, APIs, modules, classes, typeclasses, and so on.} 
\item
Semitopology not only generalises; it simplifies.
The only properties of $\AllBut{N}{f}$ that we will need, are precisely that it is a semitopology.
If all we need about $\AllBut{N}{f}$ is its semitopological structure, then we should work with that structure because that is what actually matters.
\item
Viewing $\AllBut{N}{f}$ as a semitopology, rather than just as a collection of quorums, adds real meaning and understanding.
In particular, the notions of `dense subset' and `nonempty open interior' will be relevant to understanding the axioms, and these arise naturally in the predicates for declarative axiomatisations of distributed algorithms.
These are clearly native topological notions, so the mathematics is speaking to us here, and what it has to say is inherently topological. 
\end{enumerate*} 
\end{rmrk}


\subsection{Syntax of the logic}

We now develop our three-valued modal fixedpoint logic for reasoning about distributed systems.
We so this in the following order: \emph{signatures} $\Sigma$, \emph{models} $\mathcal M$, \emph{contexts} $\acontext$, and finally \emph{syntax}. 
We develop denotation for this syntax in Subsection~\ref{subsect.denotation}.
For quick reference, the reader can just look at the core syntax and denotation in Figures~\ref{fig.predicate.syntax} and~\ref{fig.3.phi.f}.

 


\begin{defn}[Signatures]
\label{defn.signature}
\leavevmode
\begin{enumerate*}
\item\label{item.signature.varsymb}
Fix $\tf{VarSymb}$ a countably infinite set of \deffont{variable symbols} $\thea,\theb,\thec\in\tf{VarSymb}$.
This set will never change and we may assume it henceforth.
\item
Fix a \deffont{fixedpoint variable symbol} $\tf X$ (we will only need one of these, though there is inherent obstacle to having more).
\item\label{item.predicate.syntax.signature}
A \deffont{signature} $\Sigma$ is a pair 
$$
\Sigma=(\PredSymb,\arity)
$$ 
where:
\begin{enumerate*}
\item\label{item.predsymb}
$\PredSymb$ is a (possibly finite) set of \deffont{predicate symbols} $\tf P,\tf P'\in\PredSymb$.
\item
$\arity$ is an \deffont{arity function} $\arity(\tf P)\in\Ngeqz$ mapping each predicate symbol in $\PredSymb$ to a nonnegative natural number (which is intuitively its number of arguments).
\end{enumerate*}
\item
We may write a signature $\Sigma$ in list form as 
$$
\Sigma=[\tf P,\tf P':i,\ \tf Q:j,\ \dots]
$$ 
meaning that $\tf P,\tf P',\tf Q\in\PredSymb$ and $\arity(\tf P)=\arity(\tf P')=i$ and $\arity(\tf Q)=j$ and so on.
\item\label{item.predicate.syntax.empty.signature}
If the signature in part~\ref{item.predsymb} above contains no predicate symbols then we call it \deffont{empty}; otherwise we call it non-empty.
There is only one empty signature, which we will call \deffont{the empty signature}.\footnote{It is not particularly useful.} 
\end{enumerate*}
\end{defn}


\begin{defn}[Models]
\label{defn.model}
A \deffont{model} $\mathcal M$ is a tuple $\mathcal M=(\Pnt,\opens,\Val)$ where:
\begin{enumerate*}
\item
$(\Pnt,\opens)$ is a semitopology (Definition~\ref{defn.semitopology}), which we may call the \deffont{underlying semitopology} of the model.

We may call $p\in\Pnt$ a \deffont{point} or a \deffont{participant}; also \deffont{party}, \deffont{process}, or \deffont{node} would be consistent with the literature.
For our purposes, these are all synonyms: the terminology \emph{point} is from topology; \emph{participant}, \emph{party}, and \emph{process} are from distributed systems; and \emph{node} is from graph theory (thinking of a network as a graph of communicating processes).
\item\label{item.model.val}
$\Val$ is a nonempty set of \deffont{values}.\footnote{We assume values are nonempty to avoid the degenerate case where a universal quantification over values holds just because there are no values on which to test it.  See in particular the argument in Lemma~\ref{lemm.intertwined.char}(\ref{item.intertwined.char.reverse}).}

For our purposes we will not care what values are, so long as some exist, so we do not put any type system on values or endow them with structure; but we could easily do so if required. 
\end{enumerate*} 
\end{defn}


\begin{defn}[Contexts]
\label{defn.contexts}
Suppose $\Sigma=(\PredSymb,\arity)$ is a signature, and $\mathcal M=(\Pnt,\opens,\Val)$ is a model. 
\begin{enumerate*}
\item\label{item.context.stage}
A \deffont{stage} is a nonnegative number $n\in\Time$.
\item\label{item.contexts.valuation}
A \deffont{valuation} $\avaluation$ inputs a predicate symbol $\tf P\in\PredSymb$ and a stage $n\in\Time$ and a point $p\in\Pnt$ and returns a function $\avaluation(\tf P)_{n,p}:\Val^{\arity(\tf P)}\to\THREE$.

We write the set of valuations $\Valuation(\Sigma,\Val)$ or $\Valuation(\Sigma,\mathcal M)$.
\item\label{item.contexts}
A \deffont{context} $\acontext$ is a tuple $(n,p,O,\avaluation)$ where $n\in\Time$ is a stage, $p\in\Pnt$ is a point, $O\in\opensne$ is a nonempty open set, and $\avaluation\in\tf{valuation}(\Sigma,\Val)$ is a valuation. 

We write the set of contexts $\tf{Ctx}(\Sigma,\mathcal M)$; if $\Sigma$ and $\mathcal M$ are understood, we may write this just $\tf{Ctx}$.
\item
A \deffont{context valuation} is a function $\varkappa\in \THREE^{\tf{Ctx}}$.
\end{enumerate*}
\end{defn}



\begin{figure} %
$$
\begin{array}{r@{\ }l}
t ::=& 
(a\in\tf{VarSymb}) \mid (v\in\Val) %
\\
\phi ::=&  
\tvB 
\mid t\teq t \mid (\tf P(\overbrace{t,\dots,t}^{\arity(\tf P)}) : \tf P\in\PredSymb)
\\
\mid &\tneg\phi \mid \modT\phi \mid \someone\phi \mid \Quorum\phi \mid \texi a.\phi
\\
\mid & \phi\tor\phi'  %
\\
\mid & \yesterday \phi \mid \tomorrow \phi 
\\
\mid & \tf X \mid \mu\tf X.\phi \mid (\varkappa \in \THREE^{\tf{Ctx}})
\end{array}
$$ 
\caption{Term and predicate syntax of \QLogic (Definition~\ref{defn.predicate.syntax})}
\label{fig.predicate.syntax}
\end{figure}


\begin{defn}[Syntax]
\label{defn.predicate.syntax}
Suppose $\Sigma$ is a signature (Definition~\ref{defn.signature}) and $\mathcal M$ is a model (Definition~\ref{defn.model}).
\begin{enumerate*}
\item 
Define the syntax of \deffont{\QLogic} over $\Sigma$ and $\mathcal M$ as in Figure~\ref{fig.predicate.syntax}.

In this syntax, $\texi\thea.\phi$ binds $\thea$ in $\phi$, and $\mu\tf X.\phi$ binds $\tf X$ in $\phi$.
We call a term or predicate with no free variable symbols or fixedpoint variable symbols \deffont{closed}.
\item
One special case will be of particular interest: we will write 
\begin{quote}
`suppose $\texi a.\phi$ is a closed predicate',
\end{quote}
and often we do not really care about $\texi a.\phi$ \emph{per se}, except that when this is a closed predicate it means that $a\in\tf{VarSymb}$ and $\phi$ is a predicate whose set of free variables is at most $\{a\}$.
See for example Remark~\ref{rmrk.unpack.texiunique} and Lemma~\ref{lemm.unique.affine.existence}.
\item
We define substitution of variable symbols / fixedpoint variable symbols for values / context valuations as usual, and write these $\phi[a\ssm v]$ and $\phi[\tf X\ssm \varkappa]$ respectively. %
\end{enumerate*}
\end{defn}


\begin{rmrk}
We call our logic \emph{\QLogic} to reflect that it is a logic for reasoning about how coalitions of participants can work together to carry out algorithms to achieve their goals.
In this document we will apply it to Paxos, but we see this logic as having independent interest and a broader relevance for specifying and reasoning about the kinds of algorithms that get written for distributed systems.
Their hallmark is a requirement for distributed collaborative action --- action by quorums, or actionable coalitions, or winning coalitions; these things have many names --- to progress an algorithm to its next step. 
\end{rmrk}

\begin{xmpl}
We give some example predicates and suggest intuitive meanings for them.
We do not yet have the denotation (Figure~\ref{fig.3.phi.f}) nor the rich derived syntax (Figure~\ref{fig.3.derived}) that we build later, so these examples are simple and informal, but we hope they will be helpful.

Assume our signature has a unary predicate symbol $\tf{output}$ (representing the output of some process), and assume $v,v'\in\Val$ are values.
\begin{enumerate}
\item
$\tf{output}(v)$ means `where we are, today, if the process succeeds, then it outputs $v$' and $\tneg\tf{output}(v')$ means `where we are, if the process succeeds, then it does not output $v'$'.
\item
$\modT\tf{output}(v)$ means `where we are, today, the process succeeds and outputs $v$'. 
\item
$\someone\tf{output}(v)$ means `there is a place where, today, if the process succeeds, then it outputs $v$'.
\item
$\tomorrow\modT\tf{output}(v)$ means `where we are, tomorrow, the process succeeds and outputs $v$'.
\item
$\yesterday\modT\tf{output}(v)$ means `where we are, yesterday, the process succeeded and output $v$'.
This assertion is subtly stronger than it might seem; our notion of time starts at day~0, at which by convention all assertions about yesterday are false.
See Remark~\ref{rmrk.further.comments.denotation}(\ref{item.interpret.yesterday}).
\item
$\tvB$ means `where we are, today, the process (which we have not named with a predicate symbol) crashes'.
\end{enumerate}
More sophisticated examples will appear later; notably in the derived syntax in Figure~\ref{fig.3.derived}, in Remark~\ref{rmrk.comments.on.semantics}, and of course when we apply our logic to axiomatise distributed algorithms.
\end{xmpl}


\begin{rmrk}[Three pedant points]
\label{rmrk.pedant.point}
The thing about being pedantic is that it does not matter, until it does.
We now discuss three pedant points, which may not matter \dots until perhaps a moment arrives then they do:
\begin{enumerate}
\item
We may treat $\Sigma$ and $\mathcal M$ as fixed parameters, and elide them.  
In practice this means that we may write `suppose $\phi$ is a closed predicate' instead of `suppose $\phi$ is a closed predicate in the syntax of \QLogic over $\Sigma$ and $\mathcal M$'.
But note next point.
\item\label{item.model.matters}
Our syntax from Definition~\ref{defn.predicate.syntax} and Figure~\ref{fig.predicate.syntax} contains denotational elements; $v\in\Val$ and $\varkappa\in\THREE^{\tf{Ctx}(\Sigma,\mathcal M)}$.

So when we write `suppose $\phi$ is a closed predicate' then this statement is perfectly correct and well-defined but it \emph{really does matter} what the signature and model are (even if the \emph{precise} choice of model will may not matter to the theorems we prove), because closed predicates $\phi$ can include elements from the denotation such as $\varkappa\in\THREE^{\tf{Ctx}(\Sigma,\mathcal M)}$ (this function-space is determined by both the signature $\Sigma$ and the model $\mathcal M$) and $v\teq v'$ for some $v,v'\in\Val$ ($\Val$ is determined by the choice of model $\mathcal M$).

Adding denotational elements to syntax avoids us having to put this information in the context (which is already a $4$-tuple, as per Definition~\ref{defn.contexts}(\ref{item.contexts})) and simplifies the presentation.
This is a technical device with a long pedigree.\footnote{Adding elements from a model to syntax, and conversely building a model out of syntax, is a standard technical device in logic, with a pedigree that goes back (at least) to the proofs of the upward and downward L\"owenheim-Skolem theorems from the 1920s.  A clear survey of these methods is in Chapter~2 (``Models constructed from constants'') of~\cite{chang:modt}.}

Note that (by design!) the choice of model makes a difference to validity.
For example: in Definition~\ref{defn.phi.equivalent} we define equivalence of predicates (and we explicitly mention $\Sigma$ and $\mathcal M$);
if $\mathcal M$ has only a single value $v$, then $\tall a,b.a\teq b$ and $\tall a.a\teq v$ are equivalent to $\tvT$; and if $\mathcal M$ has more than one value then $\tall a,b.a\teq b$ is equivalent to $\tvF$ --- $\tall a.\phi$ is shorthand for $\tneg\texi a.\tneg\phi$; see Figure~\ref{fig.3.derived}.
In contrast: in Lemma~\ref{lemm.equiv} we elide $\Sigma$ and $\mathcal M$ for brevity, and just write `suppose $\phi$ and $\phi'$ are closed predicates and $\acontext$ is a context', without specifying over what signature or model, because this is clear in context.

The reader can mentally add $\Sigma$ and $\mathcal M$ everywhere if they like.
We will always be precise and explicit where it counts.
\item\label{item.pedant.texi.v}
We may quantify over $v\in\Val$ in our formal syntax of predicates, writing things like $\texi v.\texi v'.v\teq v'$ as predicates.
Technically this is wrong and is not syntax! 
Figure~\ref{fig.predicate.syntax} permits e.g. $\texi a.\texi a'.a\teq a'$ where $a,a'\in\tf{VarSymb}$ (Definition~\ref{defn.signature}(\ref{item.signature.varsymb})); it does not permit $\texi v$ in the formal syntax of predicates where $v\in\Val$, so writing e.g. $\texi v.\texi v'.v\teq v'$ as if it were formal syntax of \QLogic is an abuse of notation.

However, $a$ and $a'$ get instantiated to values $v$ and $v'$ in our denotation (see the line for $\texi$ in Figure~\ref{fig.3.phi.f}).
It seems not worth the disruption of maintaining a distinction between $\texi a.\phi$ (which is syntax, assuming $\phi$ is) and $\texi v.\phi[a\ssm v]$ (which looks like syntax but technically is not).
So in practice this seems to us one example of an abuse of formal notation that is positively useful and clarifies meaning.

If this bothers the reader, then every time they see $\texi v.\phi$ (and $\tall v.\phi$, and similar) in our text, they can either declare that $v$ here is \emph{actually} a variable symbol in $\tf{VarSymb}$ which we happen to have called $v$ (and not a value in $\Val$ after all), or they can pick a fresh variable symbol $a$ and replace every unbound $v$ in $\phi$ with $a$.
Either way they will get well-formed syntax and everything will be as it should be.

Notation exists to serve clarity, and not the other way around, so we will simply ignore this henceforth and carry on writing $\texi v.\phi$, just whenever we feel like it, and this will denote just what it says: `there exists a $v$ such that $\phi$'.
Meaning should always be clear. 

Note that an unbound value $v$ in a predicate is legal syntax as per Figure~\ref{fig.predicate.syntax}; 
e.g. in `$\texi a.a\teq v$', $v$ is a value, and this is fine; the (constructive and useful) abuse of notation comes when we write $\texi v$. 
\end{enumerate}
\end{rmrk}

\subsection{Denotation (semantics) of the syntax}
\label{subsect.denotation}


\begin{figure} %
$$
\begin{array}{r@{\ }l@{\quad}l}
\modellabel{\tvB}{n,p,O,\avaluation} =&\tvB
\\
\modellabel{v\teq v'}{n,p,O,\avaluation} =&
\begin{cases}
\tvT & v=v'
\\
\tvF & v\neq v'
\end{cases}
\\
\modellabel{\tf P(v_1,\dots,v_{\arity(\tf P))}}{n,p,O,\avaluation} =&\avaluation(\tf P)_{n,p}(v_1,\dots,v_{\arity(\tf P)}) %
\\[2ex]
\modellabel{\phi\tor\phi'}{n,p,O,\avaluation} =& \modellabel{\phi}{n,p,O,\avaluation}\tor\modellabel{\phi'}{n,p,O,\avaluation}
\\
\modellabel{\tneg\phi}{n,p,O,\avaluation} =& \tneg\,\modellabel{\phi}{n,p,O,\avaluation}  
\\
\modellabel{\modT\phi}{n,p,O,\avaluation} =& \modT\modellabel{\phi}{n,p,O,\avaluation}  
\\
\modellabel{\texi\thea.\phi}{n,p,O,\avaluation} =& \bigvee_{v{\in}\Val} \modellabel{\phi[\thea\ssm v]}{n,p,O,\avaluation}
\\
\modellabel{\someone\phi}{n,p,O,\avaluation} =& \bigvee_{p'\in O} \modellabel{\phi}{n,p',O,\avaluation} 
\\
\modellabel{\Quorum\phi}{n,p,O,\avaluation} =& \bigvee_{O'\in\opensne} \modellabel{\phi}{n,p,O',\avaluation} 
\\[2ex]
\modellabel{\yesterday\phi}{n,p,O,\avaluation} =&
\bigvee_{\f{max}(0,n\minus 1)\leq n'<n}\modellabel{\phi}{n',p,O,\avaluation}
\\
\modellabel{\tomorrow\phi}{n,p,O,\avaluation} =&
\modellabel{\phi}{n\plus 1,p,O,\avaluation}
\\[2ex]
\modellabel{\mu\tf X.\phi}{n,p,O,\avaluation}=&
\bigl(\bigwedge\{ \varkappa\in\THREE^{\tf{Ctx}}\mid \model{\phi[\tf X\ssm\varkappa]}\leq\varkappa\}\bigr)(n,p,O,\avaluation)
\figskip
\modellabel{\varkappa}{n,p,O,\avaluation}=&
\varkappa(n,p,O,\avaluation)
\end{array}
$$ 
\caption{Denotational semantics of \QLogic (Definition~\ref{defn.denotation})}
\label{fig.3.phi.f}
\end{figure}


\begin{figure} %
$$
\begin{array}{r@{\ }l@{\ \ }r@{\ }l}
\tvF =& \modT\tvB
&
\modellabel{\tvF}{\acontext} =&\tvF
\figskip
\tvT =& \tneg\tvF
&
\modellabel{\tvT}{\acontext} =&\tvT
\\[2ex]
\modF\phi =& \modT\tneg\phi
&
\modellabel{\modF\phi}{\acontext} =&\modF\modellabel{\phi}{\acontext}
\figskip
\modB\phi =& \tneg\modT\phi\tand\tneg\modF\phi
&
\modellabel{\modB\phi}{\acontext} =&\modB\modellabel{\phi}{\acontext}
\figskip
\modTB\phi =& \tneg\modT\tneg\phi
&
\modellabel{\modTB\phi}{\acontext} =&\modTB\modellabel{\phi}{\acontext}
\figskip
\modTF\phi =& \modT\phi\tor\modF\phi
&
\modellabel{\modTF\phi}{\acontext} =&\modTF\modellabel{\phi}{\acontext}
\figskip
\modFB\phi =& \modF\phi\tor\modB\phi
&
\modellabel{\modFB\phi}{\acontext} =&\modFB\modellabel{\phi}{\acontext}
\\[2ex]
\phi\tand\phi' =& \tneg(\tneg\phi\tor \tneg\phi')
&
\modellabel{\phi\tand\phi'}{\acontext} =&\modellabel{\phi}{\acontext}\tand\modellabel{\phi'}{\acontext}
\figskip
\phi\tnotor\phi' =& \tneg\phi\tor\phi'
&
\modellabel{\phi\tnotor\phi'}{\acontext} =&\modellabel{\phi}{\acontext}\tnotor\modellabel{\phi'}{\acontext}
\figskip
\phi\timpc\phi' =& \phi\tnotor\modT\phi'
&
\modellabel{\phi\timpc\phi'}{\acontext} =&\modellabel{\phi}{\acontext}\timpc\modellabel{\phi'}{\acontext}
\figskip
\phi\equiv\phi' =& (\modT\phi\tand\modT\phi')\tor(\modB\phi\tand\modB\phi') \tor(\modF\phi\tand\modF\phi') \hspace{-5em}
\\
&&
\modellabel{\phi\equiv\phi'}{\acontext} =&\modellabel{\phi}{\acontext}\equiv\modellabel{\phi'}{\acontext}
\\[2ex]
\tall a.\phi =& \tneg\texi a.\tneg\phi
&
\modellabel{\tall a.\phi}{\acontext} =&\bigwedge_{v{\in}\Val}\modellabel{\phi[a\ssm v]}{\acontext}
\figskip
\everyone \phi =& \tneg\someone\tneg\phi
&
\modellabel{\everyone\phi}{n,p,O,\avaluation} =&\bigwedge_{p'{\in}\Pnt}\modellabel{\phi}{n,p',O,\avaluation}
\figskip
\Coquorum\phi =& \tneg\Quorum\tneg\phi
&
\modellabel{\Coquorum\phi}{n,p,O,\avaluation} =&\bigwedge_{O'{\in}\opensne}\modellabel{\phi}{n,p,O',\avaluation}
\\[2ex]
\QuorumBox\phi =& \Quorum\everyone\phi
&
\modellabel{\QuorumBox\phi}{n,p,O,\avaluation} =&\bigvee_{O'{\in}\opensne}\bigwedge_{p'\in O'}\modellabel{\phi}{n,p',O',\avaluation}
\figskip
\CoquorumDiamond\phi =& \Coquorum\someone\phi
&
\modellabel{\CoquorumDiamond\phi}{n,p,O,\avaluation} =&\bigwedge_{O'{\in}\opensne}\bigvee_{p'\in O'}\modellabel{\phi}{n,p',O',\avaluation}
\figskip
\someoneAll\phi =& \Quorum\someone\phi
&
\modellabel{\someoneAll\phi}{n,p,O,\avaluation} =&\bigvee_{O'{\in}\opensne}\bigvee_{p'\in O'}\modellabel{\phi}{n,p',O',\avaluation}
\figskip
\everyoneAll\phi =& \Coquorum\everyone\phi
&
\modellabel{\everyoneAll\phi}{n,p,O,\avaluation} =&\bigwedge_{O'{\in}\opensne}\bigwedge_{p'\in O'}\modellabel{\phi}{n,p',O',\avaluation}
\\[2ex]
\forever\phi =& \mu\tf X.\tomorrow(\phi\tand \tf X)
&
\modellabel{\forever\phi}{n,p,O,\avaluation} =&\bigwedge_{n'> n}\modellabel{\phi}{n',p,O,\avaluation}
\figskip
\sometime\phi =& \mu\tf X.\tomorrow(\phi\tor \tf X)
&
\modellabel{\sometime\phi}{n,p,O,\avaluation} =&\bigvee_{n'> n}\modellabel{\phi}{n',p,O,\avaluation}
\figskip
\infinitely\phi =& \forever\sometime\phi
&
\modellabel{\infinitely\phi}{n,p,O,\avaluation} =&\bigwedge_{n'> n}\bigvee_{n''>n'}\modellabel{\phi}{n'',p,O,\avaluation}
\figskip
\final\phi =& \sometime\forever\phi
&
\modellabel{\final\phi}{n,p,O,\avaluation} =&\bigvee_{n'> n}\bigwedge_{n''>n'}\modellabel{\phi}{n'',p,O,\avaluation}
\figskip
\recent\phi =& \mu\tf X.\yesterday(\phi\tor\tf X)
&
\modellabel{\recent\phi}{n,p,O,\avaluation} =&\bigvee_{0\leq n'<n}\modellabel{\phi}{n',p,O,\avaluation}
\figskip
\urecent\phi =& \mu\tf X.(\phi\tor\yesterday\tf X)
&
\modellabel{\urecent\phi}{n,p,O,\avaluation} =&\bigvee_{0\leq n'\leq n}\modellabel{\phi}{n',p,O,\avaluation}
\end{array}
$$
$$
\begin{array}{r@{\ }l}
\mrup{a}{\phi}{v} =&
\mu\tf X.(\yesterday(\modT\phi[a\ssm v] \tor ((\tneg\modT\texi a.\phi) \tand \tf X)))
\figskip
\modellabel{\mrup{a}{\phi}{v}}{n,p,O,\avaluation} =&
\begin{cases}
\tvF 
&\Forall{0\leq n'<n} \modellabel{\texi a.\phi}{n',p,O,\avaluation}\neq\tvT
\\
\modT\modellabel{\phi[\thea\ssm v]}{\f{max},p,O,\avaluation}
&\Exists{0\leq n'<n} \modellabel{\texi a.\phi}{n',p,O,\avaluation}=\tvT
\end{cases}
\\
&\qquad\text{where}\quad
\f{max}=\bigvee\{0\leq n'<n \mid \modellabel{\texi a.\phi}{n',p,O,\avaluation}=\tvT\}
\end{array}
$$ 
\caption{Derived syntactic sugar expressions (Definition~\ref{defn.3.derived})}
\label{fig.3.derived}
\end{figure}



Recall from Definition~\ref{defn.signature}(\ref{item.predicate.syntax.signature}) the notion of a \emph{signature} and from Definition~\ref{defn.model} the notion of a \emph{model}:
\begin{defn}[Denotation]
\label{defn.denotation}
Suppose $\Sigma=(\PredSymb,\arity)$ is a signature, and $\mathcal M$ is a model. 
Suppose $\phi$ is a closed predicate (Definition~\ref{defn.predicate.syntax}(1)).
Define the \deffont{denotation} of $\phi$ 
$$
\model{\phi}\in\THREE^{\tf{Ctx}(\Sigma,\mathcal M)}
\quad\text{such that}\quad 
\acontext \longmapsto \modellabel{\phi}{\acontext}
$$
by the rules in Figure~\ref{fig.3.phi.f}.

We will usually treat $\Sigma$ and $\mathcal M$ as fixed parameters, so we may write $\tf{Ctx}(\Sigma,\mathcal M)$ just as $\tf{Ctx}$.
\end{defn}




\begin{defn}[Derived expressions]
\label{defn.3.derived}
The syntax of \QLogic in Figure~\ref{fig.predicate.syntax} and Definition~\ref{defn.predicate.syntax} is very expressive.
In the left-hand column of Figure~\ref{fig.3.derived} we define some derived expressions --- i.e. \emph{syntactic sugar}.
For clarity and convenience we unpack their denotation in the right-hand column.

In principle we could fully unpack our sugar down to the primitive syntax from Figure~\ref{fig.3.phi.f}, but in practice we work directly with Figure~\ref{fig.3.derived} because this is more convenient.
We will discuss syntax further below, notably in Remarks~\ref{rmrk.initial.comments.denotation}, \ref{rmrk.further.comments.denotation}, and~\ref{rmrk.comments.on.semantics}.
\end{defn}


\begin{rmrk}
Usually when we talk about denotation as defined in Definition~\ref{defn.denotation} above, the signature $\Sigma$ and model $\mathcal M$ are either understood or do not matter, so we may not mention them explicitly.
Meaning will always be clear.
\end{rmrk}

\begin{rmrk}[Initial comments on the denotation]
\label{rmrk.initial.comments.denotation}
We make some initial comments on Figure~\ref{fig.3.phi.f}:
\begin{enumerate*}
\item\label{item.first.mu}
$\mu\tf X.\phi$ is a \emph{fixedpoint operator}.
This is a standard method to include the power of recursion / iteration in a logic; an instance of $\tf X$ inside $\mu\tf X.\phi$ can be thought of as a recursive call back to $\phi$.
$\mu\tf X.\phi$ is what would be written $\mathrm{LFP}_{\tf X}\phi$ in~\cite[Chapter~4]{immerman:desc} (a Wikipedia article is also accessible and clear~\cite{wiki:Fixedpoint_logic}).
We will discuss this in Subsection~\ref{subsection.muX}.
\item
In the clause $\model{\tvB}=\tvB$ in Figure~\ref{fig.3.phi.f}, the $\tvB$ on the left is the syntax symbol from Figure~\ref{fig.predicate.syntax}, and the $\tvB$ on the right is $\tvB$-the-truth-value from Definition~\ref{defn.THREE}.
\item
In the clause for $\model{v\teq v'}$ in Figure~\ref{fig.3.phi.f}, the $\teq$ on the left is the syntax symbol from Figure~\ref{fig.predicate.syntax}, and the $\tvT$, $\tvF$, and $=$ on the right are truth-values and (actual) equality.
\item
Similarly for the clauses for $\model{\tneg\phi}$ and $\model{\phi\tor\phi'}$ and $\model{\modT\phi}$; the symbols $\tneg$, $\tor$, and $\modT$ on the left are syntax, and the symbols $\tneg$, $\tor$, and $\modT$ on the right refer to functions on truth-values as per Figure~\ref{fig.3}.
\end{enumerate*} 
\end{rmrk}


\begin{rmrk}[Further comments] 
\label{rmrk.further.comments.denotation}
\leavevmode
\begin{enumerate}
\item\label{item.interpret.yesterday}
The clause 
$\modellabel{\yesterday\phi}{n,p,O,\avaluation} = \bigvee_{\f{max}(0,n\minus 1)\leq n'<n}\modellabel{\phi}{n',p,O,\avaluation}$ is just a one-line encoding of
$$
\modellabel{\yesterday\phi}{n,p,O,\avaluation}=
\begin{cases}
\tvF & n=0
\\
\modellabel{\phi}{n\minus 1,p,O,\avaluation} & n>0 . 
\end{cases}
$$
\item
We see from the right-hand column in Figure~\ref{fig.3.derived} that $\recent\phi$ returns the greatest truth-value that $\phi$ has taken in the past, and $\urecent\phi$ returns the greatest truth-value that $\phi$ has taken today or in the past.

The notation $\urecent$ (`$\recent$', underlined) is a pun\footnote{A bad one perhaps, and the reader may or may not find this pun-acceptable.} on $<$ and $\leq$ (`$<$', underlined).
\item
Informally, we see that the only naturally \emph{non-monotone} operator in the (primitive or derived) syntax is $\tneg$, which is anti-monotone.
This design feature is what powers Lemma~\ref{lemm.positive.monotonicity}.
\end{enumerate}
\end{rmrk}


\begin{nttn}
\label{nttn.logic.terminology}
We take a moment to recall some standard conventions and terminology:
\begin{enumerate*}
\item
Implications $\tnotor$ and $\timpc$ associate to the right, so that $\phi\timpc\psi\timpc\chi$ means $\phi\timpc(\psi\timpc\chi)$.
\item
Unary connectives (like $\tneg$, $\modT$, and $\someone$) bind more strongly than binary connectives (like $\tand$ and $\tnotor$), so that $\QuorumBox\tneg\phi\tand\psi$ means $(\QuorumBox\tneg\phi)\tand\psi$.
\item
The scope of a quantification $\texi a$ and $\tall a$ extends as far to the right as possible. 
Thus $\tall\thea.\phi\tand\psi$ means $\tall\thea.(\phi\tand\psi)$.
Note however that $\QuorumBox\texi\thea.\phi\tand\psi$ still means $(\QuorumBox(\texi\thea.\phi))\tand\psi$, because the $\texi$ quantifier is trapped in the scope of the $\QuorumBox$. 
\item
All the above said, we will (as is usual) use spacing and some common sense to convey scope.
For example 
$$
\texi\thea.\phi\mathbin{\tnotor}\psi
\quad\text{means}\quad
\texi\thea.(\phi\tnotor\psi), 
$$
and
$$
\texi\thea.\phi \ \tnotor \psi
\quad\text{means}\quad
(\texi\thea.\phi)\tnotor\psi. 
$$
In practice, meaning will always be clear.
\item
\begin{enumerate*}
\item\label{item.nullary.ps}
If $\arity(\tf P)=0$ then we may call $\tf P$ \deffont{nullary} or a \deffont{predicate constant symbol}.\footnote{$\tvT$, $\tvB$, and $\tvF$ are predicate constant symbols, but they are special so we give them a distinct identity in the syntax.}
\item\label{item.unary.ps}
If $\arity(\tf P)=1$ then we may call $\tf P$ a \deffont{unary} predicate symbol.
\item\label{item.atomic.predicate}
We may call a predicate of the form $\tf P(v_1,\dots,v_n)$ or $\tneg\tf P(v_1,\dots,v_n)$ an \deffont{atomic predicate}. 
\end{enumerate*}
\end{enumerate*} 
\end{nttn}



\subsection{The fixedpoint operator $\mu \tf X$}
\label{subsection.muX}

\begin{rmrk}
\label{rmrk.comments.on.semantics}
As per Remark~\ref{rmrk.initial.comments.denotation}(\ref{item.first.mu}), our fixedpoint operator $\mu\tf X.\phi$ is a method to include the power of recursion / iteration in the logic.

A mention of $\tf X$ inside $\phi$ in $\mu\tf X.\phi$ can be thought of as a recursive call back to $\phi$, and we can think of $\mu\tf X.\phi$ as encoding an iterative loop that repeatedly executes a computation encoded in $\phi$, and loops if it encounters $\tf X$ in $\phi$.

With this intuition, we can look at some of the derived syntax in Figure~\ref{fig.3.derived} that uses $\mu$:
\begin{enumerate}
\item
$\recent\phi=\mu\tf X.\yesterday(\phi\tor\tf X)$ can be imagined as the following loop:
\begin{enumerate*}
\item
Set an accumulator to $\tvF$.\footnote{We mean `accumulator' in the sense of an `accumulator variable' inside a machine that (we imagine) is running this pseudocode.  This is not a technical term from logic.}
\item\label{recent.step.0}
If we are at day~0 then terminate and return the value of the accumulator; otherwise go to yesterday.
\item
Compute the truth-value of $\phi$, and $\tor$ this truth-value with the accumulator.
\item
Loop back to step~\ref{recent.step.0}.
\end{enumerate*} 
\item
$\forever\phi=\mu\tf X.\tomorrow(\phi\tand\tf X)$ can be viewed as the following loop:
\begin{enumerate*}
\item
Set an accumulator to $\tvT$.
\item\label{forever.step.0}
Go to tomorrow.
\item
Compute the truth-value of $\phi$, and $\tand$ this truth-value with the accumulator.
\item
Loop back to step~\ref{forever.step.0}.
\item
After infinite time (see Remark~\ref{rmrk.oo} about why that makes sense) return the limit value of the accumulator.
\end{enumerate*} 
\item
$\sometime\phi=\mu\tf X.\tomorrow(\phi\tor\tf X)$ can be viewed as executing the following loop:
\begin{enumerate*}
\item
Set an accumulator to $\tvF$.
\item\label{sometime.step.0}
Go to tomorrow.
\item
Compute the truth-value of $\phi$ and $\tor$ with the accumulator.
\item
Loop back to step~\ref{sometime.step.0}.
\item
After infinite time, return the limit value of the accumulator.
\end{enumerate*}
\item\label{item.mru.program}
The clause 
$\mrup{a}{\phi}{v} = \mu\tf X.(\yesterday(\modT\phi[a\ssm v] \tor ((\tneg\modT\texi a.\phi) \tand \tf X)))$
looks more imposing, but the idea is simply to 
\begin{itemize*}
\item
iterate back in time looking for when $\texi a.\phi$ has truth-value $\tvT$, i.e. when \emph{some} value makes $\phi(a)$ be true, and 
\item
if we find such a time then we check if \emph{at that time} $\phi[a\ssm v]$ has truth-value $\tvT$ and 
\item
only if both steps above succeed, do we return $\tvT$.
\end{itemize*}
Thus, we return $\tvT$ when $v$ is one of the values that \emph{most recently} makes $\phi(a)$ be true.

In detail:
\begin{enumerate*}
\item\label{mu.step.0}
If we are at day~0, terminate with value $\tvF$; otherwise go to yesterday.
\item
If $\phi[a\ssm v]$ has truth-value $\tvT$ then terminate with value $\tvT$; otherwise continue to the next step.
\item
If $\texi a.\phi$ has truth-value $\tvT$ then terminate with value $\tvF$; otherwise go to step~\ref{mu.step.0}.
\end{enumerate*}
\end{enumerate}
\end{rmrk}

\begin{rmrk}[Comment on the omniscient observer]
\label{rmrk.oo}
The programs in Remark~\ref{rmrk.comments.on.semantics} above are executed by an omniscient observer who stands outside of the model's notion of time and space.
This is why (for example) $\forever$ can run an infinite loop --- the omniscient observer can observe it because the observer is looking at the model from outside.\footnote{Let's make this concrete.  Consider the natural numbers $\Ngeqz$.  We can easily observe this and note that $n\plus 1\gneq n$ for every $n\in\Ngeqz$; whereas a local observer trapped on the number line and constrained to perceive this line as the passage of time can make similar observation, but from the point of view of that observer, the calculation involved is neverending.  

Note that even if we could speak with this local observer, it would be no good to appeal to it to `just realise that it is on the number line', because in order to check this to a mathematical certainty the observer would have to \dots traverse all of its local notion of time to \emph{check}, amongst other things, that $n\plus 1\gneq n$.  After all, the observer might be on $\mathbb Z/p$ for some astronomically enormous $p$, i.e. their time might be on a large loop.  

The reader interested in a branch of mathematics explicitly devoted to studying these questions, and who is not afraid by the possibility of ending up seriously doubting the nature of their own reality, can start by reading up on nonstandard models of arithmetic.} 
\end{rmrk}

\begin{rmrk}
The rest of this Subsection will show how our denotation for $\mu$ from Figure~\ref{fig.3.phi.f} does indeed yield a fixedpoint operator.
The culminating result is Proposition~\ref{prop.mu.fixedpoint} (which just says that $\mu$ is indeed a fixedpoint operator) but the key technical lemma is Lemma~\ref{lemm.positive.monotonicity} (monotonicity of positive predicates).

The reader is welcome to read the rest of this Subsection, but this is all just a standard use of the Knaster-Tarski fixedpoint theorem~\cite[Theorem~2.35, page~50]{priestley:intlo}.

In fact, the real design effort has already happened: we designed the derived syntax of Figure~\ref{fig.3.derived} so that it would respect the positivity condition of Definition~\ref{defn.positive}, without having to fully expand it out.
For instance, we defined $\modTB\phi$ to be $\tneg\modT\tneg\phi$ (which clearly puts $\phi$ in a positive position) instead of $\modT\phi\tor\modB\phi$ (which is logically equivalent, but which does not).
\end{rmrk} 



\begin{defn}
\label{defn.three.ctx.lattice}
Recall from Remark~\ref{rmrk.three.elementary}(\ref{item.three.elementary.bounded.lattice}) that $\THREE$ with ordering $\tvF<\tvB<\tvT$ is a finite lattice.
As is standard, the function space $\THREE^{\tf{Ctx}}$ (functions from $\tf{Ctx}$ to $\THREE$) inherits this ordering and lattice structure \deffont{pointwise}, as indicated in Figure~\ref{fig.chi.3.lattice}.

This makes $\THREE^{\tf{Ctx}}$ into a \deffont{bounded complete lattice}, meaning that it has a least and greatest element, and arbitrary joins (least upper bounds) and meets (greatest lower bounds).\footnote{Having a least and greatest element follows from having arbitrary joins and meets; the least element is the join of the empty set, and the greatest element is the meet of the empty set.  So in fact a complete lattice is already a bounded complete lattice.} 
\end{defn}


\begin{figure}
$$
\begin{array}{r@{\ }l}
\tbot =& \lambda\acontext.\tvF
\\
\ttop =& \lambda\acontext.\tvT
\\
\bigvee\mathcal X =& \lambda\acontext.\bigvee\{ \varkappa\,\acontext \mid \varkappa\in\mathcal X\}
\\
\bigwedge\mathcal X =& \lambda\acontext.\bigwedge\{ \varkappa\,\acontext \mid \varkappa\in\mathcal X\}
\\[2ex]
\varkappa\leq\varkappa' \liff& \Forall{\acontext} \varkappa\,\acontext\leq\varkappa'\,\acontext
\end{array}
$$
Above, $\mathcal X\subseteq \THREE^{\tf{Ctx}}$ and $\varkappa,\varkappa'\in\THREE^{\tf{Ctx}}$. 
\caption{Bounded complete lattice structure on the function space $\THREE^{\tf{Ctx}}$}
\label{fig.chi.3.lattice}
\end{figure}

\begin{rmrk}
\label{rmrk.leq.phi}
Our denotation $\model{\text{-}}$ from Definition~\ref{defn.denotation} and Figure~\ref{fig.3.phi.f} maps predicates to functions in $\THREE^{\tf{Ctx}}$.
One particular consequence of Definition~\ref{defn.three.ctx.lattice} is that (denotations of) predicates are naturally ordered, and we can write 
$$
\model{\phi}\leq\model{\psi}
\quad\text{which just means}\quad
\Forall{\acontext\in\tf{Ctx}}\modellabel{\phi}{\acontext}\leq\modellabel{\psi}{\acontext}.
$$
We use this fact for example in Lemma~\ref{lemm.positive.monotonicity}.
We can also use this fact, and the complete lattice structure on $\THREE^{\tf{Ctx}}$ from Definition~\ref{defn.three.ctx.lattice}, to simplify the presentation of the clause for $\mu\tf X.\phi$ in Figure~\ref{fig.3.phi.f} as follows:
$$
\model{\mu\tf X.\phi}=
\bigwedge\{ \varkappa\in\THREE^{\tf{Ctx}}\mid \model{\phi[\tf X\ssm\varkappa]}\leq\varkappa\} .
$$
\end{rmrk}

\begin{defn}
\label{defn.positive}
Recall the predicate syntax from Figure~\ref{fig.predicate.syntax} %
and suppose $\phi$ is a predicate.
Then:
\begin{enumerate*}
\item\label{item.positive.occurrence}
Call a free occurrence of $\tf X$ in $\phi$ 
\begin{itemize*}
\item
\deffont{positive} when the occurrence appears under an even number (which may be $0$) of negations $\tneg$, and 
\item
\deffont{negative} when the occurrence appears under an odd number of negations.
\end{itemize*}
For example: all free occurrences of $\tf X$ are positive in $\tf X$ and $\mu\tf X.\tneg\tf X$ and $\yesterday\tf X$; and the free occurrence of $\tf X$ in $\tneg\tf X$ is negative. 
\item
Call a predicate $\phi$ \deffont{positive} when every free occurence of $\tf X$ in $\phi$ is positive.
\end{enumerate*}
\end{defn}


\begin{lemm}
\label{lemm.positive.monotonicity}
Suppose $\mu\tf X.\phi$ is a closed predicate.
Then:
\begin{enumerate*}
\item
If $\phi$ is positive then $\model{\phi}$ is monotone as a function from $\THREE^{\tf{Ctx}}$ with the partial order $\leq$ from Figure~\ref{fig.chi.3.lattice}, to itself.
In symbols, for every $\varkappa,\varkappa'\in\THREE^{\tf{Ctx}}$:
$$
\varkappa\leq\varkappa'
\quad\text{implies}\quad
\model{\phi[\tf X\ssm\varkappa]}\leq\model{\phi[\tf X\ssm\varkappa']} . 
$$
\item
If $\tneg\phi$ is positive then $\model{\phi}$ is antimonotone as a function from $\THREE^{\tf{Ctx}}$ with the partial order $\leq$ from Figure~\ref{fig.chi.3.lattice}, to itself.
In symbols, for every $\varkappa,\varkappa'\in\THREE^{\tf{Ctx}}$:
$$
\varkappa\leq\varkappa'
\quad\text{implies}\quad
\model{\phi[\tf X\ssm\varkappa']}\leq\model{\phi[\tf X\ssm\varkappa]} . 
$$
\end{enumerate*}
(The meaning of $\model{\text{-}}\leq\model{\text{-}}$ is unpacked in Remark~\ref{rmrk.leq.phi}.)

\end{lemm}
\begin{proof}
Both parts follow by a routine simultaneous induction on the syntax of $\phi$, using the rules in Figure~\ref{fig.3.phi.f}.
\end{proof}

\begin{prop}
\label{prop.mu.fixedpoint}
Suppose $\mu\tf X.\phi$ is a closed predicate and $\phi$ is positive.
Then:
\begin{enumerate*}
\item
$\model{\mu\tf X.\phi}$ is a least fixedpoint of the function $\varphi:\THREE^{\tf{Ctx}}\to\THREE^{\tf{Ctx}}$ given by
$$
\varphi(\varkappa)=\lambda\acontext.\modellabel{\phi[\tf X\ssm\varkappa]}{\acontext}. 
$$
\item
If $\acontext$ is a context then
$$
\modellabel{\mu\tf X.\phi}{\acontext}=\modellabel{\phi[\tf X\ssm \model{\mu\tf X.\phi}]}{\acontext}.
$$
\end{enumerate*}
\end{prop}
\begin{proof}
By the Knaster-Tarski fixedpoint theorem for $\THREE^{\tf{Ctx}}$ as a bounded lattice as per Definition~\ref{defn.three.ctx.lattice}.
A concise but exceptionally clear presentation is in~\cite[Theorem~2.35, page~50]{priestley:intlo} (a Wikipedia article~\cite{wiki:Knaster-Tarski_theorem} is also accessible and clear). 
\end{proof}



\subsection{Validity}



We defined in Definition~\ref{defn.tv.ment} what it means for a truth-value to be \emph{valid}; now we extend this notion to predicates-in-context: 
\begin{defn}[Validity]
\label{defn.validity.judgement}
Suppose that:
\begin{itemize*}
\item
$\Sigma$ is a signature (Definition~\ref{defn.signature}(\ref{item.predicate.syntax.signature})).
We take $\Sigma$ to be fixed for the rest of this Definition, which means that we do not explicitly include it (e.g. as a subscript) in the notation.
\item
$\mathcal M$ is a model (Definition~\ref{defn.model}).
Again, we fix $\mathcal M$ as a parameter for the rest of this definition.
\end{itemize*}
Suppose further that:
\begin{itemize*}
\item
$\acontext=(n,p,O,\avaluation)$ is a context in $\tf{Ctx}(\Sigma,\mathcal M)$, so that by Definition~\ref{defn.contexts}(\ref{item.contexts}) $n\in\Time$ and $p\in\Pnt$ and $O\in\opensne$ and $\avaluation$ is a valuation to $\mathcal M$.
\item
$\phi$ is a closed predicate in $\Sigma$ and $\mathcal M$ (cf. Remark~\ref{rmrk.pedant.point}(\ref{item.model.matters})).
\end{itemize*}
Then we define: %
\begin{enumerate*}
\item\label{item.ctx.ment.phi}
$\acontext\ment\phi$ means $\ment\modellabel{\phi}{\acontext}$, which by Definition~\ref{defn.tv.ment}(\ref{item.tv.ment}) means $\modellabel{\phi}{\acontext}\in\threeValid$, which by Definition~\ref{defn.tv.ment}(\ref{item.tv.ment.correct}) means $\modellabel{\phi}{\acontext}=\tvT\lor\modellabel{\phi}{\acontext}=\tvB$.

Because this is a core definition, we sum this up for the reader's clear reference:
$$
\acontext\ment\phi
\quad\liff\quad
{\ment\modellabel{\phi}{\acontext}}
\quad\liff\quad
\modellabel{\phi}{\acontext}\in\threeValid
\quad\liff\quad
\modellabel{\phi}{\acontext}=\tvT\lor\modellabel{\phi}{\acontext}=\tvB.
$$
We read this as \deffont{$\phi$ is valid in the context $\acontext$.}
\item
$n,p\mentval\phi$ means $\Forall{O\in\opensne}(n,p,O\mentval\phi)$.
\item\label{item.ment.n}
$n\mentval\phi$ means $\Forall{p\in\Pnt}\Forall{O\in\opensne}(n,p,O\mentval\phi)$.
\item\label{item.mentval}
$\mentval\phi$ means $\Forall{n\in\Time}\Forall{p\in\Pnt}\Forall{O\in\opensne}(n,p,O\mentval\phi)$.
\item\label{item.ment.phi}
$\ment\phi$ means $\Forall{\acontext\in\tf{Ctx}(\Sigma,\mathcal M)}(\acontext\ment\phi)$.
\end{enumerate*}
\end{defn}

\begin{figure}
$$
\begin{array}{l@{\quad}l@{\quad}l}
\mentval\Theta & \text{means} & \Forall{\theta\in\Theta}({\mentval\theta})
\figskip
\Theta\ment\phi & \text{means} &
\Forall{\avaluation\in\Valuation(\Sigma,\mathcal M)}({\mentval\Theta})\limp ({\mentval\phi}) 
\end{array}
$$
\caption{Validity of $\phi$ in a theory $\Theta$ (Definition~\ref{defn.theory.axiom})}
\label{fig.validity}
\end{figure}


Later on, when we start to build declarative axiomatisations of consensus algorithms, we will express these as \emph{theories}.
We include the relevant definition here, since it follows on naturally from Definition~\ref{defn.validity.judgement}: 
\begin{defn}[Validity in a theory]
\label{defn.theory.axiom}
Suppose that:
\begin{itemize*}
\item
$\Sigma$ is a signature and $\mathcal M=(\Pnt,\opens,\Val)$ is a model.
\item
$\phi$ is a closed predicate in $\Sigma$ and $\mathcal M$.
\item
$\Theta=\{\theta_1,\dots,\theta_n\}$ is a set of closed predicates.\footnote{It would be unusual for $\Theta$, which is intended to be a collection of axioms, to mention specific elements $v\in\Val$ or $\varkappa\in\THREE^{\tf{Ctx}(\Sigma,\mathcal M)}$ from a model (cf. Remark~\ref{rmrk.pedant.point}), but the maths does not seem to mind, so we do not bother to exclude this possibility.}
\end{itemize*}
Then:
\begin{enumerate*}
\item\label{item.theory}
A \deffont{theory} is a pair $(\Sigma,\Theta)$ of 
\begin{itemize*}
\item
a signature $\Sigma$, and 
\item
a set $\Theta$ of \deffont{axioms}, which are closed predicates over $\Sigma$ and $\mathcal M$.
\end{itemize*}
\item\label{item.axioms.valid}
Given a valuation $\avaluation\in\Valuation(\Sigma,\mathcal M)$ (Definition~\ref{defn.contexts}(\ref{item.contexts.valuation})), we define 
$$
\mentval\Theta
\quad\text{and}\quad
\Theta\ment\phi
$$ 
as per Figure~\ref{fig.validity}. 
When $\Theta\ment\phi$ holds, we say that $\phi$ is \deffont{valid in the theory $\Theta$} (in $\Sigma$ and $\mathcal M$). 
\end{enumerate*}
Examples of theories are in Figures~\ref{fig.ThyPaxOne} and~\ref{fig.ThySPax}.
\end{defn}


\begin{rmrk}
\label{rmrk.thy.quantifier.scope}
Note the scope of the quantifiers in Definition~\ref{defn.theory.axiom}:
\begin{enumerate*}
\item
The assertion $\Theta\ment\phi$ does \emph{not} mean 
$$
\Forall{\avaluation}(\mentval(\theta_1\tand\dots\tand\theta_n) \timp \phi),
$$ 
which is equivalent to 
$$
\Forall{\acontext}(\acontext\ment(\theta_1\tand\dots\tand\theta_n)\limp \acontext\mentval\phi) .
$$ 
Above, the short double arrow `$\timp$' means `if the LHS is valid then so is the RHS'; see Remark~\ref{rmrk.rich.implication}(\ref{item.rich.implication.timp}).
The long double arrow `$\limp$' is just logical implication. 
\item
The assertion $\Theta\ment\phi$ means 
$$
\Forall{\avaluation}(\mentval\theta_1\tand\dots\tand\theta_n) \limp (\mentval\phi) .
$$
In words: if $\Theta$ holds everywhere in the model under the valuation $\avaluation$, then $\phi$ holds everywhere in the model under the valuation $\avaluation$.
\end{enumerate*}
We spell this out further.
The assertion $\Theta\ment\phi$ means:
$$
\Forall{\avaluation}(\Forall{n,p,O}n,p,O\mentval\theta_1\tand\dots\tand\theta_n) \limp (\Forall{n,p,O}n,P,O\mentval\phi),
$$
where above $n\in\Time$, $p\in\Pnt$, and $O\in\opensne$.

This pattern of quantifiers is what makes $\Theta$ a theory, rather than just a set of assumptions: as a theory, $\Theta$ is a restriction on the valuation $\avaluation$.
\end{rmrk}


\begin{rmrk}
One final pedant point: 
in Definition~\ref{defn.theory.axiom}(\ref{item.theory}) we defined a theory to be a pair $(\Sigma,\Theta)$ of a signature and a set of closed predicates over a signature $\Sigma$ \emph{and} over a model $\mathcal M$.
This technical convenience lets us write axioms that directly reference specific values in the model (instead of introducing explicit constant symbols to our signature, and interpreting them as \dots values in the model).

We will exploit to assume a distinguished `undefined' value $\udfn$ in axiom \rulefont{PaxPropose?} in Figure~\ref{fig.ThyPaxOne}, as per Definition~\ref{defn.ThyPaxOne}(\ref{item.ThyPaxOne.axioms}).
\end{rmrk}

\subsection{Comments on notation: eliding parameters, and functional notation}

\begin{nttn}[Eliding parameters]
\label{nttn.validity.irrelevant.context}
It will be quite common that we make validity judgements where it is clear that part of the context is irrelevant.
For example the reader can check the denotational semantics in Figure~\ref{fig.3.phi.f} and the definition of validity in Definition~\ref{defn.validity.judgement}(\ref{item.ctx.ment.phi}), to see that:
\begin{enumerate*}
\item
$\acontext \ment \tvB$ is valid; we do not need to know the values in $\acontext$. 
\item
The validity of $n,p,O\mentval \tf P(v)$ depends on $\avaluation$, $n$, and $p$ --- but not on $O$.
\item
The validity of $n,p,O\mentval\QuorumBox\phi$ depends on $n$ and $\avaluation$ (and of course on $\phi$), but it does not depend on the values of $p$ and $O$.
\end{enumerate*}
In such cases, we may omit parameters in the context if they do not matter --- we may even omit the entire context if appropriate.
Thus, we may write $\ment\tvT$, $\ment \thea\teq\thea$, and $n\mentval\QuorumBox\phi$.

This notation remains within the scope of Definition~\ref{defn.validity.judgement} --- which universally quantifies over any omitted parameter --- so omitting parameters in this way is not an abuse of notation.
However, if the reader sees us omit a value from a context, there may be two reasons for it:
\begin{enumerate*}
\item
It may indicate a universal quantification over all possible values.
For example, this is clearly what we intend when we write $\mentval\theta$ in Figure~\ref{fig.validity} (top-right).
\item
It may indicate that we have elided irrelevant parameters.
For example, this is clearly what we intend in the example of $\ment\tvT$ above.
\end{enumerate*}
In practice, it will always be clear which is intended and meaning will always be clear.\footnote{Actually, meaning will be \emph{clearer}.  

In early drafts of this paper, we dutifully included all parameters, whether relevant or not.  
But this turned out to be just too confusing: the notational overhead of carrying around values that we did not care about was burdensome; but furthermore, it was actually easier to understand the sequents once we have dropped irrelevant parameters, because this helped us to keep track of which parameters \emph{did} matter.}
\end{nttn}

\begin{figure}
$$
\begin{array}{l@{\quad}l@{\quad}l@{\quad}l@{\quad}l}
\texiunique  a.\tf{write}(a)
&
\QuorumBox\texi a.\tf{send}(a)
&
\texi a.\everyoneAll\tf{decide}(a)
&
\tall a.\modTF \tf{propose}(a)
&
\mrup{a}{(\someone\tf{accept}(a))}{v}
\\[1ex]
& \text{may be written as}
\\[1ex]
\texiunique \tf{write} 
&
\QuorumBox\texi\tf{send} 
&
\texi\everyoneAll\tf{decide}
&
\tall\modTF\tf{propose} 
&
\mru{(\someone\tf{accept})}{v}
\end{array}
$$
\caption{Abbreviations inspired by higher-order logic (Notation~\ref{nttn.hos})}
\label{fig.hos}
\end{figure}


\begin{nttn}[Functional notation]
\label{nttn.hos}
We will frequently apply quantifiers (like $\texi$, $\texiunique$, or $\tall$) to unary predicate symbols, possibly with modalities.
For example:
$$
\tall a.\QuorumBox\tf{decide}(a)
\quad\text{is}\quad\text{$\tall$ quantifying $a$ in\ $(\QuorumBox\tf{decide})$\ applied to $a$.}
$$ 
For clarity, we may use a \deffont{functional notation} inspired by higher-order logic (HOL) to elide the variable, as illustrated in Figure~\ref{fig.hos}.
Meaning will always be clear.

See also Remark~\ref{rmrk.functional.notation.in.paxos} for examples of this notation applied in an axiomatisation.
\end{nttn}






\section{Properties of predicates}
\label{sect.predicates}

This Section collects elementary technical results which we will need later.
Some of these results, and their proofs, are more technical than others, but they are not necessarily surprising if we look at what they mean in the model.
Mostly the statements are simple and it is obvious by common sense that they should work (simple is good, because it means greater productivity and fewer errors).

A typical example is Lemma~\ref{lemm.tomorrow.forever.recent}(\ref{item.tomorrow.forever.recent.1}): once we look at what it actually says about the model, this is that ``if $\phi$ is true today, then in the future it will have been true in the past''.
For saying this a mathematician earns a living? --- we might think.

What is less evident is \emph{which} elementary properties will be important later, nor is it obvious what definitions and structures make these properties all fit together nicely.
So the reader should note that the simplicity of this Section is a design feature, which was only attained in retrospect and after some work.

Broadly speaking, we start this Section with the most elementary results, and build complexity from there.
The reader is welcome to read this Section now, or to skip it and refer back to its results as they are used in later proofs.


\subsection{Equivalence of predicates, and some (de Morgan) equivalences}
\label{subsect.equivalence.of.predicates}

For the rest of this Section we make a fixed but arbitrary choice of signature $\Sigma$ (Definition~\ref{defn.signature}(\ref{item.predicate.syntax.signature})) and model $\mathcal M$ (Definition~\ref{defn.model}).
All predicates considered will be over $\Sigma$ and $\mathcal M$, unless otherwise stated.

We generalise equivalence $\equiv$ (Definition~\ref{defn.equivalent}) from truth-values, to predicates viewed as functions from contexts to truth-values:
\begin{defn}[Equivalence of predicates]
\label{defn.phi.equivalent}
Suppose $\phi$ and $\phi'$ are closed predicates.

Call $\phi$ and $\phi'$ \deffont{extensionally equivalent} when 
$$
\model{\phi}=\model{\phi'} .
$$
Unpacking Definition~\ref{defn.denotation}, this means 
$$
\Forall{\acontext\in\tf{Ctx}(\Sigma,\mathcal M)}(\modellabel{\phi}{\acontext}=\modellabel{\phi'}{\acontext}) .
$$
For brevity, we may write `extensionally equivalent' just as `equivalent' henceforth.
\end{defn}



\begin{lemm}
\label{lemm.equiv}
Suppose $\phi$ and $\phi'$ are closed predicates and $\acontext$ is a context.
Then: 
\begin{enumerate*}
\item\label{item.equiv.1}
$\acontext\ment\phi\equiv\phi'$ if and only if $\modellabel{\phi}{\acontext}=\modellabel{\phi'}{\acontext}$.
\item\label{item.equiv.2}
$\ment\phi\equiv\phi'$ if and only if $\model{\phi}=\model{\phi'}$ (i.e. $\phi$ and $\phi'$ are equivalent as per Definition~\ref{defn.phi.equivalent}).
\end{enumerate*}
\end{lemm}
\begin{proof}
We consider each part in turn:
\begin{enumerate}
\item
By Definition~\ref{defn.validity.judgement}(\ref{item.ctx.ment.phi}) $\acontext\ment\phi\equiv\phi'$ means $\modellabel{\phi\equiv\phi'}{\acontext}\in\threeValid$, and by the clause for $\equiv$ in Figure~\ref{fig.3.derived} this is if and only if $\modellabel{\phi}{\acontext}=\modellabel{\phi'}{\acontext}$.
\item
By Definition~\ref{defn.validity.judgement}(\ref{item.ment.phi}), $\ment\phi\equiv\phi'$ means precisely $\Forall{\acontext}\acontext\ment\phi\equiv\phi'$.
We use part~\ref{item.equiv.1} of this result.
\qedhere\end{enumerate}
\end{proof}


We will see many equivalences in what follows.
We start by using Definition~\ref{defn.phi.equivalent} and Lemma~\ref{lemm.equiv} to express some standard equivalences in logic, having to do with implication being expressible using negation and disjunction, and with \emph{de Morgan} dualities between modalities and connectives:
\begin{lemm}
\label{lemm.minimal.syntax}
For each of the equivalences $\f{LHS}\equiv\f{RHS}$ listed in Figure~\ref{fig.easy.equivalences}, we have that $\f{LHS}$ is extensionally equal to $\f{RHS}$ in the sense of Definition~\ref{defn.phi.equivalent}.
\end{lemm}
\begin{proof}
By routine calculations using Figures~\ref{fig.3.phi.f} and~\ref{fig.3.derived} and facts of the lattice $\THREE=(\tvF<\tvB<\tvT)$.
\end{proof}


\begin{figure}
$$
\begin{array}{r@{\ }l@{\quad}r@{\ }l@{\quad}r@{\ }l@{\quad}r@{\ }l}
\phi\tand\phi' \equiv& \tneg(\tneg\phi\tor\tneg\phi')
&
\phi\tor\phi' \equiv& \tneg(\tneg\phi\tand\tneg\phi')
&
\tneg\phi\equiv& \phi\timpc\tvF \equiv\phi\tnotor\tvF %
\\
\phi\timpc\phi' \equiv& \phi\tnotor\modT\phi'
&
\phi\tnotor\phi' \equiv& \tneg\phi\tor\phi'
&
\phi\timpc\phi' \equiv& \phi\tnotor\modT\phi'
\\
\everyone\phi\equiv&\tneg\someone\tneg\phi
&
\someone\phi\equiv&\tneg\everyone\tneg\phi
\\
\Coquorum\phi\equiv&\tneg\Quorum\tneg\phi
&
\Quorum\phi\equiv&\tneg\Coquorum\tneg\phi
\\
\forever\phi\equiv&\tneg\sometime\tneg\phi
&
\sometime\phi\equiv&\tneg\forever\tneg\phi
\\
\tall a.\phi\equiv&\tneg\texi a.\tneg\phi
&
\texi a.\phi\equiv&\tneg\tall a.\tneg\phi
&
\modTF\phi\equiv&\modT\phi\tor\modT\tneg\phi
\\
\modB\phi\equiv&\tneg\modTF\phi
&
\modT\phi\equiv&\tneg\modFB\phi
&
\modF\phi\equiv&\tneg\modTB\phi
\end{array}
$$
\caption{Some easy equivalences}
\label{fig.easy.equivalences}
\end{figure}


\subsection{Two commutation lemmas}
\label{subsect.two.commutation.results}

We note some simple commutation lemmas.
These results are all easy facts of our model, but it only takes a few lines to spell them out:

\begin{lemm}[Commuting connectives]
\label{lemm.commuting.connectives}
Suppose $\phi$ and $\phi'$ are closed predicates and suppose
$$
\bullet\in\{\tor,\tand\} 
\ \ \text{and}\ \ 
\tf M\in\{\modT,\modTB\} . 
$$
Then
$$
\ment\tf M(\phi\bullet\phi')\equiv (\tf M\phi\bullet\tf M\phi') .
$$
\end{lemm}
\begin{proof}
Routine from Lemma~\ref{lemm.equiv}(\ref{item.equiv.2}) and Figures~\ref{fig.3.phi.f} and~\ref{fig.3.derived}, and facts of $\THREE$ as a lattice.
\end{proof}


\begin{lemm}[Commuting modalities]
\label{lemm.someone.commutation}
Suppose $\phi$ is a closed predicate and suppose 
$$
\mathsf Q\in\{\texi,\tall,\someone,\everyone,\Quorum,\Coquorum,\recent,\urecent,\yesterday,\tomorrow,\forever,\sometime,\infinitely,\final\}.
$$
Then 
$$
\ment\mathsf Q\modT\phi \equiv \modT\mathsf Q\phi
\quad\text{and}\quad
\ment\mathsf Q\modTB\phi \equiv \modTB\mathsf Q\phi.
$$
\end{lemm}
\begin{proof}
Routine from Lemma~\ref{lemm.equiv}(\ref{item.equiv.2}) and Figures~\ref{fig.3.phi.f} and~\ref{fig.3.derived}, and facts of $\THREE$ as a lattice.
\end{proof}


\subsection{Lattice inequalities}

Corollary~\ref{corr.recent.everyone.commute} is in the same spirit as Lemmas~\ref{lemm.commuting.connectives} and~\ref{lemm.someone.commutation}. 
In practice we usually just use it a direct fact of the model, but we spell it out anyway.
It relies on this simple fact about lattices: 
\begin{lemm} 
\label{lemm.quantifier.swap}
Suppose $(\mathcal L,\tand,\tor)$ is a complete lattice and $X$ and $Y$ are sets and $f:(X\times Y)\to\mathcal L$ is any function. 
Then:
\begin{enumerate*}
\item\label{item.quantifier.swap.1}
$\bigvee_{x}\bigwedge_{y}f(x,y)\leq \bigwedge_{y}\bigvee_{x}f(x,y)$, 
\item\label{item.quantifier.swap.1b}
It is not in general the case that 
$\bigwedge_{y}\bigvee_{x}f(x,y) \leq \bigvee_{x}\bigwedge_{y}f(x,y)$.
\item\label{item.quantifier.swap.2}
$\bigvee_{x}\bigvee_{y}f(x,y) = \bigvee_{y}\bigvee_{x}f(x,y)$.
\end{enumerate*}
\end{lemm}
\begin{proof}
By easy lattice calculations. 
\end{proof}


\begin{corr}
\label{corr.recent.everyone.commute}
Suppose $\mathsf Q\phi$ is a closed predicate, where 
$$
\mathsf Q\in\{\texi a,\tall a \mid a{\in}\tf{VarSymb}\}\cup\{\someone,\everyone,\Quorum,\Coquorum\},
$$
and suppose $\acontext$ is a context.
Then:
\begin{enumerate*}
\item\label{item.recent.everyone.commute.1}
$\modellabel{\recent\mathsf Q\phi}{\acontext}\leq \modellabel{\mathsf Q\recent\phi}{\acontext}$. 
\item
$\modellabel{\mathsf Q\recent\phi}{\acontext}\leq \modellabel{\recent\mathsf Q\phi}{\acontext}$ for 
$\mathsf Q\in\{\texi a\mid a{\in}\tf{VarSymb}\}\cup\{\someone,\Quorum\}$.
\item
It need not be the case that $\modellabel{\mathsf Q\recent\phi}{\acontext}\leq \modellabel{\recent\mathsf Q}{\acontext}$ for 
$\mathsf Q\in\{\tall a\mid a{\in}\tf{VarSymb}\}\cup\{\everyone,\Coquorum\}$.
\end{enumerate*}
\end{corr}
\begin{proof}
Suppose $n\in\Time$ and $p\in\Pnt$ and $O\in\opensne$ and $\avaluation$ is a valuation.
We freely use \strongmodusponens.
\begin{enumerate}
\item
We consider $\mathsf Q=\Coquorum$.
We use the clause for $\recent$ in Figure~\ref{fig.3.derived}, the clause for $\Coquorum$ in Figure~\ref{fig.3.phi.f}, and then Lemma~\ref{lemm.quantifier.swap}(\ref{item.quantifier.swap.1}): 
$$
\bigvee_{0{\leq}n'{<}n}\bigwedge_{O'{\in}\opensne}\modellabel{\phi}{n',p,O',\avaluation}
\stackrel{L\ref{lemm.quantifier.swap}(\ref{item.quantifier.swap.1})}{\leq}
\bigwedge_{O'{\in}\opensne}\bigvee_{0{\leq}n'{<}n}\modellabel{\phi}{n',p,O',\avaluation} .
$$
Now we consider $\mathsf Q=\tall a$:
$$
\bigvee_{0{\leq}n'{<}n}\bigwedge_{v{\in}\Val}\modellabel{\phi[a\ssm v]}{n',p,O,\avaluation}
\stackrel{L\ref{lemm.quantifier.swap}(\ref{item.quantifier.swap.1})}{\leq}
\bigwedge_{v{\in}\Val}\bigvee_{0{\leq}n'{<}n}\modellabel{\phi[a\ssm v]}{n',p,O,\avaluation} .
$$
The other cases are precisely similar.
\item
We consider just $\mathsf Q=\someone$; the other cases are precisely similar.
We just use Lemma~\ref{lemm.quantifier.swap}(\ref{item.quantifier.swap.2}): 
$$
\bigvee_{p'{\in}\Pnt}\bigvee_{0{\leq}n'{<}n}\modellabel{\phi}{n',p',O,\avaluation}
\stackrel{L\ref{lemm.quantifier.swap}(\ref{item.quantifier.swap.2})}{=}
\bigvee_{0{\leq}n'{<}n}\bigvee_{p'{\in}\Pnt}\modellabel{\phi}{n',p',O,\avaluation} .
$$
\item
We consider just $\mathsf Q=\tall$; the other cases are precisely similar.
We note by Lemma~\ref{lemm.quantifier.swap}(\ref{item.quantifier.swap.1b}) that 
$$
\bigwedge_{v{\in}\Val}\bigvee_{0{\leq}n'_v{<}n}\modellabel{\phi[a\ssm v]}{n'_v,p,O,\avaluation}
\leq
\bigvee_{0{\leq}n'{<}n}\bigwedge_{v\in\Val}\modellabel{\phi[a\ssm v]}{n',p,O,\avaluation} 
\ \ \ \text{does not hold in general}.
$$
\qedhere\end{enumerate}
\end{proof}





\subsection{The compound modalities $\someoneAll$, $\everyoneAll$, $\QuorumBox$, and $\CoquorumDiamond$}
\label{subsect.compound.modalities}

\begin{rmrk}
\label{rmrk.compound.modalities}
We will make much use of the modalities $\someoneAll$, $\everyoneAll$, $\QuorumBox$, and $\CoquorumDiamond$.

We call these four combinations \deffont{compound modalities}, because they are built out of $\Quorum$ and $\Coquorum$ (which quantify over open sets) and $\someone$ and $\everyone$ (which quantify over points in an open set). 
In this Subsection we study some of their properties.
\end{rmrk}



Lemma~\ref{lemm.dense.char} just unpacks definitions:
\begin{lemm}
\label{lemm.dense.char}
Suppose $\phi$ is a closed predicate and $n\in\Time$ and $\avaluation$ is a valuation.
Then: 
\begin{enumerate*}
\item\label{item.dense.char.someoneAll}
$n\mentval\someoneAll\phi$ if and only if $\Exists{O\in\opensne}O\between\{p\in\Pnt \mid n,p,O\mentval\phi\}$.
\item\label{item.dense.char.everyoneAll}
$n\mentval\everyoneAll\phi$ if and only if $\Forall{O\in\opensne}O\subseteq\{p\in\Pnt \mid n,p,O\mentval\phi\}$.
\item\label{item.dense.char.QuorumBox}
$n\mentval\QuorumBox\phi$ if and only if $\Exists{O\in\opensne}O\subseteq\{p\in\Pnt \mid n,p,O\mentval\phi\}$.
\item\label{item.dense.char.CoquorumDiamond}
$n\mentval\CoquorumDiamond\phi$ if and only if $\Forall{O\in\opensne}O\between\{p\in\Pnt \mid n,p,O\mentval\phi\}$.
\end{enumerate*} 
\end{lemm}
\begin{proof}
We consider each part in turn:
\begin{enumerate}
\item
We reason as follows: 
$$
\begin{array}{r@{\ }l@{\quad}l}
n\mentval\someoneAll\phi
\liff&
\Exists{O{\in}\opensne}\Exists{p{\in}O}n,p,O\mentval\phi
&\text{Figure~\ref{fig.3.phi.f}}
\\
\liff&
\Exists{O{\in}\opensne}O\between \{p\in\Pnt \mid n,p,O\mentval\phi\}
&\text{Fact}
\end{array}
$$
\item
We reason as follows: 
$$
\begin{array}{r@{\ }l@{\quad}l}
n\mentval\everyoneAll\phi
\liff&
\Forall{O{\in}\opensne}\Forall{p{\in}O}n,p,O\mentval\phi
&\text{Figure~\ref{fig.3.phi.f}}
\\
\liff&
\Forall{O{\in}\opensne}O\subseteq\{p\in\Pnt \mid n,p,O\mentval\phi\}
&\text{Fact}
\end{array}
$$
\item
We reason as follows:
$$
\begin{array}{r@{\ }l@{\quad}l}
n\mentval\QuorumBox\phi
\liff&
\Exists{O{\in}\opensne}\Forall{p{\in}O}n,p,O\mentval\phi
&\text{Figure~\ref{fig.3.phi.f}}
\\
\liff&
\Exists{O{\in}\opensne}O\subseteq\{p\in\Pnt \mid n,p,O\mentval\phi\}
&\text{Fact}
\end{array}
$$
\item
We reason as follows:
$$
\begin{array}{r@{\ }l@{\quad}l}
n\mentval\CoquorumDiamond\phi
\liff&
\Forall{O{\in}\opensne}\Exists{p{\in}O}n,p,O\mentval\phi
&\text{Figure~\ref{fig.3.phi.f}}
\\
\liff&
\Forall{O{\in}\opensne}O\between\{p\in\Pnt \mid n,p,O\mentval\phi\}
&\text{Fact}
\end{array}
$$
\qedhere\end{enumerate}
\end{proof}



Proposition~\ref{prop.someone.implies.someoneAll} packages up some simple corollaries of Lemma~\ref{lemm.dense.char}, in forms that will be useful for our later proofs: 
\begin{prop}
\label{prop.someone.implies.someoneAll}
Suppose $\phi$ is a closed predicate and $\acontext$ is a context.
Then:
\begin{enumerate*}
\item\label{item.someone.implies.someoneAll.4}
$\acontext\ment\QuorumBox\phi\tand\CoquorumDiamond\psi\limp\acontext\ment\someoneAll(\phi\tand\psi)$.
\item\label{item.someone.implies.someoneAll.4b}
$\acontext\ment\QuorumBox\phi\tand\everyoneAll\psi\limp\acontext\ment\QuorumBox(\phi\tand\psi)$.
\item\label{item.someone.implies.someoneAll.5}
$\acontext\ment\modT(\QuorumBox\phi\tand\CoquorumDiamond\psi)\limp\acontext\ment\modT\someoneAll(\phi\tand\psi)$,\ and
\\
$\acontext\ment(\QuorumBox\phi\tand\CoquorumDiamond\psi)\timpc \someoneAll(\phi\tand\psi)$.
\item\label{item.someone.implies.someoneAll.2}
$\acontext\ment\modT\QuorumBox\phi\limp\acontext\ment\modT\someoneAll\phi$.
\item\label{item.someone.implies.someoneAll.2b}
$\acontext\ment\modT\CoquorumDiamond\phi\limp\acontext\ment\modT\someoneAll\phi$.
\end{enumerate*}
\end{prop}
\begin{proof}
We consider each part in turn:
\begin{enumerate}
\item
Routine from Lemma~\ref{lemm.dense.char}(\ref{item.dense.char.QuorumBox}\&\ref{item.dense.char.CoquorumDiamond}).
\item
Routine from Lemma~\ref{lemm.dense.char}(\ref{item.dense.char.QuorumBox}\&\ref{item.dense.char.everyoneAll}).
\item
From part~\ref{item.someone.implies.someoneAll.4} of this result taking $\modT\phi$ and $\modT\psi$ and the commutation Lemmas~\ref{lemm.someone.commutation} and~\ref{lemm.commuting.connectives}.
Then we use \strongmodusponens.
\item
From part~\ref{item.someone.implies.someoneAll.5} of this result, taking $\psi=\tvT$.
\item
From part~\ref{item.someone.implies.someoneAll.5} of this result, taking $\phi=\tvT$.
\qedhere\end{enumerate}
\end{proof}



 







\subsection{Some comments on time}

\begin{rmrk}
This Subsection notes Lemmas~\ref{lemm.tomorrow.forever.recent} and~\ref{lemm.persist.char}.
They both have to do with time, and underneath all their complicated-looking symbols they are both straightforward:
\begin{itemize*}
\item
Lemma~\ref{lemm.tomorrow.forever.recent} reflects a fact that if something holds now, then in the future it will have held.
\item
Lemma~\ref{lemm.persist.char} reflects the fact that something repeats forever if and only if it happens infinitely often.
\end{itemize*}
\end{rmrk}


\begin{lemm}
\label{lemm.tomorrow.forever.recent}
Suppose $\acontext=(n,p,O,\avaluation)$ is a context and $\phi$ is a closed predicate.
Then:
\begin{enumerate*}
\item\label{item.tomorrow.forever.recent.1}
$n,p,O\mentval\modT\phi$ implies $\Forall{n'{\gneq}n}\,n',p,O\mentval\modT\recent\phi$.
\item\label{item.rur.r}
$n,p,O\mentval\modT\urecent\phi$ implies $\Forall{n'{\geq} n}\,n',p,O\mentval\modT\urecent\phi$.
\item
As corollaries, 
$$
\acontext\ment\phi\timpc \forever\recent\phi
\quad\text{and}\quad
\acontext\ment\urecent\phi\timpc\forever\urecent\phi.
$$
\end{enumerate*}
\end{lemm}
\begin{proof}
This is just from the clauses for $\recent$ and $\urecent$ in Figure~\ref{fig.3.derived}. 
The corollary follows using \strongmodusponens.
\end{proof}



$\infinitely=\forever\sometime$ has a characterisation as `infinitely often', with the convenient feature that its denotation is independent of the stage at which it is evaluated (called $n'$ in the statement of the Lemma below):
\begin{lemm}
\label{lemm.persist.char}
Suppose $\phi$ is a closed predicate and $p\in\Pnt$ and $O\in\opensne$ and $\acontext$ is a context.
Suppose further that $n'\in\Time$ and $\f{tvs}\in\{\threeTrue,\threeValid\}$.
Then we have
$$
\modellabel{\infinitely\phi}{n',p,O,\avaluation}\in\f{tvs}
\quad\liff\quad
\{n\in\Time \mid \modellabel{\phi}{n,p,O,\avaluation}\in\f{tvs}\}\text{ is infinite} . 
$$
\end{lemm}
\begin{proof}
By routine arguments from the clause for $\infinitely$ in Figure~\ref{fig.3.derived}.
\end{proof}






\subsection{Unique existence and affine existence in three-valued logic}
\label{subsect.unique.affine.existence}

\subsubsection{The definition, and main result}

\begin{rmrk}
Distributed algorithms often require an implementation to make (at most) one choice from (zero or more) possibilities.
In our axioms, we model this using unique existence and affine existence:
\begin{itemize*}
\item
\emph{Unique existence} means `there exists precisely one'; we write this $\exists!$, as standard.
\item
\emph{Affine existence} means `there exist zero or one'; we write this $\exists_{01}$.\footnote{We are not aware of a standard symbol for affine existence.}
\end{itemize*}
See for example \rulefont{PaxPropose01} and \rulefont{PaxWrite01} in Figure~\ref{fig.logical.paxos} (a discussion of which is in Remark~\ref{rmrk.paxos.axioms.discussion}(\ref{item.01.discussion})).

So we need to think about what $\exists!$ and $\exists_{01}$ even mean in a three-valued setting.
This turns out to be not completely straightforward, because we have a third truth-value that is neither true nor false.

For concreteness, suppose $\tf P$ is a unary predicate and $\acontext$ is a context, and suppose that 
$\modellabel{\phi[a\ssm v]}{\acontext}=\tvB$ for all $v\in\Val$.
Then what truth-values should $\modellabel{\texiaffine a.\phi}{\acontext}$ and $\modellabel{\texiunique a.\phi}{\acontext}$ have?
 
Now suppose $v,v',v''\in\Val$, and suppose that $\modellabel{\phi[a\ssm v''']}{\acontext}=\tvF$ for all $v'''\in\Val\setminus\{v,v',v''\}$ and
$$
\modellabel{\phi[a\ssm v]}{\acontext}=\tvT
\quad\land\quad
\modellabel{\phi[a\ssm v']}{\acontext}=\tvB
\quad\land\quad
\modellabel{\phi[a\ssm v'']}{\acontext}=\tvB .
$$
\noindent Then what truth-values should $\modellabel{\texiaffine a.\phi}{\acontext}$ and $\modellabel{\texiunique a.\phi}{\acontext}$ have?

Now suppose that $\modellabel{\phi[a\ssm v''']}{\acontext}=\tvF$ for all $v'''\in\Val\setminus\{v\}$ and
$$
\modellabel{\phi[a\ssm v]}{\acontext}=\tvB .
$$
\noindent Then what truth-values should $\modellabel{\texiaffine a.\phi}{\acontext}$ and $\modellabel{\texiunique a.\phi}{\acontext}$ have?

Our answers to all questions above will be: $\tvB$ for both.
This is not automatic: it is a real design decision that we have to make.
\end{rmrk}

\begin{defn}
\label{defn.unique.exist}
Suppose $\phi$ is a predicate.
Define \deffont{affine existence} $\texiaffine a.\phi$ and \deffont{unique existence} $\texiunique  a.\phi$ by:
$$
\begin{array}{r@{\ }l}
\texiaffine a.\phi =& \tall a,a'.(\phi\tand\phi[a\ssm a']) \tnotor a\teq a'
\\
\texiunique  a.\phi =& (\texi a.\phi) \tand (\texiaffine a.\phi)
\end{array}
$$
\end{defn}


\begin{rmrk}
\label{rmrk.unpack.texiunique}
Suppose $\texi a.\phi$ is a closed predicate (i.e. $a\in\tf{VarSymb}$ and $\phi$ is a predicate with free variable symbols at most $\{a\}$), and suppose $\acontext$ is a context.
We spell unpack the denotations of $\texiunique$ and $\texiaffine$ from Definition~\ref{defn.unique.exist}: 
\begin{enumerate*}
\item
If $\exists v,v'\in\Val.v\neq v'\land\modellabel{\phi[a\ssm v]}{\acontext}=\modellabel{\phi[a\ssm v']}{\acontext}=\tvT$, then 
$$
\modellabel{\texiunique a.\phi}{\acontext}=\tvF=\modellabel{\texiaffine a.\phi}{\acontext}.
$$
\item
If $\exists!v\in\Val.\modellabel{\phi[a\ssm v]}{\acontext}=\tvT$ and $\Forall{v\in\Val}\modellabel{\phi[a\ssm v]}{\acontext}\in\{\tvT,\tvF\}$, then 
$$
\modellabel{\texiunique a.\phi}{\acontext}=\tvT=\modellabel{\texiaffine a.\phi}{\acontext}.
$$
\item
If $\exists!v\in\Val.\modellabel{\phi[a\ssm v]}{\acontext}=\tvT$ and $\Exists{v\in\Val}\modellabel{\phi[a\ssm v]}{\acontext}=\tvB$, then 
$$
\modellabel{\texiunique a.\phi}{\acontext}=\tvB=\modellabel{\texiaffine a.\phi}{\acontext}.
$$
\item\label{item.explain.existsunique.4}
If $\neg\Exists{v\in\Val}\modellabel{\phi[a\ssm v]}{\acontext}=\tvT$ and $\exists! v'\in\Val.\modellabel{\phi[a\ssm v]}{\acontext}=\tvB$, then 
$$
\modellabel{\texiunique a.\phi}{\acontext}=\tvB
= \modellabel{\texiaffine a.\phi}{\acontext} . %
$$
(We discuss this in Remark~\ref{rmrk.discuss.clause.4} below.)
\item\label{item.two.b}
If $\neg\Exists{v\in\Val}\modellabel{\phi[a\ssm v]}{\acontext}=\tvT$ and $\exists v,v'\in\Val.v\neq v'\land \modellabel{\phi[a\ssm v]}{\acontext}=\modellabel{\phi[a\ssm v']}{\acontext}=\tvB$, then 
$$
\modellabel{\texiunique a.\phi}{\acontext}
=\tvB=
\modellabel{\texiaffine a.\phi}{\acontext} .
$$
\item
If $\Forall{v\in\Val}\modellabel{\phi[a\ssm v]}{\acontext}=\tvF$, then
$$
\modellabel{\texiunique a.\phi}{\acontext}=\tvF
\quad\text{and}\quad
\modellabel{\texiaffine a.\phi}{\acontext}=\tvT.
$$
\end{enumerate*}
\end{rmrk}

\begin{rmrk}
\label{rmrk.discuss.texiunique}
We make two elementary observations on Remark~\ref{rmrk.unpack.texiunique}:
\begin{enumerate} 
\item
The equality $\teq$ in our logic is two-valued: $\model{v\teq v'}$ returns $\tvT$ or $\tvF$, and never $\tvB$.\footnote{In the terminology of Definition~\ref{defn.correct}(\ref{item.correct.phi}) below, equality predicates are always \emph{correct}.}
\item
Two notions of `unique existence' are in play in this discussion:
\begin{itemize*}
\item
`$\exists!v\in\Val$' refers to the actual unique existence quantifier in the real world (whatever that is).
This computes actual truth-values `true' and `false' and returns true when there actually exists a unique value.
\item
`$\texiunique a$' is part of the (derived) syntax of \QLogic, and is defined in Definition~\ref{defn.unique.exist}.
Its denotation computes elements of $\THREE=\{\tvF,\tvB,\tvT\}$.  
As is standard (though confusingly) we \emph{also} call elements of $\THREE$ truth-values; but they are truth-values of \QLogic, rather than actual truth-values in the real world (whatever that is).
\end{itemize*}
So: $\exists!$ and $\texiunique$ are both unique existence operators, but they live at different levels and they are not the same thing.\footnote{This kind of thing is far from unusual and appears also outside of papers with fancy logics.  For example in programming, a programming language may have multiple versions of the number $1$ (8-bit, floating point, 64-bit, and infinite precision) all of which are called `one', but all of which are different from one other, and also different from the actual $1\in\mathbb N$ (whatever that is).}
\end{enumerate}
\end{rmrk}


We will use unique/affine existence in our proofs via Lemma~\ref{lemm.unique.affine.existence}: 
\begin{lemm}
\label{lemm.unique.affine.existence}
Suppose $\acontext$ is a context and $\texi a.\phi$ is a closed predicate and $v,v'\in\Val$ are values.
Then:
\begin{enumerate*}
\item\label{item.unique.affine.existence.01}
$\acontext\ment\texiaffine a.\phi$ if and only if 
$\exists_{01} v\in\Val.\modellabel{\phi[a\ssm v]}{\acontext}=\tvT$.
\item\label{item.unique.affine.existence.11}
$\acontext\ment\texiunique a.\phi$ if and only if 
$\exists v\in\Val.\modellabel{\phi[a\ssm v]}{\acontext}\in\threeValid \land \exists_{01} v\in\Val.\modellabel{\phi[a\ssm v]}{\acontext}=\tvT$.
\item\label{item.unique.affine.existence.01.implies}
If $\acontext\ment\texiaffine a.\phi$ and $\acontext\ment\modT\phi[a\ssm v]$ and $\acontext\ment\modT\phi[a\ssm v']$, then $v=v'$.
\item\label{item.unique.affine.existence.11.implies}
If $\acontext\ment\texiunique a.\phi$ and $\acontext\ment\modT\phi[a\ssm v]$ and $\acontext\ment\modT\phi[a\ssm v']$, then $v=v'$.
\end{enumerate*}
\end{lemm}
\begin{proof}
We unpack Definition~\ref{defn.validity.judgement}(\ref{item.ctx.ment.phi}) %
and check the possibilities in Remark~\ref{rmrk.unpack.texiunique}.
\end{proof}


\subsubsection{Comments on an edge case}

Lemma~\ref{lemm.texiunique.all.b} notes a (perhaps surprising, at first glance) fact about our denotation, that 
if a unary predicate symbol $\tf P$ takes truth-value $\tvB$ when applied to all values (as would model a participant that has entirely crashed), then $\texiunique\tf P$ and $\texiaffine\tf P$ are valid --- even though $\tf P(v)$ is valid on every value.
This reflects the common side-condition in the literature `if the participant is correct, then \emph{stuff}', which is trivially true if the participant is \emph{not} correct: 
\begin{lemm}
\label{lemm.texiunique.all.b}
Suppose $\acontext$ is a context and $\tf P$ is a unary predicate symbol:\footnote{There is nothing mathematically special about using a \emph{unary} predicate symbol here, but in practice this is the case that we will use.}
Then
$$
\acontext\ment \tall\modB\tf P \limp (\acontext\ment \texiunique \tf P \land \acontext\ment \texiaffine \tf P).  
$$
The above statement uses Notation~\ref{nttn.hos}.
Unpacking this, we obtain:
$$
\acontext\ment \tall a.\modB\tf P(a) \limp (\acontext\ment \texiunique a.\tf P(a)\land \acontext\ment \texiaffine a.\tf P(a)) . 
$$  
\end{lemm}
\begin{proof}
Suppose $\acontext\ment \tall\modB\tf P$.
Unpacking Definition~\ref{defn.validity.judgement}(\ref{item.ctx.ment.phi}) this means $\modellabel{\tall\modB\tf P}{\acontext}\in\threeValid$, and by Figure~\ref{fig.3.derived} it follows that $\Forall{v\in\Val}\modellabel{\tf P(v)}{\acontext}=\tvB$.

By Remark~\ref{rmrk.unpack.texiunique}(\ref{item.two.b}) it follows that $\modellabel{\texiunique\tf P}{\acontext}=\tvB$.
The truth-value $\tvB$ is not true, but it is \emph{valid} by Definition~\ref{defn.validity.judgement}(\ref{item.ctx.ment.phi}), and thus $\acontext\ment\texiaffine\tf P$ and $\acontext\ment\texiunique\tf P$ as required. 
\end{proof}



\begin{rmrk}
\label{rmrk.discuss.clause.4}
Arguably, we would expect in clause~\ref{item.explain.existsunique.4} above 
(where $\neg\Exists{v\in\Val}\modellabel{\phi[a\ssm v]}{\acontext}=\tvT$ and $\exists! v'\in\Val.\modellabel{\phi[a\ssm v]}{\acontext}=\tvB$)
that $\modellabel{\texiaffine a.\phi}{\acontext}$ should equal $\tvT$, not $\tvB$. 
This is not hard to arrange, e.g. as
$$
\texiaffine a.\phi \tor \texiaffine a.\modT\phi .
$$
Either choice works with our later proofs; we never encounter an edge case where it matters.
\end{rmrk}



\subsection{Pointwise predicates (a well-behavedness property)}

For this subsection we fix a signature $\Sigma$ and a model $\mathcal M=(\Pnt,\opens,\Val)$.

\subsubsection{The definition, and basic properties}

\begin{rmrk}
\emph{Being pointwise} is a well-behavedness property on predicates, that their denotation in a context $\acontext=(n,p,O,\avaluation)$ %
does not depend on the open set parameter $O$ of the context.

Many important syntactic classes of predicates are naturally pointwise, including any predicate of the form $\tf P(v_1,\dots,v_n)$ (Lemma~\ref{lemm.atomic.O}).

Being pointwise is desirable for us because it simplifies denotation to holding `at a point';\footnote{Actually: at a stage, point, and valuation, but intuitively we tend to view the stage and valuation as fixed, and the point $p$ and open set $O$ as varying.} and also because several nice properties depend on this (see for example Proposition~\ref{prop.pointwise.dense.char}); and finally also because there is an easy syntactic condition (given in Lemma~\ref{lemm.atomic.O}(\ref{item.atomic.pointwise.2})) that guarantees that a predicate is pointwise.
\end{rmrk}
 

\begin{defn}[Pointwise predicates]
\label{defn.pointwise.meaning}
Suppose $\phi$ is a closed predicate.
\begin{enumerate*}
\item
If $n\in\Time$ and $\avaluation\in\Valuation(\Sigma,\Val)$ then call $\phi$ \deffont{pointwise in $(n,\avaluation)$} when
$$
\Forall{p\in\Pnt}\Forall{O,O'{\in}\opensne}\modellabel{\phi}{n,p,O,\avaluation}=\modellabel{\phi}{n,p,O',\avaluation} .
$$ 
In words, $\phi$ is pointwise in $(n,\avaluation)$ when at every point $p\in\Pnt$, the open set parameter $O\in\opensne$ does not affect the value of the denotation $\modellabel{\phi}{n,p,O,\avaluation}$. 
\item
Call $\phi$ \deffont{absolutely pointwise} when it is pointwise for every $n\in\Time$ and $\avaluation\in\Valuation(\Sigma,\Val)$. 
\end{enumerate*}
\end{defn}

\begin{defn}
\label{defn.pointwise.pred.nttn}
Suppose $\phi$ is closed, and $n\in\Time$ and $p\in\Pnt$ and $\avaluation$ is a valuation.
Then define $\modellabel{\phi}{n,p,\avaluation}$ by
$$
\modellabel{\phi}{n,p,\avaluation}=\bigwedge_{O\in\opensne}\modellabel{\phi}{n,p,O,\avaluation}.
$$
\end{defn}

In a manner of speaking we have already seen $\modellabel{\phi}{n,p,\avaluation}$ from Definition~\ref{defn.pointwise.pred.nttn} in our logic:
\begin{lemm}
\label{lemm.Coquorum.denote}
Suppose $n\in\Time$ and $p\in\Pnt$ and $\avaluation$ is a valuation.
Then
$$
\modellabel{\phi}{n,p,\avaluation}=\modellabel{\Coquorum\phi}{n,p,O,\avaluation}.
$$
\end{lemm}
\begin{proof}
Direct from the clause for $\Coquorum$ from Figure~\ref{fig.3.derived}.
\end{proof}

\begin{lemm}
\label{lemm.pointwise.pred}
Suppose $\phi$ is a closed predicate and $n\in\Time$ and $\avaluation$ is a valuation.
Then $\phi$ is pointwise in $(n,\avaluation)$ in the sense of Definition~\ref{defn.pointwise.meaning} 
if and only if
$$
\Forall{p\in\Pnt}\Forall{O\in\opensne}\modellabel{\phi}{n,p,O,\avaluation}=\modellabel{\phi}{n,p,\avaluation}.
$$
\end{lemm}
\begin{proof}
By a routine calculation from Definitions~\ref{defn.pointwise.meaning} and~\ref{defn.pointwise.pred.nttn}.
\end{proof}

\begin{corr}
\label{corr.pointwise.pred}
Suppose $\phi$ is a closed predicate and $n\in\Time$ and $p\in\Pnt$ and $\avaluation$ is a valuation.
Suppose $\phi$ is pointwise in $(n,\avaluation)$.
Then
$$
\Forall{O\in\opensne}n,p,O\mentval\phi\liff n,p\mentval\phi. 
$$
\end{corr}
\begin{proof}
Routine from Definition~\ref{defn.validity.judgement} and Lemma~\ref{lemm.pointwise.pred}. %
\end{proof}


Being pointwise is particularly nice because of two simple syntactic criteria that guarantee it:
\begin{lemm}[Conditions for being pointwise]
\label{lemm.atomic.O}
Suppose $\phi$ is a predicate.
Then:
\begin{enumerate*}
\item\label{item.atomic.pointwise}
If $\phi$ is atomic (Notation~\ref{nttn.logic.terminology}(\ref{item.atomic.predicate}); having the form $\tf P(v_1,\dots,v_n)$ or $\tneg\tf P(v_1,\dots,v_n)$) then it is absolutely pointwise (Definition~\ref{defn.pointwise.meaning}).
\item\label{item.atomic.pointwise.2}
If $\phi$ does not mention $\someone$ or $\everyone$ except within the scope of a $\Quorum$ or $\Coquorum$ modality, then $\phi$ is absolutely pointwise.\footnote{In particular this holds if $\Quorum$ and $\Coquorum$ and only ever mentioned in combination with $\someone$ and $\everyone$, as in the compound modalities $\someoneAll$, $\everyoneAll$, $\QuorumBox$, and $\CoquorumDiamond$.}
\end{enumerate*} 
\end{lemm}
\begin{proof}
For part~1, we just unpack the definition of $\modellabel{\phi}{}$ in Figure~\ref{fig.3.phi.f} and note that the `$\avaluation(\Pnt)_{n,p}$' used in the clause for a predicate symbol depends just on the values of $n$, $p$, and $\avaluation$, but not on the open sets $O$ and $O'$. 

Part~2 follows by a routine induction on predicate syntax, using part~1 as the base case.
\end{proof}


\subsubsection{Pointwise, as a modality in our logic}

In fact, we can express the property of being pointwise within our logic:
\begin{defn}
\label{defn.pointwise.pred}
Suppose $\phi$ is a predicate.
Then define a predicate $\pointwise{\phi}$ by
$$
\pointwise{\phi} \ =\ \everyoneAll(\Quorum\phi \equiv \Coquorum\phi) .
$$
\end{defn}


\begin{lemm}
\label{lemm.pointwise.heart}
Suppose $\phi$ is a closed predicate and $n\in\Time$ and $p\in\Pnt$ and $\avaluation$ is a valuation.
Then 
$$
n,p\mentval\Quorum\phi \equiv \Coquorum\phi
\quad\text{if and only if}\quad
\bigvee_{O'\in\opensne}\modellabel{\phi}{n,p,O',\avaluation}=\bigwedge_{O'\in\opensne}\modellabel{\phi}{n,p,O',\avaluation}.
$$
(The predicate $\Quorum\phi \equiv \Coquorum\phi$ itself satisfies the conditions for being pointwise from Lemma~\ref{lemm.atomic.O}(\ref{item.atomic.pointwise.2}), so we can sensibly write $n,p\mentval \Quorum\phi \equiv \Coquorum\phi$ above.)
\end{lemm}
\begin{proof}
From Lemma~\ref{lemm.equiv}(\ref{item.equiv.1}) and the clauses for $\Quorum$ in Figure~\ref{fig.3.phi.f}, and $\Coquorum$ from Figure~\ref{fig.3.derived} (or Lemma~\ref{lemm.Coquorum.denote}). 
\end{proof}

\begin{lemm}
Suppose $\phi$ is a closed predicate and $n\in\Time$ and $p\in\Pnt$ and $\avaluation$ is a valuation.
Then $\phi$ is pointwise in $(n,\avaluation)$ if and only if $n,p\mentval\pointwise{\phi}$. 
\end{lemm}
\begin{proof}
By routine calculations using Lemma~\ref{lemm.pointwise.heart}, and the clause for $\everyoneAll$ in Figure~\ref{fig.3.derived}, noting the elementary fact that for every $O\in\opensne$,
$$
\bigwedge_{O'\in\opensne}\modellabel{\phi}{n,p,O',\avaluation}
\leq
\modellabel{\phi}{n,p,O,\avaluation}
\leq
\bigvee_{O'\in\opensne}\modellabel{\phi}{n,p,O',\avaluation} .
$$
\end{proof}



\subsubsection{An interlude: dense sets, and sets with a nonempty open interior}


We recall a pair of standard topological definitions which will turn up repeatedly later, most notably in the context of pointwise predicates as Proposition~\ref{prop.pointwise.dense.char} immediately below, and again later in Lemma~\ref{lemm.is}:
\begin{defn}[Dense subset]
\label{defn.dense.subset}
Suppose $(\Pnt,\opens)$ is a semitopology and suppose $P\subseteq\Pnt$ is a (not necessarily open) set of points.
\begin{enumerate*}
\item\label{item.dense.subset}
Say that $P$ is \deffont{dense} when $\Forall{O{\in}\opensne}O\intersectswith P$.
\item\label{item.noi.subset}
Say that $P$ \deffont{has a nonempty open interior} when $\Exists{O{\in}\opensne}O\subseteq P$.
\item\label{item.dense.noi}
Write $\dense(P)$ when $P$ is dense, and $\noi(P)$ when $P$ has a nonempty open interior. 
\end{enumerate*} 
\end{defn}



\begin{lemm}
\label{lemm.dense.up-closed}
Suppose $(\Pnt,\opens)$ is a semitopology and $P,P'\subseteq\Pnt$.
Then:
\begin{enumerate*}
\item\label{item.dense.up-closed} 
If $P\subseteq P'$ then $\dense(P)\limp \dense(P')$ (being dense is \emph{up-closed}).
\item\label{item.noi.up-closed} 
If $P\subseteq P'$ then $\noi(P)\limp \noi(P')$ (having a non-empty open intersection is \emph{up-closed}).
\item\label{item.noi.noi} 
If $P\in\opensne$ then $\noi(P)$ (a nonempty open set has a nonempty open interior).
\end{enumerate*}
\end{lemm}
\begin{proof}
By routine sets calculations from Definition~\ref{defn.dense.subset}.
\end{proof}

\begin{lemm}
Suppose $(\Pnt,\opens)$ is a semitopology and $P\subseteq\Pnt$.
Then
$$
\dense(P) \liff \neg\noi(\Pnt\setminus P) 
\quad\text{and}\quad 
\noi(P) \liff \neg\dense(\Pnt\setminus P) . 
$$
In this sense, $\noi$ and $\dense$ form a \emph{de Morgan dual pair} of operators --- just as $\forall=\neg\exists\neg$ (as noted in Figure~\ref{fig.easy.equivalences}) and $\exists=\tneg\forall\tneg$.
\end{lemm}
\begin{proof}
We sketch the (easy) argument:
\begin{itemize}
\item
Suppose $P$ is not dense; by Definition~\ref{defn.dense.subset} this means that there exists $O\in\opensne$ such that $P\cap O=\varnothing$, thus $O\subseteq\Pnt\setminus P$, thus $\Pnt\setminus P$ has a nonempty open interior.
\item
Suppose $\Pnt\setminus P$ has a nonempty open interior $\varnothing\neq O\subseteq\Pnt\setminus P$.
Then clearly $P$ is not dense.
\qedhere\end{itemize}
\end{proof}

More on how `being dense' works in semitopologies is in~\cite[Chapter~11]{gabbay:semdca}.
The notion behaves somewhat differently than it does in topologies, but not in ways that will affect us here.

\subsubsection{Application to pointwise predicates}


Technically, Proposition~\ref{prop.pointwise.dense.char} is a corollary of Lemma~\ref{lemm.dense.char} and Corollary~\ref{corr.pointwise.pred}.
The result links the four logical modalities $\someoneAll$, $\everyoneAll$, $\QuorumBox$, and $\CoquorumDiamond$ to elementary topological properties of the sets of points at which these modalities are valid --- being nonempty, being full, having a nonempty open interior, and being dense.
We will use Proposition~\ref{prop.pointwise.dense.char}, both directly and via other lemmas that depend on it.

The most important use-case of Proposition~\ref{prop.pointwise.dense.char} will be when $\phi$ is atomic (Notation~\ref{nttn.logic.terminology}(\ref{item.atomic.predicate}) --- having the form 
$$
\tf P(v_1,\dots,v_n)
\quad\text{or}\quad
\tneg\tf P(v_1,\dots,v_n)
$$
--- since then by Lemma~\ref{lemm.atomic.O} $\phi$ is automatically pointwise:
\begin{prop}
\label{prop.pointwise.dense.char}
Suppose $\phi$ is a closed pointwise (Definition~\ref{defn.pointwise.meaning}) predicate and $n\in\Time$ and $\avaluation$ is a valuation.
Then: 
\begin{enumerate*}
\item\label{item.dense.char.pointwise.someoneAll}
$n\mentval\someoneAll\phi$ if and only if $\varnothing\neq\{p\in\Pnt \mid n,p\mentval\phi\}$. 
\item\label{item.dense.char.pointwise.everyoneAll}
$n\mentval\everyoneAll\phi$ if and only if $\Pnt=\{p\in\Pnt \mid n,p\mentval\phi\}$. 
\item\label{item.dense.char.pointwise.QuorumBox}
$n\mentval\QuorumBox\phi$ if and only if $\noi(\{p\in\Pnt \mid n,p\mentval\phi\})$.
\item\label{item.dense.char.pointwise.CoquorumDiamond}
$n\mentval\CoquorumDiamond\phi$ if and only if $\dense(\{p\in\Pnt \mid n,p\mentval\phi\})$.
\end{enumerate*} 
Recall that:
\begin{itemize*}
\item
$\dense$ and $\noi$ are from Definition~\ref{defn.dense.subset}.
\item
The notations $n,p\mentval\phi$ and $n\mentval\phi$ are from Definition~\ref{defn.validity.judgement}.
\end{itemize*}
\end{prop}
\begin{proof}
If $p$ is pointwise then by Corollary~\ref{corr.pointwise.pred} $n,p,O\mentval\phi$ if and only if $n,p\mentval\phi$, for any $O\in\opensne$, and we just simplify parts~\ref{item.dense.char.someoneAll} to~\ref{item.dense.char.CoquorumDiamond} of Lemma~\ref{lemm.dense.char}:
\begin{enumerate} 
\item
We reason as follows:
$$
\begin{array}{r@{\ }l@{\quad}l}
n\mentval\someoneAll\phi
\liff&
\Exists{O\in\opensne}O\between\{p\in\Pnt \mid n,p\mentval\phi\}
&\text{Lemma~\ref{lemm.dense.char}(\ref{item.dense.char.someoneAll})} 
\\
\liff&
\Pnt\between \{p\in\Pnt \mid n,p\mentval\phi\}
&\text{Fact}
\\
\liff&
\{p\in\Pnt \mid n,p\mentval\phi\}\neq\varnothing
&\text{Fact}
\end{array}
$$
\item
We reason as follows:
$$
\begin{array}{r@{\ }l@{\quad}l}
n\mentval\everyoneAll\phi
\liff&
\Forall{O\in\opensne}O\subseteq\{p\in\Pnt \mid n,p\mentval\phi\}
&\text{Lemma~\ref{lemm.dense.char}(\ref{item.dense.char.everyoneAll})} 
\\
\liff&
\Pnt\subseteq\{p\in\Pnt \mid n,p\mentval\phi\}
&\text{Fact}
\\
\liff&
\{p\in\Pnt \mid n,p\mentval\phi\}=\Pnt
&\text{Fact}
\end{array}
$$
\item
We reason as follows:
$$
\begin{array}{r@{\ }l@{\quad}l}
n\mentval\QuorumBox\phi
\liff&
\Exists{O\in\opensne}O\subseteq\{p\in\Pnt \mid n,p\mentval\phi\}
&\text{Lemma~\ref{lemm.dense.char}(\ref{item.dense.char.QuorumBox})} 
\\
\liff&
\noi(\{p\in\Pnt \mid n,p\mentval\phi\})
&\text{Definition~\ref{defn.dense.subset}}
\end{array}
$$
\item
We reason as follows:
$$
\begin{array}{r@{\ }l@{\quad}l}
n\mentval\CoquorumDiamond\phi
\liff&
\Forall{O\in\opensne}O\between\{p\in\Pnt \mid n,p\mentval\phi\}
&\text{Lemma~\ref{lemm.dense.char}(\ref{item.dense.char.CoquorumDiamond})} 
\\
\liff&
\dense(\{p\in\Pnt \mid n,p\mentval\phi\})
&\text{Definition~\ref{defn.dense.subset}}
\end{array}
$$
\qedhere\end{enumerate}
\end{proof}
See also Lemma~\ref{lemm.allbut.concretely}, which instantiates Proposition~\ref{prop.pointwise.dense.char}(\ref{item.dense.char.pointwise.QuorumBox}\&\ref{item.dense.char.pointwise.CoquorumDiamond}) to a special case
familiar from the literature on distributed systems.


\subsection{Correct predicates (another well-behavedness property)}

\begin{rmrk}
We defined $\threeCorrect=\{\tvT,\tvF\}$ in Definition~\ref{defn.tv.ment}(\ref{item.tv.ment.correct}) and a corresponding modality $\modTF$ in Figure~\ref{fig.3} (which returns $\tvT$ when its argument is in $\threeCorrect$ and $\tvF$ otherwise), and we studied some of their properties, most notably in Lemma~\ref{lemm.TF.correct}.
 
We now extend this to notions of \emph{correctness} on predicates $\phi$ and predicate symbols $\tf P$.

Correctness turns out to be highly relevant to how we will use our logic.
Partly this is because when a predicate is correct, we are returned to the familiar two-valued world of Boolean logic --- but more subtly, correctness is interesting because of how this concept interacts with the predicate and modal parts of our logic; in certain circumstances we can use a combination of correctness \emph{somewhere} and validity \emph{everywhere} to deduce information about truth (see Remark~\ref{rmrk.why.exciting}) and this turns out to be a good logical model for how reasoning on distributed algorithms is often done in practice.

We start with the definition and some basic lemmas, and develop from there:
\end{rmrk}


\begin{defn}[Correctness]
\label{defn.correct}
Suppose $\phi$ is a predicate and $a\in\tf{VarSymb}$ is a variable symbol. 
Then:
\begin{enumerate*}
\item\label{item.correct.phi}
Consistent with the definition of $\threeCorrect=\{\tvT,\tvF\}$ from Definition~\ref{defn.tv.ment}(\ref{item.tv.ment.correct}),
we define $\correct{\phi}$, read \deffont{$\phi$ is correct}, by
$$
\correct{\phi}=\modTF\phi .
$$
Note that $\phi$ being correct does not mean that $\phi$ returns $\tvT$; it means that $\phi$ returns $\tvT$ or $\tvF$ and in particular it does \emph{not} return $\tvB$.
\item\label{item.correct}\label{item.crashed}
Suppose $\tf P\in\PredSymb$ is a unary predicate symbol.
Define a predicate $\correct{\tf P}$ by 
$$
\correct{\tf P}\ =\ \tall a.\correct{\tf P(a)}\ \stackrel{N\ref{nttn.hos}}{=}\ \tall\modTF\tf P
 .
$$
\item
If $\tf P_1,\dots,\tf P_n$ are unary predicate symbols then we may write 
$$
\correct{\tf P_1,\dots,\tf P_n}
\quad\text{for}\quad 
\correct{\tf P_1}\tand\dots\tand\correct{\tf P_n}.
$$
(See for example axiom \rulefont{LdrCorrect} in Figure~\ref{fig.logical.paxos}.) 
\item\label{item.P-correct}
Suppose $n\in\Time$ and $p\in\Pnt$ and $\avaluation$ is a valuation, and suppose $\tf P$ is a predicate symbol.
Then when $n,p\mentval\correct{\tf P}$ we may call $p$ \deffont{$\tf P$-correct}. 
\end{enumerate*}
\end{defn}

\begin{lemm}
\label{lemm.correct.valid}
Suppose $\acontext$ is a context and $\phi$ is a closed predicate.
Then 
$$
\acontext\ment\phi
\quad\text{if and only if}\quad
\modellabel{\phi}{\acontext}\in\threeCorrect.
$$
\end{lemm}
\begin{proof}
This just re-states Lemma~\ref{lemm.tv.ment.TF}(\ref{item.tv.ment.TF.2}).
\end{proof}


\begin{rmrk}
\label{rmrk.tf.synonyms}
Definition~\ref{defn.correct}(\ref{item.correct.phi}) gives us two synonyms for the modality expressing that the truth-value returned by a predicate $\phi$ is either $\tvT$ or $\tvF$ (not $\tvB$): $\modTF$ and $\tf{correct}$.

Although these are synonymous, there is a difference in emphasis:
\begin{itemize*}
\item
$\modTF$ is a logical modality and just means `true or false'.
\item
The literature on distributed systems uses a term `correctness', which indicates that a process succeeds and returns a correct value. 
In our logical model this corresponds to returning $\tvT$ or $\tvF$ --- and thus not returning $\tvB$. 
So when we write $\tf{correct}$, the emphasis is on `not $\tvB$' with emphasis on a process `not failing'.
\end{itemize*}
We use whichever notation seems most natural in context. 
For example, in Figure~\ref{fig.logical.paxos}:
\begin{enumerate*}
\item
In \rulefont{LdrCorrect} we write $\correct{\tf{propose},\tf{write},\tf{decide}}$ because intuitively we are asserting that at a particular time the leader \emph{does not crash} when proposing, writing, or deciding.
\item
In \rulefont{LdrExist} we write $\modTF\tf{leader}$, because $\tf{leader}$ is a predicate that identifies whether a participant is the leader at a particular time, and this is either true or false.\footnote{In practice determining who the leader is require a participant to compute a function, but the computation is not distributed, in the sense that prior agreement amongst participants is assumed on what the function is and what seed data it uses, and each participant computes it locally to get (the same) definitive yes/no answer.}
\end{enumerate*} 
If the reader prefers, they can safely read $\modTF$ every time they see $\tf{correct}$, and they will be precisely correct (pun intended). 
\end{rmrk}



Many predicates of interest are correct, including correctness itself, our truth-modalities, and equality and equivalence:
\begin{lemm}
\label{lemm.equality.correct}
Suppose $\acontext$ is a context and $\phi$ and $\phi'$ are closed predicates and $v,v'\in\Val$.
Then:
\begin{enumerate*}
\item
$\modellabel{\correct{\correct{\phi}}}{\acontext}=\tvT$.
\quad We can (perhaps cheekily) read this as `correctness is correct'.
\item
$\modellabel{\correct{\tf Q\phi}}{\acontext}=\tvT$ for $\tf Q\in\{\modT,\modF,\modB,\modTB,\modTF,\modFB\}$. 
\item\label{item.equality.correct}
$\modellabel{\correct{v\teq v'}}{\acontext}=\tvT$.
In words: equality is correct.
\item\label{item.equivalence.correct}
$\modellabel{\correct{\phi\equiv \phi'}}{\acontext}=\tvT$.
In words: equivalence is correct.
\end{enumerate*}
\end{lemm}
\begin{proof}
By routine checking from Figures~\ref{fig.3}, \ref{fig.3.phi.f}, \ref{fig.3.derived}, and (for $\equiv$)~\ref{fig.3.iff}.
\end{proof}

Lemma~\ref{lemm.correct.P.Pv} and Corollary~\ref{corr.modTF.texi} are straightforward, but we spell it out:
\begin{lemm}
\label{lemm.correct.P.Pv}
Suppose $\tf P\in\PredSymb$ is a unary predicate symbol and $v\in\Val$ and $\acontext$ is a context.
Then 
$$
\modellabel{\correct{\tf P}}{\acontext} = \bigwedge_{v{\in}\Val}\modellabel{\correct{\tf P(v)}}{\acontext}
\quad\text{and}\quad 
\modellabel{\correct{\tf P}}{\acontext}\leq \modellabel{\correct{\tf P(v)}}{\acontext} .
$$
\end{lemm}
\begin{proof}
By Definition~\ref{defn.correct}(\ref{item.correct}) and Figure~\ref{fig.3.derived} 
$$
\modellabel{\correct{\tf P}}{\acontext} 
= 
\modellabel{\tall a.\correct{\tf P(a)}}{\acontext} 
= 
\bigwedge_{v{\in}\Val} \modellabel{\correct{\tf P(v)}}{\acontext} 
\leq 
\modellabel{\correct{\tf P(v)}}{\acontext} .
\qedhere$$
\end{proof}


\begin{corr}
\label{corr.modTF.texi}
Suppose $\acontext$ is a context and $\tf P\in\PredSymb$ is a unary predicate symbol and $a\in\tf{VarSymb}$ is a variable symbol. 
Then 
$$
\acontext\ment\correct{\tf P}
\quad\text{implies}\quad
\acontext\ment\correct{\texi\tf P}\ \land\ \acontext\ment\correct{\tall\tf P}.
$$
\end{corr}
\begin{proof}
Suppose $\acontext\ment\correct{\tf P}$.
By Lemma~\ref{lemm.correct.P.Pv} we see that $\Forall{v\in\Val}\acontext\ment\correct{\tf P(v)}$.
By Lemma~\ref{lemm.correct.valid} $\Forall{v\in\Val}\modellabel{\tf P(v)}{\acontext}\in\threeCorrect$.
Now the clauses for $\texi$ in Figure~\ref{fig.3.phi.f} and $\tall$ in Figure~\ref{fig.3.derived} take a supremum and infimum respectively of $\{\modellabel{\tf P(v)}{\acontext} \mid v\in\Val\}\subseteq\threeCorrect=\{\tvT,\tvF\}$, and this supremum or infimum is also clearly in $\threeCorrect=\{\tvT,\tvF\}$.
We just use Lemma~\ref{lemm.correct.valid}. 
\end{proof}


Proposition~\ref{prop.tv.ment.TF.model}(\ref{item.tv.ment.TF.phi.lr}) is a wrapper for Lemma~\ref{lemm.TF.correct}.
This deceptively simple equivalence turns out to be a key way that correctness can be usefully applied:
\begin{prop}
\label{prop.tv.ment.TF.model}
Suppose $\phi$ is a closed predicate and $\acontext$ is a context and $v\in\Val$.
Then:
\begin{enumerate*}
\item\label{item.tv.ment.TF.phi.lr}
$\acontext\ment\correct{\phi}$ if and only if $\acontext\ment\phi\equiv\modT\phi$. 

In words: correct $\phi$ are characterised precisely by validity being equivalent to truth.
\item\label{item.tv.ment.TF.P}
$\acontext\ment\correct{\tf P}$ if and only if $\Forall{v\in\Val}(\acontext\ment\tf P(v)\equiv\modT\tf P(v))$. 
\item\label{item.modTF.texi.P.valid}
$\acontext\ment\correct{\tf P}$ implies $\acontext\ment\texi\tf P\equiv \modT\texi\tf P$ and $\acontext\ment\tall\tf P\equiv\modT\tall\tf P$. 
\end{enumerate*}
\end{prop}
\begin{proof}
We consider each part in turn:
\begin{enumerate}
\item
Suppose $\acontext\ment\correct{\phi}$.
By Lemma~\ref{lemm.TF.correct} $\modT\modellabel{\phi}{\acontext}=\modellabel{\phi}{\acontext}$.
By Lemma~\ref{lemm.equiv}(\ref{item.equiv.1}) and the clause for $\modT$ in Figure~\ref{fig.3.phi.f}, $\acontext\ment \phi\equiv\modT\phi$.
The reverse implication just reverses the same reasoning.
\item
We combine Lemma~\ref{lemm.correct.P.Pv} with part~\ref{item.tv.ment.TF.phi.lr} of this result for $\phi=\tf P(v)$ for every $v\in\Val$.
\item
From Corollary~\ref{corr.modTF.texi} and Proposition~\ref{prop.tv.ment.TF.model}(\ref{item.tv.ment.TF.phi.lr}).
\qedhere\end{enumerate}
\end{proof}




Proposition~\ref{prop.correct.cobox.boxT} explores how correctness interacts with our compound modalities from Subsection~\ref{subsect.compound.modalities}.
\begin{prop}
\label{prop.correct.cobox.boxT}
Suppose $\phi$ is a closed predicate and $\tf P\in\PredSymb$ is a unary predicate symbol. 
Suppose $n\in\Time$ and $\avaluation$ is a valuation.
Then:
\begin{enumerate*}
\item\label{item.correct.cobox.boxT.1}
$n\mentval\QuorumBox\correct{\phi}\tand\everyoneAll\phi$ implies $n\mentval \modT\QuorumBox\phi$.\footnote{Alert: note that we write $\everyoneAll\phi$ in the right-hand conjunct, not $\CoquorumDiamond\phi$.}
\item\label{item.correct.cobox.boxT.2}
$n\mentval\QuorumBox\correct{\tf P}\tand\everyoneAll\texi\tf P$ implies 
$n\mentval\modT\QuorumBox\texi\tf P$.
\end{enumerate*}
\end{prop}
\begin{proof}
We consider each part in turn:
\begin{enumerate}
\item
Suppose $n\mentval\QuorumBox\correct{\phi}\tand\everyoneAll\phi$.
By Proposition~\ref{prop.someone.implies.someoneAll}(\ref{item.someone.implies.someoneAll.4b})
$n\mentval\QuorumBox(\correct{\phi}\tand\phi)$.
By Proposition~\ref{prop.tv.ment.TF.model}(\ref{item.tv.ment.TF.phi.lr}) and Lemma~\ref{lemm.someone.commutation} (for $\QuorumBox$)
we have $n\mentval\modT\QuorumBox\phi$ as required.
\item
We combine part~\ref{item.correct.cobox.boxT.1} of this result with Corollary~\ref{corr.modTF.texi}.
\qedhere\end{enumerate}
\end{proof}


\begin{rmrk}[Why this is interesting]
\label{rmrk.why.exciting}
The statement of Proposition~\ref{prop.correct.cobox.boxT} is quite heavy on the symbols.
For clarity, let us re-read it, but let us replace $\ment\QuorumBox\tf{correct}$ with the phrase `mostly correct', and $\ment\everyoneAll$ with the phrase `everywhere valid', and $\texi$ with the phrase `for some value'.
Then we get:
\begin{enumerate*}
\item
If $\phi$ is \emph{mostly correct} and \emph{everywhere valid}, then it is \emph{mostly true}.
\item
If $\tf P$ is \emph{mostly correct} and \emph{everywhere valid} for some value, then it is \emph{mostly true} for some value.
\end{enumerate*}
So what we see is that we can promote weak statements about validity everywhere, with stronger but more limited statements about correctness, to obtain limited statements about truth.

The reader familiar with how distributed algorithms work may recognise this style of reasoning.
So: here it is beginning to appear in our logic.
\end{rmrk}

\subsection{$\mrup{a}{\phi}{v}$: a most recent $v$ that makes $\phi[a\ssm v]$ true}
\label{subsect.mrup}

\begin{rmrk}
\label{rmrk.explain.mrup}
We see from Figure~\ref{fig.3.derived} and Remark~\ref{rmrk.comments.on.semantics}(\ref{item.mru.program}) that $\mrup{a}{\phi}{v}$ looks back into the past, searching for a most recent time when $\model{\texi a.\phi}=\tvT$, and checks if at that most recent time, $\model{\phi[a\ssm v]}=\tvT$.
$\mrup{a}{\phi}{v}$ returns $\tvT$ if this search succeeds, and $\tvF$ if this fails.
It never returns $\tvB$.

In words: $\mrup{a}{\phi}{v}$ checks whether $v$ is a value that has most recently made $\phi$ true.

We need $\mrup{a}{\phi}{v}$ for our axiomatisation of Paxos in Figure~\ref{fig.logical.paxos}, where 
\begin{itemize*}
\item
in \rulefont{PaxSend?} we write $\mru{\tf{accept}}{v}$\ (short for $(\mrup{a}{\tf{accept}(a))}{v}$, as per Figure~\ref{fig.hos}), and 
\item
in \rulefont{PaxWrite?} we write $\mru{(\someone\tf{accept})}{v}$\ (short for $\mrup{a}{(\someone\tf{accept}(a))}{v}$).
\end{itemize*} 
We read $\mrup{a}{\phi}{v}$ as 
\begin{quote}
\deffont{$v$ is a value that most recently makes $\phi[a\ssm v]$ true} .
\end{quote}
In the common case that $\phi=\tf P(a)$ for some unary predicate symbol $\tf P$, then we might just say
\begin{quote}
\deffont{$v$ is a value that most recently makes $\tf P$ true}.
\end{quote} 
In this Subsection we study the properties of this construct in more detail.
\end{rmrk}



\begin{lemm}
\label{lemm.mru.unambivalent}
Suppose $\texi a.\phi$ is a closed predicate and $v\in\Val$ and $\acontext$ is a context.
Then:
\begin{enumerate*}
\item
$\acontext\ment\correct{\mrup{a}{\phi}{v}}$.
\item
$\acontext\ment\mrup{a}{\phi}{v}\equiv\modT\mrup{a}{\phi}{v}$.
\end{enumerate*}
\end{lemm}
\begin{proof}
We consider each part in turn:
\begin{enumerate}
\item
By Definition~\ref{defn.correct}(\ref{item.correct.phi}) $\correct{\mrup{a}{\phi}{v}}=\modTF\mrup{a}{\phi}{v}$, and this means that $\mrup{a}{\phi}{v}$ always returns $\tvT$ or $\tvF$ (but not $\tvB$).
We just check the relevant clause in Figure~\ref{fig.3.derived} and see that this is so.
\item
We just use part~1 of this result and Proposition~\ref{prop.tv.ment.TF.model}(\ref{item.tv.ment.TF.phi.lr}).
\qedhere\end{enumerate}
\end{proof}


\begin{rmrk}
\label{rmrk.explain.mru}
\leavevmode
\begin{enumerate}
\item\label{item.explain.mru.T}
Lemma~\ref{lemm.mru.unambivalent}(2) %
notes that $\mrup{a}{\phi}{v}$ and $\modT(\mrup{a}{\phi}{v})$ are extensionally equivalent (Definition~\ref{defn.phi.equivalent}).

We may write one or the other, depending on what seems most clear in a proof.
In particular, if we are using \strongmodusponens then we might retain a $\modT$ modality for clarity.
\item 
We try to head off an easy error: note that $\mrup{a}{\phi}{v}$ is not equivalent to $\modT\recent\phi[a\ssm v]$, nor the (by Lemma~\ref{lemm.someone.commutation}) equivalent form $\recent\modT\phi[a\ssm v]$.
\begin{itemize*}
\item
$\modT\recent\phi[a\ssm v]$ returns $\tvT$ when $\phi[a\ssm v]$ has returned $\tvT$ at some point in the past.
\item
$\mrup{a}{\phi}{v}$ (`$v$ is a value that most recently makes $\phi[a\ssm v]$ true') returns $\tvT$ when $\phi[a\ssm v]$ has returned $\tvT$ in \emph{the most recent stage} such that $\texi a.\phi$ returned $\tvT$ (i.e. such that $\phi[a\ssm v']$ returned $\tvT$ for \emph{any} $v'$).
\end{itemize*}
These are not the same thing.
\end{enumerate}
\end{rmrk}

Remark~\ref{rmrk.explain.mru} notes that $\mrup{a}{\phi}{v}$ is different from $\modT\recent\phi[a\ssm v]$.
That said, there is a close link between them, which will turn out to be important for understanding our axiomatisation of Paxos (see Lemma~\ref{lemm.evaporates.simplify} and surrounding discussion): 
\begin{lemm}
\label{lemm.evaporates}
Suppose $\texi a.\phi$ is a closed predicate and $\acontext=(n,p,O,\avaluation)$ is a context. 
Then
$$
\acontext\ment(\texi v.\mrup{a}{\phi}{v})\equiv \modT\recent\texi a.\phi.
$$
\end{lemm}
\begin{proof}
First, note from Lemma~\ref{lemm.someone.commutation} that $\ment\modT\recent\texi a.\phi\equiv \recent\texi a.\modT\phi$.
And, note that $\texi a.\phi$ and $\texi v.\phi[a\ssm v]$ denote the same syntax, as per Remark~\ref{rmrk.pedant.point}(\ref{item.pedant.texi.v}).

We unpack the clause for $\modellabel{\mrup{a}{\phi}{v}}{\acontext}$ in Figure~\ref{fig.3.derived} and note that it returns $\tvT$ when there exist $v\in\Val$ and $0\leq n'<n$ such that $n',p,O\mentval \modT\phi[a\ssm v]$, and $\tvF$ otherwise.
But this is also what $\modellabel{\recent\texi v.\modT(\phi[a\ssm v])}{\acontext}$ does, so by Lemma~\ref{lemm.equiv}(\ref{item.equiv.1}) we are done.
\end{proof}

Lemma~\ref{lemm.mru.today.yesterday} connects $\mrup{a}{\phi}{v}$ to the fact that our notion of time is discrete and so has a well-defined notion of a `tomorrow' (i.e. a `next timestep'):
\begin{lemm}
\label{lemm.mru.today.yesterday}
Suppose $\texi a.\phi$ is a closed predicate and $v\in\Val$ and $\acontext$ is a context. 
.
Then
$$
\acontext\ment \tomorrow(\mrup{a}{\phi}{v}) \equiv (\modT\phi[a\ssm v] \tor ((\tneg\modT\texi a.\phi) \tand \mrup{a}{\phi}{v})).
$$
\end{lemm}
\begin{proof}
It is a fact of Figure~\ref{fig.3.derived} that $\acontext\ment\tomorrow\yesterday\psi\equiv\psi$ for any closed $\psi$.
The result follows from this observation and from the clause for $\mrup{a}{\phi}{v}$ in Figure~\ref{fig.3.derived}. 
\end{proof}



\begin{corr}
\label{corr.mru.today.to.yesterday}
Suppose $\texi a.\phi$ is a closed predicate and $v\in\Val$ and $\acontext$ is a context. 
Then:
\begin{enumerate*}
\item\label{item.mru.today.implies.tomorrow}
$\acontext\ment \phi[a\ssm v] \timpc \tomorrow(\mrup{a}{\phi}{v})$.
\item\label{item.mru.tomorrow.options}
$\acontext\ment \tomorrow(\mrup{a}{\phi}{v}) \timpc (\phi[a\ssm v]\tor \mrup{a}{\phi}{v})$. 
\end{enumerate*}
\end{corr}
\begin{proof}
Elementary from Lemma~\ref{lemm.mru.today.yesterday} using \strongmodusponens.
\end{proof}


\begin{lemm}
\label{lemm.mru.to.recent}
Suppose $\texi a.\phi$ is a closed predicate %
and $v,v'\in\Val$ and $\acontext$ is a context. 
Then:
\begin{enumerate*}
\item\label{item.mru.to.recent.2}
$\acontext\ment(\mrup{a}{\phi}{v}\tand\mrup{a}{\phi}{v'}) \timpc \recent(\phi[a\ssm v]\tand\phi[a\ssm v'])$.
\item\label{item.mru.to.recent.1}
$\acontext\ment(\mrup{a}{\phi}{v}) \timpc \recent(\phi[a\ssm v])$.

If $\phi=\modality{\tf Q}\tf P(a)$ for a unary predicate symbol $\tf P\in\PredSymb$ and some prefix of modalities $\modality{\tf Q}$, then we can elide the variable symbol $a$ as per Notation~\ref{nttn.hos} and Figure~\ref{fig.hos}, and this becomes
$$
\acontext\ment\mru{(\modality{\tf Q}\tf P)}{v} \timpc \recent\modality{\tf Q}\tf P(v).
$$
\end{enumerate*}
\end{lemm}
\begin{proof}
We consider each part in turn:
\begin{enumerate}
\item
Unpacking the clause for $\mruarrow$ in Figure~\ref{fig.3.derived}, if $v$ and $v'$ are values that most recently make $\phi[a\ssm\text{-}]$ be $\tvT$, then certainly $\recent(\phi[a\ssm v]\tand\phi[a\ssm v'])$.
\item
Follows from part~\ref{item.mru.to.recent.2} as the special case that $v'=v$.
\qedhere\end{enumerate}
\end{proof}


\section{\lowercase{n}-twined semitopologies}
\label{sect.n-twined}



\subsection{Antiseparation properties}

We are particularly interested in models over semitopologies whose open sets have certain \emph{antiseparation properties}, by which we mean properties that involve the existence of nonempty intersections between open sets.\footnote{For comparison, topology is often concerned with \emph{separation} properties, like Hausdorff separation, which have to do with existence of \emph{empty} intersections between open sets.  In this sense, the interest in this Section is the precise dual to where the emphasis traditionally lies in topology.}

The design space of antiseparation properties is rich, as the first author explored in~\cite{gabbay:semdca} --- but for our needs in this paper we will just one canonical property which we will call being \emph{$n$-twined}.
We define the $n$-twined property in Definition~\ref{defn.twined}, we explain why it matters in Remark~\ref{rmrk.motivate.n-twined}, and then we study its properties.


\begin{defn}
\label{defn.twined}
\leavevmode
Suppose $(\Pnt,\opens)$ is a semitopology and $n\in\Ngeqz$.
Then call $(\Pnt,\opens)$ \deffont{$n$-twined} when every collection of $n$ nonempty open sets has a nonempty open intersection.
In symbols:
$$
\Forall{O_1,\dots,O_n\in\opensne}O_1\cap\dots\cap O_n\neq\varnothing.
$$
\end{defn}

\begin{lemm}
\label{lemm.1-twined}
We note two edge cases:
\begin{enumerate*}
\item
A semitopology is $0$-twined if and only if it is nonempty (Notation~\ref{nttn.empty.semitopology}).
\item
Every semitopology is $1$-twined.
\end{enumerate*}
\end{lemm}
\begin{proof}
We consider each part in turn:
\begin{enumerate}
\item
If $n=0$ then we need to show that $\Pnt$ (the intersection of zero nonempty sets) is not equal to $\varnothing$.
by Notation~\ref{nttn.empty.semitopology} this is true precisely when the semitopology is nonempty. 
\item
By Definition~\ref{defn.twined}, we must show that $O\neq\varnothing$ for every $O\in\opensne$.
This is true.
\qedhere\end{enumerate}
\end{proof}

\begin{rmrk}[Discussion of the special cases $n=2$ and $n=3$]
\label{rmrk.motivate.n-twined}
Lemma~\ref{lemm.1-twined} shows that being $0$-twined and $1$-twined are straightforward properties.
In contrast, the $2$- and $3$-twined semitopologies will be very interesting to us in this paper.
We give a high-level discussion of why this is:
\begin{itemize}
\item
Suppose we are in a distributed system running on a $2$-twined semitopology.
We assume that participants are honest.
Suppose further that we have 
\begin{itemize*}
\item
a nonempty open set of participants $O\in\opensne$ who agree that some predicate $\phi$ is true (respectively: valid), and 
\item
a nonempty open set of participants $O'\in\opensne$ who agree that some predicate $\phi'$ is true (respectively: valid).
\end{itemize*}
Then because the system is $2$-twined, there exists some $p\in O\cap O'$ who agrees that both $\phi$ and $\phi'$ are true (respectively: valid).\footnote{We will render this in our logic later on, in Lemma~\ref{lemm.supertwined.qcbt}(\ref{item.intertwined.qcbt.2}).}

This simple observation underlies guarantees of consistent outputs in many distributed algorithms, because it guarantees that $\phi$ and $\phi'$ are not logically conflicting --- e.g. that we do not have something like $\phi'=\modF\phi$ --- because if they were, then an honest $p$ could not believe both.

In the special case that $\phi$ is `the output of our consensus algorithm is $v$' and $\phi'$ is `the output of our consensus algorithm is $v'$', the $2$-twined property makes the strong guarantee that where consensus algorithms succeed (i.e. where participants do not crash), the output is consistent. 
\item
Suppose we are in a distributed system running on a $3$-twined semitopology.
Some participants may be dishonest, in that they may report inconsistent predicates. 

Suppose further that we have nonempty open sets $O,O'\in\opensne$ of participants who agree on the truth (respectively: validity) of predicates $\phi$ and $\phi'$ respectively.
Suppose further that we have a nonempty open set $O''\in\opensne$ of participants who are honest (meaning that we bound dishonest participants to at most some \emph{closed} set). 
Because the system is $3$-twined, there exists some $p\in O\cap O'\cap O''$ --- this $p$ is a participant who agrees on the truth (respectively: validity) of $\phi$ and $\phi'$ and is also honest.\footnote{We will render this in our logic twice: in Lemma~\ref{lemm.supertwined.qcbt}(\ref{item.supertwined.qcbt.2}) as a general result, and then as a useful special case in Corollary~\ref{corr.two.quorums.intersect.correct}(\ref{item.two.quorums.intersect.correct.supertwined}).}

In this case, $\phi\tand\phi'$ must be consistent, because if it were not then the honest participant $p$ could not believe both $\phi$ and $\phi'$.
Just as for the $2$-twined case, this simple argument underlies many consistency arguments in distributed systems in the presence of dishonest participants.
\end{itemize}
There is no barrier to considering $n$-twined semitopologies for $n\geq 4$, and this would just give us stronger guarantees of intersection.

A brief survey of intersection properties in the literature is in Remark~\ref{rmrk.ntwined.and.the.literature}
\end{rmrk}


\subsection{Being $n$-twined and being dense / having nonempty open interior}

The property of being $n$-twined is closely related to the topological notion of being a \emph{dense set} and having a \emph{nonempty open interior} (Definition~\ref{defn.dense.subset}).
We explore this in Lemma~\ref{lemm.is} and Proposition~\ref{prop.noi.dense}; then in Proposition~\ref{prop.allbut.supertwined} we specialise the discussion to the important special case of $\AllBut{N}{f}$.

\begin{lemm}
\label{lemm.is}
Suppose $(\Pnt,\opens)$ is a semitopology and $n\in\mathbb N_{\geq 1}$.
Then:
\begin{enumerate*}
\item\label{item.is.supertwined}
$(\Pnt,\opens)$ is $n$-twined if and only if every collection of $n\minus 1$ nonempty open sets has a dense intersection: 
$$
\Forall{O_1,\dots O_{n\minus 1}\in\opensne}\dense(O_1\cap\dots\cap O_{n\minus 1}).
$$
\item\label{item.is.intertwined}
As a corollary, $(\Pnt,\opens)$ is \xxtwined if and only if every nonempty open set is dense: 
$$
\Forall{O\in\opensne}\dense(O).
$$
\end{enumerate*}
\end{lemm}
\begin{proof}
Both parts follow from Definitions~\ref{defn.dense.subset}(1) and~\ref{defn.twined}, just by pushing around definitions.

One corner case merits brief mention: if $n=1$ then we must show that every collection of $0$ sets (i.e. $\varnothing$) is dense.
So, we need to show that $\bigcap\varnothing$ is dense.
But by convention, the intersection of \emph{no} open sets is the whole of $\Pnt$, and this is indeed dense, so we are done. 
\end{proof}


\begin{prop}
\label{prop.noi.dense}
Suppose $(\Pnt,\opens)$ is a semitopology and $n\in\mathbb N_{\geq 1}$.
Then the following are equivalent:
\begin{enumerate*}
\item
$(\Pnt,\opens)$ is $n$-twined. 
\item
$\Forall{P_1,\dots,P_{n\minus 1}\subseteq\Pnt}(\bigwedge_{1{\leq}i{\leq}n\minus 1}\noi(P_i)) \limp \dense(\bigcap_{1{\leq}i{\leq}n\minus 1}P)$.
\end{enumerate*}
\end{prop}
\begin{proof}
Part~2 trivially implies part~1 because we can take $P_1,\dots,P_n\in\opensne$, so that by Definition~\ref{defn.dense.subset} $\noi(P_i)$ for $1\leq i\leq n\minus 1$ holds by construction by Lemma~\ref{lemm.dense.up-closed}(\ref{item.noi.noi}), and we use Lemma~\ref{lemm.is}(\ref{item.is.supertwined}).

Showing that part~1 implies part~2 is only slightly less trivial.
Suppose $(\Pnt,\opens)$ is $n$-twined and $P_1,\dots,P_{n\minus 1}\subseteq\Pnt$ and $\noi(P_i)$ for each $1\leq i\leq n\minus 1$, meaning by Definition~\ref{defn.dense.subset} that for each such $i$ there exists an $O_i\in\opensne$ such that $O_i\subseteq P_i$.
By Lemma~\ref{lemm.is}(\ref{item.is.supertwined}) $O_1\cap\dots\cap O_{n\minus 1}$ is dense, and by Lemma~\ref{lemm.dense.up-closed}(\ref{item.dense.up-closed}) so is $P_1\cap\dots\cap P_n$.
\end{proof}



Proposition~\ref{prop.allbut.supertwined} provides the canonical example of $n$-twined semitopologies.
It mildly generalises known properties in the literature; the case $n=2$ in Proposition~\ref{prop.allbut.supertwined}(\ref{item.allbut.intertwined}) corresponds to the $Q^2$ property from~\cite[Theorem~1]{DBLP:conf/asiacrypt/DamgardDFN07}, and the case $n=3$ corresponds to the $Q^3$ property of~\cite[Definition 2]{DBLP:journals/dc/AlposCTZ24}.
The topology-flavoured presentations in Proposition~\ref{prop.allbut.supertwined}(\ref{item.allbut.NOI.dense}\&\ref{item.allbut.NOI.dense.3}) are new.
\begin{prop}
\label{prop.allbut.supertwined}
Suppose $N,f\in\Ngeqz$ and $P\subseteq\Nset{N}$ and $n\in\mathbb N_{\geq 1}$.
Then:
\begin{enumerate*}
\item\label{item.allbut.intertwined}\label{item.allbut.supertwined}
$N>n*f$ if and only if $\AllBut{N}{f}$ is $n$-twined.
\item\label{item.allbut.NOI.dense}
$\noi(P)$ (and equivalently $P\in\opensne$) if and only if $\mycard P\geq N\minus f$.\footnote{In general semitopologies, $\noi(P)$ ($P$ having a nonempty open interior) is not the same as $P\in\opensne$ ($P$ being a nonempty open set); but in the case of $\AllBut{N}{f}$, these two properties happen to be equal.}
\item\label{item.allbut.NOI.dense.3}
$\dense(P)$ if and only if $\mycard P\geq f\plus 1$.
\end{enumerate*}
\end{prop}
\begin{proof}
We consider each part in turn:
\begin{enumerate}
\item
Suppose $N>n*f$ and suppose $O_1,\dots,O_n$ are open in $\AllBut{N}{f}$.
From Definition~\ref{defn.allbut} this means that they each contain all but $f$ elements of $\Nset{N}$.
It follows that $O_1\cap\dots \cap O_n$ contains (at least) all but $n*f$ elements of $\Nset{N}$.
Since $N\gneq n*f$, this intersection is nonempty. 

Conversely, if $N\leq n*f$ then it is routine to partition $\Nset{N}$ into $n$ sets, each of which has at most $f$ elements.
Then their complements are a set of $n$ open sets whose intersection is empty.
\item
By routine counting arguments.
\item
By routine counting arguments.
\qedhere\end{enumerate}
\end{proof}



\begin{rmrk}
\label{rmrk.ntwined.and.the.literature}
Being \xxtwined corresponds to the \emph{quorum intersection property} of~\cite[Section~2.1, Condition~B2]{DBLP:journals/tocs/Lamport98}, which requires that any two quorums intersect.
Being \xxxtwined also corresponds to a quorum intersection property with a stronger requirement that (assuming there is a quorum of honest participants) any two quorum must intersect in a participant that is honest~\cite[Definition~5.1]{DBLP:journals/dc/MalkhiR98}. 
  
These intersection properties are closely tied to necessary conditions for some consensus protocols (and other distributed algorithms) to function correctly.
Specifically for the case of the semitopology $\AllBut{N}{f}$ from Definition~\ref{defn.allbut}(\ref{item.allbut}), requiring non-empty quorum intersection is equivalent to the condition \( N > 2f \), where \( f \) denotes the number of processes that may crash during the protocol's execution.
This condition is often referred to as \( Q^2 \)~\cite[Theorem~1]{DBLP:conf/asiacrypt/DamgardDFN07}.
Similarly, being \xxxtwined corresponds to the condition \( N > 3f \), commonly referred to as \( Q^3 \)~\cite[Definition 2]{DBLP:journals/dc/AlposCTZ24}.
A more generalised version of \( Q^3 \) exists, known as \( B^3 \)~\cite[Theorem 2]{DBLP:conf/asiacrypt/DamgardDFN07}\cite[Definition 5]{DBLP:journals/dc/AlposCTZ24}, which accounts for greater heterogeneity. 
Specifically, \( Q^3 \) applies to distributed systems where failures are homogeneous, meaning they are globally known and agreed upon by all participants. 
For example, all participants understand that up to \( f \) processes may fail. 
  
In contrast, in a heterogeneous setting, participants may have differing beliefs about \emph{who} might fail, leading to a more nuanced description of trust assumptions. 
In this context, \( B^3 \) generalises \( Q^3 \) by incorporating these variations in trust and failure assumptions.
Readers interested in exploring this topic further can consult~\cite{DBLP:conf/asiacrypt/DamgardDFN07,DBLP:journals/dc/AlposCTZ24}.
\end{rmrk}
  
  
  
  

\subsection{The canonical properties}

For this subsection we fix a signature $\Sigma$ and a model $\mathcal M=(\Pnt,\opens,\Val)$.

The results in this Subsection express how the underlying semitopology of this model $(\Pnt,\opens)$ being \xxtwined or \xxxtwined manifests itself in the logic, in ways that will be relevant in our later proofs:
\begin{lemm}
\label{lemm.supertwined.qcbt}
Suppose $(\Pnt,\opens)$ is a semitopology and $n\in\Time$ and $\avaluation$ is a valuation and suppose $\phi$, $\psi$, and $\chi$ are pointwise in $(n,\avaluation)$ (Definition~\ref{defn.pointwise.meaning}). 
Then: 
\begin{enumerate*}
\item\label{item.intertwined.qcbt}
If $(\Pnt,\opens)$ is \xxtwined then we have:
\begin{enumerate*}
\item\label{item.intertwined.3}
If $n\mentval \QuorumBox\phi$ then $n\mentval\CoquorumDiamond\phi$. 
\item\label{item.intertwined.qcbt.2}
If $n\mentval \QuorumBox\phi \tand \QuorumBox\psi$ then $n\mentval\someoneAll(\phi\tand\psi)$. 
\item\label{item.intertwined.qcbt.3}
If $n\mentval \QuorumBox\correct{\phi} \tand \QuorumBox\phi$ then $n\mentval\modT\someoneAll\phi$.
\end{enumerate*}
\item\label{item.supertwined.qcbt}
If $(\Pnt,\opens)$ is \xxxtwined then we have: 
\begin{enumerate*}
\item\label{item.supertwined.qcbt.3}
If $n\mentval \QuorumBox\phi \tand \QuorumBox\psi$ then $n\mentval \CoquorumDiamond(\phi\tand\psi)$.
\item\label{item.supertwined.qcbt.2}
If $n\mentval \QuorumBox\phi \tand \QuorumBox\psi \tand \QuorumBox\chi$ then $n\mentval\someoneAll(\phi\tand\psi\tand\chi)$. 
\item\label{item.supertwined.qcbt.4}
If $n\mentval \QuorumBox\correct{\phi} \tand \QuorumBox\phi$ then $n\mentval \modT\CoquorumDiamond\phi$.
\end{enumerate*}
\end{enumerate*}
Recall that $\correct{\phi}$ is from Definition~\ref{defn.correct}(\ref{item.correct.phi}) and just means $\modTF\phi$.
\end{lemm}
\begin{proof}
We reason as follows:
\begin{enumerate}
\item
Suppose $(\Pnt,\opens)$ is \xxtwined. 
\begin{enumerate*}
\item
We reason as follows:
$$
\begin{array}{r@{\ }l@{\qquad}l}
n\mentval\QuorumBox\phi
\liff&
\noi(\{p\in\Pnt \mid n,p,\avaluation\ment\phi\})
&\text{Proposition~\ref{prop.pointwise.dense.char}(\ref{item.dense.char.pointwise.QuorumBox})}
\\
\limp&
\dense(\{p\in\Pnt \mid n,p,\avaluation\ment\phi\}) 
&\text{Proposition~\ref{prop.noi.dense}, $(\Pnt,\opens)$ $2$-twined} %
\\
\liff&
n\mentval\CoquorumDiamond\phi
&\text{Proposition~\ref{prop.pointwise.dense.char}(\ref{item.dense.char.pointwise.CoquorumDiamond})}
\end{array}
$$
\item
Suppose $n\mentval \QuorumBox\phi\tand\QuorumBox\psi$.
By part~\ref{item.intertwined.3} of this result, $n\mentval\CoquorumDiamond\phi$.
By Proposition~\ref{prop.someone.implies.someoneAll}(\ref{item.someone.implies.someoneAll.4}) $n\mentval \someoneAll(\phi\tand\psi)$ as required.
\item
Suppose $n\mentval \QuorumBox\correct{\phi} \tand \QuorumBox\phi$.
By part~\ref{item.intertwined.qcbt.2} of this result, $n\mentval\someoneAll(\correct{\phi}\tand\phi)$.
Using %
Proposition~\ref{prop.tv.ment.TF.model}(\ref{item.tv.ment.TF.phi.lr}) $n\mentval\someoneAll\modT\phi$.
By Lemma~\ref{lemm.someone.commutation} $n\mentval\modT\someoneAll\phi$ as required. 
\end{enumerate*}
\item
The reasoning for the \xxxtwined case is just as for the \xxtwined case, except we need intersection of three open sets instead of two.
\qedhere\end{enumerate} 
\end{proof}




Corollary~\ref{corr.two.quorums.intersect.correct} specialises Lemma~\ref{lemm.supertwined.qcbt} to the useful special case of a correct unary predicate symbol, such that $\phi=\correct{\tf P}$ (Definition~\ref{defn.correct}(\ref{item.correct})) and $\psi=\tf P(v)$ and $\chi=\tf P(v')$.
In our later examples, $\tf P$ will take values such as $\tf{accept}$, $\tf{send}$, $\tf{write}$, and $\tf{decide}$, representing basic actions in a consensus protocol:
\begin{corr}
\label{corr.two.quorums.intersect.correct}
Suppose $(\Pnt,\opens)$ is a semitopology and suppose $\tf P$ is a unary predicate symbol and $v,v'\in\Val$ and $n\in\Time$ and $\avaluation$ is a valuation.
Then:
\begin{enumerate*}
\item
If $(\Pnt,\opens)$ is \xxxtwined then:
\begin{enumerate*}
\item\label{item.two.quorums.intersect.correct.supertwined}
If $n\mentval\QuorumBox\correct{\tf P}\tand\QuorumBox\tf P(v)\tand\QuorumBox\tf P(v')$ then $n\mentval\modT\someoneAll(\tf P(v)\tand\tf P(v'))$.
\item\label{item.QuorumBox.correct.supertwined}
If $n\mentval\QuorumBox\correct{\tf P}\tand\QuorumBox\tf P(v)$ then $n\mentval\modT\CoquorumDiamond\tf P(v)$.
\end{enumerate*}
\item\label{item.two.quorums.intersect.correct.intertwined}
If $(\Pnt,\opens)$ is \xxtwined then
$n\mentval\QuorumBox\correct{\tf P}\tand\QuorumBox\tf P(v)$ implies $n\mentval\modT\someoneAll(\tf P(v))$.
\end{enumerate*}
\end{corr}
\begin{proof}
We consider each part in turn:
\begin{enumerate}
\item[(\ref{item.two.quorums.intersect.correct.supertwined})]
Suppose $(\Pnt,\opens)$ is \xxxtwined and $n\mentval\QuorumBox\correct{\tf P}\tand\QuorumBox\tf P(v)\tand\QuorumBox\tf P(v')$.
By Lemma~\ref{lemm.supertwined.qcbt}(\ref{item.supertwined.qcbt.2}) (since $\tf P(v)$ and $\correct{\tf P}$ are pointwise by Lemma~\ref{lemm.atomic.O}(\ref{item.atomic.pointwise.2})) $n\mentval\someoneAll(\correct{\tf P}\tand\tf P(v)\tand\tf P(v'))$.
By two applications of Proposition~\ref{prop.tv.ment.TF.model}(\ref{item.tv.ment.TF.P}) and Lemma~\ref{lemm.commuting.connectives} (for $\tand$)
$n\mentval\someoneAll\modT(\tf P(v)\tand\tf P(v'))$.
By Lemma~\ref{lemm.someone.commutation} $n\mentval\modT\someoneAll(\tf P(v)\tand\tf P(v'))$ as required.
\item[(\ref{item.QuorumBox.correct.supertwined})] 
Suppose $(\Pnt,\opens)$ is \xxxtwined and $n\mentval\QuorumBox\correct{\tf P}\tand\QuorumBox\tf P(v)$.
Using Lemma~\ref{lemm.correct.P.Pv} it follows that $n\mentval\QuorumBox\correct{\tf P(v)}\tand\QuorumBox\tf P(v)$.
We just apply Lemma~\ref{lemm.supertwined.qcbt}(\ref{item.supertwined.qcbt.4}). 
\item[(\ref{item.two.quorums.intersect.correct.intertwined})]
This follows from Lemma~\ref{lemm.supertwined.qcbt}(\ref{item.intertwined.qcbt.3}) just as part~\ref{item.QuorumBox.correct.supertwined} does from Lemma~\ref{lemm.supertwined.qcbt}(\ref{item.supertwined.qcbt.2}).
\qedhere\end{enumerate}
\end{proof}


Lemma~\ref{lemm.allbut.concretely} unpacks concretely what $\QuorumBox\phi$ and $\CoquorumDiamond\phi$ --- which are the subject of many results including of Lemmas~\ref{lemm.supertwined.qcbt} and~\ref{lemm.intertwined.char} --- mean in the semitopology $\AllBut{N}{f}$ from Definition~\ref{defn.allbut}(\ref{item.allbut}), which has $N$ points, and has open sets being either the empty set, or any set with at least $N\minus f$ elements: 
\begin{enumerate*}
\item
$\QuorumBox\phi$ means `for at least $N\minus f$ elements'; and 
\item
$\CoquorumDiamond\phi$ means `for at least $f\plus 1$ elements.
\end{enumerate*}
This special case is particularly relevant because (as the reader may know) $\AllBut{N}{f}$ is commonly used in the literature on distributed algorithms. 
\begin{lemm}
\label{lemm.allbut.concretely}
Suppose that
\begin{itemize*}
\item
$N,f\in\Ngeqz$, 
\item
$(\Pnt,\opens)=\AllBut{N}{f}$, 
\item
$n\in\Time$ and $\avaluation$ is a valuation, and
\item
$\phi$ is a closed predicate that is pointwise in $(n,\avaluation)$.
\end{itemize*}
Then:
\begin{enumerate*}
\item
$n\mentval\QuorumBox\phi$ if and only if 
$\mycard\{p\in\Pnt \mid n,p\mentval\phi\}\geq N\minus f$. 
\item
$n\mentval\CoquorumDiamond\phi$ if and only if 
$\mycard\{p\in\Pnt \mid n,p\mentval\phi\}\geq f\plus 1$. 
\item
$n\mentval\modT\QuorumBox\phi$ if and only if 
$\mycard\{p\in\Pnt \mid n,p\mentval\modT\phi\}\geq N\minus f$. 
\item
$n\mentval\modT\CoquorumDiamond\phi$ if and only if 
$\mycard\{p\in\Pnt \mid n,p\mentval\modT\phi\}\geq f\plus 1$. 
\end{enumerate*}
\end{lemm}
\begin{proof}
We consider each part in turn:
\begin{enumerate}
\item
The proof is elementary but we spell it out.
We reason as follows:
$$
\begin{array}{r@{\ }l@{\quad}l}
n\mentval\QuorumBox\phi
\liff&
\noi(\{p\in\Pnt \mid n,p\mentval \modT\phi\})
&\text{Proposition~\ref{prop.pointwise.dense.char}(\ref{item.dense.char.pointwise.QuorumBox})}
\\
\liff&
\mycard\{p\in\Pnt \mid n,p\mentval \phi\}\geq N\minus f
&\text{Proposition~\ref{prop.allbut.supertwined}(\ref{item.allbut.NOI.dense})}
\end{array}
$$
\item
Much as the previous part, but using Proposition~\ref{prop.pointwise.dense.char}(\ref{item.dense.char.pointwise.CoquorumDiamond}).
\item
By Lemma~\ref{lemm.someone.commutation}, $n\mentval\modT\QuorumBox\phi$ if and only if $n\mentval\QuorumBox\modT\phi$.
We just use part~1 of this result.
\item
Much as the previous part, using part~2 of this result.
\qedhere\end{enumerate}
\end{proof}

\subsection{Further observations on intersection properties}

For this Subsection we fix a signature $\Sigma$ and a model $\mathcal M=(\Pnt,\opens,\Val)$.
\begin{rmrk}
We collect some further results as follows:
\begin{enumerate*}
\item
Lemma~\ref{lemm.intertwined.char} proves converses to the implications in Lemma~\ref{lemm.supertwined.qcbt}, provided that the signature is nonempty.

This is interesting because it tells us that, provided the signature is nondegenerate, Lemma~\ref{lemm.supertwined.qcbt} accurately captures and reflects the intersection properties that interest us. 
\item
Proposition~\ref{prop.reformulate} is a semitopological distillation of parts~\ref{item.intertwined.qcbt.2} and~\ref{item.supertwined.qcbt.2} of Lemma~\ref{lemm.supertwined.qcbt}.
Proposition~\ref{prop.reformulate} \emph{is} Lemma~\ref{lemm.supertwined.qcbt} and Corollary~\ref{corr.two.quorums.intersect.correct}, in semitopological rather than logical form.
\end{enumerate*}
\end{rmrk}

In Lemma~\ref{lemm.intertwined.char} we could replace every instance of the expression `$\modTF\tf P$' with the (by Definition~\ref{defn.correct}(\ref{item.correct.phi}) and Remark~\ref{rmrk.tf.synonyms}) synonymous expression `$\correct{\tf P}$'.
However, we have not set up a corresponding synonym for $\modTB\tf P$ (it would be `$\tf{valid}[\tf P]$'; we just never wrote that definition) so it seems natural to stay with $\modTF$ here.
\begin{lemm}
\label{lemm.intertwined.char}
Suppose $\tf P\in\ns{PredSymb}$ is a predicate constant symbol (Notation~\ref{nttn.logic.terminology}(\ref{item.nullary.ps}); a nullary predicate symbol).\footnote{The proof works just as well for an $n$-ary predicate symbol; we just pick \emph{any} $\arity(\tf P)$-tuple of values $\vec v\in\Pnt^{\arity(\tf P)}$ --- which is possible because we assumed $\Val$ is nonempty in Definition~\ref{defn.model}(\ref{item.model.val}) --- and we use $\tf P(\vec v)$ in the proof instead.}
Then:
\begin{enumerate*}
\item
If\ \  $\ment(\QuorumBox\modTF\tf P\tand\QuorumBox\modTB\tf P)\timpc \someoneAll\tf P$ then $(\Pnt,\opens)$ is \xxtwined. 
\\
(This is a form of converse to Lemma~\ref{lemm.supertwined.qcbt}(\ref{item.intertwined.qcbt.3}).)
\item
If\ \  $\ment(\QuorumBox\modTF\tf P\tand\QuorumBox\modTB\tf P)\timpc \CoquorumDiamond\tf P$ then $(\Pnt,\opensne)$ is \xxxtwined. 
\\
(This is a form of converse to Lemma~\ref{lemm.supertwined.qcbt}(\ref{item.supertwined.qcbt.4}).) 
\item\label{item.intertwined.char.reverse}
As a corollary, provided that the signature $\Sigma$ is has at least one predicate symbol --- i.e. provided $\PredSymb$ in
it is nonempty (Definition~\ref{defn.signature}(\ref{item.predicate.syntax.empty.signature})) --- then reverse implications to those in Lemma~\ref{lemm.supertwined.qcbt} hold.\footnote{No need to also insist that $\Val$ is nonempty, because we assume this in Definition~\ref{defn.model}(\ref{item.model.val}).}
\end{enumerate*}
Recall from Definition~\ref{defn.validity.judgement}(\ref{item.ment.phi}) that $\ment\phi$ means $\Forall{\acontext\in\tf{Ctx}(\Sigma,\mathcal M)}(\acontext\ment\phi)$.
\end{lemm}
\begin{proof}
We consider each part in turn:
\begin{enumerate}
\item
We prove a contrapositive of \strongmodusponens: we suppose $\opens$ is \emph{not} \xxtwined, so there exist $O,O'\in\opensne$ with $O\cap O'=\varnothing$, and we exhibit a valuation such that $\neg(\mentval\someone\tf P)$ (equivalently: $\mentval\modF\someone\tf P$).

We choose a valuation $\avaluation$ such that for any $n\in\Time$, 
\begin{itemize*}
\item
$\avaluation(\tf P)_{n,p}(\aval)=\tvF$ for $p\in O$, and 
\item
$\avaluation(\tf P)_{n,p}(\aval)=\tvB$ for $p\in \Pnt\setminus O$.
\end{itemize*}
The reader can check that $O\mentval \everyone\modTF\tf P$ and $O'\mentval \everyone\modTB P(\aval)$, but $\mentval\modF\someone\tf P$.
\item
We prove a contrapositive: suppose $\opens$ is \emph{not} \xxxtwined, so there exist $O,O',O''\in\opensne$ with $O\cap O'\cap O''=\varnothing$.

We choose a valuation $\avaluation$ such that for any $n\in\Time$, 
\begin{itemize*}
\item
$\avaluation(\tf P)_{n,p}(\aval)=\tvT$ for $p\in O$, and 
\item
$\avaluation(\tf P)_{n,p}(\aval)=\tvB$ for $p\in O'\setminus O$, and 
\item
$\avaluation(\tf P)_{n,p}(\aval)=\tvF$ for $p\in \Pnt\setminus (O\cup O')$.
\end{itemize*}
The reader can check that $O\cup O'\mentval \everyone\modTF\tf P$ and $O\cup O''\mentval \everyone\modTB P(\aval)$, but $O'\cup O''\mentval\modF\everyone\tf P$.
\item
Just by the arguments above.
\qedhere\end{enumerate}
\end{proof}






\begin{nttn}
\label{nttn.cover}
Suppose $\ns X$ is a set.
Then:
\begin{enumerate*}
\item
A \deffont{cover} of $\ns X$ is a set of sets $\mathcal X\subseteq\powerset(\ns X)$ such that $\bigcup\mathcal X=\ns X$.
\item
A \deffont{partition} of $\ns X$ is a cover of $\ns X$ all of whose elements are pairwise disjoint (so $\Forall{X,X'\in\mathcal X}X\cap X'=\varnothing\lor X=X'$).
\end{enumerate*}
\end{nttn}

The reader might like to try giving a proof of Proposition~\ref{prop.reformulate} using direct calculations (without building on the results we have so far).
This is possible, but it is messy.
The use of logic endows the proofs with a structure and cleanness that the direct proof cannot match: 
\begin{prop}
\label{prop.reformulate}
Suppose $(\Pnt,\opens)$ is a semitopology and $n\in\Time$ and $\avaluation$ is a valuation, and suppose $\phi$ is a closed predicate that is pointwise in $(n,\avaluation)$.
Then:
\begin{enumerate*}
\item\label{item.reformulate.1}
If $(\Pnt,\opens)$ is \xxtwined then $n\mentval\CoquorumDiamond\modTB\phi \tor \CoquorumDiamond\modF\phi$. 
\item\label{item.reformulate.1b}
As a corollary, if $(\Pnt,\opens)$ is \xxtwined then any two-element partition of $\Pnt$ (Notation~\ref{nttn.cover}) contains a dense element (Definition~\ref{defn.dense.subset}(\ref{item.dense.subset})).
\item\label{item.reformulate.2}
If $(\Pnt,\opens)$ is \xxxtwined then 
$n\mentval\CoquorumDiamond\modT\phi \tor \CoquorumDiamond\modB\phi \tor \CoquorumDiamond\modF\phi$. 
\item\label{item.reformulate.2b}
As a corollary, if $(\Pnt,\opens)$ is \xxxtwined then any three-element partition of $\Pnt$ contains a dense element.
\end{enumerate*} 
\end{prop}
\begin{proof}
We consider each part in turn:
\begin{enumerate}
\item
By Lemma~\ref{lemm.supertwined.qcbt}(\ref{item.intertwined.qcbt.2})
$n\mentval\QuorumBox\modF\phi\tand\QuorumBox\modTB\phi$ implies $n\mentval\someoneAll(\modF\phi\tand\modTB\phi)$.
Unpacking what $n\mentval\someoneAll(\modF\phi\tand\modTB\phi)$ means using Definition~\ref{defn.validity.judgement} and Figures~\ref{fig.3.derived} and~\ref{fig.3.phi.f}, we see that this is impossible --- for no $p$ and $O$ can we satisfy both $n,p,O\mentval\modF\phi$ and $n,p,O\mentval\modTB\phi$ together --- so that 
$n\nmentval\QuorumBox\modF\phi\tand\QuorumBox\modTB\phi$.
By Lemma~\ref{lemm.para}(\ref{item.para.em})
$n\mentval\tneg(\QuorumBox\modF\phi\tand\QuorumBox\modTB\phi)$, and simplifying further using de Morgan equivalences as per Figure~\ref{fig.easy.equivalences}, we obtain
$n\mentval\CoquorumDiamond\modTB\phi\tor\CoquorumDiamond\modF\phi$ as required.
\item
Suppose $\{P_1,P_2\}$ partitions $\Pnt$.
Assume a $0$-ary predicate symbol $\tf P$, and let $\varsigma$ be the valuation that sets $\varsigma(\tf P)_{n,p}=\tvT$ if $p\in P_1$ and $\varsigma(\tf P)_{n,p}=\tvF$ if $p\in P_2$.\footnote{We can assume $\tf P$ and $\varsigma$ because this result is about semitopologies, so we are free to impose any model structure on it that we like.  It seems particularly nice here how the mere \emph{existence} of a model for our logic gives us a clean proof of a purely semitopological property which, on the face of it, has nothing to do with our three-valued modal logic.}
By part~\ref{item.reformulate.1} and Lemma~\ref{lemm.tand.tor} either $n\mentval \Coquorum\modTB\tf P$ or $n\mentval\Coquorum\modF\tf P$.
By Proposition~\ref{prop.pointwise.dense.char}(\ref{item.dense.char.pointwise.CoquorumDiamond}) either $P_1$ or $P_2$ is dense.
\item
We reason as for the proof of part~\ref{item.reformulate.1}, using Lemma~\ref{lemm.supertwined.qcbt}(\ref{item.supertwined.qcbt.2}) and starting from $\modFB\phi \tand \modTF\phi \tand \modTB\phi$. 
\item
We reason as for the proof of part~\ref{item.reformulate.2}, for a partition $\{P_1,P_2,P_3\}$ and setting $\varsigma(\tf P)_{n,p}=\tvT$ if $p\in P_1$ and $\varsigma(\tf P)_{n,p}=\tvB$ if $p\in P_2$ and $\varsigma(\tf P)_{n,p}=\tvF$ if $p\in P_3$.
\qedhere\end{enumerate}
\end{proof}




 





\section{Declarative Paxos}
\label{sect.declarative.paxos}

\subsection{Case study: a simple protocol}
\label{subsect.simple}

\begin{rmrk}
\label{rmrk.simple.protocol}
Before Paxos, we will warm up by sketching a simple protocol \theory{Simple}, and then we will give four \QLogic axiomatisations of \theory{Simple}, reflecting four different failure assumptions (messages can/cannot get lost; participants can/cannot crash).
\theory{Simple} is as follows:
\begin{enumerate*}
\item
Any participant $p$ may broadcast a $\tf{propose}$ message to all participants.
\item
If a participant receives a $\tf{propose}$ message, then it responds with $\tf{accept}$. 
\item
If the participant $p$, having broadcast its $\tf{propose}$ message, receives a $\tf{accept}$ message, then it declares $\tf{decide}$.
\end{enumerate*}
\end{rmrk}


\begin{rmrk}
\theory{Simple} is very simple!
Only one value can be proposed, so we do not even bother to state it; we just `propose'.
Any participant accepts immediately that it hears someone propose. 
Any participant that proposes, decides immediately that it hears that someone accepted.

For comparison, in Paxos \emph{multiple values} are available, so we need to make sure that participants all decide on the \emph{same} value if they decide on anything at all. 
This makes things much harder: so Paxos has multiple rounds; there is a designated leader participant at each round; the leader must choose a value to propose; and elaborate communication is in place to ensure that (subject to suitable precisely-stated assumptions) a single choice of agreed value will eventually be decided.

But that is for later.  
Even the simple protocol \theory{Simple} has plenty of subtlety, which will allow us to exhibit (in `toy' form) many of the techniques that we will then apply to Paxos:
\end{rmrk}


\begin{rmrk}
We will axiomatise \theory{Simple} subject to two parameters, yielding four possible axiomatisations:
\begin{enumerate*}
\item
Do all messages arrive, or might there be a finite initial period during which the network may be unreliable so that during that time, messages might not arrive?
\item
Can participants crash (and if so, how many)?
\end{enumerate*}
\end{rmrk}

\subsubsection*{All messages arrive; no participants crash}

First, we assume an ideal situation that all messages arrive, and no participants crash.
We can axiomatise this as follows:
$$
\begin{array}{l@{\qquad}l}
\rulefont{SimpAccept?}&
\tf{accept} \timpc \someoneAll\tf{propose} 
\\
\rulefont{SimpDecide?}&
\tf{decide} \timpc (\tf{propose}\tand\someoneAll\tf{accept}) 
\\
\rulefont{SimpAccept!}&
(\someoneAll\tf{propose}) \timpc \tf{accept}
\\
\rulefont{SimpDecide!}&
(\tf{propose}\tand\someoneAll\tf{accept}) \timpc \tf{decide}
\\
\rulefont{Correct}&
\correct{\tf{propose},\tf{accept},\tf{decide}}
\end{array}
$$
\noindent Recall from Definition~\ref{defn.three.imp} and Proposition~\ref{prop.mp.for.tnotor} that $\timpc$ is a \emph{strong implication} and means `if the LHS is $\tvT$ then so is the RHS', and 
recall from Proposition~\ref{prop.pointwise.dense.char}(\ref{item.dense.char.pointwise.someoneAll}) that $\someoneAll\tf P$ means `$\tf P$ somewhere', where $\tf P\in\{\tf{propose},\tf{accept},\tf{decide}\}$.

Then we explain the axioms (briefly) as follows:
\begin{itemize*}
\item
We call \rulefont{SimpAccept?} and \rulefont{SimpDecide?} \emph{backward rules}; they express backward versions of the rules in Remark~\ref{rmrk.simple.protocol}, namely that:
\begin{itemize*}
\item
if a participant asserts $\tf{accept}$ then some participant must have asserted $\tf{propose}$; and 
\item
if a participant asserts $\tf{decide}$ then it must have asserted $\tf{propose}$, and some participant must have asserted $\tf{accept}$.
\end{itemize*}
\item
We call \rulefont{SimpAccept!} and \rulefont{SimpDecide!} \emph{forward rules}.
They express forward versions of the rules in Remark~\ref{rmrk.simple.protocol}.
Because our failure assumptions are that nothing fails --- all messages arrive; no participants crash --- the forward rules are in this case just the backward rules, reversed.
\item
Axiom \rulefont{Correct} expresses that all predicates return $\tvT$ and $\tvF$ (not $\tvB$), as per Definition~\ref{defn.correct}(\ref{item.correct}). 
This reflects our assumption that no participants crash, because we model a computation crashing %
as the predicate describing it returning $\tvB$.
\end{itemize*}
Note that $\tf{propose}$, $\tf{accept}$, and $\tf{decide}$ above are not messages; they are propositions!
There is no message-passing.
$\tf{accept}$ being true at a participant represents that the participant accepts; this is implemented in the algorithm by the participant broadcasting an accept message, and represented in the logic by $\tf{accept}$ taking truth-value $\tvT$ at that participant.

Our example protocol is simple, but we hope it is already starting to illustrate how the various parts of our logic come together to express the protocol axiomatically.
We continue:

\subsubsection*{All messages arrive; participants can crash}

We assume that all messages arrive, and participants can crash --- but we assume there is a dense set of uncrashed participants.
Recall from Definition~\ref{defn.dense.subset}(\ref{item.dense.subset}) that a \emph{dense set} is one that intersects every open set; in particular, this means that every open set of participants contains a participant that is not crashed.
We model being crashed as returning the third truth-value $\tvB$.
We can axiomatise this as follows:
$$
\begin{array}{l@{\qquad}l}
\rulefont{SimpAccept?}&
\tf{accept} \timpc \someoneAll\tf{propose} 
\\
\rulefont{SimpDecide?}&
\tf{decide} \timpc (\tf{propose}\tand\someoneAll\tf{accept}) 
\\
\rulefont{SimpAccept!}&
(\someoneAll\tf{propose}) \tnotor \tf{accept}
\\
\rulefont{SimpDecide!}&
(\tf{propose}\tand\someoneAll\tf{accept}) \tnotor \tf{decide}
\\
\rulefont{Uncrashed}&
\QuorumBox\modTB\tf P \tnotor \someoneAll\modT\tf P
\quad \tf P\in\{\tf{propose},\tf{accept},\tf{decide}\}
\end{array}
$$ 
\noindent 
The backward rules are unchanged: if a participant accepts then from the form of the rules, it \emph{must} be the case that somebody proposed; and if a participant decides then from the form of the rules, it \emph{must} be the case that the participant proposed and some participant accepted.
 
The forward rules \rulefont{SimpAccept!} and \rulefont{SimpDecide!} have changed, and they now use the weak implication $\tnotor$.
By Definition~\ref{defn.three.imp} and Proposition~\ref{prop.mp.for.tnotor}, $\tnotor$ means `if the LHS is $\tvT$ then the RHS is $\tvT$ or $\tvB$'.
The $\tvB$ truth-value represents `crashed', so that (for example) \rulefont{SimpAccept!} expresses that for each participant $p$, if some (possibly other) participant $q$ asserts $\tf{propose}$, then $p$ asserts $\tf{accept}$ \emph{or} $p$ has crashed. 

Axiom \rulefont{Uncrashed} reflects that any a quorum (open set) of participants contains an uncrashed participant; e.g. in topological language, that the set of uncrashed participants is dense.
More precisely: \rulefont{Uncrashed} asserts that if an open set of participants returns $\tvT$ or $\tvB$ for $\tf P$, then at least one of them (is not crashed and so) returns $\tvT$.

Note from Lemma~\ref{lemm.someone.commutation} that $\Quorum\modTB$ is equivalent to $\modTB\Quorum$, and $\someoneAll\modT$ is equivalent to $\modT\someoneAll$, so that an equivalent form of \rulefont{Uncrashed} is $\modTB\QuorumBox\tf P \tnotor \modT\someoneAll\tf P$.
In our axiomatisations it is usually helpful to pull $\modTB$ and $\modT$ up to the topmost level possible in the predicate, so we may favour that form, but this is just a matter of convenience; meaning is unchanged. 


\subsubsection*{Messages might not arrive; no participants crash}

We assume that there is an initial period during which the network might not deliver messages, but the network eventually stabilises (in the literature this is called \emph{eventual synchrony}).
No participants crash.
We can axiomatise this as follows:
$$
\begin{array}{l@{\qquad}l}
\rulefont{SimpAccept?}&
\tf{accept} \timpc \someoneAll\tf{propose} 
\\
\rulefont{SimpDecide?}&
\tf{decide} \timpc (\tf{propose}\tand\QuorumBox\tf{accept}) 
\\
\rulefont{SimpAccept!}&
\final((\someoneAll\tf{propose}) \timpc \tf{accept})
\\
\rulefont{SimpDecide!}&
\final((\tf{propose}\tand\someoneAll\tf{accept}) \timpc \tf{decide})
\\
\rulefont{Correct}&
\correct{\tf{propose},\tf{accept},\tf{decide}}
\end{array}
$$ 
\noindent The backward rules are unchanged, but the forward rules \rulefont{SimpAccept!} and \rulefont{SimpDecide!} now have a $\final$ modality.
This is from Figure~\ref{fig.3.derived} and it means `eventually, this holds forever', reflecting an assumption that eventually the network stabilises.

Note that \QLogic does not explicitly represent message-passing, but this does not matter as much as one might suppose!
The $\final$ modality captures the essence of what we need, namely, that eventually the forward rules become valid. 

\subsubsection*{Messages might not arrive; participants can crash}

We assume that messages might not arrive, and participants can crash but there is a dense set of uncrashed participants (so that every quorum of participants contains at least one uncrashed participant).
We can axiomatise this just by combining the ideas above, as follows:
$$
\begin{array}{l@{\qquad}l}
\rulefont{SimpAccept?}&
\tf{accept} \timpc \someoneAll\tf{propose} 
\\
\rulefont{SimpDecide?}&
\tf{decide} \timpc (\tf{propose}\tand\someoneAll\tf{accept}) 
\\
\rulefont{SimpAccept!}&
\final((\someoneAll\tf{propose}) \tnotor \tf{accept})
\\
\rulefont{SimpDecide!}&
\final((\tf{propose}\tand\someoneAll\tf{accept}) \tnotor \tf{decide})
\\
\rulefont{Uncrashed}&
\modTB\QuorumBox\tf P \tnotor \modT\someoneAll\tf P
\quad \tf P\in\{\tf{propose},\tf{accept},\tf{decide}\}
\end{array}
$$ 
\noindent The axiomatisations of Declarative Paxos and Simpler Declarative Paxos in Figures~\ref{fig.ThyPaxOne} and~\ref{fig.ThySPax} are more elaborate than the axioms above, because Paxos is a more complicated protocol, but the basic ideas are as above; just scaled up. 


\subsection{A high-level view of the Paxos algorithm}



\begin{rmrk}[High-level view of Paxos]
\label{rmrk.high-level.paxos}
Paxos allows participants to crash, and messages to fail to arrive at their destination, but it assumes all participants are honest (i.e. that if they have not crashed then they follow the rules of the algorithm).
The reader can find detailed treatments of Paxos in \cite{DBLP:journals/tocs/Lamport98,lamport2001paxos,cachin2011paxos} (a Wikipedia article~\cite{wiki:Paxos_(computer_science)} is also accessible and clear).
Here is a high-level view of the algorithm:
\begin{enumerate*}
\item\label{item.l.proposes}
\deffont{Leader proposes:}\quad

A designated leader participant (if it does not crash) sends messages to all participants, \emph{proposing} some unique value $v$.

Message-passing is unreliable, so messages may fail to arrive.
\item\label{item.p.sends}
\deffont{Participant sends:}\quad

If a participant $p$ receives a propose $v$ message from the leader --- it is assumed that every participant knows who the leader is, and other participants cannot forge the leader's signature --- then $p$ sends a \emph{send} message back to the leader, which carries the value $v_p$ that $p$ most recently accepted, along with the timestamp of when $v_p$ was accepted.\footnote{For reasons of compact presentation we do not explicitly represent the timestamps in our axiomatisation; i.e. $\tf{send}$ is a unary predicate in Definition~\ref{defn.declarative.paxos}(\ref{item.ThyPaxOne.sig}), not a binary predicate.
If this bothers the reader then we can just say that, actually, $v_p$ contains a timestamp of the stage $n\in\Time$ it was accepted; i.e. that $v_p$ is actually implemented as a pair $(v_p',n_p)$ of a value and a time.} 
If $p$ has never accepted any value, so no most recently accepted value exists, then $p$ sends back a special $\udfn$ value.
\item\label{item.l.writes}
\deffont{Leader writes:}\quad

The leader participant sends a message to all participants, \emph{writing} the most recently-timestamped value $v_p$ mentioned in any \emph{send} message it received from within the first quorum of participants that it hears back from; if the leader does not hear back from any quorum, then it does not write anything.
\item\label{item.p.accepts}
\deffont{Participant accepts:}\quad

If a participant receives a write message from the leader for some value, then it sends a message back \emph{accepting} that value.
\item\label{item.l.decides}
\deffont{Leader decides:}\quad

If the leader receives accept messages for some value from a quorum of participants, it decides on that value, and sends out a \emph{decide} message to all participants with that value.
\item\label{item.p.decides}
\deffont{Participant (non-leader) decides:}\quad

If a participant receives a decide message from the leader for some value, then it decides on that value and stops.
\end{enumerate*}
\end{rmrk}

\begin{rmrk}[High-level view of our axiomatisation]
\label{rmrk.elisions}
The theory $\ThyPax$ in Definition~\ref{defn.ThyPaxOne} and Figure~\ref{fig.ThyPaxOne} is a declarative abstraction of Paxos-the-algorithm.
It stands in relation to implementations of Paxos in roughly the same way that (say) the mathematical specification of the Fibonacci function stands in relation to its concrete implementations, in the sense that we elide implementational detail to preserve (and reveal) a mathematical ideal of what the implementation does.

\ThyPax elides two features of the concrete implementation:
\begin{enumerate*}
\item
There is no explicit message-passing; there are only predicates returning truth-values.

This may seem counterintuitive --- isn't Paxos \emph{supposed} to be a message-passing algorithm? --- but we shall see that a lot of the character of Paxos survives.
Indeed, arguably (and certainly for the first author) it is much easier to see what is really going on once we have done this.
\item\label{item.elide.algorithmic.time}
There is no abstract machine, and therefore no \deffont{algorithmic time}, by which we mean no sequence of state-changes of the abstract machine. 

The only notion of time is the \emph{stage} $n\in\Time$ from the context as defined in Definition~\ref{defn.contexts}(\ref{item.context.stage}).
We may call this \deffont{logical time} and (when we are modelling Paxos) it represents the sequence of what is in the distributed consensus algorithms literature often called \emph{stages}, or \emph{ballot numbers}.\footnote{In a distributed algorithm, at any point in algorithmic time, different participants in the algorithm may be at different points in logical time.  A practical example: consider several participants racing across 100 metres. In logical time, they are both just counting to 100.  In algorithmic time, one of them will get to 100 \emph{first}, and win.  In this paper, we only care about the logical aspects of the process; there may still be a winner participant, but this would implemented just as a unary predicate $\tf{winner}$ such that $100\ment\texiaffine\tf{winner}$.} 
\end{enumerate*}
We will prove standard Paxos correctness properties for \ThyPax in Section~\ref{sect.paxos.correctness.properties}. 
\end{rmrk}


\subsection{Definition of Declarative Paxos}

\begin{figure}
$$
\begin{array}{l@{\quad}r@{\ }l}
\figunderline{Backward rules, for backward inference} 
\figskip
\rulefont{PaxPropose?}&
& \tall v.\ \tf{propose}(v)\timpc(\tf{leader}\tand v\tneq\udfn)
\figskip
\rulefont{PaxSend?}&
& \tall v.\ \tf{send}(v) \timpc 
\begin{array}[t]{l}
\someoneAll\texi\tf{propose}\ \tand 
\\
\bigl(
(\mru{\tf{accept}}{v})\ \tor  
(\tneg\modT\recent\texi\tf{accept}\tand v\teq\udfn)\bigr)
\end{array}
\figskip
\rulefont{PaxWrite?}&
& \tall v.\ \tf{write}(v) \timpc 
\begin{array}[t]{l}
\tf{leader}\tand 
\\
\Quorum(\Box\texi\tf{send} \tand 
\\
\qquad\bigl((\mru{(\someone\tf{accept})}{v}) \tor
(\tf{propose}(v)\tand\tneg\modT\recent\texi\someone\tf{accept})\bigr)) 
\end{array}
\figskip
\rulefont{PaxAccept?}&
& \tall v.\ \tf{accept}(v) \timpc \someoneAll\tf{write}(v)
\figskip
\rulefont{PaxDecideL?}&
& \tf{leader} \timpc \tf{decide}(v) \timpc \QuorumBox\tf{accept}(v)
\figskip
\rulefont{PaxDecide\tneg L?}&
& \tall v.\ \tneg\tf{leader} \timpc \tf{decide}(v) \timpc \someoneAll(\tf{leader}\tand\tf{decide}(v))
\figskip\figskip
\figunderline{Forward rules, for making progress} 
\figskip
\rulefont{PaxPropose!}&
& \tf{leader}\tnotor\texi\tf{propose}
\figskip
\rulefont{PaxSend!}&
& \final(\someoneAll\texi\tf{propose} \tnotor \texi\tf{send}) 
\figskip
\rulefont{PaxWrite!}&
& \final\bigl((\tf{leader}\tand (\QuorumBox\texi\tf{send}) \tand \texi\tf{propose}) \tnotor \texi\tf{write}\bigr)
\figskip
\rulefont{PaxAccept!}&
& \final(\texi\someoneAll\tf{write}\tnotor \texi\tf{accept}) 
\figskip
\rulefont{PaxDecideL!}&
& \final((\tf{leader} \tand\texi\QuorumBox\tf{accept}) \tnotor \texi\tf{decide}) 
\figskip
\rulefont{PaxDecide\tneg L!}&
& \final((\tneg\tf{leader} \tand \someoneAll(\tf{leader}\tand\texi\tf{decide}))\tnotor \texi\tf{decide}) 
\figskip
\figskip
\figunderline{Other rules, for correctness \& liveness}
\figskip
\rulefont{LdrExist}&
&\modTF\tf{leader}\tand\modT\someoneAll\tf{leader}
\figskip
\rulefont{LdrExt}_\phi&
& \tf{leader}\timpc\someoneAll(\tf{leader}\tand\phi)\timpc\phi
\quad\text{every closed predicate $\phi$}
\figskip
\rulefont{LdrCorrect}&
& \infinitely(\tf{leader}\timpc\correct{\tf{propose},\tf{write},\tf{decide}}) 
\figskip
\rulefont{PaxPCorrect}&
& 
\begin{array}[t]{l}
\QuorumBox\correct{\tf{propose}} \tand \QuorumBox\correct{\tf{send}} \tand \QuorumBox\correct{\tf{accept}} \tand 
\\
\qquad\QuorumBox\correct{\tf{write}} \tand \QuorumBox\correct{\tf{decide}}
\end{array}
\figskip
\rulefont{PaxPropose01}&
& \texiaffine \tf{propose}
\figskip
\rulefont{PaxWrite01}&
& \texiaffine\tf{write}
\figskip
\rulefont{Pax2Twined}&
& \tall v.\ \QuorumBox\tf{accept}(v) \timpc \CoquorumDiamond\tf{accept}(v)
\end{array}
$$
\caption{$\Theta_\ThyPax$: Axioms of Declarative Paxos (Definition~\ref{defn.declarative.paxos})}
\label{fig.logical.paxos}
\label{fig.ThyPaxOne}
\end{figure}


Recall from Definition~\ref{defn.theory.axiom}(\ref{item.theory}) that a \emph{theory} is a pair of a \emph{signature} and a set of \emph{axioms} over that signature.
Then:
\begin{defn}
\label{defn.declarative.paxos}
\label{defn.ThyPaxOne}
We define \deffont{Declarative Paxos} to be the logical theory 
$$
\ThyPax=(\Sigma_{\ThyPax},\Theta_{\ThyPax}),
$$ 
where:
\begin{enumerate*}
\item\label{item.ThyPaxOne.sig}
the signature $\Sigma_{\ThyPax}$ is 
$$
\Sigma_{\ThyPax}=[\tf{leader}:0;\ \tf{propose},\tf{send},\tf{write},\tf{accept},\tf{decide}:1] ,
$$ 
as per the notation in Definition~\ref{defn.signature}(\ref{item.predicate.syntax.signature}), and
\item\label{item.ThyPaxOne.axioms}
the axioms $\Theta_{\ThyPax}$ are as written in Figure~\ref{fig.logical.paxos}.
\end{enumerate*}
In axioms \rulefont{PaxPropose?} and \rulefont{PaxSend?} we assume an element $\udfn\in\Val$, to model the `undefined' value from Paxos (it does not matter what $\udfn\in\Val$ is; we just have to choose some element and use it consistently).
\end{defn}

\begin{rmrk}
\label{rmrk.functional.notation.in.paxos}
For convenient reference, we give pointers to help unpack the notation used in Figure~\ref{fig.logical.paxos}.
Note that this is just an overview focussing on items of notation used in the axioms, to help with reading them; in-depth discussions of the axioms will follow: 
\begin{enumerate*}
\item
We use the functional notation from Notation~\ref{nttn.hos} to elide quantified variables where convenient.
For example, axioms \rulefont{PaxPropose!} and \rulefont{PaxAccept!} written in full are 
$$
\tf{leader}\tnotor\texi\thea.\tf{propose}(\thea)
\quad\text{and}\quad
\final((\texi\thea.\someoneAll\tf{write}(\thea))\tnotor\texi\thea.\tf{accept}(\thea)) .
$$
\item
$\tnotor$ and $\timpc$ are the \emph{weak} and \emph{strong} implications respectively.
Truth-tables are in Figure~\ref{fig.3} and a discussion is in Subsection~\ref{subsect.3.imp}.
In particular, by Proposition~\ref{prop.mp.for.tnotor}, 
\begin{itemize*}
\item
$\phi\tnotor\phi'$ means `if $\phi$ is true then $\phi'$ is valid' (where `valid' means `$\tvT$ or $\tvB$') and 
\item
$\phi\timpc\phi'$ means `if $\phi$ is true then $\phi'$ is true'. 
\end{itemize*}
\item
The compound modalities --- $\someoneAll$, $\everyoneAll$, $\QuorumBox$, and $\CoquorumDiamond$ --- are discussed in Subsection~\ref{subsect.compound.modalities}.
Their meaning in full generality is in Lemma~\ref{lemm.dense.char}, but because the axioms only apply them to \emph{pointwise} predicates, their meanings simplify to `somewhere' ($\someoneAll$), `everywhere' ($\everyoneAll$), `for a set with a nonempty open interior' ($\QuorumBox$), and `for a dense set' ($\CoquorumDiamond$) respectively, as per Proposition~\ref{prop.pointwise.dense.char}. 
\item\label{item.discuss.final}
$\final$ means $\sometime\forever$ and is from Figure~\ref{fig.3.derived}.
It means `eventually, this holds for all future times'. 

In our axiomatisation this is used to reflect a standard assumption in distributed systems of \deffont{eventual network synchrony}, or synonymously \deffont{global stabilisation (of the network)}, meaning that the communications network may be unstable and fail to pass messages at first, but then after a finite but unbounded amount of time, it becomes stable and messages arrive.
Our logic does not explicitly represent messages, but the use of $\final$ in our forward rules reflects this assumption about the underlying implementation. 
\item
$\texiaffine$ in \rulefont{PaxPropose01} and \rulefont{PaxWrite01} means `for zero or one elements', as per Subsection~\ref{subsect.unique.affine.existence}, so these two axioms intuitively state that each participant proposes or writes at most one value.

Note that our logic is three-valued, so what these axioms mean \emph{precisely} is something a little richer than is the case in two-valued logic.
In particular, if $\Forall{v}\modellabel{\tf{write}(v)}{\acontext}=\tvB$ --- corresponding to the case where a participant is crashed --- then $\acontext\ment\texiaffine\tf{write}$, \emph{even though} $\acontext\ment\tf{write}(v)$ holds for every $v$. 
See the discussion and results in Subsection~\ref{subsect.unique.affine.existence}.
\item
$\correct{\tf{propose}}$ means that $\tf{propose}$ returns $\tvT$ or $\tvF$ (not $\tvB$) on every value; see Definition~\ref{defn.correct}(\ref{item.correct}).
Similarly for the other unary predicate symbols in \rulefont{PaxPCorrect}.
\item
$\mru{\tf{accept}}{v}$ means `a most recently accepted value is $v$'.
See the discussion and definition in Subsection~\ref{subsect.mrup}.
\item
Finally, note that we intend the predicates in Figure~\ref{fig.logical.paxos} to be read as axioms.

What it means for an axiom to be valid in a model --- i.e. how we judge whether a model satisfies $\Theta_\ThyPax$, or how we judge whether a predicate is valid in all models of $\Theta_\ThyPax$ --- is unpacked in Figure~\ref{fig.validity} and Definition~\ref{defn.theory.axiom}(\ref{item.axioms.valid}).
We discuss the quantifier scoping in those definitions in Remark~\ref{rmrk.thy.quantifier.scope}.
\end{enumerate*}
\end{rmrk}




\subsection{High-level discussion of the axioms}

\begin{rmrk}
The predicates in Figure~\ref{fig.logical.paxos} are intended as a \emph{theory}, in the sense of Definition~\ref{defn.theory.axiom}, which is why we call them \emph{axioms} and not just `some predicates'.

That is: we fix a signature $\Sigma$ that has (at least) the predicate symbols required to express the axioms in Figure~\ref{fig.logical.paxos}, and we fix a model $\mathcal M$, and we derive properties that hold in all the valuations $\avaluation$ to $\mathcal M$ such that the theory $\Theta_\ThyPax$ from Figure~\ref{fig.validity} is valid.
 
At least as much design effort went into crafting the axioms $\Theta_\ThyPax$ of Declarative Paxos in Figure~\ref{fig.logical.paxos}, as went into designing the three-valued modal fixedpoint logic in which the axioms are expressed.\footnote{Of course there was design feedback in both directions; we put into the logic what we needed to express the axioms, and we refined the axioms to elegantly work within the logic.} 
Therefore, we will take our time to discuss the axioms in detail.
The reader who just wants to see proofs can skip the rest of this Section (for now) and look to Section~\ref{sect.paxos.correctness.properties} onwards.
\end{rmrk}


\begin{rmrk}
\label{rmrk.axioms.paxos.three.groups}
We have organised the axioms in Figure~\ref{fig.logical.paxos} into three groups, which we now discuss:
\begin{enumerate}
\item
\emph{Backward (inference) rules.}

Axioms with names ending with $?$ represent backwards reasoning and so we call them \deffont{backward rules}.
They use the strong implication $\timpc$ (Definition~\ref{defn.three.imp}(\ref{item.three.elementary.strong.implication})), which means `if the LHS is true ($\tvT$) then the RHS is true ($\tvT$)'.

The backward rules represent a strong inference that \emph{if} some $Y$ is observed to be $\tvT$ then it \emph{must} be because some other property $X$ is $\tvT$.
For instance, \rulefont{PaxAccept?} says that if $\tf{accept}(v)$ is true, then $\someoneAll\tf{write}(v)$ must also be true.

We discuss the backward rules individually in Remark~\ref{rmrk.discuss.backward.rules} below.
\item
\emph{Forward (progress) rules.}

Axioms with names ending with $!$ represent forwards reasoning about making progress, so we call them \deffont{forward rules}.
They use $\final$ and the weak implication $\tnotor$ (Definition~\ref{defn.three.imp}(\ref{item.three.elementary.weak.implication})), which means `if the LHS is true ($\tvT$) then the RHS is valid ($\tvT$ or $\tvB$)'.

We derive the forward rules from the backward rules using a systematic recipe which we discuss in Remark~\ref{rmrk.recipe}.

The forward rules represent a weak inference that \emph{if} some precondition $X$ is $\tvT$ then --- provided the network is stable and provided that the participant has not crashed (where `crashed' is interpreted as returning $\tvB$) --- then the postcondition $Y$ is $\tvT$. 

For instance, \rulefont{PaxAccept!} expresses that, after network synchrony (Remark~\ref{rmrk.functional.notation.in.paxos}(\ref{item.discuss.final})), if $\texi\someoneAll\tf{write}$ is true (by Notation~\ref{nttn.hos} this is shorthand for $\texi a.\someoneAll\tf{write}(a)$) then if the participant has not crashed then $\texi\tf{accept}$ (i.e. $\texi a.\tf{accept}(a)$) is true. 

We discuss the forward rules in Remark~\ref{rmrk.discuss.forward.rules} below.
\item
\emph{Other rules, for correctness and liveness.}

Other rules reflect particular conditions of the algorithms.
We discuss them in Remark~\ref{rmrk.paxos.axioms.discussion}.
\end{enumerate}
\end{rmrk}

\begin{rmrk}[Recipe for deriving forward rules from backward rules]
\label{rmrk.recipe}
A design feature of Figure~\ref{fig.logical.paxos} is that each forward rule is obtained from the corresponding backward rule by a \deffont{recipe}.
Namely: for a backward rule of the form 
$$
\tall v.\ \chi \timpc \phi[a\ssm v] \timpc \psi[a\ssm v]
$$
where 
\begin{itemize*}
\item
$\chi$ expresses some possible \emph{side-conditions} (e.g. $\chi=\tf{leader}$ in \rulefont{PaxDecideL?}) and 
\item
$\phi$ expresses some \emph{preconditions}, and 
\item
$\psi$ expresses some \emph{postconditions}, 
\end{itemize*}
our recipe constructs a forward rule of the form
$$
(\texi v.\,\chi \tand \phi)\tnotor \texi v.\psi .
$$
In most of the forward rules this result also get wrapped in a $\final$ modality, representing network synchrony  (Remark~\ref{rmrk.functional.notation.in.paxos}(\ref{item.discuss.final})).

This recipe holds for every pair of rules:
\begin{itemize*}
\item
We see it cleanly in the pairs \rulefont{PaxAccept?} / \rulefont{PaxAccept!}, \rulefont{PaxDecideL?} / \rulefont{PaxDecideL!}, and \rulefont{PaxDecide\tneg L?} / \rulefont{PaxDecide\tneg L!}.
These rules follow our recipe symbol-for-symbol.
\item
We see this also with \rulefont{PaxPropose?} / \rulefont{PaxPropose!}.
Our recipe generates the following version of \rulefont{PaxPropose!}: 
$$
(\texi v.\,\tf{leader}\tand v\tneq\tf{undef})\timpc\texi v.\tf{propose}(v).
$$
The reader can check that this is logically equivalent to what is written in Figure~\ref{fig.logical.paxos}.
\item 
It is less clear how our recipe works for \rulefont{PaxSend?} and \rulefont{PaxSend!} and \rulefont{PaxWrite?} and \rulefont{PaxWrite!}; e.g. \rulefont{PaxSend!} looks different from \rulefont{PaxSend?} --- it is clearly simpler!
Remarkably, these pairs too are consistent with our recipe; see Remark~\ref{rmrk.why.evaporates} and the subsequent discussion and results.
\end{itemize*}
\end{rmrk}

\begin{rmrk}
We make some further elementary comments on our recipe from Remark~\ref{rmrk.recipe}:
\begin{enumerate*}
\item
We uncurry the forward rules --- meaning that we write $(\texi v.\,\chi \tand \phi)\tnotor \texi v.\psi$ instead of $(\texi v.\,\chi) \tnotor \phi\tnotor \texi v.\psi$ --- because this is convenient for some proofs. 
It makes no other difference and either form would be correct.
\item
A forward rule of the form $\texi v.((\chi \tand \phi)\tnotor \psi)$ would be \emph{wrong}.
This is just a different predicate and is not what is intended.
\item
A forward rule of the form $\tall v.((\chi \tand \phi)\tnotor \psi)$ would not necessarily be wrong, but it is stronger than we need for our proofs.
This reflects a general design principle in logical theories that \emph{weaker axioms are better}.
If we can prove the same result from weaker axioms then generally speaking we should, because that is a stronger result.\footnote{In practice, working from weaker axioms tends to force us to think harder and so design a better-engineered proof.}
\end{enumerate*}
\end{rmrk}


\subsection{More detailed discussion of the axioms}


\begin{rmrk}
\label{rmrk.paxos.axioms.discussion}
We discuss the \emph{axioms for correctness and liveness}:
\begin{enumerate}
\item\label{item.discuss.LdrExist}
\rulefont{LdrExist}:

In implementations of Paxos, at each stage there is a unique \emph{leader}, whose identity is computed by applying (what amounts to) a random number generator to an agreed seed.
As the level of abstraction in this paper, we just assume that a leader exists --- we do not care how it is chosen.

We use a predicate constant symbol $\tf{leader}$ to identify the leader.
Because our logic is modal --- across time $n\in\Time$ and participants $p\in\Pnt$ --- we can just assume that a constant symbol $\tf{leader}$ is assigned truth-value $\tvT$ where a participant $p$ that is a leader at a stage $n$, and is assigned truth-value $\tvF$ where a participant $p$ that is not a leader at a stage $n$.

In more technical language: $\avaluation$ is such that $n,p\mentval \modTF\tf{leader}$ for every time $n\in\Time$ and participant $p\in\Pnt$, %
and $p$ is a leader %
precisely when $n,p\mentval\tf{leader}$.
\rulefont{LdrExist} just asserts that a participant either is a leader or it is not, and some leader exists.
\item\label{item.discuss.LdrExistUniq}
\rulefont{LdrExt}:

\rulefont{LdrExt} is a rule-scheme (one for each closed $\phi$) and asserts that leaders are extensionally equal; %
a formal statement and proof are in Lemma~\ref{lemm.LeaderExtUniq}.

In an implementation this is enforced by having a unique leader, but in the logic this is unnecessary --- we don't care how many participants claim to be leaders, so long as they are all saying the same thing! --- and also it is impossible because (although we have modalities $\everyone$ and $\someone$) we do not have an explicit quantification over elements of $\Pnt$ in our logic.
This is a feature, not a bug: we want our logic to be no more expressive than it needs to be.

In fact, the axiom-scheme \rulefont{LdrExt} is stronger than we need for our proofs;
more on this in Subsection~\ref{subsect.ldrexistuniq'}.


Henceforth we may say `the leader' instead of `a leader', because up to extensional equivalence there is just one.
\item
\rulefont{LdrCorrect}:

\rulefont{LdrCorrect} asserts that for infinitely many times $n\in\Time$, the leader $l$ (which exists by the right-hand conjunct in \rulefont{LdrExist}) is $\tf{propose}$-correct, $\tf{write}$-correct, and $\tf{decide}$-correct (terminology from Definition~\ref{defn.correct}(\ref{item.P-correct}).

Unpacking Definition~\ref{defn.correct}(\ref{item.P-correct}), this means that $n,l\mentval \modTF\tf{propose}(v)$ for every $v\in\Val$, and similarly for $\tf{write}$ and $\tf{decide}$.
Intuitively, this is a \emph{liveness assumption} that there is always a future time when the leader is not crashed in the sense that it can correctly $\tvT\tvF$-propose, $\tvT\tvF$-write, and $\tvT\tvF$-decide a value. 
\item
\rulefont{PaxPCorrect}:

Given a valuation $\avaluation$, we unpack Definition~\ref{defn.validity.judgement}(\ref{item.mentval}) to see what $\mentval \rulefont{PaxPCorrect}$ actually means, as per Definition~\ref{defn.theory.axiom}:
\begin{multline*}
\Forall{n\in\Time}\Exists{O\in\opens}\Forall{p\in O}\Forall{v\in\Val}n,p,O\mentval\correct{\tf P(v)}
\\
\text{for each}\ \tf P\in\{\tf{propose},\tf{send},\tf{accept},\tf{write},\tf{decide}\} .
\end{multline*}
This corresponds to a standard correctness assumption in Paxos that at each stage there is a quorum of uncrashed participants.

However, the axiom \rulefont{PaxPCorrect} is subtly more general than the standard correctness assumption, because there is no restriction that the \emph{same} open set $O$ be used for each unary predicate symbol.\footnote{If we wanted the same $O$, then we would scope the $\QuorumBox$ modality further out and write $\QuorumBox(\correct{\tf{propose}} \tand \correct{\tf{send}} \tand \correct{\tf{accept}} \tand 
\correct{\tf{write}} \tand \correct{\tf{decide}})$.  This turns out to be unnecessary.}
It turns out that the proofs require only this assumption, and no more. 
\item\label{item.01.discussion}
\rulefont{PaxPropose01} and \rulefont{PaxWrite01}:

We justify these axioms with respect to the Paxos algorithm as summarised in Remark~\ref{rmrk.high-level.paxos}:
\begin{itemize*}
\item
Step~\ref{item.l.proposes} of Remark~\ref{rmrk.high-level.paxos} states that leader is crashed or proposes one value.
The intuitive correspondence to \rulefont{PaxPropose01} is clear.
\item
Step~\ref{item.l.writes} of Remark~\ref{rmrk.high-level.paxos} states that the leader writes a value derived from \emph{the first} quorum of participants that it hears back from (or none, if it hears back from no such quorum). 
One corollary of this, is that it writes at most one value.
Our logical theory has no state machine and notion of algorithmic time (as per Remark~\ref{rmrk.elisions}(\ref{item.elide.algorithmic.time})) --- so there is no sense in which information from one open set `arrives' at the leader `before' information from another --- so from the point of view of the logic, of all the open sets available, the leader picks at most one to act on.
This is axiom \rulefont{PaxWrite01}. 
\end{itemize*}
It is a fact that \rulefont{PaxPropose01}+\rulefont{PaxPropose!} and \rulefont{PaxWrite01}+\rulefont{PaxWrite!} give us
$$
\tf{leader}\tnotor\texiunique\tf{propose}
\quad\text{and}\quad
\final((\tf{leader}\tand\QuorumBox\texi\tf{send})\tnotor \texiunique\tf{write}) ,
$$ 
and
\rulefont{PaxPropose01}+\rulefont{PaxPropose?} and \rulefont{PaxWrite01}+\rulefont{PaxWrite?} give us
$$
\tneg\tf{leader}\tnotor\tneg\texi\tf{propose}
\quad\text{and}\quad
\tneg\tf{leader}\tnotor\tneg\texi\tf{write} .
$$
So it follows from our axioms that a leader participant proposes and writes one value (assuming it is not crashed, and if messages arrive), and a non-leader participant proposes and writes zero values. 
\item
\rulefont{Pax2Twined}:

This axiom collects a set of instances of Lemma~\ref{lemm.supertwined.qcbt}(\ref{item.intertwined.3}), namely taking $\phi=\tf{accept}(v)$ in that Lemma for every $v\in\Val$.

There is a design choice here: instead of an axiom \rulefont{Pax2Twined}, we could insist specifically that the semitopology in our underlying model be $2$-twined.
Then \rulefont{Pax2Twined} ceases to be an axiom, but we still have it as a fact of the model by directly invoking Lemma~\ref{lemm.supertwined.qcbt}(\ref{item.intertwined.3}).
Or we could be even more specific: we could insist that the underlying semitopology be $\AllBut{N}{f}$ (Definition~\ref{defn.allbut}) for $N> 2*f$; by Proposition~\ref{prop.allbut.supertwined}(\ref{item.allbut.intertwined}) this is $2$-twined, and then we can use Lemma~\ref{lemm.supertwined.qcbt}(\ref{item.intertwined.3}) to obtain \rulefont{Pax2Twined}.

We choose here to express as an axiom precisely the property that our proofs require.
This is concise, general, and yields clean proofs.
Insisting on a $2$-twined semitopology, or specifically on the semitopology $\AllBut{N}{f}$ for $N>2*f$, would also make sense, but in the context of our logical approach, using an axiom seems nice. %
\end{enumerate}
\end{rmrk}

\begin{rmrk}
\label{rmrk.discuss.backward.rules}
We discuss the \emph{axioms for the backward rules}:
\begin{enumerate}
\item
\rulefont{PaxPropose?}:

If we propose some value $v$ then 
\begin{enumerate}
\item
we are the leader and 
\item
$v$ is not the undefined value $\udfn$.
\end{enumerate}
\item
\rulefont{PaxSend?}:

If we send some value $v$ then 
\begin{enumerate}
\item
somebody must have proposed some value and 
\item
\begin{enumerate*}
\item
\emph{either} $v$ is the value we have most recently accepted, 
\item
\emph{or} we never accepted any value and $v$ is the undefined value.
\end{enumerate*}
\end{enumerate}
\item
\rulefont{PaxWrite?}:

If we write some value $v$ then 
\begin{enumerate}
\item
we are a leader, and
\item
there exists an open set of points (= participants) that have all sent us values, and
\item
\begin{enumerate*}
\item
\emph{either} some point in that open set has most recently accepted $v$, 
\item
\emph{or} 
$v$ is the value that we proposed, and no point in that open set has ever accepted any value. 
\end{enumerate*}
\end{enumerate}
\item
\rulefont{PaxAccept?}:

If we accept $v$ then somebody somewhere must have written $v$.
\item
\rulefont{PaxDecideL?}:

If a leader decides $v$ then an open set of points must have all accepted $v$.
\item
\rulefont{PaxDecide{\tneg}L?}:

If we are not a leader and we decide $v$, then some leader somewhere must have decided $v$. 
\end{enumerate}
\end{rmrk}

\begin{rmrk}
\label{rmrk.discuss.forward.rules}
We discuss the \emph{axioms for the forward rules}.

We noted in Remark~\ref{rmrk.recipe} our design methodology of deriving the forward rules from the backward rules; forward rules are rules for making progress, so they are existentially quantified over the variable $v$ and they use the weak implication $\tnotor$ to emulate possible crashes.

The other subtlety is the $\final$ modality wrapping all of the forward rules, except for \rulefont{PaxPropose!}.
\begin{itemize*}
\item
\emph{What $\final$ does for us in most of the forward rules.}

It models possible network failure, and eventual network synchrony.
As touched on in Remark~\ref{rmrk.functional.notation.in.paxos}(\ref{item.discuss.final})), the $\final$ modality reflects the possibility that the network might be down for a while, but eventually --- `finally' --- the network stabilises and messages arrive.

Predicates in our logic are \emph{not} messages, but we get the effect of eventual network synchrony in our logic (even though messages are not explicitly present) by wrapping the implications in our forward rules in the $\final$ modality.

More on this when we introduce the notion of \emph{global stabilisation of logical time} in Subsection~\ref{subsect.gslt}.
\item
\emph{Why $\final$ is absent from \rulefont{PaxPropose!}.}

Because this rule models a property of the implementation that involves no network communication.
The leader does not need to communicate with any other participant to pick a value to propose.

It is true that the leader's propose message might not arrive if the network is down --- but this is reflected by the use of $\final$ in the \rulefont{PaxSend!} rule.
\end{itemize*}
\end{rmrk}


\subsection{Some asides}

\subsubsection{An aside on \rulefont{PaxWrite?}}

\rulefont{PaxWrite?} contains this predicate:
$$
\mru{(\someone\tf{accept})}{v} .
$$
Unpacking what this means (cf. Remark~\ref{rmrk.explain.mrup}), this looks into the past to check whether $v$ is a value that was most recently accepted by some participant.

Note that this slightly differently-scoped predicate would be \emph{wrong}:
$$
\someone(\mru{\tf{accept}}{v}) .
$$
This finds some participant such that $v$ is a value that was most recently accepted by \emph{that} participant (but there might be another participant who more recently accepted some other value).

The reader with a background in logic will know all about how important scoping can be (e.g. $\forall x.\exists y.\phi$ vs. $\exists x.\forall y.\phi$, and $\exists x.(\phi\land \phi')$ vs. $\exists x.\phi\land\exists x.\phi'$), but perhaps the point bears repeating.


\subsubsection{An aside on \rulefont{LdrExt}}

In a Paxos implementation, there is a unique leader at each stage. 
We mentioned in Remark~\ref{rmrk.paxos.axioms.discussion}(\ref{item.discuss.LdrExist}\&\ref{item.discuss.LdrExistUniq}) that \rulefont{LdrExist} asserts that a leader exists and that the \rulefont{LdrExt} axiom-scheme asserts that all leader participants are extensionally equal.
We now make this formal:\footnote{Later on in Subsection~\ref{subsect.ldrexistuniq'} we will note weaker, but more abstract, forms of the axiom.}
\begin{lemm}[Leaders are extensionally equivalent]
\label{lemm.LeaderExtUniq}
Suppose $n\in\Time$ and $\avaluation$ is a valuation, and suppose $n\mentval\rulefont{LdrExt}_\phi$ holds for every closed predicate $\phi$.\footnote{By Definition~\ref{defn.validity.judgement}(\ref{item.ment.n}) this means $\Forall{p\in\Pnt}\Forall{O\in\opensne}(\modellabel{\rulefont{LdrExt}_\phi}{n,p,O,\avaluation}\in\{\tvT,\tvB\})$.} 
Suppose that $l,l'\in\Pnt$ and $n,l\mentval\modT\tf{leader}$ and $n,l'\mentval\modT\tf{leader}$.

Then $l$ and $l'$ are \deffont{extensionally equivalent} at $(n,\avaluation)$, by which we mean that for every closed predicate $\phi$ and every $O\in\opensne$,
$$
\modellabel{\phi}{n,l,O,\avaluation} = \modellabel{\phi}{n,l',O,\avaluation} .
$$
\end{lemm}
\begin{proof}
Suppose $\phi$ is a closed predicate and $O\in\opensne$, and suppose for the sake of argument that $\modellabel{\phi}{n,l,O,\avaluation}=\tvB$.

It is routine to check from Figures~\ref{fig.3.derived} and~\ref{fig.3} and from Lemma~\ref{lemm.dense.char}(\ref{item.dense.char.someoneAll}) that $\modellabel{\someoneAll\modB\phi}{n,l,O,\avaluation}=\tvT$. 
By assumption $n,l',O\mentval \rulefont{LdrExt}_{\modB\phi}$, and by two applications of \strongmodusponens (once for $\modT\tf{leader}$, and another for $\modT\someoneAll(\tf{leader}\tand\modB\phi)$) we conclude that $\modellabel{\modB\phi}{n,l',O,\avaluation}=\tvT$, and therefore $\modellabel{\phi}{n,l',O,\avaluation}=\tvB$ as required. 
 
The cases $\modellabel{\phi}{n,l,O,\avaluation}=\tvT$ and $\modellabel{\phi}{n,l,O,\avaluation}=\tvF$ are be precisely similar, and this suffices to prove the result.
\end{proof}


\subsubsection{An aside on \rulefont{PaxSend?} and \rulefont{PaxSend!}}

\begin{rmrk}
\label{rmrk.why.evaporates}
In Remark~\ref{rmrk.recipe} we proposed a recipe for deriving a forward rule from a backward one: a backward rule of the form
$$
\rulefont{Rule?} 
\qquad
\tall v.\ \chi \timpc \phi[a\ssm v] \timpc \psi[a\ssm v]
$$
generates a forward rule (possibly wrapped in a $\final$ modality) of the form
$$
\rulefont{Rule!}
\qquad
(\texi v.\, \chi \tand \phi)\tnotor \texi v.\psi .
$$
We noted in Remark~\ref{rmrk.recipe} that the recipe works cleanly for some rules --- but it was less clear how it can be applied to derive \rulefont{PaxSend!} from \rulefont{PaxSend?}, and \rulefont{PaxWrite!} from \rulefont{PaxWrite?}.
In these cases, the left-hand sides of \rulefont{PaxSend!} and \rulefont{PaxWrite!} seem much simpler than the right-hand sides of \rulefont{PaxSend?} and \rulefont{PaxWrite?} respectively. 

We now show why, in fact, \rulefont{PaxSend!} and \rulefont{PaxWrite!} are correct and make sense in the context of our recipe, culminating with Corollary~\ref{corr.back.ok}.
\end{rmrk}


\begin{lemm}
\label{lemm.evaporates.simplify}
We have:
\begin{enumerate*}
\item
$\ment \texi v.\bigl((\mru{\tf{accept}}{v})\ \tor (\tneg\modT\recent\texi\tf{accept}\tand v\teq\udfn)\bigr) \ \equiv\ \tvT$.
\item
$\ment \texi v.\bigl(\ (\mru{(\someone\tf{accept})}{v})\ \tor\ (\tneg\modT\recent\texi\someone\tf{accept}\tand \tf{propose}(v))\ \bigr)\ \equiv\ \texi\tf{propose}$.
\end{enumerate*}
Note that:
\begin{itemize*}
\item
The notation $\texi\tf{accept}$ and $\texi\tf{propose}$ and $\texi\someone\tf{accept}$ is from Notation~\ref{nttn.hos}.
\item
We do not restrict to contexts and valuations that satisfy \ThyPax; the meaning of the $\ment$ symbol above is as in Definition~\ref{defn.validity.judgement}(\ref{item.ment.phi}).
\end{itemize*}
\end{lemm}
\begin{proof}
We consider each part in turn:
\begin{enumerate}
\item
We distribute the $\texi v$ through the $\tor$, and then through the $\tand$ on the right-hand side (we can do this because $v$ is not free in the left-hand conjunct) and we see that it suffices to show 
$$
\ment \texi v.(\mru{\tf{accept}}{v})\ \tor\ (\tneg\modT\recent\texi\tf{accept}\tand \texi v.v\teq\udfn)\ \equiv\ \tvT
$$
Using Lemma~\ref{lemm.evaporates} this simplifies to just
$$
\ment \modT\recent\texi\tf{accept}\ \tor\ (\tneg\modT\recent\texi\tf{accept}\tand \texi v.v\teq\udfn) \ \equiv\ \tvT.
$$
This clearly holds, just taking $v\teq\udfn$.
\item
We distribute $\texi v$ into the predicate as we did for the proof of part~1.
We see that it suffices to show 
$$
\ment \bigl(\ \texi v.(\mru{(\someone\tf{accept})}{v})\ \tor\ (\tneg\modT\recent\texi\someone\tf{accept}\tand \texi\tf{propose})\ \bigr)\ \equiv\ \texi\tf{propose}.
$$
We simplify using Lemma~\ref{lemm.evaporates} to
$\ment \texi\tf{propose}\ \equiv\ \texi\tf{propose}$,
which is true.\footnote{Yes, we have a lemma for that: Lemma~\ref{lemm.equiv}(\ref{item.equiv.2}) and reflexivity of equality.}
\qedhere
\end{enumerate}
\end{proof}

We can now justify the left-hand sides of \rulefont{PaxSend!} and \rulefont{PaxWrite!} as being consistent with our recipe as discussed above in Remark~\ref{rmrk.why.evaporates}:
\begin{corr}
\label{corr.back.ok}
We have:
$$
\begin{array}{r@{\ }l}
\ment\texi v.&\bigl(\someoneAll\texi\tf{propose}\ \tand \bigl((\mru{\tf{accept}}{v})\ \tor (\tneg\modT\recent\texi\tf{accept}\tand v\teq\udfn)\bigr)\bigr) 
\\
&\equiv \someoneAll\texi\tf{propose}
\\[2ex]
\ment\texi v.&\bigl(\tf{leader}\tand 
\Quorum(\Box\texi\tf{send} \tand 
\bigl((\mru{(\someone\tf{accept})}{v}) \tor
(\tneg\modT\recent\texi\someone\tf{accept}\tand \tf{propose}(v))\bigr)) 
\bigr)
\\
&\equiv
(\tf{leader}\tand \QuorumBox\texi\tf{send} \tand \texi\tf{propose})
\end{array}
$$
\end{corr}
\begin{proof}
By elementary predicate manipulations, distributing the $\texi v$ quantifier down into the predicate and using Lemma~\ref{lemm.evaporates.simplify}.
\end{proof}

 




\section{Correctness properties}
\label{sect.paxos.correctness.properties}

\subsection{Summary of the properties}

\begin{rmrk}
\label{rmrk.standard.paxos.correctness}
We will derive the following standard correctness properties for Paxos (as laid out in~\cite[Subsection~2.2.1]{cachin2011paxos}) from the theory \ThyPax in Figure~\ref{fig.logical.paxos}:
\begin{enumerate*}
\item
\textbf{Validity:} \quad If a participant decides some value $v$, then $v$ was proposed by some participant.

A formal statement and proof is in Theorem~\ref{thrm.validity}.
\item
\textbf{Agreement:}\quad 
If two participants decide values, then those values are equal.

A formal statement and proof is in Theorem~\ref{thrm.full.agreement}.
\item
\textbf{Termination:}\quad 
Every correct participant eventually decides some value.

A formal statement and proof is in Theorem~\ref{thrm.termination}.
\item
\textbf{Integrity:}\quad 
Every correct (non-crashing) participant decides at most once.

Our logic abstracts away from a direct notion of execution, by design, but this property is represented indirectly in how Theorem~\ref{thrm.termination} is expressed.
See Remark~\ref{rmrk.paxos.integrity} below. 
\end{enumerate*}
\end{rmrk}

\begin{rmrk}[On integrity]
\label{rmrk.paxos.integrity}
In implementations of Paxos, once a participant has decided, it stops.
But in our logic, `stop' is not even a concept.
There is no abstract machine to stop; there are just predicates, with denotations.

Nevertheless, the integrity property is echoed in our maths, as the \emph{infinite repetition} in Theorem~\ref{thrm.termination}.
Once a predicate $\tf{decide}(v)$ returns $\tvT$, it does so \emph{forever} --- and because of Agreement, it never decides any other value.
A predicate eventually always returning the same truth-value is a model of an abstract machine that has computed its result, and then just sits there forever, displaying that same result, unchanged, every time we query its display.
\end{rmrk}

\subsection{Validity}


\begin{rmrk}
\label{rmrk.validity.in.english}
\emph{Validity} is expressed in English as:
\begin{quote}
\emph{If a participant decides some value $v$, then $v$ was proposed by some participant.}
\end{quote}
We render this in our logic in Theorem~\ref{thrm.full.agreement} as
$$
\ThyPax\ment\tall v.(\tf{decide}(v) \timpc \urecent\someoneAll(\tf{leader}\tand\tf{propose}(v))).
$$

In words we can read this as: 
\begin{quote}
If someone decides $v$ now, then at some stage now or in the past a leader of that stage proposed $v$. 
\end{quote}
We will now set about proving this property from our axioms, culminating with Theorem~\ref{thrm.validity}.
\end{rmrk}


\begin{lemm}
\label{lemm.validity.helper}
Suppose $\avaluation$ is a valuation such that $\mentval\ThyPax$ ($\mentval\ThyPax$ is from Figure~\ref{fig.validity}).
Suppose $n\in\Time$ and $v\in\Val$.
Then:
\begin{enumerate*}
\item\label{item.validity.helper.3}
$n\mentval \someoneAll\tf{write}(v) \timpc \urecent\someoneAll\tf{propose}(v)$. 
\item\label{item.validity.helper.2}
$n\mentval \someoneAll\tf{accept}(v) \timpc \urecent\someoneAll\tf{propose}(v)$. 
\end{enumerate*}
\end{lemm}
\begin{proof}
We prove both parts by a simultaneous induction on $n$.
We freely use \strongmodusponens:
\begin{enumerate}
\item
Suppose $n\mentval\modT\someoneAll\tf{write}(v)$.
Then using \rulefont{PaxWrite?} %
$$
n\mentval \modT\Quorum(\mru{\someone\tf{accept}}{v}) 
\quad\text{or}\quad
n\mentval \modT\someoneAll\tf{propose}(v) .
$$
We consider the two cases:
\begin{itemize*}
\item
\emph{If $n\mentval \modT\someoneAll\tf{propose}(v)$,}\ then 
$n\mentval\modT\urecent\someoneAll\tf{propose}(v)$ follows immediately from the clause for %
$\urecent$ in Figure~\ref{fig.3.derived}.
\item
\emph{If $n\mentval \modT\Quorum(\mru{\someone\tf{accept}}{v})$,}\ 
then by Lemma~\ref{lemm.mru.to.recent}(\ref{item.mru.to.recent.1}) $n\mentval\modT\recent\someoneAll\tf{accept}(v)$,
so that by Figure~\ref{fig.3.derived} $n'\mentval\modT\someoneAll\tf{accept}(v)$ for some $n'<n$.
By part~\ref{item.validity.helper.2} of this result for $n'<n$, we have that $n'\mentval\modT\urecent\someoneAll\tf{propose}(v)$, and thus using Lemma~\ref{lemm.tomorrow.forever.recent}(\ref{item.rur.r}) $n\mentval\modT\urecent\someoneAll\tf{propose}(v)$. 
\end{itemize*}
In either case, we are done.
\item
Suppose $n\mentval\modT\someoneAll\tf{accept}(v)$.
By \rulefont{PaxAccept?} $n\mentval\modT\someoneAll\tf{write}(v)$.
We use part~\ref{item.validity.helper.3} of the inductive hypothesis for $n$. 
\qedhere\end{enumerate}
\end{proof}



\begin{thrm}[Validity]
\label{thrm.validity}
We have:
$$
\ThyPax\ment\tall v.(\tf{decide}(v) \timpc \urecent\someoneAll(\tf{leader}\tand\tf{propose}(v))) .
$$
\end{thrm}
\begin{proof}
Suppose $\avaluation$ is a valuation such that $\mentval\ThyPax$. 
We reason using \strongmodusponens. 
Suppose $v\in\Val$ and $n\in\Time$ and $p\in\Pnt$, and suppose $n,p\mentval\modT\tf{decide}(v)$.
By \rulefont{PaxDecideL?} (and \rulefont{PaxDecide\tneg L?} if $p\neq\f{leader}(n)$), $n\mentval\modT\QuorumBox\tf{accept}(v)$.
It follows by Proposition~\ref{prop.someone.implies.someoneAll}(\ref{item.someone.implies.someoneAll.2}) that $n\mentval\modT\someoneAll\tf{accept}(v)$.
By Lemma~\ref{lemm.validity.helper}(\ref{item.validity.helper.2}) $n\mentval\modT\urecent\someoneAll\tf{propose}(v)$, and using \rulefont{PaxPropose?} we conclude that $n\mentval\modT\urecent\someoneAll(\tf{leader}\tand\tf{propose}(v))$ as required. 
\end{proof}




\subsection{Agreement}

\begin{rmrk}
\emph{Agreement} is expressed in English as:
\begin{quote}
\emph{If two participants decide values, then those values are equal.}
\end{quote}
We render this in our logic in Theorem~\ref{thrm.full.agreement} as
$$
\ThyPax\ment\tall v,v'.(\urecent\someoneAll\tf{decide}(v) \timpc \someoneAll\tf{decide}(v') \timpc v'\teq v) .
$$
In words we can read this as: 
\begin{quote}
If someone decides $v$ now or in the past, and someone decides $v'$ now, then $v'=v$.
\end{quote}
It is convenient to prove the `now' and the `in the past' parts of this as two sub-cases: 
\begin{enumerate*}
\item
Proposition~\ref{prop.decide.n.agree} handles the case of $\someoneAll\tf{decide}(v)\tand\someoneAll\tf{decide}(v')$, i.e. when the two decisions are made at the same time.
\item
Theorem~\ref{thrm.full.agreement} then extends this with the case of $\recent\someoneAll\tf{decide}(v)\tand\someoneAll\tf{decide}(v')$, i.e. when one decision precedes the other.
\end{enumerate*}
\end{rmrk}


\subsubsection{Agreement at time $n$}



\begin{lemm}
\label{lemm.various.unique}
Suppose $\avaluation$ is a valuation such that $\mentval\ThyPax$. 
Suppose $n\in\Time$ and $v,v'\in\Val$. 
Then:
\begin{enumerate*}
\item\label{item.various.unique.write}
$n\mentval (\someoneAll\tf{write}(v)\tand \someoneAll\tf{write}(v')) \timpc v\teq v'$. 
\item\label{item.various.unique.accept}
$n\mentval (\someoneAll\tf{accept}(v)\tand \someoneAll\tf{accept}(v')) \timpc v\teq v'$. 
\item\label{item.various.unique.recent.accept}
$n\mentval \bigl(\mru{(\someoneAll\tf{accept})}{v}\tand \mru{(\someoneAll\tf{accept})}{v'}\bigr) \timpc v\teq v'$.
\end{enumerate*}
\end{lemm}
\begin{proof}
We consider each part in turn, using \strongmodusponens (we may also silently commute the $\modT$ modality as convenient):
\begin{enumerate}
\item
Suppose $n\mentval\modT\someoneAll\tf{write}(v)\tand\modT\someoneAll\tf{write}(v')$.
Using \rulefont{PaxWrite?} 
$n\mentval\modT\someoneAll(\tf{write}(v)\tand\tf{leader})\tand\modT\someoneAll(\tf{write}(v')\tand\tf{leader})$. 
Using \rulefont{LdrExt} therefore $n\mentval\modT\someoneAll(\tf{write}(v)\tand\tf{write}(v')\tand\tf{leader})$.
By \rulefont{PaxWrite01} and Lemma~\ref{lemm.unique.affine.existence}(\ref{item.unique.affine.existence.01.implies}) we have $v=v'$ as required.
\item
Suppose $n\mentval \modT\someoneAll\tf{accept}(v)\tand\modT\someoneAll\tf{accept}(v')$.
By \rulefont{PaxAccept?} $n\mentval\modT\someoneAll\tf{write}(v)\tand\modT\someoneAll\tf{write}(v')$.
We use part~\ref{item.various.unique.write} of this result.
\item
Suppose $n\mentval\modT(\mru{(\someoneAll\tf{accept})}{v})\tand\modT(\mru{(\someoneAll\tf{accept})}{v'}\bigr)$.\footnote{By Lemma~\ref{lemm.mru.unambivalent}(2) we do not need to write the $\modT$ here but we do anyway; see Remark~\ref{rmrk.explain.mru}(\ref{item.explain.mru.T}).}
Using Lemma~\ref{lemm.mru.to.recent}(\ref{item.mru.to.recent.2}) (taking $a.\phi=\someoneAll\tf{accept}$ in that Lemma) there exists $n'<n$ such that 
$n'\mentval\modT\someoneAll\tf{accept}(v)\tand\modT\someoneAll\tf{accept}(v')$.
We use part~\ref{item.various.unique.accept} of this result.
\qedhere\end{enumerate}
\end{proof}



\begin{prop}
\label{prop.decide.n.agree}
Suppose $\avaluation$ is a valuation such that $\mentval\ThyPax$. 
Suppose $n\in\Time$ and $v,v'\in\Val$.
Then
$$
n\mentval (\someoneAll\tf{decide}(v) \tand \someoneAll\tf{decide}(v')) \timpc v\teq v' . 
$$
\end{prop}
\begin{proof}
We use \strongmodusponens.
Suppose 
$n\mentval \modT(\someoneAll\tf{decide}(v)\tand\someoneAll\tf{decide}(v'))$.
Using \rulefont{PaxDecide\tneg L?} if needed, and then \rulefont{PaxDecideL?} twice, we have
$n\mentval \modT(\QuorumBox\tf{accept}(v)\tand\QuorumBox\tf{accept}(v'))$ and thus by Proposition~\ref{prop.someone.implies.someoneAll}(\ref{item.someone.implies.someoneAll.2}) and Lemma~\ref{lemm.commuting.connectives}
also $n\mentval \modT(\someoneAll\tf{accept}(v)\tand\someoneAll\tf{accept}(v'))$.
We can now use Lemma~\ref{lemm.various.unique}(\ref{item.various.unique.accept}).
\end{proof}



\begin{rmrk}
Proposition~\ref{prop.decide.n.agree} proves part of the agreement property; that no two correct participants decide differently \emph{at a given time}.
To prove the result across different times will require a bit more work, culminating with Theorem~\ref{thrm.full.agreement}.
\end{rmrk}


\subsubsection{Some fine structure of the rules}
\label{subsect.ldrexistuniq'}

We continue a discussion from Remark~\ref{rmrk.paxos.axioms.discussion}(\ref{item.discuss.LdrExistUniq}) and also Lemma~\ref{lemm.LeaderExtUniq}:

Note that we use the axiom-scheme \rulefont{LdrExt} from Figure~\ref{fig.logical.paxos} just once: in the proof of Lemma~\ref{lemm.various.unique}(\ref{item.various.unique.write}).

Examining how we use this axiom-scheme in the proof above, we see that for our specific purposes it would suffice to assume a single axiom as follows:
$$
\hspace{-6em}\rulefont{LdrExt'}
\qquad\qquad
\tf{leader}\timpc \someoneAll(\tf{leader}\tand\tf{write}(v))\timpc \tf{write}(v) .
$$ 
The axiom scheme \rulefont{LdrExt} asserts that there is a unique leader at each stage, up to extensional equivalence, which is a fair reflection of how the algorithms work where there \emph{is} a unique leader at each stage.

\rulefont{LdrExt'} asserts that there is a unique leader at each stage, up to extensional equivalence \emph{for writing values}.
In the proofs, this turns out to be all we need.

In fact, we can simplify this just to 
$$
\hspace{-6em}\rulefont{SPaxWrite01}
\qquad\qquad
\texiaffine\someoneAll\tf{write} .
$$
In full, this is $\texiaffine v.\someoneAll\tf{write}(v)$, and in words it expresses: 
\begin{quote}
`at each stage, at most one value is written'.
\end{quote}
In the presence of \rulefont{PaxWrite?}, \rulefont{SPaxWrite01} is implied by \rulefont{LdrExt'} --- but it does not imply \rulefont{LdrExt'}, because it permits the possibility that one leader participant writes a value and another leader participant crashes.
The reader can check that \rulefont{SPaxWrite01} is sufficient to complete the proof of Lemma~\ref{lemm.various.unique}(\ref{item.various.unique.write}).

Note also that that \rulefont{SPaxWrite01} implies \rulefont{PaxWrite01}, so that --- for the proofs we care about --- \rulefont{SPaxWrite01} can replace both \rulefont{LdrExt} and \rulefont{PaxWrite01}. 
We will do exactly this, later: the job of $\Theta_\ThyPax$ is to be `declarative Paxos', that is, its purpose is to be declarative \emph{and} to be recognisably Paxos. 
Later on in Figure~\ref{fig.ThySPax} we will define a \emph{simpler} declarative Paxos in (whence the name `SPaxWrite01').
This simplified theory is derived by looking at the proofs here, and seeing what we actually require.










 



\subsubsection{Agreement at possibly different times}



\begin{prop}
\label{prop.integrity.helper}
Suppose $\avaluation$ is a valuation such that $\mentval\ThyPax$, and suppose $v,v'\in\Val$ and $\avaluation$. 
Then:
\begin{enumerate*}
\item\label{item.integrity.helper.accept}
$\mentval\recent\QuorumBox\tf{accept}(v) \timpc \someoneAll\tf{accept}(v')\timpc v'\teq v$.
\item\label{item.integrity.helper.write}
$\mentval\recent\QuorumBox\tf{accept}(v) \timpc \someoneAll\tf{write}(v')\timpc v'\teq v$.
\end{enumerate*}
\end{prop}
\begin{proof}
Suppose $n\in\Time$.
We will reason using \strongmodusponens, so suppose 
\begin{equation}
\label{eq.n.1}
n\mentval\modT\QuorumBox\tf{accept}(v). 
\end{equation}
We will prove by induction on $n'>n$ that for every $v'\in\Val$,
$$
\begin{array}{r@{\ }l}
n'\mentval&\modT(\mru{(\someoneAll\tf{accept})}{v'}\timpc v'\teq v)
\quad\text{and}\quad
\\
n'\mentval&\modT(\someoneAll\tf{write}(v')\timpc v'\teq v)
\end{array}
$$
Clearly, by \strongmodusponens the $\tf{write}$ part of the inductive hypothesis is enough to prove part~\ref{item.integrity.helper.write} of this result.
For the $\tf{accept}$ part of the inductive hypothesis, we see that this implies part~\ref{item.integrity.helper.accept} of this result either by a routine argument using \rulefont{PaxAccept?} and part~\ref{item.integrity.helper.write} of this result --- or using Corollary~\ref{corr.mru.today.to.yesterday}(\ref{item.mru.today.implies.tomorrow}) (for $a.\phi=a.\someoneAll\tf{accept}(a)$ in that result).

We consider cases:
\begin{itemize}
\item
\emph{The base case for $\tf{accept}$, where $n'=n+1$.}

Suppose $n'\mentval\modT(\mru{\someoneAll\tf{accept}}{v'})$.

We assumed in equation~\eqref{eq.n.1} that $n\mentval\modT\QuorumBox\tf{accept}(v)$, so by Proposition~\ref{prop.someone.implies.someoneAll}(\ref{item.someone.implies.someoneAll.2}) also $n\mentval\modT\someoneAll\tf{accept}(v)$, and by Corollary~\ref{corr.mru.today.to.yesterday}(\ref{item.mru.today.implies.tomorrow}) $n'\mentval\modT(\mru{(\someoneAll\tf{accept})}{v})$. 
By Lemma~\ref{lemm.various.unique}(\ref{item.various.unique.recent.accept}) $v'=v$.

\item
\emph{The inductive step for $\tf{accept}$, where $n'>n+1$.}

Suppose $n'\mentval\modT(\mru{\someoneAll\tf{accept}}{v'})$.
By Corollary~\ref{corr.mru.today.to.yesterday}(\ref{item.mru.tomorrow.options}) either $n'\minus 1\mentval\modT\someoneAll\tf{accept}(v')$ or $n'\minus 1\mentval\modT(\mru{\someoneAll\tf{accept}}{v'})$.

In the former case, using \rulefont{PaxAccept?} $n'\minus 1\mentval\someoneAll\tf{write}(v')$, and by the $\tf{write}$ part of the inductive hypothesis, $v'=v$.
In the latter case, we just use the $\tf{accept}$ part of the inductive hypothesis. 
\item
\emph{The case of $\tf{write}$.}

Suppose $p\in\Pnt$ and $n',p\mentval\modT\tf{write}(v')$.
Using \rulefont{PaxWrite?}, 
there exists an $O\in\opensne$ such that 
$$
n',O\mentval \modT\bigl( 
(\mru{\someone\tf{accept}}{v})\ \tor
(\tneg\modT\recent\texi\someone\tf{accept}\tand \tf{propose}(v)) \bigr),
$$
from which it follows using Lemmas~\ref{lemm.commuting.connectives} and~\ref{lemm.tand.tor} and some basic reasoning on truth-values that 
$$
\text{either}\quad
n',O\mentval \modT\mru{\someone\tf{accept}}{v}
\quad\text{or}\quad
n',O\mentval \tneg\modT\recent\texi\someone\tf{accept} .
$$
Now $n<n'$ and we assumed in equation~\eqref{eq.n.1} that $n\mentval\modT\QuorumBox\tf{accept}(v)$, so using 
\rulefont{Pax2Twined} and \strongmodusponens $n\mentval\modT\CoquorumDiamond\tf{accept}(v)$, and in particular (by routine reasoning on the clauses for $\CoquorumDiamond$ in Figure~\ref{fig.3.derived} and for $\someone$ in Figure~\ref{fig.3.phi.f})
$n,O\mentval\modT\someone\tf{accept}(v)$.
Thus $n',O\mentval\modT\texi\someone\tf{accept}$ and by Lemma~\ref{lemm.tomorrow.forever.recent}(\ref{item.tomorrow.forever.recent.1}) (taking $\phi=\texi\someone\tf{accept}$ in that Lemma) 
$n',O\mentval\modT\recent\texi\someone\tf{accept}$.
It follows that we are not in the right-hand (`or') branch above, so we must be in the left-hand (`either') branch, and thus $n',O\mentval\modT\mru{\someone\tf{accept}}{v}$.

By the $\tf{accept}$ part of the inductive hypothesis, $v'=v$.
\qedhere\end{itemize}

\end{proof}


\begin{thrm}[Agreement]
\label{thrm.full.agreement}
We have:
$$
\ThyPax\ment\tall v,v'.(\urecent\someoneAll\tf{decide}(v) \timpc \someoneAll\tf{decide}(v') \timpc v'\teq v).
$$
\end{thrm}
\begin{proof}
Suppose $v,v'\in\Val$ and $n\in\Time$ and $\avaluation$ is a valuation.
We use \strongmodusponens, so suppose
$$
n\mentval\modT\urecent\someoneAll\tf{decide}(v)
\quad\text{and}\quad
n\mentval\modT\someoneAll\tf{decide}(v') .
$$
If $n\mentval\modT\someoneAll\tf{decide}(v)$ then we use Proposition~\ref{prop.decide.n.agree}.

So suppose $n\mentval\modT\recent\someoneAll\tf{decide}(v)$.
By \rulefont{PaxDecideL?} (and \rulefont{PaxDecide\tneg L?} if needed), $n\mentval\modT\recent\QuorumBox\tf{accept}(v)$ and $n\mentval\modT\QuorumBox\tf{accept}(v')$, so that by Proposition~\ref{prop.someone.implies.someoneAll}(\ref{item.someone.implies.someoneAll.2}) also %
$n\mentval\modT\someoneAll\tf{accept}(v')$.
By Proposition~\ref{prop.integrity.helper}(\ref{item.integrity.helper.accept}) $v'=v$ as required.
\end{proof}

\subsection{Termination}
\label{subsect.pax.termination}

\begin{rmrk}
\emph{Termination} is expressed in English as:
\begin{quote}
\emph{Every correct participant eventually decides some value.}
\end{quote}
We render this in our logic in Theorem~\ref{thrm.termination} as 
$$
\ThyPax\ment\infinitely\everyoneAll\texi v.\tf{decide}(v).
$$
In words: 
\begin{quote}
There is always some future time when all participants decide.
\end{quote}
Note that there is no need to explicitly specify `all correct/uncrashed participants' above.
This is taken care of by our logic's in-built notion of validity, since incorrect participants return $\tvB$, which is valid as per Remark~\ref{rmrk.motivate.validity}(\ref{item.tb.valid}) (but not true or correct; cf. Remark~\ref{rmrk.motivate.validity}).
\end{rmrk}


\subsubsection{Global stabilisation of logical time (GSLT)}
\label{subsect.gslt}

Our forward rules are wrapped in a $\final$ modality, to reflect an assumption that while there may be some initial network instability, if we wait long enough the network will stabilise (see the discussion in Remark~\ref{rmrk.discuss.forward.rules}).
This leads us to the notion of \emph{global stabilisation of logical time}, or \emph{GSLT} for short: 
\begin{defn}[GSLT]
\label{defn.gslt}
Suppose $\phi$ is a closed predicate and $\avaluation$ is a valuation.
Then:
\begin{enumerate*}
\item\label{item.gslt.synchronous}
Say that \ $\final\phi$\ (i.e. $\phi$ under a $\final$ modality) is \deffont{synchronous at $n$ in $\avaluation$} when $n\mentval\forever\phi$.\footnote{%
By Definition~\ref{defn.validity.judgement}, $n\ment\final\phi$ means that $n,p,O\mentval\final\phi$ for every $p\in\Pnt$ and $O\in\opens$ and valuation $\varsigma$.} %
That is:
$$
\final\phi\ \ \text{ is synchronous at $n$ in $\avaluation$}
\quad\text{when}\quad
n\mentval\forever\phi .
$$
In words, $\final\phi$ is synchronous when the `$\final$' becomes a `$\forever$'.
\item\label{item.gslt.gslt}
Define $\gslt(\avaluation)\in\Time$ to be the least stage $n$ such that all of the rules in Figure~\ref{fig.3.phi.f} that start with $\final$ --- namely, the forward rules \rulefont{PaxSend!}, \rulefont{PaxAccept!}, \rulefont{PaxDecideL!}, and \rulefont{PaxDecide\tneg L!} --- are synchronous at $n$ in the valuation $\avaluation$.

We call $\gslt(\avaluation)$ the \deffont{Global Stabilisation of Logical Time} in $\avaluation$. 
\end{enumerate*}
\end{defn}

\begin{rmrk}
We make some observations about Definition~\ref{defn.gslt}:
\begin{enumerate}
\item
Unpacking Definition~\ref{defn.gslt}(\ref{item.gslt.synchronous}) using the clause for $\forever$ in Figure~\ref{fig.3.derived}, 
$$
\final\phi \text{ is synchronous at $n$ in $\avaluation$}
\quad\liff\quad
n\mentval\forever\phi 
\quad\liff\quad
\Forall{n'{\geq}n}(n'\mentval\phi),
$$
\noindent
and $\gslt(\avaluation)$ is least such that $\gslt(\avaluation)\mentval\forever(\someoneAll\texi\tf{propose} \tnotor \texi\tf{send})$, and $\gslt(\avaluation)\mentval\forever\bigl((\tf{leader}\tand \QuorumBox\texi\tf{send}) \tnotor \texi\tf{write}\bigr)$, and so on.
\item  
$\gslt$ is well-defined in Definition~\ref{defn.gslt}(\ref{item.gslt.gslt}) because there are finitely many axioms, mentioning finitely many $\final$ modalities, and a finite set of natural numbers has a finite least upper bound.
This statement has real content, e.g. if we had written \rulefont{PaxAccept!} as $\tall v.\final(\someoneAll\tf{write}(v) \tnotor \tf{accept}(v))$ --- which is a subtly weaker assertion than the form used in Figure~\ref{fig.logical.paxos} --- then $\gslt$ would \emph{not} be defined. 
\end{enumerate}
\end{rmrk}




\begin{lemm}
\label{lemm.someone.writes.after.gslt.accept}
Suppose $\avaluation$ is a valuation such that $\mentval\ThyPax$. 
Suppose $n\in\Time$ and $n\geq\gslt(\avaluation)$ and $v\in\Val$. 
Then 
\begin{multline*}
n\mentval\modT\someoneAll\tf{write}(v) 
\liff
n\mentval\everyoneAll\tf{accept}(v) 
\liff
\\
n\mentval\modT\QuorumBox\tf{accept}(v) 
\liff
n\mentval\modT\someoneAll\tf{accept}(v)  .
\end{multline*}
\end{lemm}
\begin{proof}
We prove a cycle of implications, freely using weak and strong modus ponens (Proposition~\ref{prop.mp.for.tnotor}(\ref{item.mp.for.tnotor}\&\ref{item.mp.for.timpc})).

Suppose $n\mentval\modT\someoneAll\tf{write}(v)$. 
By \rulefont{PaxAccept!} (since $n\geq\gslt(\avaluation)$), $\Forall{p\in\Pnt}n\mentval\tf{accept}(v)$.
By Lemma~\ref{lemm.atomic.O}(\ref{item.atomic.pointwise}) $\tf{accept}(v)$ is pointwise, so by Proposition~\ref{prop.pointwise.dense.char}(\ref{item.dense.char.pointwise.everyoneAll}) $n\mentval\everyoneAll\texi\tf{accept}(v)$.


By Proposition~\ref{prop.correct.cobox.boxT}(\ref{item.correct.cobox.boxT.1}) and \rulefont{PaxPCorrect} (for $\tf{accept}$),
$n\mentval\modT\QuorumBox\tf{accept}(v)$, and by Proposition~\ref{prop.someone.implies.someoneAll}(\ref{item.someone.implies.someoneAll.2}) $n\mentval\modT\someoneAll\tf{accept}(v)$.

Now suppose $n\mentval\modT\someoneAll\tf{accept}(v)$.
By \rulefont{PaxAccept?} %
$n\mentval\modT\someoneAll\tf{write}(v)$.  
\end{proof}


\begin{lemm}
\label{lemm.leaderlive.quorumsend}
Suppose $\avaluation$ is a valuation such that $\mentval\ThyPax$. 
Suppose $n\in\Time$ and $n\geq\gslt(\avaluation)$. 
Then:
\begin{enumerate*}
\item\label{item.leaderlive.quorumsend.1}
$n\mentval\someoneAll\texi\tf{propose}\tnotor\everyoneAll\texi\tf{send}$.
\item\label{item.leaderlive.quorumsend.2}
$n\mentval\someoneAll\texi\tf{propose}\timpc\QuorumBox\texi\tf{send}$.
\end{enumerate*}
\end{lemm}
\begin{proof}
We consider each part in turn:
\begin{enumerate*}
\item
We use \weakmodusponens.
Suppose $n\mentval\modT\someoneAll\texi\tf{propose}$.
By \rulefont{PaxSend!} (since $n\geq\gslt(\avaluation)$) 
$$
\Forall{p\in\Pnt}n,p\mentval\texi\tf{send}.
$$
By Lemma~\ref{lemm.atomic.O}(\ref{item.atomic.pointwise.2}) $\texi\tf{send}$ is pointwise, so by Proposition~\ref{prop.pointwise.dense.char}(\ref{item.dense.char.pointwise.everyoneAll}) $n\mentval\everyoneAll\texi\tf{send}$ as required.
\item
We use \strongmodusponens.
Suppose $n\mentval\modT\someoneAll\texi\tf{propose}$.
By part~\ref{item.leaderlive.quorumsend.1} of this result and \weakmodusponens, $n\mentval\everyoneAll\texi\tf{send}$.
By \rulefont{PaxPCorrect} (for $\tf{send}$) $\acontext\ment\QuorumBox\correct{\tf{send}}$, so by 
Proposition~\ref{prop.correct.cobox.boxT}(\ref{item.correct.cobox.boxT.2})
$\acontext\ment\modT\QuorumBox\texi\tf{send}$ as required. 
\qedhere\end{enumerate*}
\end{proof}




\subsubsection{GSLT implies termination}


Recall from Definition~\ref{defn.correct}(\ref{item.correct}) that $\correct{\tf{decide}}=\tall v.\modTF\tf{decide}(v)$. 
\begin{lemm}
\label{lemm.someone.propose.someone.write}
Suppose $\avaluation$ is a valuation such that $\mentval\ThyPax$. 
Suppose $n\in\Time$ and $n\geq\gslt(\avaluation)$, and $p\in\Pnt$.
Then:
\begin{enumerate*}
\item\label{item.someone.propose.someone.write}
If $n,p\mentval\modT\texi\tf{propose}$ and $n,p\mentval\correct{\tf{write}}$ then $n,p\mentval\modT(\tf{leader}\tand\texi\tf{write})$.
\item\label{item.someone.propose.someone.write.live}
$\mentval\infinitely(\correct{\tf{decide}}\tand\modT\someoneAll\texi\tf{write})$.  
\item\label{item.someone.propose.someone.accept}
$\mentval\infinitely(\someoneAll(\tf{leader}\tand\correct{\tf{decide}})\tand\modT\texi\QuorumBox\tf{accept})$.
\end{enumerate*}
\end{lemm}
\begin{proof}
We consider each part in turn:
\begin{enumerate}
\item
Suppose $n,p\mentval\modT\texi\tf{propose}$.
By \rulefont{PaxPropose?} and \strongmodusponens $n,p\mentval\modT\tf{leader}$.
By Lemma~\ref{lemm.leaderlive.quorumsend}(\ref{item.leaderlive.quorumsend.2}) and strong modus ponens $n\mentval\modT\QuorumBox\texi\tf{send}$.
Using \weakmodusponens and \rulefont{PaxWrite!} (since $n\geq\gslt(\avaluation)$) $n,p\mentval\texi\tf{write}$.
By Proposition~\ref{prop.tv.ment.TF.model}(\ref{item.modTF.texi.P.valid}) (since $n,p\mentval\correct{\tf{write}}$) $n,p\mentval\modT\texi\tf{write}$. 
\item
By \rulefont{LdrExist}, and \rulefont{LdrCorrect} with Lemma~\ref{lemm.persist.char} (for $\f{tvs}=\{\tvT\}$ in that Lemma), and \strongmodusponens, 
for infinitely many $n\in\Time$ there exists $l_n\in\Pnt$ such that 
$$
n,l_n\ment\modT\tf{leader}\tand\correct{\tf{propose},\tf{write},\tf{decide}}. 
$$
So for each $n$ and $l_n$ we have $n,l_n\mentval\correct{\tf{decide}}$.
We will now check that $n,l_n\mentval\modT\texi\tf{write}$.
We reason as follows:
\begin{itemize*}
\item
By \rulefont{PaxPropose!} and \weakmodusponens we have that $n,l_n\ment\texi\tf{propose}$.
\item
By Proposition~\ref{prop.tv.ment.TF.model}(\ref{item.modTF.texi.P.valid}) (since $n,l_n\mentval\correct{\tf{propose}}$ and $n,l_n\mentval\texi\tf{propose}$) $n,l_n\mentval\modT\texi\tf{propose}$. 
\item
By part~\ref{item.someone.propose.someone.write} of this result (since $n,l_n\mentval\correct{\tf{write}}$) $n,l_n\mentval\modT\texi\tf{write}$.
\end{itemize*}
\item
We just combine part~\ref{item.someone.propose.someone.write.live} of this result with Lemma~\ref{lemm.someone.writes.after.gslt.accept}.
\qedhere\end{enumerate}
\end{proof}


Recall that Theorem~\ref{thrm.termination} expresses termination (as the name suggests) but also an abstracted form of the Integrity correctness property by virtue of the use of the $\infinitely$ modality, as per Remark~\ref{rmrk.paxos.integrity}:
\begin{thrm}[Termination]
\label{thrm.termination}
We have:
$$
\ThyPax\ment\infinitely\everyoneAll\texi\tf{decide} .
$$ 
\end{thrm}
\begin{proof}
Suppose $\avaluation$ is a valuation and $\mentval\ThyPax$.
By Lemmas~\ref{lemm.someone.propose.someone.write}(\ref{item.someone.propose.someone.accept}) and~\ref{lemm.persist.char} (for $\f{tvs}=\{\tvT\}$ in that Lemma), for infinitely many $n\in\Time$ there exists an $l_n\in\Pnt$ such that 
\begin{enumerate*}
\item
$n,l_n\mentval\tf{leader}$,
\item
$n,l_n\mentval\correct{\tf{decide}}$ (Definition~\ref{defn.correct}(\ref{item.correct})), and
\item
$n\mentval\modT\texi\QuorumBox\tf{accept}$.
\end{enumerate*}
By \rulefont{PaxDecideL!} (since $n\geq\gslt(\avaluation)$) 
and \weakmodusponens $n,l_n\mentval\texi\tf{decide}$.
By Proposition~\ref{prop.tv.ment.TF.model}(\ref{item.modTF.texi.P.valid}) (since $n,l_n\mentval\correct{\tf{decide}}$ and $n,l_n\mentval\texi\tf{decide}$) $n,l_n\mentval\modT\texi\tf{decide}$. 
Then by \rulefont{PaxDecide\tneg L!} and \weakmodusponens also $n,p\mentval\texi\tf{decide}$ for every $p\in\Pnt$, and so by (Lemma~\ref{lemm.atomic.O}(\ref{item.atomic.pointwise.2}) and) Proposition~\ref{prop.pointwise.dense.char}(\ref{item.dense.char.pointwise.everyoneAll}) $n\mentval\everyoneAll\texi\tf{decide}$. 

The above holds for infinitely many $n\in\Time$, so by Lemma~\ref{lemm.persist.char} (for $\f{tvs}=\{\tvT,\tvB\}$ in that Lemma) we have that $\mentval\infinitely\everyoneAll\texi\tf{decide}$ as required. 
\end{proof}


\begin{rmrk}
Theorem~\ref{thrm.termination} proves that $\ThyPax\ment\infinitely\everyoneAll\texi\tf{decide}$ holds.
Note that $\tvB$ (the truth-value associated to a crashed process) is a valid truth-value, so $n\ment\everyoneAll\texi\tf{decide}$ can hold in the trivial case that everyone is crashed.
Thus, what Theorem~\ref{thrm.termination} asserts in English is that infinitely often, everyone who is not crashed, decides (but it might be that everyone is crashed) --- from which it obviously follows that if a participant never crashes then eventually it decides a value, as per Remark~\ref{rmrk.standard.paxos.correctness}.

But in fact, we can prove a bit more:
$$
\ThyPax\ment\infinitely\QuorumBox\modT\texi\tf{decide}
$$
In words: infinitely often, there exists a \emph{quorum of uncrashed participants}, who decide.

This follows from Theorem~\ref{thrm.termination} by easy reasoning using \rulefont{PaxPCorrect}, \rulefont{PaxDecide\tneg L!}, and Proposition~\ref{prop.decide.n.agree}; we leave this to the reader to check. 
\end{rmrk}


\section{Simpler Declarative Paxos}
\label{sect.simpler.paxos}

Logic is useful for simplification and abstraction in many ways.
In particular, we can axiomatise something, prove some properties --- and then look at what we actually needed from the axioms to make the proofs work.
This can feed back to suggest refinements and generalisations of the original axioms.
(For example: many properties of natural numbers are actually properties of fields, or rings, or integral domains, and so on.) 

We can do something similar here, and look for the most elementary set of axioms that are actually required for our correctness proofs as written: 

\subsection{What we used in the correctness proofs for \ThyPax}
\label{subsect.reflection.on.correctness}
 
\begin{rmrk}
\label{rmrk.spax.explanation}
We note a curious thing about our correctness proofs in Section~\ref{sect.paxos.correctness.properties}, that we never used axioms \rulefont{PaxSend?} or \rulefont{PaxPropose01}.

How can this be?
Is this an error?

No.
This is an opportunity to further simplify the Declarative Paxos from Figure~\ref{fig.ThyPaxOne} to the \emph{Simpler Declarative Paxos} in Figure~\ref{fig.ThySPax}.
To understand how and why this works, we need to highlight some details about how the Paxos algorithm works.
\end{rmrk}

\begin{rmrk}[Why we did not use \rulefont{PaxSend?}]
Recall from Remark~\ref{rmrk.high-level.paxos} our high-level view of the Paxos algorithm, and recall that \ThyPax in Figure~\ref{fig.ThyPaxOne} is a declarative version of that algorithm (obtained by abstracting away explicit message-passing and algorithmic time, as per Remark~\ref{rmrk.elisions}).

Note of the algorithm that 
\begin{itemize*}
\item
in step~\ref{item.l.proposes} of Remark~\ref{rmrk.high-level.paxos}, the leader broadcasts a `propose' message; and 
\item
in step~\ref{item.p.sends}, a participant $p$ that receives that `propose' message, replies with a `send' message carrying as its payload a value $v_p$, which is $p$'s most recently accepted value (if any); and
\item
in step~\ref{item.l.writes}, the leader gathers the $v_p$ values from these responses and, if it has a quorum (in topological language: an open set) of responses, it selects a most recent value to use in its write step.
\end{itemize*}
In the algorithm, the only way that the leader $l$ can know of a participant $p$'s value of $v_p$, is if $l$ receives a send message from $p$ with payload $v_p$.
If $l$ receives no such message, then it does not know the value of $v_p$.

In contrast, predicates of \QLogic are not assigned truth-values by participants in an algorithm; they are assigned truth-values as if by an omniscient observer who is standing outside of the model and who has full (omniscient) access to it.
This is reflected in \rulefont{PaxWrite?}, where we write $\mru{(\someone\tf{accept})}{v}$.
This directly examines the model to find the most recently accepted value.
No reference here is required to the payload of the quorum of $\tf{send}$ messages.
\rulefont{PaxWrite?} still insists that such a quorum exists, as per the $\everyone\texi\tf{send}$ that appears in \rulefont{PaxWrite?}, so in this case the value $\mru{(\someone\tf{accept})}{v}$ is computable by the leader, and we could specify a more complex predicate that checks the set of $v_p$ payloads in the send messages of that quorum to select a most recently-accepted value.
In the implementation, this is necessary because the leader $l$ is not omniscient.
But in the logic we have no reason to go to this trouble, because we can just read the information directly from the model using $\mru{(\someone\tf{accept})}{v}$.\footnote{For example: if we want to compute $100!$ in the real world, we have to build a computer to do this for us, which involves many practical details including the construction of a semiconductor industry, compiler design and so on.  In logic, however, we can just say that it is $\prod_{i=1}^{100}i$, because that is what $100!$ \emph{is}.}

So our axioms are just doing their job, which is to be a declarative abstraction of the implementation. 
But since we do not use \rulefont{PaxSend?}, this is evidence that we may treat the $\tf{send}$ step as an implementational detail having to do with how information moves around the network, and elide it.
\end{rmrk}

\begin{rmrk}[Why we did not use \rulefont{PaxPropose01}]
This is due to another particularity of how the Paxos algorithm works.
The leader chooses one value in the propose step (or no value, if the leader has crashed), but it only uses that value in the case that it finds a quorum of participants who have never previously accepted any value (this is the right-hand disjunct in \rulefont{PaxWrite?}).
Just in this case, it writes the value that it chose in the propose step.

So what really matters is that the leader \emph{only writes one value}.
If the reader had proposed two values, this would be fine --- provided that it only writes one of them.

Leaders in the Paxos algorithm do not bother with multiple values in the propose step, because this would be a waste of time; they know they will only ever need at most one.
Thus, \rulefont{PaxPropose01} is right to reflect a true fact of the algorithm that it does not propose any more than the value than it will need, but this an algorithmic efficiency not a logical necessity. 
Likewise, the logic is right to point this out to us by \emph{not} requiring \rulefont{PaxPropose01} in the proofs: so long as only at most one value gets \emph{written} --- which is handled by \rulefont{PaxWrite01} --- none of our correctness properties are affected.
\end{rmrk}

\begin{rmrk}
By examining the proofs further, some further possible simplifications are revealed:
\begin{enumerate*}
\item
We can combine rules \rulefont{SPaxWrite01} and \rulefont{PaxWrite01} and \rulefont{LdrExt}, as per the discussion in Subsection~\ref{subsect.ldrexistuniq'}, to obtain an elegantly concise rule \rulefont{SPaxWrite01}.
\item
We used \rulefont{PaxPCorrect} in our proofs to obtain $\QuorumBox\correct{\tf P}$ only twice: 
\begin{itemize*}
\item
in Lemma~\ref{lemm.someone.writes.after.gslt.accept} for $\tf P=\tf{accept}$, and 
\item
in Lemma~\ref{lemm.leaderlive.quorumsend} for $\tf P=\tf{send}$.
\end{itemize*}
We have elided the \emph{send} step in $\Theta_\ThySPax$, so we only need to consider $\tf{accept}$.
Accordingly we take take \rulefont{SPaxPCorrect} to be just $\QuorumBox\correct{\tf{accept}}$.
\end{enumerate*}
\end{rmrk}

\begin{rmrk}
In conclusion, we see that the proofs in Section~\ref{sect.paxos.correctness.properties} are not only correctness proofs, but can also be read as an active commentary on the axioms.
Our logic has done what logic does, and has separated what is logically necessary from what is implementationally required or implementationally convenient.
This shows us how to simplify the axiomatisation (if we we wish) further away from the implementation, and closer towards what is logically essential: 
\end{rmrk}


\subsection{The definition}

\begin{figure}
$$
\begin{array}{l@{\quad}r@{\ }l}
\figunderline{Backward rules, for backward inference} 
\figskip
\rulefont{SPaxPropose?}&
& \tall v.\ \tf{propose}(v)\timpc\tf{leader}
\figskip
\rulefont{SPaxWrite?}&
& \tall v.\ \tf{write}(v) \timpc 
\begin{array}[t]{l}
\tf{leader}\tand\Quorum( 
\\
\quad\bigl((\mru{(\someone\tf{accept})}{v}) \tor
(\tf{propose}(v)\tand\tneg\modT\recent\texi\someone\tf{accept})\bigr)) 
\end{array}
\figskip
\rulefont{SPaxAccept?}&
& \tall v.\ \tf{accept}(v) \timpc \someoneAll\tf{write}(v)
\figskip
\rulefont{SPaxDecide?}&
& \tall v.\ \tf{decide}(v) \timpc \QuorumBox\tf{accept}(v)
\figskip\figskip
\figunderline{Forward rules, for making progress} 
\figskip
\rulefont{SPaxWrite!}&
& \final(\tf{leader}\tnotor\texi\tf{write})
\figskip
\rulefont{SPaxAccept!}&
& \final(\texi\someoneAll\tf{write}\tnotor \texi\tf{accept})
\figskip
\rulefont{SPaxDecide!}&
& \final(\texi\QuorumBox\tf{accept} \tnotor \texi\tf{decide})
\figskip
\figskip
\figunderline{Other rules, for correctness \& liveness}
\figskip
\rulefont{SLdrExist}&
&\modTF\tf{leader}\tand \modT\someoneAll\tf{leader}
\figskip
\rulefont{SLdrCorrect}&
& \infinitely\someoneAll(\tf{leader}\tand \correct{\tf{propose},\tf{write},\tf{decide}}) 
\figskip
\rulefont{SPaxPCorrect}&
& \QuorumBox\correct{\tf{accept}}
\figskip
\rulefont{SPaxWrite01}&
& \texiaffine\someoneAll\tf{write}
\end{array}
$$
\caption{Simpler Declarative Paxos (Definition~\ref{defn.thyspax})}
\label{fig.logical.paxos'}
\label{fig.ThySPax}
\end{figure}


\begin{defn}
\label{defn.thyspax}
We define \deffont{Simpler Declarative Paxos} to be the logical theory 
$$
\ThySPax=(\Sigma_{\ThySPax},\Theta_{\ThySPax}),
$$ 
where:
\begin{enumerate*}
\item
the signature $\Sigma_{\ThySPax}$ is 
$$
\Sigma_{\ThySPax}=[\tf{leader}:0;\ \tf{propose},\tf{write},\tf{accept},\tf{decide}:1] ,
$$ 
as per the notation in Definition~\ref{defn.signature}(\ref{item.predicate.syntax.signature}), and
\item
the axioms $\Theta_{\ThySPax}$ are as written in Figure~\ref{fig.ThySPax}.
\end{enumerate*}
\end{defn}

The same validity properties hold of \ThySPax from Figure~\ref{fig.ThySPax} as hold for \ThyPax from Figure~\ref{fig.ThyPaxOne}, namely: Validity (Theorem~\ref{thrm.validity}), Agreement (Theorem~\ref{thrm.full.agreement}), and Termination and Integrity (Theorem~\ref{thrm.termination}):
\begin{thrm}
\label{thrm.simpler.dp}
Suppose $\Sigma$ is a signature that includes $\Sigma_\ThySPax$, and suppose $\mathcal M$ is a model.
Then we have:
\begin{enumerate*}
\item
$\ThyPax\ment\tall v.(\tf{decide}(v) \timpc \urecent\someoneAll(\tf{leader}\tand\tf{propose}(v)))$.
\item
$\ThyPax\ment\tall v,v'.(\urecent\someoneAll\tf{decide}(v) \timpc \someoneAll\tf{decide}(v') \timpc v'\teq v)$.
\item
$\ThyPax\ment\infinitely\everyoneAll\texi\tf{decide}$.
\end{enumerate*}
\end{thrm}
\begin{proof}
This is by construction: from the discussion in Subsection~\ref{subsect.reflection.on.correctness}, the theory \rulefont{ThySPax} consists of just those parts of \rulefont{ThyPax} that are actually required in the correctness proofs.
\end{proof}

Thus, Theorem~\ref{thrm.simpler.dp} just notes that Simpler Declarative Paxos extracts from the proofs in Section~\ref{sect.paxos.correctness.properties} the logical essence of Figure~\ref{fig.ThyPaxOne}, as applied to the standard correctness properties of Paxos.
There is no right or wrong about whether we prefer $\ThyPax$ or $\ThySPax$; it depends on what level of abstraction we want to work with at any particular moment.
But, implicit in Theorem~\ref{thrm.simpler.dp} is an observation that \ThySPax is \emph{minimal}, because we put into it only that which we know we need to make some specific correctness proofs work.


\section{Conclusions and future work}
\label{sect.conclusions}

\subsection{Conclusions}







We have presented \QLogic, a three-valued modal fixedpoint logic for declarative specification of consensus algorithms, and we have shown how to use this to axiomatise Paxos, a canonical and nontrivial (though still relatively straightforward) consensus algorithm that can handle benign failures of participants.

We do not presented more complex consensus algorithms in this paper but we have looked into these in preliminary investigations, and they do seem to work. 

This paper is aimed at 
\begin{itemize*}
\item
logicians who might be interested in exploring a logic with a new and unusual combination of features inspired by strong practical motivations, and 
\item
researchers in distributed systems, who might consider (and benefit from) adopting such a framework.
\end{itemize*}
\QLogic offers benefits that could make it a useful addition to the toolkit of anyone working in distributed systems:
\begin{enumerate}
\item 
Much of what is shown in this paper just reflects existing practices, formalised into a new logic.
Therefore, readers may already be reasoning in ways that align with our logic, without explicitly realising it. 
For example, the third truth-value ($\tvB$) in \QLogic corresponds to typical side conditions like ``for every correct participant,'' and the backward axioms mirror the reverse inferences that practitioners often use. 

By formalising these familiar patterns, \QLogic provides a structured and precise way to represent and reason about concepts that are already informally expressed in the literature.
\item 
Logic is a concise, powerful, and unambiguous language. 
It distills paragraphs of English prose into a single line of precise symbols, enabling clear and rigorous reasoning. 

While explanatory English is helpful for accessibility and context, formal logic is the superior tool for working through the details of complex algorithms. 
It ensures precision and eliminates the ambiguities that often arise in informal descriptions.
\item 
\QLogic operates at a new (and higher) level of abstraction, which complements existing logic-based techniques. 
TLA+~\cite{DBLP:journals/toplas/Lamport94} provides a state-based view of distributed algorithms, whereas \QLogic abstracts lower-level implementation details while retaining the essential behaviours and properties of the algorithms.
This abstraction does not obscure the logic of the system; on the contrary we have found that it clarifies it, making the underlying mechanisms and invariants more apparent. 
\item 
One strength of logic is its ability to abstract away what is irrelevant while maintaining precision about what is important. 
In \QLogic, constructs like strong and weak implications ($\timpc$ and $\tnotor$) exemplify this principle, capturing subtle distinctions in meaning with precise formal symbols. 
These constructs distill complex ideas into simple yet powerful tools that can be manipulated symbolically and reasoned about systematically. 

While mastering these constructs requires some initial effort, the benefits of precision and clarity justify the investment.
\item 
Learning \QLogic requires some effort, but for researchers already familiar with (for example) TLA+, the transition to \QLogic should be eminently manageable. 

While axiomatisations in \QLogic are certainly subtle, this is not gratuitous: they reflect the inherent nature of the distributed systems they are designed to model. 
By formalising these complexities, \QLogic captures their structure and reveals and addresses the challenges that are \emph{already present} in the underlying algorithms. 

This transparency makes it a powerful tool for understanding and reasoning about distributed protocols.\footnote{Certainly, the first author (who has a background in logic) has found this study invaluable for gaining a better understanding of consensus algorithms.  It seems likely that other readers might benefit similarly from using \QLogic.}
\item 
Distributed algorithms, and consensus protocols in particular, are foundational to the functionality of blockchain systems and other critical infrastructures. 

These algorithms are recognised to be notoriously challenging to design and implement correctly. 
Ambiguities and informal specifications can lead to vulnerabilities and failures. 

\QLogic bridges the gap between declarative specifications of invariants, and practical verification; enabling researchers to clarify, reason about, and verify the properties of distributed protocols rigorously. 
This is especially valuable in blockchain systems, where clear and robust specifications are essential for trust and reliability.
\end{enumerate}

In conclusion: by providing a formal framework that combines the abstraction of declarative reasoning with the rigour of logical precision, \QLogic offers researchers in distributed systems a powerful tool for modeling, understanding, and verifying complex algorithms. 

\subsection{Related and future work}

In this paper, our focus has been to lay the foundational theory and framework of \QLogic and to demonstrate its application to one of the most well-known consensus protocols, Paxos~\cite{DBLP:journals/tocs/Lamport98}.

Paxos is a sensible starting point, being a cornerstone of distributed systems, but (certainly by modern standards!) it is relatively straightforward.
For example, Paxos only addresses crash failures, limiting its applicability to systems that assume benign (rather than hostile) faults.

A natural direction for future work involves extending our framework to handle more complex and realistic fault models, such as Byzantine faults, where participants can deviate arbitrarily from the algorithm.
Consensus protocols that address Byzantine faults, such as Practical Byzantine Fault Tolerance (PBFT)~\cite{DBLP:conf/osdi/CastroL99} and HotStuff~\cite{DBLP:conf/podc/YinMRGA19}, provide robust mechanisms for fault tolerance and serve as the foundation for many modern blockchain systems, including Tendermint (\url{https://tendermint.com})~\cite{DBLP:journals/corr/abs-1807-04938}. 

Applying \QLogic to these protocols would see how it handles the complexity introduced by arbitrary faults, which is necessary to model and verify more sophisticated algorithms.
Preliminary exploration by the authors in this direction indicates that \QLogic works very well in these more complex situations --- it can indeed axiomatise algorithms that take account of arbitrary, including hostile, faults.
Unpacking this in full is for future papers.

Recent research has explored DAG-based consensus protocols, which deviate from the traditional linear blockchain structure. 
In these protocols, transactions are organized into a directed acyclic graph (DAG) rather than a strict chain, enabling greater parallelism and scalability. 
Protocols such as DAG-rider~\cite{DBLP:conf/podc/KeidarKNS21}, Narwhal and Tusk~\cite{DBLP:conf/eurosys/DanezisKSS22}, and Bullshark~\cite{DBLP:conf/ccs/SpiegelmanGSK22} represent cutting-edge advancements in this domain. 
Extending \QLogic to analyse DAG-based protocols would be another advance, as these systems often involve more intricate interactions between components and require reasoning about concurrent and asynchronous behaviour at a much higher level of complexity.

Another promising avenue for future research involves leveraging the modularity and expressiveness of \QLogic to study hybrid protocols~\cite{DBLP:conf/sp/NeuTT21} that combine elements of classical consensus with newer techniques, such as Gasper~\cite{DBLP:journals/corr/abs-2003-03052} and others~\cite{DBLP:conf/esorics/DAmatoZ23,DBLP:journals/corr/abs-2411-00558}

It would be helpful in future work to develop automated tools for \QLogic: the tools developed for TLA+~\cite{DBLP:books/aw/Lamport2002} could show the way.
TLA+ is a formal specification language and framework designed to model, specify and verify properties of complex systems; particularly distributed algorithms. 
Its TLC model checker~\cite{DBLP:conf/charme/YuML99} provides exhaustive state exploration to verify properties such as safety and liveness, and the PlusCal translator simplifies specification by offering a high-level pseudocode-like interface.
Additionally, the TLA+ Toolbox~\cite{DBLP:journals/corr/abs-1912-10633} integrates these capabilities into a development environment, supporting workflows that are valuable in both research and industrial applications.

Adapting similar tools to \QLogic could enable automated verification of correctness properties for consensus protocols, including invariants and liveness guarantees. 
For example: a higher-level language akin to PlusCal could simplify the specification of protocols using the axiomatic foundations of \QLogic.
Furthermore, leveraging techniques such as the symbolic model checking enhancements presented by Otoni et al.~\cite{DBLP:conf/tacas/OtoniKKES23} or integrating modular proof systems like TLAPS~\cite{DBLP:conf/fm/CousineauDLMRV12} could further enhance the usability and robustness of the framework.

Such tools would streamline the application of \QLogic to real-world systems, and also make it accessible to a broader audience of researchers and practitioners. 
This direction represents an exciting convergence of theoretical advancements and practical tool development, providing a robust framework for analysing both classical and emerging consensus protocols.

Adopting \QLogic requires some investment in time and effort, but the benefits in clarity, correctness, and reliability may be worthwhile.
Consensus algorithms are known to be particularly challenging, so our relatively high-level declarative approach could help to manage (and debug) the complexity of these algorithms.

We hope this work will inspire readers %
to explore and adopt this new perspective, contributing to more robust and trustworthy distributed protocols in the future.














\documentclass[acmjacm]{acmsmall}

\usepackage[shortlabels]{enumitem}
 
\usepackage{latexsym,xspace,calc,amsthm}
\usepackage{amsmath,amssymb} %
\usepackage{centernot}

\setcounter{tocdepth}{4}
\setcounter{secnumdepth}{3}


\usepackage{lmodern}
\usepackage{tgtermes}
\usepackage{fix-cm}


\usepackage{cancel}

\usepackage{savesym}
\savesymbol{subcaption}
\usepackage{subcaption}

\usepackage{graphbox} %

\makeatletter
\def\@biblabel#1{[#1]} %
\def\thebibliography#1{%
    \footnotesize
    \refsection*{{\refname}
        \@mkboth{\uppercase{\refname}}{\uppercase{\refname}}%
    }
    \list{\@biblabel{\@arabic\c@enumiv}}%
       {\settowidth\labelwidth{\@biblabel{#1}}%
        \leftmargin\labelwidth
        \advance\leftmargin\bibindent
        \itemindent-\bibindent
        \itemsep2pt
        \parsep \z@
        \usecounter{enumiv}%
        \let\p@enumiv\@empty
        \renewcommand\theenumiv{\@arabic\c@enumiv}%
    }%
    \let\newblock\@empty
    \sloppy
    \sfcode`\.=1000\relax
}
\makeatother

\let\odlthebibliography=\thebibliography
 \let\endodlthebibliography=\endthebibliography
 \renewenvironment{thebibliography}[1]{%
   \begin{odlthebibliography}{#1}%
     \setlength{\parskip}{0ex}%
     \setlength{\itemsep}{3pt}%
     \fontsize{9.5}{9.5} %
     \selectfont
}%
 {%
   \end{odlthebibliography}%
 }


\usepackage{tikz}
\usepackage{tikz-cd}
\usetikzlibrary{cd,decorations.markings}
\tikzset{
    dharrow/.style={
        <->,
        postaction={decorate,-},
        }
}
\tikzset{
    dhdashedarrow/.style={
        <->,
        dashed,
        postaction={decorate,-},
        }
    }
\tikzset{
    lrharpoonarrow/.style={
        <[harpoon]->[harpoon],
        postaction={decorate,-},
        }
}
\tikzset{
    lrharpoondashedarrow/.style={
        <[harpoon]->[harpoon],
        dashed, %
        postaction={decorate,-},
        }
}
\usetikzlibrary{arrows}
\usetikzlibrary {arrows.meta} 
\usetikzlibrary{calc}
\usetikzlibrary{positioning}
\usetikzlibrary{snakes,automata,chains}
\usetikzlibrary{graphs}

\usepackage{binarytree}

\usepackage{amssymb,stmaryrd,amsmath}
\usepackage{mdwlist} %
\usepackage{float}   %
\usepackage{centernot} %


\usepackage[colorlinks]{hyperref}
\usepackage{breakurl}  %


\usepackage{bm} %





\newtheoremstyle{jamiestyle}%
  {4pt}%
  {0pt}%
  {\it}%
  {0pt}%
  {\sc}%
  {.}%
  { }%
  {}%
\theoremstyle{jamiestyle}
\newtheorem{thrm}{Theorem}[subsection]
\newtheorem{prop}[thrm]{Proposition}
\newtheorem{lemm}[thrm]{Lemma}
\newtheorem{corr}[thrm]{Corollary}
\newtheorem{conj}[thrm]{Conjecture}
\newtheoremstyle{jamienfstyle}%
  {4pt}%
  {0pt}%
  {\normalfont}%
  {0pt}%
  {\sc}%
  {.}%
  { }%
  {}%
\theoremstyle{jamienfstyle}
\newtheorem{nttn}[thrm]{Notation}
\newtheorem{defn}[thrm]{Definition}
\newtheorem{xmpl}[thrm]{Example}
\newtheorem{rmrk}[thrm]{Remark}


\usepackage{color}
\definecolor{mygreen}{rgb}{0,0.6,0}
\definecolor{mygray}{rgb}{0.5,0.5,0.5}
\definecolor{mymauve}{rgb}{0.58,0,0.82}

\usepackage{listings}
 
\definecolor{gray}{RGB}{128, 128, 128}
\definecolor{lightgray}{RGB}{200, 200, 200}
\definecolor{cyan}{RGB}{0, 255, 255}
\definecolor{blue}{RGB}{0, 0, 255}
\definecolor{red}{RGB}{255, 0, 0}
\definecolor{pink}{RGB}{255, 128, 128}
\definecolor{green}{RGB}{0, 128, 0}
\definecolor{lightyellow}{RGB}{255, 255, 200}
\definecolor{purple}{RGB}{128, 0, 128}

\lstdefinestyle{all}
    {basicstyle=\ttfamily\scriptsize,
     keywordstyle=\color{blue}\ttfamily\scriptsize,
     commentstyle=\color{green}\ttfamily\scriptsize,
     stringstyle=\color{red}\ttfamily\scriptsize}

\lstdefinelanguage{hask}{%
    frame=none,
    xleftmargin=2pt,
    belowcaptionskip=\bigskipamount,
    captionpos=b,
    tabsize=2,
    numbers=left,
    numberstyle=\tiny\color{gray},
    emphstyle={\bf},
	morecomment=[s][\color{green}]{\{-}{-\}},
    stringstyle=\mdseries\rmfamily,
    commentstyle=\color{green},
    keywords={},
    keywords=[1]{case, of, data, if, then, else, where, let, in, do},
    keywords=[2]{Chip, Config, CurrencySymbol, TokenName, PubKeyHash, Integer, Value, State, Action, Text, Maybe, Void, TxConstraints,  Contract},
    keywords=[3]{HasNative},
    keywords=[4]{=>},
    keywords=[5]{Just, Nothing, MkChip, MkConfig, SetPrice, Buy},
    keywordstyle=[1]\mdseries\sffamily\color{red},
    keywordstyle=[2]\mdseries\sffamily\color{blue},
    keywordstyle=[3]\mdseries\sffamily\color{green},
    keywordstyle=[4]\mdseries\sffamily,
    keywordstyle=[5]\mdseries\sffamily\color{purple},
    columns=flexible,
    basicstyle=\small\sffamily,
    showstringspaces=false,
    breaklines=false,
    showspaces=false,
    escapeinside={--}{\^^M},escapebegin={\color{green}--},escapeend={},
    literate= {+}{{$+$}}1 {/}{{$/$}}1 {*}{{$*$}}1 {=}{{$=$}}1
              {>}{{$>$}}1 {<}{{$<$}}1 {\\}{{$\lambda$}}1
              {\\\\}{{\char`\\\char`\\}}1
              {->}{{$\rightarrow$}}2 {>=}{{$\geq$}}2 {<-}{{$\leftarrow$}}2
              {<=}{{$\leq$}}2 {=>}{{$\Rightarrow$}}2
              {\ .}{{$\circ$}}2 {\ .\ }{{$\circ$}}2
              {>>}{{>>}}2 {>>=}{{>>=}}2
              {|}{{$\mid$}}1
              {\_}{{\underline{\hspace{2mm}}}}2
}

\lstdefinelanguage{solidity}{%
    frame=none,
    xleftmargin=2pt,
    belowcaptionskip=\bigskipamount,
    captionpos=b,
    tabsize=2,
    numbers=left,
    numberstyle=\tiny\color{gray},
    emphstyle={\bf},
	morecomment=[s][\color{green}]{\{-}{-\}},
    stringstyle=\mdseries\rmfamily,
    commentstyle=\color{green},
    keywords={},
    keywords=[1]{pragma, solidity, contract, event, constructor, require, function, return, emit},
    keywords=[2]{address, uint, mapping},
    keywords=[3]{public, payable, external, view, returns},
    keywords=[4]{=>, +=, -=, =, <=, ==},
    keywords=[5]{msg, sender, transfer, value},
    keywordstyle=[1]\mdseries\sffamily\color{red},
    keywordstyle=[2]\mdseries\sffamily\color{blue},
    keywordstyle=[3]\mdseries\sffamily\color{green},
    keywordstyle=[4]\mdseries\sffamily,
    keywordstyle=[5]\mdseries\sffamily\color{purple},
    columns=flexible,
    basicstyle=\small\sffamily,
    showstringspaces=false,
    breaklines=false,
    showspaces=false,
    escapeinside={--}{\^^M},escapebegin={\color{green}--},escapeend={},
    literate= {+}{{$+$}}1 {/}{{$/$}}1 {*}{{$*$}}1 {=}{{$=$}}1
              {>}{{$>$}}1 {<}{{$<$}}1 {\\}{{$\lambda$}}1
              {\\\\}{{\char`\\\char`\\}}1
              {->}{{$\rightarrow$}}2 {>=}{{$\geq$}}2 {<-}{{$\leftarrow$}}2
              {<=}{{$\leq$}}2 {=>}{{$\Rightarrow$}}2
              {\ .}{{$\circ$}}2 {\ .\ }{{$\circ$}}2
              {>>}{{>>}}2 {>>=}{{>>=}}2
              {|}{{$\mid$}}1
              {\_}{{\underline{\hspace{2mm}}}}2
}




\newcommand\flanks{\ltimes}
\newcommand\setminusskip{\hskip 2pt plus 1pt minus 2pt}
\newcommand\mysetminus{\setminusskip{{\setminus}}\setminusskip}

\newcommand\declaresoundness[2]{The case of \rulefont{#1} with $#2$.}
\newcommand\declaresoundnessshort[1]{The case of \rulefont{#1}.}


\makeatletter
\newcommand\hpn[2][]{%
  \ext@arrow 9999{\hpnfill@}{#1}{#2}}
\newcommand\hpnfill@{%
  \arrowfill@\leftharpoonup\relbar\rightharpoondown}
\makeatother

\newcommand\arrowlength{0.5}
\newcommand\arrowr{{\tikz \draw [semithick,->] (0,0) -- (\arrowlength*0.8,0);}}
\newcommand\Arrowr{{\tikz \draw [semithick,->,decorate,decoration={zigzag,segment length=3,amplitude=0.5}] (0,0) -- (\arrowlength*0.8,0);}}
\newcommand\Arrowrr{{\tikz \draw [semithick,->>,decorate,decoration={zigzag,segment length=3,amplitude=0.5}] (0,0) -- (\arrowlength*0.8,0);}}
\newcommand\arrowlr{{\tikz \draw [semithick,<->] (0,0) -- (\arrowlength*0.8,0);}}
\newcommand\arrowllrr{{\tikz \draw [semithick,<<->>] (0,0) -- (\arrowlength,0);}}
\newcommand\arrowrr{{\tikz \draw [semithick,->>] (0,0) -- (\arrowlength*0.8,0);}}
\newcommand\arrowrrr{{\tikz \draw [thick,-Stealth] (0,0) -- (\arrowlength*0.8,0);}}
\newcommand\arrowll{{\tikz \draw [semithick,<<-] (0,0) -- (\arrowlength*0.8,0);}}
\newcommand\arrowerr{{\tikz \draw [semithick,-.>>] (0,0) -- (\arrowlength*0.8,0);}}
\newcommand\arrowlrr{{\tikz \draw [semithick,<->>] (0,0) -- (\arrowlength,0);}}
\newcommand\arrowllr{{\tikz \draw [semithick,<<->] (0,0) -- (\arrowlength,0);}}
\newcommand\arrowelrr{{\tikz \draw [semithick,<.->>] (0,0) -- (\arrowlength,0);}}
\newcommand\arrowelr{{\tikz \draw [semithick,<.->] (0,0) -- (\arrowlength,0);}}
\newcommand\arrower{{\tikz \draw [semithick,-.>] (0,0) -- (\arrowlength*0.8,0);}}


\newcommand\everyone{\modality{\Box}}
\newcommand\someone{\modality{\Diamond}}
\newcommand\someoneAll{{\mathring{\exists}_{\hspace{-0pt}\someone}}}
\newcommand\everyoneAll{{\mathring{\forall}_{\hspace{-1pt}\everyone}}}

\newcommand\Quorum{\modality{\mathring{\exists}}}
\newcommand\QuorumBox{\modality{\mathring{\exists}_{\Box}}} %
\newcommand\Coquorum{\modality{\mathring{\forall}}}
\newcommand\CoquorumDiamond{\modality{\mathring{\forall}_{\hspace{-1pt}\Diamond}}} %


\newcommand{\ite}[2]{[#1,\ \text{\it o\hspace{-1pt}'wise}\ #2]}
\newcommand\mruarrow{\mathbin{\rlap{\scalebox{0.73}{$\recentsimple$}}\hspace{-2pt}\rightarrow}}
\newcommand{\mru}[2]{#1\mruarrow#2}
\newcommand{\mrup}[3]{#2[#1\mruarrow #3]}
\newcommand{\mruX}[2]{{#1\mathbin{\recent}#2}}
\newcommand{\established}[2]{#1 \mathbin{\Downarrow_{\QuorumBox}} #2}
\newcommand{\lrharpoon}{\mathbin{\hpn{}}}
\newcommand{\nlrharpoon}{\mathbin{\not\hpn{}}}
\newcommand\id{\f{id}}
\newcommand\st{{:}}

\usepackage{tikz}
\newcommand*{\myhash}{%
  \begin{tikzpicture}
    \pgfmathsetlengthmacro\myWidth{.8*width("=")}%
    \pgfmathsetlengthmacro\myHeight{height("H")}%
    \pgfmathsetlengthmacro\mySepY{.3333*\myWidth}%
    \pgfmathsetlengthmacro\mySideBearing{.1*\myWidth}%
    \def\myAngle{70}%
    \pgfmathsetlengthmacro\mySepX{\mySepY/sin(\myAngle)}%
    \pgfmathsetlengthmacro\mySlantX{\myHeight/tan(\myAngle)}%
    \draw[line cap=round]
      (0, {(\myHeight - \mySepY)/2}) -- ++(\myWidth, 0)
      (0, {(\myHeight + \mySepY)/2}) -- ++(\myWidth, 0)
      ({(\myWidth - \mySepX - \mySlantX)/2}, 0)
      -- ({(\myWidth - \mySepX + \mySlantX)/2}, \myHeight)
      ({(\myWidth + \mySepX - \mySlantX)/2}, 0)
      -- ({(\myWidth + \mySepX + \mySlantX)/2}, \myHeight)
    ;%
    \useasboundingbox
      (-\mySideBearing, 0)
      (\myWidth + \mySideBearing, \myHeight)
    ;%
  \end{tikzpicture}%
}

\newcommand\THREE{{\mathbf 3}}
\newcommand\threeValid{\THREE_{\f{valid}}}
\newcommand\threeCorrect{\THREE_{\f{correct}}}
\newcommand\threeTrue{\THREE_{\f{true}}}
\newcommand\threeFalse{\THREE_{\f{false}}}
\newcommand\binaryconnectives{\ensuremath{\{\tand,\tor,\tnotor,\timpc\}}} %
\newcommand\unaryconnectives{\ensuremath{\{\tneg,\modT,\modB,\modF,\modTB,\modTF\}}}
\newcommand\modalityconnectives{\ensuremath{\{\someone,\everyone,\Quorum,\Coquorum,\sometime,\forever,\tomorrow\}}} %
\newcommand\doublemodalities{\ensuremath{\{\QuorumBox,\CoquorumDiamond,\someoneAll,\everyoneAll\}}}
\newcommand\singlemodalities{\ensuremath{\{\someone,\everyone,\Quorum,\Coquorum\}}}
\newcommand\Pnt{\ns{Point}}
\newcommand\Val{\ns{Value}}
\newcommand\Valuation{\ns{Valuation}}
\newcommand\PredSymb{\ns{PredSymb}}

\newcommand\FOUR{{\mathbf 4}}
\newcommand\tvT{{\mathbf t}}
\newcommand\tvF{{\mathbf f}}
\newcommand\tvB{{\mathbf b}}
\newcommand\tvN{{\mathbf n}}
\newcommand\tvsTB{{\mathit{tb}}}
\newcommand\tvsFB{{\mathit{fb}}}
\newcommand\tvsTT{{\mathit{tt}}}
\newcommand\tvsFF{{\mathit{ff}}}
\newcommand*\Rcirc{\ensuremath{\smash{\overset{\circ}{\smash{\mathsf{R}}\rule{0pt}{1.2ex}}\vphantom{\mathbb{R}}}}}
\newcommand\recentsimple{{\reflectbox{$\mathsf R$}}} 
\newcommand\recent{\recentsimple} %
\newcommand\urecent{{\ubar\recentsimple}} %
\newcommand\mrecent[2]{#1\recentsimple #2}

\newcommand\figunderline[1]{\text{\rlap{\underline{\emph{#1}}}}}




\usepackage{stackengine}
\newcommand\ubar[1]{\stackunder[0.6pt]{$#1$}{\rule{1ex}{.18ex}\hspace{0.5pt}}}


\newcommand\figskip{\\[1ex]}
\newcommand\from[1]{{\uparrow}_{#1}}
\newcommand\fs[1]{\mathtt{#1}}
\newcommand\ps[1]{\mathtt{#1}}
\newcommand\thea{a}
\newcommand\theb{b}
\newcommand\thec{c}
\newcommand\aval{v}
\newcommand\xor{\mathbin{\mathsf{\small xor}}}
\newcommand\modality[1]{#1\hspace{1.5pt}}
\newcommand\modT{\modality{\tf T}}
\newcommand\modTB{\modality{\tf T\hspace{-2.5pt}\tf B}}
\newcommand\modFB{\modality{\tf F\hspace{-1.5pt}\tf B}}
\newcommand\modB{\modality{\tf B}}
\newcommand\modTF{\modality{\tf T\hspace{-2.5pt}\tf F}}
\newcommand\modF{\modality{\tf F}}
\newcommand\witno{\prec}
\newcommand\atopen{T}
\newcommand\afilter{F}
\newcommand\apoint{P}
\newcommand\acontext{\f{ctx}}
\newcommand\avaluation{\varsigma} %
\newcommand\indicator[1]{\delta_{#1}} 
\newcommand\charop[1]{\f{char}(#1)}
\newcommand\resilient{R}
\newcommand\rulefont[1]{\ensuremath{{\mathsf{(#1)}}}}
\newcommand\leftopeninterval[1]{(#1]}
\newcommand\rightopeninterval[1]{[#1)}
\newcommand\openinterval[1]{(#1)}
\newcommand\closedinterval[1]{[#1]}
\newcommand\opens{\ns{Open}}
\newcommand\opensne{\opens_{\hspace{-1pt}{\neq}\varnothing}}
\newcommand\regularOpens{\tf{Open}_{\f{reg}}}
\newcommand\topens{\tf{Topen}}
\newcommand\closed{\tf{Closed}}
\newcommand\regularClosed{\tf{Closed}_{\f{reg}}}
\newcommand\thenext[0]{\f{ad}}
\newcommand\closure[1]{|#1|}
\newcommand\dotclosure[1]{\closure{#1}^{\bigcdot}}
\newcommand{\dotarrow}{%
   \mathrel{\ooalign{\hss\raise.85ex\hbox{\scalebox{1.25}{\normalfont .}}%
   \kern0.35ex\hss\cr$\rightarrow$}}}
\newcommand{\dotleftarrow}{\mathrel{\reflectbox{$\dotarrow$}}}
\newcommand\mycard{\myhash}
\newcommand\narrowmath[1]{\scalebox{0.8}[1.0]{$\mathsf{#1}$}}
\newcommand\sometime{\modality{\narrowmath{sometime}}}
\newcommand\forever{\modality{\narrowmath{forever}}}
\newcommand\tomorrow{\modality{\narrowmath{tomorrow}}}
\newcommand\yesterday{\modality{\narrowmath{yesterday}}}
\newcommand\infinitely{\modality{\narrowmath{infinitely}}}
\newcommand\final{\modality{\narrowmath{finally}}}
\newcommand\arity{\f{arity}}
\newcommand\correct[1]{\tf{correct}[#1]}
\newcommand\correctterm{correct\xspace}
\newcommand\pointwise[1]{\tf{pointwise}[#1]}
\newcommand\pointwiseterm{pointwise\xspace}
\newcommand\incorrect[1]{\tf{incorrect}[#1]}
\newcommand\uniform[1]{\tf{uniform}[#1]}
\newcommand\udfn{\tf{undef}}
\newcommand\gslt{\f{GSLT}}
\newcommand\dense{\f{dense}}
\newcommand\noi{\f{noi}}

\newcommand{\upperset}[2]{
  \underset{
            \text{\raisebox{0.9ex}{\smash{\scalebox{0.8}{$#1$}}\ \ }}%
           }
           {
            \text{\raisebox{0ex}{\smash{$#2$}}}
           }
                       }
\newcommand{\convergesto}[1]{\mathrel{\upperset{#1}{\dotarrow}}}
\newcommand\betwixt{\asymp}
\newcommand\notbetween{\mathbin{\cancel{\between}}}
\newcommand\notintertwinedwith{\mathrel{\notbetween}}
\newcommand\intersectswith{\between}
\newcommand\notintersectswith{\notbetween}
\newcommand\stronglybetween{\mathbin{\rlap{$\between$}\hspace{-1.5pt}\between}}
\newcommand\xxtwined{$2$-twined\xspace}
\newcommand\xxxtwined{$3$-twined\xspace}
\newcommand\intertwined[1]{#1_{\between}}
\newcommand\intertwinedwith{\mathrel{\between}}
\newcommand\stronglyintertwinedwith{\mathrel{\stronglybetween}}
\newcommand\leqk{\leq_{\hspace{-.7pt}\intertwinedwith}}
\newcommand\geqk{\geq_{\hspace{-.7pt}\intertwinedwith}}
\newcommand\cw{\leftrightarrow}  %
\newcommand\ctr{\rightarrow}  %
\newcommand\notctr{{\not\rightarrow}}  %
\newcommand\ctl{\leftarrow}  %
\newcommand\cti{\leq}  %
\newcommand\ctieq{{\stackrel{=}{\hookrightarrow}}}  %
\newcommand\notcti{{\not\hookrightarrow}}  %
\newcommand\nbhd[0]{\f{nbhd}}
\newcommand\interior[0]{\f{interior}}
\newcommand\kiss[0]{\f{kiss}}
\newcommand\community[0]{\f{K}}
\newcommand\framecommunity[0]{\f{k}}
\newcommand\kernel[0]{\f{ker}}
\newcommand\Kernel[0]{\f{Ker}}
\newcommand\witness[0]{\f{W}}
\newcommand\Kmod[1]{\tf K_{\hspace{-2pt}#1}}
\newcommand\cast[1]{#1^{\ast c}}
\newcommand\cclo[1]{#1^c}
\newcommand\mbot{\mathsf{e}}
\newcommand\mtop{\mathsf{f}}
\newcommand\isdef[2]{#1\bullet #2}
\newcommand\deffont[1]{{\bfseries #1}}
\newcommand\powerset{\f{pow}}
\newcommand\comppowerset{\f{pow}_{\hspace{-2pt}c}}
\newcommand\finpow{\f{fin}}
\newcommand\finto{\stackrel{\f{fin}}{\rightharpoonup}}
\newcommand\f[1]{\mathit{#1}}
\newcommand\tf[1]{\mathsf{#1}}
\DeclareMathAlphabet{\mathsb}{OT1}{cmss}{sbc}{n}
\newcommand\ns[1]{{\mathsb{#1}}}
\usepackage{slantsc}
\newcommand\theory[1]{\ensuremath{\text{\itshape\scshape #1}}\xspace}
\newcommand\ThyPax{\theory{TyPax}}
\newcommand\ThySPax{\theory{TySPax}}
\newcommand\ThyBB{\theory{TyBB}}
\newcommand\finsubseteq{\mathbin{\subseteq_{\text{\it fin}}}}
\newcommand\at{\text{@}}
\newcommand\tx{\f{tx}}
\newcommand\ty{\f{ty}}
\newcommand\ctx{\f{ctx}}
\newcommand\txs{\f{txs}}
\newcommand\utxi{\f{utxi}}
\newcommand\utxo{\f{utxo}}
\newcommand\stx{\f{stx}}
\newcommand\utxio{\f{utxio}}
\newcommand\UTxIO{\f{UTxIO}}
\newcommand\inlinehask[1]{\lstinline[language=hask]{#1}}
\newcommand\inlinesolidity[1]{\lstinline[language=solidity]{#1}}
\newcommand\valid{\f{valid}}
\newcommand\liff{\Longleftrightarrow}
\newcommand\limp{\Longrightarrow}
\newcommand\nlimp{\centernot{\Longrightarrow}}
\newcommand\aeq{\mathrel{=_{\alpha}}}
\newcommand\ssm{{{:}\text{=}}}
\DeclareMathSymbol{\shortminus}{\mathbin}{AMSa}{"39}
\newcommand\minus{{\shortminus}}
\newcommand\plus{{+}}
\newcommand\Forall[1]{\forall #1.}
\newcommand\Exists[1]{\exists #1.}
\newcommand\sm{\mapsto}
\newcommand\rewritesto{\longrightarrow}
\newcommand\cent{\vdash}
\newcommand\ncent{\not\vdash}
\newcommand\ment{\vDash}
\newcommand\nment{\not\vDash}
\newcommand\mentval{\ment_{\avaluation}}
\newcommand\nmentval{\nment_{\avaluation}}
\newcommand\supp{\f{supp}}
\newcommand\boundary{\f{boundary}}
\newcommand\lmodel{[\hspace{-0.2em}[}
\newcommand\rmodel{]\hspace{-0.2em}]}
\newcommand\model[1]{{\lmodel #1 \rmodel}}
\newcommand\modellabel[2]{{\lmodel #1 \rmodel}_{#2}}
\newcommand\synmodel[1]{{\langle #1 \rangle}}
\newcommand\mone{{\text{-}1}}
\newcommand\fv{\f{fv}}
\newcommand\Ngeqz{\mathbb N_{\scriptscriptstyle\geq 0}}
\newcommand\Time{\Ngeqz}
\newcommand\weakmodusponens{weak modus ponens (Proposition~\ref{prop.mp.for.tnotor}(\ref{item.mp.for.tnotor}))\xspace}
\newcommand\strongmodusponens{strong modus ponens (Proposition~\ref{prop.mp.for.tnotor}(\ref{item.mp.for.timpc}))\xspace}
\newcommand\weakmodusponensnoref{weak modus ponens\xspace}
\newcommand\strongmodusponensnoref{strong modus ponens\xspace}

\makeatletter
\DeclareRobustCommand{\barcent}{\mathbin{\mathpalette\barcent@@\relax}}
\newcommand{\barcent@@}[2]{%
  \vbox{\offinterlineskip
    \sbox\z@{$\m@th#1\cent$}%
    \ialign{%
      \hfil##\hfil\cr
      $\m@th#1{}_{\minus}\kern-\scriptspace$\cr
      \noalign{\kern-.3\ht\z@}
      \box\z@\cr
    }%
  }%
}
\makeatother

\makeatletter
\def\pmb@#1#2{\setbox8\hbox{$\m@th#1{#2}$}%
  \setboxz@h{$\m@th#1\mkern-.1mu$}\pmbraise@\wdz@
  \binrel@{#2}%
  \dimen@-\wd8 %
  \binrel@@{%
    \mkern-.1mu\copy8 %
    \kern\dimen@\mkern-.2mu\copy8 %
    \kern\dimen@\mkern-.3mu\copy8 %
    \kern\dimen@\mkern-.4mu\copy8 %
    \kern\dimen@\mkern.1mu\copy8 %
    \kern\dimen@\mkern.2mu\copy8 %
    \kern\dimen@\mkern.3mu\copy8 %
    \kern\dimen@\mkern.0mu\raise\pmbraise@\copy8 %
    \kern\dimen@\mkern.4mu\box8 %
           }%
}
\makeatother


\newcommand\overbracebig[2]{\overbrace{#2}^{\mbox{$#1$}}}
\newcommand\compressthislight[1]{{\hspace{.8pt}\raisebox{.5pt}{\scalebox{.85}{$#1$}}\hspace{.2pt}}}
\newcommand\compressthis[1]{\pmb{\compressthislight{#1}}}
\newcommand\tempty{\compressthis{\varnothing}}
\newcommand\tcap{{\compressthis{\cap}}}
\newcommand\tcup{{\compressthis{\cup}}}
\newcommand\tsubseteq{{\compressthis\subseteq}}
\newcommand\tneg{{\pmb\neg}}
\newcommand\tnegnf{{\pmb{\sim}}}
\newcommand\tnegneg{{\pmb{{\sim}{\sim}}}}
\newcommand\ttop{{\pmb\top}}
\newcommand\tbot{{\pmb\bot}}
\newcommand\teq{{\pmb{\text{=}}}}
\newcommand\tneq{{\pmb{\neq}}}
\newcommand\tand{\mathrel{\pmb\wedge}}
\newcommand\tor{\mathbin{\pmb\vee}}
\newcommand\tgg{{\compressthis\gg}}
\newcommand\tleq{\compressthis{\leq}}
\newcommand\tlesssim{\compressthis{\lesssim}}
\newcommand\tapprox{\compressthis{\approx}}
\newcommand\timp{{\compressthis{\Rightarrow}}}
\newcommand\ntimp{{\pmb\not\Rightarrow}}
\newcommand\tnotor{\mathrel{\arrowr}}
\newcommand\tnotore{\mathrel{\arrower}}
\newcommand\tlatticeiff{\mathrel{\arrowlr}}
\newcommand\tlatticeiffec{\mathbin{\arrowelr}}
\newcommand\timpc{\mathrel{\arrowrr}}
\newcommand\timpcc{\mathrel{\arrowrrr}}
\newcommand\timpcleft{\mathrel{\arrowll}}
\newcommand\timpec{\mathrel{\arrowerr}}
\newcommand\tlatticeiffc{{\pmb{\bar\equiv}}}
\newcommand\tiff{\compressthis{\Leftrightarrow}} %
\newcommand\tiffc{\mathbin{\arrowlrr}}
\newcommand\tiffcc{\mathbin{\arrowllrr}}
\newcommand\tiffec{\mathbin{\arrowelrr}}
\newcommand\tlam{{\pmb\lambda}}
\newcommand\tsim{\compressthis{{\approx}}}
\newcommand\tin{{\pmb{\in}}}
\newcommand\ntin{{\pmb{\not\in}}}
\newcommand\qall{{\compressthis{\mathsf Q}}}
\newcommand\tall{{\compressthis{\forall}}}
\newcommand\texi{{\compressthis{\exists}}}
\newcommand\texiunique{{\compressthis{\exists\hspace{0.5pt}!}}}
\newcommand\texiaffine{{\compressthis{\exists_{01}}}}
\newcommand\texiR{{\compressthis{\exists_{\mathsf R}}}}
\newcommand\tlot{\textbf{?}}
\newcommand\AllBut[2]{\tf{AllBut}_{#2}(#1)}
\newcommand\Nset[1]{\{0\dots #1\minus 1\}}

\newcommand\oimplies{\mathbin{\rlap{$+$}\hspace{.7pt}{\supset}}}

\newcommand\QLogic{Coalition Logic\xspace}

\makeatletter
\newcommand{\circlearrow}{}%
\DeclareRobustCommand{\circlearrow}{%
  \mathrel{\vphantom{\shortrightarrow}\mathpalette\circle@arrow\relax}%
}
\newcommand{\circle@arrow}[2]{%
  \m@th
  \ooalign{%
    \hidewidth$#1\circ\mkern1mu$\hidewidth\cr
    $#1\longrightarrow$\cr}%
}
\makeatother
\newcommand\minusoplus{{\pmb{\circlearrow}}} %

\makeatletter
\newcommand*\bigcdot{\mathpalette\bigcdot@{.5}}
\newcommand*\bigcdot@[2]{\mathbin{\vcenter{\hbox{\scalebox{#2}{$\m@th#1\bullet$}}}}}
\makeatother

\newcommand\act{{\cdot}}
\newcommand\pact{{\cdot}}
\newcommand\mact{\mathbin{\bigcdot}}

\usepackage{datetime}
\yyyymmdddate

\begin{document}
\title{A declarative approach to specifying distributed algorithms using three-valued modal logic}
\newcommand\titlerunning{\emph{Declarative distributed algorithms via modal logic}} %
\newcommand\authorrunning{\emph{Murdoch J. Gabbay and Luca Zanolini}}
\author{Murdoch J. Gabbay \affil{Heriot-Watt University, Edinburgh, UK}
Luca Zanolini \affil{Ethereum Foundation, London, UK}}

\begin{abstract}
  We present \emph{\QLogic}, a three-valued modal fixed-point logic designed for declaratively specifying and reasoning about distributed algorithms, such as the Paxos consensus algorithm.

  Our methodology represents a distributed algorithm as a logical theory, enabling correctness properties to be derived directly within the framework—or revealing logical errors in the algorithm's design when they exist.
  
  \QLogic adopts a declarative approach, specifying the overall logic of computation without prescribing control flow. Notably, message-passing is not explicitly modeled, distinguishing our framework from approaches like TLA+. This abstraction emphasises the logical essence of distributed algorithms, offering a novel perspective on their specification and reasoning.
  
  We define the syntax and semantics of \QLogic, explore its theoretical properties, and demonstrate its applicability through a detailed treatment of the Paxos consensus algorithm. By presenting Paxos as a logical theory and deriving its standard correctness properties, we showcase the framework’s capacity to handle non-trivial distributed systems.
  
  We envision \QLogic as a versatile tool for specifying and reasoning about distributed algorithms. The Paxos example highlights the framework’s ability to capture intricate details, offering a new lens through which distributed algorithms can be specified, studied, and checked.


\keywords{Modal logic, distributed consensus} 
\end{abstract}


\maketitle
\thispagestyle{empty}

\tableofcontents

\section{Introduction}
\label{sect.intro}

\subsection{Initial overview}

\subsubsection*{Introducing \QLogic}
\emph{\QLogic}is a logical framework for studying distributed algorithms in a \emph{declarative} manner.
In our framework an algorithm is presented as a logical theory (a signature and a set of axioms), and then its correctness properties are derived from those axioms.
We apply our framework to present and study an axiomatisation of Paxos, a canonical distributed consensus algorithm. 

Logic \emph{per se} is widely-applied in distributed algorithms: for example the logic TLA+ (based on Linear Temporal Logic) is an industry standard to model and verify distributed protocols~\cite{DBLP:journals/toplas/Lamport94,DBLP:books/aw/Lamport2002,DBLP:conf/charme/YuML99}. 
However, TLA+ specifications describe state changes in a sequential, step-by-step manner.
A \QLogic axiomatisation feels more abstract: we dispense with step-by-step descriptions of the algorithm and with explicit message-passing, to capture an abstract but still entirely identifiable essence of the algorithm.

Abstracting an algorithm to a logical theory \emph{per se} is also not new: \emph{declarative programming} is based on this.
However, turning a distributed algorithm (like Paxos) into an axiomatic theory, is new, and this paper shows that it can be done.
Examples follow below and in the body of the text.\footnote{Or, if the reader wants to skip the preliminaries and jump to a toy example that illustrates our method, then see Subsection~\ref{subsect.simple}.}

For reasons of length, we focus on the example of Paxos.
This is a canonical example with sufficient complexity to be interesting, but we have checked that \QLogic can be adapted for other algorithms; it is intended as a general logical tool that can be tailored to particular case studies.
In the long run, we hope this will help do all the things that declarative techniques are supposed to: help find errors, decrease development times, increase productivity, and lead to new mathematics.
 
\subsubsection*{What is declarative programming}
By `declarative' we mean that we specify the logic of a computation without specifying its control flow.

Thus our axiomatisation of Paxos (in Figure~\ref{fig.logical.paxos}) relates to a concrete Paxos implementation and/or its TLA+ specification, much as this axiomatic, algebraic specification of the factorial function\footnote{By `axiomatic' we mean that the function is specified just by asserting properties that are true of it.  By `algebraic' we mean that these assertions are purely equational, i.e. we just assert some equalities between terms.}
$$
0!=1
\qquad 
(n+1)!=(n+1)*n!
$$
relates to this concrete implementation of factorial in BASH
\begin{quote}
{\tt\small\hspace{-2em} value=1; while [ \$number -gt 1 ]; do value=\$((value * number)); number=\$((number - 1)); done}
\end{quote}
The reader may be familiar with programming languages based on declarative principles, such as Haskell, Erlang, or Prolog.
Their design philosophy is that we should assert things that we want to be true, and leave it to the compiler to figure out how best to make this happen.
Contrast with C, which is a canonical imperative language.\footnote{We are not interested here in taking sides.  With a variety of tools to choose from, we pick the one that is best for the job.  What matters is that the widest possible \emph{variety} of levels of abstraction is made available.}

This paper %
explores a methodology for specifying distributed algorithms which sits decidedly at the `Haskell/Erlang/Prolog' side of the spectrum, compared to TLA+ which (while still quite high-level) sits more at the C side of the spectrum. 

Consider a tiny concrete example of a simple voting system: we will briefly exhibit this in \QLogic, just to give a flavour of how it works.
Consider a voting system such that 
\begin{quote}
`if a quorum of voters vote for candidate $c$, then $c$ wins the election'.
\end{quote}
We write this axiomatically simply as
$$
(\QuorumBox\tf{vote}(c)) \timpc \tf{wins}(c) .
$$ 
The $\QuorumBox$ is a logical modality that is true when there exists a quorum of participants at which $\tf{vote}(v)$ holds.\footnote{Details of this modality in our framework are in Proposition~\ref{prop.pointwise.dense.char}(\ref{item.dense.char.pointwise.QuorumBox}) and Lemma~\ref{lemm.allbut.concretely}.}

Even without going into the details of what these symbols mean, it is clear that details of message-passing, of how votes are communicated, and of how they are counted up, are absent.
In fact in this case the predicate is just a direct rendering of the English statement above it, without specifying how it is concretely implemented.
Thus, in keeping with a declarative philosophy, we specify what we want the algorithm to do, but not how it should do it: it is the point of the axiom to express formally and precisely our idea --- and to add as little as possible to that, beyond what is necessary to make it logically precise.


\subsubsection*{Who should read this document}
We have in mind two communities: 
\begin{enumerate*}
\item
\emph{logicians}, who might be interested in exploring a logic with an interesting combination of features (three-valued modal fixedpoint logic) with specific practical motivations in programming distributed systems, and 
\item
\emph{researchers in distributed systems}, who might find \QLogic useful for analysing and communicating their designs.
\end{enumerate*}
We will be particularly interested in consensus algorithms as applied to blockchain, because this is an important and safety-critical use case (the second author is part of the core team developing consensus algorithms for Ethereum).

\subsubsection*{References to key points in the text}

The reader who might want to jump to where the action is and read backwards and forwards from there, is welcome to proceed as follows:
\begin{itemize*}
\item
The axioms of Declarative Paxos are in Figure~\ref{fig.logical.paxos}.
\item
The syntax and semantics of \QLogic are in Figures~\ref{fig.predicate.syntax}, \ref{fig.3.phi.f}, and~\ref{fig.3.derived}.
\item
Truth-tables for three-valued logic and its connectives are in Figure~\ref{fig.3}.
\item
An outline of the Paxos algorithm (for the reader who might like to be reminded what it is) is in Remark~\ref{rmrk.high-level.paxos}.
\item
See also toy axiomatisations for an imaginary `simple' protocol in Subsection~\ref{subsect.simple}, designed to be vaguely reminiscent of Paxos, but simplified almost to triviality for the sake of exposition.
\end{itemize*}
 

\subsection{Motivation and background (further details)}

Let us unpack what consensus in distributed systems (like blockchains) is, why it matters, and why it is hard:

\subsubsection*{Distributed systems}
A distributed computing system is is one that distributes its computation over several machines.
Distributed systems are essential to computing: they enable a network of parties to collaborate to achieve common goals~\cite{cachinbook}. 

Being distributed is attractive to software engineers because it can give us desirable properties, including some or all of: scalability, redundancy, lower latency, reliability, and resilience.
For this reason, almost any nontrivial commercial software system that the reader might care to name is implemented in a distributed manner.
Furthermore, blockchain systems (like Bitcoin and Ethereum\footnote{\url{https://bitcoin.org/en/} and \url{https://ethereum.org/en/}}) and peer-to-peer systems are inherently distributed by design; if they were centralised, then they would not even make sense.

\subsubsection*{The challenge of consensus}
Among the many challenges in distributed systems, \emph{consensus} represents a fundamental abstraction.
Consensus captures the problem of 
\begin{quote}
reaching agreement among multiple participants on a common value, in the face of possibly unreliable communication, and in the presence of possibly faulty participants.
\end{quote}
Consensus algorithms are an essential prerequisite for coordination in any distributed system, and they play a particularly critical role in blockchain networks.
In such networks, consensus is the mechanism that ensures all participants agree on the state of the ledger --- the blockchain's source of truth --- thereby solving problems like \emph{double spending} (manipulating the blockchain ledger to spend the same token twice).
Consensus protocols also enable the secure ordering of transactions, and they ensure the integrity and consistency of the blockchain, even in the face of adversarial actors.

Consensus protocols in blockchain systems are particularly challenging to design and implement correctly. 
Indeed, there are numerous examples in which incomplete or insufficiently-specified protocols have led to serious issues, such as network halts or vulnerabilities that adversaries could exploit.
For example:
\begin{enumerate*}
\item
Solana (\url{https://solana.com}) has experienced instances where its consensus mechanism contributed to temporary network halts~\cite{yakovenko2018solana,shoup2022poh}.
\item 
Ripple's (\url{https://ripple.com}) consensus protocol has been analysed for conditions in which safety and liveness may not be guaranteed~\cite{DBLP:conf/opodis/Amores-SesarCM20}. 
\item
Avalanche's (\url{https://www.avax.network}) consensus protocol has faced scrutiny regarding its assumptions and performance in adversarial scenarios~\cite{DBLP:conf/opodis/Amores-SesarCT22}.
\end{enumerate*}
Even mature blockchains like Ethereum are not immune to challenges.
Attacks such as \emph{balancing attacks}~\cite{DBLP:conf/sp/NeuTT21} and \emph{ex-ante reorganisation} (\emph{reorg}) attacks~\cite{DBLP:conf/fc/Schwarz-Schilling22} have highlighted vulnerabilities in Ethereum's protocol.

These examples underscore how subtle flaws or ambiguities in protocol design can have significant consequences, emphasising the need for rigorous analysis and well-defined specifications in this complex and safety-critical field.

Compounding the fact that consensus algorithms are foundational to the security and functionality of blockchain systems, is the fact that designing and implementing these algorithms is notoriously challenging.
Distributed algorithms --- including consensus protocols --- are difficult (at least for human brains) because of the inherent difficulty of coordinating multiple participants, which may be faulty, hostile, drifting in and out of network connectivity~\cite{sleepy}, and/or have different local clocks, states, agendas, and so forth.

In short, this is a notoriously difficult field which is full of `gotchas', and within this difficult field, consensus protocols are amongst the most intricate, demanding, and practically important objects of study. 

\subsubsection*{Introducing a logic-based declarative approach (applied to Paxos)}
In this paper, we develop a logic-based declarative approach to modeling distributed algorithms.
It stands in relation to existing techniques in roughly the same way as declarative programming languages stand to imperative ones. %

Technically speaking, \QLogic is a \emph{three-valued modal fixedpoint logic}. 
This means:
\begin{itemize*}
\item
There are three truth-values: $\tvT$ (true), $\tvB$ (no value returned, i.e. `crashed'), and $\tvF$ (false)~\cite{sep-logic-manyvalued}.
\item
Modalities let us express things like `some participant has property $\phi$' or `a quorum of participants has property $\psi$' or `$\phi$ has been true in the past' or `$\psi$ will be true in the future'.\footnote{In order, these are as follows in our logic: $\someone\phi$, $\QuorumBox\phi$, $\recent\phi$, and $\sometime\phi$.  Precise semantics are in Figures~\ref{fig.3.phi.f} and~\ref{fig.3.derived}.} 
\item
A fixedpoint operator $\mu$ imports the expressive power of iteration~\cite{wiki:Fixedpoint_logic}.
Consensus algorithms are often iterative, so it is natural to include this in our logic.
\end{itemize*} 
We will apply these constructs to give axiomatic models of the Paxos consensus protocol~\cite{DBLP:journals/tocs/Lamport98,lamport2001paxos,cachin2011paxos} and then we will state and derive standard correctness properties from the axioms.
Our example is elementary but nontrivial and strikes (we think) a good balance between being compact but still exhibiting enough richness of structure to illustrate the essential features of how our logic works. 


The first half of this paper will focus on the logic, and the second half of this paper will apply the logic to the examples.
The reader from a logic background might like to start with the first half and refer to examples for motivation; the reader from a distributed systems background might like to do the opposite and start with the examples in the second half, referring back to the logic as required to understand what is being done.
 
\subsubsection*{Related work using logic}
Logic-based approaches have long been used to model and verify distributed protocols, with TLA+~\cite{DBLP:journals/toplas/Lamport94,DBLP:books/aw/Lamport2002,DBLP:conf/charme/YuML99} being one of the most prominent frameworks in the field. 
Introduced by Leslie Lamport in 1999, TLA+ uses Linear Temporal Logic (LTL)~\cite{DBLP:conf/focs/Pnueli77} to describe state changes in a sequential, step-by-step manner, which has become an industry standard for specifying and verifying distributed algorithms.

But TLA+ does things in a very particular way, which is not declarative.
We would argue that it is good for specifying \emph{algorithms}, and so is not particularly \emph{declarative}.
Specifically, TLA+ excels at specifying and verifying system behaviours through state transitions; it is operational and state-based.
This can sometimes distance TLA+ protocol specifications from the invariants that they are intended to uphold.

This separation arises because TLA+ primarily describes how states evolve \emph{over time}, which can make it challenging to directly reason about the essential properties that protocols are designed to maintain. 
So, although TLA+ provides powerful tools for specifying and checking invariants, the process of verifying that these invariants are preserved across all possible state transitions can be complex and unintuitive, particularly when dealing with the intricacies of distributed algorithms.

\subsubsection*{Our approach}
Our novel framework leverages the mathematics of semitopology and three-valued logic. 

Semitopology is from previous work~\cite{gabbay:semdca,gabbay:semtad}.
It provides a natural way to describe quorums --- sets of participants that work together to progress the system state --- by representing them as open sets.
This perspective is particularly suited to distributed systems, where the behaviour of the system emerges from the interactions of coalitions its components, rather than from a sequence of steps by a central controller.

Three-valued logic (a good overview is in \cite[Chapter~5]{bergmann:intmvf}) helps us to reason about the (possibility for) faults that is inherent in distributed systems.
Unlike classical logic with its binary truth-values $\tvT$ and $\tvF$, three-valued logic introduces a third truth value $\tvB$ that captures the indeterminate states often encountered when dealing with faulty parties.
This turns out to give particularly natural and compact descriptions of distributed protocols, distinguishing between normal operation and faulty behaviour. 
It also provides a formal foundation that would be amenable to verification through tools like SAT solvers and theorem provers (building tooling for what we see in this paper would be future work), enhancing both the expressiveness and verifiability of protocol specifications.

Our approach is particularly relevant in the context of blockchain consensus protocols, which are distributed and decentralised systems that require robust and clear specifications to function correctly.
It is not an exaggeration to say that blockchains in particular have put a rocket underneath the development of new consensus protocols~\cite{DBLP:conf/wdag/CachinV17}.
Many of these are not fully specified or formally verified. 
Often, they are even described informally in blog posts and other non-academic sources --- leading to confusion about their properties and behaviours. 
This lack of formalisation creates a fragmented landscape, where the guarantees and limitations of the protocols remain unclear, hindering their adoption and undermining trustworthiness.

We hope that by providing a logic-based declarative language, and the ability to axiomatise consensus protocols, our framework will offer a way to formally specify, reason about, and verify the properties of the many distributed systems under development --- thus lowering costs, increasing reliability, and reducing development time. 
We would further argue that effective formalisation is essential in reducing ambiguity in the design and evaluation of the many new consensus mechanisms being developed in the blockchain space in particular.

Our approach bridges the gap between the declarative specification of invariants and the practical verification of protocol correctness. 

\subsubsection*{We will explain everything, and that's a feature}

Different readers will find some parts of this paper easy, and other parts less easy --- but which parts these are depends on the reader and their background.\footnote{Not a hypothetical: the first author (with a logic background) took a long time to get his head around how Paxos works, whereas the second author (who knows Paxos well and is familiar with logic in general terms, but is not a logician) was not immediately familiar with modal fixedpoint logic and models of axiomatic logical theories.  The community of readers who approach this paper and are already well-versed in both fields may not be the largest.  We hope this paper will help ameliorate this lamentable state of affairs.}

For example:
\begin{itemize*}
\item
three-valued modal logic and its models, and technical results like the Knaster-Taski fixedpoint theorem used in Proposition~\ref{prop.mu.fixedpoint} may be known to the reader with a background in logic (like the first author), but perhaps not to a distributed systems theorist (like the second author); conversely, 
\item
the Paxos algorithm, and standard technical devices like the discussion motivating $n$-twined semitopologies in Remark~\ref{rmrk.motivate.n-twined}, are the bread-and-butter of a distributed systems expert's work --- but this application might be new to someone from a logic background. 
\end{itemize*}
We play safe and explain everything. 
If the reader finds we are going a bit slow in places, then better that than going too fast and skipping details that might help some other reader follow the argument.
 


\subsection{Structure of the paper}

\begin{enumerate*}
\item
Section~\ref{sect.intro} is the Introduction; You are Here.
The rest of the paper splits into two halves: Sections~\ref{sect.three.truth.values} to Section~\ref{sect.n-twined} have to do with logic and semitopologies; and Sections~\ref{sect.declarative.paxos} to~\ref{sect.simpler.paxos} use \QLogic to declaratively specify and study Paxos.
\ \\
\item
Section~\ref{sect.three.truth.values} introduces the basics of three-valued logic, including notable connectives ($\tnotor$ and $\timpc$) and modalities ($\modT$ and $\modTB$).
We also introduce a notion of validity for truth-values.
\item
Section~\ref{sect.quorum.logic} defines semitopologies (our model of the quorums that are ubiquitous in consensus algorithms), then it builds the syntax and semantics of \QLogic.
\item
Section~\ref{sect.predicates} studies \QLogic predicates; notably developing some logical equivalences, and well-behavedness criteria like being \emph{pointwise}.
\item 
Section~\ref{sect.n-twined} returns to semitopologies, and shows how antiseparation properties can be represented in logical form in \QLogic.
One particular antiseparation property corresponds to the \emph{quorum intersection property} which is used in the correctness proofs for Paxos (see Remark~\ref{rmrk.ntwined.and.the.literature}), but we will put this fact in a more general logical/semitopological context.
\ \\
\item
Section~\ref{sect.declarative.paxos} introduces Declarative Paxos; this is representative for the kind of algorithm that we built \QLogic to express.
We take our time to explain how Paxos works, state the axiomatisation, and discuss it in detail.
\item
Section~\ref{sect.paxos.correctness.properties} states the standard correctness properties for Paxos, phrases them in declarative form, and then proves them from the axioms of Declarative Paxos.
What all our work has accomplished is to convert proving correctness properties into a (relatively straightforward) argument of proving facts of an abstract model from its axioms.
\item
Section~\ref{sect.simpler.paxos} revisits the axioms of Declarative Paxos from Section~\ref{sect.declarative.paxos}, in the light of the proofs in Section~\ref{sect.paxos.correctness.properties} --- and then prunes them down to a minimal theory required to prove correctness.

We call this minimal theory \emph{Simpler Declarative Paxos}: it is simpler, and because further details have been elided, it is more abstract.
The reader may prefer Declarative Paxos to Simpler Declarative Paxos because it is more identifiably related to the Paxos algorithm, but if we want to prove correctness properties, then Simpler Declarative Paxos distils the logical essence of what makes those proofs work. 
\ \\
\item
Section~\ref{sect.conclusions} concludes with conclusions and reflections on future work.
\end{enumerate*} 


\section{Three truth-values}
\label{sect.three.truth.values}

\subsection{Basic definition of $\THREE$, and some truth-tables}

\begin{defn}[Three-valued truth-values]
\label{defn.THREE}
Let $\THREE$ be the poset of \deffont{truth-values} with elements $\tvF$ (false), $\tvB$ (both), and $\tvT$ (true) ordered as 
$$
\tvF<\tvB<\tvT.
$$
\end{defn}

\begin{figure} %
$$
\begin{array}{c}
\begin{array}{c@{\qquad}c@{\qquad}c@{\qquad}c@{\qquad}c@{\qquad}c}
\begin{array}{c c c c}
p\tand q & \tvT & \tvB & \tvF
\\
\tvT &  \tvT & \tvB & \tvF
\\
\tvB &  \tvB & \tvB & \tvF 
\\
\tvF &  \tvF & \tvF & \tvF 
\end{array}
&
\begin{array}{c c c c}
p\tor q & \tvT & \tvB & \tvF
\\
\tvT &  \tvT & \tvT & \tvT
\\
\tvB &  \tvT & \tvB & \tvB
\\
\tvF &  \tvT & \tvB & \tvF
\end{array}
&
\begin{array}{c c c c}
p\tnotor q & \tvT & \tvB & \tvF
\\
\tvT &  \tvT & \tvB & \tvF
\\
\tvB &  \tvT & \tvB & \tvB
\\
\tvF &  \tvT & \tvT & \tvT
\end{array}
&
\begin{array}{c c c c}
p \timpc q & \tvT & \tvB & \tvF
\\
\tvT       & \tvT & \tvF & \tvF
\\
\tvB       & \tvT & \tvB & \tvB
\\
\tvF       & \tvT & \tvT & \tvT
\end{array}
\end{array}
\\[7ex]
\begin{array}{c@{\qquad}c@{\qquad}c@{\qquad}c@{\qquad}c@{\qquad}c@{\qquad}c}
\begin{array}{cc}
\tneg p
\\
\tvT & \tvF
\\
\tvB & \tvB
\\
\tvF & \tvT
\end{array}
&
\begin{array}{cc}
\modT p 
\\
\tvT & \tvT
\\
\tvB & \tvF
\\
\tvF & \tvF
\end{array}
&
\begin{array}{cc}
\modF p 
\\
\tvT & \tvF
\\
\tvB & \tvF
\\
\tvF & \tvT
\end{array}
&
\begin{array}{cc}
\modB p
\\
\tvT & \tvF
\\
\tvB & \tvT
\\
\tvF & \tvF
\end{array}
&
\begin{array}{cc}
\modTB p  
\\
\tvT & \tvT
\\
\tvB & \tvT
\\
\tvF & \tvF
\end{array}
&
\begin{array}{cc}
\modTF p  
\\
\tvT & \tvT
\\
\tvB & \tvF
\\
\tvF & \tvT
\end{array}
&
\begin{array}{cc}
\modFB p  
\\
\tvT & \tvF
\\
\tvB & \tvT
\\
\tvF & \tvT
\end{array}
\end{array}
\\[7ex]
\begin{array}{c@{\qquad}c@{\qquad}c@{\qquad}c@{\qquad}c@{\qquad}c}
\begin{array}{c c c c}
p\timp q & \tvT & \tvB & \tvF
\\
\tvT &  \tvT & \tvB & \tvF
\\
\tvB &  \tvT & \tvB & \tvF 
\\
\tvF &  \tvT & \tvT & \tvT 
\end{array}
&
\begin{array}{c c c c}
p\tleq q & \tvT & \tvB & \tvF
\\
\tvT &  \tvT & \tvF & \tvF
\\
\tvB &  \tvT & \tvB & \tvF 
\\
\tvF &  \tvT & \tvT & \tvT 
\end{array}
\end{array}
\end{array}
$$
\emph{Above, the vertical axis of a table indicates values for $p$; the horizontal axis (if nontrivial) denotes values for $q$.}
\caption{Truth-tables for some operations on $\THREE$ (Remark~\ref{rmrk.three.elementary})}
\label{fig.3}
\end{figure}


\begin{figure} %
$$
\begin{array}{c@{\qquad}c@{\qquad}c@{\qquad}c@{\qquad}c@{\qquad}c}
\begin{array}{c c c c}
p\tlatticeiff q & \tvT & \tvB & \tvF
\\
\tvT &  \tvT & \tvB & \tvF
\\
\tvB &  \tvB & \tvB & \tvB
\\
\tvF &  \tvF & \tvB & \tvT
\end{array}
&
\begin{array}{c c c c}
p \tiffcc q & \tvT & \tvB & \tvF
\\
\tvT       & \tvT & \tvF & \tvF
\\
\tvB       & \tvF & \tvB & \tvB
\\
\tvF       & \tvF & \tvB & \tvT
\end{array}
&
\begin{array}{c c c c}
p \equiv q & \tvT & \tvB & \tvF
\\
\tvT       & \tvT & \tvF & \tvF
\\
\tvB       & \tvF & \tvT & \tvF
\\
\tvF       & \tvF & \tvF & \tvT
\end{array}
\end{array}
$$
\caption{Truth-tables for notions of equivalence (Definition~\ref{defn.equivalent})}
\label{fig.3.iff}
\end{figure}


\begin{rmrk}
\label{rmrk.three.elementary}
We start with some elementary observations about $\THREE$.
\begin{enumerate*}
\item\label{item.three.elementary.bounded.lattice}
$\THREE$ %
has bottom element $\tvF$, top element $\tvT$, and meets $\tand$ and joins $\tor$.
This makes it a \emph{finite lattice}.
\item
$\THREE$ has a negation operation $\tneg$.
\item
$\THREE$ has several natural modalities that identify truth-values, including $\modT$, $\modF$, $\modB$, $\modTB$, and $\modTF$.

For example, the reader can examine the truth-table for $\modTF$ in Figure~\ref{fig.3} and see that it identifies elements of $\{\tvT,\tvF\}$.\footnote{One could also call these modalities \emph{characteristic functions} for various subsets of $\THREE$, by which we mean functions which (literally) return $\tvT$ on elements in the set and $\tvF$ on elements not in the set.}
\item
$\THREE$ also has multiple notions of implication and equivalence, some of which are illustrated in Figures~\ref{fig.3} and~\ref{fig.3.iff}.
We will discuss these implications in Subsection~\ref{subsect.3.imp}.
\end{enumerate*}
\end{rmrk}


\begin{rmrk}
We will use the third truth-value $\tvB$ to represent participants in a distributed algorithm \emph{crashing}.
The method is simple: the truth-value of a predicate describing the outcome of a crashed computation, is deemed to be $\tvB$.

This is similar in spirit to how in domain theory an additional value --- usually written $\bot$ --- is added to a denotation to represent nontermination; the function-value of a nonterminating computation is deemed to be $\bot$~\cite{abramsky:domt}.
The nonterminating computation does not return a value, but we represent this absence of a value in the denotation, using $\bot$. 

In our declarative axiomatisation of Paxos, we use $\tvB$ in much the same way, as a truth-value to describe the \emph{absence} of a return value (either $\tvT$ or $\tvF$) from participants that have crashed.
\end{rmrk}



\subsection{Valid and correct truth-values}



\begin{defn}
\label{defn.tv.ment}
\leavevmode
\begin{enumerate*}
\item\label{item.tv.ment.correct}
Define subsets $\threeValid,\threeCorrect,\threeTrue,\threeFalse\subseteq\THREE$ by
$$
\threeValid = \{\tvT,\tvB\}
\quad\text{and}\quad
\threeCorrect = \{\tvT,\tvF\} 
\quad\text{and}\quad
\threeTrue = \{\tvT\} 
\quad\text{and}\quad
\threeFalse = \{\tvF\} .
$$ 
\item
We refer to $\f{tv}\in\THREE$ as being \deffont{valid} / \deffont{invalid},\ \deffont{correct} / \deffont{incorrect},\ and \deffont{true} / \deffont{untrue},\ and \deffont{false} / \deffont{unfalse}\footnote{\dots `unfalse' is `$\tvT$-or-$\tvB$', not `$\tvT$'.  Similarly, `untrue' is $\tvB$-or-$\tvF$.  We cannot resist stating the following gem of three-valued clarity: not-$\tvT$, is not $\tneg\tvT$, and not-$\tvF$ is not $\tneg\tvF$, and not-$\tvB$ \emph{is} not-$\tneg\tvB$, because $\tvB$ is $\tneg\tvB$.} depending on whether it is / is not in $\threeValid$, $\threeCorrect$, $\threeTrue$, and $\threeFalse$ respectively.
\item\label{item.tv.ment}
Define $\ment\f{tv}$ \ (read `\deffont{$\f{tv}$ is valid}') by
$$
\ment\f{tv} 
\quad\text{when}\quad
\f{tv}\in\threeValid ,
\ \ \text{which just means}\ \ 
\f{tv}\in\{\tvT,\tvB\} .
$$
Thus, `$\f{tv}$ is valid' holds precisely when $\f{tv}$ is a valid truth-value.
\end{enumerate*}
\end{defn}

\begin{rmrk}
\label{rmrk.motivate.validity}
We take a moment to discuss Definition~\ref{defn.tv.ment}:
\begin{enumerate}
\item
We can think of $\tvT$ and $\tvF$ as specific `yes/no' results from the successful completion of some computation.
Therefore, we will call $\tvT$ and $\tvF$ \emph{correct}.
\item
We can think of $\tvB$ as representing failure of a computation to complete correctly --- e.g. because a participant has crashed and so can return neither $\tvT$ nor $\tvF$.
\item\label{item.tb.valid}
In Definition~\ref{defn.tv.ment} we set $\threeValid=\{\tvT,\tvB\}$ and define $\ment\f{tv}$ when $\f{tv}\in\threeValid$, and we call $\tvT$ and $\tvB$ \emph{valid}.
\end{enumerate}
It might seem strange at first to let $\tvB$ be valid; if $\tvB$ represents the (non)behaviour of a crashed participant, then surely crashing should \emph{not} be valid?

Not necessarily.
To see why, consider that many of assertions of interest that we will consider will have the form 
\begin{quote}
`if $p$ is not crashed, then $\ment\phi$ at $p$', 
\end{quote}
or equivalently 
\begin{quote}
`either $p$ is crashed, or $\ment\phi$ at $p$'.
\end{quote}
If we let $\tvB$ be valid (i.e. if we define $\ment\tvB$ to hold), then we can simplify the above to the conveniently succinct assertion 
\begin{quote}
`$\ment\phi$ at $p$'.
\end{quote}
The `if $p$ is not crashed' precondition gets swallowed up into the fact that if $p$ is crashed then $\phi$ returns $\tvB$, and $\ment\tvB$ so there is nothing to prove. 
\end{rmrk}


\begin{rmrk}
Later on in Definition~\ref{defn.validity.judgement}(\ref{item.ctx.ment.phi}) we will extend validity to predicates-in-context, just by evaluating a predicate-in-context to a truth-value using a denotation function, and checking that truth-value is valid in the sense we define here. 
\end{rmrk}


We make some elementary observations about $\THREE$ and Definition~\ref{defn.tv.ment}:
\begin{lemm}
\label{lemm.tand.tor}
Suppose $\f{tv}\in\THREE$ is a truth-value.
Then:
\begin{enumerate*}
\item
$\ment\tvT$ and $\ment\tvB$.
\item
$\nment\tvF$.
\item\label{item.tand.tor.tor}
$\ment\f{tv}\tor\f{tv}'$ if and only if ${\ment\f{tv}}\lor{\ment\f{tv}'}$.
\item
$\ment\f{tv}\tand\f{tv}'$ if and only if ${\ment\f{tv}}\land{\ment\f{tv}'}$.
\end{enumerate*}
\end{lemm}
\begin{proof}
Fact of Definition~\ref{defn.tv.ment} and the truth-tables in Figure~\ref{fig.3}.
\end{proof} 

Negation is conspicuous by its absence in Lemma~\ref{lemm.tand.tor}.
We consider this next:
\begin{lemm}
\label{lemm.para}
Suppose $\f{tv}\in\THREE$ is a truth-value.
Then:
\begin{enumerate*}
\item\label{item.para.em}
$\ment\f{tv}\tor\tneg\f{tv}$, and as a corollary $\nment\f{tv}$ implies $\ment\tneg\f{tv}$.
\item\label{item.para.para}
$\ment\f{tv}\tand\tneg\f{tv}$ is possible, and as corollaries: $\ment\f{tv}$ does not necessarily imply $\nment\tneg\f{tv}$; and it is possible that $\ment\f{tv}$ and $\ment\tneg\f{tv}$ both hold.
\end{enumerate*}
\end{lemm}
\begin{proof}
We consider each part in turn:
\begin{enumerate}
\item
Using Lemma~\ref{lemm.tand.tor}(3) it would suffice to show that $\ment\f{tv}$ or $\ment\tneg\f{tv}$.
So suppose $\nment\f{tv}$.
By Definition~\ref{defn.tv.ment}(\ref{item.tv.ment}) $\f{tv}\not\in\threeValid$, thus $\f{tv}=\tvF$ and by Figure~\ref{fig.3} $\tneg\f{tv}=\tvT\in\threeValid$ and so $\ment\tneg\f{tv}$.
\item
We just take $\f{tv}=\tvB$ and check from Figure~\ref{fig.3} that $\tneg\tvB=\tvB$.
We use Lemma~\ref{lemm.tand.tor}(1\&4).
\qedhere\end{enumerate}
\end{proof}

\begin{rmrk}
Some comments on Lemma~\ref{lemm.para}:
\begin{itemize}
\item
Lemma~\ref{lemm.para}(\ref{item.para.para}) notes a typical fact of validity in three-valued logic, that it is naturally \deffont{paraconsistent}, meaning that it is possible for a thing, and the negation of that thing, to both be valid.
This is a feature of $\tvB\in\THREE$, not a bug, and we will make good use of it later when we use $\tvB$ %
as a return value for `I return no correct truth-value'. 
\item
Lemma~\ref{lemm.para}(\ref{item.para.em}) asserts validity of the \emph{excluded middle} (`$\phi$-or-negation-of-$\phi$'), but its corollaries implicitly illustrate another fact of three-valued logic, by what they do \emph{not} say: 
looking at the proof of Lemma~\ref{lemm.para}(\ref{item.para.em}), it is clearly also true that $\ment\f{tv}\tor\modF\f{tv}$, and that $\ment\f{tv}\tor\modT\tneg\f{tv}$, and that $\nment\f{tv}$ implies $\ment\modF\f{tv}$, and also that $\nment\f{tv}$ implies $\ment\modT\tneg\f{tv}$.  

These are stronger assertions, so why did we not state them in the lemma instead?
Because we will not need the result in those forms later.  

In two-valued logic, if something is true then it is probably a useful lemma.
In contrast, in three-valued logic, there is so much more structure with the third truth-value that more things are \emph{true} than are necessarily \emph{useful}.
So, we have to pick and choose our lemmas, and it is not always the case that the strongest form of a result is the most useful.
\end{itemize}
\end{rmrk}


Lemmas~\ref{lemm.tv.ment.TF} and~\ref{lemm.TF.correct} are easy, but they spell out some useful facts about truth-values and validity judgements:
\begin{lemm}
\label{lemm.tv.ment.TF}
Suppose $\f{tv}\in\THREE$ is a truth-value.
Then:
\begin{enumerate*}
\item\label{item.tv.ment.TF.1}
$\threeCorrect\cap\threeValid=\threeTrue$.

In words: a truth-value is true if and only if it is correct and valid.
\item\label{item.tv.ment.TF.2b}
$\ment \modT\f{tv}$ if and only if $\f{tv}\in\threeTrue$. 

In words: `$\modT\f{tv}$ is valid' holds if and only if $\f{tv}$ is true.
\item\label{item.tv.ment.TF.2b.F}
$\ment \modF\f{tv}$ if and only if $\f{tv}\in\threeFalse$. 

In words: `$\modT\f{tv}$ is valid' holds if and only if $\f{tv}$ is true.
\item\label{item.tv.ment.TF.TB}
$\ment \modTB\f{tv}$ if and only if $\f{tv}\in\threeValid$ if and only if $\ment\f{tv}$. 

In words: `$\modTB\f{tv}$ is valid' if and only if $\f{tv}$ is valid.\footnote{Note that $\modTB\f{tv}$ is not necessarily \emph{literally equal} to $\f{tv}$, in the sense of $\modT\f{tv}\equiv\f{tv}$ from Definition~\ref{defn.equivalent}.  For example, $\modTB\tvB=\tvT\neq\tvB$.}
\item\label{item.tv.ment.TF.2}
$\ment \modTF\f{tv}$ if and only if $\f{tv}\in\threeCorrect$. 

In words: `$\modTF\f{tv}$ is valid' if and only if $\f{tv}$ is correct.
\end{enumerate*}
\end{lemm}
\begin{proof}
Part~\ref{item.tv.ment.TF.1} is a fact of sets.
Parts~\ref{item.tv.ment.TF.2b} and~\ref{item.tv.ment.TF.TB} and~\ref{item.tv.ment.TF.2} are just from the truth-table for $\modT$, $\modTB$, and $\modTF$ in Figure~\ref{fig.3} and Definition~\ref{defn.tv.ment}. 
\end{proof}


\begin{lemm}
\label{lemm.TF.correct}
Suppose $\f{tv}\in\THREE$ is a truth-value.
Then the following are all equivalent:
$$
\ment\modTF\f{tv}
\ \liff\ 
\modTB\f{tv}=\modT\f{tv}
\ \liff\ 
\modT\f{tv}=\f{tv} 
\ \liff\ 
\modFB\f{tv}=\modF\f{tv}
\ \liff\ 
\modF\f{tv}=\tneg\f{tv} 
$$ 
Informally --- we do not build the machinery to make this formal, but it may still be a helpful observation --- we can write that the modality $\modTF$ is equivalent to the modal equivalences $\modTB=\modT=\f{id}$ and $\modFB=\modF=\tneg$.
\end{lemm}
\begin{proof}
By Lemma~\ref{lemm.tv.ment.TF}(\ref{item.tv.ment.TF.2}) and Definition~\ref{defn.tv.ment}(\ref{item.tv.ment.correct}),
$\ment\modTF\f{tv}$ when 
$$
\f{tv}=\tvT \ \lor\ \f{tv}=\tvF
\quad\text{and when}\quad
\f{tv}\neq\tvB .
$$
We check the truth-tables in Figure~\ref{fig.3}, and see that all of the other statements above precisely characterise the same condition.
\end{proof}


\subsection{The rich implication and equivalence structure of $\THREE$}
\label{subsect.3.imp}

\subsubsection{Weak and strong implications, and others}


Three-valued logic supports sixteen different kinds of implication~\cite[note~5, page~22]{arieli:idepl}.
We discuss two of them in Definition~\ref{defn.three.imp}, and mention two more in Remark~\ref{rmrk.rich.implication}:
\begin{defn}
\label{defn.three.imp}
\leavevmode
\begin{enumerate*}
\item\label{item.three.elementary.weak.implication}
Define \deffont{weak implication} $\tnotor$, using $\tneg$ and $\tor$, as follows:
$$
\f{tv}\tnotor\f{tv}'=(\tneg\f{tv})\tor\f{tv}'.
$$
\item\label{item.three.elementary.strong.implication}
Define \deffont{strong implication} $\timpc$, using $\tnotor$ and $\modT$, as follows:
$$
\f{tv}\timpc\f{tv}'=\f{tv}\tnotor \modT\f{tv}'. 
$$
\end{enumerate*}
These and other encodings are assembled in Figure~\ref{fig.3.derived}.
\end{defn}

We motivate the weak and strong implications from Definition~\ref{defn.three.imp}(\ref{item.three.elementary.weak.implication}\&\ref{item.three.elementary.strong.implication}):
\begin{rmrk}[How to interpret $\tnotor$ and $\timpc$]
\label{rmrk.interpret.implication}
\leavevmode
For this Remark, let us interpret $\f{tv}\in\THREE$ as the result of a computation process which either succeeds and returns true (truth-value $\tvT$), or it succeeds and returns false (truth-value $\tvF$), or the process does not succeed (e.g. it crashes) and it returns an invalid value / no value (represented by the truth-value $\tvB$).
Recall also from Definition~\ref{defn.tv.ment}(\ref{item.tv.ment}) that $\ment\f{tv}$ means $\f{tv}\in\{\tvT,\tvB\}$.
Then by this reading:
\begin{enumerate*}
\item
$\ment\f{tv}$ means \emph{``the process $\f{tv}$ returns $\tvT$, if valid''}. 
\item
$\ment\f{tv}\tnotor\f{tv}'$ means \emph{``\emph{if} process~$\f{tv}$ returns true \emph{then} process~$\f{tv}'$ returns $\tvT$, if valid''}.

Implication does not imply sequentiality: we are just making a logical assertion about outcomes. 
\item
$\ment\f{tv}\timpc\f{tv}'$ means \emph{``\emph{if} process~$\f{tv}$ returns true \emph{then} process~$\f{tv}'$ is \emph{also} valid, and it returns $\tvT$''}.
\end{enumerate*}
\end{rmrk}


We can sum this up as a \emph{modus ponens} result, which relates the validity $\ment$ from Definition~\ref{defn.tv.ment}, the weak implication $\tnotor$ and strong implication $\timpc$ from Figure~\ref{fig.3}, to real-world implication:
\begin{prop}[Weak and strong modus ponens]
\label{prop.mp.for.tnotor}
Suppose $\f{tv},\f{tv}'\in\THREE$ are truth-values.
Then:
\begin{enumerate*}
\item\label{item.mp.for.tnotor}
$\ment\f{tv}\tnotor\f{tv}'$ if and only if\ \ $\ment\modT\f{tv}$ implies $\ment\f{tv}'$. 

In words: `$\f{tv}$ weakly implies $\f{tv}'$' is valid, if and only if $\f{tv}$ is true implies $\f{tv}'$ is valid (but $\f{tv}'$ need not be true). 

We may call this logical equivalence %
\deffont{weak modus ponens}.
\item\label{item.mp.for.timpc}
$\ment\f{tv}\timpc \f{tv}'$ if and only if $\ment\modT\f{tv}$ implies $\ment\modT\f{tv}'$. 

In words: `$\f{tv}$ strongly implies $\f{tv}'$' is valid if and only if $\f{tv}$ is true implies $\f{tv}'$ is true. 

We may call this logical equivalence \deffont{strong modus ponens}. 
\end{enumerate*}
\end{prop}
\begin{proof}
We consider each part in turn:
\begin{enumerate}
\item
By direct inspection of the truth-table for $\tnotor$ in Figure~\ref{fig.3}, using Definition~\ref{defn.tv.ment}(\ref{item.tv.ment}).
\item
By a similar direct inspection.
\qedhere\end{enumerate}
\end{proof}


\begin{rmrk}[Explicit and implicit implications]
\label{rmrk.rich.implication}
So far, the structure we have seen (the lattice structure, the negation, the modalities, and the two implications $\tnotor$ and $\timpc$) is reflected explicitly in the syntax of our logic in Figures~\ref{fig.predicate.syntax} and~\ref{fig.3.derived}.
For instance, Figure~\ref{fig.predicate.syntax} has a unary negation connective $\tneg$ whose denotation as per Figure~\ref{fig.3.phi.f} is precisely the negation operator $\tneg$ in Figure~\ref{fig.3}.
So far, so explicit.

However, $\THREE$ has further implication structure, which we do not mention explicitly in our logic, but which is still there, is important, and which appears implicitly in the proofs.
Notably: 
\begin{enumerate*}
\item\label{item.rich.implication.timp}
An implication $\timp$ defined by $\f{tv}\timp \f{tv}'=\tneg\f{tv}'\timpc\tneg\f{tv}=(\modTB\f{tv})\tnotor\f{tv}'$ has a truth-table spelled out in Figure~\ref{fig.3} (bottom left).
This can be read as ``if $\f{tv}$ is valid then $\f{tv}'$ is valid''.

$\timp$ is clearly implicit in results like Proposition~\ref{prop.someone.implies.someoneAll} -- see all the parts whose statement uses logical implication $\limp$ --- and in Proposition~\ref{prop.tv.ment.TF.model}.
\item
Another implication is implicit in the use of $\leq$ in results like Lemma~\ref{lemm.positive.monotonicity} and Corollary~\ref{corr.recent.everyone.commute}.
That symbol refers to the lattice order on $\THREE$, but this could easily be internalised in our logic --- e.g. we can set $\f{tv}\tleq \f{tv'}$ to be $(\f{tv}\timpc\f{tv}')\tand(\f{tv}\timp\f{tv}')$, to obtain the connective $\tleq$ also illustrated in Figure~\ref{fig.3} (bottom right).
\end{enumerate*}
So in fact, we use all four of $\timp$ and $\tleq$, and $\tnotor$ and $\timpc$.

The reason that $\tnotor$ and $\timpc$ make it into our logic as (sugared) syntax, and $\timp$ and $\leq$ do not, is that $\tnotor$ and $\timpc$ naturally arise in our axioms (see most of the axioms in Figures~\ref{fig.logical.paxos}) %
whereas we do not need $\timp$ or $\tleq$ in the same way for the axiomatisations in this paper.
Different axiomatisations of different algorithms might make different demands on our logic, and it would be straightforward to include (sugar for) connectives accordingly. 
\end{rmrk}

\subsubsection{Equivalence $\equiv$}

Our three-valued domain of truth-values $\THREE$ has rich notions of equivalence, to match its rich notions of implication.
We mention three of them:
\begin{defn}[Notions of equivalence in $\THREE$]
\label{defn.equivalent}
Define $\phi\tlatticeiff\phi'$ and $\phi\tiffcc\phi'$ and $\phi\equiv\phi'$ by:
$$
\begin{array}{r@{\ }l}
\phi\tlatticeiff\phi' =& (\phi\tnotor\phi')\tand(\phi'\tnotor\phi)
\\
\phi\tiffcc\phi' =& (\phi\timpc\phi')\tand(\phi'\timpc\phi)
\\
\phi\equiv\phi' =& (\modT\phi\tand\modT\phi')\tor(\modB\phi\tand\modB\phi')\tor(\modF\phi\tand\modF\phi')
\end{array}
$$
Truth-tables for $\tlatticeiff$, $\tiffcc$, and $\equiv$ are given in Figure~\ref{fig.3.iff}, and we unpack this further in Remark~\ref{rmrk.equiv}.
\end{defn}


\begin{rmrk}
\label{rmrk.equiv}
We can read 
\begin{itemize*}
\item
$\f{tv}\tlatticeiff \f{tv}'$ as `$\f{tv}$ is valid if and only if $\f{tv}'$ is valid', 
and 
\item
$\f{tv}\tiffcc \f{tv}'$ as `$\f{tv}$ is true if and only if $\f{tv}'$ is true', 
and 
\item
$\f{tv}\equiv \f{tv}'$ as `$\f{tv}$ is literally equal to $\f{tv}'$'.
\end{itemize*}
Of these three notions of equivalence, $\equiv$ will be most useful to us later.
We will consider it in our predicate language in Subsection~\ref{subsect.equivalence.of.predicates}, and (for instance) it will help us to express notions of extensional equivalence in Figure~\ref{fig.easy.equivalences}.

We will call $\equiv$ \deffont{literal equivalence}, and $\equiv$ appears in the sugared syntax of Figure~\ref{fig.3.derived}.
As per the discussion in Remark~\ref{rmrk.rich.implication}, there would be no barrier to including $\tlatticeiff$ or $\tiffcc$ but we do not need to for now.
\end{rmrk}








\section{\QLogic}
\label{sect.quorum.logic}

\subsection{A preliminary: semitopologies}

\begin{rmrk}
A semitopology is like a topology, but generalised to not require that intersections of open sets necessarily be opens.

The intuition of an open set is an \emph{actionable coalition}; a set of participants with the voting power (or other resources) to make updates to their local state.
The reader might be familiar with the notion of \emph{quorum} from literature like~\cite{DBLP:journals/tocs/Lamport98,DBLP:journals/dc/MalkhiR98,naor:loacaq}, or the notion of \emph{winning coalition} in social choice theory (the theory of voting)~\cite[Item~5, page~40]{riker:thepc}.
Semitopologies are similar to this, except that open sets are always closed under arbitrary (possibly empty) unions, whereas 
\begin{itemize*}
\item
quorums may be closed under arbitrary nonempty unions (in which case the correspondence is precise, except that $\varnothing$ is an open set but is not a quorum), or 
\item
quorums are not closed under arbitrary unions (in which case a quorum corresponds to a minimal nonempty open set).
\end{itemize*}
Our presentation of semitopologies will be quite brief --- we just define the parts we need to give semantics to the modalities in \QLogic --- and the reader who wants to know more is referred to previous work~\cite{gabbay:semdca}. 
\end{rmrk} 

\begin{defn}
\label{defn.semitopology}
A \deffont{semitopology} is a pair $(\Pnt,\opens)$ where:
\begin{enumerate*}
\item
$\Pnt$ is a set of \deffont{points}, which following terminology from distributed systems one might also call \deffont{participants}, \deffont{processes}, \deffont{parties}, or \deffont{nodes}.
\item
$\opens\subseteq\powerset(\Pnt)$ is a set of \deffont{open sets} such that:
\begin{enumerate*}
\item
If $\mathcal O\subseteq\powerset(\Pnt)$ then $\bigcup\mathcal O\in\Pnt$.
In words: $\opens$ is closed under arbitrary (possibly empty) unions.
\item
$\Pnt\in\opens$.
\end{enumerate*} 
\end{enumerate*} 
We may write $\opensne$ for the set of nonempty open sets, so
$$
\opensne = \{O\in\opens \mid O\neq\varnothing\} .
$$
\end{defn}


\begin{nttn}
\label{nttn.empty.semitopology}
Suppose $(\Pnt,\opens)$ is a semitopology.
\begin{enumerate*}
\item
We call $(\Pnt,\opens)$ the \deffont{empty semitopology} when its set of points is empty --- i.e. when $\Pnt=\varnothing$ --- so that necessarily $\opens=\{\varnothing\}$.
\item
We call $(\Pnt,\opens)$ a \deffont{nonempty semitopology} when its set of points is nonempty --- i.e. when $\Pnt\neq\varnothing$.
\end{enumerate*}
\end{nttn}


\begin{defn}[All-but-$f$-many semitopology]
\label{defn.allbut}
\leavevmode
\begin{enumerate*}
\item
Suppose $X$ is a finite set and $N\in\Ngeqz$.
Define $\powerset_{\geq N}(X)$ to be the set of subsets of $X$ that contain at least $N$ elements.
In symbols:
$$
\powerset_{\geq N}(X) = \{O\subseteq X \mid \mycard O\geq N\} .
$$
\item\label{item.allbut}
Suppose $N,f\in\Ngeqz$ and $N\geq f$.
Define a semitopology $\AllBut{N}{f}$ by
$$
\AllBut{N}{f} = \bigl(\Nset{N}, \{\varnothing\}\cup\powerset_{\geq N\minus f}(\Nset{N}\bigr)\bigr).
$$
Thus: points are numbers $0\leq n<N$; and open sets are sets of points of cardinality at least $N\minus f$.
\end{enumerate*}
\end{defn}

\begin{lemm}
Suppose $N,f\in\Ngeqz$ and $N\geq f$.
Then:
\begin{enumerate*}
\item
$\AllBut{N}{f}$ is indeed a semitopology.
\item
$\AllBut{N}{f}$ is a topology precisely when $f=N$ or $f=0$.
\end{enumerate*}
\end{lemm}
\begin{proof}
We consider each part in turn:
\begin{enumerate}
\item
It is routine to check that the conditions on open sets in Definition~\ref{defn.semitopology} are satisfied.
The key observation is just the elementary fact that if $\mycard O\geq N\minus f$ and $\mycard O'\geq N\minus f$ then $\mycard(O\cup O')\geq N\minus f$.
\item
For $\AllBut{N}{f}$ to be a topology, open sets need to be closed under intersections.
It is routine to check that this can happen only when $\AllBut{N}{f}$ is the full powerset (so $f=N$) or when open sets are just $\varnothing$ and $\Pnt$ (so $f=0$).
\qedhere\end{enumerate}
\end{proof}

\begin{rmrk}
The original Paxos paper by Lamport~\cite{DBLP:journals/tocs/Lamport98} used semitopologies of the form $\AllBut{N}{f}$ where $N=2n$ and $f=n\minus 1$, for $n\geq 1$ (i.e. open sets are strict majorities of points).
This invites a question: why bother with semitopologies at all?
\begin{enumerate*}
\item
Because we can. 
More specifically, our axiomatisation is parametric over the underlying semitopology.
As good mathematicians, if we can generalise at no cost, then we should.\footnote{\dots and this applies not just in mathematics papers.  Software engineers follow a similar principle and have many names for it: e.g. encapsulation, APIs, modules, classes, typeclasses, and so on.} 
\item
Semitopology not only generalises; it simplifies.
The only properties of $\AllBut{N}{f}$ that we will need, are precisely that it is a semitopology.
If all we need about $\AllBut{N}{f}$ is its semitopological structure, then we should work with that structure because that is what actually matters.
\item
Viewing $\AllBut{N}{f}$ as a semitopology, rather than just as a collection of quorums, adds real meaning and understanding.
In particular, the notions of `dense subset' and `nonempty open interior' will be relevant to understanding the axioms, and these arise naturally in the predicates for declarative axiomatisations of distributed algorithms.
These are clearly native topological notions, so the mathematics is speaking to us here, and what it has to say is inherently topological. 
\end{enumerate*} 
\end{rmrk}


\subsection{Syntax of the logic}

We now develop our three-valued modal fixedpoint logic for reasoning about distributed systems.
We so this in the following order: \emph{signatures} $\Sigma$, \emph{models} $\mathcal M$, \emph{contexts} $\acontext$, and finally \emph{syntax}. 
We develop denotation for this syntax in Subsection~\ref{subsect.denotation}.
For quick reference, the reader can just look at the core syntax and denotation in Figures~\ref{fig.predicate.syntax} and~\ref{fig.3.phi.f}.

 


\begin{defn}[Signatures]
\label{defn.signature}
\leavevmode
\begin{enumerate*}
\item\label{item.signature.varsymb}
Fix $\tf{VarSymb}$ a countably infinite set of \deffont{variable symbols} $\thea,\theb,\thec\in\tf{VarSymb}$.
This set will never change and we may assume it henceforth.
\item
Fix a \deffont{fixedpoint variable symbol} $\tf X$ (we will only need one of these, though there is inherent obstacle to having more).
\item\label{item.predicate.syntax.signature}
A \deffont{signature} $\Sigma$ is a pair 
$$
\Sigma=(\PredSymb,\arity)
$$ 
where:
\begin{enumerate*}
\item\label{item.predsymb}
$\PredSymb$ is a (possibly finite) set of \deffont{predicate symbols} $\tf P,\tf P'\in\PredSymb$.
\item
$\arity$ is an \deffont{arity function} $\arity(\tf P)\in\Ngeqz$ mapping each predicate symbol in $\PredSymb$ to a nonnegative natural number (which is intuitively its number of arguments).
\end{enumerate*}
\item
We may write a signature $\Sigma$ in list form as 
$$
\Sigma=[\tf P,\tf P':i,\ \tf Q:j,\ \dots]
$$ 
meaning that $\tf P,\tf P',\tf Q\in\PredSymb$ and $\arity(\tf P)=\arity(\tf P')=i$ and $\arity(\tf Q)=j$ and so on.
\item\label{item.predicate.syntax.empty.signature}
If the signature in part~\ref{item.predsymb} above contains no predicate symbols then we call it \deffont{empty}; otherwise we call it non-empty.
There is only one empty signature, which we will call \deffont{the empty signature}.\footnote{It is not particularly useful.} 
\end{enumerate*}
\end{defn}


\begin{defn}[Models]
\label{defn.model}
A \deffont{model} $\mathcal M$ is a tuple $\mathcal M=(\Pnt,\opens,\Val)$ where:
\begin{enumerate*}
\item
$(\Pnt,\opens)$ is a semitopology (Definition~\ref{defn.semitopology}), which we may call the \deffont{underlying semitopology} of the model.

We may call $p\in\Pnt$ a \deffont{point} or a \deffont{participant}; also \deffont{party}, \deffont{process}, or \deffont{node} would be consistent with the literature.
For our purposes, these are all synonyms: the terminology \emph{point} is from topology; \emph{participant}, \emph{party}, and \emph{process} are from distributed systems; and \emph{node} is from graph theory (thinking of a network as a graph of communicating processes).
\item\label{item.model.val}
$\Val$ is a nonempty set of \deffont{values}.\footnote{We assume values are nonempty to avoid the degenerate case where a universal quantification over values holds just because there are no values on which to test it.  See in particular the argument in Lemma~\ref{lemm.intertwined.char}(\ref{item.intertwined.char.reverse}).}

For our purposes we will not care what values are, so long as some exist, so we do not put any type system on values or endow them with structure; but we could easily do so if required. 
\end{enumerate*} 
\end{defn}


\begin{defn}[Contexts]
\label{defn.contexts}
Suppose $\Sigma=(\PredSymb,\arity)$ is a signature, and $\mathcal M=(\Pnt,\opens,\Val)$ is a model. 
\begin{enumerate*}
\item\label{item.context.stage}
A \deffont{stage} is a nonnegative number $n\in\Time$.
\item\label{item.contexts.valuation}
A \deffont{valuation} $\avaluation$ inputs a predicate symbol $\tf P\in\PredSymb$ and a stage $n\in\Time$ and a point $p\in\Pnt$ and returns a function $\avaluation(\tf P)_{n,p}:\Val^{\arity(\tf P)}\to\THREE$.

We write the set of valuations $\Valuation(\Sigma,\Val)$ or $\Valuation(\Sigma,\mathcal M)$.
\item\label{item.contexts}
A \deffont{context} $\acontext$ is a tuple $(n,p,O,\avaluation)$ where $n\in\Time$ is a stage, $p\in\Pnt$ is a point, $O\in\opensne$ is a nonempty open set, and $\avaluation\in\tf{valuation}(\Sigma,\Val)$ is a valuation. 

We write the set of contexts $\tf{Ctx}(\Sigma,\mathcal M)$; if $\Sigma$ and $\mathcal M$ are understood, we may write this just $\tf{Ctx}$.
\item
A \deffont{context valuation} is a function $\varkappa\in \THREE^{\tf{Ctx}}$.
\end{enumerate*}
\end{defn}



\begin{figure} %
$$
\begin{array}{r@{\ }l}
t ::=& 
(a\in\tf{VarSymb}) \mid (v\in\Val) %
\\
\phi ::=&  
\tvB 
\mid t\teq t \mid (\tf P(\overbrace{t,\dots,t}^{\arity(\tf P)}) : \tf P\in\PredSymb)
\\
\mid &\tneg\phi \mid \modT\phi \mid \someone\phi \mid \Quorum\phi \mid \texi a.\phi
\\
\mid & \phi\tor\phi'  %
\\
\mid & \yesterday \phi \mid \tomorrow \phi 
\\
\mid & \tf X \mid \mu\tf X.\phi \mid (\varkappa \in \THREE^{\tf{Ctx}})
\end{array}
$$ 
\caption{Term and predicate syntax of \QLogic (Definition~\ref{defn.predicate.syntax})}
\label{fig.predicate.syntax}
\end{figure}


\begin{defn}[Syntax]
\label{defn.predicate.syntax}
Suppose $\Sigma$ is a signature (Definition~\ref{defn.signature}) and $\mathcal M$ is a model (Definition~\ref{defn.model}).
\begin{enumerate*}
\item 
Define the syntax of \deffont{\QLogic} over $\Sigma$ and $\mathcal M$ as in Figure~\ref{fig.predicate.syntax}.

In this syntax, $\texi\thea.\phi$ binds $\thea$ in $\phi$, and $\mu\tf X.\phi$ binds $\tf X$ in $\phi$.
We call a term or predicate with no free variable symbols or fixedpoint variable symbols \deffont{closed}.
\item
One special case will be of particular interest: we will write 
\begin{quote}
`suppose $\texi a.\phi$ is a closed predicate',
\end{quote}
and often we do not really care about $\texi a.\phi$ \emph{per se}, except that when this is a closed predicate it means that $a\in\tf{VarSymb}$ and $\phi$ is a predicate whose set of free variables is at most $\{a\}$.
See for example Remark~\ref{rmrk.unpack.texiunique} and Lemma~\ref{lemm.unique.affine.existence}.
\item
We define substitution of variable symbols / fixedpoint variable symbols for values / context valuations as usual, and write these $\phi[a\ssm v]$ and $\phi[\tf X\ssm \varkappa]$ respectively. %
\end{enumerate*}
\end{defn}


\begin{rmrk}
We call our logic \emph{\QLogic} to reflect that it is a logic for reasoning about how coalitions of participants can work together to carry out algorithms to achieve their goals.
In this document we will apply it to Paxos, but we see this logic as having independent interest and a broader relevance for specifying and reasoning about the kinds of algorithms that get written for distributed systems.
Their hallmark is a requirement for distributed collaborative action --- action by quorums, or actionable coalitions, or winning coalitions; these things have many names --- to progress an algorithm to its next step. 
\end{rmrk}

\begin{xmpl}
We give some example predicates and suggest intuitive meanings for them.
We do not yet have the denotation (Figure~\ref{fig.3.phi.f}) nor the rich derived syntax (Figure~\ref{fig.3.derived}) that we build later, so these examples are simple and informal, but we hope they will be helpful.

Assume our signature has a unary predicate symbol $\tf{output}$ (representing the output of some process), and assume $v,v'\in\Val$ are values.
\begin{enumerate}
\item
$\tf{output}(v)$ means `where we are, today, if the process succeeds, then it outputs $v$' and $\tneg\tf{output}(v')$ means `where we are, if the process succeeds, then it does not output $v'$'.
\item
$\modT\tf{output}(v)$ means `where we are, today, the process succeeds and outputs $v$'. 
\item
$\someone\tf{output}(v)$ means `there is a place where, today, if the process succeeds, then it outputs $v$'.
\item
$\tomorrow\modT\tf{output}(v)$ means `where we are, tomorrow, the process succeeds and outputs $v$'.
\item
$\yesterday\modT\tf{output}(v)$ means `where we are, yesterday, the process succeeded and output $v$'.
This assertion is subtly stronger than it might seem; our notion of time starts at day~0, at which by convention all assertions about yesterday are false.
See Remark~\ref{rmrk.further.comments.denotation}(\ref{item.interpret.yesterday}).
\item
$\tvB$ means `where we are, today, the process (which we have not named with a predicate symbol) crashes'.
\end{enumerate}
More sophisticated examples will appear later; notably in the derived syntax in Figure~\ref{fig.3.derived}, in Remark~\ref{rmrk.comments.on.semantics}, and of course when we apply our logic to axiomatise distributed algorithms.
\end{xmpl}


\begin{rmrk}[Three pedant points]
\label{rmrk.pedant.point}
The thing about being pedantic is that it does not matter, until it does.
We now discuss three pedant points, which may not matter \dots until perhaps a moment arrives then they do:
\begin{enumerate}
\item
We may treat $\Sigma$ and $\mathcal M$ as fixed parameters, and elide them.  
In practice this means that we may write `suppose $\phi$ is a closed predicate' instead of `suppose $\phi$ is a closed predicate in the syntax of \QLogic over $\Sigma$ and $\mathcal M$'.
But note next point.
\item\label{item.model.matters}
Our syntax from Definition~\ref{defn.predicate.syntax} and Figure~\ref{fig.predicate.syntax} contains denotational elements; $v\in\Val$ and $\varkappa\in\THREE^{\tf{Ctx}(\Sigma,\mathcal M)}$.

So when we write `suppose $\phi$ is a closed predicate' then this statement is perfectly correct and well-defined but it \emph{really does matter} what the signature and model are (even if the \emph{precise} choice of model will may not matter to the theorems we prove), because closed predicates $\phi$ can include elements from the denotation such as $\varkappa\in\THREE^{\tf{Ctx}(\Sigma,\mathcal M)}$ (this function-space is determined by both the signature $\Sigma$ and the model $\mathcal M$) and $v\teq v'$ for some $v,v'\in\Val$ ($\Val$ is determined by the choice of model $\mathcal M$).

Adding denotational elements to syntax avoids us having to put this information in the context (which is already a $4$-tuple, as per Definition~\ref{defn.contexts}(\ref{item.contexts})) and simplifies the presentation.
This is a technical device with a long pedigree.\footnote{Adding elements from a model to syntax, and conversely building a model out of syntax, is a standard technical device in logic, with a pedigree that goes back (at least) to the proofs of the upward and downward L\"owenheim-Skolem theorems from the 1920s.  A clear survey of these methods is in Chapter~2 (``Models constructed from constants'') of~\cite{chang:modt}.}

Note that (by design!) the choice of model makes a difference to validity.
For example: in Definition~\ref{defn.phi.equivalent} we define equivalence of predicates (and we explicitly mention $\Sigma$ and $\mathcal M$);
if $\mathcal M$ has only a single value $v$, then $\tall a,b.a\teq b$ and $\tall a.a\teq v$ are equivalent to $\tvT$; and if $\mathcal M$ has more than one value then $\tall a,b.a\teq b$ is equivalent to $\tvF$ --- $\tall a.\phi$ is shorthand for $\tneg\texi a.\tneg\phi$; see Figure~\ref{fig.3.derived}.
In contrast: in Lemma~\ref{lemm.equiv} we elide $\Sigma$ and $\mathcal M$ for brevity, and just write `suppose $\phi$ and $\phi'$ are closed predicates and $\acontext$ is a context', without specifying over what signature or model, because this is clear in context.

The reader can mentally add $\Sigma$ and $\mathcal M$ everywhere if they like.
We will always be precise and explicit where it counts.
\item\label{item.pedant.texi.v}
We may quantify over $v\in\Val$ in our formal syntax of predicates, writing things like $\texi v.\texi v'.v\teq v'$ as predicates.
Technically this is wrong and is not syntax! 
Figure~\ref{fig.predicate.syntax} permits e.g. $\texi a.\texi a'.a\teq a'$ where $a,a'\in\tf{VarSymb}$ (Definition~\ref{defn.signature}(\ref{item.signature.varsymb})); it does not permit $\texi v$ in the formal syntax of predicates where $v\in\Val$, so writing e.g. $\texi v.\texi v'.v\teq v'$ as if it were formal syntax of \QLogic is an abuse of notation.

However, $a$ and $a'$ get instantiated to values $v$ and $v'$ in our denotation (see the line for $\texi$ in Figure~\ref{fig.3.phi.f}).
It seems not worth the disruption of maintaining a distinction between $\texi a.\phi$ (which is syntax, assuming $\phi$ is) and $\texi v.\phi[a\ssm v]$ (which looks like syntax but technically is not).
So in practice this seems to us one example of an abuse of formal notation that is positively useful and clarifies meaning.

If this bothers the reader, then every time they see $\texi v.\phi$ (and $\tall v.\phi$, and similar) in our text, they can either declare that $v$ here is \emph{actually} a variable symbol in $\tf{VarSymb}$ which we happen to have called $v$ (and not a value in $\Val$ after all), or they can pick a fresh variable symbol $a$ and replace every unbound $v$ in $\phi$ with $a$.
Either way they will get well-formed syntax and everything will be as it should be.

Notation exists to serve clarity, and not the other way around, so we will simply ignore this henceforth and carry on writing $\texi v.\phi$, just whenever we feel like it, and this will denote just what it says: `there exists a $v$ such that $\phi$'.
Meaning should always be clear. 

Note that an unbound value $v$ in a predicate is legal syntax as per Figure~\ref{fig.predicate.syntax}; 
e.g. in `$\texi a.a\teq v$', $v$ is a value, and this is fine; the (constructive and useful) abuse of notation comes when we write $\texi v$. 
\end{enumerate}
\end{rmrk}

\subsection{Denotation (semantics) of the syntax}
\label{subsect.denotation}


\begin{figure} %
$$
\begin{array}{r@{\ }l@{\quad}l}
\modellabel{\tvB}{n,p,O,\avaluation} =&\tvB
\\
\modellabel{v\teq v'}{n,p,O,\avaluation} =&
\begin{cases}
\tvT & v=v'
\\
\tvF & v\neq v'
\end{cases}
\\
\modellabel{\tf P(v_1,\dots,v_{\arity(\tf P))}}{n,p,O,\avaluation} =&\avaluation(\tf P)_{n,p}(v_1,\dots,v_{\arity(\tf P)}) %
\\[2ex]
\modellabel{\phi\tor\phi'}{n,p,O,\avaluation} =& \modellabel{\phi}{n,p,O,\avaluation}\tor\modellabel{\phi'}{n,p,O,\avaluation}
\\
\modellabel{\tneg\phi}{n,p,O,\avaluation} =& \tneg\,\modellabel{\phi}{n,p,O,\avaluation}  
\\
\modellabel{\modT\phi}{n,p,O,\avaluation} =& \modT\modellabel{\phi}{n,p,O,\avaluation}  
\\
\modellabel{\texi\thea.\phi}{n,p,O,\avaluation} =& \bigvee_{v{\in}\Val} \modellabel{\phi[\thea\ssm v]}{n,p,O,\avaluation}
\\
\modellabel{\someone\phi}{n,p,O,\avaluation} =& \bigvee_{p'\in O} \modellabel{\phi}{n,p',O,\avaluation} 
\\
\modellabel{\Quorum\phi}{n,p,O,\avaluation} =& \bigvee_{O'\in\opensne} \modellabel{\phi}{n,p,O',\avaluation} 
\\[2ex]
\modellabel{\yesterday\phi}{n,p,O,\avaluation} =&
\bigvee_{\f{max}(0,n\minus 1)\leq n'<n}\modellabel{\phi}{n',p,O,\avaluation}
\\
\modellabel{\tomorrow\phi}{n,p,O,\avaluation} =&
\modellabel{\phi}{n\plus 1,p,O,\avaluation}
\\[2ex]
\modellabel{\mu\tf X.\phi}{n,p,O,\avaluation}=&
\bigl(\bigwedge\{ \varkappa\in\THREE^{\tf{Ctx}}\mid \model{\phi[\tf X\ssm\varkappa]}\leq\varkappa\}\bigr)(n,p,O,\avaluation)
\figskip
\modellabel{\varkappa}{n,p,O,\avaluation}=&
\varkappa(n,p,O,\avaluation)
\end{array}
$$ 
\caption{Denotational semantics of \QLogic (Definition~\ref{defn.denotation})}
\label{fig.3.phi.f}
\end{figure}


\begin{figure} %
$$
\begin{array}{r@{\ }l@{\ \ }r@{\ }l}
\tvF =& \modT\tvB
&
\modellabel{\tvF}{\acontext} =&\tvF
\figskip
\tvT =& \tneg\tvF
&
\modellabel{\tvT}{\acontext} =&\tvT
\\[2ex]
\modF\phi =& \modT\tneg\phi
&
\modellabel{\modF\phi}{\acontext} =&\modF\modellabel{\phi}{\acontext}
\figskip
\modB\phi =& \tneg\modT\phi\tand\tneg\modF\phi
&
\modellabel{\modB\phi}{\acontext} =&\modB\modellabel{\phi}{\acontext}
\figskip
\modTB\phi =& \tneg\modT\tneg\phi
&
\modellabel{\modTB\phi}{\acontext} =&\modTB\modellabel{\phi}{\acontext}
\figskip
\modTF\phi =& \modT\phi\tor\modF\phi
&
\modellabel{\modTF\phi}{\acontext} =&\modTF\modellabel{\phi}{\acontext}
\figskip
\modFB\phi =& \modF\phi\tor\modB\phi
&
\modellabel{\modFB\phi}{\acontext} =&\modFB\modellabel{\phi}{\acontext}
\\[2ex]
\phi\tand\phi' =& \tneg(\tneg\phi\tor \tneg\phi')
&
\modellabel{\phi\tand\phi'}{\acontext} =&\modellabel{\phi}{\acontext}\tand\modellabel{\phi'}{\acontext}
\figskip
\phi\tnotor\phi' =& \tneg\phi\tor\phi'
&
\modellabel{\phi\tnotor\phi'}{\acontext} =&\modellabel{\phi}{\acontext}\tnotor\modellabel{\phi'}{\acontext}
\figskip
\phi\timpc\phi' =& \phi\tnotor\modT\phi'
&
\modellabel{\phi\timpc\phi'}{\acontext} =&\modellabel{\phi}{\acontext}\timpc\modellabel{\phi'}{\acontext}
\figskip
\phi\equiv\phi' =& (\modT\phi\tand\modT\phi')\tor(\modB\phi\tand\modB\phi') \tor(\modF\phi\tand\modF\phi') \hspace{-5em}
\\
&&
\modellabel{\phi\equiv\phi'}{\acontext} =&\modellabel{\phi}{\acontext}\equiv\modellabel{\phi'}{\acontext}
\\[2ex]
\tall a.\phi =& \tneg\texi a.\tneg\phi
&
\modellabel{\tall a.\phi}{\acontext} =&\bigwedge_{v{\in}\Val}\modellabel{\phi[a\ssm v]}{\acontext}
\figskip
\everyone \phi =& \tneg\someone\tneg\phi
&
\modellabel{\everyone\phi}{n,p,O,\avaluation} =&\bigwedge_{p'{\in}\Pnt}\modellabel{\phi}{n,p',O,\avaluation}
\figskip
\Coquorum\phi =& \tneg\Quorum\tneg\phi
&
\modellabel{\Coquorum\phi}{n,p,O,\avaluation} =&\bigwedge_{O'{\in}\opensne}\modellabel{\phi}{n,p,O',\avaluation}
\\[2ex]
\QuorumBox\phi =& \Quorum\everyone\phi
&
\modellabel{\QuorumBox\phi}{n,p,O,\avaluation} =&\bigvee_{O'{\in}\opensne}\bigwedge_{p'\in O'}\modellabel{\phi}{n,p',O',\avaluation}
\figskip
\CoquorumDiamond\phi =& \Coquorum\someone\phi
&
\modellabel{\CoquorumDiamond\phi}{n,p,O,\avaluation} =&\bigwedge_{O'{\in}\opensne}\bigvee_{p'\in O'}\modellabel{\phi}{n,p',O',\avaluation}
\figskip
\someoneAll\phi =& \Quorum\someone\phi
&
\modellabel{\someoneAll\phi}{n,p,O,\avaluation} =&\bigvee_{O'{\in}\opensne}\bigvee_{p'\in O'}\modellabel{\phi}{n,p',O',\avaluation}
\figskip
\everyoneAll\phi =& \Coquorum\everyone\phi
&
\modellabel{\everyoneAll\phi}{n,p,O,\avaluation} =&\bigwedge_{O'{\in}\opensne}\bigwedge_{p'\in O'}\modellabel{\phi}{n,p',O',\avaluation}
\\[2ex]
\forever\phi =& \mu\tf X.\tomorrow(\phi\tand \tf X)
&
\modellabel{\forever\phi}{n,p,O,\avaluation} =&\bigwedge_{n'> n}\modellabel{\phi}{n',p,O,\avaluation}
\figskip
\sometime\phi =& \mu\tf X.\tomorrow(\phi\tor \tf X)
&
\modellabel{\sometime\phi}{n,p,O,\avaluation} =&\bigvee_{n'> n}\modellabel{\phi}{n',p,O,\avaluation}
\figskip
\infinitely\phi =& \forever\sometime\phi
&
\modellabel{\infinitely\phi}{n,p,O,\avaluation} =&\bigwedge_{n'> n}\bigvee_{n''>n'}\modellabel{\phi}{n'',p,O,\avaluation}
\figskip
\final\phi =& \sometime\forever\phi
&
\modellabel{\final\phi}{n,p,O,\avaluation} =&\bigvee_{n'> n}\bigwedge_{n''>n'}\modellabel{\phi}{n'',p,O,\avaluation}
\figskip
\recent\phi =& \mu\tf X.\yesterday(\phi\tor\tf X)
&
\modellabel{\recent\phi}{n,p,O,\avaluation} =&\bigvee_{0\leq n'<n}\modellabel{\phi}{n',p,O,\avaluation}
\figskip
\urecent\phi =& \mu\tf X.(\phi\tor\yesterday\tf X)
&
\modellabel{\urecent\phi}{n,p,O,\avaluation} =&\bigvee_{0\leq n'\leq n}\modellabel{\phi}{n',p,O,\avaluation}
\end{array}
$$
$$
\begin{array}{r@{\ }l}
\mrup{a}{\phi}{v} =&
\mu\tf X.(\yesterday(\modT\phi[a\ssm v] \tor ((\tneg\modT\texi a.\phi) \tand \tf X)))
\figskip
\modellabel{\mrup{a}{\phi}{v}}{n,p,O,\avaluation} =&
\begin{cases}
\tvF 
&\Forall{0\leq n'<n} \modellabel{\texi a.\phi}{n',p,O,\avaluation}\neq\tvT
\\
\modT\modellabel{\phi[\thea\ssm v]}{\f{max},p,O,\avaluation}
&\Exists{0\leq n'<n} \modellabel{\texi a.\phi}{n',p,O,\avaluation}=\tvT
\end{cases}
\\
&\qquad\text{where}\quad
\f{max}=\bigvee\{0\leq n'<n \mid \modellabel{\texi a.\phi}{n',p,O,\avaluation}=\tvT\}
\end{array}
$$ 
\caption{Derived syntactic sugar expressions (Definition~\ref{defn.3.derived})}
\label{fig.3.derived}
\end{figure}



Recall from Definition~\ref{defn.signature}(\ref{item.predicate.syntax.signature}) the notion of a \emph{signature} and from Definition~\ref{defn.model} the notion of a \emph{model}:
\begin{defn}[Denotation]
\label{defn.denotation}
Suppose $\Sigma=(\PredSymb,\arity)$ is a signature, and $\mathcal M$ is a model. 
Suppose $\phi$ is a closed predicate (Definition~\ref{defn.predicate.syntax}(1)).
Define the \deffont{denotation} of $\phi$ 
$$
\model{\phi}\in\THREE^{\tf{Ctx}(\Sigma,\mathcal M)}
\quad\text{such that}\quad 
\acontext \longmapsto \modellabel{\phi}{\acontext}
$$
by the rules in Figure~\ref{fig.3.phi.f}.

We will usually treat $\Sigma$ and $\mathcal M$ as fixed parameters, so we may write $\tf{Ctx}(\Sigma,\mathcal M)$ just as $\tf{Ctx}$.
\end{defn}




\begin{defn}[Derived expressions]
\label{defn.3.derived}
The syntax of \QLogic in Figure~\ref{fig.predicate.syntax} and Definition~\ref{defn.predicate.syntax} is very expressive.
In the left-hand column of Figure~\ref{fig.3.derived} we define some derived expressions --- i.e. \emph{syntactic sugar}.
For clarity and convenience we unpack their denotation in the right-hand column.

In principle we could fully unpack our sugar down to the primitive syntax from Figure~\ref{fig.3.phi.f}, but in practice we work directly with Figure~\ref{fig.3.derived} because this is more convenient.
We will discuss syntax further below, notably in Remarks~\ref{rmrk.initial.comments.denotation}, \ref{rmrk.further.comments.denotation}, and~\ref{rmrk.comments.on.semantics}.
\end{defn}


\begin{rmrk}
Usually when we talk about denotation as defined in Definition~\ref{defn.denotation} above, the signature $\Sigma$ and model $\mathcal M$ are either understood or do not matter, so we may not mention them explicitly.
Meaning will always be clear.
\end{rmrk}

\begin{rmrk}[Initial comments on the denotation]
\label{rmrk.initial.comments.denotation}
We make some initial comments on Figure~\ref{fig.3.phi.f}:
\begin{enumerate*}
\item\label{item.first.mu}
$\mu\tf X.\phi$ is a \emph{fixedpoint operator}.
This is a standard method to include the power of recursion / iteration in a logic; an instance of $\tf X$ inside $\mu\tf X.\phi$ can be thought of as a recursive call back to $\phi$.
$\mu\tf X.\phi$ is what would be written $\mathrm{LFP}_{\tf X}\phi$ in~\cite[Chapter~4]{immerman:desc} (a Wikipedia article is also accessible and clear~\cite{wiki:Fixedpoint_logic}).
We will discuss this in Subsection~\ref{subsection.muX}.
\item
In the clause $\model{\tvB}=\tvB$ in Figure~\ref{fig.3.phi.f}, the $\tvB$ on the left is the syntax symbol from Figure~\ref{fig.predicate.syntax}, and the $\tvB$ on the right is $\tvB$-the-truth-value from Definition~\ref{defn.THREE}.
\item
In the clause for $\model{v\teq v'}$ in Figure~\ref{fig.3.phi.f}, the $\teq$ on the left is the syntax symbol from Figure~\ref{fig.predicate.syntax}, and the $\tvT$, $\tvF$, and $=$ on the right are truth-values and (actual) equality.
\item
Similarly for the clauses for $\model{\tneg\phi}$ and $\model{\phi\tor\phi'}$ and $\model{\modT\phi}$; the symbols $\tneg$, $\tor$, and $\modT$ on the left are syntax, and the symbols $\tneg$, $\tor$, and $\modT$ on the right refer to functions on truth-values as per Figure~\ref{fig.3}.
\end{enumerate*} 
\end{rmrk}


\begin{rmrk}[Further comments] 
\label{rmrk.further.comments.denotation}
\leavevmode
\begin{enumerate}
\item\label{item.interpret.yesterday}
The clause 
$\modellabel{\yesterday\phi}{n,p,O,\avaluation} = \bigvee_{\f{max}(0,n\minus 1)\leq n'<n}\modellabel{\phi}{n',p,O,\avaluation}$ is just a one-line encoding of
$$
\modellabel{\yesterday\phi}{n,p,O,\avaluation}=
\begin{cases}
\tvF & n=0
\\
\modellabel{\phi}{n\minus 1,p,O,\avaluation} & n>0 . 
\end{cases}
$$
\item
We see from the right-hand column in Figure~\ref{fig.3.derived} that $\recent\phi$ returns the greatest truth-value that $\phi$ has taken in the past, and $\urecent\phi$ returns the greatest truth-value that $\phi$ has taken today or in the past.

The notation $\urecent$ (`$\recent$', underlined) is a pun\footnote{A bad one perhaps, and the reader may or may not find this pun-acceptable.} on $<$ and $\leq$ (`$<$', underlined).
\item
Informally, we see that the only naturally \emph{non-monotone} operator in the (primitive or derived) syntax is $\tneg$, which is anti-monotone.
This design feature is what powers Lemma~\ref{lemm.positive.monotonicity}.
\end{enumerate}
\end{rmrk}


\begin{nttn}
\label{nttn.logic.terminology}
We take a moment to recall some standard conventions and terminology:
\begin{enumerate*}
\item
Implications $\tnotor$ and $\timpc$ associate to the right, so that $\phi\timpc\psi\timpc\chi$ means $\phi\timpc(\psi\timpc\chi)$.
\item
Unary connectives (like $\tneg$, $\modT$, and $\someone$) bind more strongly than binary connectives (like $\tand$ and $\tnotor$), so that $\QuorumBox\tneg\phi\tand\psi$ means $(\QuorumBox\tneg\phi)\tand\psi$.
\item
The scope of a quantification $\texi a$ and $\tall a$ extends as far to the right as possible. 
Thus $\tall\thea.\phi\tand\psi$ means $\tall\thea.(\phi\tand\psi)$.
Note however that $\QuorumBox\texi\thea.\phi\tand\psi$ still means $(\QuorumBox(\texi\thea.\phi))\tand\psi$, because the $\texi$ quantifier is trapped in the scope of the $\QuorumBox$. 
\item
All the above said, we will (as is usual) use spacing and some common sense to convey scope.
For example 
$$
\texi\thea.\phi\mathbin{\tnotor}\psi
\quad\text{means}\quad
\texi\thea.(\phi\tnotor\psi), 
$$
and
$$
\texi\thea.\phi \ \tnotor \psi
\quad\text{means}\quad
(\texi\thea.\phi)\tnotor\psi. 
$$
In practice, meaning will always be clear.
\item
\begin{enumerate*}
\item\label{item.nullary.ps}
If $\arity(\tf P)=0$ then we may call $\tf P$ \deffont{nullary} or a \deffont{predicate constant symbol}.\footnote{$\tvT$, $\tvB$, and $\tvF$ are predicate constant symbols, but they are special so we give them a distinct identity in the syntax.}
\item\label{item.unary.ps}
If $\arity(\tf P)=1$ then we may call $\tf P$ a \deffont{unary} predicate symbol.
\item\label{item.atomic.predicate}
We may call a predicate of the form $\tf P(v_1,\dots,v_n)$ or $\tneg\tf P(v_1,\dots,v_n)$ an \deffont{atomic predicate}. 
\end{enumerate*}
\end{enumerate*} 
\end{nttn}



\subsection{The fixedpoint operator $\mu \tf X$}
\label{subsection.muX}

\begin{rmrk}
\label{rmrk.comments.on.semantics}
As per Remark~\ref{rmrk.initial.comments.denotation}(\ref{item.first.mu}), our fixedpoint operator $\mu\tf X.\phi$ is a method to include the power of recursion / iteration in the logic.

A mention of $\tf X$ inside $\phi$ in $\mu\tf X.\phi$ can be thought of as a recursive call back to $\phi$, and we can think of $\mu\tf X.\phi$ as encoding an iterative loop that repeatedly executes a computation encoded in $\phi$, and loops if it encounters $\tf X$ in $\phi$.

With this intuition, we can look at some of the derived syntax in Figure~\ref{fig.3.derived} that uses $\mu$:
\begin{enumerate}
\item
$\recent\phi=\mu\tf X.\yesterday(\phi\tor\tf X)$ can be imagined as the following loop:
\begin{enumerate*}
\item
Set an accumulator to $\tvF$.\footnote{We mean `accumulator' in the sense of an `accumulator variable' inside a machine that (we imagine) is running this pseudocode.  This is not a technical term from logic.}
\item\label{recent.step.0}
If we are at day~0 then terminate and return the value of the accumulator; otherwise go to yesterday.
\item
Compute the truth-value of $\phi$, and $\tor$ this truth-value with the accumulator.
\item
Loop back to step~\ref{recent.step.0}.
\end{enumerate*} 
\item
$\forever\phi=\mu\tf X.\tomorrow(\phi\tand\tf X)$ can be viewed as the following loop:
\begin{enumerate*}
\item
Set an accumulator to $\tvT$.
\item\label{forever.step.0}
Go to tomorrow.
\item
Compute the truth-value of $\phi$, and $\tand$ this truth-value with the accumulator.
\item
Loop back to step~\ref{forever.step.0}.
\item
After infinite time (see Remark~\ref{rmrk.oo} about why that makes sense) return the limit value of the accumulator.
\end{enumerate*} 
\item
$\sometime\phi=\mu\tf X.\tomorrow(\phi\tor\tf X)$ can be viewed as executing the following loop:
\begin{enumerate*}
\item
Set an accumulator to $\tvF$.
\item\label{sometime.step.0}
Go to tomorrow.
\item
Compute the truth-value of $\phi$ and $\tor$ with the accumulator.
\item
Loop back to step~\ref{sometime.step.0}.
\item
After infinite time, return the limit value of the accumulator.
\end{enumerate*}
\item\label{item.mru.program}
The clause 
$\mrup{a}{\phi}{v} = \mu\tf X.(\yesterday(\modT\phi[a\ssm v] \tor ((\tneg\modT\texi a.\phi) \tand \tf X)))$
looks more imposing, but the idea is simply to 
\begin{itemize*}
\item
iterate back in time looking for when $\texi a.\phi$ has truth-value $\tvT$, i.e. when \emph{some} value makes $\phi(a)$ be true, and 
\item
if we find such a time then we check if \emph{at that time} $\phi[a\ssm v]$ has truth-value $\tvT$ and 
\item
only if both steps above succeed, do we return $\tvT$.
\end{itemize*}
Thus, we return $\tvT$ when $v$ is one of the values that \emph{most recently} makes $\phi(a)$ be true.

In detail:
\begin{enumerate*}
\item\label{mu.step.0}
If we are at day~0, terminate with value $\tvF$; otherwise go to yesterday.
\item
If $\phi[a\ssm v]$ has truth-value $\tvT$ then terminate with value $\tvT$; otherwise continue to the next step.
\item
If $\texi a.\phi$ has truth-value $\tvT$ then terminate with value $\tvF$; otherwise go to step~\ref{mu.step.0}.
\end{enumerate*}
\end{enumerate}
\end{rmrk}

\begin{rmrk}[Comment on the omniscient observer]
\label{rmrk.oo}
The programs in Remark~\ref{rmrk.comments.on.semantics} above are executed by an omniscient observer who stands outside of the model's notion of time and space.
This is why (for example) $\forever$ can run an infinite loop --- the omniscient observer can observe it because the observer is looking at the model from outside.\footnote{Let's make this concrete.  Consider the natural numbers $\Ngeqz$.  We can easily observe this and note that $n\plus 1\gneq n$ for every $n\in\Ngeqz$; whereas a local observer trapped on the number line and constrained to perceive this line as the passage of time can make similar observation, but from the point of view of that observer, the calculation involved is neverending.  

Note that even if we could speak with this local observer, it would be no good to appeal to it to `just realise that it is on the number line', because in order to check this to a mathematical certainty the observer would have to \dots traverse all of its local notion of time to \emph{check}, amongst other things, that $n\plus 1\gneq n$.  After all, the observer might be on $\mathbb Z/p$ for some astronomically enormous $p$, i.e. their time might be on a large loop.  

The reader interested in a branch of mathematics explicitly devoted to studying these questions, and who is not afraid by the possibility of ending up seriously doubting the nature of their own reality, can start by reading up on nonstandard models of arithmetic.} 
\end{rmrk}

\begin{rmrk}
The rest of this Subsection will show how our denotation for $\mu$ from Figure~\ref{fig.3.phi.f} does indeed yield a fixedpoint operator.
The culminating result is Proposition~\ref{prop.mu.fixedpoint} (which just says that $\mu$ is indeed a fixedpoint operator) but the key technical lemma is Lemma~\ref{lemm.positive.monotonicity} (monotonicity of positive predicates).

The reader is welcome to read the rest of this Subsection, but this is all just a standard use of the Knaster-Tarski fixedpoint theorem~\cite[Theorem~2.35, page~50]{priestley:intlo}.

In fact, the real design effort has already happened: we designed the derived syntax of Figure~\ref{fig.3.derived} so that it would respect the positivity condition of Definition~\ref{defn.positive}, without having to fully expand it out.
For instance, we defined $\modTB\phi$ to be $\tneg\modT\tneg\phi$ (which clearly puts $\phi$ in a positive position) instead of $\modT\phi\tor\modB\phi$ (which is logically equivalent, but which does not).
\end{rmrk} 



\begin{defn}
\label{defn.three.ctx.lattice}
Recall from Remark~\ref{rmrk.three.elementary}(\ref{item.three.elementary.bounded.lattice}) that $\THREE$ with ordering $\tvF<\tvB<\tvT$ is a finite lattice.
As is standard, the function space $\THREE^{\tf{Ctx}}$ (functions from $\tf{Ctx}$ to $\THREE$) inherits this ordering and lattice structure \deffont{pointwise}, as indicated in Figure~\ref{fig.chi.3.lattice}.

This makes $\THREE^{\tf{Ctx}}$ into a \deffont{bounded complete lattice}, meaning that it has a least and greatest element, and arbitrary joins (least upper bounds) and meets (greatest lower bounds).\footnote{Having a least and greatest element follows from having arbitrary joins and meets; the least element is the join of the empty set, and the greatest element is the meet of the empty set.  So in fact a complete lattice is already a bounded complete lattice.} 
\end{defn}


\begin{figure}
$$
\begin{array}{r@{\ }l}
\tbot =& \lambda\acontext.\tvF
\\
\ttop =& \lambda\acontext.\tvT
\\
\bigvee\mathcal X =& \lambda\acontext.\bigvee\{ \varkappa\,\acontext \mid \varkappa\in\mathcal X\}
\\
\bigwedge\mathcal X =& \lambda\acontext.\bigwedge\{ \varkappa\,\acontext \mid \varkappa\in\mathcal X\}
\\[2ex]
\varkappa\leq\varkappa' \liff& \Forall{\acontext} \varkappa\,\acontext\leq\varkappa'\,\acontext
\end{array}
$$
Above, $\mathcal X\subseteq \THREE^{\tf{Ctx}}$ and $\varkappa,\varkappa'\in\THREE^{\tf{Ctx}}$. 
\caption{Bounded complete lattice structure on the function space $\THREE^{\tf{Ctx}}$}
\label{fig.chi.3.lattice}
\end{figure}

\begin{rmrk}
\label{rmrk.leq.phi}
Our denotation $\model{\text{-}}$ from Definition~\ref{defn.denotation} and Figure~\ref{fig.3.phi.f} maps predicates to functions in $\THREE^{\tf{Ctx}}$.
One particular consequence of Definition~\ref{defn.three.ctx.lattice} is that (denotations of) predicates are naturally ordered, and we can write 
$$
\model{\phi}\leq\model{\psi}
\quad\text{which just means}\quad
\Forall{\acontext\in\tf{Ctx}}\modellabel{\phi}{\acontext}\leq\modellabel{\psi}{\acontext}.
$$
We use this fact for example in Lemma~\ref{lemm.positive.monotonicity}.
We can also use this fact, and the complete lattice structure on $\THREE^{\tf{Ctx}}$ from Definition~\ref{defn.three.ctx.lattice}, to simplify the presentation of the clause for $\mu\tf X.\phi$ in Figure~\ref{fig.3.phi.f} as follows:
$$
\model{\mu\tf X.\phi}=
\bigwedge\{ \varkappa\in\THREE^{\tf{Ctx}}\mid \model{\phi[\tf X\ssm\varkappa]}\leq\varkappa\} .
$$
\end{rmrk}

\begin{defn}
\label{defn.positive}
Recall the predicate syntax from Figure~\ref{fig.predicate.syntax} %
and suppose $\phi$ is a predicate.
Then:
\begin{enumerate*}
\item\label{item.positive.occurrence}
Call a free occurrence of $\tf X$ in $\phi$ 
\begin{itemize*}
\item
\deffont{positive} when the occurrence appears under an even number (which may be $0$) of negations $\tneg$, and 
\item
\deffont{negative} when the occurrence appears under an odd number of negations.
\end{itemize*}
For example: all free occurrences of $\tf X$ are positive in $\tf X$ and $\mu\tf X.\tneg\tf X$ and $\yesterday\tf X$; and the free occurrence of $\tf X$ in $\tneg\tf X$ is negative. 
\item
Call a predicate $\phi$ \deffont{positive} when every free occurence of $\tf X$ in $\phi$ is positive.
\end{enumerate*}
\end{defn}


\begin{lemm}
\label{lemm.positive.monotonicity}
Suppose $\mu\tf X.\phi$ is a closed predicate.
Then:
\begin{enumerate*}
\item
If $\phi$ is positive then $\model{\phi}$ is monotone as a function from $\THREE^{\tf{Ctx}}$ with the partial order $\leq$ from Figure~\ref{fig.chi.3.lattice}, to itself.
In symbols, for every $\varkappa,\varkappa'\in\THREE^{\tf{Ctx}}$:
$$
\varkappa\leq\varkappa'
\quad\text{implies}\quad
\model{\phi[\tf X\ssm\varkappa]}\leq\model{\phi[\tf X\ssm\varkappa']} . 
$$
\item
If $\tneg\phi$ is positive then $\model{\phi}$ is antimonotone as a function from $\THREE^{\tf{Ctx}}$ with the partial order $\leq$ from Figure~\ref{fig.chi.3.lattice}, to itself.
In symbols, for every $\varkappa,\varkappa'\in\THREE^{\tf{Ctx}}$:
$$
\varkappa\leq\varkappa'
\quad\text{implies}\quad
\model{\phi[\tf X\ssm\varkappa']}\leq\model{\phi[\tf X\ssm\varkappa]} . 
$$
\end{enumerate*}
(The meaning of $\model{\text{-}}\leq\model{\text{-}}$ is unpacked in Remark~\ref{rmrk.leq.phi}.)

\end{lemm}
\begin{proof}
Both parts follow by a routine simultaneous induction on the syntax of $\phi$, using the rules in Figure~\ref{fig.3.phi.f}.
\end{proof}

\begin{prop}
\label{prop.mu.fixedpoint}
Suppose $\mu\tf X.\phi$ is a closed predicate and $\phi$ is positive.
Then:
\begin{enumerate*}
\item
$\model{\mu\tf X.\phi}$ is a least fixedpoint of the function $\varphi:\THREE^{\tf{Ctx}}\to\THREE^{\tf{Ctx}}$ given by
$$
\varphi(\varkappa)=\lambda\acontext.\modellabel{\phi[\tf X\ssm\varkappa]}{\acontext}. 
$$
\item
If $\acontext$ is a context then
$$
\modellabel{\mu\tf X.\phi}{\acontext}=\modellabel{\phi[\tf X\ssm \model{\mu\tf X.\phi}]}{\acontext}.
$$
\end{enumerate*}
\end{prop}
\begin{proof}
By the Knaster-Tarski fixedpoint theorem for $\THREE^{\tf{Ctx}}$ as a bounded lattice as per Definition~\ref{defn.three.ctx.lattice}.
A concise but exceptionally clear presentation is in~\cite[Theorem~2.35, page~50]{priestley:intlo} (a Wikipedia article~\cite{wiki:Knaster-Tarski_theorem} is also accessible and clear). 
\end{proof}



\subsection{Validity}



We defined in Definition~\ref{defn.tv.ment} what it means for a truth-value to be \emph{valid}; now we extend this notion to predicates-in-context: 
\begin{defn}[Validity]
\label{defn.validity.judgement}
Suppose that:
\begin{itemize*}
\item
$\Sigma$ is a signature (Definition~\ref{defn.signature}(\ref{item.predicate.syntax.signature})).
We take $\Sigma$ to be fixed for the rest of this Definition, which means that we do not explicitly include it (e.g. as a subscript) in the notation.
\item
$\mathcal M$ is a model (Definition~\ref{defn.model}).
Again, we fix $\mathcal M$ as a parameter for the rest of this definition.
\end{itemize*}
Suppose further that:
\begin{itemize*}
\item
$\acontext=(n,p,O,\avaluation)$ is a context in $\tf{Ctx}(\Sigma,\mathcal M)$, so that by Definition~\ref{defn.contexts}(\ref{item.contexts}) $n\in\Time$ and $p\in\Pnt$ and $O\in\opensne$ and $\avaluation$ is a valuation to $\mathcal M$.
\item
$\phi$ is a closed predicate in $\Sigma$ and $\mathcal M$ (cf. Remark~\ref{rmrk.pedant.point}(\ref{item.model.matters})).
\end{itemize*}
Then we define: %
\begin{enumerate*}
\item\label{item.ctx.ment.phi}
$\acontext\ment\phi$ means $\ment\modellabel{\phi}{\acontext}$, which by Definition~\ref{defn.tv.ment}(\ref{item.tv.ment}) means $\modellabel{\phi}{\acontext}\in\threeValid$, which by Definition~\ref{defn.tv.ment}(\ref{item.tv.ment.correct}) means $\modellabel{\phi}{\acontext}=\tvT\lor\modellabel{\phi}{\acontext}=\tvB$.

Because this is a core definition, we sum this up for the reader's clear reference:
$$
\acontext\ment\phi
\quad\liff\quad
{\ment\modellabel{\phi}{\acontext}}
\quad\liff\quad
\modellabel{\phi}{\acontext}\in\threeValid
\quad\liff\quad
\modellabel{\phi}{\acontext}=\tvT\lor\modellabel{\phi}{\acontext}=\tvB.
$$
We read this as \deffont{$\phi$ is valid in the context $\acontext$.}
\item
$n,p\mentval\phi$ means $\Forall{O\in\opensne}(n,p,O\mentval\phi)$.
\item\label{item.ment.n}
$n\mentval\phi$ means $\Forall{p\in\Pnt}\Forall{O\in\opensne}(n,p,O\mentval\phi)$.
\item\label{item.mentval}
$\mentval\phi$ means $\Forall{n\in\Time}\Forall{p\in\Pnt}\Forall{O\in\opensne}(n,p,O\mentval\phi)$.
\item\label{item.ment.phi}
$\ment\phi$ means $\Forall{\acontext\in\tf{Ctx}(\Sigma,\mathcal M)}(\acontext\ment\phi)$.
\end{enumerate*}
\end{defn}

\begin{figure}
$$
\begin{array}{l@{\quad}l@{\quad}l}
\mentval\Theta & \text{means} & \Forall{\theta\in\Theta}({\mentval\theta})
\figskip
\Theta\ment\phi & \text{means} &
\Forall{\avaluation\in\Valuation(\Sigma,\mathcal M)}({\mentval\Theta})\limp ({\mentval\phi}) 
\end{array}
$$
\caption{Validity of $\phi$ in a theory $\Theta$ (Definition~\ref{defn.theory.axiom})}
\label{fig.validity}
\end{figure}


Later on, when we start to build declarative axiomatisations of consensus algorithms, we will express these as \emph{theories}.
We include the relevant definition here, since it follows on naturally from Definition~\ref{defn.validity.judgement}: 
\begin{defn}[Validity in a theory]
\label{defn.theory.axiom}
Suppose that:
\begin{itemize*}
\item
$\Sigma$ is a signature and $\mathcal M=(\Pnt,\opens,\Val)$ is a model.
\item
$\phi$ is a closed predicate in $\Sigma$ and $\mathcal M$.
\item
$\Theta=\{\theta_1,\dots,\theta_n\}$ is a set of closed predicates.\footnote{It would be unusual for $\Theta$, which is intended to be a collection of axioms, to mention specific elements $v\in\Val$ or $\varkappa\in\THREE^{\tf{Ctx}(\Sigma,\mathcal M)}$ from a model (cf. Remark~\ref{rmrk.pedant.point}), but the maths does not seem to mind, so we do not bother to exclude this possibility.}
\end{itemize*}
Then:
\begin{enumerate*}
\item\label{item.theory}
A \deffont{theory} is a pair $(\Sigma,\Theta)$ of 
\begin{itemize*}
\item
a signature $\Sigma$, and 
\item
a set $\Theta$ of \deffont{axioms}, which are closed predicates over $\Sigma$ and $\mathcal M$.
\end{itemize*}
\item\label{item.axioms.valid}
Given a valuation $\avaluation\in\Valuation(\Sigma,\mathcal M)$ (Definition~\ref{defn.contexts}(\ref{item.contexts.valuation})), we define 
$$
\mentval\Theta
\quad\text{and}\quad
\Theta\ment\phi
$$ 
as per Figure~\ref{fig.validity}. 
When $\Theta\ment\phi$ holds, we say that $\phi$ is \deffont{valid in the theory $\Theta$} (in $\Sigma$ and $\mathcal M$). 
\end{enumerate*}
Examples of theories are in Figures~\ref{fig.ThyPaxOne} and~\ref{fig.ThySPax}.
\end{defn}


\begin{rmrk}
\label{rmrk.thy.quantifier.scope}
Note the scope of the quantifiers in Definition~\ref{defn.theory.axiom}:
\begin{enumerate*}
\item
The assertion $\Theta\ment\phi$ does \emph{not} mean 
$$
\Forall{\avaluation}(\mentval(\theta_1\tand\dots\tand\theta_n) \timp \phi),
$$ 
which is equivalent to 
$$
\Forall{\acontext}(\acontext\ment(\theta_1\tand\dots\tand\theta_n)\limp \acontext\mentval\phi) .
$$ 
Above, the short double arrow `$\timp$' means `if the LHS is valid then so is the RHS'; see Remark~\ref{rmrk.rich.implication}(\ref{item.rich.implication.timp}).
The long double arrow `$\limp$' is just logical implication. 
\item
The assertion $\Theta\ment\phi$ means 
$$
\Forall{\avaluation}(\mentval\theta_1\tand\dots\tand\theta_n) \limp (\mentval\phi) .
$$
In words: if $\Theta$ holds everywhere in the model under the valuation $\avaluation$, then $\phi$ holds everywhere in the model under the valuation $\avaluation$.
\end{enumerate*}
We spell this out further.
The assertion $\Theta\ment\phi$ means:
$$
\Forall{\avaluation}(\Forall{n,p,O}n,p,O\mentval\theta_1\tand\dots\tand\theta_n) \limp (\Forall{n,p,O}n,P,O\mentval\phi),
$$
where above $n\in\Time$, $p\in\Pnt$, and $O\in\opensne$.

This pattern of quantifiers is what makes $\Theta$ a theory, rather than just a set of assumptions: as a theory, $\Theta$ is a restriction on the valuation $\avaluation$.
\end{rmrk}


\begin{rmrk}
One final pedant point: 
in Definition~\ref{defn.theory.axiom}(\ref{item.theory}) we defined a theory to be a pair $(\Sigma,\Theta)$ of a signature and a set of closed predicates over a signature $\Sigma$ \emph{and} over a model $\mathcal M$.
This technical convenience lets us write axioms that directly reference specific values in the model (instead of introducing explicit constant symbols to our signature, and interpreting them as \dots values in the model).

We will exploit to assume a distinguished `undefined' value $\udfn$ in axiom \rulefont{PaxPropose?} in Figure~\ref{fig.ThyPaxOne}, as per Definition~\ref{defn.ThyPaxOne}(\ref{item.ThyPaxOne.axioms}).
\end{rmrk}

\subsection{Comments on notation: eliding parameters, and functional notation}

\begin{nttn}[Eliding parameters]
\label{nttn.validity.irrelevant.context}
It will be quite common that we make validity judgements where it is clear that part of the context is irrelevant.
For example the reader can check the denotational semantics in Figure~\ref{fig.3.phi.f} and the definition of validity in Definition~\ref{defn.validity.judgement}(\ref{item.ctx.ment.phi}), to see that:
\begin{enumerate*}
\item
$\acontext \ment \tvB$ is valid; we do not need to know the values in $\acontext$. 
\item
The validity of $n,p,O\mentval \tf P(v)$ depends on $\avaluation$, $n$, and $p$ --- but not on $O$.
\item
The validity of $n,p,O\mentval\QuorumBox\phi$ depends on $n$ and $\avaluation$ (and of course on $\phi$), but it does not depend on the values of $p$ and $O$.
\end{enumerate*}
In such cases, we may omit parameters in the context if they do not matter --- we may even omit the entire context if appropriate.
Thus, we may write $\ment\tvT$, $\ment \thea\teq\thea$, and $n\mentval\QuorumBox\phi$.

This notation remains within the scope of Definition~\ref{defn.validity.judgement} --- which universally quantifies over any omitted parameter --- so omitting parameters in this way is not an abuse of notation.
However, if the reader sees us omit a value from a context, there may be two reasons for it:
\begin{enumerate*}
\item
It may indicate a universal quantification over all possible values.
For example, this is clearly what we intend when we write $\mentval\theta$ in Figure~\ref{fig.validity} (top-right).
\item
It may indicate that we have elided irrelevant parameters.
For example, this is clearly what we intend in the example of $\ment\tvT$ above.
\end{enumerate*}
In practice, it will always be clear which is intended and meaning will always be clear.\footnote{Actually, meaning will be \emph{clearer}.  

In early drafts of this paper, we dutifully included all parameters, whether relevant or not.  
But this turned out to be just too confusing: the notational overhead of carrying around values that we did not care about was burdensome; but furthermore, it was actually easier to understand the sequents once we have dropped irrelevant parameters, because this helped us to keep track of which parameters \emph{did} matter.}
\end{nttn}

\begin{figure}
$$
\begin{array}{l@{\quad}l@{\quad}l@{\quad}l@{\quad}l}
\texiunique  a.\tf{write}(a)
&
\QuorumBox\texi a.\tf{send}(a)
&
\texi a.\everyoneAll\tf{decide}(a)
&
\tall a.\modTF \tf{propose}(a)
&
\mrup{a}{(\someone\tf{accept}(a))}{v}
\\[1ex]
& \text{may be written as}
\\[1ex]
\texiunique \tf{write} 
&
\QuorumBox\texi\tf{send} 
&
\texi\everyoneAll\tf{decide}
&
\tall\modTF\tf{propose} 
&
\mru{(\someone\tf{accept})}{v}
\end{array}
$$
\caption{Abbreviations inspired by higher-order logic (Notation~\ref{nttn.hos})}
\label{fig.hos}
\end{figure}


\begin{nttn}[Functional notation]
\label{nttn.hos}
We will frequently apply quantifiers (like $\texi$, $\texiunique$, or $\tall$) to unary predicate symbols, possibly with modalities.
For example:
$$
\tall a.\QuorumBox\tf{decide}(a)
\quad\text{is}\quad\text{$\tall$ quantifying $a$ in\ $(\QuorumBox\tf{decide})$\ applied to $a$.}
$$ 
For clarity, we may use a \deffont{functional notation} inspired by higher-order logic (HOL) to elide the variable, as illustrated in Figure~\ref{fig.hos}.
Meaning will always be clear.

See also Remark~\ref{rmrk.functional.notation.in.paxos} for examples of this notation applied in an axiomatisation.
\end{nttn}






\section{Properties of predicates}
\label{sect.predicates}

This Section collects elementary technical results which we will need later.
Some of these results, and their proofs, are more technical than others, but they are not necessarily surprising if we look at what they mean in the model.
Mostly the statements are simple and it is obvious by common sense that they should work (simple is good, because it means greater productivity and fewer errors).

A typical example is Lemma~\ref{lemm.tomorrow.forever.recent}(\ref{item.tomorrow.forever.recent.1}): once we look at what it actually says about the model, this is that ``if $\phi$ is true today, then in the future it will have been true in the past''.
For saying this a mathematician earns a living? --- we might think.

What is less evident is \emph{which} elementary properties will be important later, nor is it obvious what definitions and structures make these properties all fit together nicely.
So the reader should note that the simplicity of this Section is a design feature, which was only attained in retrospect and after some work.

Broadly speaking, we start this Section with the most elementary results, and build complexity from there.
The reader is welcome to read this Section now, or to skip it and refer back to its results as they are used in later proofs.


\subsection{Equivalence of predicates, and some (de Morgan) equivalences}
\label{subsect.equivalence.of.predicates}

For the rest of this Section we make a fixed but arbitrary choice of signature $\Sigma$ (Definition~\ref{defn.signature}(\ref{item.predicate.syntax.signature})) and model $\mathcal M$ (Definition~\ref{defn.model}).
All predicates considered will be over $\Sigma$ and $\mathcal M$, unless otherwise stated.

We generalise equivalence $\equiv$ (Definition~\ref{defn.equivalent}) from truth-values, to predicates viewed as functions from contexts to truth-values:
\begin{defn}[Equivalence of predicates]
\label{defn.phi.equivalent}
Suppose $\phi$ and $\phi'$ are closed predicates.

Call $\phi$ and $\phi'$ \deffont{extensionally equivalent} when 
$$
\model{\phi}=\model{\phi'} .
$$
Unpacking Definition~\ref{defn.denotation}, this means 
$$
\Forall{\acontext\in\tf{Ctx}(\Sigma,\mathcal M)}(\modellabel{\phi}{\acontext}=\modellabel{\phi'}{\acontext}) .
$$
For brevity, we may write `extensionally equivalent' just as `equivalent' henceforth.
\end{defn}



\begin{lemm}
\label{lemm.equiv}
Suppose $\phi$ and $\phi'$ are closed predicates and $\acontext$ is a context.
Then: 
\begin{enumerate*}
\item\label{item.equiv.1}
$\acontext\ment\phi\equiv\phi'$ if and only if $\modellabel{\phi}{\acontext}=\modellabel{\phi'}{\acontext}$.
\item\label{item.equiv.2}
$\ment\phi\equiv\phi'$ if and only if $\model{\phi}=\model{\phi'}$ (i.e. $\phi$ and $\phi'$ are equivalent as per Definition~\ref{defn.phi.equivalent}).
\end{enumerate*}
\end{lemm}
\begin{proof}
We consider each part in turn:
\begin{enumerate}
\item
By Definition~\ref{defn.validity.judgement}(\ref{item.ctx.ment.phi}) $\acontext\ment\phi\equiv\phi'$ means $\modellabel{\phi\equiv\phi'}{\acontext}\in\threeValid$, and by the clause for $\equiv$ in Figure~\ref{fig.3.derived} this is if and only if $\modellabel{\phi}{\acontext}=\modellabel{\phi'}{\acontext}$.
\item
By Definition~\ref{defn.validity.judgement}(\ref{item.ment.phi}), $\ment\phi\equiv\phi'$ means precisely $\Forall{\acontext}\acontext\ment\phi\equiv\phi'$.
We use part~\ref{item.equiv.1} of this result.
\qedhere\end{enumerate}
\end{proof}


We will see many equivalences in what follows.
We start by using Definition~\ref{defn.phi.equivalent} and Lemma~\ref{lemm.equiv} to express some standard equivalences in logic, having to do with implication being expressible using negation and disjunction, and with \emph{de Morgan} dualities between modalities and connectives:
\begin{lemm}
\label{lemm.minimal.syntax}
For each of the equivalences $\f{LHS}\equiv\f{RHS}$ listed in Figure~\ref{fig.easy.equivalences}, we have that $\f{LHS}$ is extensionally equal to $\f{RHS}$ in the sense of Definition~\ref{defn.phi.equivalent}.
\end{lemm}
\begin{proof}
By routine calculations using Figures~\ref{fig.3.phi.f} and~\ref{fig.3.derived} and facts of the lattice $\THREE=(\tvF<\tvB<\tvT)$.
\end{proof}


\begin{figure}
$$
\begin{array}{r@{\ }l@{\quad}r@{\ }l@{\quad}r@{\ }l@{\quad}r@{\ }l}
\phi\tand\phi' \equiv& \tneg(\tneg\phi\tor\tneg\phi')
&
\phi\tor\phi' \equiv& \tneg(\tneg\phi\tand\tneg\phi')
&
\tneg\phi\equiv& \phi\timpc\tvF \equiv\phi\tnotor\tvF %
\\
\phi\timpc\phi' \equiv& \phi\tnotor\modT\phi'
&
\phi\tnotor\phi' \equiv& \tneg\phi\tor\phi'
&
\phi\timpc\phi' \equiv& \phi\tnotor\modT\phi'
\\
\everyone\phi\equiv&\tneg\someone\tneg\phi
&
\someone\phi\equiv&\tneg\everyone\tneg\phi
\\
\Coquorum\phi\equiv&\tneg\Quorum\tneg\phi
&
\Quorum\phi\equiv&\tneg\Coquorum\tneg\phi
\\
\forever\phi\equiv&\tneg\sometime\tneg\phi
&
\sometime\phi\equiv&\tneg\forever\tneg\phi
\\
\tall a.\phi\equiv&\tneg\texi a.\tneg\phi
&
\texi a.\phi\equiv&\tneg\tall a.\tneg\phi
&
\modTF\phi\equiv&\modT\phi\tor\modT\tneg\phi
\\
\modB\phi\equiv&\tneg\modTF\phi
&
\modT\phi\equiv&\tneg\modFB\phi
&
\modF\phi\equiv&\tneg\modTB\phi
\end{array}
$$
\caption{Some easy equivalences}
\label{fig.easy.equivalences}
\end{figure}


\subsection{Two commutation lemmas}
\label{subsect.two.commutation.results}

We note some simple commutation lemmas.
These results are all easy facts of our model, but it only takes a few lines to spell them out:

\begin{lemm}[Commuting connectives]
\label{lemm.commuting.connectives}
Suppose $\phi$ and $\phi'$ are closed predicates and suppose
$$
\bullet\in\{\tor,\tand\} 
\ \ \text{and}\ \ 
\tf M\in\{\modT,\modTB\} . 
$$
Then
$$
\ment\tf M(\phi\bullet\phi')\equiv (\tf M\phi\bullet\tf M\phi') .
$$
\end{lemm}
\begin{proof}
Routine from Lemma~\ref{lemm.equiv}(\ref{item.equiv.2}) and Figures~\ref{fig.3.phi.f} and~\ref{fig.3.derived}, and facts of $\THREE$ as a lattice.
\end{proof}


\begin{lemm}[Commuting modalities]
\label{lemm.someone.commutation}
Suppose $\phi$ is a closed predicate and suppose 
$$
\mathsf Q\in\{\texi,\tall,\someone,\everyone,\Quorum,\Coquorum,\recent,\urecent,\yesterday,\tomorrow,\forever,\sometime,\infinitely,\final\}.
$$
Then 
$$
\ment\mathsf Q\modT\phi \equiv \modT\mathsf Q\phi
\quad\text{and}\quad
\ment\mathsf Q\modTB\phi \equiv \modTB\mathsf Q\phi.
$$
\end{lemm}
\begin{proof}
Routine from Lemma~\ref{lemm.equiv}(\ref{item.equiv.2}) and Figures~\ref{fig.3.phi.f} and~\ref{fig.3.derived}, and facts of $\THREE$ as a lattice.
\end{proof}


\subsection{Lattice inequalities}

Corollary~\ref{corr.recent.everyone.commute} is in the same spirit as Lemmas~\ref{lemm.commuting.connectives} and~\ref{lemm.someone.commutation}. 
In practice we usually just use it a direct fact of the model, but we spell it out anyway.
It relies on this simple fact about lattices: 
\begin{lemm} 
\label{lemm.quantifier.swap}
Suppose $(\mathcal L,\tand,\tor)$ is a complete lattice and $X$ and $Y$ are sets and $f:(X\times Y)\to\mathcal L$ is any function. 
Then:
\begin{enumerate*}
\item\label{item.quantifier.swap.1}
$\bigvee_{x}\bigwedge_{y}f(x,y)\leq \bigwedge_{y}\bigvee_{x}f(x,y)$, 
\item\label{item.quantifier.swap.1b}
It is not in general the case that 
$\bigwedge_{y}\bigvee_{x}f(x,y) \leq \bigvee_{x}\bigwedge_{y}f(x,y)$.
\item\label{item.quantifier.swap.2}
$\bigvee_{x}\bigvee_{y}f(x,y) = \bigvee_{y}\bigvee_{x}f(x,y)$.
\end{enumerate*}
\end{lemm}
\begin{proof}
By easy lattice calculations. 
\end{proof}


\begin{corr}
\label{corr.recent.everyone.commute}
Suppose $\mathsf Q\phi$ is a closed predicate, where 
$$
\mathsf Q\in\{\texi a,\tall a \mid a{\in}\tf{VarSymb}\}\cup\{\someone,\everyone,\Quorum,\Coquorum\},
$$
and suppose $\acontext$ is a context.
Then:
\begin{enumerate*}
\item\label{item.recent.everyone.commute.1}
$\modellabel{\recent\mathsf Q\phi}{\acontext}\leq \modellabel{\mathsf Q\recent\phi}{\acontext}$. 
\item
$\modellabel{\mathsf Q\recent\phi}{\acontext}\leq \modellabel{\recent\mathsf Q\phi}{\acontext}$ for 
$\mathsf Q\in\{\texi a\mid a{\in}\tf{VarSymb}\}\cup\{\someone,\Quorum\}$.
\item
It need not be the case that $\modellabel{\mathsf Q\recent\phi}{\acontext}\leq \modellabel{\recent\mathsf Q}{\acontext}$ for 
$\mathsf Q\in\{\tall a\mid a{\in}\tf{VarSymb}\}\cup\{\everyone,\Coquorum\}$.
\end{enumerate*}
\end{corr}
\begin{proof}
Suppose $n\in\Time$ and $p\in\Pnt$ and $O\in\opensne$ and $\avaluation$ is a valuation.
We freely use \strongmodusponens.
\begin{enumerate}
\item
We consider $\mathsf Q=\Coquorum$.
We use the clause for $\recent$ in Figure~\ref{fig.3.derived}, the clause for $\Coquorum$ in Figure~\ref{fig.3.phi.f}, and then Lemma~\ref{lemm.quantifier.swap}(\ref{item.quantifier.swap.1}): 
$$
\bigvee_{0{\leq}n'{<}n}\bigwedge_{O'{\in}\opensne}\modellabel{\phi}{n',p,O',\avaluation}
\stackrel{L\ref{lemm.quantifier.swap}(\ref{item.quantifier.swap.1})}{\leq}
\bigwedge_{O'{\in}\opensne}\bigvee_{0{\leq}n'{<}n}\modellabel{\phi}{n',p,O',\avaluation} .
$$
Now we consider $\mathsf Q=\tall a$:
$$
\bigvee_{0{\leq}n'{<}n}\bigwedge_{v{\in}\Val}\modellabel{\phi[a\ssm v]}{n',p,O,\avaluation}
\stackrel{L\ref{lemm.quantifier.swap}(\ref{item.quantifier.swap.1})}{\leq}
\bigwedge_{v{\in}\Val}\bigvee_{0{\leq}n'{<}n}\modellabel{\phi[a\ssm v]}{n',p,O,\avaluation} .
$$
The other cases are precisely similar.
\item
We consider just $\mathsf Q=\someone$; the other cases are precisely similar.
We just use Lemma~\ref{lemm.quantifier.swap}(\ref{item.quantifier.swap.2}): 
$$
\bigvee_{p'{\in}\Pnt}\bigvee_{0{\leq}n'{<}n}\modellabel{\phi}{n',p',O,\avaluation}
\stackrel{L\ref{lemm.quantifier.swap}(\ref{item.quantifier.swap.2})}{=}
\bigvee_{0{\leq}n'{<}n}\bigvee_{p'{\in}\Pnt}\modellabel{\phi}{n',p',O,\avaluation} .
$$
\item
We consider just $\mathsf Q=\tall$; the other cases are precisely similar.
We note by Lemma~\ref{lemm.quantifier.swap}(\ref{item.quantifier.swap.1b}) that 
$$
\bigwedge_{v{\in}\Val}\bigvee_{0{\leq}n'_v{<}n}\modellabel{\phi[a\ssm v]}{n'_v,p,O,\avaluation}
\leq
\bigvee_{0{\leq}n'{<}n}\bigwedge_{v\in\Val}\modellabel{\phi[a\ssm v]}{n',p,O,\avaluation} 
\ \ \ \text{does not hold in general}.
$$
\qedhere\end{enumerate}
\end{proof}





\subsection{The compound modalities $\someoneAll$, $\everyoneAll$, $\QuorumBox$, and $\CoquorumDiamond$}
\label{subsect.compound.modalities}

\begin{rmrk}
\label{rmrk.compound.modalities}
We will make much use of the modalities $\someoneAll$, $\everyoneAll$, $\QuorumBox$, and $\CoquorumDiamond$.

We call these four combinations \deffont{compound modalities}, because they are built out of $\Quorum$ and $\Coquorum$ (which quantify over open sets) and $\someone$ and $\everyone$ (which quantify over points in an open set). 
In this Subsection we study some of their properties.
\end{rmrk}



Lemma~\ref{lemm.dense.char} just unpacks definitions:
\begin{lemm}
\label{lemm.dense.char}
Suppose $\phi$ is a closed predicate and $n\in\Time$ and $\avaluation$ is a valuation.
Then: 
\begin{enumerate*}
\item\label{item.dense.char.someoneAll}
$n\mentval\someoneAll\phi$ if and only if $\Exists{O\in\opensne}O\between\{p\in\Pnt \mid n,p,O\mentval\phi\}$.
\item\label{item.dense.char.everyoneAll}
$n\mentval\everyoneAll\phi$ if and only if $\Forall{O\in\opensne}O\subseteq\{p\in\Pnt \mid n,p,O\mentval\phi\}$.
\item\label{item.dense.char.QuorumBox}
$n\mentval\QuorumBox\phi$ if and only if $\Exists{O\in\opensne}O\subseteq\{p\in\Pnt \mid n,p,O\mentval\phi\}$.
\item\label{item.dense.char.CoquorumDiamond}
$n\mentval\CoquorumDiamond\phi$ if and only if $\Forall{O\in\opensne}O\between\{p\in\Pnt \mid n,p,O\mentval\phi\}$.
\end{enumerate*} 
\end{lemm}
\begin{proof}
We consider each part in turn:
\begin{enumerate}
\item
We reason as follows: 
$$
\begin{array}{r@{\ }l@{\quad}l}
n\mentval\someoneAll\phi
\liff&
\Exists{O{\in}\opensne}\Exists{p{\in}O}n,p,O\mentval\phi
&\text{Figure~\ref{fig.3.phi.f}}
\\
\liff&
\Exists{O{\in}\opensne}O\between \{p\in\Pnt \mid n,p,O\mentval\phi\}
&\text{Fact}
\end{array}
$$
\item
We reason as follows: 
$$
\begin{array}{r@{\ }l@{\quad}l}
n\mentval\everyoneAll\phi
\liff&
\Forall{O{\in}\opensne}\Forall{p{\in}O}n,p,O\mentval\phi
&\text{Figure~\ref{fig.3.phi.f}}
\\
\liff&
\Forall{O{\in}\opensne}O\subseteq\{p\in\Pnt \mid n,p,O\mentval\phi\}
&\text{Fact}
\end{array}
$$
\item
We reason as follows:
$$
\begin{array}{r@{\ }l@{\quad}l}
n\mentval\QuorumBox\phi
\liff&
\Exists{O{\in}\opensne}\Forall{p{\in}O}n,p,O\mentval\phi
&\text{Figure~\ref{fig.3.phi.f}}
\\
\liff&
\Exists{O{\in}\opensne}O\subseteq\{p\in\Pnt \mid n,p,O\mentval\phi\}
&\text{Fact}
\end{array}
$$
\item
We reason as follows:
$$
\begin{array}{r@{\ }l@{\quad}l}
n\mentval\CoquorumDiamond\phi
\liff&
\Forall{O{\in}\opensne}\Exists{p{\in}O}n,p,O\mentval\phi
&\text{Figure~\ref{fig.3.phi.f}}
\\
\liff&
\Forall{O{\in}\opensne}O\between\{p\in\Pnt \mid n,p,O\mentval\phi\}
&\text{Fact}
\end{array}
$$
\qedhere\end{enumerate}
\end{proof}



Proposition~\ref{prop.someone.implies.someoneAll} packages up some simple corollaries of Lemma~\ref{lemm.dense.char}, in forms that will be useful for our later proofs: 
\begin{prop}
\label{prop.someone.implies.someoneAll}
Suppose $\phi$ is a closed predicate and $\acontext$ is a context.
Then:
\begin{enumerate*}
\item\label{item.someone.implies.someoneAll.4}
$\acontext\ment\QuorumBox\phi\tand\CoquorumDiamond\psi\limp\acontext\ment\someoneAll(\phi\tand\psi)$.
\item\label{item.someone.implies.someoneAll.4b}
$\acontext\ment\QuorumBox\phi\tand\everyoneAll\psi\limp\acontext\ment\QuorumBox(\phi\tand\psi)$.
\item\label{item.someone.implies.someoneAll.5}
$\acontext\ment\modT(\QuorumBox\phi\tand\CoquorumDiamond\psi)\limp\acontext\ment\modT\someoneAll(\phi\tand\psi)$,\ and
\\
$\acontext\ment(\QuorumBox\phi\tand\CoquorumDiamond\psi)\timpc \someoneAll(\phi\tand\psi)$.
\item\label{item.someone.implies.someoneAll.2}
$\acontext\ment\modT\QuorumBox\phi\limp\acontext\ment\modT\someoneAll\phi$.
\item\label{item.someone.implies.someoneAll.2b}
$\acontext\ment\modT\CoquorumDiamond\phi\limp\acontext\ment\modT\someoneAll\phi$.
\end{enumerate*}
\end{prop}
\begin{proof}
We consider each part in turn:
\begin{enumerate}
\item
Routine from Lemma~\ref{lemm.dense.char}(\ref{item.dense.char.QuorumBox}\&\ref{item.dense.char.CoquorumDiamond}).
\item
Routine from Lemma~\ref{lemm.dense.char}(\ref{item.dense.char.QuorumBox}\&\ref{item.dense.char.everyoneAll}).
\item
From part~\ref{item.someone.implies.someoneAll.4} of this result taking $\modT\phi$ and $\modT\psi$ and the commutation Lemmas~\ref{lemm.someone.commutation} and~\ref{lemm.commuting.connectives}.
Then we use \strongmodusponens.
\item
From part~\ref{item.someone.implies.someoneAll.5} of this result, taking $\psi=\tvT$.
\item
From part~\ref{item.someone.implies.someoneAll.5} of this result, taking $\phi=\tvT$.
\qedhere\end{enumerate}
\end{proof}



 







\subsection{Some comments on time}

\begin{rmrk}
This Subsection notes Lemmas~\ref{lemm.tomorrow.forever.recent} and~\ref{lemm.persist.char}.
They both have to do with time, and underneath all their complicated-looking symbols they are both straightforward:
\begin{itemize*}
\item
Lemma~\ref{lemm.tomorrow.forever.recent} reflects a fact that if something holds now, then in the future it will have held.
\item
Lemma~\ref{lemm.persist.char} reflects the fact that something repeats forever if and only if it happens infinitely often.
\end{itemize*}
\end{rmrk}


\begin{lemm}
\label{lemm.tomorrow.forever.recent}
Suppose $\acontext=(n,p,O,\avaluation)$ is a context and $\phi$ is a closed predicate.
Then:
\begin{enumerate*}
\item\label{item.tomorrow.forever.recent.1}
$n,p,O\mentval\modT\phi$ implies $\Forall{n'{\gneq}n}\,n',p,O\mentval\modT\recent\phi$.
\item\label{item.rur.r}
$n,p,O\mentval\modT\urecent\phi$ implies $\Forall{n'{\geq} n}\,n',p,O\mentval\modT\urecent\phi$.
\item
As corollaries, 
$$
\acontext\ment\phi\timpc \forever\recent\phi
\quad\text{and}\quad
\acontext\ment\urecent\phi\timpc\forever\urecent\phi.
$$
\end{enumerate*}
\end{lemm}
\begin{proof}
This is just from the clauses for $\recent$ and $\urecent$ in Figure~\ref{fig.3.derived}. 
The corollary follows using \strongmodusponens.
\end{proof}



$\infinitely=\forever\sometime$ has a characterisation as `infinitely often', with the convenient feature that its denotation is independent of the stage at which it is evaluated (called $n'$ in the statement of the Lemma below):
\begin{lemm}
\label{lemm.persist.char}
Suppose $\phi$ is a closed predicate and $p\in\Pnt$ and $O\in\opensne$ and $\acontext$ is a context.
Suppose further that $n'\in\Time$ and $\f{tvs}\in\{\threeTrue,\threeValid\}$.
Then we have
$$
\modellabel{\infinitely\phi}{n',p,O,\avaluation}\in\f{tvs}
\quad\liff\quad
\{n\in\Time \mid \modellabel{\phi}{n,p,O,\avaluation}\in\f{tvs}\}\text{ is infinite} . 
$$
\end{lemm}
\begin{proof}
By routine arguments from the clause for $\infinitely$ in Figure~\ref{fig.3.derived}.
\end{proof}






\subsection{Unique existence and affine existence in three-valued logic}
\label{subsect.unique.affine.existence}

\subsubsection{The definition, and main result}

\begin{rmrk}
Distributed algorithms often require an implementation to make (at most) one choice from (zero or more) possibilities.
In our axioms, we model this using unique existence and affine existence:
\begin{itemize*}
\item
\emph{Unique existence} means `there exists precisely one'; we write this $\exists!$, as standard.
\item
\emph{Affine existence} means `there exist zero or one'; we write this $\exists_{01}$.\footnote{We are not aware of a standard symbol for affine existence.}
\end{itemize*}
See for example \rulefont{PaxPropose01} and \rulefont{PaxWrite01} in Figure~\ref{fig.logical.paxos} (a discussion of which is in Remark~\ref{rmrk.paxos.axioms.discussion}(\ref{item.01.discussion})).

So we need to think about what $\exists!$ and $\exists_{01}$ even mean in a three-valued setting.
This turns out to be not completely straightforward, because we have a third truth-value that is neither true nor false.

For concreteness, suppose $\tf P$ is a unary predicate and $\acontext$ is a context, and suppose that 
$\modellabel{\phi[a\ssm v]}{\acontext}=\tvB$ for all $v\in\Val$.
Then what truth-values should $\modellabel{\texiaffine a.\phi}{\acontext}$ and $\modellabel{\texiunique a.\phi}{\acontext}$ have?
 
Now suppose $v,v',v''\in\Val$, and suppose that $\modellabel{\phi[a\ssm v''']}{\acontext}=\tvF$ for all $v'''\in\Val\setminus\{v,v',v''\}$ and
$$
\modellabel{\phi[a\ssm v]}{\acontext}=\tvT
\quad\land\quad
\modellabel{\phi[a\ssm v']}{\acontext}=\tvB
\quad\land\quad
\modellabel{\phi[a\ssm v'']}{\acontext}=\tvB .
$$
\noindent Then what truth-values should $\modellabel{\texiaffine a.\phi}{\acontext}$ and $\modellabel{\texiunique a.\phi}{\acontext}$ have?

Now suppose that $\modellabel{\phi[a\ssm v''']}{\acontext}=\tvF$ for all $v'''\in\Val\setminus\{v\}$ and
$$
\modellabel{\phi[a\ssm v]}{\acontext}=\tvB .
$$
\noindent Then what truth-values should $\modellabel{\texiaffine a.\phi}{\acontext}$ and $\modellabel{\texiunique a.\phi}{\acontext}$ have?

Our answers to all questions above will be: $\tvB$ for both.
This is not automatic: it is a real design decision that we have to make.
\end{rmrk}

\begin{defn}
\label{defn.unique.exist}
Suppose $\phi$ is a predicate.
Define \deffont{affine existence} $\texiaffine a.\phi$ and \deffont{unique existence} $\texiunique  a.\phi$ by:
$$
\begin{array}{r@{\ }l}
\texiaffine a.\phi =& \tall a,a'.(\phi\tand\phi[a\ssm a']) \tnotor a\teq a'
\\
\texiunique  a.\phi =& (\texi a.\phi) \tand (\texiaffine a.\phi)
\end{array}
$$
\end{defn}


\begin{rmrk}
\label{rmrk.unpack.texiunique}
Suppose $\texi a.\phi$ is a closed predicate (i.e. $a\in\tf{VarSymb}$ and $\phi$ is a predicate with free variable symbols at most $\{a\}$), and suppose $\acontext$ is a context.
We spell unpack the denotations of $\texiunique$ and $\texiaffine$ from Definition~\ref{defn.unique.exist}: 
\begin{enumerate*}
\item
If $\exists v,v'\in\Val.v\neq v'\land\modellabel{\phi[a\ssm v]}{\acontext}=\modellabel{\phi[a\ssm v']}{\acontext}=\tvT$, then 
$$
\modellabel{\texiunique a.\phi}{\acontext}=\tvF=\modellabel{\texiaffine a.\phi}{\acontext}.
$$
\item
If $\exists!v\in\Val.\modellabel{\phi[a\ssm v]}{\acontext}=\tvT$ and $\Forall{v\in\Val}\modellabel{\phi[a\ssm v]}{\acontext}\in\{\tvT,\tvF\}$, then 
$$
\modellabel{\texiunique a.\phi}{\acontext}=\tvT=\modellabel{\texiaffine a.\phi}{\acontext}.
$$
\item
If $\exists!v\in\Val.\modellabel{\phi[a\ssm v]}{\acontext}=\tvT$ and $\Exists{v\in\Val}\modellabel{\phi[a\ssm v]}{\acontext}=\tvB$, then 
$$
\modellabel{\texiunique a.\phi}{\acontext}=\tvB=\modellabel{\texiaffine a.\phi}{\acontext}.
$$
\item\label{item.explain.existsunique.4}
If $\neg\Exists{v\in\Val}\modellabel{\phi[a\ssm v]}{\acontext}=\tvT$ and $\exists! v'\in\Val.\modellabel{\phi[a\ssm v]}{\acontext}=\tvB$, then 
$$
\modellabel{\texiunique a.\phi}{\acontext}=\tvB
= \modellabel{\texiaffine a.\phi}{\acontext} . %
$$
(We discuss this in Remark~\ref{rmrk.discuss.clause.4} below.)
\item\label{item.two.b}
If $\neg\Exists{v\in\Val}\modellabel{\phi[a\ssm v]}{\acontext}=\tvT$ and $\exists v,v'\in\Val.v\neq v'\land \modellabel{\phi[a\ssm v]}{\acontext}=\modellabel{\phi[a\ssm v']}{\acontext}=\tvB$, then 
$$
\modellabel{\texiunique a.\phi}{\acontext}
=\tvB=
\modellabel{\texiaffine a.\phi}{\acontext} .
$$
\item
If $\Forall{v\in\Val}\modellabel{\phi[a\ssm v]}{\acontext}=\tvF$, then
$$
\modellabel{\texiunique a.\phi}{\acontext}=\tvF
\quad\text{and}\quad
\modellabel{\texiaffine a.\phi}{\acontext}=\tvT.
$$
\end{enumerate*}
\end{rmrk}

\begin{rmrk}
\label{rmrk.discuss.texiunique}
We make two elementary observations on Remark~\ref{rmrk.unpack.texiunique}:
\begin{enumerate} 
\item
The equality $\teq$ in our logic is two-valued: $\model{v\teq v'}$ returns $\tvT$ or $\tvF$, and never $\tvB$.\footnote{In the terminology of Definition~\ref{defn.correct}(\ref{item.correct.phi}) below, equality predicates are always \emph{correct}.}
\item
Two notions of `unique existence' are in play in this discussion:
\begin{itemize*}
\item
`$\exists!v\in\Val$' refers to the actual unique existence quantifier in the real world (whatever that is).
This computes actual truth-values `true' and `false' and returns true when there actually exists a unique value.
\item
`$\texiunique a$' is part of the (derived) syntax of \QLogic, and is defined in Definition~\ref{defn.unique.exist}.
Its denotation computes elements of $\THREE=\{\tvF,\tvB,\tvT\}$.  
As is standard (though confusingly) we \emph{also} call elements of $\THREE$ truth-values; but they are truth-values of \QLogic, rather than actual truth-values in the real world (whatever that is).
\end{itemize*}
So: $\exists!$ and $\texiunique$ are both unique existence operators, but they live at different levels and they are not the same thing.\footnote{This kind of thing is far from unusual and appears also outside of papers with fancy logics.  For example in programming, a programming language may have multiple versions of the number $1$ (8-bit, floating point, 64-bit, and infinite precision) all of which are called `one', but all of which are different from one other, and also different from the actual $1\in\mathbb N$ (whatever that is).}
\end{enumerate}
\end{rmrk}


We will use unique/affine existence in our proofs via Lemma~\ref{lemm.unique.affine.existence}: 
\begin{lemm}
\label{lemm.unique.affine.existence}
Suppose $\acontext$ is a context and $\texi a.\phi$ is a closed predicate and $v,v'\in\Val$ are values.
Then:
\begin{enumerate*}
\item\label{item.unique.affine.existence.01}
$\acontext\ment\texiaffine a.\phi$ if and only if 
$\exists_{01} v\in\Val.\modellabel{\phi[a\ssm v]}{\acontext}=\tvT$.
\item\label{item.unique.affine.existence.11}
$\acontext\ment\texiunique a.\phi$ if and only if 
$\exists v\in\Val.\modellabel{\phi[a\ssm v]}{\acontext}\in\threeValid \land \exists_{01} v\in\Val.\modellabel{\phi[a\ssm v]}{\acontext}=\tvT$.
\item\label{item.unique.affine.existence.01.implies}
If $\acontext\ment\texiaffine a.\phi$ and $\acontext\ment\modT\phi[a\ssm v]$ and $\acontext\ment\modT\phi[a\ssm v']$, then $v=v'$.
\item\label{item.unique.affine.existence.11.implies}
If $\acontext\ment\texiunique a.\phi$ and $\acontext\ment\modT\phi[a\ssm v]$ and $\acontext\ment\modT\phi[a\ssm v']$, then $v=v'$.
\end{enumerate*}
\end{lemm}
\begin{proof}
We unpack Definition~\ref{defn.validity.judgement}(\ref{item.ctx.ment.phi}) %
and check the possibilities in Remark~\ref{rmrk.unpack.texiunique}.
\end{proof}


\subsubsection{Comments on an edge case}

Lemma~\ref{lemm.texiunique.all.b} notes a (perhaps surprising, at first glance) fact about our denotation, that 
if a unary predicate symbol $\tf P$ takes truth-value $\tvB$ when applied to all values (as would model a participant that has entirely crashed), then $\texiunique\tf P$ and $\texiaffine\tf P$ are valid --- even though $\tf P(v)$ is valid on every value.
This reflects the common side-condition in the literature `if the participant is correct, then \emph{stuff}', which is trivially true if the participant is \emph{not} correct: 
\begin{lemm}
\label{lemm.texiunique.all.b}
Suppose $\acontext$ is a context and $\tf P$ is a unary predicate symbol:\footnote{There is nothing mathematically special about using a \emph{unary} predicate symbol here, but in practice this is the case that we will use.}
Then
$$
\acontext\ment \tall\modB\tf P \limp (\acontext\ment \texiunique \tf P \land \acontext\ment \texiaffine \tf P).  
$$
The above statement uses Notation~\ref{nttn.hos}.
Unpacking this, we obtain:
$$
\acontext\ment \tall a.\modB\tf P(a) \limp (\acontext\ment \texiunique a.\tf P(a)\land \acontext\ment \texiaffine a.\tf P(a)) . 
$$  
\end{lemm}
\begin{proof}
Suppose $\acontext\ment \tall\modB\tf P$.
Unpacking Definition~\ref{defn.validity.judgement}(\ref{item.ctx.ment.phi}) this means $\modellabel{\tall\modB\tf P}{\acontext}\in\threeValid$, and by Figure~\ref{fig.3.derived} it follows that $\Forall{v\in\Val}\modellabel{\tf P(v)}{\acontext}=\tvB$.

By Remark~\ref{rmrk.unpack.texiunique}(\ref{item.two.b}) it follows that $\modellabel{\texiunique\tf P}{\acontext}=\tvB$.
The truth-value $\tvB$ is not true, but it is \emph{valid} by Definition~\ref{defn.validity.judgement}(\ref{item.ctx.ment.phi}), and thus $\acontext\ment\texiaffine\tf P$ and $\acontext\ment\texiunique\tf P$ as required. 
\end{proof}



\begin{rmrk}
\label{rmrk.discuss.clause.4}
Arguably, we would expect in clause~\ref{item.explain.existsunique.4} above 
(where $\neg\Exists{v\in\Val}\modellabel{\phi[a\ssm v]}{\acontext}=\tvT$ and $\exists! v'\in\Val.\modellabel{\phi[a\ssm v]}{\acontext}=\tvB$)
that $\modellabel{\texiaffine a.\phi}{\acontext}$ should equal $\tvT$, not $\tvB$. 
This is not hard to arrange, e.g. as
$$
\texiaffine a.\phi \tor \texiaffine a.\modT\phi .
$$
Either choice works with our later proofs; we never encounter an edge case where it matters.
\end{rmrk}



\subsection{Pointwise predicates (a well-behavedness property)}

For this subsection we fix a signature $\Sigma$ and a model $\mathcal M=(\Pnt,\opens,\Val)$.

\subsubsection{The definition, and basic properties}

\begin{rmrk}
\emph{Being pointwise} is a well-behavedness property on predicates, that their denotation in a context $\acontext=(n,p,O,\avaluation)$ %
does not depend on the open set parameter $O$ of the context.

Many important syntactic classes of predicates are naturally pointwise, including any predicate of the form $\tf P(v_1,\dots,v_n)$ (Lemma~\ref{lemm.atomic.O}).

Being pointwise is desirable for us because it simplifies denotation to holding `at a point';\footnote{Actually: at a stage, point, and valuation, but intuitively we tend to view the stage and valuation as fixed, and the point $p$ and open set $O$ as varying.} and also because several nice properties depend on this (see for example Proposition~\ref{prop.pointwise.dense.char}); and finally also because there is an easy syntactic condition (given in Lemma~\ref{lemm.atomic.O}(\ref{item.atomic.pointwise.2})) that guarantees that a predicate is pointwise.
\end{rmrk}
 

\begin{defn}[Pointwise predicates]
\label{defn.pointwise.meaning}
Suppose $\phi$ is a closed predicate.
\begin{enumerate*}
\item
If $n\in\Time$ and $\avaluation\in\Valuation(\Sigma,\Val)$ then call $\phi$ \deffont{pointwise in $(n,\avaluation)$} when
$$
\Forall{p\in\Pnt}\Forall{O,O'{\in}\opensne}\modellabel{\phi}{n,p,O,\avaluation}=\modellabel{\phi}{n,p,O',\avaluation} .
$$ 
In words, $\phi$ is pointwise in $(n,\avaluation)$ when at every point $p\in\Pnt$, the open set parameter $O\in\opensne$ does not affect the value of the denotation $\modellabel{\phi}{n,p,O,\avaluation}$. 
\item
Call $\phi$ \deffont{absolutely pointwise} when it is pointwise for every $n\in\Time$ and $\avaluation\in\Valuation(\Sigma,\Val)$. 
\end{enumerate*}
\end{defn}

\begin{defn}
\label{defn.pointwise.pred.nttn}
Suppose $\phi$ is closed, and $n\in\Time$ and $p\in\Pnt$ and $\avaluation$ is a valuation.
Then define $\modellabel{\phi}{n,p,\avaluation}$ by
$$
\modellabel{\phi}{n,p,\avaluation}=\bigwedge_{O\in\opensne}\modellabel{\phi}{n,p,O,\avaluation}.
$$
\end{defn}

In a manner of speaking we have already seen $\modellabel{\phi}{n,p,\avaluation}$ from Definition~\ref{defn.pointwise.pred.nttn} in our logic:
\begin{lemm}
\label{lemm.Coquorum.denote}
Suppose $n\in\Time$ and $p\in\Pnt$ and $\avaluation$ is a valuation.
Then
$$
\modellabel{\phi}{n,p,\avaluation}=\modellabel{\Coquorum\phi}{n,p,O,\avaluation}.
$$
\end{lemm}
\begin{proof}
Direct from the clause for $\Coquorum$ from Figure~\ref{fig.3.derived}.
\end{proof}

\begin{lemm}
\label{lemm.pointwise.pred}
Suppose $\phi$ is a closed predicate and $n\in\Time$ and $\avaluation$ is a valuation.
Then $\phi$ is pointwise in $(n,\avaluation)$ in the sense of Definition~\ref{defn.pointwise.meaning} 
if and only if
$$
\Forall{p\in\Pnt}\Forall{O\in\opensne}\modellabel{\phi}{n,p,O,\avaluation}=\modellabel{\phi}{n,p,\avaluation}.
$$
\end{lemm}
\begin{proof}
By a routine calculation from Definitions~\ref{defn.pointwise.meaning} and~\ref{defn.pointwise.pred.nttn}.
\end{proof}

\begin{corr}
\label{corr.pointwise.pred}
Suppose $\phi$ is a closed predicate and $n\in\Time$ and $p\in\Pnt$ and $\avaluation$ is a valuation.
Suppose $\phi$ is pointwise in $(n,\avaluation)$.
Then
$$
\Forall{O\in\opensne}n,p,O\mentval\phi\liff n,p\mentval\phi. 
$$
\end{corr}
\begin{proof}
Routine from Definition~\ref{defn.validity.judgement} and Lemma~\ref{lemm.pointwise.pred}. %
\end{proof}


Being pointwise is particularly nice because of two simple syntactic criteria that guarantee it:
\begin{lemm}[Conditions for being pointwise]
\label{lemm.atomic.O}
Suppose $\phi$ is a predicate.
Then:
\begin{enumerate*}
\item\label{item.atomic.pointwise}
If $\phi$ is atomic (Notation~\ref{nttn.logic.terminology}(\ref{item.atomic.predicate}); having the form $\tf P(v_1,\dots,v_n)$ or $\tneg\tf P(v_1,\dots,v_n)$) then it is absolutely pointwise (Definition~\ref{defn.pointwise.meaning}).
\item\label{item.atomic.pointwise.2}
If $\phi$ does not mention $\someone$ or $\everyone$ except within the scope of a $\Quorum$ or $\Coquorum$ modality, then $\phi$ is absolutely pointwise.\footnote{In particular this holds if $\Quorum$ and $\Coquorum$ and only ever mentioned in combination with $\someone$ and $\everyone$, as in the compound modalities $\someoneAll$, $\everyoneAll$, $\QuorumBox$, and $\CoquorumDiamond$.}
\end{enumerate*} 
\end{lemm}
\begin{proof}
For part~1, we just unpack the definition of $\modellabel{\phi}{}$ in Figure~\ref{fig.3.phi.f} and note that the `$\avaluation(\Pnt)_{n,p}$' used in the clause for a predicate symbol depends just on the values of $n$, $p$, and $\avaluation$, but not on the open sets $O$ and $O'$. 

Part~2 follows by a routine induction on predicate syntax, using part~1 as the base case.
\end{proof}


\subsubsection{Pointwise, as a modality in our logic}

In fact, we can express the property of being pointwise within our logic:
\begin{defn}
\label{defn.pointwise.pred}
Suppose $\phi$ is a predicate.
Then define a predicate $\pointwise{\phi}$ by
$$
\pointwise{\phi} \ =\ \everyoneAll(\Quorum\phi \equiv \Coquorum\phi) .
$$
\end{defn}


\begin{lemm}
\label{lemm.pointwise.heart}
Suppose $\phi$ is a closed predicate and $n\in\Time$ and $p\in\Pnt$ and $\avaluation$ is a valuation.
Then 
$$
n,p\mentval\Quorum\phi \equiv \Coquorum\phi
\quad\text{if and only if}\quad
\bigvee_{O'\in\opensne}\modellabel{\phi}{n,p,O',\avaluation}=\bigwedge_{O'\in\opensne}\modellabel{\phi}{n,p,O',\avaluation}.
$$
(The predicate $\Quorum\phi \equiv \Coquorum\phi$ itself satisfies the conditions for being pointwise from Lemma~\ref{lemm.atomic.O}(\ref{item.atomic.pointwise.2}), so we can sensibly write $n,p\mentval \Quorum\phi \equiv \Coquorum\phi$ above.)
\end{lemm}
\begin{proof}
From Lemma~\ref{lemm.equiv}(\ref{item.equiv.1}) and the clauses for $\Quorum$ in Figure~\ref{fig.3.phi.f}, and $\Coquorum$ from Figure~\ref{fig.3.derived} (or Lemma~\ref{lemm.Coquorum.denote}). 
\end{proof}

\begin{lemm}
Suppose $\phi$ is a closed predicate and $n\in\Time$ and $p\in\Pnt$ and $\avaluation$ is a valuation.
Then $\phi$ is pointwise in $(n,\avaluation)$ if and only if $n,p\mentval\pointwise{\phi}$. 
\end{lemm}
\begin{proof}
By routine calculations using Lemma~\ref{lemm.pointwise.heart}, and the clause for $\everyoneAll$ in Figure~\ref{fig.3.derived}, noting the elementary fact that for every $O\in\opensne$,
$$
\bigwedge_{O'\in\opensne}\modellabel{\phi}{n,p,O',\avaluation}
\leq
\modellabel{\phi}{n,p,O,\avaluation}
\leq
\bigvee_{O'\in\opensne}\modellabel{\phi}{n,p,O',\avaluation} .
$$
\end{proof}



\subsubsection{An interlude: dense sets, and sets with a nonempty open interior}


We recall a pair of standard topological definitions which will turn up repeatedly later, most notably in the context of pointwise predicates as Proposition~\ref{prop.pointwise.dense.char} immediately below, and again later in Lemma~\ref{lemm.is}:
\begin{defn}[Dense subset]
\label{defn.dense.subset}
Suppose $(\Pnt,\opens)$ is a semitopology and suppose $P\subseteq\Pnt$ is a (not necessarily open) set of points.
\begin{enumerate*}
\item\label{item.dense.subset}
Say that $P$ is \deffont{dense} when $\Forall{O{\in}\opensne}O\intersectswith P$.
\item\label{item.noi.subset}
Say that $P$ \deffont{has a nonempty open interior} when $\Exists{O{\in}\opensne}O\subseteq P$.
\item\label{item.dense.noi}
Write $\dense(P)$ when $P$ is dense, and $\noi(P)$ when $P$ has a nonempty open interior. 
\end{enumerate*} 
\end{defn}



\begin{lemm}
\label{lemm.dense.up-closed}
Suppose $(\Pnt,\opens)$ is a semitopology and $P,P'\subseteq\Pnt$.
Then:
\begin{enumerate*}
\item\label{item.dense.up-closed} 
If $P\subseteq P'$ then $\dense(P)\limp \dense(P')$ (being dense is \emph{up-closed}).
\item\label{item.noi.up-closed} 
If $P\subseteq P'$ then $\noi(P)\limp \noi(P')$ (having a non-empty open intersection is \emph{up-closed}).
\item\label{item.noi.noi} 
If $P\in\opensne$ then $\noi(P)$ (a nonempty open set has a nonempty open interior).
\end{enumerate*}
\end{lemm}
\begin{proof}
By routine sets calculations from Definition~\ref{defn.dense.subset}.
\end{proof}

\begin{lemm}
Suppose $(\Pnt,\opens)$ is a semitopology and $P\subseteq\Pnt$.
Then
$$
\dense(P) \liff \neg\noi(\Pnt\setminus P) 
\quad\text{and}\quad 
\noi(P) \liff \neg\dense(\Pnt\setminus P) . 
$$
In this sense, $\noi$ and $\dense$ form a \emph{de Morgan dual pair} of operators --- just as $\forall=\neg\exists\neg$ (as noted in Figure~\ref{fig.easy.equivalences}) and $\exists=\tneg\forall\tneg$.
\end{lemm}
\begin{proof}
We sketch the (easy) argument:
\begin{itemize}
\item
Suppose $P$ is not dense; by Definition~\ref{defn.dense.subset} this means that there exists $O\in\opensne$ such that $P\cap O=\varnothing$, thus $O\subseteq\Pnt\setminus P$, thus $\Pnt\setminus P$ has a nonempty open interior.
\item
Suppose $\Pnt\setminus P$ has a nonempty open interior $\varnothing\neq O\subseteq\Pnt\setminus P$.
Then clearly $P$ is not dense.
\qedhere\end{itemize}
\end{proof}

More on how `being dense' works in semitopologies is in~\cite[Chapter~11]{gabbay:semdca}.
The notion behaves somewhat differently than it does in topologies, but not in ways that will affect us here.

\subsubsection{Application to pointwise predicates}


Technically, Proposition~\ref{prop.pointwise.dense.char} is a corollary of Lemma~\ref{lemm.dense.char} and Corollary~\ref{corr.pointwise.pred}.
The result links the four logical modalities $\someoneAll$, $\everyoneAll$, $\QuorumBox$, and $\CoquorumDiamond$ to elementary topological properties of the sets of points at which these modalities are valid --- being nonempty, being full, having a nonempty open interior, and being dense.
We will use Proposition~\ref{prop.pointwise.dense.char}, both directly and via other lemmas that depend on it.

The most important use-case of Proposition~\ref{prop.pointwise.dense.char} will be when $\phi$ is atomic (Notation~\ref{nttn.logic.terminology}(\ref{item.atomic.predicate}) --- having the form 
$$
\tf P(v_1,\dots,v_n)
\quad\text{or}\quad
\tneg\tf P(v_1,\dots,v_n)
$$
--- since then by Lemma~\ref{lemm.atomic.O} $\phi$ is automatically pointwise:
\begin{prop}
\label{prop.pointwise.dense.char}
Suppose $\phi$ is a closed pointwise (Definition~\ref{defn.pointwise.meaning}) predicate and $n\in\Time$ and $\avaluation$ is a valuation.
Then: 
\begin{enumerate*}
\item\label{item.dense.char.pointwise.someoneAll}
$n\mentval\someoneAll\phi$ if and only if $\varnothing\neq\{p\in\Pnt \mid n,p\mentval\phi\}$. 
\item\label{item.dense.char.pointwise.everyoneAll}
$n\mentval\everyoneAll\phi$ if and only if $\Pnt=\{p\in\Pnt \mid n,p\mentval\phi\}$. 
\item\label{item.dense.char.pointwise.QuorumBox}
$n\mentval\QuorumBox\phi$ if and only if $\noi(\{p\in\Pnt \mid n,p\mentval\phi\})$.
\item\label{item.dense.char.pointwise.CoquorumDiamond}
$n\mentval\CoquorumDiamond\phi$ if and only if $\dense(\{p\in\Pnt \mid n,p\mentval\phi\})$.
\end{enumerate*} 
Recall that:
\begin{itemize*}
\item
$\dense$ and $\noi$ are from Definition~\ref{defn.dense.subset}.
\item
The notations $n,p\mentval\phi$ and $n\mentval\phi$ are from Definition~\ref{defn.validity.judgement}.
\end{itemize*}
\end{prop}
\begin{proof}
If $p$ is pointwise then by Corollary~\ref{corr.pointwise.pred} $n,p,O\mentval\phi$ if and only if $n,p\mentval\phi$, for any $O\in\opensne$, and we just simplify parts~\ref{item.dense.char.someoneAll} to~\ref{item.dense.char.CoquorumDiamond} of Lemma~\ref{lemm.dense.char}:
\begin{enumerate} 
\item
We reason as follows:
$$
\begin{array}{r@{\ }l@{\quad}l}
n\mentval\someoneAll\phi
\liff&
\Exists{O\in\opensne}O\between\{p\in\Pnt \mid n,p\mentval\phi\}
&\text{Lemma~\ref{lemm.dense.char}(\ref{item.dense.char.someoneAll})} 
\\
\liff&
\Pnt\between \{p\in\Pnt \mid n,p\mentval\phi\}
&\text{Fact}
\\
\liff&
\{p\in\Pnt \mid n,p\mentval\phi\}\neq\varnothing
&\text{Fact}
\end{array}
$$
\item
We reason as follows:
$$
\begin{array}{r@{\ }l@{\quad}l}
n\mentval\everyoneAll\phi
\liff&
\Forall{O\in\opensne}O\subseteq\{p\in\Pnt \mid n,p\mentval\phi\}
&\text{Lemma~\ref{lemm.dense.char}(\ref{item.dense.char.everyoneAll})} 
\\
\liff&
\Pnt\subseteq\{p\in\Pnt \mid n,p\mentval\phi\}
&\text{Fact}
\\
\liff&
\{p\in\Pnt \mid n,p\mentval\phi\}=\Pnt
&\text{Fact}
\end{array}
$$
\item
We reason as follows:
$$
\begin{array}{r@{\ }l@{\quad}l}
n\mentval\QuorumBox\phi
\liff&
\Exists{O\in\opensne}O\subseteq\{p\in\Pnt \mid n,p\mentval\phi\}
&\text{Lemma~\ref{lemm.dense.char}(\ref{item.dense.char.QuorumBox})} 
\\
\liff&
\noi(\{p\in\Pnt \mid n,p\mentval\phi\})
&\text{Definition~\ref{defn.dense.subset}}
\end{array}
$$
\item
We reason as follows:
$$
\begin{array}{r@{\ }l@{\quad}l}
n\mentval\CoquorumDiamond\phi
\liff&
\Forall{O\in\opensne}O\between\{p\in\Pnt \mid n,p\mentval\phi\}
&\text{Lemma~\ref{lemm.dense.char}(\ref{item.dense.char.CoquorumDiamond})} 
\\
\liff&
\dense(\{p\in\Pnt \mid n,p\mentval\phi\})
&\text{Definition~\ref{defn.dense.subset}}
\end{array}
$$
\qedhere\end{enumerate}
\end{proof}
See also Lemma~\ref{lemm.allbut.concretely}, which instantiates Proposition~\ref{prop.pointwise.dense.char}(\ref{item.dense.char.pointwise.QuorumBox}\&\ref{item.dense.char.pointwise.CoquorumDiamond}) to a special case
familiar from the literature on distributed systems.


\subsection{Correct predicates (another well-behavedness property)}

\begin{rmrk}
We defined $\threeCorrect=\{\tvT,\tvF\}$ in Definition~\ref{defn.tv.ment}(\ref{item.tv.ment.correct}) and a corresponding modality $\modTF$ in Figure~\ref{fig.3} (which returns $\tvT$ when its argument is in $\threeCorrect$ and $\tvF$ otherwise), and we studied some of their properties, most notably in Lemma~\ref{lemm.TF.correct}.
 
We now extend this to notions of \emph{correctness} on predicates $\phi$ and predicate symbols $\tf P$.

Correctness turns out to be highly relevant to how we will use our logic.
Partly this is because when a predicate is correct, we are returned to the familiar two-valued world of Boolean logic --- but more subtly, correctness is interesting because of how this concept interacts with the predicate and modal parts of our logic; in certain circumstances we can use a combination of correctness \emph{somewhere} and validity \emph{everywhere} to deduce information about truth (see Remark~\ref{rmrk.why.exciting}) and this turns out to be a good logical model for how reasoning on distributed algorithms is often done in practice.

We start with the definition and some basic lemmas, and develop from there:
\end{rmrk}


\begin{defn}[Correctness]
\label{defn.correct}
Suppose $\phi$ is a predicate and $a\in\tf{VarSymb}$ is a variable symbol. 
Then:
\begin{enumerate*}
\item\label{item.correct.phi}
Consistent with the definition of $\threeCorrect=\{\tvT,\tvF\}$ from Definition~\ref{defn.tv.ment}(\ref{item.tv.ment.correct}),
we define $\correct{\phi}$, read \deffont{$\phi$ is correct}, by
$$
\correct{\phi}=\modTF\phi .
$$
Note that $\phi$ being correct does not mean that $\phi$ returns $\tvT$; it means that $\phi$ returns $\tvT$ or $\tvF$ and in particular it does \emph{not} return $\tvB$.
\item\label{item.correct}\label{item.crashed}
Suppose $\tf P\in\PredSymb$ is a unary predicate symbol.
Define a predicate $\correct{\tf P}$ by 
$$
\correct{\tf P}\ =\ \tall a.\correct{\tf P(a)}\ \stackrel{N\ref{nttn.hos}}{=}\ \tall\modTF\tf P
 .
$$
\item
If $\tf P_1,\dots,\tf P_n$ are unary predicate symbols then we may write 
$$
\correct{\tf P_1,\dots,\tf P_n}
\quad\text{for}\quad 
\correct{\tf P_1}\tand\dots\tand\correct{\tf P_n}.
$$
(See for example axiom \rulefont{LdrCorrect} in Figure~\ref{fig.logical.paxos}.) 
\item\label{item.P-correct}
Suppose $n\in\Time$ and $p\in\Pnt$ and $\avaluation$ is a valuation, and suppose $\tf P$ is a predicate symbol.
Then when $n,p\mentval\correct{\tf P}$ we may call $p$ \deffont{$\tf P$-correct}. 
\end{enumerate*}
\end{defn}

\begin{lemm}
\label{lemm.correct.valid}
Suppose $\acontext$ is a context and $\phi$ is a closed predicate.
Then 
$$
\acontext\ment\phi
\quad\text{if and only if}\quad
\modellabel{\phi}{\acontext}\in\threeCorrect.
$$
\end{lemm}
\begin{proof}
This just re-states Lemma~\ref{lemm.tv.ment.TF}(\ref{item.tv.ment.TF.2}).
\end{proof}


\begin{rmrk}
\label{rmrk.tf.synonyms}
Definition~\ref{defn.correct}(\ref{item.correct.phi}) gives us two synonyms for the modality expressing that the truth-value returned by a predicate $\phi$ is either $\tvT$ or $\tvF$ (not $\tvB$): $\modTF$ and $\tf{correct}$.

Although these are synonymous, there is a difference in emphasis:
\begin{itemize*}
\item
$\modTF$ is a logical modality and just means `true or false'.
\item
The literature on distributed systems uses a term `correctness', which indicates that a process succeeds and returns a correct value. 
In our logical model this corresponds to returning $\tvT$ or $\tvF$ --- and thus not returning $\tvB$. 
So when we write $\tf{correct}$, the emphasis is on `not $\tvB$' with emphasis on a process `not failing'.
\end{itemize*}
We use whichever notation seems most natural in context. 
For example, in Figure~\ref{fig.logical.paxos}:
\begin{enumerate*}
\item
In \rulefont{LdrCorrect} we write $\correct{\tf{propose},\tf{write},\tf{decide}}$ because intuitively we are asserting that at a particular time the leader \emph{does not crash} when proposing, writing, or deciding.
\item
In \rulefont{LdrExist} we write $\modTF\tf{leader}$, because $\tf{leader}$ is a predicate that identifies whether a participant is the leader at a particular time, and this is either true or false.\footnote{In practice determining who the leader is require a participant to compute a function, but the computation is not distributed, in the sense that prior agreement amongst participants is assumed on what the function is and what seed data it uses, and each participant computes it locally to get (the same) definitive yes/no answer.}
\end{enumerate*} 
If the reader prefers, they can safely read $\modTF$ every time they see $\tf{correct}$, and they will be precisely correct (pun intended). 
\end{rmrk}



Many predicates of interest are correct, including correctness itself, our truth-modalities, and equality and equivalence:
\begin{lemm}
\label{lemm.equality.correct}
Suppose $\acontext$ is a context and $\phi$ and $\phi'$ are closed predicates and $v,v'\in\Val$.
Then:
\begin{enumerate*}
\item
$\modellabel{\correct{\correct{\phi}}}{\acontext}=\tvT$.
\quad We can (perhaps cheekily) read this as `correctness is correct'.
\item
$\modellabel{\correct{\tf Q\phi}}{\acontext}=\tvT$ for $\tf Q\in\{\modT,\modF,\modB,\modTB,\modTF,\modFB\}$. 
\item\label{item.equality.correct}
$\modellabel{\correct{v\teq v'}}{\acontext}=\tvT$.
In words: equality is correct.
\item\label{item.equivalence.correct}
$\modellabel{\correct{\phi\equiv \phi'}}{\acontext}=\tvT$.
In words: equivalence is correct.
\end{enumerate*}
\end{lemm}
\begin{proof}
By routine checking from Figures~\ref{fig.3}, \ref{fig.3.phi.f}, \ref{fig.3.derived}, and (for $\equiv$)~\ref{fig.3.iff}.
\end{proof}

Lemma~\ref{lemm.correct.P.Pv} and Corollary~\ref{corr.modTF.texi} are straightforward, but we spell it out:
\begin{lemm}
\label{lemm.correct.P.Pv}
Suppose $\tf P\in\PredSymb$ is a unary predicate symbol and $v\in\Val$ and $\acontext$ is a context.
Then 
$$
\modellabel{\correct{\tf P}}{\acontext} = \bigwedge_{v{\in}\Val}\modellabel{\correct{\tf P(v)}}{\acontext}
\quad\text{and}\quad 
\modellabel{\correct{\tf P}}{\acontext}\leq \modellabel{\correct{\tf P(v)}}{\acontext} .
$$
\end{lemm}
\begin{proof}
By Definition~\ref{defn.correct}(\ref{item.correct}) and Figure~\ref{fig.3.derived} 
$$
\modellabel{\correct{\tf P}}{\acontext} 
= 
\modellabel{\tall a.\correct{\tf P(a)}}{\acontext} 
= 
\bigwedge_{v{\in}\Val} \modellabel{\correct{\tf P(v)}}{\acontext} 
\leq 
\modellabel{\correct{\tf P(v)}}{\acontext} .
\qedhere$$
\end{proof}


\begin{corr}
\label{corr.modTF.texi}
Suppose $\acontext$ is a context and $\tf P\in\PredSymb$ is a unary predicate symbol and $a\in\tf{VarSymb}$ is a variable symbol. 
Then 
$$
\acontext\ment\correct{\tf P}
\quad\text{implies}\quad
\acontext\ment\correct{\texi\tf P}\ \land\ \acontext\ment\correct{\tall\tf P}.
$$
\end{corr}
\begin{proof}
Suppose $\acontext\ment\correct{\tf P}$.
By Lemma~\ref{lemm.correct.P.Pv} we see that $\Forall{v\in\Val}\acontext\ment\correct{\tf P(v)}$.
By Lemma~\ref{lemm.correct.valid} $\Forall{v\in\Val}\modellabel{\tf P(v)}{\acontext}\in\threeCorrect$.
Now the clauses for $\texi$ in Figure~\ref{fig.3.phi.f} and $\tall$ in Figure~\ref{fig.3.derived} take a supremum and infimum respectively of $\{\modellabel{\tf P(v)}{\acontext} \mid v\in\Val\}\subseteq\threeCorrect=\{\tvT,\tvF\}$, and this supremum or infimum is also clearly in $\threeCorrect=\{\tvT,\tvF\}$.
We just use Lemma~\ref{lemm.correct.valid}. 
\end{proof}


Proposition~\ref{prop.tv.ment.TF.model}(\ref{item.tv.ment.TF.phi.lr}) is a wrapper for Lemma~\ref{lemm.TF.correct}.
This deceptively simple equivalence turns out to be a key way that correctness can be usefully applied:
\begin{prop}
\label{prop.tv.ment.TF.model}
Suppose $\phi$ is a closed predicate and $\acontext$ is a context and $v\in\Val$.
Then:
\begin{enumerate*}
\item\label{item.tv.ment.TF.phi.lr}
$\acontext\ment\correct{\phi}$ if and only if $\acontext\ment\phi\equiv\modT\phi$. 

In words: correct $\phi$ are characterised precisely by validity being equivalent to truth.
\item\label{item.tv.ment.TF.P}
$\acontext\ment\correct{\tf P}$ if and only if $\Forall{v\in\Val}(\acontext\ment\tf P(v)\equiv\modT\tf P(v))$. 
\item\label{item.modTF.texi.P.valid}
$\acontext\ment\correct{\tf P}$ implies $\acontext\ment\texi\tf P\equiv \modT\texi\tf P$ and $\acontext\ment\tall\tf P\equiv\modT\tall\tf P$. 
\end{enumerate*}
\end{prop}
\begin{proof}
We consider each part in turn:
\begin{enumerate}
\item
Suppose $\acontext\ment\correct{\phi}$.
By Lemma~\ref{lemm.TF.correct} $\modT\modellabel{\phi}{\acontext}=\modellabel{\phi}{\acontext}$.
By Lemma~\ref{lemm.equiv}(\ref{item.equiv.1}) and the clause for $\modT$ in Figure~\ref{fig.3.phi.f}, $\acontext\ment \phi\equiv\modT\phi$.
The reverse implication just reverses the same reasoning.
\item
We combine Lemma~\ref{lemm.correct.P.Pv} with part~\ref{item.tv.ment.TF.phi.lr} of this result for $\phi=\tf P(v)$ for every $v\in\Val$.
\item
From Corollary~\ref{corr.modTF.texi} and Proposition~\ref{prop.tv.ment.TF.model}(\ref{item.tv.ment.TF.phi.lr}).
\qedhere\end{enumerate}
\end{proof}




Proposition~\ref{prop.correct.cobox.boxT} explores how correctness interacts with our compound modalities from Subsection~\ref{subsect.compound.modalities}.
\begin{prop}
\label{prop.correct.cobox.boxT}
Suppose $\phi$ is a closed predicate and $\tf P\in\PredSymb$ is a unary predicate symbol. 
Suppose $n\in\Time$ and $\avaluation$ is a valuation.
Then:
\begin{enumerate*}
\item\label{item.correct.cobox.boxT.1}
$n\mentval\QuorumBox\correct{\phi}\tand\everyoneAll\phi$ implies $n\mentval \modT\QuorumBox\phi$.\footnote{Alert: note that we write $\everyoneAll\phi$ in the right-hand conjunct, not $\CoquorumDiamond\phi$.}
\item\label{item.correct.cobox.boxT.2}
$n\mentval\QuorumBox\correct{\tf P}\tand\everyoneAll\texi\tf P$ implies 
$n\mentval\modT\QuorumBox\texi\tf P$.
\end{enumerate*}
\end{prop}
\begin{proof}
We consider each part in turn:
\begin{enumerate}
\item
Suppose $n\mentval\QuorumBox\correct{\phi}\tand\everyoneAll\phi$.
By Proposition~\ref{prop.someone.implies.someoneAll}(\ref{item.someone.implies.someoneAll.4b})
$n\mentval\QuorumBox(\correct{\phi}\tand\phi)$.
By Proposition~\ref{prop.tv.ment.TF.model}(\ref{item.tv.ment.TF.phi.lr}) and Lemma~\ref{lemm.someone.commutation} (for $\QuorumBox$)
we have $n\mentval\modT\QuorumBox\phi$ as required.
\item
We combine part~\ref{item.correct.cobox.boxT.1} of this result with Corollary~\ref{corr.modTF.texi}.
\qedhere\end{enumerate}
\end{proof}


\begin{rmrk}[Why this is interesting]
\label{rmrk.why.exciting}
The statement of Proposition~\ref{prop.correct.cobox.boxT} is quite heavy on the symbols.
For clarity, let us re-read it, but let us replace $\ment\QuorumBox\tf{correct}$ with the phrase `mostly correct', and $\ment\everyoneAll$ with the phrase `everywhere valid', and $\texi$ with the phrase `for some value'.
Then we get:
\begin{enumerate*}
\item
If $\phi$ is \emph{mostly correct} and \emph{everywhere valid}, then it is \emph{mostly true}.
\item
If $\tf P$ is \emph{mostly correct} and \emph{everywhere valid} for some value, then it is \emph{mostly true} for some value.
\end{enumerate*}
So what we see is that we can promote weak statements about validity everywhere, with stronger but more limited statements about correctness, to obtain limited statements about truth.

The reader familiar with how distributed algorithms work may recognise this style of reasoning.
So: here it is beginning to appear in our logic.
\end{rmrk}

\subsection{$\mrup{a}{\phi}{v}$: a most recent $v$ that makes $\phi[a\ssm v]$ true}
\label{subsect.mrup}

\begin{rmrk}
\label{rmrk.explain.mrup}
We see from Figure~\ref{fig.3.derived} and Remark~\ref{rmrk.comments.on.semantics}(\ref{item.mru.program}) that $\mrup{a}{\phi}{v}$ looks back into the past, searching for a most recent time when $\model{\texi a.\phi}=\tvT$, and checks if at that most recent time, $\model{\phi[a\ssm v]}=\tvT$.
$\mrup{a}{\phi}{v}$ returns $\tvT$ if this search succeeds, and $\tvF$ if this fails.
It never returns $\tvB$.

In words: $\mrup{a}{\phi}{v}$ checks whether $v$ is a value that has most recently made $\phi$ true.

We need $\mrup{a}{\phi}{v}$ for our axiomatisation of Paxos in Figure~\ref{fig.logical.paxos}, where 
\begin{itemize*}
\item
in \rulefont{PaxSend?} we write $\mru{\tf{accept}}{v}$\ (short for $(\mrup{a}{\tf{accept}(a))}{v}$, as per Figure~\ref{fig.hos}), and 
\item
in \rulefont{PaxWrite?} we write $\mru{(\someone\tf{accept})}{v}$\ (short for $\mrup{a}{(\someone\tf{accept}(a))}{v}$).
\end{itemize*} 
We read $\mrup{a}{\phi}{v}$ as 
\begin{quote}
\deffont{$v$ is a value that most recently makes $\phi[a\ssm v]$ true} .
\end{quote}
In the common case that $\phi=\tf P(a)$ for some unary predicate symbol $\tf P$, then we might just say
\begin{quote}
\deffont{$v$ is a value that most recently makes $\tf P$ true}.
\end{quote} 
In this Subsection we study the properties of this construct in more detail.
\end{rmrk}



\begin{lemm}
\label{lemm.mru.unambivalent}
Suppose $\texi a.\phi$ is a closed predicate and $v\in\Val$ and $\acontext$ is a context.
Then:
\begin{enumerate*}
\item
$\acontext\ment\correct{\mrup{a}{\phi}{v}}$.
\item
$\acontext\ment\mrup{a}{\phi}{v}\equiv\modT\mrup{a}{\phi}{v}$.
\end{enumerate*}
\end{lemm}
\begin{proof}
We consider each part in turn:
\begin{enumerate}
\item
By Definition~\ref{defn.correct}(\ref{item.correct.phi}) $\correct{\mrup{a}{\phi}{v}}=\modTF\mrup{a}{\phi}{v}$, and this means that $\mrup{a}{\phi}{v}$ always returns $\tvT$ or $\tvF$ (but not $\tvB$).
We just check the relevant clause in Figure~\ref{fig.3.derived} and see that this is so.
\item
We just use part~1 of this result and Proposition~\ref{prop.tv.ment.TF.model}(\ref{item.tv.ment.TF.phi.lr}).
\qedhere\end{enumerate}
\end{proof}


\begin{rmrk}
\label{rmrk.explain.mru}
\leavevmode
\begin{enumerate}
\item\label{item.explain.mru.T}
Lemma~\ref{lemm.mru.unambivalent}(2) %
notes that $\mrup{a}{\phi}{v}$ and $\modT(\mrup{a}{\phi}{v})$ are extensionally equivalent (Definition~\ref{defn.phi.equivalent}).

We may write one or the other, depending on what seems most clear in a proof.
In particular, if we are using \strongmodusponens then we might retain a $\modT$ modality for clarity.
\item 
We try to head off an easy error: note that $\mrup{a}{\phi}{v}$ is not equivalent to $\modT\recent\phi[a\ssm v]$, nor the (by Lemma~\ref{lemm.someone.commutation}) equivalent form $\recent\modT\phi[a\ssm v]$.
\begin{itemize*}
\item
$\modT\recent\phi[a\ssm v]$ returns $\tvT$ when $\phi[a\ssm v]$ has returned $\tvT$ at some point in the past.
\item
$\mrup{a}{\phi}{v}$ (`$v$ is a value that most recently makes $\phi[a\ssm v]$ true') returns $\tvT$ when $\phi[a\ssm v]$ has returned $\tvT$ in \emph{the most recent stage} such that $\texi a.\phi$ returned $\tvT$ (i.e. such that $\phi[a\ssm v']$ returned $\tvT$ for \emph{any} $v'$).
\end{itemize*}
These are not the same thing.
\end{enumerate}
\end{rmrk}

Remark~\ref{rmrk.explain.mru} notes that $\mrup{a}{\phi}{v}$ is different from $\modT\recent\phi[a\ssm v]$.
That said, there is a close link between them, which will turn out to be important for understanding our axiomatisation of Paxos (see Lemma~\ref{lemm.evaporates.simplify} and surrounding discussion): 
\begin{lemm}
\label{lemm.evaporates}
Suppose $\texi a.\phi$ is a closed predicate and $\acontext=(n,p,O,\avaluation)$ is a context. 
Then
$$
\acontext\ment(\texi v.\mrup{a}{\phi}{v})\equiv \modT\recent\texi a.\phi.
$$
\end{lemm}
\begin{proof}
First, note from Lemma~\ref{lemm.someone.commutation} that $\ment\modT\recent\texi a.\phi\equiv \recent\texi a.\modT\phi$.
And, note that $\texi a.\phi$ and $\texi v.\phi[a\ssm v]$ denote the same syntax, as per Remark~\ref{rmrk.pedant.point}(\ref{item.pedant.texi.v}).

We unpack the clause for $\modellabel{\mrup{a}{\phi}{v}}{\acontext}$ in Figure~\ref{fig.3.derived} and note that it returns $\tvT$ when there exist $v\in\Val$ and $0\leq n'<n$ such that $n',p,O\mentval \modT\phi[a\ssm v]$, and $\tvF$ otherwise.
But this is also what $\modellabel{\recent\texi v.\modT(\phi[a\ssm v])}{\acontext}$ does, so by Lemma~\ref{lemm.equiv}(\ref{item.equiv.1}) we are done.
\end{proof}

Lemma~\ref{lemm.mru.today.yesterday} connects $\mrup{a}{\phi}{v}$ to the fact that our notion of time is discrete and so has a well-defined notion of a `tomorrow' (i.e. a `next timestep'):
\begin{lemm}
\label{lemm.mru.today.yesterday}
Suppose $\texi a.\phi$ is a closed predicate and $v\in\Val$ and $\acontext$ is a context. 
.
Then
$$
\acontext\ment \tomorrow(\mrup{a}{\phi}{v}) \equiv (\modT\phi[a\ssm v] \tor ((\tneg\modT\texi a.\phi) \tand \mrup{a}{\phi}{v})).
$$
\end{lemm}
\begin{proof}
It is a fact of Figure~\ref{fig.3.derived} that $\acontext\ment\tomorrow\yesterday\psi\equiv\psi$ for any closed $\psi$.
The result follows from this observation and from the clause for $\mrup{a}{\phi}{v}$ in Figure~\ref{fig.3.derived}. 
\end{proof}



\begin{corr}
\label{corr.mru.today.to.yesterday}
Suppose $\texi a.\phi$ is a closed predicate and $v\in\Val$ and $\acontext$ is a context. 
Then:
\begin{enumerate*}
\item\label{item.mru.today.implies.tomorrow}
$\acontext\ment \phi[a\ssm v] \timpc \tomorrow(\mrup{a}{\phi}{v})$.
\item\label{item.mru.tomorrow.options}
$\acontext\ment \tomorrow(\mrup{a}{\phi}{v}) \timpc (\phi[a\ssm v]\tor \mrup{a}{\phi}{v})$. 
\end{enumerate*}
\end{corr}
\begin{proof}
Elementary from Lemma~\ref{lemm.mru.today.yesterday} using \strongmodusponens.
\end{proof}


\begin{lemm}
\label{lemm.mru.to.recent}
Suppose $\texi a.\phi$ is a closed predicate %
and $v,v'\in\Val$ and $\acontext$ is a context. 
Then:
\begin{enumerate*}
\item\label{item.mru.to.recent.2}
$\acontext\ment(\mrup{a}{\phi}{v}\tand\mrup{a}{\phi}{v'}) \timpc \recent(\phi[a\ssm v]\tand\phi[a\ssm v'])$.
\item\label{item.mru.to.recent.1}
$\acontext\ment(\mrup{a}{\phi}{v}) \timpc \recent(\phi[a\ssm v])$.

If $\phi=\modality{\tf Q}\tf P(a)$ for a unary predicate symbol $\tf P\in\PredSymb$ and some prefix of modalities $\modality{\tf Q}$, then we can elide the variable symbol $a$ as per Notation~\ref{nttn.hos} and Figure~\ref{fig.hos}, and this becomes
$$
\acontext\ment\mru{(\modality{\tf Q}\tf P)}{v} \timpc \recent\modality{\tf Q}\tf P(v).
$$
\end{enumerate*}
\end{lemm}
\begin{proof}
We consider each part in turn:
\begin{enumerate}
\item
Unpacking the clause for $\mruarrow$ in Figure~\ref{fig.3.derived}, if $v$ and $v'$ are values that most recently make $\phi[a\ssm\text{-}]$ be $\tvT$, then certainly $\recent(\phi[a\ssm v]\tand\phi[a\ssm v'])$.
\item
Follows from part~\ref{item.mru.to.recent.2} as the special case that $v'=v$.
\qedhere\end{enumerate}
\end{proof}


\section{\lowercase{n}-twined semitopologies}
\label{sect.n-twined}



\subsection{Antiseparation properties}

We are particularly interested in models over semitopologies whose open sets have certain \emph{antiseparation properties}, by which we mean properties that involve the existence of nonempty intersections between open sets.\footnote{For comparison, topology is often concerned with \emph{separation} properties, like Hausdorff separation, which have to do with existence of \emph{empty} intersections between open sets.  In this sense, the interest in this Section is the precise dual to where the emphasis traditionally lies in topology.}

The design space of antiseparation properties is rich, as the first author explored in~\cite{gabbay:semdca} --- but for our needs in this paper we will just one canonical property which we will call being \emph{$n$-twined}.
We define the $n$-twined property in Definition~\ref{defn.twined}, we explain why it matters in Remark~\ref{rmrk.motivate.n-twined}, and then we study its properties.


\begin{defn}
\label{defn.twined}
\leavevmode
Suppose $(\Pnt,\opens)$ is a semitopology and $n\in\Ngeqz$.
Then call $(\Pnt,\opens)$ \deffont{$n$-twined} when every collection of $n$ nonempty open sets has a nonempty open intersection.
In symbols:
$$
\Forall{O_1,\dots,O_n\in\opensne}O_1\cap\dots\cap O_n\neq\varnothing.
$$
\end{defn}

\begin{lemm}
\label{lemm.1-twined}
We note two edge cases:
\begin{enumerate*}
\item
A semitopology is $0$-twined if and only if it is nonempty (Notation~\ref{nttn.empty.semitopology}).
\item
Every semitopology is $1$-twined.
\end{enumerate*}
\end{lemm}
\begin{proof}
We consider each part in turn:
\begin{enumerate}
\item
If $n=0$ then we need to show that $\Pnt$ (the intersection of zero nonempty sets) is not equal to $\varnothing$.
by Notation~\ref{nttn.empty.semitopology} this is true precisely when the semitopology is nonempty. 
\item
By Definition~\ref{defn.twined}, we must show that $O\neq\varnothing$ for every $O\in\opensne$.
This is true.
\qedhere\end{enumerate}
\end{proof}

\begin{rmrk}[Discussion of the special cases $n=2$ and $n=3$]
\label{rmrk.motivate.n-twined}
Lemma~\ref{lemm.1-twined} shows that being $0$-twined and $1$-twined are straightforward properties.
In contrast, the $2$- and $3$-twined semitopologies will be very interesting to us in this paper.
We give a high-level discussion of why this is:
\begin{itemize}
\item
Suppose we are in a distributed system running on a $2$-twined semitopology.
We assume that participants are honest.
Suppose further that we have 
\begin{itemize*}
\item
a nonempty open set of participants $O\in\opensne$ who agree that some predicate $\phi$ is true (respectively: valid), and 
\item
a nonempty open set of participants $O'\in\opensne$ who agree that some predicate $\phi'$ is true (respectively: valid).
\end{itemize*}
Then because the system is $2$-twined, there exists some $p\in O\cap O'$ who agrees that both $\phi$ and $\phi'$ are true (respectively: valid).\footnote{We will render this in our logic later on, in Lemma~\ref{lemm.supertwined.qcbt}(\ref{item.intertwined.qcbt.2}).}

This simple observation underlies guarantees of consistent outputs in many distributed algorithms, because it guarantees that $\phi$ and $\phi'$ are not logically conflicting --- e.g. that we do not have something like $\phi'=\modF\phi$ --- because if they were, then an honest $p$ could not believe both.

In the special case that $\phi$ is `the output of our consensus algorithm is $v$' and $\phi'$ is `the output of our consensus algorithm is $v'$', the $2$-twined property makes the strong guarantee that where consensus algorithms succeed (i.e. where participants do not crash), the output is consistent. 
\item
Suppose we are in a distributed system running on a $3$-twined semitopology.
Some participants may be dishonest, in that they may report inconsistent predicates. 

Suppose further that we have nonempty open sets $O,O'\in\opensne$ of participants who agree on the truth (respectively: validity) of predicates $\phi$ and $\phi'$ respectively.
Suppose further that we have a nonempty open set $O''\in\opensne$ of participants who are honest (meaning that we bound dishonest participants to at most some \emph{closed} set). 
Because the system is $3$-twined, there exists some $p\in O\cap O'\cap O''$ --- this $p$ is a participant who agrees on the truth (respectively: validity) of $\phi$ and $\phi'$ and is also honest.\footnote{We will render this in our logic twice: in Lemma~\ref{lemm.supertwined.qcbt}(\ref{item.supertwined.qcbt.2}) as a general result, and then as a useful special case in Corollary~\ref{corr.two.quorums.intersect.correct}(\ref{item.two.quorums.intersect.correct.supertwined}).}

In this case, $\phi\tand\phi'$ must be consistent, because if it were not then the honest participant $p$ could not believe both $\phi$ and $\phi'$.
Just as for the $2$-twined case, this simple argument underlies many consistency arguments in distributed systems in the presence of dishonest participants.
\end{itemize}
There is no barrier to considering $n$-twined semitopologies for $n\geq 4$, and this would just give us stronger guarantees of intersection.

A brief survey of intersection properties in the literature is in Remark~\ref{rmrk.ntwined.and.the.literature}
\end{rmrk}


\subsection{Being $n$-twined and being dense / having nonempty open interior}

The property of being $n$-twined is closely related to the topological notion of being a \emph{dense set} and having a \emph{nonempty open interior} (Definition~\ref{defn.dense.subset}).
We explore this in Lemma~\ref{lemm.is} and Proposition~\ref{prop.noi.dense}; then in Proposition~\ref{prop.allbut.supertwined} we specialise the discussion to the important special case of $\AllBut{N}{f}$.

\begin{lemm}
\label{lemm.is}
Suppose $(\Pnt,\opens)$ is a semitopology and $n\in\mathbb N_{\geq 1}$.
Then:
\begin{enumerate*}
\item\label{item.is.supertwined}
$(\Pnt,\opens)$ is $n$-twined if and only if every collection of $n\minus 1$ nonempty open sets has a dense intersection: 
$$
\Forall{O_1,\dots O_{n\minus 1}\in\opensne}\dense(O_1\cap\dots\cap O_{n\minus 1}).
$$
\item\label{item.is.intertwined}
As a corollary, $(\Pnt,\opens)$ is \xxtwined if and only if every nonempty open set is dense: 
$$
\Forall{O\in\opensne}\dense(O).
$$
\end{enumerate*}
\end{lemm}
\begin{proof}
Both parts follow from Definitions~\ref{defn.dense.subset}(1) and~\ref{defn.twined}, just by pushing around definitions.

One corner case merits brief mention: if $n=1$ then we must show that every collection of $0$ sets (i.e. $\varnothing$) is dense.
So, we need to show that $\bigcap\varnothing$ is dense.
But by convention, the intersection of \emph{no} open sets is the whole of $\Pnt$, and this is indeed dense, so we are done. 
\end{proof}


\begin{prop}
\label{prop.noi.dense}
Suppose $(\Pnt,\opens)$ is a semitopology and $n\in\mathbb N_{\geq 1}$.
Then the following are equivalent:
\begin{enumerate*}
\item
$(\Pnt,\opens)$ is $n$-twined. 
\item
$\Forall{P_1,\dots,P_{n\minus 1}\subseteq\Pnt}(\bigwedge_{1{\leq}i{\leq}n\minus 1}\noi(P_i)) \limp \dense(\bigcap_{1{\leq}i{\leq}n\minus 1}P)$.
\end{enumerate*}
\end{prop}
\begin{proof}
Part~2 trivially implies part~1 because we can take $P_1,\dots,P_n\in\opensne$, so that by Definition~\ref{defn.dense.subset} $\noi(P_i)$ for $1\leq i\leq n\minus 1$ holds by construction by Lemma~\ref{lemm.dense.up-closed}(\ref{item.noi.noi}), and we use Lemma~\ref{lemm.is}(\ref{item.is.supertwined}).

Showing that part~1 implies part~2 is only slightly less trivial.
Suppose $(\Pnt,\opens)$ is $n$-twined and $P_1,\dots,P_{n\minus 1}\subseteq\Pnt$ and $\noi(P_i)$ for each $1\leq i\leq n\minus 1$, meaning by Definition~\ref{defn.dense.subset} that for each such $i$ there exists an $O_i\in\opensne$ such that $O_i\subseteq P_i$.
By Lemma~\ref{lemm.is}(\ref{item.is.supertwined}) $O_1\cap\dots\cap O_{n\minus 1}$ is dense, and by Lemma~\ref{lemm.dense.up-closed}(\ref{item.dense.up-closed}) so is $P_1\cap\dots\cap P_n$.
\end{proof}



Proposition~\ref{prop.allbut.supertwined} provides the canonical example of $n$-twined semitopologies.
It mildly generalises known properties in the literature; the case $n=2$ in Proposition~\ref{prop.allbut.supertwined}(\ref{item.allbut.intertwined}) corresponds to the $Q^2$ property from~\cite[Theorem~1]{DBLP:conf/asiacrypt/DamgardDFN07}, and the case $n=3$ corresponds to the $Q^3$ property of~\cite[Definition 2]{DBLP:journals/dc/AlposCTZ24}.
The topology-flavoured presentations in Proposition~\ref{prop.allbut.supertwined}(\ref{item.allbut.NOI.dense}\&\ref{item.allbut.NOI.dense.3}) are new.
\begin{prop}
\label{prop.allbut.supertwined}
Suppose $N,f\in\Ngeqz$ and $P\subseteq\Nset{N}$ and $n\in\mathbb N_{\geq 1}$.
Then:
\begin{enumerate*}
\item\label{item.allbut.intertwined}\label{item.allbut.supertwined}
$N>n*f$ if and only if $\AllBut{N}{f}$ is $n$-twined.
\item\label{item.allbut.NOI.dense}
$\noi(P)$ (and equivalently $P\in\opensne$) if and only if $\mycard P\geq N\minus f$.\footnote{In general semitopologies, $\noi(P)$ ($P$ having a nonempty open interior) is not the same as $P\in\opensne$ ($P$ being a nonempty open set); but in the case of $\AllBut{N}{f}$, these two properties happen to be equal.}
\item\label{item.allbut.NOI.dense.3}
$\dense(P)$ if and only if $\mycard P\geq f\plus 1$.
\end{enumerate*}
\end{prop}
\begin{proof}
We consider each part in turn:
\begin{enumerate}
\item
Suppose $N>n*f$ and suppose $O_1,\dots,O_n$ are open in $\AllBut{N}{f}$.
From Definition~\ref{defn.allbut} this means that they each contain all but $f$ elements of $\Nset{N}$.
It follows that $O_1\cap\dots \cap O_n$ contains (at least) all but $n*f$ elements of $\Nset{N}$.
Since $N\gneq n*f$, this intersection is nonempty. 

Conversely, if $N\leq n*f$ then it is routine to partition $\Nset{N}$ into $n$ sets, each of which has at most $f$ elements.
Then their complements are a set of $n$ open sets whose intersection is empty.
\item
By routine counting arguments.
\item
By routine counting arguments.
\qedhere\end{enumerate}
\end{proof}



\begin{rmrk}
\label{rmrk.ntwined.and.the.literature}
Being \xxtwined corresponds to the \emph{quorum intersection property} of~\cite[Section~2.1, Condition~B2]{DBLP:journals/tocs/Lamport98}, which requires that any two quorums intersect.
Being \xxxtwined also corresponds to a quorum intersection property with a stronger requirement that (assuming there is a quorum of honest participants) any two quorum must intersect in a participant that is honest~\cite[Definition~5.1]{DBLP:journals/dc/MalkhiR98}. 
  
These intersection properties are closely tied to necessary conditions for some consensus protocols (and other distributed algorithms) to function correctly.
Specifically for the case of the semitopology $\AllBut{N}{f}$ from Definition~\ref{defn.allbut}(\ref{item.allbut}), requiring non-empty quorum intersection is equivalent to the condition \( N > 2f \), where \( f \) denotes the number of processes that may crash during the protocol's execution.
This condition is often referred to as \( Q^2 \)~\cite[Theorem~1]{DBLP:conf/asiacrypt/DamgardDFN07}.
Similarly, being \xxxtwined corresponds to the condition \( N > 3f \), commonly referred to as \( Q^3 \)~\cite[Definition 2]{DBLP:journals/dc/AlposCTZ24}.
A more generalised version of \( Q^3 \) exists, known as \( B^3 \)~\cite[Theorem 2]{DBLP:conf/asiacrypt/DamgardDFN07}\cite[Definition 5]{DBLP:journals/dc/AlposCTZ24}, which accounts for greater heterogeneity. 
Specifically, \( Q^3 \) applies to distributed systems where failures are homogeneous, meaning they are globally known and agreed upon by all participants. 
For example, all participants understand that up to \( f \) processes may fail. 
  
In contrast, in a heterogeneous setting, participants may have differing beliefs about \emph{who} might fail, leading to a more nuanced description of trust assumptions. 
In this context, \( B^3 \) generalises \( Q^3 \) by incorporating these variations in trust and failure assumptions.
Readers interested in exploring this topic further can consult~\cite{DBLP:conf/asiacrypt/DamgardDFN07,DBLP:journals/dc/AlposCTZ24}.
\end{rmrk}
  
  
  
  

\subsection{The canonical properties}

For this subsection we fix a signature $\Sigma$ and a model $\mathcal M=(\Pnt,\opens,\Val)$.

The results in this Subsection express how the underlying semitopology of this model $(\Pnt,\opens)$ being \xxtwined or \xxxtwined manifests itself in the logic, in ways that will be relevant in our later proofs:
\begin{lemm}
\label{lemm.supertwined.qcbt}
Suppose $(\Pnt,\opens)$ is a semitopology and $n\in\Time$ and $\avaluation$ is a valuation and suppose $\phi$, $\psi$, and $\chi$ are pointwise in $(n,\avaluation)$ (Definition~\ref{defn.pointwise.meaning}). 
Then: 
\begin{enumerate*}
\item\label{item.intertwined.qcbt}
If $(\Pnt,\opens)$ is \xxtwined then we have:
\begin{enumerate*}
\item\label{item.intertwined.3}
If $n\mentval \QuorumBox\phi$ then $n\mentval\CoquorumDiamond\phi$. 
\item\label{item.intertwined.qcbt.2}
If $n\mentval \QuorumBox\phi \tand \QuorumBox\psi$ then $n\mentval\someoneAll(\phi\tand\psi)$. 
\item\label{item.intertwined.qcbt.3}
If $n\mentval \QuorumBox\correct{\phi} \tand \QuorumBox\phi$ then $n\mentval\modT\someoneAll\phi$.
\end{enumerate*}
\item\label{item.supertwined.qcbt}
If $(\Pnt,\opens)$ is \xxxtwined then we have: 
\begin{enumerate*}
\item\label{item.supertwined.qcbt.3}
If $n\mentval \QuorumBox\phi \tand \QuorumBox\psi$ then $n\mentval \CoquorumDiamond(\phi\tand\psi)$.
\item\label{item.supertwined.qcbt.2}
If $n\mentval \QuorumBox\phi \tand \QuorumBox\psi \tand \QuorumBox\chi$ then $n\mentval\someoneAll(\phi\tand\psi\tand\chi)$. 
\item\label{item.supertwined.qcbt.4}
If $n\mentval \QuorumBox\correct{\phi} \tand \QuorumBox\phi$ then $n\mentval \modT\CoquorumDiamond\phi$.
\end{enumerate*}
\end{enumerate*}
Recall that $\correct{\phi}$ is from Definition~\ref{defn.correct}(\ref{item.correct.phi}) and just means $\modTF\phi$.
\end{lemm}
\begin{proof}
We reason as follows:
\begin{enumerate}
\item
Suppose $(\Pnt,\opens)$ is \xxtwined. 
\begin{enumerate*}
\item
We reason as follows:
$$
\begin{array}{r@{\ }l@{\qquad}l}
n\mentval\QuorumBox\phi
\liff&
\noi(\{p\in\Pnt \mid n,p,\avaluation\ment\phi\})
&\text{Proposition~\ref{prop.pointwise.dense.char}(\ref{item.dense.char.pointwise.QuorumBox})}
\\
\limp&
\dense(\{p\in\Pnt \mid n,p,\avaluation\ment\phi\}) 
&\text{Proposition~\ref{prop.noi.dense}, $(\Pnt,\opens)$ $2$-twined} %
\\
\liff&
n\mentval\CoquorumDiamond\phi
&\text{Proposition~\ref{prop.pointwise.dense.char}(\ref{item.dense.char.pointwise.CoquorumDiamond})}
\end{array}
$$
\item
Suppose $n\mentval \QuorumBox\phi\tand\QuorumBox\psi$.
By part~\ref{item.intertwined.3} of this result, $n\mentval\CoquorumDiamond\phi$.
By Proposition~\ref{prop.someone.implies.someoneAll}(\ref{item.someone.implies.someoneAll.4}) $n\mentval \someoneAll(\phi\tand\psi)$ as required.
\item
Suppose $n\mentval \QuorumBox\correct{\phi} \tand \QuorumBox\phi$.
By part~\ref{item.intertwined.qcbt.2} of this result, $n\mentval\someoneAll(\correct{\phi}\tand\phi)$.
Using %
Proposition~\ref{prop.tv.ment.TF.model}(\ref{item.tv.ment.TF.phi.lr}) $n\mentval\someoneAll\modT\phi$.
By Lemma~\ref{lemm.someone.commutation} $n\mentval\modT\someoneAll\phi$ as required. 
\end{enumerate*}
\item
The reasoning for the \xxxtwined case is just as for the \xxtwined case, except we need intersection of three open sets instead of two.
\qedhere\end{enumerate} 
\end{proof}




Corollary~\ref{corr.two.quorums.intersect.correct} specialises Lemma~\ref{lemm.supertwined.qcbt} to the useful special case of a correct unary predicate symbol, such that $\phi=\correct{\tf P}$ (Definition~\ref{defn.correct}(\ref{item.correct})) and $\psi=\tf P(v)$ and $\chi=\tf P(v')$.
In our later examples, $\tf P$ will take values such as $\tf{accept}$, $\tf{send}$, $\tf{write}$, and $\tf{decide}$, representing basic actions in a consensus protocol:
\begin{corr}
\label{corr.two.quorums.intersect.correct}
Suppose $(\Pnt,\opens)$ is a semitopology and suppose $\tf P$ is a unary predicate symbol and $v,v'\in\Val$ and $n\in\Time$ and $\avaluation$ is a valuation.
Then:
\begin{enumerate*}
\item
If $(\Pnt,\opens)$ is \xxxtwined then:
\begin{enumerate*}
\item\label{item.two.quorums.intersect.correct.supertwined}
If $n\mentval\QuorumBox\correct{\tf P}\tand\QuorumBox\tf P(v)\tand\QuorumBox\tf P(v')$ then $n\mentval\modT\someoneAll(\tf P(v)\tand\tf P(v'))$.
\item\label{item.QuorumBox.correct.supertwined}
If $n\mentval\QuorumBox\correct{\tf P}\tand\QuorumBox\tf P(v)$ then $n\mentval\modT\CoquorumDiamond\tf P(v)$.
\end{enumerate*}
\item\label{item.two.quorums.intersect.correct.intertwined}
If $(\Pnt,\opens)$ is \xxtwined then
$n\mentval\QuorumBox\correct{\tf P}\tand\QuorumBox\tf P(v)$ implies $n\mentval\modT\someoneAll(\tf P(v))$.
\end{enumerate*}
\end{corr}
\begin{proof}
We consider each part in turn:
\begin{enumerate}
\item[(\ref{item.two.quorums.intersect.correct.supertwined})]
Suppose $(\Pnt,\opens)$ is \xxxtwined and $n\mentval\QuorumBox\correct{\tf P}\tand\QuorumBox\tf P(v)\tand\QuorumBox\tf P(v')$.
By Lemma~\ref{lemm.supertwined.qcbt}(\ref{item.supertwined.qcbt.2}) (since $\tf P(v)$ and $\correct{\tf P}$ are pointwise by Lemma~\ref{lemm.atomic.O}(\ref{item.atomic.pointwise.2})) $n\mentval\someoneAll(\correct{\tf P}\tand\tf P(v)\tand\tf P(v'))$.
By two applications of Proposition~\ref{prop.tv.ment.TF.model}(\ref{item.tv.ment.TF.P}) and Lemma~\ref{lemm.commuting.connectives} (for $\tand$)
$n\mentval\someoneAll\modT(\tf P(v)\tand\tf P(v'))$.
By Lemma~\ref{lemm.someone.commutation} $n\mentval\modT\someoneAll(\tf P(v)\tand\tf P(v'))$ as required.
\item[(\ref{item.QuorumBox.correct.supertwined})] 
Suppose $(\Pnt,\opens)$ is \xxxtwined and $n\mentval\QuorumBox\correct{\tf P}\tand\QuorumBox\tf P(v)$.
Using Lemma~\ref{lemm.correct.P.Pv} it follows that $n\mentval\QuorumBox\correct{\tf P(v)}\tand\QuorumBox\tf P(v)$.
We just apply Lemma~\ref{lemm.supertwined.qcbt}(\ref{item.supertwined.qcbt.4}). 
\item[(\ref{item.two.quorums.intersect.correct.intertwined})]
This follows from Lemma~\ref{lemm.supertwined.qcbt}(\ref{item.intertwined.qcbt.3}) just as part~\ref{item.QuorumBox.correct.supertwined} does from Lemma~\ref{lemm.supertwined.qcbt}(\ref{item.supertwined.qcbt.2}).
\qedhere\end{enumerate}
\end{proof}


Lemma~\ref{lemm.allbut.concretely} unpacks concretely what $\QuorumBox\phi$ and $\CoquorumDiamond\phi$ --- which are the subject of many results including of Lemmas~\ref{lemm.supertwined.qcbt} and~\ref{lemm.intertwined.char} --- mean in the semitopology $\AllBut{N}{f}$ from Definition~\ref{defn.allbut}(\ref{item.allbut}), which has $N$ points, and has open sets being either the empty set, or any set with at least $N\minus f$ elements: 
\begin{enumerate*}
\item
$\QuorumBox\phi$ means `for at least $N\minus f$ elements'; and 
\item
$\CoquorumDiamond\phi$ means `for at least $f\plus 1$ elements.
\end{enumerate*}
This special case is particularly relevant because (as the reader may know) $\AllBut{N}{f}$ is commonly used in the literature on distributed algorithms. 
\begin{lemm}
\label{lemm.allbut.concretely}
Suppose that
\begin{itemize*}
\item
$N,f\in\Ngeqz$, 
\item
$(\Pnt,\opens)=\AllBut{N}{f}$, 
\item
$n\in\Time$ and $\avaluation$ is a valuation, and
\item
$\phi$ is a closed predicate that is pointwise in $(n,\avaluation)$.
\end{itemize*}
Then:
\begin{enumerate*}
\item
$n\mentval\QuorumBox\phi$ if and only if 
$\mycard\{p\in\Pnt \mid n,p\mentval\phi\}\geq N\minus f$. 
\item
$n\mentval\CoquorumDiamond\phi$ if and only if 
$\mycard\{p\in\Pnt \mid n,p\mentval\phi\}\geq f\plus 1$. 
\item
$n\mentval\modT\QuorumBox\phi$ if and only if 
$\mycard\{p\in\Pnt \mid n,p\mentval\modT\phi\}\geq N\minus f$. 
\item
$n\mentval\modT\CoquorumDiamond\phi$ if and only if 
$\mycard\{p\in\Pnt \mid n,p\mentval\modT\phi\}\geq f\plus 1$. 
\end{enumerate*}
\end{lemm}
\begin{proof}
We consider each part in turn:
\begin{enumerate}
\item
The proof is elementary but we spell it out.
We reason as follows:
$$
\begin{array}{r@{\ }l@{\quad}l}
n\mentval\QuorumBox\phi
\liff&
\noi(\{p\in\Pnt \mid n,p\mentval \modT\phi\})
&\text{Proposition~\ref{prop.pointwise.dense.char}(\ref{item.dense.char.pointwise.QuorumBox})}
\\
\liff&
\mycard\{p\in\Pnt \mid n,p\mentval \phi\}\geq N\minus f
&\text{Proposition~\ref{prop.allbut.supertwined}(\ref{item.allbut.NOI.dense})}
\end{array}
$$
\item
Much as the previous part, but using Proposition~\ref{prop.pointwise.dense.char}(\ref{item.dense.char.pointwise.CoquorumDiamond}).
\item
By Lemma~\ref{lemm.someone.commutation}, $n\mentval\modT\QuorumBox\phi$ if and only if $n\mentval\QuorumBox\modT\phi$.
We just use part~1 of this result.
\item
Much as the previous part, using part~2 of this result.
\qedhere\end{enumerate}
\end{proof}

\subsection{Further observations on intersection properties}

For this Subsection we fix a signature $\Sigma$ and a model $\mathcal M=(\Pnt,\opens,\Val)$.
\begin{rmrk}
We collect some further results as follows:
\begin{enumerate*}
\item
Lemma~\ref{lemm.intertwined.char} proves converses to the implications in Lemma~\ref{lemm.supertwined.qcbt}, provided that the signature is nonempty.

This is interesting because it tells us that, provided the signature is nondegenerate, Lemma~\ref{lemm.supertwined.qcbt} accurately captures and reflects the intersection properties that interest us. 
\item
Proposition~\ref{prop.reformulate} is a semitopological distillation of parts~\ref{item.intertwined.qcbt.2} and~\ref{item.supertwined.qcbt.2} of Lemma~\ref{lemm.supertwined.qcbt}.
Proposition~\ref{prop.reformulate} \emph{is} Lemma~\ref{lemm.supertwined.qcbt} and Corollary~\ref{corr.two.quorums.intersect.correct}, in semitopological rather than logical form.
\end{enumerate*}
\end{rmrk}

In Lemma~\ref{lemm.intertwined.char} we could replace every instance of the expression `$\modTF\tf P$' with the (by Definition~\ref{defn.correct}(\ref{item.correct.phi}) and Remark~\ref{rmrk.tf.synonyms}) synonymous expression `$\correct{\tf P}$'.
However, we have not set up a corresponding synonym for $\modTB\tf P$ (it would be `$\tf{valid}[\tf P]$'; we just never wrote that definition) so it seems natural to stay with $\modTF$ here.
\begin{lemm}
\label{lemm.intertwined.char}
Suppose $\tf P\in\ns{PredSymb}$ is a predicate constant symbol (Notation~\ref{nttn.logic.terminology}(\ref{item.nullary.ps}); a nullary predicate symbol).\footnote{The proof works just as well for an $n$-ary predicate symbol; we just pick \emph{any} $\arity(\tf P)$-tuple of values $\vec v\in\Pnt^{\arity(\tf P)}$ --- which is possible because we assumed $\Val$ is nonempty in Definition~\ref{defn.model}(\ref{item.model.val}) --- and we use $\tf P(\vec v)$ in the proof instead.}
Then:
\begin{enumerate*}
\item
If\ \  $\ment(\QuorumBox\modTF\tf P\tand\QuorumBox\modTB\tf P)\timpc \someoneAll\tf P$ then $(\Pnt,\opens)$ is \xxtwined. 
\\
(This is a form of converse to Lemma~\ref{lemm.supertwined.qcbt}(\ref{item.intertwined.qcbt.3}).)
\item
If\ \  $\ment(\QuorumBox\modTF\tf P\tand\QuorumBox\modTB\tf P)\timpc \CoquorumDiamond\tf P$ then $(\Pnt,\opensne)$ is \xxxtwined. 
\\
(This is a form of converse to Lemma~\ref{lemm.supertwined.qcbt}(\ref{item.supertwined.qcbt.4}).) 
\item\label{item.intertwined.char.reverse}
As a corollary, provided that the signature $\Sigma$ is has at least one predicate symbol --- i.e. provided $\PredSymb$ in
it is nonempty (Definition~\ref{defn.signature}(\ref{item.predicate.syntax.empty.signature})) --- then reverse implications to those in Lemma~\ref{lemm.supertwined.qcbt} hold.\footnote{No need to also insist that $\Val$ is nonempty, because we assume this in Definition~\ref{defn.model}(\ref{item.model.val}).}
\end{enumerate*}
Recall from Definition~\ref{defn.validity.judgement}(\ref{item.ment.phi}) that $\ment\phi$ means $\Forall{\acontext\in\tf{Ctx}(\Sigma,\mathcal M)}(\acontext\ment\phi)$.
\end{lemm}
\begin{proof}
We consider each part in turn:
\begin{enumerate}
\item
We prove a contrapositive of \strongmodusponens: we suppose $\opens$ is \emph{not} \xxtwined, so there exist $O,O'\in\opensne$ with $O\cap O'=\varnothing$, and we exhibit a valuation such that $\neg(\mentval\someone\tf P)$ (equivalently: $\mentval\modF\someone\tf P$).

We choose a valuation $\avaluation$ such that for any $n\in\Time$, 
\begin{itemize*}
\item
$\avaluation(\tf P)_{n,p}(\aval)=\tvF$ for $p\in O$, and 
\item
$\avaluation(\tf P)_{n,p}(\aval)=\tvB$ for $p\in \Pnt\setminus O$.
\end{itemize*}
The reader can check that $O\mentval \everyone\modTF\tf P$ and $O'\mentval \everyone\modTB P(\aval)$, but $\mentval\modF\someone\tf P$.
\item
We prove a contrapositive: suppose $\opens$ is \emph{not} \xxxtwined, so there exist $O,O',O''\in\opensne$ with $O\cap O'\cap O''=\varnothing$.

We choose a valuation $\avaluation$ such that for any $n\in\Time$, 
\begin{itemize*}
\item
$\avaluation(\tf P)_{n,p}(\aval)=\tvT$ for $p\in O$, and 
\item
$\avaluation(\tf P)_{n,p}(\aval)=\tvB$ for $p\in O'\setminus O$, and 
\item
$\avaluation(\tf P)_{n,p}(\aval)=\tvF$ for $p\in \Pnt\setminus (O\cup O')$.
\end{itemize*}
The reader can check that $O\cup O'\mentval \everyone\modTF\tf P$ and $O\cup O''\mentval \everyone\modTB P(\aval)$, but $O'\cup O''\mentval\modF\everyone\tf P$.
\item
Just by the arguments above.
\qedhere\end{enumerate}
\end{proof}






\begin{nttn}
\label{nttn.cover}
Suppose $\ns X$ is a set.
Then:
\begin{enumerate*}
\item
A \deffont{cover} of $\ns X$ is a set of sets $\mathcal X\subseteq\powerset(\ns X)$ such that $\bigcup\mathcal X=\ns X$.
\item
A \deffont{partition} of $\ns X$ is a cover of $\ns X$ all of whose elements are pairwise disjoint (so $\Forall{X,X'\in\mathcal X}X\cap X'=\varnothing\lor X=X'$).
\end{enumerate*}
\end{nttn}

The reader might like to try giving a proof of Proposition~\ref{prop.reformulate} using direct calculations (without building on the results we have so far).
This is possible, but it is messy.
The use of logic endows the proofs with a structure and cleanness that the direct proof cannot match: 
\begin{prop}
\label{prop.reformulate}
Suppose $(\Pnt,\opens)$ is a semitopology and $n\in\Time$ and $\avaluation$ is a valuation, and suppose $\phi$ is a closed predicate that is pointwise in $(n,\avaluation)$.
Then:
\begin{enumerate*}
\item\label{item.reformulate.1}
If $(\Pnt,\opens)$ is \xxtwined then $n\mentval\CoquorumDiamond\modTB\phi \tor \CoquorumDiamond\modF\phi$. 
\item\label{item.reformulate.1b}
As a corollary, if $(\Pnt,\opens)$ is \xxtwined then any two-element partition of $\Pnt$ (Notation~\ref{nttn.cover}) contains a dense element (Definition~\ref{defn.dense.subset}(\ref{item.dense.subset})).
\item\label{item.reformulate.2}
If $(\Pnt,\opens)$ is \xxxtwined then 
$n\mentval\CoquorumDiamond\modT\phi \tor \CoquorumDiamond\modB\phi \tor \CoquorumDiamond\modF\phi$. 
\item\label{item.reformulate.2b}
As a corollary, if $(\Pnt,\opens)$ is \xxxtwined then any three-element partition of $\Pnt$ contains a dense element.
\end{enumerate*} 
\end{prop}
\begin{proof}
We consider each part in turn:
\begin{enumerate}
\item
By Lemma~\ref{lemm.supertwined.qcbt}(\ref{item.intertwined.qcbt.2})
$n\mentval\QuorumBox\modF\phi\tand\QuorumBox\modTB\phi$ implies $n\mentval\someoneAll(\modF\phi\tand\modTB\phi)$.
Unpacking what $n\mentval\someoneAll(\modF\phi\tand\modTB\phi)$ means using Definition~\ref{defn.validity.judgement} and Figures~\ref{fig.3.derived} and~\ref{fig.3.phi.f}, we see that this is impossible --- for no $p$ and $O$ can we satisfy both $n,p,O\mentval\modF\phi$ and $n,p,O\mentval\modTB\phi$ together --- so that 
$n\nmentval\QuorumBox\modF\phi\tand\QuorumBox\modTB\phi$.
By Lemma~\ref{lemm.para}(\ref{item.para.em})
$n\mentval\tneg(\QuorumBox\modF\phi\tand\QuorumBox\modTB\phi)$, and simplifying further using de Morgan equivalences as per Figure~\ref{fig.easy.equivalences}, we obtain
$n\mentval\CoquorumDiamond\modTB\phi\tor\CoquorumDiamond\modF\phi$ as required.
\item
Suppose $\{P_1,P_2\}$ partitions $\Pnt$.
Assume a $0$-ary predicate symbol $\tf P$, and let $\varsigma$ be the valuation that sets $\varsigma(\tf P)_{n,p}=\tvT$ if $p\in P_1$ and $\varsigma(\tf P)_{n,p}=\tvF$ if $p\in P_2$.\footnote{We can assume $\tf P$ and $\varsigma$ because this result is about semitopologies, so we are free to impose any model structure on it that we like.  It seems particularly nice here how the mere \emph{existence} of a model for our logic gives us a clean proof of a purely semitopological property which, on the face of it, has nothing to do with our three-valued modal logic.}
By part~\ref{item.reformulate.1} and Lemma~\ref{lemm.tand.tor} either $n\mentval \Coquorum\modTB\tf P$ or $n\mentval\Coquorum\modF\tf P$.
By Proposition~\ref{prop.pointwise.dense.char}(\ref{item.dense.char.pointwise.CoquorumDiamond}) either $P_1$ or $P_2$ is dense.
\item
We reason as for the proof of part~\ref{item.reformulate.1}, using Lemma~\ref{lemm.supertwined.qcbt}(\ref{item.supertwined.qcbt.2}) and starting from $\modFB\phi \tand \modTF\phi \tand \modTB\phi$. 
\item
We reason as for the proof of part~\ref{item.reformulate.2}, for a partition $\{P_1,P_2,P_3\}$ and setting $\varsigma(\tf P)_{n,p}=\tvT$ if $p\in P_1$ and $\varsigma(\tf P)_{n,p}=\tvB$ if $p\in P_2$ and $\varsigma(\tf P)_{n,p}=\tvF$ if $p\in P_3$.
\qedhere\end{enumerate}
\end{proof}




 





\section{Declarative Paxos}
\label{sect.declarative.paxos}

\subsection{Case study: a simple protocol}
\label{subsect.simple}

\begin{rmrk}
\label{rmrk.simple.protocol}
Before Paxos, we will warm up by sketching a simple protocol \theory{Simple}, and then we will give four \QLogic axiomatisations of \theory{Simple}, reflecting four different failure assumptions (messages can/cannot get lost; participants can/cannot crash).
\theory{Simple} is as follows:
\begin{enumerate*}
\item
Any participant $p$ may broadcast a $\tf{propose}$ message to all participants.
\item
If a participant receives a $\tf{propose}$ message, then it responds with $\tf{accept}$. 
\item
If the participant $p$, having broadcast its $\tf{propose}$ message, receives a $\tf{accept}$ message, then it declares $\tf{decide}$.
\end{enumerate*}
\end{rmrk}


\begin{rmrk}
\theory{Simple} is very simple!
Only one value can be proposed, so we do not even bother to state it; we just `propose'.
Any participant accepts immediately that it hears someone propose. 
Any participant that proposes, decides immediately that it hears that someone accepted.

For comparison, in Paxos \emph{multiple values} are available, so we need to make sure that participants all decide on the \emph{same} value if they decide on anything at all. 
This makes things much harder: so Paxos has multiple rounds; there is a designated leader participant at each round; the leader must choose a value to propose; and elaborate communication is in place to ensure that (subject to suitable precisely-stated assumptions) a single choice of agreed value will eventually be decided.

But that is for later.  
Even the simple protocol \theory{Simple} has plenty of subtlety, which will allow us to exhibit (in `toy' form) many of the techniques that we will then apply to Paxos:
\end{rmrk}


\begin{rmrk}
We will axiomatise \theory{Simple} subject to two parameters, yielding four possible axiomatisations:
\begin{enumerate*}
\item
Do all messages arrive, or might there be a finite initial period during which the network may be unreliable so that during that time, messages might not arrive?
\item
Can participants crash (and if so, how many)?
\end{enumerate*}
\end{rmrk}

\subsubsection*{All messages arrive; no participants crash}

First, we assume an ideal situation that all messages arrive, and no participants crash.
We can axiomatise this as follows:
$$
\begin{array}{l@{\qquad}l}
\rulefont{SimpAccept?}&
\tf{accept} \timpc \someoneAll\tf{propose} 
\\
\rulefont{SimpDecide?}&
\tf{decide} \timpc (\tf{propose}\tand\someoneAll\tf{accept}) 
\\
\rulefont{SimpAccept!}&
(\someoneAll\tf{propose}) \timpc \tf{accept}
\\
\rulefont{SimpDecide!}&
(\tf{propose}\tand\someoneAll\tf{accept}) \timpc \tf{decide}
\\
\rulefont{Correct}&
\correct{\tf{propose},\tf{accept},\tf{decide}}
\end{array}
$$
\noindent Recall from Definition~\ref{defn.three.imp} and Proposition~\ref{prop.mp.for.tnotor} that $\timpc$ is a \emph{strong implication} and means `if the LHS is $\tvT$ then so is the RHS', and 
recall from Proposition~\ref{prop.pointwise.dense.char}(\ref{item.dense.char.pointwise.someoneAll}) that $\someoneAll\tf P$ means `$\tf P$ somewhere', where $\tf P\in\{\tf{propose},\tf{accept},\tf{decide}\}$.

Then we explain the axioms (briefly) as follows:
\begin{itemize*}
\item
We call \rulefont{SimpAccept?} and \rulefont{SimpDecide?} \emph{backward rules}; they express backward versions of the rules in Remark~\ref{rmrk.simple.protocol}, namely that:
\begin{itemize*}
\item
if a participant asserts $\tf{accept}$ then some participant must have asserted $\tf{propose}$; and 
\item
if a participant asserts $\tf{decide}$ then it must have asserted $\tf{propose}$, and some participant must have asserted $\tf{accept}$.
\end{itemize*}
\item
We call \rulefont{SimpAccept!} and \rulefont{SimpDecide!} \emph{forward rules}.
They express forward versions of the rules in Remark~\ref{rmrk.simple.protocol}.
Because our failure assumptions are that nothing fails --- all messages arrive; no participants crash --- the forward rules are in this case just the backward rules, reversed.
\item
Axiom \rulefont{Correct} expresses that all predicates return $\tvT$ and $\tvF$ (not $\tvB$), as per Definition~\ref{defn.correct}(\ref{item.correct}). 
This reflects our assumption that no participants crash, because we model a computation crashing %
as the predicate describing it returning $\tvB$.
\end{itemize*}
Note that $\tf{propose}$, $\tf{accept}$, and $\tf{decide}$ above are not messages; they are propositions!
There is no message-passing.
$\tf{accept}$ being true at a participant represents that the participant accepts; this is implemented in the algorithm by the participant broadcasting an accept message, and represented in the logic by $\tf{accept}$ taking truth-value $\tvT$ at that participant.

Our example protocol is simple, but we hope it is already starting to illustrate how the various parts of our logic come together to express the protocol axiomatically.
We continue:

\subsubsection*{All messages arrive; participants can crash}

We assume that all messages arrive, and participants can crash --- but we assume there is a dense set of uncrashed participants.
Recall from Definition~\ref{defn.dense.subset}(\ref{item.dense.subset}) that a \emph{dense set} is one that intersects every open set; in particular, this means that every open set of participants contains a participant that is not crashed.
We model being crashed as returning the third truth-value $\tvB$.
We can axiomatise this as follows:
$$
\begin{array}{l@{\qquad}l}
\rulefont{SimpAccept?}&
\tf{accept} \timpc \someoneAll\tf{propose} 
\\
\rulefont{SimpDecide?}&
\tf{decide} \timpc (\tf{propose}\tand\someoneAll\tf{accept}) 
\\
\rulefont{SimpAccept!}&
(\someoneAll\tf{propose}) \tnotor \tf{accept}
\\
\rulefont{SimpDecide!}&
(\tf{propose}\tand\someoneAll\tf{accept}) \tnotor \tf{decide}
\\
\rulefont{Uncrashed}&
\QuorumBox\modTB\tf P \tnotor \someoneAll\modT\tf P
\quad \tf P\in\{\tf{propose},\tf{accept},\tf{decide}\}
\end{array}
$$ 
\noindent 
The backward rules are unchanged: if a participant accepts then from the form of the rules, it \emph{must} be the case that somebody proposed; and if a participant decides then from the form of the rules, it \emph{must} be the case that the participant proposed and some participant accepted.
 
The forward rules \rulefont{SimpAccept!} and \rulefont{SimpDecide!} have changed, and they now use the weak implication $\tnotor$.
By Definition~\ref{defn.three.imp} and Proposition~\ref{prop.mp.for.tnotor}, $\tnotor$ means `if the LHS is $\tvT$ then the RHS is $\tvT$ or $\tvB$'.
The $\tvB$ truth-value represents `crashed', so that (for example) \rulefont{SimpAccept!} expresses that for each participant $p$, if some (possibly other) participant $q$ asserts $\tf{propose}$, then $p$ asserts $\tf{accept}$ \emph{or} $p$ has crashed. 

Axiom \rulefont{Uncrashed} reflects that any a quorum (open set) of participants contains an uncrashed participant; e.g. in topological language, that the set of uncrashed participants is dense.
More precisely: \rulefont{Uncrashed} asserts that if an open set of participants returns $\tvT$ or $\tvB$ for $\tf P$, then at least one of them (is not crashed and so) returns $\tvT$.

Note from Lemma~\ref{lemm.someone.commutation} that $\Quorum\modTB$ is equivalent to $\modTB\Quorum$, and $\someoneAll\modT$ is equivalent to $\modT\someoneAll$, so that an equivalent form of \rulefont{Uncrashed} is $\modTB\QuorumBox\tf P \tnotor \modT\someoneAll\tf P$.
In our axiomatisations it is usually helpful to pull $\modTB$ and $\modT$ up to the topmost level possible in the predicate, so we may favour that form, but this is just a matter of convenience; meaning is unchanged. 


\subsubsection*{Messages might not arrive; no participants crash}

We assume that there is an initial period during which the network might not deliver messages, but the network eventually stabilises (in the literature this is called \emph{eventual synchrony}).
No participants crash.
We can axiomatise this as follows:
$$
\begin{array}{l@{\qquad}l}
\rulefont{SimpAccept?}&
\tf{accept} \timpc \someoneAll\tf{propose} 
\\
\rulefont{SimpDecide?}&
\tf{decide} \timpc (\tf{propose}\tand\QuorumBox\tf{accept}) 
\\
\rulefont{SimpAccept!}&
\final((\someoneAll\tf{propose}) \timpc \tf{accept})
\\
\rulefont{SimpDecide!}&
\final((\tf{propose}\tand\someoneAll\tf{accept}) \timpc \tf{decide})
\\
\rulefont{Correct}&
\correct{\tf{propose},\tf{accept},\tf{decide}}
\end{array}
$$ 
\noindent The backward rules are unchanged, but the forward rules \rulefont{SimpAccept!} and \rulefont{SimpDecide!} now have a $\final$ modality.
This is from Figure~\ref{fig.3.derived} and it means `eventually, this holds forever', reflecting an assumption that eventually the network stabilises.

Note that \QLogic does not explicitly represent message-passing, but this does not matter as much as one might suppose!
The $\final$ modality captures the essence of what we need, namely, that eventually the forward rules become valid. 

\subsubsection*{Messages might not arrive; participants can crash}

We assume that messages might not arrive, and participants can crash but there is a dense set of uncrashed participants (so that every quorum of participants contains at least one uncrashed participant).
We can axiomatise this just by combining the ideas above, as follows:
$$
\begin{array}{l@{\qquad}l}
\rulefont{SimpAccept?}&
\tf{accept} \timpc \someoneAll\tf{propose} 
\\
\rulefont{SimpDecide?}&
\tf{decide} \timpc (\tf{propose}\tand\someoneAll\tf{accept}) 
\\
\rulefont{SimpAccept!}&
\final((\someoneAll\tf{propose}) \tnotor \tf{accept})
\\
\rulefont{SimpDecide!}&
\final((\tf{propose}\tand\someoneAll\tf{accept}) \tnotor \tf{decide})
\\
\rulefont{Uncrashed}&
\modTB\QuorumBox\tf P \tnotor \modT\someoneAll\tf P
\quad \tf P\in\{\tf{propose},\tf{accept},\tf{decide}\}
\end{array}
$$ 
\noindent The axiomatisations of Declarative Paxos and Simpler Declarative Paxos in Figures~\ref{fig.ThyPaxOne} and~\ref{fig.ThySPax} are more elaborate than the axioms above, because Paxos is a more complicated protocol, but the basic ideas are as above; just scaled up. 


\subsection{A high-level view of the Paxos algorithm}



\begin{rmrk}[High-level view of Paxos]
\label{rmrk.high-level.paxos}
Paxos allows participants to crash, and messages to fail to arrive at their destination, but it assumes all participants are honest (i.e. that if they have not crashed then they follow the rules of the algorithm).
The reader can find detailed treatments of Paxos in \cite{DBLP:journals/tocs/Lamport98,lamport2001paxos,cachin2011paxos} (a Wikipedia article~\cite{wiki:Paxos_(computer_science)} is also accessible and clear).
Here is a high-level view of the algorithm:
\begin{enumerate*}
\item\label{item.l.proposes}
\deffont{Leader proposes:}\quad

A designated leader participant (if it does not crash) sends messages to all participants, \emph{proposing} some unique value $v$.

Message-passing is unreliable, so messages may fail to arrive.
\item\label{item.p.sends}
\deffont{Participant sends:}\quad

If a participant $p$ receives a propose $v$ message from the leader --- it is assumed that every participant knows who the leader is, and other participants cannot forge the leader's signature --- then $p$ sends a \emph{send} message back to the leader, which carries the value $v_p$ that $p$ most recently accepted, along with the timestamp of when $v_p$ was accepted.\footnote{For reasons of compact presentation we do not explicitly represent the timestamps in our axiomatisation; i.e. $\tf{send}$ is a unary predicate in Definition~\ref{defn.declarative.paxos}(\ref{item.ThyPaxOne.sig}), not a binary predicate.
If this bothers the reader then we can just say that, actually, $v_p$ contains a timestamp of the stage $n\in\Time$ it was accepted; i.e. that $v_p$ is actually implemented as a pair $(v_p',n_p)$ of a value and a time.} 
If $p$ has never accepted any value, so no most recently accepted value exists, then $p$ sends back a special $\udfn$ value.
\item\label{item.l.writes}
\deffont{Leader writes:}\quad

The leader participant sends a message to all participants, \emph{writing} the most recently-timestamped value $v_p$ mentioned in any \emph{send} message it received from within the first quorum of participants that it hears back from; if the leader does not hear back from any quorum, then it does not write anything.
\item\label{item.p.accepts}
\deffont{Participant accepts:}\quad

If a participant receives a write message from the leader for some value, then it sends a message back \emph{accepting} that value.
\item\label{item.l.decides}
\deffont{Leader decides:}\quad

If the leader receives accept messages for some value from a quorum of participants, it decides on that value, and sends out a \emph{decide} message to all participants with that value.
\item\label{item.p.decides}
\deffont{Participant (non-leader) decides:}\quad

If a participant receives a decide message from the leader for some value, then it decides on that value and stops.
\end{enumerate*}
\end{rmrk}

\begin{rmrk}[High-level view of our axiomatisation]
\label{rmrk.elisions}
The theory $\ThyPax$ in Definition~\ref{defn.ThyPaxOne} and Figure~\ref{fig.ThyPaxOne} is a declarative abstraction of Paxos-the-algorithm.
It stands in relation to implementations of Paxos in roughly the same way that (say) the mathematical specification of the Fibonacci function stands in relation to its concrete implementations, in the sense that we elide implementational detail to preserve (and reveal) a mathematical ideal of what the implementation does.

\ThyPax elides two features of the concrete implementation:
\begin{enumerate*}
\item
There is no explicit message-passing; there are only predicates returning truth-values.

This may seem counterintuitive --- isn't Paxos \emph{supposed} to be a message-passing algorithm? --- but we shall see that a lot of the character of Paxos survives.
Indeed, arguably (and certainly for the first author) it is much easier to see what is really going on once we have done this.
\item\label{item.elide.algorithmic.time}
There is no abstract machine, and therefore no \deffont{algorithmic time}, by which we mean no sequence of state-changes of the abstract machine. 

The only notion of time is the \emph{stage} $n\in\Time$ from the context as defined in Definition~\ref{defn.contexts}(\ref{item.context.stage}).
We may call this \deffont{logical time} and (when we are modelling Paxos) it represents the sequence of what is in the distributed consensus algorithms literature often called \emph{stages}, or \emph{ballot numbers}.\footnote{In a distributed algorithm, at any point in algorithmic time, different participants in the algorithm may be at different points in logical time.  A practical example: consider several participants racing across 100 metres. In logical time, they are both just counting to 100.  In algorithmic time, one of them will get to 100 \emph{first}, and win.  In this paper, we only care about the logical aspects of the process; there may still be a winner participant, but this would implemented just as a unary predicate $\tf{winner}$ such that $100\ment\texiaffine\tf{winner}$.} 
\end{enumerate*}
We will prove standard Paxos correctness properties for \ThyPax in Section~\ref{sect.paxos.correctness.properties}. 
\end{rmrk}


\subsection{Definition of Declarative Paxos}

\begin{figure}
$$
\begin{array}{l@{\quad}r@{\ }l}
\figunderline{Backward rules, for backward inference} 
\figskip
\rulefont{PaxPropose?}&
& \tall v.\ \tf{propose}(v)\timpc(\tf{leader}\tand v\tneq\udfn)
\figskip
\rulefont{PaxSend?}&
& \tall v.\ \tf{send}(v) \timpc 
\begin{array}[t]{l}
\someoneAll\texi\tf{propose}\ \tand 
\\
\bigl(
(\mru{\tf{accept}}{v})\ \tor  
(\tneg\modT\recent\texi\tf{accept}\tand v\teq\udfn)\bigr)
\end{array}
\figskip
\rulefont{PaxWrite?}&
& \tall v.\ \tf{write}(v) \timpc 
\begin{array}[t]{l}
\tf{leader}\tand 
\\
\Quorum(\Box\texi\tf{send} \tand 
\\
\qquad\bigl((\mru{(\someone\tf{accept})}{v}) \tor
(\tf{propose}(v)\tand\tneg\modT\recent\texi\someone\tf{accept})\bigr)) 
\end{array}
\figskip
\rulefont{PaxAccept?}&
& \tall v.\ \tf{accept}(v) \timpc \someoneAll\tf{write}(v)
\figskip
\rulefont{PaxDecideL?}&
& \tf{leader} \timpc \tf{decide}(v) \timpc \QuorumBox\tf{accept}(v)
\figskip
\rulefont{PaxDecide\tneg L?}&
& \tall v.\ \tneg\tf{leader} \timpc \tf{decide}(v) \timpc \someoneAll(\tf{leader}\tand\tf{decide}(v))
\figskip\figskip
\figunderline{Forward rules, for making progress} 
\figskip
\rulefont{PaxPropose!}&
& \tf{leader}\tnotor\texi\tf{propose}
\figskip
\rulefont{PaxSend!}&
& \final(\someoneAll\texi\tf{propose} \tnotor \texi\tf{send}) 
\figskip
\rulefont{PaxWrite!}&
& \final\bigl((\tf{leader}\tand (\QuorumBox\texi\tf{send}) \tand \texi\tf{propose}) \tnotor \texi\tf{write}\bigr)
\figskip
\rulefont{PaxAccept!}&
& \final(\texi\someoneAll\tf{write}\tnotor \texi\tf{accept}) 
\figskip
\rulefont{PaxDecideL!}&
& \final((\tf{leader} \tand\texi\QuorumBox\tf{accept}) \tnotor \texi\tf{decide}) 
\figskip
\rulefont{PaxDecide\tneg L!}&
& \final((\tneg\tf{leader} \tand \someoneAll(\tf{leader}\tand\texi\tf{decide}))\tnotor \texi\tf{decide}) 
\figskip
\figskip
\figunderline{Other rules, for correctness \& liveness}
\figskip
\rulefont{LdrExist}&
&\modTF\tf{leader}\tand\modT\someoneAll\tf{leader}
\figskip
\rulefont{LdrExt}_\phi&
& \tf{leader}\timpc\someoneAll(\tf{leader}\tand\phi)\timpc\phi
\quad\text{every closed predicate $\phi$}
\figskip
\rulefont{LdrCorrect}&
& \infinitely(\tf{leader}\timpc\correct{\tf{propose},\tf{write},\tf{decide}}) 
\figskip
\rulefont{PaxPCorrect}&
& 
\begin{array}[t]{l}
\QuorumBox\correct{\tf{propose}} \tand \QuorumBox\correct{\tf{send}} \tand \QuorumBox\correct{\tf{accept}} \tand 
\\
\qquad\QuorumBox\correct{\tf{write}} \tand \QuorumBox\correct{\tf{decide}}
\end{array}
\figskip
\rulefont{PaxPropose01}&
& \texiaffine \tf{propose}
\figskip
\rulefont{PaxWrite01}&
& \texiaffine\tf{write}
\figskip
\rulefont{Pax2Twined}&
& \tall v.\ \QuorumBox\tf{accept}(v) \timpc \CoquorumDiamond\tf{accept}(v)
\end{array}
$$
\caption{$\Theta_\ThyPax$: Axioms of Declarative Paxos (Definition~\ref{defn.declarative.paxos})}
\label{fig.logical.paxos}
\label{fig.ThyPaxOne}
\end{figure}


Recall from Definition~\ref{defn.theory.axiom}(\ref{item.theory}) that a \emph{theory} is a pair of a \emph{signature} and a set of \emph{axioms} over that signature.
Then:
\begin{defn}
\label{defn.declarative.paxos}
\label{defn.ThyPaxOne}
We define \deffont{Declarative Paxos} to be the logical theory 
$$
\ThyPax=(\Sigma_{\ThyPax},\Theta_{\ThyPax}),
$$ 
where:
\begin{enumerate*}
\item\label{item.ThyPaxOne.sig}
the signature $\Sigma_{\ThyPax}$ is 
$$
\Sigma_{\ThyPax}=[\tf{leader}:0;\ \tf{propose},\tf{send},\tf{write},\tf{accept},\tf{decide}:1] ,
$$ 
as per the notation in Definition~\ref{defn.signature}(\ref{item.predicate.syntax.signature}), and
\item\label{item.ThyPaxOne.axioms}
the axioms $\Theta_{\ThyPax}$ are as written in Figure~\ref{fig.logical.paxos}.
\end{enumerate*}
In axioms \rulefont{PaxPropose?} and \rulefont{PaxSend?} we assume an element $\udfn\in\Val$, to model the `undefined' value from Paxos (it does not matter what $\udfn\in\Val$ is; we just have to choose some element and use it consistently).
\end{defn}

\begin{rmrk}
\label{rmrk.functional.notation.in.paxos}
For convenient reference, we give pointers to help unpack the notation used in Figure~\ref{fig.logical.paxos}.
Note that this is just an overview focussing on items of notation used in the axioms, to help with reading them; in-depth discussions of the axioms will follow: 
\begin{enumerate*}
\item
We use the functional notation from Notation~\ref{nttn.hos} to elide quantified variables where convenient.
For example, axioms \rulefont{PaxPropose!} and \rulefont{PaxAccept!} written in full are 
$$
\tf{leader}\tnotor\texi\thea.\tf{propose}(\thea)
\quad\text{and}\quad
\final((\texi\thea.\someoneAll\tf{write}(\thea))\tnotor\texi\thea.\tf{accept}(\thea)) .
$$
\item
$\tnotor$ and $\timpc$ are the \emph{weak} and \emph{strong} implications respectively.
Truth-tables are in Figure~\ref{fig.3} and a discussion is in Subsection~\ref{subsect.3.imp}.
In particular, by Proposition~\ref{prop.mp.for.tnotor}, 
\begin{itemize*}
\item
$\phi\tnotor\phi'$ means `if $\phi$ is true then $\phi'$ is valid' (where `valid' means `$\tvT$ or $\tvB$') and 
\item
$\phi\timpc\phi'$ means `if $\phi$ is true then $\phi'$ is true'. 
\end{itemize*}
\item
The compound modalities --- $\someoneAll$, $\everyoneAll$, $\QuorumBox$, and $\CoquorumDiamond$ --- are discussed in Subsection~\ref{subsect.compound.modalities}.
Their meaning in full generality is in Lemma~\ref{lemm.dense.char}, but because the axioms only apply them to \emph{pointwise} predicates, their meanings simplify to `somewhere' ($\someoneAll$), `everywhere' ($\everyoneAll$), `for a set with a nonempty open interior' ($\QuorumBox$), and `for a dense set' ($\CoquorumDiamond$) respectively, as per Proposition~\ref{prop.pointwise.dense.char}. 
\item\label{item.discuss.final}
$\final$ means $\sometime\forever$ and is from Figure~\ref{fig.3.derived}.
It means `eventually, this holds for all future times'. 

In our axiomatisation this is used to reflect a standard assumption in distributed systems of \deffont{eventual network synchrony}, or synonymously \deffont{global stabilisation (of the network)}, meaning that the communications network may be unstable and fail to pass messages at first, but then after a finite but unbounded amount of time, it becomes stable and messages arrive.
Our logic does not explicitly represent messages, but the use of $\final$ in our forward rules reflects this assumption about the underlying implementation. 
\item
$\texiaffine$ in \rulefont{PaxPropose01} and \rulefont{PaxWrite01} means `for zero or one elements', as per Subsection~\ref{subsect.unique.affine.existence}, so these two axioms intuitively state that each participant proposes or writes at most one value.

Note that our logic is three-valued, so what these axioms mean \emph{precisely} is something a little richer than is the case in two-valued logic.
In particular, if $\Forall{v}\modellabel{\tf{write}(v)}{\acontext}=\tvB$ --- corresponding to the case where a participant is crashed --- then $\acontext\ment\texiaffine\tf{write}$, \emph{even though} $\acontext\ment\tf{write}(v)$ holds for every $v$. 
See the discussion and results in Subsection~\ref{subsect.unique.affine.existence}.
\item
$\correct{\tf{propose}}$ means that $\tf{propose}$ returns $\tvT$ or $\tvF$ (not $\tvB$) on every value; see Definition~\ref{defn.correct}(\ref{item.correct}).
Similarly for the other unary predicate symbols in \rulefont{PaxPCorrect}.
\item
$\mru{\tf{accept}}{v}$ means `a most recently accepted value is $v$'.
See the discussion and definition in Subsection~\ref{subsect.mrup}.
\item
Finally, note that we intend the predicates in Figure~\ref{fig.logical.paxos} to be read as axioms.

What it means for an axiom to be valid in a model --- i.e. how we judge whether a model satisfies $\Theta_\ThyPax$, or how we judge whether a predicate is valid in all models of $\Theta_\ThyPax$ --- is unpacked in Figure~\ref{fig.validity} and Definition~\ref{defn.theory.axiom}(\ref{item.axioms.valid}).
We discuss the quantifier scoping in those definitions in Remark~\ref{rmrk.thy.quantifier.scope}.
\end{enumerate*}
\end{rmrk}




\subsection{High-level discussion of the axioms}

\begin{rmrk}
The predicates in Figure~\ref{fig.logical.paxos} are intended as a \emph{theory}, in the sense of Definition~\ref{defn.theory.axiom}, which is why we call them \emph{axioms} and not just `some predicates'.

That is: we fix a signature $\Sigma$ that has (at least) the predicate symbols required to express the axioms in Figure~\ref{fig.logical.paxos}, and we fix a model $\mathcal M$, and we derive properties that hold in all the valuations $\avaluation$ to $\mathcal M$ such that the theory $\Theta_\ThyPax$ from Figure~\ref{fig.validity} is valid.
 
At least as much design effort went into crafting the axioms $\Theta_\ThyPax$ of Declarative Paxos in Figure~\ref{fig.logical.paxos}, as went into designing the three-valued modal fixedpoint logic in which the axioms are expressed.\footnote{Of course there was design feedback in both directions; we put into the logic what we needed to express the axioms, and we refined the axioms to elegantly work within the logic.} 
Therefore, we will take our time to discuss the axioms in detail.
The reader who just wants to see proofs can skip the rest of this Section (for now) and look to Section~\ref{sect.paxos.correctness.properties} onwards.
\end{rmrk}


\begin{rmrk}
\label{rmrk.axioms.paxos.three.groups}
We have organised the axioms in Figure~\ref{fig.logical.paxos} into three groups, which we now discuss:
\begin{enumerate}
\item
\emph{Backward (inference) rules.}

Axioms with names ending with $?$ represent backwards reasoning and so we call them \deffont{backward rules}.
They use the strong implication $\timpc$ (Definition~\ref{defn.three.imp}(\ref{item.three.elementary.strong.implication})), which means `if the LHS is true ($\tvT$) then the RHS is true ($\tvT$)'.

The backward rules represent a strong inference that \emph{if} some $Y$ is observed to be $\tvT$ then it \emph{must} be because some other property $X$ is $\tvT$.
For instance, \rulefont{PaxAccept?} says that if $\tf{accept}(v)$ is true, then $\someoneAll\tf{write}(v)$ must also be true.

We discuss the backward rules individually in Remark~\ref{rmrk.discuss.backward.rules} below.
\item
\emph{Forward (progress) rules.}

Axioms with names ending with $!$ represent forwards reasoning about making progress, so we call them \deffont{forward rules}.
They use $\final$ and the weak implication $\tnotor$ (Definition~\ref{defn.three.imp}(\ref{item.three.elementary.weak.implication})), which means `if the LHS is true ($\tvT$) then the RHS is valid ($\tvT$ or $\tvB$)'.

We derive the forward rules from the backward rules using a systematic recipe which we discuss in Remark~\ref{rmrk.recipe}.

The forward rules represent a weak inference that \emph{if} some precondition $X$ is $\tvT$ then --- provided the network is stable and provided that the participant has not crashed (where `crashed' is interpreted as returning $\tvB$) --- then the postcondition $Y$ is $\tvT$. 

For instance, \rulefont{PaxAccept!} expresses that, after network synchrony (Remark~\ref{rmrk.functional.notation.in.paxos}(\ref{item.discuss.final})), if $\texi\someoneAll\tf{write}$ is true (by Notation~\ref{nttn.hos} this is shorthand for $\texi a.\someoneAll\tf{write}(a)$) then if the participant has not crashed then $\texi\tf{accept}$ (i.e. $\texi a.\tf{accept}(a)$) is true. 

We discuss the forward rules in Remark~\ref{rmrk.discuss.forward.rules} below.
\item
\emph{Other rules, for correctness and liveness.}

Other rules reflect particular conditions of the algorithms.
We discuss them in Remark~\ref{rmrk.paxos.axioms.discussion}.
\end{enumerate}
\end{rmrk}

\begin{rmrk}[Recipe for deriving forward rules from backward rules]
\label{rmrk.recipe}
A design feature of Figure~\ref{fig.logical.paxos} is that each forward rule is obtained from the corresponding backward rule by a \deffont{recipe}.
Namely: for a backward rule of the form 
$$
\tall v.\ \chi \timpc \phi[a\ssm v] \timpc \psi[a\ssm v]
$$
where 
\begin{itemize*}
\item
$\chi$ expresses some possible \emph{side-conditions} (e.g. $\chi=\tf{leader}$ in \rulefont{PaxDecideL?}) and 
\item
$\phi$ expresses some \emph{preconditions}, and 
\item
$\psi$ expresses some \emph{postconditions}, 
\end{itemize*}
our recipe constructs a forward rule of the form
$$
(\texi v.\,\chi \tand \phi)\tnotor \texi v.\psi .
$$
In most of the forward rules this result also get wrapped in a $\final$ modality, representing network synchrony  (Remark~\ref{rmrk.functional.notation.in.paxos}(\ref{item.discuss.final})).

This recipe holds for every pair of rules:
\begin{itemize*}
\item
We see it cleanly in the pairs \rulefont{PaxAccept?} / \rulefont{PaxAccept!}, \rulefont{PaxDecideL?} / \rulefont{PaxDecideL!}, and \rulefont{PaxDecide\tneg L?} / \rulefont{PaxDecide\tneg L!}.
These rules follow our recipe symbol-for-symbol.
\item
We see this also with \rulefont{PaxPropose?} / \rulefont{PaxPropose!}.
Our recipe generates the following version of \rulefont{PaxPropose!}: 
$$
(\texi v.\,\tf{leader}\tand v\tneq\tf{undef})\timpc\texi v.\tf{propose}(v).
$$
The reader can check that this is logically equivalent to what is written in Figure~\ref{fig.logical.paxos}.
\item 
It is less clear how our recipe works for \rulefont{PaxSend?} and \rulefont{PaxSend!} and \rulefont{PaxWrite?} and \rulefont{PaxWrite!}; e.g. \rulefont{PaxSend!} looks different from \rulefont{PaxSend?} --- it is clearly simpler!
Remarkably, these pairs too are consistent with our recipe; see Remark~\ref{rmrk.why.evaporates} and the subsequent discussion and results.
\end{itemize*}
\end{rmrk}

\begin{rmrk}
We make some further elementary comments on our recipe from Remark~\ref{rmrk.recipe}:
\begin{enumerate*}
\item
We uncurry the forward rules --- meaning that we write $(\texi v.\,\chi \tand \phi)\tnotor \texi v.\psi$ instead of $(\texi v.\,\chi) \tnotor \phi\tnotor \texi v.\psi$ --- because this is convenient for some proofs. 
It makes no other difference and either form would be correct.
\item
A forward rule of the form $\texi v.((\chi \tand \phi)\tnotor \psi)$ would be \emph{wrong}.
This is just a different predicate and is not what is intended.
\item
A forward rule of the form $\tall v.((\chi \tand \phi)\tnotor \psi)$ would not necessarily be wrong, but it is stronger than we need for our proofs.
This reflects a general design principle in logical theories that \emph{weaker axioms are better}.
If we can prove the same result from weaker axioms then generally speaking we should, because that is a stronger result.\footnote{In practice, working from weaker axioms tends to force us to think harder and so design a better-engineered proof.}
\end{enumerate*}
\end{rmrk}


\subsection{More detailed discussion of the axioms}


\begin{rmrk}
\label{rmrk.paxos.axioms.discussion}
We discuss the \emph{axioms for correctness and liveness}:
\begin{enumerate}
\item\label{item.discuss.LdrExist}
\rulefont{LdrExist}:

In implementations of Paxos, at each stage there is a unique \emph{leader}, whose identity is computed by applying (what amounts to) a random number generator to an agreed seed.
As the level of abstraction in this paper, we just assume that a leader exists --- we do not care how it is chosen.

We use a predicate constant symbol $\tf{leader}$ to identify the leader.
Because our logic is modal --- across time $n\in\Time$ and participants $p\in\Pnt$ --- we can just assume that a constant symbol $\tf{leader}$ is assigned truth-value $\tvT$ where a participant $p$ that is a leader at a stage $n$, and is assigned truth-value $\tvF$ where a participant $p$ that is not a leader at a stage $n$.

In more technical language: $\avaluation$ is such that $n,p\mentval \modTF\tf{leader}$ for every time $n\in\Time$ and participant $p\in\Pnt$, %
and $p$ is a leader %
precisely when $n,p\mentval\tf{leader}$.
\rulefont{LdrExist} just asserts that a participant either is a leader or it is not, and some leader exists.
\item\label{item.discuss.LdrExistUniq}
\rulefont{LdrExt}:

\rulefont{LdrExt} is a rule-scheme (one for each closed $\phi$) and asserts that leaders are extensionally equal; %
a formal statement and proof are in Lemma~\ref{lemm.LeaderExtUniq}.

In an implementation this is enforced by having a unique leader, but in the logic this is unnecessary --- we don't care how many participants claim to be leaders, so long as they are all saying the same thing! --- and also it is impossible because (although we have modalities $\everyone$ and $\someone$) we do not have an explicit quantification over elements of $\Pnt$ in our logic.
This is a feature, not a bug: we want our logic to be no more expressive than it needs to be.

In fact, the axiom-scheme \rulefont{LdrExt} is stronger than we need for our proofs;
more on this in Subsection~\ref{subsect.ldrexistuniq'}.


Henceforth we may say `the leader' instead of `a leader', because up to extensional equivalence there is just one.
\item
\rulefont{LdrCorrect}:

\rulefont{LdrCorrect} asserts that for infinitely many times $n\in\Time$, the leader $l$ (which exists by the right-hand conjunct in \rulefont{LdrExist}) is $\tf{propose}$-correct, $\tf{write}$-correct, and $\tf{decide}$-correct (terminology from Definition~\ref{defn.correct}(\ref{item.P-correct}).

Unpacking Definition~\ref{defn.correct}(\ref{item.P-correct}), this means that $n,l\mentval \modTF\tf{propose}(v)$ for every $v\in\Val$, and similarly for $\tf{write}$ and $\tf{decide}$.
Intuitively, this is a \emph{liveness assumption} that there is always a future time when the leader is not crashed in the sense that it can correctly $\tvT\tvF$-propose, $\tvT\tvF$-write, and $\tvT\tvF$-decide a value. 
\item
\rulefont{PaxPCorrect}:

Given a valuation $\avaluation$, we unpack Definition~\ref{defn.validity.judgement}(\ref{item.mentval}) to see what $\mentval \rulefont{PaxPCorrect}$ actually means, as per Definition~\ref{defn.theory.axiom}:
\begin{multline*}
\Forall{n\in\Time}\Exists{O\in\opens}\Forall{p\in O}\Forall{v\in\Val}n,p,O\mentval\correct{\tf P(v)}
\\
\text{for each}\ \tf P\in\{\tf{propose},\tf{send},\tf{accept},\tf{write},\tf{decide}\} .
\end{multline*}
This corresponds to a standard correctness assumption in Paxos that at each stage there is a quorum of uncrashed participants.

However, the axiom \rulefont{PaxPCorrect} is subtly more general than the standard correctness assumption, because there is no restriction that the \emph{same} open set $O$ be used for each unary predicate symbol.\footnote{If we wanted the same $O$, then we would scope the $\QuorumBox$ modality further out and write $\QuorumBox(\correct{\tf{propose}} \tand \correct{\tf{send}} \tand \correct{\tf{accept}} \tand 
\correct{\tf{write}} \tand \correct{\tf{decide}})$.  This turns out to be unnecessary.}
It turns out that the proofs require only this assumption, and no more. 
\item\label{item.01.discussion}
\rulefont{PaxPropose01} and \rulefont{PaxWrite01}:

We justify these axioms with respect to the Paxos algorithm as summarised in Remark~\ref{rmrk.high-level.paxos}:
\begin{itemize*}
\item
Step~\ref{item.l.proposes} of Remark~\ref{rmrk.high-level.paxos} states that leader is crashed or proposes one value.
The intuitive correspondence to \rulefont{PaxPropose01} is clear.
\item
Step~\ref{item.l.writes} of Remark~\ref{rmrk.high-level.paxos} states that the leader writes a value derived from \emph{the first} quorum of participants that it hears back from (or none, if it hears back from no such quorum). 
One corollary of this, is that it writes at most one value.
Our logical theory has no state machine and notion of algorithmic time (as per Remark~\ref{rmrk.elisions}(\ref{item.elide.algorithmic.time})) --- so there is no sense in which information from one open set `arrives' at the leader `before' information from another --- so from the point of view of the logic, of all the open sets available, the leader picks at most one to act on.
This is axiom \rulefont{PaxWrite01}. 
\end{itemize*}
It is a fact that \rulefont{PaxPropose01}+\rulefont{PaxPropose!} and \rulefont{PaxWrite01}+\rulefont{PaxWrite!} give us
$$
\tf{leader}\tnotor\texiunique\tf{propose}
\quad\text{and}\quad
\final((\tf{leader}\tand\QuorumBox\texi\tf{send})\tnotor \texiunique\tf{write}) ,
$$ 
and
\rulefont{PaxPropose01}+\rulefont{PaxPropose?} and \rulefont{PaxWrite01}+\rulefont{PaxWrite?} give us
$$
\tneg\tf{leader}\tnotor\tneg\texi\tf{propose}
\quad\text{and}\quad
\tneg\tf{leader}\tnotor\tneg\texi\tf{write} .
$$
So it follows from our axioms that a leader participant proposes and writes one value (assuming it is not crashed, and if messages arrive), and a non-leader participant proposes and writes zero values. 
\item
\rulefont{Pax2Twined}:

This axiom collects a set of instances of Lemma~\ref{lemm.supertwined.qcbt}(\ref{item.intertwined.3}), namely taking $\phi=\tf{accept}(v)$ in that Lemma for every $v\in\Val$.

There is a design choice here: instead of an axiom \rulefont{Pax2Twined}, we could insist specifically that the semitopology in our underlying model be $2$-twined.
Then \rulefont{Pax2Twined} ceases to be an axiom, but we still have it as a fact of the model by directly invoking Lemma~\ref{lemm.supertwined.qcbt}(\ref{item.intertwined.3}).
Or we could be even more specific: we could insist that the underlying semitopology be $\AllBut{N}{f}$ (Definition~\ref{defn.allbut}) for $N> 2*f$; by Proposition~\ref{prop.allbut.supertwined}(\ref{item.allbut.intertwined}) this is $2$-twined, and then we can use Lemma~\ref{lemm.supertwined.qcbt}(\ref{item.intertwined.3}) to obtain \rulefont{Pax2Twined}.

We choose here to express as an axiom precisely the property that our proofs require.
This is concise, general, and yields clean proofs.
Insisting on a $2$-twined semitopology, or specifically on the semitopology $\AllBut{N}{f}$ for $N>2*f$, would also make sense, but in the context of our logical approach, using an axiom seems nice. %
\end{enumerate}
\end{rmrk}

\begin{rmrk}
\label{rmrk.discuss.backward.rules}
We discuss the \emph{axioms for the backward rules}:
\begin{enumerate}
\item
\rulefont{PaxPropose?}:

If we propose some value $v$ then 
\begin{enumerate}
\item
we are the leader and 
\item
$v$ is not the undefined value $\udfn$.
\end{enumerate}
\item
\rulefont{PaxSend?}:

If we send some value $v$ then 
\begin{enumerate}
\item
somebody must have proposed some value and 
\item
\begin{enumerate*}
\item
\emph{either} $v$ is the value we have most recently accepted, 
\item
\emph{or} we never accepted any value and $v$ is the undefined value.
\end{enumerate*}
\end{enumerate}
\item
\rulefont{PaxWrite?}:

If we write some value $v$ then 
\begin{enumerate}
\item
we are a leader, and
\item
there exists an open set of points (= participants) that have all sent us values, and
\item
\begin{enumerate*}
\item
\emph{either} some point in that open set has most recently accepted $v$, 
\item
\emph{or} 
$v$ is the value that we proposed, and no point in that open set has ever accepted any value. 
\end{enumerate*}
\end{enumerate}
\item
\rulefont{PaxAccept?}:

If we accept $v$ then somebody somewhere must have written $v$.
\item
\rulefont{PaxDecideL?}:

If a leader decides $v$ then an open set of points must have all accepted $v$.
\item
\rulefont{PaxDecide{\tneg}L?}:

If we are not a leader and we decide $v$, then some leader somewhere must have decided $v$. 
\end{enumerate}
\end{rmrk}

\begin{rmrk}
\label{rmrk.discuss.forward.rules}
We discuss the \emph{axioms for the forward rules}.

We noted in Remark~\ref{rmrk.recipe} our design methodology of deriving the forward rules from the backward rules; forward rules are rules for making progress, so they are existentially quantified over the variable $v$ and they use the weak implication $\tnotor$ to emulate possible crashes.

The other subtlety is the $\final$ modality wrapping all of the forward rules, except for \rulefont{PaxPropose!}.
\begin{itemize*}
\item
\emph{What $\final$ does for us in most of the forward rules.}

It models possible network failure, and eventual network synchrony.
As touched on in Remark~\ref{rmrk.functional.notation.in.paxos}(\ref{item.discuss.final})), the $\final$ modality reflects the possibility that the network might be down for a while, but eventually --- `finally' --- the network stabilises and messages arrive.

Predicates in our logic are \emph{not} messages, but we get the effect of eventual network synchrony in our logic (even though messages are not explicitly present) by wrapping the implications in our forward rules in the $\final$ modality.

More on this when we introduce the notion of \emph{global stabilisation of logical time} in Subsection~\ref{subsect.gslt}.
\item
\emph{Why $\final$ is absent from \rulefont{PaxPropose!}.}

Because this rule models a property of the implementation that involves no network communication.
The leader does not need to communicate with any other participant to pick a value to propose.

It is true that the leader's propose message might not arrive if the network is down --- but this is reflected by the use of $\final$ in the \rulefont{PaxSend!} rule.
\end{itemize*}
\end{rmrk}


\subsection{Some asides}

\subsubsection{An aside on \rulefont{PaxWrite?}}

\rulefont{PaxWrite?} contains this predicate:
$$
\mru{(\someone\tf{accept})}{v} .
$$
Unpacking what this means (cf. Remark~\ref{rmrk.explain.mrup}), this looks into the past to check whether $v$ is a value that was most recently accepted by some participant.

Note that this slightly differently-scoped predicate would be \emph{wrong}:
$$
\someone(\mru{\tf{accept}}{v}) .
$$
This finds some participant such that $v$ is a value that was most recently accepted by \emph{that} participant (but there might be another participant who more recently accepted some other value).

The reader with a background in logic will know all about how important scoping can be (e.g. $\forall x.\exists y.\phi$ vs. $\exists x.\forall y.\phi$, and $\exists x.(\phi\land \phi')$ vs. $\exists x.\phi\land\exists x.\phi'$), but perhaps the point bears repeating.


\subsubsection{An aside on \rulefont{LdrExt}}

In a Paxos implementation, there is a unique leader at each stage. 
We mentioned in Remark~\ref{rmrk.paxos.axioms.discussion}(\ref{item.discuss.LdrExist}\&\ref{item.discuss.LdrExistUniq}) that \rulefont{LdrExist} asserts that a leader exists and that the \rulefont{LdrExt} axiom-scheme asserts that all leader participants are extensionally equal.
We now make this formal:\footnote{Later on in Subsection~\ref{subsect.ldrexistuniq'} we will note weaker, but more abstract, forms of the axiom.}
\begin{lemm}[Leaders are extensionally equivalent]
\label{lemm.LeaderExtUniq}
Suppose $n\in\Time$ and $\avaluation$ is a valuation, and suppose $n\mentval\rulefont{LdrExt}_\phi$ holds for every closed predicate $\phi$.\footnote{By Definition~\ref{defn.validity.judgement}(\ref{item.ment.n}) this means $\Forall{p\in\Pnt}\Forall{O\in\opensne}(\modellabel{\rulefont{LdrExt}_\phi}{n,p,O,\avaluation}\in\{\tvT,\tvB\})$.} 
Suppose that $l,l'\in\Pnt$ and $n,l\mentval\modT\tf{leader}$ and $n,l'\mentval\modT\tf{leader}$.

Then $l$ and $l'$ are \deffont{extensionally equivalent} at $(n,\avaluation)$, by which we mean that for every closed predicate $\phi$ and every $O\in\opensne$,
$$
\modellabel{\phi}{n,l,O,\avaluation} = \modellabel{\phi}{n,l',O,\avaluation} .
$$
\end{lemm}
\begin{proof}
Suppose $\phi$ is a closed predicate and $O\in\opensne$, and suppose for the sake of argument that $\modellabel{\phi}{n,l,O,\avaluation}=\tvB$.

It is routine to check from Figures~\ref{fig.3.derived} and~\ref{fig.3} and from Lemma~\ref{lemm.dense.char}(\ref{item.dense.char.someoneAll}) that $\modellabel{\someoneAll\modB\phi}{n,l,O,\avaluation}=\tvT$. 
By assumption $n,l',O\mentval \rulefont{LdrExt}_{\modB\phi}$, and by two applications of \strongmodusponens (once for $\modT\tf{leader}$, and another for $\modT\someoneAll(\tf{leader}\tand\modB\phi)$) we conclude that $\modellabel{\modB\phi}{n,l',O,\avaluation}=\tvT$, and therefore $\modellabel{\phi}{n,l',O,\avaluation}=\tvB$ as required. 
 
The cases $\modellabel{\phi}{n,l,O,\avaluation}=\tvT$ and $\modellabel{\phi}{n,l,O,\avaluation}=\tvF$ are be precisely similar, and this suffices to prove the result.
\end{proof}


\subsubsection{An aside on \rulefont{PaxSend?} and \rulefont{PaxSend!}}

\begin{rmrk}
\label{rmrk.why.evaporates}
In Remark~\ref{rmrk.recipe} we proposed a recipe for deriving a forward rule from a backward one: a backward rule of the form
$$
\rulefont{Rule?} 
\qquad
\tall v.\ \chi \timpc \phi[a\ssm v] \timpc \psi[a\ssm v]
$$
generates a forward rule (possibly wrapped in a $\final$ modality) of the form
$$
\rulefont{Rule!}
\qquad
(\texi v.\, \chi \tand \phi)\tnotor \texi v.\psi .
$$
We noted in Remark~\ref{rmrk.recipe} that the recipe works cleanly for some rules --- but it was less clear how it can be applied to derive \rulefont{PaxSend!} from \rulefont{PaxSend?}, and \rulefont{PaxWrite!} from \rulefont{PaxWrite?}.
In these cases, the left-hand sides of \rulefont{PaxSend!} and \rulefont{PaxWrite!} seem much simpler than the right-hand sides of \rulefont{PaxSend?} and \rulefont{PaxWrite?} respectively. 

We now show why, in fact, \rulefont{PaxSend!} and \rulefont{PaxWrite!} are correct and make sense in the context of our recipe, culminating with Corollary~\ref{corr.back.ok}.
\end{rmrk}


\begin{lemm}
\label{lemm.evaporates.simplify}
We have:
\begin{enumerate*}
\item
$\ment \texi v.\bigl((\mru{\tf{accept}}{v})\ \tor (\tneg\modT\recent\texi\tf{accept}\tand v\teq\udfn)\bigr) \ \equiv\ \tvT$.
\item
$\ment \texi v.\bigl(\ (\mru{(\someone\tf{accept})}{v})\ \tor\ (\tneg\modT\recent\texi\someone\tf{accept}\tand \tf{propose}(v))\ \bigr)\ \equiv\ \texi\tf{propose}$.
\end{enumerate*}
Note that:
\begin{itemize*}
\item
The notation $\texi\tf{accept}$ and $\texi\tf{propose}$ and $\texi\someone\tf{accept}$ is from Notation~\ref{nttn.hos}.
\item
We do not restrict to contexts and valuations that satisfy \ThyPax; the meaning of the $\ment$ symbol above is as in Definition~\ref{defn.validity.judgement}(\ref{item.ment.phi}).
\end{itemize*}
\end{lemm}
\begin{proof}
We consider each part in turn:
\begin{enumerate}
\item
We distribute the $\texi v$ through the $\tor$, and then through the $\tand$ on the right-hand side (we can do this because $v$ is not free in the left-hand conjunct) and we see that it suffices to show 
$$
\ment \texi v.(\mru{\tf{accept}}{v})\ \tor\ (\tneg\modT\recent\texi\tf{accept}\tand \texi v.v\teq\udfn)\ \equiv\ \tvT
$$
Using Lemma~\ref{lemm.evaporates} this simplifies to just
$$
\ment \modT\recent\texi\tf{accept}\ \tor\ (\tneg\modT\recent\texi\tf{accept}\tand \texi v.v\teq\udfn) \ \equiv\ \tvT.
$$
This clearly holds, just taking $v\teq\udfn$.
\item
We distribute $\texi v$ into the predicate as we did for the proof of part~1.
We see that it suffices to show 
$$
\ment \bigl(\ \texi v.(\mru{(\someone\tf{accept})}{v})\ \tor\ (\tneg\modT\recent\texi\someone\tf{accept}\tand \texi\tf{propose})\ \bigr)\ \equiv\ \texi\tf{propose}.
$$
We simplify using Lemma~\ref{lemm.evaporates} to
$\ment \texi\tf{propose}\ \equiv\ \texi\tf{propose}$,
which is true.\footnote{Yes, we have a lemma for that: Lemma~\ref{lemm.equiv}(\ref{item.equiv.2}) and reflexivity of equality.}
\qedhere
\end{enumerate}
\end{proof}

We can now justify the left-hand sides of \rulefont{PaxSend!} and \rulefont{PaxWrite!} as being consistent with our recipe as discussed above in Remark~\ref{rmrk.why.evaporates}:
\begin{corr}
\label{corr.back.ok}
We have:
$$
\begin{array}{r@{\ }l}
\ment\texi v.&\bigl(\someoneAll\texi\tf{propose}\ \tand \bigl((\mru{\tf{accept}}{v})\ \tor (\tneg\modT\recent\texi\tf{accept}\tand v\teq\udfn)\bigr)\bigr) 
\\
&\equiv \someoneAll\texi\tf{propose}
\\[2ex]
\ment\texi v.&\bigl(\tf{leader}\tand 
\Quorum(\Box\texi\tf{send} \tand 
\bigl((\mru{(\someone\tf{accept})}{v}) \tor
(\tneg\modT\recent\texi\someone\tf{accept}\tand \tf{propose}(v))\bigr)) 
\bigr)
\\
&\equiv
(\tf{leader}\tand \QuorumBox\texi\tf{send} \tand \texi\tf{propose})
\end{array}
$$
\end{corr}
\begin{proof}
By elementary predicate manipulations, distributing the $\texi v$ quantifier down into the predicate and using Lemma~\ref{lemm.evaporates.simplify}.
\end{proof}

 




\section{Correctness properties}
\label{sect.paxos.correctness.properties}

\subsection{Summary of the properties}

\begin{rmrk}
\label{rmrk.standard.paxos.correctness}
We will derive the following standard correctness properties for Paxos (as laid out in~\cite[Subsection~2.2.1]{cachin2011paxos}) from the theory \ThyPax in Figure~\ref{fig.logical.paxos}:
\begin{enumerate*}
\item
\textbf{Validity:} \quad If a participant decides some value $v$, then $v$ was proposed by some participant.

A formal statement and proof is in Theorem~\ref{thrm.validity}.
\item
\textbf{Agreement:}\quad 
If two participants decide values, then those values are equal.

A formal statement and proof is in Theorem~\ref{thrm.full.agreement}.
\item
\textbf{Termination:}\quad 
Every correct participant eventually decides some value.

A formal statement and proof is in Theorem~\ref{thrm.termination}.
\item
\textbf{Integrity:}\quad 
Every correct (non-crashing) participant decides at most once.

Our logic abstracts away from a direct notion of execution, by design, but this property is represented indirectly in how Theorem~\ref{thrm.termination} is expressed.
See Remark~\ref{rmrk.paxos.integrity} below. 
\end{enumerate*}
\end{rmrk}

\begin{rmrk}[On integrity]
\label{rmrk.paxos.integrity}
In implementations of Paxos, once a participant has decided, it stops.
But in our logic, `stop' is not even a concept.
There is no abstract machine to stop; there are just predicates, with denotations.

Nevertheless, the integrity property is echoed in our maths, as the \emph{infinite repetition} in Theorem~\ref{thrm.termination}.
Once a predicate $\tf{decide}(v)$ returns $\tvT$, it does so \emph{forever} --- and because of Agreement, it never decides any other value.
A predicate eventually always returning the same truth-value is a model of an abstract machine that has computed its result, and then just sits there forever, displaying that same result, unchanged, every time we query its display.
\end{rmrk}

\subsection{Validity}


\begin{rmrk}
\label{rmrk.validity.in.english}
\emph{Validity} is expressed in English as:
\begin{quote}
\emph{If a participant decides some value $v$, then $v$ was proposed by some participant.}
\end{quote}
We render this in our logic in Theorem~\ref{thrm.full.agreement} as
$$
\ThyPax\ment\tall v.(\tf{decide}(v) \timpc \urecent\someoneAll(\tf{leader}\tand\tf{propose}(v))).
$$

In words we can read this as: 
\begin{quote}
If someone decides $v$ now, then at some stage now or in the past a leader of that stage proposed $v$. 
\end{quote}
We will now set about proving this property from our axioms, culminating with Theorem~\ref{thrm.validity}.
\end{rmrk}


\begin{lemm}
\label{lemm.validity.helper}
Suppose $\avaluation$ is a valuation such that $\mentval\ThyPax$ ($\mentval\ThyPax$ is from Figure~\ref{fig.validity}).
Suppose $n\in\Time$ and $v\in\Val$.
Then:
\begin{enumerate*}
\item\label{item.validity.helper.3}
$n\mentval \someoneAll\tf{write}(v) \timpc \urecent\someoneAll\tf{propose}(v)$. 
\item\label{item.validity.helper.2}
$n\mentval \someoneAll\tf{accept}(v) \timpc \urecent\someoneAll\tf{propose}(v)$. 
\end{enumerate*}
\end{lemm}
\begin{proof}
We prove both parts by a simultaneous induction on $n$.
We freely use \strongmodusponens:
\begin{enumerate}
\item
Suppose $n\mentval\modT\someoneAll\tf{write}(v)$.
Then using \rulefont{PaxWrite?} %
$$
n\mentval \modT\Quorum(\mru{\someone\tf{accept}}{v}) 
\quad\text{or}\quad
n\mentval \modT\someoneAll\tf{propose}(v) .
$$
We consider the two cases:
\begin{itemize*}
\item
\emph{If $n\mentval \modT\someoneAll\tf{propose}(v)$,}\ then 
$n\mentval\modT\urecent\someoneAll\tf{propose}(v)$ follows immediately from the clause for %
$\urecent$ in Figure~\ref{fig.3.derived}.
\item
\emph{If $n\mentval \modT\Quorum(\mru{\someone\tf{accept}}{v})$,}\ 
then by Lemma~\ref{lemm.mru.to.recent}(\ref{item.mru.to.recent.1}) $n\mentval\modT\recent\someoneAll\tf{accept}(v)$,
so that by Figure~\ref{fig.3.derived} $n'\mentval\modT\someoneAll\tf{accept}(v)$ for some $n'<n$.
By part~\ref{item.validity.helper.2} of this result for $n'<n$, we have that $n'\mentval\modT\urecent\someoneAll\tf{propose}(v)$, and thus using Lemma~\ref{lemm.tomorrow.forever.recent}(\ref{item.rur.r}) $n\mentval\modT\urecent\someoneAll\tf{propose}(v)$. 
\end{itemize*}
In either case, we are done.
\item
Suppose $n\mentval\modT\someoneAll\tf{accept}(v)$.
By \rulefont{PaxAccept?} $n\mentval\modT\someoneAll\tf{write}(v)$.
We use part~\ref{item.validity.helper.3} of the inductive hypothesis for $n$. 
\qedhere\end{enumerate}
\end{proof}



\begin{thrm}[Validity]
\label{thrm.validity}
We have:
$$
\ThyPax\ment\tall v.(\tf{decide}(v) \timpc \urecent\someoneAll(\tf{leader}\tand\tf{propose}(v))) .
$$
\end{thrm}
\begin{proof}
Suppose $\avaluation$ is a valuation such that $\mentval\ThyPax$. 
We reason using \strongmodusponens. 
Suppose $v\in\Val$ and $n\in\Time$ and $p\in\Pnt$, and suppose $n,p\mentval\modT\tf{decide}(v)$.
By \rulefont{PaxDecideL?} (and \rulefont{PaxDecide\tneg L?} if $p\neq\f{leader}(n)$), $n\mentval\modT\QuorumBox\tf{accept}(v)$.
It follows by Proposition~\ref{prop.someone.implies.someoneAll}(\ref{item.someone.implies.someoneAll.2}) that $n\mentval\modT\someoneAll\tf{accept}(v)$.
By Lemma~\ref{lemm.validity.helper}(\ref{item.validity.helper.2}) $n\mentval\modT\urecent\someoneAll\tf{propose}(v)$, and using \rulefont{PaxPropose?} we conclude that $n\mentval\modT\urecent\someoneAll(\tf{leader}\tand\tf{propose}(v))$ as required. 
\end{proof}




\subsection{Agreement}

\begin{rmrk}
\emph{Agreement} is expressed in English as:
\begin{quote}
\emph{If two participants decide values, then those values are equal.}
\end{quote}
We render this in our logic in Theorem~\ref{thrm.full.agreement} as
$$
\ThyPax\ment\tall v,v'.(\urecent\someoneAll\tf{decide}(v) \timpc \someoneAll\tf{decide}(v') \timpc v'\teq v) .
$$
In words we can read this as: 
\begin{quote}
If someone decides $v$ now or in the past, and someone decides $v'$ now, then $v'=v$.
\end{quote}
It is convenient to prove the `now' and the `in the past' parts of this as two sub-cases: 
\begin{enumerate*}
\item
Proposition~\ref{prop.decide.n.agree} handles the case of $\someoneAll\tf{decide}(v)\tand\someoneAll\tf{decide}(v')$, i.e. when the two decisions are made at the same time.
\item
Theorem~\ref{thrm.full.agreement} then extends this with the case of $\recent\someoneAll\tf{decide}(v)\tand\someoneAll\tf{decide}(v')$, i.e. when one decision precedes the other.
\end{enumerate*}
\end{rmrk}


\subsubsection{Agreement at time $n$}



\begin{lemm}
\label{lemm.various.unique}
Suppose $\avaluation$ is a valuation such that $\mentval\ThyPax$. 
Suppose $n\in\Time$ and $v,v'\in\Val$. 
Then:
\begin{enumerate*}
\item\label{item.various.unique.write}
$n\mentval (\someoneAll\tf{write}(v)\tand \someoneAll\tf{write}(v')) \timpc v\teq v'$. 
\item\label{item.various.unique.accept}
$n\mentval (\someoneAll\tf{accept}(v)\tand \someoneAll\tf{accept}(v')) \timpc v\teq v'$. 
\item\label{item.various.unique.recent.accept}
$n\mentval \bigl(\mru{(\someoneAll\tf{accept})}{v}\tand \mru{(\someoneAll\tf{accept})}{v'}\bigr) \timpc v\teq v'$.
\end{enumerate*}
\end{lemm}
\begin{proof}
We consider each part in turn, using \strongmodusponens (we may also silently commute the $\modT$ modality as convenient):
\begin{enumerate}
\item
Suppose $n\mentval\modT\someoneAll\tf{write}(v)\tand\modT\someoneAll\tf{write}(v')$.
Using \rulefont{PaxWrite?} 
$n\mentval\modT\someoneAll(\tf{write}(v)\tand\tf{leader})\tand\modT\someoneAll(\tf{write}(v')\tand\tf{leader})$. 
Using \rulefont{LdrExt} therefore $n\mentval\modT\someoneAll(\tf{write}(v)\tand\tf{write}(v')\tand\tf{leader})$.
By \rulefont{PaxWrite01} and Lemma~\ref{lemm.unique.affine.existence}(\ref{item.unique.affine.existence.01.implies}) we have $v=v'$ as required.
\item
Suppose $n\mentval \modT\someoneAll\tf{accept}(v)\tand\modT\someoneAll\tf{accept}(v')$.
By \rulefont{PaxAccept?} $n\mentval\modT\someoneAll\tf{write}(v)\tand\modT\someoneAll\tf{write}(v')$.
We use part~\ref{item.various.unique.write} of this result.
\item
Suppose $n\mentval\modT(\mru{(\someoneAll\tf{accept})}{v})\tand\modT(\mru{(\someoneAll\tf{accept})}{v'}\bigr)$.\footnote{By Lemma~\ref{lemm.mru.unambivalent}(2) we do not need to write the $\modT$ here but we do anyway; see Remark~\ref{rmrk.explain.mru}(\ref{item.explain.mru.T}).}
Using Lemma~\ref{lemm.mru.to.recent}(\ref{item.mru.to.recent.2}) (taking $a.\phi=\someoneAll\tf{accept}$ in that Lemma) there exists $n'<n$ such that 
$n'\mentval\modT\someoneAll\tf{accept}(v)\tand\modT\someoneAll\tf{accept}(v')$.
We use part~\ref{item.various.unique.accept} of this result.
\qedhere\end{enumerate}
\end{proof}



\begin{prop}
\label{prop.decide.n.agree}
Suppose $\avaluation$ is a valuation such that $\mentval\ThyPax$. 
Suppose $n\in\Time$ and $v,v'\in\Val$.
Then
$$
n\mentval (\someoneAll\tf{decide}(v) \tand \someoneAll\tf{decide}(v')) \timpc v\teq v' . 
$$
\end{prop}
\begin{proof}
We use \strongmodusponens.
Suppose 
$n\mentval \modT(\someoneAll\tf{decide}(v)\tand\someoneAll\tf{decide}(v'))$.
Using \rulefont{PaxDecide\tneg L?} if needed, and then \rulefont{PaxDecideL?} twice, we have
$n\mentval \modT(\QuorumBox\tf{accept}(v)\tand\QuorumBox\tf{accept}(v'))$ and thus by Proposition~\ref{prop.someone.implies.someoneAll}(\ref{item.someone.implies.someoneAll.2}) and Lemma~\ref{lemm.commuting.connectives}
also $n\mentval \modT(\someoneAll\tf{accept}(v)\tand\someoneAll\tf{accept}(v'))$.
We can now use Lemma~\ref{lemm.various.unique}(\ref{item.various.unique.accept}).
\end{proof}



\begin{rmrk}
Proposition~\ref{prop.decide.n.agree} proves part of the agreement property; that no two correct participants decide differently \emph{at a given time}.
To prove the result across different times will require a bit more work, culminating with Theorem~\ref{thrm.full.agreement}.
\end{rmrk}


\subsubsection{Some fine structure of the rules}
\label{subsect.ldrexistuniq'}

We continue a discussion from Remark~\ref{rmrk.paxos.axioms.discussion}(\ref{item.discuss.LdrExistUniq}) and also Lemma~\ref{lemm.LeaderExtUniq}:

Note that we use the axiom-scheme \rulefont{LdrExt} from Figure~\ref{fig.logical.paxos} just once: in the proof of Lemma~\ref{lemm.various.unique}(\ref{item.various.unique.write}).

Examining how we use this axiom-scheme in the proof above, we see that for our specific purposes it would suffice to assume a single axiom as follows:
$$
\hspace{-6em}\rulefont{LdrExt'}
\qquad\qquad
\tf{leader}\timpc \someoneAll(\tf{leader}\tand\tf{write}(v))\timpc \tf{write}(v) .
$$ 
The axiom scheme \rulefont{LdrExt} asserts that there is a unique leader at each stage, up to extensional equivalence, which is a fair reflection of how the algorithms work where there \emph{is} a unique leader at each stage.

\rulefont{LdrExt'} asserts that there is a unique leader at each stage, up to extensional equivalence \emph{for writing values}.
In the proofs, this turns out to be all we need.

In fact, we can simplify this just to 
$$
\hspace{-6em}\rulefont{SPaxWrite01}
\qquad\qquad
\texiaffine\someoneAll\tf{write} .
$$
In full, this is $\texiaffine v.\someoneAll\tf{write}(v)$, and in words it expresses: 
\begin{quote}
`at each stage, at most one value is written'.
\end{quote}
In the presence of \rulefont{PaxWrite?}, \rulefont{SPaxWrite01} is implied by \rulefont{LdrExt'} --- but it does not imply \rulefont{LdrExt'}, because it permits the possibility that one leader participant writes a value and another leader participant crashes.
The reader can check that \rulefont{SPaxWrite01} is sufficient to complete the proof of Lemma~\ref{lemm.various.unique}(\ref{item.various.unique.write}).

Note also that that \rulefont{SPaxWrite01} implies \rulefont{PaxWrite01}, so that --- for the proofs we care about --- \rulefont{SPaxWrite01} can replace both \rulefont{LdrExt} and \rulefont{PaxWrite01}. 
We will do exactly this, later: the job of $\Theta_\ThyPax$ is to be `declarative Paxos', that is, its purpose is to be declarative \emph{and} to be recognisably Paxos. 
Later on in Figure~\ref{fig.ThySPax} we will define a \emph{simpler} declarative Paxos in (whence the name `SPaxWrite01').
This simplified theory is derived by looking at the proofs here, and seeing what we actually require.










 



\subsubsection{Agreement at possibly different times}



\begin{prop}
\label{prop.integrity.helper}
Suppose $\avaluation$ is a valuation such that $\mentval\ThyPax$, and suppose $v,v'\in\Val$ and $\avaluation$. 
Then:
\begin{enumerate*}
\item\label{item.integrity.helper.accept}
$\mentval\recent\QuorumBox\tf{accept}(v) \timpc \someoneAll\tf{accept}(v')\timpc v'\teq v$.
\item\label{item.integrity.helper.write}
$\mentval\recent\QuorumBox\tf{accept}(v) \timpc \someoneAll\tf{write}(v')\timpc v'\teq v$.
\end{enumerate*}
\end{prop}
\begin{proof}
Suppose $n\in\Time$.
We will reason using \strongmodusponens, so suppose 
\begin{equation}
\label{eq.n.1}
n\mentval\modT\QuorumBox\tf{accept}(v). 
\end{equation}
We will prove by induction on $n'>n$ that for every $v'\in\Val$,
$$
\begin{array}{r@{\ }l}
n'\mentval&\modT(\mru{(\someoneAll\tf{accept})}{v'}\timpc v'\teq v)
\quad\text{and}\quad
\\
n'\mentval&\modT(\someoneAll\tf{write}(v')\timpc v'\teq v)
\end{array}
$$
Clearly, by \strongmodusponens the $\tf{write}$ part of the inductive hypothesis is enough to prove part~\ref{item.integrity.helper.write} of this result.
For the $\tf{accept}$ part of the inductive hypothesis, we see that this implies part~\ref{item.integrity.helper.accept} of this result either by a routine argument using \rulefont{PaxAccept?} and part~\ref{item.integrity.helper.write} of this result --- or using Corollary~\ref{corr.mru.today.to.yesterday}(\ref{item.mru.today.implies.tomorrow}) (for $a.\phi=a.\someoneAll\tf{accept}(a)$ in that result).

We consider cases:
\begin{itemize}
\item
\emph{The base case for $\tf{accept}$, where $n'=n+1$.}

Suppose $n'\mentval\modT(\mru{\someoneAll\tf{accept}}{v'})$.

We assumed in equation~\eqref{eq.n.1} that $n\mentval\modT\QuorumBox\tf{accept}(v)$, so by Proposition~\ref{prop.someone.implies.someoneAll}(\ref{item.someone.implies.someoneAll.2}) also $n\mentval\modT\someoneAll\tf{accept}(v)$, and by Corollary~\ref{corr.mru.today.to.yesterday}(\ref{item.mru.today.implies.tomorrow}) $n'\mentval\modT(\mru{(\someoneAll\tf{accept})}{v})$. 
By Lemma~\ref{lemm.various.unique}(\ref{item.various.unique.recent.accept}) $v'=v$.

\item
\emph{The inductive step for $\tf{accept}$, where $n'>n+1$.}

Suppose $n'\mentval\modT(\mru{\someoneAll\tf{accept}}{v'})$.
By Corollary~\ref{corr.mru.today.to.yesterday}(\ref{item.mru.tomorrow.options}) either $n'\minus 1\mentval\modT\someoneAll\tf{accept}(v')$ or $n'\minus 1\mentval\modT(\mru{\someoneAll\tf{accept}}{v'})$.

In the former case, using \rulefont{PaxAccept?} $n'\minus 1\mentval\someoneAll\tf{write}(v')$, and by the $\tf{write}$ part of the inductive hypothesis, $v'=v$.
In the latter case, we just use the $\tf{accept}$ part of the inductive hypothesis. 
\item
\emph{The case of $\tf{write}$.}

Suppose $p\in\Pnt$ and $n',p\mentval\modT\tf{write}(v')$.
Using \rulefont{PaxWrite?}, 
there exists an $O\in\opensne$ such that 
$$
n',O\mentval \modT\bigl( 
(\mru{\someone\tf{accept}}{v})\ \tor
(\tneg\modT\recent\texi\someone\tf{accept}\tand \tf{propose}(v)) \bigr),
$$
from which it follows using Lemmas~\ref{lemm.commuting.connectives} and~\ref{lemm.tand.tor} and some basic reasoning on truth-values that 
$$
\text{either}\quad
n',O\mentval \modT\mru{\someone\tf{accept}}{v}
\quad\text{or}\quad
n',O\mentval \tneg\modT\recent\texi\someone\tf{accept} .
$$
Now $n<n'$ and we assumed in equation~\eqref{eq.n.1} that $n\mentval\modT\QuorumBox\tf{accept}(v)$, so using 
\rulefont{Pax2Twined} and \strongmodusponens $n\mentval\modT\CoquorumDiamond\tf{accept}(v)$, and in particular (by routine reasoning on the clauses for $\CoquorumDiamond$ in Figure~\ref{fig.3.derived} and for $\someone$ in Figure~\ref{fig.3.phi.f})
$n,O\mentval\modT\someone\tf{accept}(v)$.
Thus $n',O\mentval\modT\texi\someone\tf{accept}$ and by Lemma~\ref{lemm.tomorrow.forever.recent}(\ref{item.tomorrow.forever.recent.1}) (taking $\phi=\texi\someone\tf{accept}$ in that Lemma) 
$n',O\mentval\modT\recent\texi\someone\tf{accept}$.
It follows that we are not in the right-hand (`or') branch above, so we must be in the left-hand (`either') branch, and thus $n',O\mentval\modT\mru{\someone\tf{accept}}{v}$.

By the $\tf{accept}$ part of the inductive hypothesis, $v'=v$.
\qedhere\end{itemize}

\end{proof}


\begin{thrm}[Agreement]
\label{thrm.full.agreement}
We have:
$$
\ThyPax\ment\tall v,v'.(\urecent\someoneAll\tf{decide}(v) \timpc \someoneAll\tf{decide}(v') \timpc v'\teq v).
$$
\end{thrm}
\begin{proof}
Suppose $v,v'\in\Val$ and $n\in\Time$ and $\avaluation$ is a valuation.
We use \strongmodusponens, so suppose
$$
n\mentval\modT\urecent\someoneAll\tf{decide}(v)
\quad\text{and}\quad
n\mentval\modT\someoneAll\tf{decide}(v') .
$$
If $n\mentval\modT\someoneAll\tf{decide}(v)$ then we use Proposition~\ref{prop.decide.n.agree}.

So suppose $n\mentval\modT\recent\someoneAll\tf{decide}(v)$.
By \rulefont{PaxDecideL?} (and \rulefont{PaxDecide\tneg L?} if needed), $n\mentval\modT\recent\QuorumBox\tf{accept}(v)$ and $n\mentval\modT\QuorumBox\tf{accept}(v')$, so that by Proposition~\ref{prop.someone.implies.someoneAll}(\ref{item.someone.implies.someoneAll.2}) also %
$n\mentval\modT\someoneAll\tf{accept}(v')$.
By Proposition~\ref{prop.integrity.helper}(\ref{item.integrity.helper.accept}) $v'=v$ as required.
\end{proof}

\subsection{Termination}
\label{subsect.pax.termination}

\begin{rmrk}
\emph{Termination} is expressed in English as:
\begin{quote}
\emph{Every correct participant eventually decides some value.}
\end{quote}
We render this in our logic in Theorem~\ref{thrm.termination} as 
$$
\ThyPax\ment\infinitely\everyoneAll\texi v.\tf{decide}(v).
$$
In words: 
\begin{quote}
There is always some future time when all participants decide.
\end{quote}
Note that there is no need to explicitly specify `all correct/uncrashed participants' above.
This is taken care of by our logic's in-built notion of validity, since incorrect participants return $\tvB$, which is valid as per Remark~\ref{rmrk.motivate.validity}(\ref{item.tb.valid}) (but not true or correct; cf. Remark~\ref{rmrk.motivate.validity}).
\end{rmrk}


\subsubsection{Global stabilisation of logical time (GSLT)}
\label{subsect.gslt}

Our forward rules are wrapped in a $\final$ modality, to reflect an assumption that while there may be some initial network instability, if we wait long enough the network will stabilise (see the discussion in Remark~\ref{rmrk.discuss.forward.rules}).
This leads us to the notion of \emph{global stabilisation of logical time}, or \emph{GSLT} for short: 
\begin{defn}[GSLT]
\label{defn.gslt}
Suppose $\phi$ is a closed predicate and $\avaluation$ is a valuation.
Then:
\begin{enumerate*}
\item\label{item.gslt.synchronous}
Say that \ $\final\phi$\ (i.e. $\phi$ under a $\final$ modality) is \deffont{synchronous at $n$ in $\avaluation$} when $n\mentval\forever\phi$.\footnote{%
By Definition~\ref{defn.validity.judgement}, $n\ment\final\phi$ means that $n,p,O\mentval\final\phi$ for every $p\in\Pnt$ and $O\in\opens$ and valuation $\varsigma$.} %
That is:
$$
\final\phi\ \ \text{ is synchronous at $n$ in $\avaluation$}
\quad\text{when}\quad
n\mentval\forever\phi .
$$
In words, $\final\phi$ is synchronous when the `$\final$' becomes a `$\forever$'.
\item\label{item.gslt.gslt}
Define $\gslt(\avaluation)\in\Time$ to be the least stage $n$ such that all of the rules in Figure~\ref{fig.3.phi.f} that start with $\final$ --- namely, the forward rules \rulefont{PaxSend!}, \rulefont{PaxAccept!}, \rulefont{PaxDecideL!}, and \rulefont{PaxDecide\tneg L!} --- are synchronous at $n$ in the valuation $\avaluation$.

We call $\gslt(\avaluation)$ the \deffont{Global Stabilisation of Logical Time} in $\avaluation$. 
\end{enumerate*}
\end{defn}

\begin{rmrk}
We make some observations about Definition~\ref{defn.gslt}:
\begin{enumerate}
\item
Unpacking Definition~\ref{defn.gslt}(\ref{item.gslt.synchronous}) using the clause for $\forever$ in Figure~\ref{fig.3.derived}, 
$$
\final\phi \text{ is synchronous at $n$ in $\avaluation$}
\quad\liff\quad
n\mentval\forever\phi 
\quad\liff\quad
\Forall{n'{\geq}n}(n'\mentval\phi),
$$
\noindent
and $\gslt(\avaluation)$ is least such that $\gslt(\avaluation)\mentval\forever(\someoneAll\texi\tf{propose} \tnotor \texi\tf{send})$, and $\gslt(\avaluation)\mentval\forever\bigl((\tf{leader}\tand \QuorumBox\texi\tf{send}) \tnotor \texi\tf{write}\bigr)$, and so on.
\item  
$\gslt$ is well-defined in Definition~\ref{defn.gslt}(\ref{item.gslt.gslt}) because there are finitely many axioms, mentioning finitely many $\final$ modalities, and a finite set of natural numbers has a finite least upper bound.
This statement has real content, e.g. if we had written \rulefont{PaxAccept!} as $\tall v.\final(\someoneAll\tf{write}(v) \tnotor \tf{accept}(v))$ --- which is a subtly weaker assertion than the form used in Figure~\ref{fig.logical.paxos} --- then $\gslt$ would \emph{not} be defined. 
\end{enumerate}
\end{rmrk}




\begin{lemm}
\label{lemm.someone.writes.after.gslt.accept}
Suppose $\avaluation$ is a valuation such that $\mentval\ThyPax$. 
Suppose $n\in\Time$ and $n\geq\gslt(\avaluation)$ and $v\in\Val$. 
Then 
\begin{multline*}
n\mentval\modT\someoneAll\tf{write}(v) 
\liff
n\mentval\everyoneAll\tf{accept}(v) 
\liff
\\
n\mentval\modT\QuorumBox\tf{accept}(v) 
\liff
n\mentval\modT\someoneAll\tf{accept}(v)  .
\end{multline*}
\end{lemm}
\begin{proof}
We prove a cycle of implications, freely using weak and strong modus ponens (Proposition~\ref{prop.mp.for.tnotor}(\ref{item.mp.for.tnotor}\&\ref{item.mp.for.timpc})).

Suppose $n\mentval\modT\someoneAll\tf{write}(v)$. 
By \rulefont{PaxAccept!} (since $n\geq\gslt(\avaluation)$), $\Forall{p\in\Pnt}n\mentval\tf{accept}(v)$.
By Lemma~\ref{lemm.atomic.O}(\ref{item.atomic.pointwise}) $\tf{accept}(v)$ is pointwise, so by Proposition~\ref{prop.pointwise.dense.char}(\ref{item.dense.char.pointwise.everyoneAll}) $n\mentval\everyoneAll\texi\tf{accept}(v)$.


By Proposition~\ref{prop.correct.cobox.boxT}(\ref{item.correct.cobox.boxT.1}) and \rulefont{PaxPCorrect} (for $\tf{accept}$),
$n\mentval\modT\QuorumBox\tf{accept}(v)$, and by Proposition~\ref{prop.someone.implies.someoneAll}(\ref{item.someone.implies.someoneAll.2}) $n\mentval\modT\someoneAll\tf{accept}(v)$.

Now suppose $n\mentval\modT\someoneAll\tf{accept}(v)$.
By \rulefont{PaxAccept?} %
$n\mentval\modT\someoneAll\tf{write}(v)$.  
\end{proof}


\begin{lemm}
\label{lemm.leaderlive.quorumsend}
Suppose $\avaluation$ is a valuation such that $\mentval\ThyPax$. 
Suppose $n\in\Time$ and $n\geq\gslt(\avaluation)$. 
Then:
\begin{enumerate*}
\item\label{item.leaderlive.quorumsend.1}
$n\mentval\someoneAll\texi\tf{propose}\tnotor\everyoneAll\texi\tf{send}$.
\item\label{item.leaderlive.quorumsend.2}
$n\mentval\someoneAll\texi\tf{propose}\timpc\QuorumBox\texi\tf{send}$.
\end{enumerate*}
\end{lemm}
\begin{proof}
We consider each part in turn:
\begin{enumerate*}
\item
We use \weakmodusponens.
Suppose $n\mentval\modT\someoneAll\texi\tf{propose}$.
By \rulefont{PaxSend!} (since $n\geq\gslt(\avaluation)$) 
$$
\Forall{p\in\Pnt}n,p\mentval\texi\tf{send}.
$$
By Lemma~\ref{lemm.atomic.O}(\ref{item.atomic.pointwise.2}) $\texi\tf{send}$ is pointwise, so by Proposition~\ref{prop.pointwise.dense.char}(\ref{item.dense.char.pointwise.everyoneAll}) $n\mentval\everyoneAll\texi\tf{send}$ as required.
\item
We use \strongmodusponens.
Suppose $n\mentval\modT\someoneAll\texi\tf{propose}$.
By part~\ref{item.leaderlive.quorumsend.1} of this result and \weakmodusponens, $n\mentval\everyoneAll\texi\tf{send}$.
By \rulefont{PaxPCorrect} (for $\tf{send}$) $\acontext\ment\QuorumBox\correct{\tf{send}}$, so by 
Proposition~\ref{prop.correct.cobox.boxT}(\ref{item.correct.cobox.boxT.2})
$\acontext\ment\modT\QuorumBox\texi\tf{send}$ as required. 
\qedhere\end{enumerate*}
\end{proof}




\subsubsection{GSLT implies termination}


Recall from Definition~\ref{defn.correct}(\ref{item.correct}) that $\correct{\tf{decide}}=\tall v.\modTF\tf{decide}(v)$. 
\begin{lemm}
\label{lemm.someone.propose.someone.write}
Suppose $\avaluation$ is a valuation such that $\mentval\ThyPax$. 
Suppose $n\in\Time$ and $n\geq\gslt(\avaluation)$, and $p\in\Pnt$.
Then:
\begin{enumerate*}
\item\label{item.someone.propose.someone.write}
If $n,p\mentval\modT\texi\tf{propose}$ and $n,p\mentval\correct{\tf{write}}$ then $n,p\mentval\modT(\tf{leader}\tand\texi\tf{write})$.
\item\label{item.someone.propose.someone.write.live}
$\mentval\infinitely(\correct{\tf{decide}}\tand\modT\someoneAll\texi\tf{write})$.  
\item\label{item.someone.propose.someone.accept}
$\mentval\infinitely(\someoneAll(\tf{leader}\tand\correct{\tf{decide}})\tand\modT\texi\QuorumBox\tf{accept})$.
\end{enumerate*}
\end{lemm}
\begin{proof}
We consider each part in turn:
\begin{enumerate}
\item
Suppose $n,p\mentval\modT\texi\tf{propose}$.
By \rulefont{PaxPropose?} and \strongmodusponens $n,p\mentval\modT\tf{leader}$.
By Lemma~\ref{lemm.leaderlive.quorumsend}(\ref{item.leaderlive.quorumsend.2}) and strong modus ponens $n\mentval\modT\QuorumBox\texi\tf{send}$.
Using \weakmodusponens and \rulefont{PaxWrite!} (since $n\geq\gslt(\avaluation)$) $n,p\mentval\texi\tf{write}$.
By Proposition~\ref{prop.tv.ment.TF.model}(\ref{item.modTF.texi.P.valid}) (since $n,p\mentval\correct{\tf{write}}$) $n,p\mentval\modT\texi\tf{write}$. 
\item
By \rulefont{LdrExist}, and \rulefont{LdrCorrect} with Lemma~\ref{lemm.persist.char} (for $\f{tvs}=\{\tvT\}$ in that Lemma), and \strongmodusponens, 
for infinitely many $n\in\Time$ there exists $l_n\in\Pnt$ such that 
$$
n,l_n\ment\modT\tf{leader}\tand\correct{\tf{propose},\tf{write},\tf{decide}}. 
$$
So for each $n$ and $l_n$ we have $n,l_n\mentval\correct{\tf{decide}}$.
We will now check that $n,l_n\mentval\modT\texi\tf{write}$.
We reason as follows:
\begin{itemize*}
\item
By \rulefont{PaxPropose!} and \weakmodusponens we have that $n,l_n\ment\texi\tf{propose}$.
\item
By Proposition~\ref{prop.tv.ment.TF.model}(\ref{item.modTF.texi.P.valid}) (since $n,l_n\mentval\correct{\tf{propose}}$ and $n,l_n\mentval\texi\tf{propose}$) $n,l_n\mentval\modT\texi\tf{propose}$. 
\item
By part~\ref{item.someone.propose.someone.write} of this result (since $n,l_n\mentval\correct{\tf{write}}$) $n,l_n\mentval\modT\texi\tf{write}$.
\end{itemize*}
\item
We just combine part~\ref{item.someone.propose.someone.write.live} of this result with Lemma~\ref{lemm.someone.writes.after.gslt.accept}.
\qedhere\end{enumerate}
\end{proof}


Recall that Theorem~\ref{thrm.termination} expresses termination (as the name suggests) but also an abstracted form of the Integrity correctness property by virtue of the use of the $\infinitely$ modality, as per Remark~\ref{rmrk.paxos.integrity}:
\begin{thrm}[Termination]
\label{thrm.termination}
We have:
$$
\ThyPax\ment\infinitely\everyoneAll\texi\tf{decide} .
$$ 
\end{thrm}
\begin{proof}
Suppose $\avaluation$ is a valuation and $\mentval\ThyPax$.
By Lemmas~\ref{lemm.someone.propose.someone.write}(\ref{item.someone.propose.someone.accept}) and~\ref{lemm.persist.char} (for $\f{tvs}=\{\tvT\}$ in that Lemma), for infinitely many $n\in\Time$ there exists an $l_n\in\Pnt$ such that 
\begin{enumerate*}
\item
$n,l_n\mentval\tf{leader}$,
\item
$n,l_n\mentval\correct{\tf{decide}}$ (Definition~\ref{defn.correct}(\ref{item.correct})), and
\item
$n\mentval\modT\texi\QuorumBox\tf{accept}$.
\end{enumerate*}
By \rulefont{PaxDecideL!} (since $n\geq\gslt(\avaluation)$) 
and \weakmodusponens $n,l_n\mentval\texi\tf{decide}$.
By Proposition~\ref{prop.tv.ment.TF.model}(\ref{item.modTF.texi.P.valid}) (since $n,l_n\mentval\correct{\tf{decide}}$ and $n,l_n\mentval\texi\tf{decide}$) $n,l_n\mentval\modT\texi\tf{decide}$. 
Then by \rulefont{PaxDecide\tneg L!} and \weakmodusponens also $n,p\mentval\texi\tf{decide}$ for every $p\in\Pnt$, and so by (Lemma~\ref{lemm.atomic.O}(\ref{item.atomic.pointwise.2}) and) Proposition~\ref{prop.pointwise.dense.char}(\ref{item.dense.char.pointwise.everyoneAll}) $n\mentval\everyoneAll\texi\tf{decide}$. 

The above holds for infinitely many $n\in\Time$, so by Lemma~\ref{lemm.persist.char} (for $\f{tvs}=\{\tvT,\tvB\}$ in that Lemma) we have that $\mentval\infinitely\everyoneAll\texi\tf{decide}$ as required. 
\end{proof}


\begin{rmrk}
Theorem~\ref{thrm.termination} proves that $\ThyPax\ment\infinitely\everyoneAll\texi\tf{decide}$ holds.
Note that $\tvB$ (the truth-value associated to a crashed process) is a valid truth-value, so $n\ment\everyoneAll\texi\tf{decide}$ can hold in the trivial case that everyone is crashed.
Thus, what Theorem~\ref{thrm.termination} asserts in English is that infinitely often, everyone who is not crashed, decides (but it might be that everyone is crashed) --- from which it obviously follows that if a participant never crashes then eventually it decides a value, as per Remark~\ref{rmrk.standard.paxos.correctness}.

But in fact, we can prove a bit more:
$$
\ThyPax\ment\infinitely\QuorumBox\modT\texi\tf{decide}
$$
In words: infinitely often, there exists a \emph{quorum of uncrashed participants}, who decide.

This follows from Theorem~\ref{thrm.termination} by easy reasoning using \rulefont{PaxPCorrect}, \rulefont{PaxDecide\tneg L!}, and Proposition~\ref{prop.decide.n.agree}; we leave this to the reader to check. 
\end{rmrk}


\section{Simpler Declarative Paxos}
\label{sect.simpler.paxos}

Logic is useful for simplification and abstraction in many ways.
In particular, we can axiomatise something, prove some properties --- and then look at what we actually needed from the axioms to make the proofs work.
This can feed back to suggest refinements and generalisations of the original axioms.
(For example: many properties of natural numbers are actually properties of fields, or rings, or integral domains, and so on.) 

We can do something similar here, and look for the most elementary set of axioms that are actually required for our correctness proofs as written: 

\subsection{What we used in the correctness proofs for \ThyPax}
\label{subsect.reflection.on.correctness}
 
\begin{rmrk}
\label{rmrk.spax.explanation}
We note a curious thing about our correctness proofs in Section~\ref{sect.paxos.correctness.properties}, that we never used axioms \rulefont{PaxSend?} or \rulefont{PaxPropose01}.

How can this be?
Is this an error?

No.
This is an opportunity to further simplify the Declarative Paxos from Figure~\ref{fig.ThyPaxOne} to the \emph{Simpler Declarative Paxos} in Figure~\ref{fig.ThySPax}.
To understand how and why this works, we need to highlight some details about how the Paxos algorithm works.
\end{rmrk}

\begin{rmrk}[Why we did not use \rulefont{PaxSend?}]
Recall from Remark~\ref{rmrk.high-level.paxos} our high-level view of the Paxos algorithm, and recall that \ThyPax in Figure~\ref{fig.ThyPaxOne} is a declarative version of that algorithm (obtained by abstracting away explicit message-passing and algorithmic time, as per Remark~\ref{rmrk.elisions}).

Note of the algorithm that 
\begin{itemize*}
\item
in step~\ref{item.l.proposes} of Remark~\ref{rmrk.high-level.paxos}, the leader broadcasts a `propose' message; and 
\item
in step~\ref{item.p.sends}, a participant $p$ that receives that `propose' message, replies with a `send' message carrying as its payload a value $v_p$, which is $p$'s most recently accepted value (if any); and
\item
in step~\ref{item.l.writes}, the leader gathers the $v_p$ values from these responses and, if it has a quorum (in topological language: an open set) of responses, it selects a most recent value to use in its write step.
\end{itemize*}
In the algorithm, the only way that the leader $l$ can know of a participant $p$'s value of $v_p$, is if $l$ receives a send message from $p$ with payload $v_p$.
If $l$ receives no such message, then it does not know the value of $v_p$.

In contrast, predicates of \QLogic are not assigned truth-values by participants in an algorithm; they are assigned truth-values as if by an omniscient observer who is standing outside of the model and who has full (omniscient) access to it.
This is reflected in \rulefont{PaxWrite?}, where we write $\mru{(\someone\tf{accept})}{v}$.
This directly examines the model to find the most recently accepted value.
No reference here is required to the payload of the quorum of $\tf{send}$ messages.
\rulefont{PaxWrite?} still insists that such a quorum exists, as per the $\everyone\texi\tf{send}$ that appears in \rulefont{PaxWrite?}, so in this case the value $\mru{(\someone\tf{accept})}{v}$ is computable by the leader, and we could specify a more complex predicate that checks the set of $v_p$ payloads in the send messages of that quorum to select a most recently-accepted value.
In the implementation, this is necessary because the leader $l$ is not omniscient.
But in the logic we have no reason to go to this trouble, because we can just read the information directly from the model using $\mru{(\someone\tf{accept})}{v}$.\footnote{For example: if we want to compute $100!$ in the real world, we have to build a computer to do this for us, which involves many practical details including the construction of a semiconductor industry, compiler design and so on.  In logic, however, we can just say that it is $\prod_{i=1}^{100}i$, because that is what $100!$ \emph{is}.}

So our axioms are just doing their job, which is to be a declarative abstraction of the implementation. 
But since we do not use \rulefont{PaxSend?}, this is evidence that we may treat the $\tf{send}$ step as an implementational detail having to do with how information moves around the network, and elide it.
\end{rmrk}

\begin{rmrk}[Why we did not use \rulefont{PaxPropose01}]
This is due to another particularity of how the Paxos algorithm works.
The leader chooses one value in the propose step (or no value, if the leader has crashed), but it only uses that value in the case that it finds a quorum of participants who have never previously accepted any value (this is the right-hand disjunct in \rulefont{PaxWrite?}).
Just in this case, it writes the value that it chose in the propose step.

So what really matters is that the leader \emph{only writes one value}.
If the reader had proposed two values, this would be fine --- provided that it only writes one of them.

Leaders in the Paxos algorithm do not bother with multiple values in the propose step, because this would be a waste of time; they know they will only ever need at most one.
Thus, \rulefont{PaxPropose01} is right to reflect a true fact of the algorithm that it does not propose any more than the value than it will need, but this an algorithmic efficiency not a logical necessity. 
Likewise, the logic is right to point this out to us by \emph{not} requiring \rulefont{PaxPropose01} in the proofs: so long as only at most one value gets \emph{written} --- which is handled by \rulefont{PaxWrite01} --- none of our correctness properties are affected.
\end{rmrk}

\begin{rmrk}
By examining the proofs further, some further possible simplifications are revealed:
\begin{enumerate*}
\item
We can combine rules \rulefont{SPaxWrite01} and \rulefont{PaxWrite01} and \rulefont{LdrExt}, as per the discussion in Subsection~\ref{subsect.ldrexistuniq'}, to obtain an elegantly concise rule \rulefont{SPaxWrite01}.
\item
We used \rulefont{PaxPCorrect} in our proofs to obtain $\QuorumBox\correct{\tf P}$ only twice: 
\begin{itemize*}
\item
in Lemma~\ref{lemm.someone.writes.after.gslt.accept} for $\tf P=\tf{accept}$, and 
\item
in Lemma~\ref{lemm.leaderlive.quorumsend} for $\tf P=\tf{send}$.
\end{itemize*}
We have elided the \emph{send} step in $\Theta_\ThySPax$, so we only need to consider $\tf{accept}$.
Accordingly we take take \rulefont{SPaxPCorrect} to be just $\QuorumBox\correct{\tf{accept}}$.
\end{enumerate*}
\end{rmrk}

\begin{rmrk}
In conclusion, we see that the proofs in Section~\ref{sect.paxos.correctness.properties} are not only correctness proofs, but can also be read as an active commentary on the axioms.
Our logic has done what logic does, and has separated what is logically necessary from what is implementationally required or implementationally convenient.
This shows us how to simplify the axiomatisation (if we we wish) further away from the implementation, and closer towards what is logically essential: 
\end{rmrk}


\subsection{The definition}

\begin{figure}
$$
\begin{array}{l@{\quad}r@{\ }l}
\figunderline{Backward rules, for backward inference} 
\figskip
\rulefont{SPaxPropose?}&
& \tall v.\ \tf{propose}(v)\timpc\tf{leader}
\figskip
\rulefont{SPaxWrite?}&
& \tall v.\ \tf{write}(v) \timpc 
\begin{array}[t]{l}
\tf{leader}\tand\Quorum( 
\\
\quad\bigl((\mru{(\someone\tf{accept})}{v}) \tor
(\tf{propose}(v)\tand\tneg\modT\recent\texi\someone\tf{accept})\bigr)) 
\end{array}
\figskip
\rulefont{SPaxAccept?}&
& \tall v.\ \tf{accept}(v) \timpc \someoneAll\tf{write}(v)
\figskip
\rulefont{SPaxDecide?}&
& \tall v.\ \tf{decide}(v) \timpc \QuorumBox\tf{accept}(v)
\figskip\figskip
\figunderline{Forward rules, for making progress} 
\figskip
\rulefont{SPaxWrite!}&
& \final(\tf{leader}\tnotor\texi\tf{write})
\figskip
\rulefont{SPaxAccept!}&
& \final(\texi\someoneAll\tf{write}\tnotor \texi\tf{accept})
\figskip
\rulefont{SPaxDecide!}&
& \final(\texi\QuorumBox\tf{accept} \tnotor \texi\tf{decide})
\figskip
\figskip
\figunderline{Other rules, for correctness \& liveness}
\figskip
\rulefont{SLdrExist}&
&\modTF\tf{leader}\tand \modT\someoneAll\tf{leader}
\figskip
\rulefont{SLdrCorrect}&
& \infinitely\someoneAll(\tf{leader}\tand \correct{\tf{propose},\tf{write},\tf{decide}}) 
\figskip
\rulefont{SPaxPCorrect}&
& \QuorumBox\correct{\tf{accept}}
\figskip
\rulefont{SPaxWrite01}&
& \texiaffine\someoneAll\tf{write}
\end{array}
$$
\caption{Simpler Declarative Paxos (Definition~\ref{defn.thyspax})}
\label{fig.logical.paxos'}
\label{fig.ThySPax}
\end{figure}


\begin{defn}
\label{defn.thyspax}
We define \deffont{Simpler Declarative Paxos} to be the logical theory 
$$
\ThySPax=(\Sigma_{\ThySPax},\Theta_{\ThySPax}),
$$ 
where:
\begin{enumerate*}
\item
the signature $\Sigma_{\ThySPax}$ is 
$$
\Sigma_{\ThySPax}=[\tf{leader}:0;\ \tf{propose},\tf{write},\tf{accept},\tf{decide}:1] ,
$$ 
as per the notation in Definition~\ref{defn.signature}(\ref{item.predicate.syntax.signature}), and
\item
the axioms $\Theta_{\ThySPax}$ are as written in Figure~\ref{fig.ThySPax}.
\end{enumerate*}
\end{defn}

The same validity properties hold of \ThySPax from Figure~\ref{fig.ThySPax} as hold for \ThyPax from Figure~\ref{fig.ThyPaxOne}, namely: Validity (Theorem~\ref{thrm.validity}), Agreement (Theorem~\ref{thrm.full.agreement}), and Termination and Integrity (Theorem~\ref{thrm.termination}):
\begin{thrm}
\label{thrm.simpler.dp}
Suppose $\Sigma$ is a signature that includes $\Sigma_\ThySPax$, and suppose $\mathcal M$ is a model.
Then we have:
\begin{enumerate*}
\item
$\ThyPax\ment\tall v.(\tf{decide}(v) \timpc \urecent\someoneAll(\tf{leader}\tand\tf{propose}(v)))$.
\item
$\ThyPax\ment\tall v,v'.(\urecent\someoneAll\tf{decide}(v) \timpc \someoneAll\tf{decide}(v') \timpc v'\teq v)$.
\item
$\ThyPax\ment\infinitely\everyoneAll\texi\tf{decide}$.
\end{enumerate*}
\end{thrm}
\begin{proof}
This is by construction: from the discussion in Subsection~\ref{subsect.reflection.on.correctness}, the theory \rulefont{ThySPax} consists of just those parts of \rulefont{ThyPax} that are actually required in the correctness proofs.
\end{proof}

Thus, Theorem~\ref{thrm.simpler.dp} just notes that Simpler Declarative Paxos extracts from the proofs in Section~\ref{sect.paxos.correctness.properties} the logical essence of Figure~\ref{fig.ThyPaxOne}, as applied to the standard correctness properties of Paxos.
There is no right or wrong about whether we prefer $\ThyPax$ or $\ThySPax$; it depends on what level of abstraction we want to work with at any particular moment.
But, implicit in Theorem~\ref{thrm.simpler.dp} is an observation that \ThySPax is \emph{minimal}, because we put into it only that which we know we need to make some specific correctness proofs work.


\section{Conclusions and future work}
\label{sect.conclusions}

\subsection{Conclusions}







We have presented \QLogic, a three-valued modal fixedpoint logic for declarative specification of consensus algorithms, and we have shown how to use this to axiomatise Paxos, a canonical and nontrivial (though still relatively straightforward) consensus algorithm that can handle benign failures of participants.

We do not presented more complex consensus algorithms in this paper but we have looked into these in preliminary investigations, and they do seem to work. 

This paper is aimed at 
\begin{itemize*}
\item
logicians who might be interested in exploring a logic with a new and unusual combination of features inspired by strong practical motivations, and 
\item
researchers in distributed systems, who might consider (and benefit from) adopting such a framework.
\end{itemize*}
\QLogic offers benefits that could make it a useful addition to the toolkit of anyone working in distributed systems:
\begin{enumerate}
\item 
Much of what is shown in this paper just reflects existing practices, formalised into a new logic.
Therefore, readers may already be reasoning in ways that align with our logic, without explicitly realising it. 
For example, the third truth-value ($\tvB$) in \QLogic corresponds to typical side conditions like ``for every correct participant,'' and the backward axioms mirror the reverse inferences that practitioners often use. 

By formalising these familiar patterns, \QLogic provides a structured and precise way to represent and reason about concepts that are already informally expressed in the literature.
\item 
Logic is a concise, powerful, and unambiguous language. 
It distills paragraphs of English prose into a single line of precise symbols, enabling clear and rigorous reasoning. 

While explanatory English is helpful for accessibility and context, formal logic is the superior tool for working through the details of complex algorithms. 
It ensures precision and eliminates the ambiguities that often arise in informal descriptions.
\item 
\QLogic operates at a new (and higher) level of abstraction, which complements existing logic-based techniques. 
TLA+~\cite{DBLP:journals/toplas/Lamport94} provides a state-based view of distributed algorithms, whereas \QLogic abstracts lower-level implementation details while retaining the essential behaviours and properties of the algorithms.
This abstraction does not obscure the logic of the system; on the contrary we have found that it clarifies it, making the underlying mechanisms and invariants more apparent. 
\item 
One strength of logic is its ability to abstract away what is irrelevant while maintaining precision about what is important. 
In \QLogic, constructs like strong and weak implications ($\timpc$ and $\tnotor$) exemplify this principle, capturing subtle distinctions in meaning with precise formal symbols. 
These constructs distill complex ideas into simple yet powerful tools that can be manipulated symbolically and reasoned about systematically. 

While mastering these constructs requires some initial effort, the benefits of precision and clarity justify the investment.
\item 
Learning \QLogic requires some effort, but for researchers already familiar with (for example) TLA+, the transition to \QLogic should be eminently manageable. 

While axiomatisations in \QLogic are certainly subtle, this is not gratuitous: they reflect the inherent nature of the distributed systems they are designed to model. 
By formalising these complexities, \QLogic captures their structure and reveals and addresses the challenges that are \emph{already present} in the underlying algorithms. 

This transparency makes it a powerful tool for understanding and reasoning about distributed protocols.\footnote{Certainly, the first author (who has a background in logic) has found this study invaluable for gaining a better understanding of consensus algorithms.  It seems likely that other readers might benefit similarly from using \QLogic.}
\item 
Distributed algorithms, and consensus protocols in particular, are foundational to the functionality of blockchain systems and other critical infrastructures. 

These algorithms are recognised to be notoriously challenging to design and implement correctly. 
Ambiguities and informal specifications can lead to vulnerabilities and failures. 

\QLogic bridges the gap between declarative specifications of invariants, and practical verification; enabling researchers to clarify, reason about, and verify the properties of distributed protocols rigorously. 
This is especially valuable in blockchain systems, where clear and robust specifications are essential for trust and reliability.
\end{enumerate}

In conclusion: by providing a formal framework that combines the abstraction of declarative reasoning with the rigour of logical precision, \QLogic offers researchers in distributed systems a powerful tool for modeling, understanding, and verifying complex algorithms. 

\subsection{Related and future work}

In this paper, our focus has been to lay the foundational theory and framework of \QLogic and to demonstrate its application to one of the most well-known consensus protocols, Paxos~\cite{DBLP:journals/tocs/Lamport98}.

Paxos is a sensible starting point, being a cornerstone of distributed systems, but (certainly by modern standards!) it is relatively straightforward.
For example, Paxos only addresses crash failures, limiting its applicability to systems that assume benign (rather than hostile) faults.

A natural direction for future work involves extending our framework to handle more complex and realistic fault models, such as Byzantine faults, where participants can deviate arbitrarily from the algorithm.
Consensus protocols that address Byzantine faults, such as Practical Byzantine Fault Tolerance (PBFT)~\cite{DBLP:conf/osdi/CastroL99} and HotStuff~\cite{DBLP:conf/podc/YinMRGA19}, provide robust mechanisms for fault tolerance and serve as the foundation for many modern blockchain systems, including Tendermint (\url{https://tendermint.com})~\cite{DBLP:journals/corr/abs-1807-04938}. 

Applying \QLogic to these protocols would see how it handles the complexity introduced by arbitrary faults, which is necessary to model and verify more sophisticated algorithms.
Preliminary exploration by the authors in this direction indicates that \QLogic works very well in these more complex situations --- it can indeed axiomatise algorithms that take account of arbitrary, including hostile, faults.
Unpacking this in full is for future papers.

Recent research has explored DAG-based consensus protocols, which deviate from the traditional linear blockchain structure. 
In these protocols, transactions are organized into a directed acyclic graph (DAG) rather than a strict chain, enabling greater parallelism and scalability. 
Protocols such as DAG-rider~\cite{DBLP:conf/podc/KeidarKNS21}, Narwhal and Tusk~\cite{DBLP:conf/eurosys/DanezisKSS22}, and Bullshark~\cite{DBLP:conf/ccs/SpiegelmanGSK22} represent cutting-edge advancements in this domain. 
Extending \QLogic to analyse DAG-based protocols would be another advance, as these systems often involve more intricate interactions between components and require reasoning about concurrent and asynchronous behaviour at a much higher level of complexity.

Another promising avenue for future research involves leveraging the modularity and expressiveness of \QLogic to study hybrid protocols~\cite{DBLP:conf/sp/NeuTT21} that combine elements of classical consensus with newer techniques, such as Gasper~\cite{DBLP:journals/corr/abs-2003-03052} and others~\cite{DBLP:conf/esorics/DAmatoZ23,DBLP:journals/corr/abs-2411-00558}

It would be helpful in future work to develop automated tools for \QLogic: the tools developed for TLA+~\cite{DBLP:books/aw/Lamport2002} could show the way.
TLA+ is a formal specification language and framework designed to model, specify and verify properties of complex systems; particularly distributed algorithms. 
Its TLC model checker~\cite{DBLP:conf/charme/YuML99} provides exhaustive state exploration to verify properties such as safety and liveness, and the PlusCal translator simplifies specification by offering a high-level pseudocode-like interface.
Additionally, the TLA+ Toolbox~\cite{DBLP:journals/corr/abs-1912-10633} integrates these capabilities into a development environment, supporting workflows that are valuable in both research and industrial applications.

Adapting similar tools to \QLogic could enable automated verification of correctness properties for consensus protocols, including invariants and liveness guarantees. 
For example: a higher-level language akin to PlusCal could simplify the specification of protocols using the axiomatic foundations of \QLogic.
Furthermore, leveraging techniques such as the symbolic model checking enhancements presented by Otoni et al.~\cite{DBLP:conf/tacas/OtoniKKES23} or integrating modular proof systems like TLAPS~\cite{DBLP:conf/fm/CousineauDLMRV12} could further enhance the usability and robustness of the framework.

Such tools would streamline the application of \QLogic to real-world systems, and also make it accessible to a broader audience of researchers and practitioners. 
This direction represents an exciting convergence of theoretical advancements and practical tool development, providing a robust framework for analysing both classical and emerging consensus protocols.

Adopting \QLogic requires some investment in time and effort, but the benefits in clarity, correctness, and reliability may be worthwhile.
Consensus algorithms are known to be particularly challenging, so our relatively high-level declarative approach could help to manage (and debug) the complexity of these algorithms.

We hope this work will inspire readers %
to explore and adopt this new perspective, contributing to more robust and trustworthy distributed protocols in the future.














\documentclass[acmjacm]{acmsmall}

\usepackage[shortlabels]{enumitem}
 
\usepackage{latexsym,xspace,calc,amsthm}
\usepackage{amsmath,amssymb} %
\usepackage{centernot}

\setcounter{tocdepth}{4}
\setcounter{secnumdepth}{3}


\usepackage{lmodern}
\usepackage{tgtermes}
\usepackage{fix-cm}


\usepackage{cancel}

\usepackage{savesym}
\savesymbol{subcaption}
\usepackage{subcaption}

\usepackage{graphbox} %

\makeatletter
\def\@biblabel#1{[#1]} %
\def\thebibliography#1{%
    \footnotesize
    \refsection*{{\refname}
        \@mkboth{\uppercase{\refname}}{\uppercase{\refname}}%
    }
    \list{\@biblabel{\@arabic\c@enumiv}}%
       {\settowidth\labelwidth{\@biblabel{#1}}%
        \leftmargin\labelwidth
        \advance\leftmargin\bibindent
        \itemindent-\bibindent
        \itemsep2pt
        \parsep \z@
        \usecounter{enumiv}%
        \let\p@enumiv\@empty
        \renewcommand\theenumiv{\@arabic\c@enumiv}%
    }%
    \let\newblock\@empty
    \sloppy
    \sfcode`\.=1000\relax
}
\makeatother

\let\odlthebibliography=\thebibliography
 \let\endodlthebibliography=\endthebibliography
 \renewenvironment{thebibliography}[1]{%
   \begin{odlthebibliography}{#1}%
     \setlength{\parskip}{0ex}%
     \setlength{\itemsep}{3pt}%
     \fontsize{9.5}{9.5} %
     \selectfont
}%
 {%
   \end{odlthebibliography}%
 }


\usepackage{tikz}
\usepackage{tikz-cd}
\usetikzlibrary{cd,decorations.markings}
\tikzset{
    dharrow/.style={
        <->,
        postaction={decorate,-},
        }
}
\tikzset{
    dhdashedarrow/.style={
        <->,
        dashed,
        postaction={decorate,-},
        }
    }
\tikzset{
    lrharpoonarrow/.style={
        <[harpoon]->[harpoon],
        postaction={decorate,-},
        }
}
\tikzset{
    lrharpoondashedarrow/.style={
        <[harpoon]->[harpoon],
        dashed, %
        postaction={decorate,-},
        }
}
\usetikzlibrary{arrows}
\usetikzlibrary {arrows.meta} 
\usetikzlibrary{calc}
\usetikzlibrary{positioning}
\usetikzlibrary{snakes,automata,chains}
\usetikzlibrary{graphs}

\usepackage{binarytree}

\usepackage{amssymb,stmaryrd,amsmath}
\usepackage{mdwlist} %
\usepackage{float}   %
\usepackage{centernot} %


\usepackage[colorlinks]{hyperref}
\usepackage{breakurl}  %


\usepackage{bm} %





\newtheoremstyle{jamiestyle}%
  {4pt}%
  {0pt}%
  {\it}%
  {0pt}%
  {\sc}%
  {.}%
  { }%
  {}%
\theoremstyle{jamiestyle}
\newtheorem{thrm}{Theorem}[subsection]
\newtheorem{prop}[thrm]{Proposition}
\newtheorem{lemm}[thrm]{Lemma}
\newtheorem{corr}[thrm]{Corollary}
\newtheorem{conj}[thrm]{Conjecture}
\newtheoremstyle{jamienfstyle}%
  {4pt}%
  {0pt}%
  {\normalfont}%
  {0pt}%
  {\sc}%
  {.}%
  { }%
  {}%
\theoremstyle{jamienfstyle}
\newtheorem{nttn}[thrm]{Notation}
\newtheorem{defn}[thrm]{Definition}
\newtheorem{xmpl}[thrm]{Example}
\newtheorem{rmrk}[thrm]{Remark}


\usepackage{color}
\definecolor{mygreen}{rgb}{0,0.6,0}
\definecolor{mygray}{rgb}{0.5,0.5,0.5}
\definecolor{mymauve}{rgb}{0.58,0,0.82}

\usepackage{listings}
 
\definecolor{gray}{RGB}{128, 128, 128}
\definecolor{lightgray}{RGB}{200, 200, 200}
\definecolor{cyan}{RGB}{0, 255, 255}
\definecolor{blue}{RGB}{0, 0, 255}
\definecolor{red}{RGB}{255, 0, 0}
\definecolor{pink}{RGB}{255, 128, 128}
\definecolor{green}{RGB}{0, 128, 0}
\definecolor{lightyellow}{RGB}{255, 255, 200}
\definecolor{purple}{RGB}{128, 0, 128}

\lstdefinestyle{all}
    {basicstyle=\ttfamily\scriptsize,
     keywordstyle=\color{blue}\ttfamily\scriptsize,
     commentstyle=\color{green}\ttfamily\scriptsize,
     stringstyle=\color{red}\ttfamily\scriptsize}

\lstdefinelanguage{hask}{%
    frame=none,
    xleftmargin=2pt,
    belowcaptionskip=\bigskipamount,
    captionpos=b,
    tabsize=2,
    numbers=left,
    numberstyle=\tiny\color{gray},
    emphstyle={\bf},
	morecomment=[s][\color{green}]{\{-}{-\}},
    stringstyle=\mdseries\rmfamily,
    commentstyle=\color{green},
    keywords={},
    keywords=[1]{case, of, data, if, then, else, where, let, in, do},
    keywords=[2]{Chip, Config, CurrencySymbol, TokenName, PubKeyHash, Integer, Value, State, Action, Text, Maybe, Void, TxConstraints,  Contract},
    keywords=[3]{HasNative},
    keywords=[4]{=>},
    keywords=[5]{Just, Nothing, MkChip, MkConfig, SetPrice, Buy},
    keywordstyle=[1]\mdseries\sffamily\color{red},
    keywordstyle=[2]\mdseries\sffamily\color{blue},
    keywordstyle=[3]\mdseries\sffamily\color{green},
    keywordstyle=[4]\mdseries\sffamily,
    keywordstyle=[5]\mdseries\sffamily\color{purple},
    columns=flexible,
    basicstyle=\small\sffamily,
    showstringspaces=false,
    breaklines=false,
    showspaces=false,
    escapeinside={--}{\^^M},escapebegin={\color{green}--},escapeend={},
    literate= {+}{{$+$}}1 {/}{{$/$}}1 {*}{{$*$}}1 {=}{{$=$}}1
              {>}{{$>$}}1 {<}{{$<$}}1 {\\}{{$\lambda$}}1
              {\\\\}{{\char`\\\char`\\}}1
              {->}{{$\rightarrow$}}2 {>=}{{$\geq$}}2 {<-}{{$\leftarrow$}}2
              {<=}{{$\leq$}}2 {=>}{{$\Rightarrow$}}2
              {\ .}{{$\circ$}}2 {\ .\ }{{$\circ$}}2
              {>>}{{>>}}2 {>>=}{{>>=}}2
              {|}{{$\mid$}}1
              {\_}{{\underline{\hspace{2mm}}}}2
}

\lstdefinelanguage{solidity}{%
    frame=none,
    xleftmargin=2pt,
    belowcaptionskip=\bigskipamount,
    captionpos=b,
    tabsize=2,
    numbers=left,
    numberstyle=\tiny\color{gray},
    emphstyle={\bf},
	morecomment=[s][\color{green}]{\{-}{-\}},
    stringstyle=\mdseries\rmfamily,
    commentstyle=\color{green},
    keywords={},
    keywords=[1]{pragma, solidity, contract, event, constructor, require, function, return, emit},
    keywords=[2]{address, uint, mapping},
    keywords=[3]{public, payable, external, view, returns},
    keywords=[4]{=>, +=, -=, =, <=, ==},
    keywords=[5]{msg, sender, transfer, value},
    keywordstyle=[1]\mdseries\sffamily\color{red},
    keywordstyle=[2]\mdseries\sffamily\color{blue},
    keywordstyle=[3]\mdseries\sffamily\color{green},
    keywordstyle=[4]\mdseries\sffamily,
    keywordstyle=[5]\mdseries\sffamily\color{purple},
    columns=flexible,
    basicstyle=\small\sffamily,
    showstringspaces=false,
    breaklines=false,
    showspaces=false,
    escapeinside={--}{\^^M},escapebegin={\color{green}--},escapeend={},
    literate= {+}{{$+$}}1 {/}{{$/$}}1 {*}{{$*$}}1 {=}{{$=$}}1
              {>}{{$>$}}1 {<}{{$<$}}1 {\\}{{$\lambda$}}1
              {\\\\}{{\char`\\\char`\\}}1
              {->}{{$\rightarrow$}}2 {>=}{{$\geq$}}2 {<-}{{$\leftarrow$}}2
              {<=}{{$\leq$}}2 {=>}{{$\Rightarrow$}}2
              {\ .}{{$\circ$}}2 {\ .\ }{{$\circ$}}2
              {>>}{{>>}}2 {>>=}{{>>=}}2
              {|}{{$\mid$}}1
              {\_}{{\underline{\hspace{2mm}}}}2
}




\newcommand\flanks{\ltimes}
\newcommand\setminusskip{\hskip 2pt plus 1pt minus 2pt}
\newcommand\mysetminus{\setminusskip{{\setminus}}\setminusskip}

\newcommand\declaresoundness[2]{The case of \rulefont{#1} with $#2$.}
\newcommand\declaresoundnessshort[1]{The case of \rulefont{#1}.}


\makeatletter
\newcommand\hpn[2][]{%
  \ext@arrow 9999{\hpnfill@}{#1}{#2}}
\newcommand\hpnfill@{%
  \arrowfill@\leftharpoonup\relbar\rightharpoondown}
\makeatother

\newcommand\arrowlength{0.5}
\newcommand\arrowr{{\tikz \draw [semithick,->] (0,0) -- (\arrowlength*0.8,0);}}
\newcommand\Arrowr{{\tikz \draw [semithick,->,decorate,decoration={zigzag,segment length=3,amplitude=0.5}] (0,0) -- (\arrowlength*0.8,0);}}
\newcommand\Arrowrr{{\tikz \draw [semithick,->>,decorate,decoration={zigzag,segment length=3,amplitude=0.5}] (0,0) -- (\arrowlength*0.8,0);}}
\newcommand\arrowlr{{\tikz \draw [semithick,<->] (0,0) -- (\arrowlength*0.8,0);}}
\newcommand\arrowllrr{{\tikz \draw [semithick,<<->>] (0,0) -- (\arrowlength,0);}}
\newcommand\arrowrr{{\tikz \draw [semithick,->>] (0,0) -- (\arrowlength*0.8,0);}}
\newcommand\arrowrrr{{\tikz \draw [thick,-Stealth] (0,0) -- (\arrowlength*0.8,0);}}
\newcommand\arrowll{{\tikz \draw [semithick,<<-] (0,0) -- (\arrowlength*0.8,0);}}
\newcommand\arrowerr{{\tikz \draw [semithick,-.>>] (0,0) -- (\arrowlength*0.8,0);}}
\newcommand\arrowlrr{{\tikz \draw [semithick,<->>] (0,0) -- (\arrowlength,0);}}
\newcommand\arrowllr{{\tikz \draw [semithick,<<->] (0,0) -- (\arrowlength,0);}}
\newcommand\arrowelrr{{\tikz \draw [semithick,<.->>] (0,0) -- (\arrowlength,0);}}
\newcommand\arrowelr{{\tikz \draw [semithick,<.->] (0,0) -- (\arrowlength,0);}}
\newcommand\arrower{{\tikz \draw [semithick,-.>] (0,0) -- (\arrowlength*0.8,0);}}


\newcommand\everyone{\modality{\Box}}
\newcommand\someone{\modality{\Diamond}}
\newcommand\someoneAll{{\mathring{\exists}_{\hspace{-0pt}\someone}}}
\newcommand\everyoneAll{{\mathring{\forall}_{\hspace{-1pt}\everyone}}}

\newcommand\Quorum{\modality{\mathring{\exists}}}
\newcommand\QuorumBox{\modality{\mathring{\exists}_{\Box}}} %
\newcommand\Coquorum{\modality{\mathring{\forall}}}
\newcommand\CoquorumDiamond{\modality{\mathring{\forall}_{\hspace{-1pt}\Diamond}}} %


\newcommand{\ite}[2]{[#1,\ \text{\it o\hspace{-1pt}'wise}\ #2]}
\newcommand\mruarrow{\mathbin{\rlap{\scalebox{0.73}{$\recentsimple$}}\hspace{-2pt}\rightarrow}}
\newcommand{\mru}[2]{#1\mruarrow#2}
\newcommand{\mrup}[3]{#2[#1\mruarrow #3]}
\newcommand{\mruX}[2]{{#1\mathbin{\recent}#2}}
\newcommand{\established}[2]{#1 \mathbin{\Downarrow_{\QuorumBox}} #2}
\newcommand{\lrharpoon}{\mathbin{\hpn{}}}
\newcommand{\nlrharpoon}{\mathbin{\not\hpn{}}}
\newcommand\id{\f{id}}
\newcommand\st{{:}}

\usepackage{tikz}
\newcommand*{\myhash}{%
  \begin{tikzpicture}
    \pgfmathsetlengthmacro\myWidth{.8*width("=")}%
    \pgfmathsetlengthmacro\myHeight{height("H")}%
    \pgfmathsetlengthmacro\mySepY{.3333*\myWidth}%
    \pgfmathsetlengthmacro\mySideBearing{.1*\myWidth}%
    \def\myAngle{70}%
    \pgfmathsetlengthmacro\mySepX{\mySepY/sin(\myAngle)}%
    \pgfmathsetlengthmacro\mySlantX{\myHeight/tan(\myAngle)}%
    \draw[line cap=round]
      (0, {(\myHeight - \mySepY)/2}) -- ++(\myWidth, 0)
      (0, {(\myHeight + \mySepY)/2}) -- ++(\myWidth, 0)
      ({(\myWidth - \mySepX - \mySlantX)/2}, 0)
      -- ({(\myWidth - \mySepX + \mySlantX)/2}, \myHeight)
      ({(\myWidth + \mySepX - \mySlantX)/2}, 0)
      -- ({(\myWidth + \mySepX + \mySlantX)/2}, \myHeight)
    ;%
    \useasboundingbox
      (-\mySideBearing, 0)
      (\myWidth + \mySideBearing, \myHeight)
    ;%
  \end{tikzpicture}%
}

\newcommand\THREE{{\mathbf 3}}
\newcommand\threeValid{\THREE_{\f{valid}}}
\newcommand\threeCorrect{\THREE_{\f{correct}}}
\newcommand\threeTrue{\THREE_{\f{true}}}
\newcommand\threeFalse{\THREE_{\f{false}}}
\newcommand\binaryconnectives{\ensuremath{\{\tand,\tor,\tnotor,\timpc\}}} %
\newcommand\unaryconnectives{\ensuremath{\{\tneg,\modT,\modB,\modF,\modTB,\modTF\}}}
\newcommand\modalityconnectives{\ensuremath{\{\someone,\everyone,\Quorum,\Coquorum,\sometime,\forever,\tomorrow\}}} %
\newcommand\doublemodalities{\ensuremath{\{\QuorumBox,\CoquorumDiamond,\someoneAll,\everyoneAll\}}}
\newcommand\singlemodalities{\ensuremath{\{\someone,\everyone,\Quorum,\Coquorum\}}}
\newcommand\Pnt{\ns{Point}}
\newcommand\Val{\ns{Value}}
\newcommand\Valuation{\ns{Valuation}}
\newcommand\PredSymb{\ns{PredSymb}}

\newcommand\FOUR{{\mathbf 4}}
\newcommand\tvT{{\mathbf t}}
\newcommand\tvF{{\mathbf f}}
\newcommand\tvB{{\mathbf b}}
\newcommand\tvN{{\mathbf n}}
\newcommand\tvsTB{{\mathit{tb}}}
\newcommand\tvsFB{{\mathit{fb}}}
\newcommand\tvsTT{{\mathit{tt}}}
\newcommand\tvsFF{{\mathit{ff}}}
\newcommand*\Rcirc{\ensuremath{\smash{\overset{\circ}{\smash{\mathsf{R}}\rule{0pt}{1.2ex}}\vphantom{\mathbb{R}}}}}
\newcommand\recentsimple{{\reflectbox{$\mathsf R$}}} 
\newcommand\recent{\recentsimple} %
\newcommand\urecent{{\ubar\recentsimple}} %
\newcommand\mrecent[2]{#1\recentsimple #2}

\newcommand\figunderline[1]{\text{\rlap{\underline{\emph{#1}}}}}




\usepackage{stackengine}
\newcommand\ubar[1]{\stackunder[0.6pt]{$#1$}{\rule{1ex}{.18ex}\hspace{0.5pt}}}


\newcommand\figskip{\\[1ex]}
\newcommand\from[1]{{\uparrow}_{#1}}
\newcommand\fs[1]{\mathtt{#1}}
\newcommand\ps[1]{\mathtt{#1}}
\newcommand\thea{a}
\newcommand\theb{b}
\newcommand\thec{c}
\newcommand\aval{v}
\newcommand\xor{\mathbin{\mathsf{\small xor}}}
\newcommand\modality[1]{#1\hspace{1.5pt}}
\newcommand\modT{\modality{\tf T}}
\newcommand\modTB{\modality{\tf T\hspace{-2.5pt}\tf B}}
\newcommand\modFB{\modality{\tf F\hspace{-1.5pt}\tf B}}
\newcommand\modB{\modality{\tf B}}
\newcommand\modTF{\modality{\tf T\hspace{-2.5pt}\tf F}}
\newcommand\modF{\modality{\tf F}}
\newcommand\witno{\prec}
\newcommand\atopen{T}
\newcommand\afilter{F}
\newcommand\apoint{P}
\newcommand\acontext{\f{ctx}}
\newcommand\avaluation{\varsigma} %
\newcommand\indicator[1]{\delta_{#1}} 
\newcommand\charop[1]{\f{char}(#1)}
\newcommand\resilient{R}
\newcommand\rulefont[1]{\ensuremath{{\mathsf{(#1)}}}}
\newcommand\leftopeninterval[1]{(#1]}
\newcommand\rightopeninterval[1]{[#1)}
\newcommand\openinterval[1]{(#1)}
\newcommand\closedinterval[1]{[#1]}
\newcommand\opens{\ns{Open}}
\newcommand\opensne{\opens_{\hspace{-1pt}{\neq}\varnothing}}
\newcommand\regularOpens{\tf{Open}_{\f{reg}}}
\newcommand\topens{\tf{Topen}}
\newcommand\closed{\tf{Closed}}
\newcommand\regularClosed{\tf{Closed}_{\f{reg}}}
\newcommand\thenext[0]{\f{ad}}
\newcommand\closure[1]{|#1|}
\newcommand\dotclosure[1]{\closure{#1}^{\bigcdot}}
\newcommand{\dotarrow}{%
   \mathrel{\ooalign{\hss\raise.85ex\hbox{\scalebox{1.25}{\normalfont .}}%
   \kern0.35ex\hss\cr$\rightarrow$}}}
\newcommand{\dotleftarrow}{\mathrel{\reflectbox{$\dotarrow$}}}
\newcommand\mycard{\myhash}
\newcommand\narrowmath[1]{\scalebox{0.8}[1.0]{$\mathsf{#1}$}}
\newcommand\sometime{\modality{\narrowmath{sometime}}}
\newcommand\forever{\modality{\narrowmath{forever}}}
\newcommand\tomorrow{\modality{\narrowmath{tomorrow}}}
\newcommand\yesterday{\modality{\narrowmath{yesterday}}}
\newcommand\infinitely{\modality{\narrowmath{infinitely}}}
\newcommand\final{\modality{\narrowmath{finally}}}
\newcommand\arity{\f{arity}}
\newcommand\correct[1]{\tf{correct}[#1]}
\newcommand\correctterm{correct\xspace}
\newcommand\pointwise[1]{\tf{pointwise}[#1]}
\newcommand\pointwiseterm{pointwise\xspace}
\newcommand\incorrect[1]{\tf{incorrect}[#1]}
\newcommand\uniform[1]{\tf{uniform}[#1]}
\newcommand\udfn{\tf{undef}}
\newcommand\gslt{\f{GSLT}}
\newcommand\dense{\f{dense}}
\newcommand\noi{\f{noi}}

\newcommand{\upperset}[2]{
  \underset{
            \text{\raisebox{0.9ex}{\smash{\scalebox{0.8}{$#1$}}\ \ }}%
           }
           {
            \text{\raisebox{0ex}{\smash{$#2$}}}
           }
                       }
\newcommand{\convergesto}[1]{\mathrel{\upperset{#1}{\dotarrow}}}
\newcommand\betwixt{\asymp}
\newcommand\notbetween{\mathbin{\cancel{\between}}}
\newcommand\notintertwinedwith{\mathrel{\notbetween}}
\newcommand\intersectswith{\between}
\newcommand\notintersectswith{\notbetween}
\newcommand\stronglybetween{\mathbin{\rlap{$\between$}\hspace{-1.5pt}\between}}
\newcommand\xxtwined{$2$-twined\xspace}
\newcommand\xxxtwined{$3$-twined\xspace}
\newcommand\intertwined[1]{#1_{\between}}
\newcommand\intertwinedwith{\mathrel{\between}}
\newcommand\stronglyintertwinedwith{\mathrel{\stronglybetween}}
\newcommand\leqk{\leq_{\hspace{-.7pt}\intertwinedwith}}
\newcommand\geqk{\geq_{\hspace{-.7pt}\intertwinedwith}}
\newcommand\cw{\leftrightarrow}  %
\newcommand\ctr{\rightarrow}  %
\newcommand\notctr{{\not\rightarrow}}  %
\newcommand\ctl{\leftarrow}  %
\newcommand\cti{\leq}  %
\newcommand\ctieq{{\stackrel{=}{\hookrightarrow}}}  %
\newcommand\notcti{{\not\hookrightarrow}}  %
\newcommand\nbhd[0]{\f{nbhd}}
\newcommand\interior[0]{\f{interior}}
\newcommand\kiss[0]{\f{kiss}}
\newcommand\community[0]{\f{K}}
\newcommand\framecommunity[0]{\f{k}}
\newcommand\kernel[0]{\f{ker}}
\newcommand\Kernel[0]{\f{Ker}}
\newcommand\witness[0]{\f{W}}
\newcommand\Kmod[1]{\tf K_{\hspace{-2pt}#1}}
\newcommand\cast[1]{#1^{\ast c}}
\newcommand\cclo[1]{#1^c}
\newcommand\mbot{\mathsf{e}}
\newcommand\mtop{\mathsf{f}}
\newcommand\isdef[2]{#1\bullet #2}
\newcommand\deffont[1]{{\bfseries #1}}
\newcommand\powerset{\f{pow}}
\newcommand\comppowerset{\f{pow}_{\hspace{-2pt}c}}
\newcommand\finpow{\f{fin}}
\newcommand\finto{\stackrel{\f{fin}}{\rightharpoonup}}
\newcommand\f[1]{\mathit{#1}}
\newcommand\tf[1]{\mathsf{#1}}
\DeclareMathAlphabet{\mathsb}{OT1}{cmss}{sbc}{n}
\newcommand\ns[1]{{\mathsb{#1}}}
\usepackage{slantsc}
\newcommand\theory[1]{\ensuremath{\text{\itshape\scshape #1}}\xspace}
\newcommand\ThyPax{\theory{TyPax}}
\newcommand\ThySPax{\theory{TySPax}}
\newcommand\ThyBB{\theory{TyBB}}
\newcommand\finsubseteq{\mathbin{\subseteq_{\text{\it fin}}}}
\newcommand\at{\text{@}}
\newcommand\tx{\f{tx}}
\newcommand\ty{\f{ty}}
\newcommand\ctx{\f{ctx}}
\newcommand\txs{\f{txs}}
\newcommand\utxi{\f{utxi}}
\newcommand\utxo{\f{utxo}}
\newcommand\stx{\f{stx}}
\newcommand\utxio{\f{utxio}}
\newcommand\UTxIO{\f{UTxIO}}
\newcommand\inlinehask[1]{\lstinline[language=hask]{#1}}
\newcommand\inlinesolidity[1]{\lstinline[language=solidity]{#1}}
\newcommand\valid{\f{valid}}
\newcommand\liff{\Longleftrightarrow}
\newcommand\limp{\Longrightarrow}
\newcommand\nlimp{\centernot{\Longrightarrow}}
\newcommand\aeq{\mathrel{=_{\alpha}}}
\newcommand\ssm{{{:}\text{=}}}
\DeclareMathSymbol{\shortminus}{\mathbin}{AMSa}{"39}
\newcommand\minus{{\shortminus}}
\newcommand\plus{{+}}
\newcommand\Forall[1]{\forall #1.}
\newcommand\Exists[1]{\exists #1.}
\newcommand\sm{\mapsto}
\newcommand\rewritesto{\longrightarrow}
\newcommand\cent{\vdash}
\newcommand\ncent{\not\vdash}
\newcommand\ment{\vDash}
\newcommand\nment{\not\vDash}
\newcommand\mentval{\ment_{\avaluation}}
\newcommand\nmentval{\nment_{\avaluation}}
\newcommand\supp{\f{supp}}
\newcommand\boundary{\f{boundary}}
\newcommand\lmodel{[\hspace{-0.2em}[}
\newcommand\rmodel{]\hspace{-0.2em}]}
\newcommand\model[1]{{\lmodel #1 \rmodel}}
\newcommand\modellabel[2]{{\lmodel #1 \rmodel}_{#2}}
\newcommand\synmodel[1]{{\langle #1 \rangle}}
\newcommand\mone{{\text{-}1}}
\newcommand\fv{\f{fv}}
\newcommand\Ngeqz{\mathbb N_{\scriptscriptstyle\geq 0}}
\newcommand\Time{\Ngeqz}
\newcommand\weakmodusponens{weak modus ponens (Proposition~\ref{prop.mp.for.tnotor}(\ref{item.mp.for.tnotor}))\xspace}
\newcommand\strongmodusponens{strong modus ponens (Proposition~\ref{prop.mp.for.tnotor}(\ref{item.mp.for.timpc}))\xspace}
\newcommand\weakmodusponensnoref{weak modus ponens\xspace}
\newcommand\strongmodusponensnoref{strong modus ponens\xspace}

\makeatletter
\DeclareRobustCommand{\barcent}{\mathbin{\mathpalette\barcent@@\relax}}
\newcommand{\barcent@@}[2]{%
  \vbox{\offinterlineskip
    \sbox\z@{$\m@th#1\cent$}%
    \ialign{%
      \hfil##\hfil\cr
      $\m@th#1{}_{\minus}\kern-\scriptspace$\cr
      \noalign{\kern-.3\ht\z@}
      \box\z@\cr
    }%
  }%
}
\makeatother

\makeatletter
\def\pmb@#1#2{\setbox8\hbox{$\m@th#1{#2}$}%
  \setboxz@h{$\m@th#1\mkern-.1mu$}\pmbraise@\wdz@
  \binrel@{#2}%
  \dimen@-\wd8 %
  \binrel@@{%
    \mkern-.1mu\copy8 %
    \kern\dimen@\mkern-.2mu\copy8 %
    \kern\dimen@\mkern-.3mu\copy8 %
    \kern\dimen@\mkern-.4mu\copy8 %
    \kern\dimen@\mkern.1mu\copy8 %
    \kern\dimen@\mkern.2mu\copy8 %
    \kern\dimen@\mkern.3mu\copy8 %
    \kern\dimen@\mkern.0mu\raise\pmbraise@\copy8 %
    \kern\dimen@\mkern.4mu\box8 %
           }%
}
\makeatother


\newcommand\overbracebig[2]{\overbrace{#2}^{\mbox{$#1$}}}
\newcommand\compressthislight[1]{{\hspace{.8pt}\raisebox{.5pt}{\scalebox{.85}{$#1$}}\hspace{.2pt}}}
\newcommand\compressthis[1]{\pmb{\compressthislight{#1}}}
\newcommand\tempty{\compressthis{\varnothing}}
\newcommand\tcap{{\compressthis{\cap}}}
\newcommand\tcup{{\compressthis{\cup}}}
\newcommand\tsubseteq{{\compressthis\subseteq}}
\newcommand\tneg{{\pmb\neg}}
\newcommand\tnegnf{{\pmb{\sim}}}
\newcommand\tnegneg{{\pmb{{\sim}{\sim}}}}
\newcommand\ttop{{\pmb\top}}
\newcommand\tbot{{\pmb\bot}}
\newcommand\teq{{\pmb{\text{=}}}}
\newcommand\tneq{{\pmb{\neq}}}
\newcommand\tand{\mathrel{\pmb\wedge}}
\newcommand\tor{\mathbin{\pmb\vee}}
\newcommand\tgg{{\compressthis\gg}}
\newcommand\tleq{\compressthis{\leq}}
\newcommand\tlesssim{\compressthis{\lesssim}}
\newcommand\tapprox{\compressthis{\approx}}
\newcommand\timp{{\compressthis{\Rightarrow}}}
\newcommand\ntimp{{\pmb\not\Rightarrow}}
\newcommand\tnotor{\mathrel{\arrowr}}
\newcommand\tnotore{\mathrel{\arrower}}
\newcommand\tlatticeiff{\mathrel{\arrowlr}}
\newcommand\tlatticeiffec{\mathbin{\arrowelr}}
\newcommand\timpc{\mathrel{\arrowrr}}
\newcommand\timpcc{\mathrel{\arrowrrr}}
\newcommand\timpcleft{\mathrel{\arrowll}}
\newcommand\timpec{\mathrel{\arrowerr}}
\newcommand\tlatticeiffc{{\pmb{\bar\equiv}}}
\newcommand\tiff{\compressthis{\Leftrightarrow}} %
\newcommand\tiffc{\mathbin{\arrowlrr}}
\newcommand\tiffcc{\mathbin{\arrowllrr}}
\newcommand\tiffec{\mathbin{\arrowelrr}}
\newcommand\tlam{{\pmb\lambda}}
\newcommand\tsim{\compressthis{{\approx}}}
\newcommand\tin{{\pmb{\in}}}
\newcommand\ntin{{\pmb{\not\in}}}
\newcommand\qall{{\compressthis{\mathsf Q}}}
\newcommand\tall{{\compressthis{\forall}}}
\newcommand\texi{{\compressthis{\exists}}}
\newcommand\texiunique{{\compressthis{\exists\hspace{0.5pt}!}}}
\newcommand\texiaffine{{\compressthis{\exists_{01}}}}
\newcommand\texiR{{\compressthis{\exists_{\mathsf R}}}}
\newcommand\tlot{\textbf{?}}
\newcommand\AllBut[2]{\tf{AllBut}_{#2}(#1)}
\newcommand\Nset[1]{\{0\dots #1\minus 1\}}

\newcommand\oimplies{\mathbin{\rlap{$+$}\hspace{.7pt}{\supset}}}

\newcommand\QLogic{Coalition Logic\xspace}

\makeatletter
\newcommand{\circlearrow}{}%
\DeclareRobustCommand{\circlearrow}{%
  \mathrel{\vphantom{\shortrightarrow}\mathpalette\circle@arrow\relax}%
}
\newcommand{\circle@arrow}[2]{%
  \m@th
  \ooalign{%
    \hidewidth$#1\circ\mkern1mu$\hidewidth\cr
    $#1\longrightarrow$\cr}%
}
\makeatother
\newcommand\minusoplus{{\pmb{\circlearrow}}} %

\makeatletter
\newcommand*\bigcdot{\mathpalette\bigcdot@{.5}}
\newcommand*\bigcdot@[2]{\mathbin{\vcenter{\hbox{\scalebox{#2}{$\m@th#1\bullet$}}}}}
\makeatother

\newcommand\act{{\cdot}}
\newcommand\pact{{\cdot}}
\newcommand\mact{\mathbin{\bigcdot}}

\usepackage{datetime}
\yyyymmdddate

\begin{document}
\title{A declarative approach to specifying distributed algorithms using three-valued modal logic}
\newcommand\titlerunning{\emph{Declarative distributed algorithms via modal logic}} %
\newcommand\authorrunning{\emph{Murdoch J. Gabbay and Luca Zanolini}}
\author{Murdoch J. Gabbay \affil{Heriot-Watt University, Edinburgh, UK}
Luca Zanolini \affil{Ethereum Foundation, London, UK}}

\begin{abstract}
  We present \emph{\QLogic}, a three-valued modal fixed-point logic designed for declaratively specifying and reasoning about distributed algorithms, such as the Paxos consensus algorithm.

  Our methodology represents a distributed algorithm as a logical theory, enabling correctness properties to be derived directly within the framework—or revealing logical errors in the algorithm's design when they exist.
  
  \QLogic adopts a declarative approach, specifying the overall logic of computation without prescribing control flow. Notably, message-passing is not explicitly modeled, distinguishing our framework from approaches like TLA+. This abstraction emphasises the logical essence of distributed algorithms, offering a novel perspective on their specification and reasoning.
  
  We define the syntax and semantics of \QLogic, explore its theoretical properties, and demonstrate its applicability through a detailed treatment of the Paxos consensus algorithm. By presenting Paxos as a logical theory and deriving its standard correctness properties, we showcase the framework’s capacity to handle non-trivial distributed systems.
  
  We envision \QLogic as a versatile tool for specifying and reasoning about distributed algorithms. The Paxos example highlights the framework’s ability to capture intricate details, offering a new lens through which distributed algorithms can be specified, studied, and checked.


\keywords{Modal logic, distributed consensus} 
\end{abstract}


\maketitle
\thispagestyle{empty}

\tableofcontents

\section{Introduction}
\label{sect.intro}

\subsection{Initial overview}

\subsubsection*{Introducing \QLogic}
\emph{\QLogic}is a logical framework for studying distributed algorithms in a \emph{declarative} manner.
In our framework an algorithm is presented as a logical theory (a signature and a set of axioms), and then its correctness properties are derived from those axioms.
We apply our framework to present and study an axiomatisation of Paxos, a canonical distributed consensus algorithm. 

Logic \emph{per se} is widely-applied in distributed algorithms: for example the logic TLA+ (based on Linear Temporal Logic) is an industry standard to model and verify distributed protocols~\cite{DBLP:journals/toplas/Lamport94,DBLP:books/aw/Lamport2002,DBLP:conf/charme/YuML99}. 
However, TLA+ specifications describe state changes in a sequential, step-by-step manner.
A \QLogic axiomatisation feels more abstract: we dispense with step-by-step descriptions of the algorithm and with explicit message-passing, to capture an abstract but still entirely identifiable essence of the algorithm.

Abstracting an algorithm to a logical theory \emph{per se} is also not new: \emph{declarative programming} is based on this.
However, turning a distributed algorithm (like Paxos) into an axiomatic theory, is new, and this paper shows that it can be done.
Examples follow below and in the body of the text.\footnote{Or, if the reader wants to skip the preliminaries and jump to a toy example that illustrates our method, then see Subsection~\ref{subsect.simple}.}

For reasons of length, we focus on the example of Paxos.
This is a canonical example with sufficient complexity to be interesting, but we have checked that \QLogic can be adapted for other algorithms; it is intended as a general logical tool that can be tailored to particular case studies.
In the long run, we hope this will help do all the things that declarative techniques are supposed to: help find errors, decrease development times, increase productivity, and lead to new mathematics.
 
\subsubsection*{What is declarative programming}
By `declarative' we mean that we specify the logic of a computation without specifying its control flow.

Thus our axiomatisation of Paxos (in Figure~\ref{fig.logical.paxos}) relates to a concrete Paxos implementation and/or its TLA+ specification, much as this axiomatic, algebraic specification of the factorial function\footnote{By `axiomatic' we mean that the function is specified just by asserting properties that are true of it.  By `algebraic' we mean that these assertions are purely equational, i.e. we just assert some equalities between terms.}
$$
0!=1
\qquad 
(n+1)!=(n+1)*n!
$$
relates to this concrete implementation of factorial in BASH
\begin{quote}
{\tt\small\hspace{-2em} value=1; while [ \$number -gt 1 ]; do value=\$((value * number)); number=\$((number - 1)); done}
\end{quote}
The reader may be familiar with programming languages based on declarative principles, such as Haskell, Erlang, or Prolog.
Their design philosophy is that we should assert things that we want to be true, and leave it to the compiler to figure out how best to make this happen.
Contrast with C, which is a canonical imperative language.\footnote{We are not interested here in taking sides.  With a variety of tools to choose from, we pick the one that is best for the job.  What matters is that the widest possible \emph{variety} of levels of abstraction is made available.}

This paper %
explores a methodology for specifying distributed algorithms which sits decidedly at the `Haskell/Erlang/Prolog' side of the spectrum, compared to TLA+ which (while still quite high-level) sits more at the C side of the spectrum. 

Consider a tiny concrete example of a simple voting system: we will briefly exhibit this in \QLogic, just to give a flavour of how it works.
Consider a voting system such that 
\begin{quote}
`if a quorum of voters vote for candidate $c$, then $c$ wins the election'.
\end{quote}
We write this axiomatically simply as
$$
(\QuorumBox\tf{vote}(c)) \timpc \tf{wins}(c) .
$$ 
The $\QuorumBox$ is a logical modality that is true when there exists a quorum of participants at which $\tf{vote}(v)$ holds.\footnote{Details of this modality in our framework are in Proposition~\ref{prop.pointwise.dense.char}(\ref{item.dense.char.pointwise.QuorumBox}) and Lemma~\ref{lemm.allbut.concretely}.}

Even without going into the details of what these symbols mean, it is clear that details of message-passing, of how votes are communicated, and of how they are counted up, are absent.
In fact in this case the predicate is just a direct rendering of the English statement above it, without specifying how it is concretely implemented.
Thus, in keeping with a declarative philosophy, we specify what we want the algorithm to do, but not how it should do it: it is the point of the axiom to express formally and precisely our idea --- and to add as little as possible to that, beyond what is necessary to make it logically precise.


\subsubsection*{Who should read this document}
We have in mind two communities: 
\begin{enumerate*}
\item
\emph{logicians}, who might be interested in exploring a logic with an interesting combination of features (three-valued modal fixedpoint logic) with specific practical motivations in programming distributed systems, and 
\item
\emph{researchers in distributed systems}, who might find \QLogic useful for analysing and communicating their designs.
\end{enumerate*}
We will be particularly interested in consensus algorithms as applied to blockchain, because this is an important and safety-critical use case (the second author is part of the core team developing consensus algorithms for Ethereum).

\subsubsection*{References to key points in the text}

The reader who might want to jump to where the action is and read backwards and forwards from there, is welcome to proceed as follows:
\begin{itemize*}
\item
The axioms of Declarative Paxos are in Figure~\ref{fig.logical.paxos}.
\item
The syntax and semantics of \QLogic are in Figures~\ref{fig.predicate.syntax}, \ref{fig.3.phi.f}, and~\ref{fig.3.derived}.
\item
Truth-tables for three-valued logic and its connectives are in Figure~\ref{fig.3}.
\item
An outline of the Paxos algorithm (for the reader who might like to be reminded what it is) is in Remark~\ref{rmrk.high-level.paxos}.
\item
See also toy axiomatisations for an imaginary `simple' protocol in Subsection~\ref{subsect.simple}, designed to be vaguely reminiscent of Paxos, but simplified almost to triviality for the sake of exposition.
\end{itemize*}
 

\subsection{Motivation and background (further details)}

Let us unpack what consensus in distributed systems (like blockchains) is, why it matters, and why it is hard:

\subsubsection*{Distributed systems}
A distributed computing system is is one that distributes its computation over several machines.
Distributed systems are essential to computing: they enable a network of parties to collaborate to achieve common goals~\cite{cachinbook}. 

Being distributed is attractive to software engineers because it can give us desirable properties, including some or all of: scalability, redundancy, lower latency, reliability, and resilience.
For this reason, almost any nontrivial commercial software system that the reader might care to name is implemented in a distributed manner.
Furthermore, blockchain systems (like Bitcoin and Ethereum\footnote{\url{https://bitcoin.org/en/} and \url{https://ethereum.org/en/}}) and peer-to-peer systems are inherently distributed by design; if they were centralised, then they would not even make sense.

\subsubsection*{The challenge of consensus}
Among the many challenges in distributed systems, \emph{consensus} represents a fundamental abstraction.
Consensus captures the problem of 
\begin{quote}
reaching agreement among multiple participants on a common value, in the face of possibly unreliable communication, and in the presence of possibly faulty participants.
\end{quote}
Consensus algorithms are an essential prerequisite for coordination in any distributed system, and they play a particularly critical role in blockchain networks.
In such networks, consensus is the mechanism that ensures all participants agree on the state of the ledger --- the blockchain's source of truth --- thereby solving problems like \emph{double spending} (manipulating the blockchain ledger to spend the same token twice).
Consensus protocols also enable the secure ordering of transactions, and they ensure the integrity and consistency of the blockchain, even in the face of adversarial actors.

Consensus protocols in blockchain systems are particularly challenging to design and implement correctly. 
Indeed, there are numerous examples in which incomplete or insufficiently-specified protocols have led to serious issues, such as network halts or vulnerabilities that adversaries could exploit.
For example:
\begin{enumerate*}
\item
Solana (\url{https://solana.com}) has experienced instances where its consensus mechanism contributed to temporary network halts~\cite{yakovenko2018solana,shoup2022poh}.
\item 
Ripple's (\url{https://ripple.com}) consensus protocol has been analysed for conditions in which safety and liveness may not be guaranteed~\cite{DBLP:conf/opodis/Amores-SesarCM20}. 
\item
Avalanche's (\url{https://www.avax.network}) consensus protocol has faced scrutiny regarding its assumptions and performance in adversarial scenarios~\cite{DBLP:conf/opodis/Amores-SesarCT22}.
\end{enumerate*}
Even mature blockchains like Ethereum are not immune to challenges.
Attacks such as \emph{balancing attacks}~\cite{DBLP:conf/sp/NeuTT21} and \emph{ex-ante reorganisation} (\emph{reorg}) attacks~\cite{DBLP:conf/fc/Schwarz-Schilling22} have highlighted vulnerabilities in Ethereum's protocol.

These examples underscore how subtle flaws or ambiguities in protocol design can have significant consequences, emphasising the need for rigorous analysis and well-defined specifications in this complex and safety-critical field.

Compounding the fact that consensus algorithms are foundational to the security and functionality of blockchain systems, is the fact that designing and implementing these algorithms is notoriously challenging.
Distributed algorithms --- including consensus protocols --- are difficult (at least for human brains) because of the inherent difficulty of coordinating multiple participants, which may be faulty, hostile, drifting in and out of network connectivity~\cite{sleepy}, and/or have different local clocks, states, agendas, and so forth.

In short, this is a notoriously difficult field which is full of `gotchas', and within this difficult field, consensus protocols are amongst the most intricate, demanding, and practically important objects of study. 

\subsubsection*{Introducing a logic-based declarative approach (applied to Paxos)}
In this paper, we develop a logic-based declarative approach to modeling distributed algorithms.
It stands in relation to existing techniques in roughly the same way as declarative programming languages stand to imperative ones. %

Technically speaking, \QLogic is a \emph{three-valued modal fixedpoint logic}. 
This means:
\begin{itemize*}
\item
There are three truth-values: $\tvT$ (true), $\tvB$ (no value returned, i.e. `crashed'), and $\tvF$ (false)~\cite{sep-logic-manyvalued}.
\item
Modalities let us express things like `some participant has property $\phi$' or `a quorum of participants has property $\psi$' or `$\phi$ has been true in the past' or `$\psi$ will be true in the future'.\footnote{In order, these are as follows in our logic: $\someone\phi$, $\QuorumBox\phi$, $\recent\phi$, and $\sometime\phi$.  Precise semantics are in Figures~\ref{fig.3.phi.f} and~\ref{fig.3.derived}.} 
\item
A fixedpoint operator $\mu$ imports the expressive power of iteration~\cite{wiki:Fixedpoint_logic}.
Consensus algorithms are often iterative, so it is natural to include this in our logic.
\end{itemize*} 
We will apply these constructs to give axiomatic models of the Paxos consensus protocol~\cite{DBLP:journals/tocs/Lamport98,lamport2001paxos,cachin2011paxos} and then we will state and derive standard correctness properties from the axioms.
Our example is elementary but nontrivial and strikes (we think) a good balance between being compact but still exhibiting enough richness of structure to illustrate the essential features of how our logic works. 


The first half of this paper will focus on the logic, and the second half of this paper will apply the logic to the examples.
The reader from a logic background might like to start with the first half and refer to examples for motivation; the reader from a distributed systems background might like to do the opposite and start with the examples in the second half, referring back to the logic as required to understand what is being done.
 
\subsubsection*{Related work using logic}
Logic-based approaches have long been used to model and verify distributed protocols, with TLA+~\cite{DBLP:journals/toplas/Lamport94,DBLP:books/aw/Lamport2002,DBLP:conf/charme/YuML99} being one of the most prominent frameworks in the field. 
Introduced by Leslie Lamport in 1999, TLA+ uses Linear Temporal Logic (LTL)~\cite{DBLP:conf/focs/Pnueli77} to describe state changes in a sequential, step-by-step manner, which has become an industry standard for specifying and verifying distributed algorithms.

But TLA+ does things in a very particular way, which is not declarative.
We would argue that it is good for specifying \emph{algorithms}, and so is not particularly \emph{declarative}.
Specifically, TLA+ excels at specifying and verifying system behaviours through state transitions; it is operational and state-based.
This can sometimes distance TLA+ protocol specifications from the invariants that they are intended to uphold.

This separation arises because TLA+ primarily describes how states evolve \emph{over time}, which can make it challenging to directly reason about the essential properties that protocols are designed to maintain. 
So, although TLA+ provides powerful tools for specifying and checking invariants, the process of verifying that these invariants are preserved across all possible state transitions can be complex and unintuitive, particularly when dealing with the intricacies of distributed algorithms.

\subsubsection*{Our approach}
Our novel framework leverages the mathematics of semitopology and three-valued logic. 

Semitopology is from previous work~\cite{gabbay:semdca,gabbay:semtad}.
It provides a natural way to describe quorums --- sets of participants that work together to progress the system state --- by representing them as open sets.
This perspective is particularly suited to distributed systems, where the behaviour of the system emerges from the interactions of coalitions its components, rather than from a sequence of steps by a central controller.

Three-valued logic (a good overview is in \cite[Chapter~5]{bergmann:intmvf}) helps us to reason about the (possibility for) faults that is inherent in distributed systems.
Unlike classical logic with its binary truth-values $\tvT$ and $\tvF$, three-valued logic introduces a third truth value $\tvB$ that captures the indeterminate states often encountered when dealing with faulty parties.
This turns out to give particularly natural and compact descriptions of distributed protocols, distinguishing between normal operation and faulty behaviour. 
It also provides a formal foundation that would be amenable to verification through tools like SAT solvers and theorem provers (building tooling for what we see in this paper would be future work), enhancing both the expressiveness and verifiability of protocol specifications.

Our approach is particularly relevant in the context of blockchain consensus protocols, which are distributed and decentralised systems that require robust and clear specifications to function correctly.
It is not an exaggeration to say that blockchains in particular have put a rocket underneath the development of new consensus protocols~\cite{DBLP:conf/wdag/CachinV17}.
Many of these are not fully specified or formally verified. 
Often, they are even described informally in blog posts and other non-academic sources --- leading to confusion about their properties and behaviours. 
This lack of formalisation creates a fragmented landscape, where the guarantees and limitations of the protocols remain unclear, hindering their adoption and undermining trustworthiness.

We hope that by providing a logic-based declarative language, and the ability to axiomatise consensus protocols, our framework will offer a way to formally specify, reason about, and verify the properties of the many distributed systems under development --- thus lowering costs, increasing reliability, and reducing development time. 
We would further argue that effective formalisation is essential in reducing ambiguity in the design and evaluation of the many new consensus mechanisms being developed in the blockchain space in particular.

Our approach bridges the gap between the declarative specification of invariants and the practical verification of protocol correctness. 

\subsubsection*{We will explain everything, and that's a feature}

Different readers will find some parts of this paper easy, and other parts less easy --- but which parts these are depends on the reader and their background.\footnote{Not a hypothetical: the first author (with a logic background) took a long time to get his head around how Paxos works, whereas the second author (who knows Paxos well and is familiar with logic in general terms, but is not a logician) was not immediately familiar with modal fixedpoint logic and models of axiomatic logical theories.  The community of readers who approach this paper and are already well-versed in both fields may not be the largest.  We hope this paper will help ameliorate this lamentable state of affairs.}

For example:
\begin{itemize*}
\item
three-valued modal logic and its models, and technical results like the Knaster-Taski fixedpoint theorem used in Proposition~\ref{prop.mu.fixedpoint} may be known to the reader with a background in logic (like the first author), but perhaps not to a distributed systems theorist (like the second author); conversely, 
\item
the Paxos algorithm, and standard technical devices like the discussion motivating $n$-twined semitopologies in Remark~\ref{rmrk.motivate.n-twined}, are the bread-and-butter of a distributed systems expert's work --- but this application might be new to someone from a logic background. 
\end{itemize*}
We play safe and explain everything. 
If the reader finds we are going a bit slow in places, then better that than going too fast and skipping details that might help some other reader follow the argument.
 


\subsection{Structure of the paper}

\begin{enumerate*}
\item
Section~\ref{sect.intro} is the Introduction; You are Here.
The rest of the paper splits into two halves: Sections~\ref{sect.three.truth.values} to Section~\ref{sect.n-twined} have to do with logic and semitopologies; and Sections~\ref{sect.declarative.paxos} to~\ref{sect.simpler.paxos} use \QLogic to declaratively specify and study Paxos.
\ \\
\item
Section~\ref{sect.three.truth.values} introduces the basics of three-valued logic, including notable connectives ($\tnotor$ and $\timpc$) and modalities ($\modT$ and $\modTB$).
We also introduce a notion of validity for truth-values.
\item
Section~\ref{sect.quorum.logic} defines semitopologies (our model of the quorums that are ubiquitous in consensus algorithms), then it builds the syntax and semantics of \QLogic.
\item
Section~\ref{sect.predicates} studies \QLogic predicates; notably developing some logical equivalences, and well-behavedness criteria like being \emph{pointwise}.
\item 
Section~\ref{sect.n-twined} returns to semitopologies, and shows how antiseparation properties can be represented in logical form in \QLogic.
One particular antiseparation property corresponds to the \emph{quorum intersection property} which is used in the correctness proofs for Paxos (see Remark~\ref{rmrk.ntwined.and.the.literature}), but we will put this fact in a more general logical/semitopological context.
\ \\
\item
Section~\ref{sect.declarative.paxos} introduces Declarative Paxos; this is representative for the kind of algorithm that we built \QLogic to express.
We take our time to explain how Paxos works, state the axiomatisation, and discuss it in detail.
\item
Section~\ref{sect.paxos.correctness.properties} states the standard correctness properties for Paxos, phrases them in declarative form, and then proves them from the axioms of Declarative Paxos.
What all our work has accomplished is to convert proving correctness properties into a (relatively straightforward) argument of proving facts of an abstract model from its axioms.
\item
Section~\ref{sect.simpler.paxos} revisits the axioms of Declarative Paxos from Section~\ref{sect.declarative.paxos}, in the light of the proofs in Section~\ref{sect.paxos.correctness.properties} --- and then prunes them down to a minimal theory required to prove correctness.

We call this minimal theory \emph{Simpler Declarative Paxos}: it is simpler, and because further details have been elided, it is more abstract.
The reader may prefer Declarative Paxos to Simpler Declarative Paxos because it is more identifiably related to the Paxos algorithm, but if we want to prove correctness properties, then Simpler Declarative Paxos distils the logical essence of what makes those proofs work. 
\ \\
\item
Section~\ref{sect.conclusions} concludes with conclusions and reflections on future work.
\end{enumerate*} 


\section{Three truth-values}
\label{sect.three.truth.values}

\subsection{Basic definition of $\THREE$, and some truth-tables}

\begin{defn}[Three-valued truth-values]
\label{defn.THREE}
Let $\THREE$ be the poset of \deffont{truth-values} with elements $\tvF$ (false), $\tvB$ (both), and $\tvT$ (true) ordered as 
$$
\tvF<\tvB<\tvT.
$$
\end{defn}

\begin{figure} %
$$
\begin{array}{c}
\begin{array}{c@{\qquad}c@{\qquad}c@{\qquad}c@{\qquad}c@{\qquad}c}
\begin{array}{c c c c}
p\tand q & \tvT & \tvB & \tvF
\\
\tvT &  \tvT & \tvB & \tvF
\\
\tvB &  \tvB & \tvB & \tvF 
\\
\tvF &  \tvF & \tvF & \tvF 
\end{array}
&
\begin{array}{c c c c}
p\tor q & \tvT & \tvB & \tvF
\\
\tvT &  \tvT & \tvT & \tvT
\\
\tvB &  \tvT & \tvB & \tvB
\\
\tvF &  \tvT & \tvB & \tvF
\end{array}
&
\begin{array}{c c c c}
p\tnotor q & \tvT & \tvB & \tvF
\\
\tvT &  \tvT & \tvB & \tvF
\\
\tvB &  \tvT & \tvB & \tvB
\\
\tvF &  \tvT & \tvT & \tvT
\end{array}
&
\begin{array}{c c c c}
p \timpc q & \tvT & \tvB & \tvF
\\
\tvT       & \tvT & \tvF & \tvF
\\
\tvB       & \tvT & \tvB & \tvB
\\
\tvF       & \tvT & \tvT & \tvT
\end{array}
\end{array}
\\[7ex]
\begin{array}{c@{\qquad}c@{\qquad}c@{\qquad}c@{\qquad}c@{\qquad}c@{\qquad}c}
\begin{array}{cc}
\tneg p
\\
\tvT & \tvF
\\
\tvB & \tvB
\\
\tvF & \tvT
\end{array}
&
\begin{array}{cc}
\modT p 
\\
\tvT & \tvT
\\
\tvB & \tvF
\\
\tvF & \tvF
\end{array}
&
\begin{array}{cc}
\modF p 
\\
\tvT & \tvF
\\
\tvB & \tvF
\\
\tvF & \tvT
\end{array}
&
\begin{array}{cc}
\modB p
\\
\tvT & \tvF
\\
\tvB & \tvT
\\
\tvF & \tvF
\end{array}
&
\begin{array}{cc}
\modTB p  
\\
\tvT & \tvT
\\
\tvB & \tvT
\\
\tvF & \tvF
\end{array}
&
\begin{array}{cc}
\modTF p  
\\
\tvT & \tvT
\\
\tvB & \tvF
\\
\tvF & \tvT
\end{array}
&
\begin{array}{cc}
\modFB p  
\\
\tvT & \tvF
\\
\tvB & \tvT
\\
\tvF & \tvT
\end{array}
\end{array}
\\[7ex]
\begin{array}{c@{\qquad}c@{\qquad}c@{\qquad}c@{\qquad}c@{\qquad}c}
\begin{array}{c c c c}
p\timp q & \tvT & \tvB & \tvF
\\
\tvT &  \tvT & \tvB & \tvF
\\
\tvB &  \tvT & \tvB & \tvF 
\\
\tvF &  \tvT & \tvT & \tvT 
\end{array}
&
\begin{array}{c c c c}
p\tleq q & \tvT & \tvB & \tvF
\\
\tvT &  \tvT & \tvF & \tvF
\\
\tvB &  \tvT & \tvB & \tvF 
\\
\tvF &  \tvT & \tvT & \tvT 
\end{array}
\end{array}
\end{array}
$$
\emph{Above, the vertical axis of a table indicates values for $p$; the horizontal axis (if nontrivial) denotes values for $q$.}
\caption{Truth-tables for some operations on $\THREE$ (Remark~\ref{rmrk.three.elementary})}
\label{fig.3}
\end{figure}


\begin{figure} %
$$
\begin{array}{c@{\qquad}c@{\qquad}c@{\qquad}c@{\qquad}c@{\qquad}c}
\begin{array}{c c c c}
p\tlatticeiff q & \tvT & \tvB & \tvF
\\
\tvT &  \tvT & \tvB & \tvF
\\
\tvB &  \tvB & \tvB & \tvB
\\
\tvF &  \tvF & \tvB & \tvT
\end{array}
&
\begin{array}{c c c c}
p \tiffcc q & \tvT & \tvB & \tvF
\\
\tvT       & \tvT & \tvF & \tvF
\\
\tvB       & \tvF & \tvB & \tvB
\\
\tvF       & \tvF & \tvB & \tvT
\end{array}
&
\begin{array}{c c c c}
p \equiv q & \tvT & \tvB & \tvF
\\
\tvT       & \tvT & \tvF & \tvF
\\
\tvB       & \tvF & \tvT & \tvF
\\
\tvF       & \tvF & \tvF & \tvT
\end{array}
\end{array}
$$
\caption{Truth-tables for notions of equivalence (Definition~\ref{defn.equivalent})}
\label{fig.3.iff}
\end{figure}


\begin{rmrk}
\label{rmrk.three.elementary}
We start with some elementary observations about $\THREE$.
\begin{enumerate*}
\item\label{item.three.elementary.bounded.lattice}
$\THREE$ %
has bottom element $\tvF$, top element $\tvT$, and meets $\tand$ and joins $\tor$.
This makes it a \emph{finite lattice}.
\item
$\THREE$ has a negation operation $\tneg$.
\item
$\THREE$ has several natural modalities that identify truth-values, including $\modT$, $\modF$, $\modB$, $\modTB$, and $\modTF$.

For example, the reader can examine the truth-table for $\modTF$ in Figure~\ref{fig.3} and see that it identifies elements of $\{\tvT,\tvF\}$.\footnote{One could also call these modalities \emph{characteristic functions} for various subsets of $\THREE$, by which we mean functions which (literally) return $\tvT$ on elements in the set and $\tvF$ on elements not in the set.}
\item
$\THREE$ also has multiple notions of implication and equivalence, some of which are illustrated in Figures~\ref{fig.3} and~\ref{fig.3.iff}.
We will discuss these implications in Subsection~\ref{subsect.3.imp}.
\end{enumerate*}
\end{rmrk}


\begin{rmrk}
We will use the third truth-value $\tvB$ to represent participants in a distributed algorithm \emph{crashing}.
The method is simple: the truth-value of a predicate describing the outcome of a crashed computation, is deemed to be $\tvB$.

This is similar in spirit to how in domain theory an additional value --- usually written $\bot$ --- is added to a denotation to represent nontermination; the function-value of a nonterminating computation is deemed to be $\bot$~\cite{abramsky:domt}.
The nonterminating computation does not return a value, but we represent this absence of a value in the denotation, using $\bot$. 

In our declarative axiomatisation of Paxos, we use $\tvB$ in much the same way, as a truth-value to describe the \emph{absence} of a return value (either $\tvT$ or $\tvF$) from participants that have crashed.
\end{rmrk}



\subsection{Valid and correct truth-values}



\begin{defn}
\label{defn.tv.ment}
\leavevmode
\begin{enumerate*}
\item\label{item.tv.ment.correct}
Define subsets $\threeValid,\threeCorrect,\threeTrue,\threeFalse\subseteq\THREE$ by
$$
\threeValid = \{\tvT,\tvB\}
\quad\text{and}\quad
\threeCorrect = \{\tvT,\tvF\} 
\quad\text{and}\quad
\threeTrue = \{\tvT\} 
\quad\text{and}\quad
\threeFalse = \{\tvF\} .
$$ 
\item
We refer to $\f{tv}\in\THREE$ as being \deffont{valid} / \deffont{invalid},\ \deffont{correct} / \deffont{incorrect},\ and \deffont{true} / \deffont{untrue},\ and \deffont{false} / \deffont{unfalse}\footnote{\dots `unfalse' is `$\tvT$-or-$\tvB$', not `$\tvT$'.  Similarly, `untrue' is $\tvB$-or-$\tvF$.  We cannot resist stating the following gem of three-valued clarity: not-$\tvT$, is not $\tneg\tvT$, and not-$\tvF$ is not $\tneg\tvF$, and not-$\tvB$ \emph{is} not-$\tneg\tvB$, because $\tvB$ is $\tneg\tvB$.} depending on whether it is / is not in $\threeValid$, $\threeCorrect$, $\threeTrue$, and $\threeFalse$ respectively.
\item\label{item.tv.ment}
Define $\ment\f{tv}$ \ (read `\deffont{$\f{tv}$ is valid}') by
$$
\ment\f{tv} 
\quad\text{when}\quad
\f{tv}\in\threeValid ,
\ \ \text{which just means}\ \ 
\f{tv}\in\{\tvT,\tvB\} .
$$
Thus, `$\f{tv}$ is valid' holds precisely when $\f{tv}$ is a valid truth-value.
\end{enumerate*}
\end{defn}

\begin{rmrk}
\label{rmrk.motivate.validity}
We take a moment to discuss Definition~\ref{defn.tv.ment}:
\begin{enumerate}
\item
We can think of $\tvT$ and $\tvF$ as specific `yes/no' results from the successful completion of some computation.
Therefore, we will call $\tvT$ and $\tvF$ \emph{correct}.
\item
We can think of $\tvB$ as representing failure of a computation to complete correctly --- e.g. because a participant has crashed and so can return neither $\tvT$ nor $\tvF$.
\item\label{item.tb.valid}
In Definition~\ref{defn.tv.ment} we set $\threeValid=\{\tvT,\tvB\}$ and define $\ment\f{tv}$ when $\f{tv}\in\threeValid$, and we call $\tvT$ and $\tvB$ \emph{valid}.
\end{enumerate}
It might seem strange at first to let $\tvB$ be valid; if $\tvB$ represents the (non)behaviour of a crashed participant, then surely crashing should \emph{not} be valid?

Not necessarily.
To see why, consider that many of assertions of interest that we will consider will have the form 
\begin{quote}
`if $p$ is not crashed, then $\ment\phi$ at $p$', 
\end{quote}
or equivalently 
\begin{quote}
`either $p$ is crashed, or $\ment\phi$ at $p$'.
\end{quote}
If we let $\tvB$ be valid (i.e. if we define $\ment\tvB$ to hold), then we can simplify the above to the conveniently succinct assertion 
\begin{quote}
`$\ment\phi$ at $p$'.
\end{quote}
The `if $p$ is not crashed' precondition gets swallowed up into the fact that if $p$ is crashed then $\phi$ returns $\tvB$, and $\ment\tvB$ so there is nothing to prove. 
\end{rmrk}


\begin{rmrk}
Later on in Definition~\ref{defn.validity.judgement}(\ref{item.ctx.ment.phi}) we will extend validity to predicates-in-context, just by evaluating a predicate-in-context to a truth-value using a denotation function, and checking that truth-value is valid in the sense we define here. 
\end{rmrk}


We make some elementary observations about $\THREE$ and Definition~\ref{defn.tv.ment}:
\begin{lemm}
\label{lemm.tand.tor}
Suppose $\f{tv}\in\THREE$ is a truth-value.
Then:
\begin{enumerate*}
\item
$\ment\tvT$ and $\ment\tvB$.
\item
$\nment\tvF$.
\item\label{item.tand.tor.tor}
$\ment\f{tv}\tor\f{tv}'$ if and only if ${\ment\f{tv}}\lor{\ment\f{tv}'}$.
\item
$\ment\f{tv}\tand\f{tv}'$ if and only if ${\ment\f{tv}}\land{\ment\f{tv}'}$.
\end{enumerate*}
\end{lemm}
\begin{proof}
Fact of Definition~\ref{defn.tv.ment} and the truth-tables in Figure~\ref{fig.3}.
\end{proof} 

Negation is conspicuous by its absence in Lemma~\ref{lemm.tand.tor}.
We consider this next:
\begin{lemm}
\label{lemm.para}
Suppose $\f{tv}\in\THREE$ is a truth-value.
Then:
\begin{enumerate*}
\item\label{item.para.em}
$\ment\f{tv}\tor\tneg\f{tv}$, and as a corollary $\nment\f{tv}$ implies $\ment\tneg\f{tv}$.
\item\label{item.para.para}
$\ment\f{tv}\tand\tneg\f{tv}$ is possible, and as corollaries: $\ment\f{tv}$ does not necessarily imply $\nment\tneg\f{tv}$; and it is possible that $\ment\f{tv}$ and $\ment\tneg\f{tv}$ both hold.
\end{enumerate*}
\end{lemm}
\begin{proof}
We consider each part in turn:
\begin{enumerate}
\item
Using Lemma~\ref{lemm.tand.tor}(3) it would suffice to show that $\ment\f{tv}$ or $\ment\tneg\f{tv}$.
So suppose $\nment\f{tv}$.
By Definition~\ref{defn.tv.ment}(\ref{item.tv.ment}) $\f{tv}\not\in\threeValid$, thus $\f{tv}=\tvF$ and by Figure~\ref{fig.3} $\tneg\f{tv}=\tvT\in\threeValid$ and so $\ment\tneg\f{tv}$.
\item
We just take $\f{tv}=\tvB$ and check from Figure~\ref{fig.3} that $\tneg\tvB=\tvB$.
We use Lemma~\ref{lemm.tand.tor}(1\&4).
\qedhere\end{enumerate}
\end{proof}

\begin{rmrk}
Some comments on Lemma~\ref{lemm.para}:
\begin{itemize}
\item
Lemma~\ref{lemm.para}(\ref{item.para.para}) notes a typical fact of validity in three-valued logic, that it is naturally \deffont{paraconsistent}, meaning that it is possible for a thing, and the negation of that thing, to both be valid.
This is a feature of $\tvB\in\THREE$, not a bug, and we will make good use of it later when we use $\tvB$ %
as a return value for `I return no correct truth-value'. 
\item
Lemma~\ref{lemm.para}(\ref{item.para.em}) asserts validity of the \emph{excluded middle} (`$\phi$-or-negation-of-$\phi$'), but its corollaries implicitly illustrate another fact of three-valued logic, by what they do \emph{not} say: 
looking at the proof of Lemma~\ref{lemm.para}(\ref{item.para.em}), it is clearly also true that $\ment\f{tv}\tor\modF\f{tv}$, and that $\ment\f{tv}\tor\modT\tneg\f{tv}$, and that $\nment\f{tv}$ implies $\ment\modF\f{tv}$, and also that $\nment\f{tv}$ implies $\ment\modT\tneg\f{tv}$.  

These are stronger assertions, so why did we not state them in the lemma instead?
Because we will not need the result in those forms later.  

In two-valued logic, if something is true then it is probably a useful lemma.
In contrast, in three-valued logic, there is so much more structure with the third truth-value that more things are \emph{true} than are necessarily \emph{useful}.
So, we have to pick and choose our lemmas, and it is not always the case that the strongest form of a result is the most useful.
\end{itemize}
\end{rmrk}


Lemmas~\ref{lemm.tv.ment.TF} and~\ref{lemm.TF.correct} are easy, but they spell out some useful facts about truth-values and validity judgements:
\begin{lemm}
\label{lemm.tv.ment.TF}
Suppose $\f{tv}\in\THREE$ is a truth-value.
Then:
\begin{enumerate*}
\item\label{item.tv.ment.TF.1}
$\threeCorrect\cap\threeValid=\threeTrue$.

In words: a truth-value is true if and only if it is correct and valid.
\item\label{item.tv.ment.TF.2b}
$\ment \modT\f{tv}$ if and only if $\f{tv}\in\threeTrue$. 

In words: `$\modT\f{tv}$ is valid' holds if and only if $\f{tv}$ is true.
\item\label{item.tv.ment.TF.2b.F}
$\ment \modF\f{tv}$ if and only if $\f{tv}\in\threeFalse$. 

In words: `$\modT\f{tv}$ is valid' holds if and only if $\f{tv}$ is true.
\item\label{item.tv.ment.TF.TB}
$\ment \modTB\f{tv}$ if and only if $\f{tv}\in\threeValid$ if and only if $\ment\f{tv}$. 

In words: `$\modTB\f{tv}$ is valid' if and only if $\f{tv}$ is valid.\footnote{Note that $\modTB\f{tv}$ is not necessarily \emph{literally equal} to $\f{tv}$, in the sense of $\modT\f{tv}\equiv\f{tv}$ from Definition~\ref{defn.equivalent}.  For example, $\modTB\tvB=\tvT\neq\tvB$.}
\item\label{item.tv.ment.TF.2}
$\ment \modTF\f{tv}$ if and only if $\f{tv}\in\threeCorrect$. 

In words: `$\modTF\f{tv}$ is valid' if and only if $\f{tv}$ is correct.
\end{enumerate*}
\end{lemm}
\begin{proof}
Part~\ref{item.tv.ment.TF.1} is a fact of sets.
Parts~\ref{item.tv.ment.TF.2b} and~\ref{item.tv.ment.TF.TB} and~\ref{item.tv.ment.TF.2} are just from the truth-table for $\modT$, $\modTB$, and $\modTF$ in Figure~\ref{fig.3} and Definition~\ref{defn.tv.ment}. 
\end{proof}


\begin{lemm}
\label{lemm.TF.correct}
Suppose $\f{tv}\in\THREE$ is a truth-value.
Then the following are all equivalent:
$$
\ment\modTF\f{tv}
\ \liff\ 
\modTB\f{tv}=\modT\f{tv}
\ \liff\ 
\modT\f{tv}=\f{tv} 
\ \liff\ 
\modFB\f{tv}=\modF\f{tv}
\ \liff\ 
\modF\f{tv}=\tneg\f{tv} 
$$ 
Informally --- we do not build the machinery to make this formal, but it may still be a helpful observation --- we can write that the modality $\modTF$ is equivalent to the modal equivalences $\modTB=\modT=\f{id}$ and $\modFB=\modF=\tneg$.
\end{lemm}
\begin{proof}
By Lemma~\ref{lemm.tv.ment.TF}(\ref{item.tv.ment.TF.2}) and Definition~\ref{defn.tv.ment}(\ref{item.tv.ment.correct}),
$\ment\modTF\f{tv}$ when 
$$
\f{tv}=\tvT \ \lor\ \f{tv}=\tvF
\quad\text{and when}\quad
\f{tv}\neq\tvB .
$$
We check the truth-tables in Figure~\ref{fig.3}, and see that all of the other statements above precisely characterise the same condition.
\end{proof}


\subsection{The rich implication and equivalence structure of $\THREE$}
\label{subsect.3.imp}

\subsubsection{Weak and strong implications, and others}


Three-valued logic supports sixteen different kinds of implication~\cite[note~5, page~22]{arieli:idepl}.
We discuss two of them in Definition~\ref{defn.three.imp}, and mention two more in Remark~\ref{rmrk.rich.implication}:
\begin{defn}
\label{defn.three.imp}
\leavevmode
\begin{enumerate*}
\item\label{item.three.elementary.weak.implication}
Define \deffont{weak implication} $\tnotor$, using $\tneg$ and $\tor$, as follows:
$$
\f{tv}\tnotor\f{tv}'=(\tneg\f{tv})\tor\f{tv}'.
$$
\item\label{item.three.elementary.strong.implication}
Define \deffont{strong implication} $\timpc$, using $\tnotor$ and $\modT$, as follows:
$$
\f{tv}\timpc\f{tv}'=\f{tv}\tnotor \modT\f{tv}'. 
$$
\end{enumerate*}
These and other encodings are assembled in Figure~\ref{fig.3.derived}.
\end{defn}

We motivate the weak and strong implications from Definition~\ref{defn.three.imp}(\ref{item.three.elementary.weak.implication}\&\ref{item.three.elementary.strong.implication}):
\begin{rmrk}[How to interpret $\tnotor$ and $\timpc$]
\label{rmrk.interpret.implication}
\leavevmode
For this Remark, let us interpret $\f{tv}\in\THREE$ as the result of a computation process which either succeeds and returns true (truth-value $\tvT$), or it succeeds and returns false (truth-value $\tvF$), or the process does not succeed (e.g. it crashes) and it returns an invalid value / no value (represented by the truth-value $\tvB$).
Recall also from Definition~\ref{defn.tv.ment}(\ref{item.tv.ment}) that $\ment\f{tv}$ means $\f{tv}\in\{\tvT,\tvB\}$.
Then by this reading:
\begin{enumerate*}
\item
$\ment\f{tv}$ means \emph{``the process $\f{tv}$ returns $\tvT$, if valid''}. 
\item
$\ment\f{tv}\tnotor\f{tv}'$ means \emph{``\emph{if} process~$\f{tv}$ returns true \emph{then} process~$\f{tv}'$ returns $\tvT$, if valid''}.

Implication does not imply sequentiality: we are just making a logical assertion about outcomes. 
\item
$\ment\f{tv}\timpc\f{tv}'$ means \emph{``\emph{if} process~$\f{tv}$ returns true \emph{then} process~$\f{tv}'$ is \emph{also} valid, and it returns $\tvT$''}.
\end{enumerate*}
\end{rmrk}


We can sum this up as a \emph{modus ponens} result, which relates the validity $\ment$ from Definition~\ref{defn.tv.ment}, the weak implication $\tnotor$ and strong implication $\timpc$ from Figure~\ref{fig.3}, to real-world implication:
\begin{prop}[Weak and strong modus ponens]
\label{prop.mp.for.tnotor}
Suppose $\f{tv},\f{tv}'\in\THREE$ are truth-values.
Then:
\begin{enumerate*}
\item\label{item.mp.for.tnotor}
$\ment\f{tv}\tnotor\f{tv}'$ if and only if\ \ $\ment\modT\f{tv}$ implies $\ment\f{tv}'$. 

In words: `$\f{tv}$ weakly implies $\f{tv}'$' is valid, if and only if $\f{tv}$ is true implies $\f{tv}'$ is valid (but $\f{tv}'$ need not be true). 

We may call this logical equivalence %
\deffont{weak modus ponens}.
\item\label{item.mp.for.timpc}
$\ment\f{tv}\timpc \f{tv}'$ if and only if $\ment\modT\f{tv}$ implies $\ment\modT\f{tv}'$. 

In words: `$\f{tv}$ strongly implies $\f{tv}'$' is valid if and only if $\f{tv}$ is true implies $\f{tv}'$ is true. 

We may call this logical equivalence \deffont{strong modus ponens}. 
\end{enumerate*}
\end{prop}
\begin{proof}
We consider each part in turn:
\begin{enumerate}
\item
By direct inspection of the truth-table for $\tnotor$ in Figure~\ref{fig.3}, using Definition~\ref{defn.tv.ment}(\ref{item.tv.ment}).
\item
By a similar direct inspection.
\qedhere\end{enumerate}
\end{proof}


\begin{rmrk}[Explicit and implicit implications]
\label{rmrk.rich.implication}
So far, the structure we have seen (the lattice structure, the negation, the modalities, and the two implications $\tnotor$ and $\timpc$) is reflected explicitly in the syntax of our logic in Figures~\ref{fig.predicate.syntax} and~\ref{fig.3.derived}.
For instance, Figure~\ref{fig.predicate.syntax} has a unary negation connective $\tneg$ whose denotation as per Figure~\ref{fig.3.phi.f} is precisely the negation operator $\tneg$ in Figure~\ref{fig.3}.
So far, so explicit.

However, $\THREE$ has further implication structure, which we do not mention explicitly in our logic, but which is still there, is important, and which appears implicitly in the proofs.
Notably: 
\begin{enumerate*}
\item\label{item.rich.implication.timp}
An implication $\timp$ defined by $\f{tv}\timp \f{tv}'=\tneg\f{tv}'\timpc\tneg\f{tv}=(\modTB\f{tv})\tnotor\f{tv}'$ has a truth-table spelled out in Figure~\ref{fig.3} (bottom left).
This can be read as ``if $\f{tv}$ is valid then $\f{tv}'$ is valid''.

$\timp$ is clearly implicit in results like Proposition~\ref{prop.someone.implies.someoneAll} -- see all the parts whose statement uses logical implication $\limp$ --- and in Proposition~\ref{prop.tv.ment.TF.model}.
\item
Another implication is implicit in the use of $\leq$ in results like Lemma~\ref{lemm.positive.monotonicity} and Corollary~\ref{corr.recent.everyone.commute}.
That symbol refers to the lattice order on $\THREE$, but this could easily be internalised in our logic --- e.g. we can set $\f{tv}\tleq \f{tv'}$ to be $(\f{tv}\timpc\f{tv}')\tand(\f{tv}\timp\f{tv}')$, to obtain the connective $\tleq$ also illustrated in Figure~\ref{fig.3} (bottom right).
\end{enumerate*}
So in fact, we use all four of $\timp$ and $\tleq$, and $\tnotor$ and $\timpc$.

The reason that $\tnotor$ and $\timpc$ make it into our logic as (sugared) syntax, and $\timp$ and $\leq$ do not, is that $\tnotor$ and $\timpc$ naturally arise in our axioms (see most of the axioms in Figures~\ref{fig.logical.paxos}) %
whereas we do not need $\timp$ or $\tleq$ in the same way for the axiomatisations in this paper.
Different axiomatisations of different algorithms might make different demands on our logic, and it would be straightforward to include (sugar for) connectives accordingly. 
\end{rmrk}

\subsubsection{Equivalence $\equiv$}

Our three-valued domain of truth-values $\THREE$ has rich notions of equivalence, to match its rich notions of implication.
We mention three of them:
\begin{defn}[Notions of equivalence in $\THREE$]
\label{defn.equivalent}
Define $\phi\tlatticeiff\phi'$ and $\phi\tiffcc\phi'$ and $\phi\equiv\phi'$ by:
$$
\begin{array}{r@{\ }l}
\phi\tlatticeiff\phi' =& (\phi\tnotor\phi')\tand(\phi'\tnotor\phi)
\\
\phi\tiffcc\phi' =& (\phi\timpc\phi')\tand(\phi'\timpc\phi)
\\
\phi\equiv\phi' =& (\modT\phi\tand\modT\phi')\tor(\modB\phi\tand\modB\phi')\tor(\modF\phi\tand\modF\phi')
\end{array}
$$
Truth-tables for $\tlatticeiff$, $\tiffcc$, and $\equiv$ are given in Figure~\ref{fig.3.iff}, and we unpack this further in Remark~\ref{rmrk.equiv}.
\end{defn}


\begin{rmrk}
\label{rmrk.equiv}
We can read 
\begin{itemize*}
\item
$\f{tv}\tlatticeiff \f{tv}'$ as `$\f{tv}$ is valid if and only if $\f{tv}'$ is valid', 
and 
\item
$\f{tv}\tiffcc \f{tv}'$ as `$\f{tv}$ is true if and only if $\f{tv}'$ is true', 
and 
\item
$\f{tv}\equiv \f{tv}'$ as `$\f{tv}$ is literally equal to $\f{tv}'$'.
\end{itemize*}
Of these three notions of equivalence, $\equiv$ will be most useful to us later.
We will consider it in our predicate language in Subsection~\ref{subsect.equivalence.of.predicates}, and (for instance) it will help us to express notions of extensional equivalence in Figure~\ref{fig.easy.equivalences}.

We will call $\equiv$ \deffont{literal equivalence}, and $\equiv$ appears in the sugared syntax of Figure~\ref{fig.3.derived}.
As per the discussion in Remark~\ref{rmrk.rich.implication}, there would be no barrier to including $\tlatticeiff$ or $\tiffcc$ but we do not need to for now.
\end{rmrk}








\section{\QLogic}
\label{sect.quorum.logic}

\subsection{A preliminary: semitopologies}

\begin{rmrk}
A semitopology is like a topology, but generalised to not require that intersections of open sets necessarily be opens.

The intuition of an open set is an \emph{actionable coalition}; a set of participants with the voting power (or other resources) to make updates to their local state.
The reader might be familiar with the notion of \emph{quorum} from literature like~\cite{DBLP:journals/tocs/Lamport98,DBLP:journals/dc/MalkhiR98,naor:loacaq}, or the notion of \emph{winning coalition} in social choice theory (the theory of voting)~\cite[Item~5, page~40]{riker:thepc}.
Semitopologies are similar to this, except that open sets are always closed under arbitrary (possibly empty) unions, whereas 
\begin{itemize*}
\item
quorums may be closed under arbitrary nonempty unions (in which case the correspondence is precise, except that $\varnothing$ is an open set but is not a quorum), or 
\item
quorums are not closed under arbitrary unions (in which case a quorum corresponds to a minimal nonempty open set).
\end{itemize*}
Our presentation of semitopologies will be quite brief --- we just define the parts we need to give semantics to the modalities in \QLogic --- and the reader who wants to know more is referred to previous work~\cite{gabbay:semdca}. 
\end{rmrk} 

\begin{defn}
\label{defn.semitopology}
A \deffont{semitopology} is a pair $(\Pnt,\opens)$ where:
\begin{enumerate*}
\item
$\Pnt$ is a set of \deffont{points}, which following terminology from distributed systems one might also call \deffont{participants}, \deffont{processes}, \deffont{parties}, or \deffont{nodes}.
\item
$\opens\subseteq\powerset(\Pnt)$ is a set of \deffont{open sets} such that:
\begin{enumerate*}
\item
If $\mathcal O\subseteq\powerset(\Pnt)$ then $\bigcup\mathcal O\in\Pnt$.
In words: $\opens$ is closed under arbitrary (possibly empty) unions.
\item
$\Pnt\in\opens$.
\end{enumerate*} 
\end{enumerate*} 
We may write $\opensne$ for the set of nonempty open sets, so
$$
\opensne = \{O\in\opens \mid O\neq\varnothing\} .
$$
\end{defn}


\begin{nttn}
\label{nttn.empty.semitopology}
Suppose $(\Pnt,\opens)$ is a semitopology.
\begin{enumerate*}
\item
We call $(\Pnt,\opens)$ the \deffont{empty semitopology} when its set of points is empty --- i.e. when $\Pnt=\varnothing$ --- so that necessarily $\opens=\{\varnothing\}$.
\item
We call $(\Pnt,\opens)$ a \deffont{nonempty semitopology} when its set of points is nonempty --- i.e. when $\Pnt\neq\varnothing$.
\end{enumerate*}
\end{nttn}


\begin{defn}[All-but-$f$-many semitopology]
\label{defn.allbut}
\leavevmode
\begin{enumerate*}
\item
Suppose $X$ is a finite set and $N\in\Ngeqz$.
Define $\powerset_{\geq N}(X)$ to be the set of subsets of $X$ that contain at least $N$ elements.
In symbols:
$$
\powerset_{\geq N}(X) = \{O\subseteq X \mid \mycard O\geq N\} .
$$
\item\label{item.allbut}
Suppose $N,f\in\Ngeqz$ and $N\geq f$.
Define a semitopology $\AllBut{N}{f}$ by
$$
\AllBut{N}{f} = \bigl(\Nset{N}, \{\varnothing\}\cup\powerset_{\geq N\minus f}(\Nset{N}\bigr)\bigr).
$$
Thus: points are numbers $0\leq n<N$; and open sets are sets of points of cardinality at least $N\minus f$.
\end{enumerate*}
\end{defn}

\begin{lemm}
Suppose $N,f\in\Ngeqz$ and $N\geq f$.
Then:
\begin{enumerate*}
\item
$\AllBut{N}{f}$ is indeed a semitopology.
\item
$\AllBut{N}{f}$ is a topology precisely when $f=N$ or $f=0$.
\end{enumerate*}
\end{lemm}
\begin{proof}
We consider each part in turn:
\begin{enumerate}
\item
It is routine to check that the conditions on open sets in Definition~\ref{defn.semitopology} are satisfied.
The key observation is just the elementary fact that if $\mycard O\geq N\minus f$ and $\mycard O'\geq N\minus f$ then $\mycard(O\cup O')\geq N\minus f$.
\item
For $\AllBut{N}{f}$ to be a topology, open sets need to be closed under intersections.
It is routine to check that this can happen only when $\AllBut{N}{f}$ is the full powerset (so $f=N$) or when open sets are just $\varnothing$ and $\Pnt$ (so $f=0$).
\qedhere\end{enumerate}
\end{proof}

\begin{rmrk}
The original Paxos paper by Lamport~\cite{DBLP:journals/tocs/Lamport98} used semitopologies of the form $\AllBut{N}{f}$ where $N=2n$ and $f=n\minus 1$, for $n\geq 1$ (i.e. open sets are strict majorities of points).
This invites a question: why bother with semitopologies at all?
\begin{enumerate*}
\item
Because we can. 
More specifically, our axiomatisation is parametric over the underlying semitopology.
As good mathematicians, if we can generalise at no cost, then we should.\footnote{\dots and this applies not just in mathematics papers.  Software engineers follow a similar principle and have many names for it: e.g. encapsulation, APIs, modules, classes, typeclasses, and so on.} 
\item
Semitopology not only generalises; it simplifies.
The only properties of $\AllBut{N}{f}$ that we will need, are precisely that it is a semitopology.
If all we need about $\AllBut{N}{f}$ is its semitopological structure, then we should work with that structure because that is what actually matters.
\item
Viewing $\AllBut{N}{f}$ as a semitopology, rather than just as a collection of quorums, adds real meaning and understanding.
In particular, the notions of `dense subset' and `nonempty open interior' will be relevant to understanding the axioms, and these arise naturally in the predicates for declarative axiomatisations of distributed algorithms.
These are clearly native topological notions, so the mathematics is speaking to us here, and what it has to say is inherently topological. 
\end{enumerate*} 
\end{rmrk}


\subsection{Syntax of the logic}

We now develop our three-valued modal fixedpoint logic for reasoning about distributed systems.
We so this in the following order: \emph{signatures} $\Sigma$, \emph{models} $\mathcal M$, \emph{contexts} $\acontext$, and finally \emph{syntax}. 
We develop denotation for this syntax in Subsection~\ref{subsect.denotation}.
For quick reference, the reader can just look at the core syntax and denotation in Figures~\ref{fig.predicate.syntax} and~\ref{fig.3.phi.f}.

 


\begin{defn}[Signatures]
\label{defn.signature}
\leavevmode
\begin{enumerate*}
\item\label{item.signature.varsymb}
Fix $\tf{VarSymb}$ a countably infinite set of \deffont{variable symbols} $\thea,\theb,\thec\in\tf{VarSymb}$.
This set will never change and we may assume it henceforth.
\item
Fix a \deffont{fixedpoint variable symbol} $\tf X$ (we will only need one of these, though there is inherent obstacle to having more).
\item\label{item.predicate.syntax.signature}
A \deffont{signature} $\Sigma$ is a pair 
$$
\Sigma=(\PredSymb,\arity)
$$ 
where:
\begin{enumerate*}
\item\label{item.predsymb}
$\PredSymb$ is a (possibly finite) set of \deffont{predicate symbols} $\tf P,\tf P'\in\PredSymb$.
\item
$\arity$ is an \deffont{arity function} $\arity(\tf P)\in\Ngeqz$ mapping each predicate symbol in $\PredSymb$ to a nonnegative natural number (which is intuitively its number of arguments).
\end{enumerate*}
\item
We may write a signature $\Sigma$ in list form as 
$$
\Sigma=[\tf P,\tf P':i,\ \tf Q:j,\ \dots]
$$ 
meaning that $\tf P,\tf P',\tf Q\in\PredSymb$ and $\arity(\tf P)=\arity(\tf P')=i$ and $\arity(\tf Q)=j$ and so on.
\item\label{item.predicate.syntax.empty.signature}
If the signature in part~\ref{item.predsymb} above contains no predicate symbols then we call it \deffont{empty}; otherwise we call it non-empty.
There is only one empty signature, which we will call \deffont{the empty signature}.\footnote{It is not particularly useful.} 
\end{enumerate*}
\end{defn}


\begin{defn}[Models]
\label{defn.model}
A \deffont{model} $\mathcal M$ is a tuple $\mathcal M=(\Pnt,\opens,\Val)$ where:
\begin{enumerate*}
\item
$(\Pnt,\opens)$ is a semitopology (Definition~\ref{defn.semitopology}), which we may call the \deffont{underlying semitopology} of the model.

We may call $p\in\Pnt$ a \deffont{point} or a \deffont{participant}; also \deffont{party}, \deffont{process}, or \deffont{node} would be consistent with the literature.
For our purposes, these are all synonyms: the terminology \emph{point} is from topology; \emph{participant}, \emph{party}, and \emph{process} are from distributed systems; and \emph{node} is from graph theory (thinking of a network as a graph of communicating processes).
\item\label{item.model.val}
$\Val$ is a nonempty set of \deffont{values}.\footnote{We assume values are nonempty to avoid the degenerate case where a universal quantification over values holds just because there are no values on which to test it.  See in particular the argument in Lemma~\ref{lemm.intertwined.char}(\ref{item.intertwined.char.reverse}).}

For our purposes we will not care what values are, so long as some exist, so we do not put any type system on values or endow them with structure; but we could easily do so if required. 
\end{enumerate*} 
\end{defn}


\begin{defn}[Contexts]
\label{defn.contexts}
Suppose $\Sigma=(\PredSymb,\arity)$ is a signature, and $\mathcal M=(\Pnt,\opens,\Val)$ is a model. 
\begin{enumerate*}
\item\label{item.context.stage}
A \deffont{stage} is a nonnegative number $n\in\Time$.
\item\label{item.contexts.valuation}
A \deffont{valuation} $\avaluation$ inputs a predicate symbol $\tf P\in\PredSymb$ and a stage $n\in\Time$ and a point $p\in\Pnt$ and returns a function $\avaluation(\tf P)_{n,p}:\Val^{\arity(\tf P)}\to\THREE$.

We write the set of valuations $\Valuation(\Sigma,\Val)$ or $\Valuation(\Sigma,\mathcal M)$.
\item\label{item.contexts}
A \deffont{context} $\acontext$ is a tuple $(n,p,O,\avaluation)$ where $n\in\Time$ is a stage, $p\in\Pnt$ is a point, $O\in\opensne$ is a nonempty open set, and $\avaluation\in\tf{valuation}(\Sigma,\Val)$ is a valuation. 

We write the set of contexts $\tf{Ctx}(\Sigma,\mathcal M)$; if $\Sigma$ and $\mathcal M$ are understood, we may write this just $\tf{Ctx}$.
\item
A \deffont{context valuation} is a function $\varkappa\in \THREE^{\tf{Ctx}}$.
\end{enumerate*}
\end{defn}



\begin{figure} %
$$
\begin{array}{r@{\ }l}
t ::=& 
(a\in\tf{VarSymb}) \mid (v\in\Val) %
\\
\phi ::=&  
\tvB 
\mid t\teq t \mid (\tf P(\overbrace{t,\dots,t}^{\arity(\tf P)}) : \tf P\in\PredSymb)
\\
\mid &\tneg\phi \mid \modT\phi \mid \someone\phi \mid \Quorum\phi \mid \texi a.\phi
\\
\mid & \phi\tor\phi'  %
\\
\mid & \yesterday \phi \mid \tomorrow \phi 
\\
\mid & \tf X \mid \mu\tf X.\phi \mid (\varkappa \in \THREE^{\tf{Ctx}})
\end{array}
$$ 
\caption{Term and predicate syntax of \QLogic (Definition~\ref{defn.predicate.syntax})}
\label{fig.predicate.syntax}
\end{figure}


\begin{defn}[Syntax]
\label{defn.predicate.syntax}
Suppose $\Sigma$ is a signature (Definition~\ref{defn.signature}) and $\mathcal M$ is a model (Definition~\ref{defn.model}).
\begin{enumerate*}
\item 
Define the syntax of \deffont{\QLogic} over $\Sigma$ and $\mathcal M$ as in Figure~\ref{fig.predicate.syntax}.

In this syntax, $\texi\thea.\phi$ binds $\thea$ in $\phi$, and $\mu\tf X.\phi$ binds $\tf X$ in $\phi$.
We call a term or predicate with no free variable symbols or fixedpoint variable symbols \deffont{closed}.
\item
One special case will be of particular interest: we will write 
\begin{quote}
`suppose $\texi a.\phi$ is a closed predicate',
\end{quote}
and often we do not really care about $\texi a.\phi$ \emph{per se}, except that when this is a closed predicate it means that $a\in\tf{VarSymb}$ and $\phi$ is a predicate whose set of free variables is at most $\{a\}$.
See for example Remark~\ref{rmrk.unpack.texiunique} and Lemma~\ref{lemm.unique.affine.existence}.
\item
We define substitution of variable symbols / fixedpoint variable symbols for values / context valuations as usual, and write these $\phi[a\ssm v]$ and $\phi[\tf X\ssm \varkappa]$ respectively. %
\end{enumerate*}
\end{defn}


\begin{rmrk}
We call our logic \emph{\QLogic} to reflect that it is a logic for reasoning about how coalitions of participants can work together to carry out algorithms to achieve their goals.
In this document we will apply it to Paxos, but we see this logic as having independent interest and a broader relevance for specifying and reasoning about the kinds of algorithms that get written for distributed systems.
Their hallmark is a requirement for distributed collaborative action --- action by quorums, or actionable coalitions, or winning coalitions; these things have many names --- to progress an algorithm to its next step. 
\end{rmrk}

\begin{xmpl}
We give some example predicates and suggest intuitive meanings for them.
We do not yet have the denotation (Figure~\ref{fig.3.phi.f}) nor the rich derived syntax (Figure~\ref{fig.3.derived}) that we build later, so these examples are simple and informal, but we hope they will be helpful.

Assume our signature has a unary predicate symbol $\tf{output}$ (representing the output of some process), and assume $v,v'\in\Val$ are values.
\begin{enumerate}
\item
$\tf{output}(v)$ means `where we are, today, if the process succeeds, then it outputs $v$' and $\tneg\tf{output}(v')$ means `where we are, if the process succeeds, then it does not output $v'$'.
\item
$\modT\tf{output}(v)$ means `where we are, today, the process succeeds and outputs $v$'. 
\item
$\someone\tf{output}(v)$ means `there is a place where, today, if the process succeeds, then it outputs $v$'.
\item
$\tomorrow\modT\tf{output}(v)$ means `where we are, tomorrow, the process succeeds and outputs $v$'.
\item
$\yesterday\modT\tf{output}(v)$ means `where we are, yesterday, the process succeeded and output $v$'.
This assertion is subtly stronger than it might seem; our notion of time starts at day~0, at which by convention all assertions about yesterday are false.
See Remark~\ref{rmrk.further.comments.denotation}(\ref{item.interpret.yesterday}).
\item
$\tvB$ means `where we are, today, the process (which we have not named with a predicate symbol) crashes'.
\end{enumerate}
More sophisticated examples will appear later; notably in the derived syntax in Figure~\ref{fig.3.derived}, in Remark~\ref{rmrk.comments.on.semantics}, and of course when we apply our logic to axiomatise distributed algorithms.
\end{xmpl}


\begin{rmrk}[Three pedant points]
\label{rmrk.pedant.point}
The thing about being pedantic is that it does not matter, until it does.
We now discuss three pedant points, which may not matter \dots until perhaps a moment arrives then they do:
\begin{enumerate}
\item
We may treat $\Sigma$ and $\mathcal M$ as fixed parameters, and elide them.  
In practice this means that we may write `suppose $\phi$ is a closed predicate' instead of `suppose $\phi$ is a closed predicate in the syntax of \QLogic over $\Sigma$ and $\mathcal M$'.
But note next point.
\item\label{item.model.matters}
Our syntax from Definition~\ref{defn.predicate.syntax} and Figure~\ref{fig.predicate.syntax} contains denotational elements; $v\in\Val$ and $\varkappa\in\THREE^{\tf{Ctx}(\Sigma,\mathcal M)}$.

So when we write `suppose $\phi$ is a closed predicate' then this statement is perfectly correct and well-defined but it \emph{really does matter} what the signature and model are (even if the \emph{precise} choice of model will may not matter to the theorems we prove), because closed predicates $\phi$ can include elements from the denotation such as $\varkappa\in\THREE^{\tf{Ctx}(\Sigma,\mathcal M)}$ (this function-space is determined by both the signature $\Sigma$ and the model $\mathcal M$) and $v\teq v'$ for some $v,v'\in\Val$ ($\Val$ is determined by the choice of model $\mathcal M$).

Adding denotational elements to syntax avoids us having to put this information in the context (which is already a $4$-tuple, as per Definition~\ref{defn.contexts}(\ref{item.contexts})) and simplifies the presentation.
This is a technical device with a long pedigree.\footnote{Adding elements from a model to syntax, and conversely building a model out of syntax, is a standard technical device in logic, with a pedigree that goes back (at least) to the proofs of the upward and downward L\"owenheim-Skolem theorems from the 1920s.  A clear survey of these methods is in Chapter~2 (``Models constructed from constants'') of~\cite{chang:modt}.}

Note that (by design!) the choice of model makes a difference to validity.
For example: in Definition~\ref{defn.phi.equivalent} we define equivalence of predicates (and we explicitly mention $\Sigma$ and $\mathcal M$);
if $\mathcal M$ has only a single value $v$, then $\tall a,b.a\teq b$ and $\tall a.a\teq v$ are equivalent to $\tvT$; and if $\mathcal M$ has more than one value then $\tall a,b.a\teq b$ is equivalent to $\tvF$ --- $\tall a.\phi$ is shorthand for $\tneg\texi a.\tneg\phi$; see Figure~\ref{fig.3.derived}.
In contrast: in Lemma~\ref{lemm.equiv} we elide $\Sigma$ and $\mathcal M$ for brevity, and just write `suppose $\phi$ and $\phi'$ are closed predicates and $\acontext$ is a context', without specifying over what signature or model, because this is clear in context.

The reader can mentally add $\Sigma$ and $\mathcal M$ everywhere if they like.
We will always be precise and explicit where it counts.
\item\label{item.pedant.texi.v}
We may quantify over $v\in\Val$ in our formal syntax of predicates, writing things like $\texi v.\texi v'.v\teq v'$ as predicates.
Technically this is wrong and is not syntax! 
Figure~\ref{fig.predicate.syntax} permits e.g. $\texi a.\texi a'.a\teq a'$ where $a,a'\in\tf{VarSymb}$ (Definition~\ref{defn.signature}(\ref{item.signature.varsymb})); it does not permit $\texi v$ in the formal syntax of predicates where $v\in\Val$, so writing e.g. $\texi v.\texi v'.v\teq v'$ as if it were formal syntax of \QLogic is an abuse of notation.

However, $a$ and $a'$ get instantiated to values $v$ and $v'$ in our denotation (see the line for $\texi$ in Figure~\ref{fig.3.phi.f}).
It seems not worth the disruption of maintaining a distinction between $\texi a.\phi$ (which is syntax, assuming $\phi$ is) and $\texi v.\phi[a\ssm v]$ (which looks like syntax but technically is not).
So in practice this seems to us one example of an abuse of formal notation that is positively useful and clarifies meaning.

If this bothers the reader, then every time they see $\texi v.\phi$ (and $\tall v.\phi$, and similar) in our text, they can either declare that $v$ here is \emph{actually} a variable symbol in $\tf{VarSymb}$ which we happen to have called $v$ (and not a value in $\Val$ after all), or they can pick a fresh variable symbol $a$ and replace every unbound $v$ in $\phi$ with $a$.
Either way they will get well-formed syntax and everything will be as it should be.

Notation exists to serve clarity, and not the other way around, so we will simply ignore this henceforth and carry on writing $\texi v.\phi$, just whenever we feel like it, and this will denote just what it says: `there exists a $v$ such that $\phi$'.
Meaning should always be clear. 

Note that an unbound value $v$ in a predicate is legal syntax as per Figure~\ref{fig.predicate.syntax}; 
e.g. in `$\texi a.a\teq v$', $v$ is a value, and this is fine; the (constructive and useful) abuse of notation comes when we write $\texi v$. 
\end{enumerate}
\end{rmrk}

\subsection{Denotation (semantics) of the syntax}
\label{subsect.denotation}


\begin{figure} %
$$
\begin{array}{r@{\ }l@{\quad}l}
\modellabel{\tvB}{n,p,O,\avaluation} =&\tvB
\\
\modellabel{v\teq v'}{n,p,O,\avaluation} =&
\begin{cases}
\tvT & v=v'
\\
\tvF & v\neq v'
\end{cases}
\\
\modellabel{\tf P(v_1,\dots,v_{\arity(\tf P))}}{n,p,O,\avaluation} =&\avaluation(\tf P)_{n,p}(v_1,\dots,v_{\arity(\tf P)}) %
\\[2ex]
\modellabel{\phi\tor\phi'}{n,p,O,\avaluation} =& \modellabel{\phi}{n,p,O,\avaluation}\tor\modellabel{\phi'}{n,p,O,\avaluation}
\\
\modellabel{\tneg\phi}{n,p,O,\avaluation} =& \tneg\,\modellabel{\phi}{n,p,O,\avaluation}  
\\
\modellabel{\modT\phi}{n,p,O,\avaluation} =& \modT\modellabel{\phi}{n,p,O,\avaluation}  
\\
\modellabel{\texi\thea.\phi}{n,p,O,\avaluation} =& \bigvee_{v{\in}\Val} \modellabel{\phi[\thea\ssm v]}{n,p,O,\avaluation}
\\
\modellabel{\someone\phi}{n,p,O,\avaluation} =& \bigvee_{p'\in O} \modellabel{\phi}{n,p',O,\avaluation} 
\\
\modellabel{\Quorum\phi}{n,p,O,\avaluation} =& \bigvee_{O'\in\opensne} \modellabel{\phi}{n,p,O',\avaluation} 
\\[2ex]
\modellabel{\yesterday\phi}{n,p,O,\avaluation} =&
\bigvee_{\f{max}(0,n\minus 1)\leq n'<n}\modellabel{\phi}{n',p,O,\avaluation}
\\
\modellabel{\tomorrow\phi}{n,p,O,\avaluation} =&
\modellabel{\phi}{n\plus 1,p,O,\avaluation}
\\[2ex]
\modellabel{\mu\tf X.\phi}{n,p,O,\avaluation}=&
\bigl(\bigwedge\{ \varkappa\in\THREE^{\tf{Ctx}}\mid \model{\phi[\tf X\ssm\varkappa]}\leq\varkappa\}\bigr)(n,p,O,\avaluation)
\figskip
\modellabel{\varkappa}{n,p,O,\avaluation}=&
\varkappa(n,p,O,\avaluation)
\end{array}
$$ 
\caption{Denotational semantics of \QLogic (Definition~\ref{defn.denotation})}
\label{fig.3.phi.f}
\end{figure}


\begin{figure} %
$$
\begin{array}{r@{\ }l@{\ \ }r@{\ }l}
\tvF =& \modT\tvB
&
\modellabel{\tvF}{\acontext} =&\tvF
\figskip
\tvT =& \tneg\tvF
&
\modellabel{\tvT}{\acontext} =&\tvT
\\[2ex]
\modF\phi =& \modT\tneg\phi
&
\modellabel{\modF\phi}{\acontext} =&\modF\modellabel{\phi}{\acontext}
\figskip
\modB\phi =& \tneg\modT\phi\tand\tneg\modF\phi
&
\modellabel{\modB\phi}{\acontext} =&\modB\modellabel{\phi}{\acontext}
\figskip
\modTB\phi =& \tneg\modT\tneg\phi
&
\modellabel{\modTB\phi}{\acontext} =&\modTB\modellabel{\phi}{\acontext}
\figskip
\modTF\phi =& \modT\phi\tor\modF\phi
&
\modellabel{\modTF\phi}{\acontext} =&\modTF\modellabel{\phi}{\acontext}
\figskip
\modFB\phi =& \modF\phi\tor\modB\phi
&
\modellabel{\modFB\phi}{\acontext} =&\modFB\modellabel{\phi}{\acontext}
\\[2ex]
\phi\tand\phi' =& \tneg(\tneg\phi\tor \tneg\phi')
&
\modellabel{\phi\tand\phi'}{\acontext} =&\modellabel{\phi}{\acontext}\tand\modellabel{\phi'}{\acontext}
\figskip
\phi\tnotor\phi' =& \tneg\phi\tor\phi'
&
\modellabel{\phi\tnotor\phi'}{\acontext} =&\modellabel{\phi}{\acontext}\tnotor\modellabel{\phi'}{\acontext}
\figskip
\phi\timpc\phi' =& \phi\tnotor\modT\phi'
&
\modellabel{\phi\timpc\phi'}{\acontext} =&\modellabel{\phi}{\acontext}\timpc\modellabel{\phi'}{\acontext}
\figskip
\phi\equiv\phi' =& (\modT\phi\tand\modT\phi')\tor(\modB\phi\tand\modB\phi') \tor(\modF\phi\tand\modF\phi') \hspace{-5em}
\\
&&
\modellabel{\phi\equiv\phi'}{\acontext} =&\modellabel{\phi}{\acontext}\equiv\modellabel{\phi'}{\acontext}
\\[2ex]
\tall a.\phi =& \tneg\texi a.\tneg\phi
&
\modellabel{\tall a.\phi}{\acontext} =&\bigwedge_{v{\in}\Val}\modellabel{\phi[a\ssm v]}{\acontext}
\figskip
\everyone \phi =& \tneg\someone\tneg\phi
&
\modellabel{\everyone\phi}{n,p,O,\avaluation} =&\bigwedge_{p'{\in}\Pnt}\modellabel{\phi}{n,p',O,\avaluation}
\figskip
\Coquorum\phi =& \tneg\Quorum\tneg\phi
&
\modellabel{\Coquorum\phi}{n,p,O,\avaluation} =&\bigwedge_{O'{\in}\opensne}\modellabel{\phi}{n,p,O',\avaluation}
\\[2ex]
\QuorumBox\phi =& \Quorum\everyone\phi
&
\modellabel{\QuorumBox\phi}{n,p,O,\avaluation} =&\bigvee_{O'{\in}\opensne}\bigwedge_{p'\in O'}\modellabel{\phi}{n,p',O',\avaluation}
\figskip
\CoquorumDiamond\phi =& \Coquorum\someone\phi
&
\modellabel{\CoquorumDiamond\phi}{n,p,O,\avaluation} =&\bigwedge_{O'{\in}\opensne}\bigvee_{p'\in O'}\modellabel{\phi}{n,p',O',\avaluation}
\figskip
\someoneAll\phi =& \Quorum\someone\phi
&
\modellabel{\someoneAll\phi}{n,p,O,\avaluation} =&\bigvee_{O'{\in}\opensne}\bigvee_{p'\in O'}\modellabel{\phi}{n,p',O',\avaluation}
\figskip
\everyoneAll\phi =& \Coquorum\everyone\phi
&
\modellabel{\everyoneAll\phi}{n,p,O,\avaluation} =&\bigwedge_{O'{\in}\opensne}\bigwedge_{p'\in O'}\modellabel{\phi}{n,p',O',\avaluation}
\\[2ex]
\forever\phi =& \mu\tf X.\tomorrow(\phi\tand \tf X)
&
\modellabel{\forever\phi}{n,p,O,\avaluation} =&\bigwedge_{n'> n}\modellabel{\phi}{n',p,O,\avaluation}
\figskip
\sometime\phi =& \mu\tf X.\tomorrow(\phi\tor \tf X)
&
\modellabel{\sometime\phi}{n,p,O,\avaluation} =&\bigvee_{n'> n}\modellabel{\phi}{n',p,O,\avaluation}
\figskip
\infinitely\phi =& \forever\sometime\phi
&
\modellabel{\infinitely\phi}{n,p,O,\avaluation} =&\bigwedge_{n'> n}\bigvee_{n''>n'}\modellabel{\phi}{n'',p,O,\avaluation}
\figskip
\final\phi =& \sometime\forever\phi
&
\modellabel{\final\phi}{n,p,O,\avaluation} =&\bigvee_{n'> n}\bigwedge_{n''>n'}\modellabel{\phi}{n'',p,O,\avaluation}
\figskip
\recent\phi =& \mu\tf X.\yesterday(\phi\tor\tf X)
&
\modellabel{\recent\phi}{n,p,O,\avaluation} =&\bigvee_{0\leq n'<n}\modellabel{\phi}{n',p,O,\avaluation}
\figskip
\urecent\phi =& \mu\tf X.(\phi\tor\yesterday\tf X)
&
\modellabel{\urecent\phi}{n,p,O,\avaluation} =&\bigvee_{0\leq n'\leq n}\modellabel{\phi}{n',p,O,\avaluation}
\end{array}
$$
$$
\begin{array}{r@{\ }l}
\mrup{a}{\phi}{v} =&
\mu\tf X.(\yesterday(\modT\phi[a\ssm v] \tor ((\tneg\modT\texi a.\phi) \tand \tf X)))
\figskip
\modellabel{\mrup{a}{\phi}{v}}{n,p,O,\avaluation} =&
\begin{cases}
\tvF 
&\Forall{0\leq n'<n} \modellabel{\texi a.\phi}{n',p,O,\avaluation}\neq\tvT
\\
\modT\modellabel{\phi[\thea\ssm v]}{\f{max},p,O,\avaluation}
&\Exists{0\leq n'<n} \modellabel{\texi a.\phi}{n',p,O,\avaluation}=\tvT
\end{cases}
\\
&\qquad\text{where}\quad
\f{max}=\bigvee\{0\leq n'<n \mid \modellabel{\texi a.\phi}{n',p,O,\avaluation}=\tvT\}
\end{array}
$$ 
\caption{Derived syntactic sugar expressions (Definition~\ref{defn.3.derived})}
\label{fig.3.derived}
\end{figure}



Recall from Definition~\ref{defn.signature}(\ref{item.predicate.syntax.signature}) the notion of a \emph{signature} and from Definition~\ref{defn.model} the notion of a \emph{model}:
\begin{defn}[Denotation]
\label{defn.denotation}
Suppose $\Sigma=(\PredSymb,\arity)$ is a signature, and $\mathcal M$ is a model. 
Suppose $\phi$ is a closed predicate (Definition~\ref{defn.predicate.syntax}(1)).
Define the \deffont{denotation} of $\phi$ 
$$
\model{\phi}\in\THREE^{\tf{Ctx}(\Sigma,\mathcal M)}
\quad\text{such that}\quad 
\acontext \longmapsto \modellabel{\phi}{\acontext}
$$
by the rules in Figure~\ref{fig.3.phi.f}.

We will usually treat $\Sigma$ and $\mathcal M$ as fixed parameters, so we may write $\tf{Ctx}(\Sigma,\mathcal M)$ just as $\tf{Ctx}$.
\end{defn}




\begin{defn}[Derived expressions]
\label{defn.3.derived}
The syntax of \QLogic in Figure~\ref{fig.predicate.syntax} and Definition~\ref{defn.predicate.syntax} is very expressive.
In the left-hand column of Figure~\ref{fig.3.derived} we define some derived expressions --- i.e. \emph{syntactic sugar}.
For clarity and convenience we unpack their denotation in the right-hand column.

In principle we could fully unpack our sugar down to the primitive syntax from Figure~\ref{fig.3.phi.f}, but in practice we work directly with Figure~\ref{fig.3.derived} because this is more convenient.
We will discuss syntax further below, notably in Remarks~\ref{rmrk.initial.comments.denotation}, \ref{rmrk.further.comments.denotation}, and~\ref{rmrk.comments.on.semantics}.
\end{defn}


\begin{rmrk}
Usually when we talk about denotation as defined in Definition~\ref{defn.denotation} above, the signature $\Sigma$ and model $\mathcal M$ are either understood or do not matter, so we may not mention them explicitly.
Meaning will always be clear.
\end{rmrk}

\begin{rmrk}[Initial comments on the denotation]
\label{rmrk.initial.comments.denotation}
We make some initial comments on Figure~\ref{fig.3.phi.f}:
\begin{enumerate*}
\item\label{item.first.mu}
$\mu\tf X.\phi$ is a \emph{fixedpoint operator}.
This is a standard method to include the power of recursion / iteration in a logic; an instance of $\tf X$ inside $\mu\tf X.\phi$ can be thought of as a recursive call back to $\phi$.
$\mu\tf X.\phi$ is what would be written $\mathrm{LFP}_{\tf X}\phi$ in~\cite[Chapter~4]{immerman:desc} (a Wikipedia article is also accessible and clear~\cite{wiki:Fixedpoint_logic}).
We will discuss this in Subsection~\ref{subsection.muX}.
\item
In the clause $\model{\tvB}=\tvB$ in Figure~\ref{fig.3.phi.f}, the $\tvB$ on the left is the syntax symbol from Figure~\ref{fig.predicate.syntax}, and the $\tvB$ on the right is $\tvB$-the-truth-value from Definition~\ref{defn.THREE}.
\item
In the clause for $\model{v\teq v'}$ in Figure~\ref{fig.3.phi.f}, the $\teq$ on the left is the syntax symbol from Figure~\ref{fig.predicate.syntax}, and the $\tvT$, $\tvF$, and $=$ on the right are truth-values and (actual) equality.
\item
Similarly for the clauses for $\model{\tneg\phi}$ and $\model{\phi\tor\phi'}$ and $\model{\modT\phi}$; the symbols $\tneg$, $\tor$, and $\modT$ on the left are syntax, and the symbols $\tneg$, $\tor$, and $\modT$ on the right refer to functions on truth-values as per Figure~\ref{fig.3}.
\end{enumerate*} 
\end{rmrk}


\begin{rmrk}[Further comments] 
\label{rmrk.further.comments.denotation}
\leavevmode
\begin{enumerate}
\item\label{item.interpret.yesterday}
The clause 
$\modellabel{\yesterday\phi}{n,p,O,\avaluation} = \bigvee_{\f{max}(0,n\minus 1)\leq n'<n}\modellabel{\phi}{n',p,O,\avaluation}$ is just a one-line encoding of
$$
\modellabel{\yesterday\phi}{n,p,O,\avaluation}=
\begin{cases}
\tvF & n=0
\\
\modellabel{\phi}{n\minus 1,p,O,\avaluation} & n>0 . 
\end{cases}
$$
\item
We see from the right-hand column in Figure~\ref{fig.3.derived} that $\recent\phi$ returns the greatest truth-value that $\phi$ has taken in the past, and $\urecent\phi$ returns the greatest truth-value that $\phi$ has taken today or in the past.

The notation $\urecent$ (`$\recent$', underlined) is a pun\footnote{A bad one perhaps, and the reader may or may not find this pun-acceptable.} on $<$ and $\leq$ (`$<$', underlined).
\item
Informally, we see that the only naturally \emph{non-monotone} operator in the (primitive or derived) syntax is $\tneg$, which is anti-monotone.
This design feature is what powers Lemma~\ref{lemm.positive.monotonicity}.
\end{enumerate}
\end{rmrk}


\begin{nttn}
\label{nttn.logic.terminology}
We take a moment to recall some standard conventions and terminology:
\begin{enumerate*}
\item
Implications $\tnotor$ and $\timpc$ associate to the right, so that $\phi\timpc\psi\timpc\chi$ means $\phi\timpc(\psi\timpc\chi)$.
\item
Unary connectives (like $\tneg$, $\modT$, and $\someone$) bind more strongly than binary connectives (like $\tand$ and $\tnotor$), so that $\QuorumBox\tneg\phi\tand\psi$ means $(\QuorumBox\tneg\phi)\tand\psi$.
\item
The scope of a quantification $\texi a$ and $\tall a$ extends as far to the right as possible. 
Thus $\tall\thea.\phi\tand\psi$ means $\tall\thea.(\phi\tand\psi)$.
Note however that $\QuorumBox\texi\thea.\phi\tand\psi$ still means $(\QuorumBox(\texi\thea.\phi))\tand\psi$, because the $\texi$ quantifier is trapped in the scope of the $\QuorumBox$. 
\item
All the above said, we will (as is usual) use spacing and some common sense to convey scope.
For example 
$$
\texi\thea.\phi\mathbin{\tnotor}\psi
\quad\text{means}\quad
\texi\thea.(\phi\tnotor\psi), 
$$
and
$$
\texi\thea.\phi \ \tnotor \psi
\quad\text{means}\quad
(\texi\thea.\phi)\tnotor\psi. 
$$
In practice, meaning will always be clear.
\item
\begin{enumerate*}
\item\label{item.nullary.ps}
If $\arity(\tf P)=0$ then we may call $\tf P$ \deffont{nullary} or a \deffont{predicate constant symbol}.\footnote{$\tvT$, $\tvB$, and $\tvF$ are predicate constant symbols, but they are special so we give them a distinct identity in the syntax.}
\item\label{item.unary.ps}
If $\arity(\tf P)=1$ then we may call $\tf P$ a \deffont{unary} predicate symbol.
\item\label{item.atomic.predicate}
We may call a predicate of the form $\tf P(v_1,\dots,v_n)$ or $\tneg\tf P(v_1,\dots,v_n)$ an \deffont{atomic predicate}. 
\end{enumerate*}
\end{enumerate*} 
\end{nttn}



\subsection{The fixedpoint operator $\mu \tf X$}
\label{subsection.muX}

\begin{rmrk}
\label{rmrk.comments.on.semantics}
As per Remark~\ref{rmrk.initial.comments.denotation}(\ref{item.first.mu}), our fixedpoint operator $\mu\tf X.\phi$ is a method to include the power of recursion / iteration in the logic.

A mention of $\tf X$ inside $\phi$ in $\mu\tf X.\phi$ can be thought of as a recursive call back to $\phi$, and we can think of $\mu\tf X.\phi$ as encoding an iterative loop that repeatedly executes a computation encoded in $\phi$, and loops if it encounters $\tf X$ in $\phi$.

With this intuition, we can look at some of the derived syntax in Figure~\ref{fig.3.derived} that uses $\mu$:
\begin{enumerate}
\item
$\recent\phi=\mu\tf X.\yesterday(\phi\tor\tf X)$ can be imagined as the following loop:
\begin{enumerate*}
\item
Set an accumulator to $\tvF$.\footnote{We mean `accumulator' in the sense of an `accumulator variable' inside a machine that (we imagine) is running this pseudocode.  This is not a technical term from logic.}
\item\label{recent.step.0}
If we are at day~0 then terminate and return the value of the accumulator; otherwise go to yesterday.
\item
Compute the truth-value of $\phi$, and $\tor$ this truth-value with the accumulator.
\item
Loop back to step~\ref{recent.step.0}.
\end{enumerate*} 
\item
$\forever\phi=\mu\tf X.\tomorrow(\phi\tand\tf X)$ can be viewed as the following loop:
\begin{enumerate*}
\item
Set an accumulator to $\tvT$.
\item\label{forever.step.0}
Go to tomorrow.
\item
Compute the truth-value of $\phi$, and $\tand$ this truth-value with the accumulator.
\item
Loop back to step~\ref{forever.step.0}.
\item
After infinite time (see Remark~\ref{rmrk.oo} about why that makes sense) return the limit value of the accumulator.
\end{enumerate*} 
\item
$\sometime\phi=\mu\tf X.\tomorrow(\phi\tor\tf X)$ can be viewed as executing the following loop:
\begin{enumerate*}
\item
Set an accumulator to $\tvF$.
\item\label{sometime.step.0}
Go to tomorrow.
\item
Compute the truth-value of $\phi$ and $\tor$ with the accumulator.
\item
Loop back to step~\ref{sometime.step.0}.
\item
After infinite time, return the limit value of the accumulator.
\end{enumerate*}
\item\label{item.mru.program}
The clause 
$\mrup{a}{\phi}{v} = \mu\tf X.(\yesterday(\modT\phi[a\ssm v] \tor ((\tneg\modT\texi a.\phi) \tand \tf X)))$
looks more imposing, but the idea is simply to 
\begin{itemize*}
\item
iterate back in time looking for when $\texi a.\phi$ has truth-value $\tvT$, i.e. when \emph{some} value makes $\phi(a)$ be true, and 
\item
if we find such a time then we check if \emph{at that time} $\phi[a\ssm v]$ has truth-value $\tvT$ and 
\item
only if both steps above succeed, do we return $\tvT$.
\end{itemize*}
Thus, we return $\tvT$ when $v$ is one of the values that \emph{most recently} makes $\phi(a)$ be true.

In detail:
\begin{enumerate*}
\item\label{mu.step.0}
If we are at day~0, terminate with value $\tvF$; otherwise go to yesterday.
\item
If $\phi[a\ssm v]$ has truth-value $\tvT$ then terminate with value $\tvT$; otherwise continue to the next step.
\item
If $\texi a.\phi$ has truth-value $\tvT$ then terminate with value $\tvF$; otherwise go to step~\ref{mu.step.0}.
\end{enumerate*}
\end{enumerate}
\end{rmrk}

\begin{rmrk}[Comment on the omniscient observer]
\label{rmrk.oo}
The programs in Remark~\ref{rmrk.comments.on.semantics} above are executed by an omniscient observer who stands outside of the model's notion of time and space.
This is why (for example) $\forever$ can run an infinite loop --- the omniscient observer can observe it because the observer is looking at the model from outside.\footnote{Let's make this concrete.  Consider the natural numbers $\Ngeqz$.  We can easily observe this and note that $n\plus 1\gneq n$ for every $n\in\Ngeqz$; whereas a local observer trapped on the number line and constrained to perceive this line as the passage of time can make similar observation, but from the point of view of that observer, the calculation involved is neverending.  

Note that even if we could speak with this local observer, it would be no good to appeal to it to `just realise that it is on the number line', because in order to check this to a mathematical certainty the observer would have to \dots traverse all of its local notion of time to \emph{check}, amongst other things, that $n\plus 1\gneq n$.  After all, the observer might be on $\mathbb Z/p$ for some astronomically enormous $p$, i.e. their time might be on a large loop.  

The reader interested in a branch of mathematics explicitly devoted to studying these questions, and who is not afraid by the possibility of ending up seriously doubting the nature of their own reality, can start by reading up on nonstandard models of arithmetic.} 
\end{rmrk}

\begin{rmrk}
The rest of this Subsection will show how our denotation for $\mu$ from Figure~\ref{fig.3.phi.f} does indeed yield a fixedpoint operator.
The culminating result is Proposition~\ref{prop.mu.fixedpoint} (which just says that $\mu$ is indeed a fixedpoint operator) but the key technical lemma is Lemma~\ref{lemm.positive.monotonicity} (monotonicity of positive predicates).

The reader is welcome to read the rest of this Subsection, but this is all just a standard use of the Knaster-Tarski fixedpoint theorem~\cite[Theorem~2.35, page~50]{priestley:intlo}.

In fact, the real design effort has already happened: we designed the derived syntax of Figure~\ref{fig.3.derived} so that it would respect the positivity condition of Definition~\ref{defn.positive}, without having to fully expand it out.
For instance, we defined $\modTB\phi$ to be $\tneg\modT\tneg\phi$ (which clearly puts $\phi$ in a positive position) instead of $\modT\phi\tor\modB\phi$ (which is logically equivalent, but which does not).
\end{rmrk} 



\begin{defn}
\label{defn.three.ctx.lattice}
Recall from Remark~\ref{rmrk.three.elementary}(\ref{item.three.elementary.bounded.lattice}) that $\THREE$ with ordering $\tvF<\tvB<\tvT$ is a finite lattice.
As is standard, the function space $\THREE^{\tf{Ctx}}$ (functions from $\tf{Ctx}$ to $\THREE$) inherits this ordering and lattice structure \deffont{pointwise}, as indicated in Figure~\ref{fig.chi.3.lattice}.

This makes $\THREE^{\tf{Ctx}}$ into a \deffont{bounded complete lattice}, meaning that it has a least and greatest element, and arbitrary joins (least upper bounds) and meets (greatest lower bounds).\footnote{Having a least and greatest element follows from having arbitrary joins and meets; the least element is the join of the empty set, and the greatest element is the meet of the empty set.  So in fact a complete lattice is already a bounded complete lattice.} 
\end{defn}


\begin{figure}
$$
\begin{array}{r@{\ }l}
\tbot =& \lambda\acontext.\tvF
\\
\ttop =& \lambda\acontext.\tvT
\\
\bigvee\mathcal X =& \lambda\acontext.\bigvee\{ \varkappa\,\acontext \mid \varkappa\in\mathcal X\}
\\
\bigwedge\mathcal X =& \lambda\acontext.\bigwedge\{ \varkappa\,\acontext \mid \varkappa\in\mathcal X\}
\\[2ex]
\varkappa\leq\varkappa' \liff& \Forall{\acontext} \varkappa\,\acontext\leq\varkappa'\,\acontext
\end{array}
$$
Above, $\mathcal X\subseteq \THREE^{\tf{Ctx}}$ and $\varkappa,\varkappa'\in\THREE^{\tf{Ctx}}$. 
\caption{Bounded complete lattice structure on the function space $\THREE^{\tf{Ctx}}$}
\label{fig.chi.3.lattice}
\end{figure}

\begin{rmrk}
\label{rmrk.leq.phi}
Our denotation $\model{\text{-}}$ from Definition~\ref{defn.denotation} and Figure~\ref{fig.3.phi.f} maps predicates to functions in $\THREE^{\tf{Ctx}}$.
One particular consequence of Definition~\ref{defn.three.ctx.lattice} is that (denotations of) predicates are naturally ordered, and we can write 
$$
\model{\phi}\leq\model{\psi}
\quad\text{which just means}\quad
\Forall{\acontext\in\tf{Ctx}}\modellabel{\phi}{\acontext}\leq\modellabel{\psi}{\acontext}.
$$
We use this fact for example in Lemma~\ref{lemm.positive.monotonicity}.
We can also use this fact, and the complete lattice structure on $\THREE^{\tf{Ctx}}$ from Definition~\ref{defn.three.ctx.lattice}, to simplify the presentation of the clause for $\mu\tf X.\phi$ in Figure~\ref{fig.3.phi.f} as follows:
$$
\model{\mu\tf X.\phi}=
\bigwedge\{ \varkappa\in\THREE^{\tf{Ctx}}\mid \model{\phi[\tf X\ssm\varkappa]}\leq\varkappa\} .
$$
\end{rmrk}

\begin{defn}
\label{defn.positive}
Recall the predicate syntax from Figure~\ref{fig.predicate.syntax} %
and suppose $\phi$ is a predicate.
Then:
\begin{enumerate*}
\item\label{item.positive.occurrence}
Call a free occurrence of $\tf X$ in $\phi$ 
\begin{itemize*}
\item
\deffont{positive} when the occurrence appears under an even number (which may be $0$) of negations $\tneg$, and 
\item
\deffont{negative} when the occurrence appears under an odd number of negations.
\end{itemize*}
For example: all free occurrences of $\tf X$ are positive in $\tf X$ and $\mu\tf X.\tneg\tf X$ and $\yesterday\tf X$; and the free occurrence of $\tf X$ in $\tneg\tf X$ is negative. 
\item
Call a predicate $\phi$ \deffont{positive} when every free occurence of $\tf X$ in $\phi$ is positive.
\end{enumerate*}
\end{defn}


\begin{lemm}
\label{lemm.positive.monotonicity}
Suppose $\mu\tf X.\phi$ is a closed predicate.
Then:
\begin{enumerate*}
\item
If $\phi$ is positive then $\model{\phi}$ is monotone as a function from $\THREE^{\tf{Ctx}}$ with the partial order $\leq$ from Figure~\ref{fig.chi.3.lattice}, to itself.
In symbols, for every $\varkappa,\varkappa'\in\THREE^{\tf{Ctx}}$:
$$
\varkappa\leq\varkappa'
\quad\text{implies}\quad
\model{\phi[\tf X\ssm\varkappa]}\leq\model{\phi[\tf X\ssm\varkappa']} . 
$$
\item
If $\tneg\phi$ is positive then $\model{\phi}$ is antimonotone as a function from $\THREE^{\tf{Ctx}}$ with the partial order $\leq$ from Figure~\ref{fig.chi.3.lattice}, to itself.
In symbols, for every $\varkappa,\varkappa'\in\THREE^{\tf{Ctx}}$:
$$
\varkappa\leq\varkappa'
\quad\text{implies}\quad
\model{\phi[\tf X\ssm\varkappa']}\leq\model{\phi[\tf X\ssm\varkappa]} . 
$$
\end{enumerate*}
(The meaning of $\model{\text{-}}\leq\model{\text{-}}$ is unpacked in Remark~\ref{rmrk.leq.phi}.)

\end{lemm}
\begin{proof}
Both parts follow by a routine simultaneous induction on the syntax of $\phi$, using the rules in Figure~\ref{fig.3.phi.f}.
\end{proof}

\begin{prop}
\label{prop.mu.fixedpoint}
Suppose $\mu\tf X.\phi$ is a closed predicate and $\phi$ is positive.
Then:
\begin{enumerate*}
\item
$\model{\mu\tf X.\phi}$ is a least fixedpoint of the function $\varphi:\THREE^{\tf{Ctx}}\to\THREE^{\tf{Ctx}}$ given by
$$
\varphi(\varkappa)=\lambda\acontext.\modellabel{\phi[\tf X\ssm\varkappa]}{\acontext}. 
$$
\item
If $\acontext$ is a context then
$$
\modellabel{\mu\tf X.\phi}{\acontext}=\modellabel{\phi[\tf X\ssm \model{\mu\tf X.\phi}]}{\acontext}.
$$
\end{enumerate*}
\end{prop}
\begin{proof}
By the Knaster-Tarski fixedpoint theorem for $\THREE^{\tf{Ctx}}$ as a bounded lattice as per Definition~\ref{defn.three.ctx.lattice}.
A concise but exceptionally clear presentation is in~\cite[Theorem~2.35, page~50]{priestley:intlo} (a Wikipedia article~\cite{wiki:Knaster-Tarski_theorem} is also accessible and clear). 
\end{proof}



\subsection{Validity}



We defined in Definition~\ref{defn.tv.ment} what it means for a truth-value to be \emph{valid}; now we extend this notion to predicates-in-context: 
\begin{defn}[Validity]
\label{defn.validity.judgement}
Suppose that:
\begin{itemize*}
\item
$\Sigma$ is a signature (Definition~\ref{defn.signature}(\ref{item.predicate.syntax.signature})).
We take $\Sigma$ to be fixed for the rest of this Definition, which means that we do not explicitly include it (e.g. as a subscript) in the notation.
\item
$\mathcal M$ is a model (Definition~\ref{defn.model}).
Again, we fix $\mathcal M$ as a parameter for the rest of this definition.
\end{itemize*}
Suppose further that:
\begin{itemize*}
\item
$\acontext=(n,p,O,\avaluation)$ is a context in $\tf{Ctx}(\Sigma,\mathcal M)$, so that by Definition~\ref{defn.contexts}(\ref{item.contexts}) $n\in\Time$ and $p\in\Pnt$ and $O\in\opensne$ and $\avaluation$ is a valuation to $\mathcal M$.
\item
$\phi$ is a closed predicate in $\Sigma$ and $\mathcal M$ (cf. Remark~\ref{rmrk.pedant.point}(\ref{item.model.matters})).
\end{itemize*}
Then we define: %
\begin{enumerate*}
\item\label{item.ctx.ment.phi}
$\acontext\ment\phi$ means $\ment\modellabel{\phi}{\acontext}$, which by Definition~\ref{defn.tv.ment}(\ref{item.tv.ment}) means $\modellabel{\phi}{\acontext}\in\threeValid$, which by Definition~\ref{defn.tv.ment}(\ref{item.tv.ment.correct}) means $\modellabel{\phi}{\acontext}=\tvT\lor\modellabel{\phi}{\acontext}=\tvB$.

Because this is a core definition, we sum this up for the reader's clear reference:
$$
\acontext\ment\phi
\quad\liff\quad
{\ment\modellabel{\phi}{\acontext}}
\quad\liff\quad
\modellabel{\phi}{\acontext}\in\threeValid
\quad\liff\quad
\modellabel{\phi}{\acontext}=\tvT\lor\modellabel{\phi}{\acontext}=\tvB.
$$
We read this as \deffont{$\phi$ is valid in the context $\acontext$.}
\item
$n,p\mentval\phi$ means $\Forall{O\in\opensne}(n,p,O\mentval\phi)$.
\item\label{item.ment.n}
$n\mentval\phi$ means $\Forall{p\in\Pnt}\Forall{O\in\opensne}(n,p,O\mentval\phi)$.
\item\label{item.mentval}
$\mentval\phi$ means $\Forall{n\in\Time}\Forall{p\in\Pnt}\Forall{O\in\opensne}(n,p,O\mentval\phi)$.
\item\label{item.ment.phi}
$\ment\phi$ means $\Forall{\acontext\in\tf{Ctx}(\Sigma,\mathcal M)}(\acontext\ment\phi)$.
\end{enumerate*}
\end{defn}

\begin{figure}
$$
\begin{array}{l@{\quad}l@{\quad}l}
\mentval\Theta & \text{means} & \Forall{\theta\in\Theta}({\mentval\theta})
\figskip
\Theta\ment\phi & \text{means} &
\Forall{\avaluation\in\Valuation(\Sigma,\mathcal M)}({\mentval\Theta})\limp ({\mentval\phi}) 
\end{array}
$$
\caption{Validity of $\phi$ in a theory $\Theta$ (Definition~\ref{defn.theory.axiom})}
\label{fig.validity}
\end{figure}


Later on, when we start to build declarative axiomatisations of consensus algorithms, we will express these as \emph{theories}.
We include the relevant definition here, since it follows on naturally from Definition~\ref{defn.validity.judgement}: 
\begin{defn}[Validity in a theory]
\label{defn.theory.axiom}
Suppose that:
\begin{itemize*}
\item
$\Sigma$ is a signature and $\mathcal M=(\Pnt,\opens,\Val)$ is a model.
\item
$\phi$ is a closed predicate in $\Sigma$ and $\mathcal M$.
\item
$\Theta=\{\theta_1,\dots,\theta_n\}$ is a set of closed predicates.\footnote{It would be unusual for $\Theta$, which is intended to be a collection of axioms, to mention specific elements $v\in\Val$ or $\varkappa\in\THREE^{\tf{Ctx}(\Sigma,\mathcal M)}$ from a model (cf. Remark~\ref{rmrk.pedant.point}), but the maths does not seem to mind, so we do not bother to exclude this possibility.}
\end{itemize*}
Then:
\begin{enumerate*}
\item\label{item.theory}
A \deffont{theory} is a pair $(\Sigma,\Theta)$ of 
\begin{itemize*}
\item
a signature $\Sigma$, and 
\item
a set $\Theta$ of \deffont{axioms}, which are closed predicates over $\Sigma$ and $\mathcal M$.
\end{itemize*}
\item\label{item.axioms.valid}
Given a valuation $\avaluation\in\Valuation(\Sigma,\mathcal M)$ (Definition~\ref{defn.contexts}(\ref{item.contexts.valuation})), we define 
$$
\mentval\Theta
\quad\text{and}\quad
\Theta\ment\phi
$$ 
as per Figure~\ref{fig.validity}. 
When $\Theta\ment\phi$ holds, we say that $\phi$ is \deffont{valid in the theory $\Theta$} (in $\Sigma$ and $\mathcal M$). 
\end{enumerate*}
Examples of theories are in Figures~\ref{fig.ThyPaxOne} and~\ref{fig.ThySPax}.
\end{defn}


\begin{rmrk}
\label{rmrk.thy.quantifier.scope}
Note the scope of the quantifiers in Definition~\ref{defn.theory.axiom}:
\begin{enumerate*}
\item
The assertion $\Theta\ment\phi$ does \emph{not} mean 
$$
\Forall{\avaluation}(\mentval(\theta_1\tand\dots\tand\theta_n) \timp \phi),
$$ 
which is equivalent to 
$$
\Forall{\acontext}(\acontext\ment(\theta_1\tand\dots\tand\theta_n)\limp \acontext\mentval\phi) .
$$ 
Above, the short double arrow `$\timp$' means `if the LHS is valid then so is the RHS'; see Remark~\ref{rmrk.rich.implication}(\ref{item.rich.implication.timp}).
The long double arrow `$\limp$' is just logical implication. 
\item
The assertion $\Theta\ment\phi$ means 
$$
\Forall{\avaluation}(\mentval\theta_1\tand\dots\tand\theta_n) \limp (\mentval\phi) .
$$
In words: if $\Theta$ holds everywhere in the model under the valuation $\avaluation$, then $\phi$ holds everywhere in the model under the valuation $\avaluation$.
\end{enumerate*}
We spell this out further.
The assertion $\Theta\ment\phi$ means:
$$
\Forall{\avaluation}(\Forall{n,p,O}n,p,O\mentval\theta_1\tand\dots\tand\theta_n) \limp (\Forall{n,p,O}n,P,O\mentval\phi),
$$
where above $n\in\Time$, $p\in\Pnt$, and $O\in\opensne$.

This pattern of quantifiers is what makes $\Theta$ a theory, rather than just a set of assumptions: as a theory, $\Theta$ is a restriction on the valuation $\avaluation$.
\end{rmrk}


\begin{rmrk}
One final pedant point: 
in Definition~\ref{defn.theory.axiom}(\ref{item.theory}) we defined a theory to be a pair $(\Sigma,\Theta)$ of a signature and a set of closed predicates over a signature $\Sigma$ \emph{and} over a model $\mathcal M$.
This technical convenience lets us write axioms that directly reference specific values in the model (instead of introducing explicit constant symbols to our signature, and interpreting them as \dots values in the model).

We will exploit to assume a distinguished `undefined' value $\udfn$ in axiom \rulefont{PaxPropose?} in Figure~\ref{fig.ThyPaxOne}, as per Definition~\ref{defn.ThyPaxOne}(\ref{item.ThyPaxOne.axioms}).
\end{rmrk}

\subsection{Comments on notation: eliding parameters, and functional notation}

\begin{nttn}[Eliding parameters]
\label{nttn.validity.irrelevant.context}
It will be quite common that we make validity judgements where it is clear that part of the context is irrelevant.
For example the reader can check the denotational semantics in Figure~\ref{fig.3.phi.f} and the definition of validity in Definition~\ref{defn.validity.judgement}(\ref{item.ctx.ment.phi}), to see that:
\begin{enumerate*}
\item
$\acontext \ment \tvB$ is valid; we do not need to know the values in $\acontext$. 
\item
The validity of $n,p,O\mentval \tf P(v)$ depends on $\avaluation$, $n$, and $p$ --- but not on $O$.
\item
The validity of $n,p,O\mentval\QuorumBox\phi$ depends on $n$ and $\avaluation$ (and of course on $\phi$), but it does not depend on the values of $p$ and $O$.
\end{enumerate*}
In such cases, we may omit parameters in the context if they do not matter --- we may even omit the entire context if appropriate.
Thus, we may write $\ment\tvT$, $\ment \thea\teq\thea$, and $n\mentval\QuorumBox\phi$.

This notation remains within the scope of Definition~\ref{defn.validity.judgement} --- which universally quantifies over any omitted parameter --- so omitting parameters in this way is not an abuse of notation.
However, if the reader sees us omit a value from a context, there may be two reasons for it:
\begin{enumerate*}
\item
It may indicate a universal quantification over all possible values.
For example, this is clearly what we intend when we write $\mentval\theta$ in Figure~\ref{fig.validity} (top-right).
\item
It may indicate that we have elided irrelevant parameters.
For example, this is clearly what we intend in the example of $\ment\tvT$ above.
\end{enumerate*}
In practice, it will always be clear which is intended and meaning will always be clear.\footnote{Actually, meaning will be \emph{clearer}.  

In early drafts of this paper, we dutifully included all parameters, whether relevant or not.  
But this turned out to be just too confusing: the notational overhead of carrying around values that we did not care about was burdensome; but furthermore, it was actually easier to understand the sequents once we have dropped irrelevant parameters, because this helped us to keep track of which parameters \emph{did} matter.}
\end{nttn}

\begin{figure}
$$
\begin{array}{l@{\quad}l@{\quad}l@{\quad}l@{\quad}l}
\texiunique  a.\tf{write}(a)
&
\QuorumBox\texi a.\tf{send}(a)
&
\texi a.\everyoneAll\tf{decide}(a)
&
\tall a.\modTF \tf{propose}(a)
&
\mrup{a}{(\someone\tf{accept}(a))}{v}
\\[1ex]
& \text{may be written as}
\\[1ex]
\texiunique \tf{write} 
&
\QuorumBox\texi\tf{send} 
&
\texi\everyoneAll\tf{decide}
&
\tall\modTF\tf{propose} 
&
\mru{(\someone\tf{accept})}{v}
\end{array}
$$
\caption{Abbreviations inspired by higher-order logic (Notation~\ref{nttn.hos})}
\label{fig.hos}
\end{figure}


\begin{nttn}[Functional notation]
\label{nttn.hos}
We will frequently apply quantifiers (like $\texi$, $\texiunique$, or $\tall$) to unary predicate symbols, possibly with modalities.
For example:
$$
\tall a.\QuorumBox\tf{decide}(a)
\quad\text{is}\quad\text{$\tall$ quantifying $a$ in\ $(\QuorumBox\tf{decide})$\ applied to $a$.}
$$ 
For clarity, we may use a \deffont{functional notation} inspired by higher-order logic (HOL) to elide the variable, as illustrated in Figure~\ref{fig.hos}.
Meaning will always be clear.

See also Remark~\ref{rmrk.functional.notation.in.paxos} for examples of this notation applied in an axiomatisation.
\end{nttn}






\section{Properties of predicates}
\label{sect.predicates}

This Section collects elementary technical results which we will need later.
Some of these results, and their proofs, are more technical than others, but they are not necessarily surprising if we look at what they mean in the model.
Mostly the statements are simple and it is obvious by common sense that they should work (simple is good, because it means greater productivity and fewer errors).

A typical example is Lemma~\ref{lemm.tomorrow.forever.recent}(\ref{item.tomorrow.forever.recent.1}): once we look at what it actually says about the model, this is that ``if $\phi$ is true today, then in the future it will have been true in the past''.
For saying this a mathematician earns a living? --- we might think.

What is less evident is \emph{which} elementary properties will be important later, nor is it obvious what definitions and structures make these properties all fit together nicely.
So the reader should note that the simplicity of this Section is a design feature, which was only attained in retrospect and after some work.

Broadly speaking, we start this Section with the most elementary results, and build complexity from there.
The reader is welcome to read this Section now, or to skip it and refer back to its results as they are used in later proofs.


\subsection{Equivalence of predicates, and some (de Morgan) equivalences}
\label{subsect.equivalence.of.predicates}

For the rest of this Section we make a fixed but arbitrary choice of signature $\Sigma$ (Definition~\ref{defn.signature}(\ref{item.predicate.syntax.signature})) and model $\mathcal M$ (Definition~\ref{defn.model}).
All predicates considered will be over $\Sigma$ and $\mathcal M$, unless otherwise stated.

We generalise equivalence $\equiv$ (Definition~\ref{defn.equivalent}) from truth-values, to predicates viewed as functions from contexts to truth-values:
\begin{defn}[Equivalence of predicates]
\label{defn.phi.equivalent}
Suppose $\phi$ and $\phi'$ are closed predicates.

Call $\phi$ and $\phi'$ \deffont{extensionally equivalent} when 
$$
\model{\phi}=\model{\phi'} .
$$
Unpacking Definition~\ref{defn.denotation}, this means 
$$
\Forall{\acontext\in\tf{Ctx}(\Sigma,\mathcal M)}(\modellabel{\phi}{\acontext}=\modellabel{\phi'}{\acontext}) .
$$
For brevity, we may write `extensionally equivalent' just as `equivalent' henceforth.
\end{defn}



\begin{lemm}
\label{lemm.equiv}
Suppose $\phi$ and $\phi'$ are closed predicates and $\acontext$ is a context.
Then: 
\begin{enumerate*}
\item\label{item.equiv.1}
$\acontext\ment\phi\equiv\phi'$ if and only if $\modellabel{\phi}{\acontext}=\modellabel{\phi'}{\acontext}$.
\item\label{item.equiv.2}
$\ment\phi\equiv\phi'$ if and only if $\model{\phi}=\model{\phi'}$ (i.e. $\phi$ and $\phi'$ are equivalent as per Definition~\ref{defn.phi.equivalent}).
\end{enumerate*}
\end{lemm}
\begin{proof}
We consider each part in turn:
\begin{enumerate}
\item
By Definition~\ref{defn.validity.judgement}(\ref{item.ctx.ment.phi}) $\acontext\ment\phi\equiv\phi'$ means $\modellabel{\phi\equiv\phi'}{\acontext}\in\threeValid$, and by the clause for $\equiv$ in Figure~\ref{fig.3.derived} this is if and only if $\modellabel{\phi}{\acontext}=\modellabel{\phi'}{\acontext}$.
\item
By Definition~\ref{defn.validity.judgement}(\ref{item.ment.phi}), $\ment\phi\equiv\phi'$ means precisely $\Forall{\acontext}\acontext\ment\phi\equiv\phi'$.
We use part~\ref{item.equiv.1} of this result.
\qedhere\end{enumerate}
\end{proof}


We will see many equivalences in what follows.
We start by using Definition~\ref{defn.phi.equivalent} and Lemma~\ref{lemm.equiv} to express some standard equivalences in logic, having to do with implication being expressible using negation and disjunction, and with \emph{de Morgan} dualities between modalities and connectives:
\begin{lemm}
\label{lemm.minimal.syntax}
For each of the equivalences $\f{LHS}\equiv\f{RHS}$ listed in Figure~\ref{fig.easy.equivalences}, we have that $\f{LHS}$ is extensionally equal to $\f{RHS}$ in the sense of Definition~\ref{defn.phi.equivalent}.
\end{lemm}
\begin{proof}
By routine calculations using Figures~\ref{fig.3.phi.f} and~\ref{fig.3.derived} and facts of the lattice $\THREE=(\tvF<\tvB<\tvT)$.
\end{proof}


\begin{figure}
$$
\begin{array}{r@{\ }l@{\quad}r@{\ }l@{\quad}r@{\ }l@{\quad}r@{\ }l}
\phi\tand\phi' \equiv& \tneg(\tneg\phi\tor\tneg\phi')
&
\phi\tor\phi' \equiv& \tneg(\tneg\phi\tand\tneg\phi')
&
\tneg\phi\equiv& \phi\timpc\tvF \equiv\phi\tnotor\tvF %
\\
\phi\timpc\phi' \equiv& \phi\tnotor\modT\phi'
&
\phi\tnotor\phi' \equiv& \tneg\phi\tor\phi'
&
\phi\timpc\phi' \equiv& \phi\tnotor\modT\phi'
\\
\everyone\phi\equiv&\tneg\someone\tneg\phi
&
\someone\phi\equiv&\tneg\everyone\tneg\phi
\\
\Coquorum\phi\equiv&\tneg\Quorum\tneg\phi
&
\Quorum\phi\equiv&\tneg\Coquorum\tneg\phi
\\
\forever\phi\equiv&\tneg\sometime\tneg\phi
&
\sometime\phi\equiv&\tneg\forever\tneg\phi
\\
\tall a.\phi\equiv&\tneg\texi a.\tneg\phi
&
\texi a.\phi\equiv&\tneg\tall a.\tneg\phi
&
\modTF\phi\equiv&\modT\phi\tor\modT\tneg\phi
\\
\modB\phi\equiv&\tneg\modTF\phi
&
\modT\phi\equiv&\tneg\modFB\phi
&
\modF\phi\equiv&\tneg\modTB\phi
\end{array}
$$
\caption{Some easy equivalences}
\label{fig.easy.equivalences}
\end{figure}


\subsection{Two commutation lemmas}
\label{subsect.two.commutation.results}

We note some simple commutation lemmas.
These results are all easy facts of our model, but it only takes a few lines to spell them out:

\begin{lemm}[Commuting connectives]
\label{lemm.commuting.connectives}
Suppose $\phi$ and $\phi'$ are closed predicates and suppose
$$
\bullet\in\{\tor,\tand\} 
\ \ \text{and}\ \ 
\tf M\in\{\modT,\modTB\} . 
$$
Then
$$
\ment\tf M(\phi\bullet\phi')\equiv (\tf M\phi\bullet\tf M\phi') .
$$
\end{lemm}
\begin{proof}
Routine from Lemma~\ref{lemm.equiv}(\ref{item.equiv.2}) and Figures~\ref{fig.3.phi.f} and~\ref{fig.3.derived}, and facts of $\THREE$ as a lattice.
\end{proof}


\begin{lemm}[Commuting modalities]
\label{lemm.someone.commutation}
Suppose $\phi$ is a closed predicate and suppose 
$$
\mathsf Q\in\{\texi,\tall,\someone,\everyone,\Quorum,\Coquorum,\recent,\urecent,\yesterday,\tomorrow,\forever,\sometime,\infinitely,\final\}.
$$
Then 
$$
\ment\mathsf Q\modT\phi \equiv \modT\mathsf Q\phi
\quad\text{and}\quad
\ment\mathsf Q\modTB\phi \equiv \modTB\mathsf Q\phi.
$$
\end{lemm}
\begin{proof}
Routine from Lemma~\ref{lemm.equiv}(\ref{item.equiv.2}) and Figures~\ref{fig.3.phi.f} and~\ref{fig.3.derived}, and facts of $\THREE$ as a lattice.
\end{proof}


\subsection{Lattice inequalities}

Corollary~\ref{corr.recent.everyone.commute} is in the same spirit as Lemmas~\ref{lemm.commuting.connectives} and~\ref{lemm.someone.commutation}. 
In practice we usually just use it a direct fact of the model, but we spell it out anyway.
It relies on this simple fact about lattices: 
\begin{lemm} 
\label{lemm.quantifier.swap}
Suppose $(\mathcal L,\tand,\tor)$ is a complete lattice and $X$ and $Y$ are sets and $f:(X\times Y)\to\mathcal L$ is any function. 
Then:
\begin{enumerate*}
\item\label{item.quantifier.swap.1}
$\bigvee_{x}\bigwedge_{y}f(x,y)\leq \bigwedge_{y}\bigvee_{x}f(x,y)$, 
\item\label{item.quantifier.swap.1b}
It is not in general the case that 
$\bigwedge_{y}\bigvee_{x}f(x,y) \leq \bigvee_{x}\bigwedge_{y}f(x,y)$.
\item\label{item.quantifier.swap.2}
$\bigvee_{x}\bigvee_{y}f(x,y) = \bigvee_{y}\bigvee_{x}f(x,y)$.
\end{enumerate*}
\end{lemm}
\begin{proof}
By easy lattice calculations. 
\end{proof}


\begin{corr}
\label{corr.recent.everyone.commute}
Suppose $\mathsf Q\phi$ is a closed predicate, where 
$$
\mathsf Q\in\{\texi a,\tall a \mid a{\in}\tf{VarSymb}\}\cup\{\someone,\everyone,\Quorum,\Coquorum\},
$$
and suppose $\acontext$ is a context.
Then:
\begin{enumerate*}
\item\label{item.recent.everyone.commute.1}
$\modellabel{\recent\mathsf Q\phi}{\acontext}\leq \modellabel{\mathsf Q\recent\phi}{\acontext}$. 
\item
$\modellabel{\mathsf Q\recent\phi}{\acontext}\leq \modellabel{\recent\mathsf Q\phi}{\acontext}$ for 
$\mathsf Q\in\{\texi a\mid a{\in}\tf{VarSymb}\}\cup\{\someone,\Quorum\}$.
\item
It need not be the case that $\modellabel{\mathsf Q\recent\phi}{\acontext}\leq \modellabel{\recent\mathsf Q}{\acontext}$ for 
$\mathsf Q\in\{\tall a\mid a{\in}\tf{VarSymb}\}\cup\{\everyone,\Coquorum\}$.
\end{enumerate*}
\end{corr}
\begin{proof}
Suppose $n\in\Time$ and $p\in\Pnt$ and $O\in\opensne$ and $\avaluation$ is a valuation.
We freely use \strongmodusponens.
\begin{enumerate}
\item
We consider $\mathsf Q=\Coquorum$.
We use the clause for $\recent$ in Figure~\ref{fig.3.derived}, the clause for $\Coquorum$ in Figure~\ref{fig.3.phi.f}, and then Lemma~\ref{lemm.quantifier.swap}(\ref{item.quantifier.swap.1}): 
$$
\bigvee_{0{\leq}n'{<}n}\bigwedge_{O'{\in}\opensne}\modellabel{\phi}{n',p,O',\avaluation}
\stackrel{L\ref{lemm.quantifier.swap}(\ref{item.quantifier.swap.1})}{\leq}
\bigwedge_{O'{\in}\opensne}\bigvee_{0{\leq}n'{<}n}\modellabel{\phi}{n',p,O',\avaluation} .
$$
Now we consider $\mathsf Q=\tall a$:
$$
\bigvee_{0{\leq}n'{<}n}\bigwedge_{v{\in}\Val}\modellabel{\phi[a\ssm v]}{n',p,O,\avaluation}
\stackrel{L\ref{lemm.quantifier.swap}(\ref{item.quantifier.swap.1})}{\leq}
\bigwedge_{v{\in}\Val}\bigvee_{0{\leq}n'{<}n}\modellabel{\phi[a\ssm v]}{n',p,O,\avaluation} .
$$
The other cases are precisely similar.
\item
We consider just $\mathsf Q=\someone$; the other cases are precisely similar.
We just use Lemma~\ref{lemm.quantifier.swap}(\ref{item.quantifier.swap.2}): 
$$
\bigvee_{p'{\in}\Pnt}\bigvee_{0{\leq}n'{<}n}\modellabel{\phi}{n',p',O,\avaluation}
\stackrel{L\ref{lemm.quantifier.swap}(\ref{item.quantifier.swap.2})}{=}
\bigvee_{0{\leq}n'{<}n}\bigvee_{p'{\in}\Pnt}\modellabel{\phi}{n',p',O,\avaluation} .
$$
\item
We consider just $\mathsf Q=\tall$; the other cases are precisely similar.
We note by Lemma~\ref{lemm.quantifier.swap}(\ref{item.quantifier.swap.1b}) that 
$$
\bigwedge_{v{\in}\Val}\bigvee_{0{\leq}n'_v{<}n}\modellabel{\phi[a\ssm v]}{n'_v,p,O,\avaluation}
\leq
\bigvee_{0{\leq}n'{<}n}\bigwedge_{v\in\Val}\modellabel{\phi[a\ssm v]}{n',p,O,\avaluation} 
\ \ \ \text{does not hold in general}.
$$
\qedhere\end{enumerate}
\end{proof}





\subsection{The compound modalities $\someoneAll$, $\everyoneAll$, $\QuorumBox$, and $\CoquorumDiamond$}
\label{subsect.compound.modalities}

\begin{rmrk}
\label{rmrk.compound.modalities}
We will make much use of the modalities $\someoneAll$, $\everyoneAll$, $\QuorumBox$, and $\CoquorumDiamond$.

We call these four combinations \deffont{compound modalities}, because they are built out of $\Quorum$ and $\Coquorum$ (which quantify over open sets) and $\someone$ and $\everyone$ (which quantify over points in an open set). 
In this Subsection we study some of their properties.
\end{rmrk}



Lemma~\ref{lemm.dense.char} just unpacks definitions:
\begin{lemm}
\label{lemm.dense.char}
Suppose $\phi$ is a closed predicate and $n\in\Time$ and $\avaluation$ is a valuation.
Then: 
\begin{enumerate*}
\item\label{item.dense.char.someoneAll}
$n\mentval\someoneAll\phi$ if and only if $\Exists{O\in\opensne}O\between\{p\in\Pnt \mid n,p,O\mentval\phi\}$.
\item\label{item.dense.char.everyoneAll}
$n\mentval\everyoneAll\phi$ if and only if $\Forall{O\in\opensne}O\subseteq\{p\in\Pnt \mid n,p,O\mentval\phi\}$.
\item\label{item.dense.char.QuorumBox}
$n\mentval\QuorumBox\phi$ if and only if $\Exists{O\in\opensne}O\subseteq\{p\in\Pnt \mid n,p,O\mentval\phi\}$.
\item\label{item.dense.char.CoquorumDiamond}
$n\mentval\CoquorumDiamond\phi$ if and only if $\Forall{O\in\opensne}O\between\{p\in\Pnt \mid n,p,O\mentval\phi\}$.
\end{enumerate*} 
\end{lemm}
\begin{proof}
We consider each part in turn:
\begin{enumerate}
\item
We reason as follows: 
$$
\begin{array}{r@{\ }l@{\quad}l}
n\mentval\someoneAll\phi
\liff&
\Exists{O{\in}\opensne}\Exists{p{\in}O}n,p,O\mentval\phi
&\text{Figure~\ref{fig.3.phi.f}}
\\
\liff&
\Exists{O{\in}\opensne}O\between \{p\in\Pnt \mid n,p,O\mentval\phi\}
&\text{Fact}
\end{array}
$$
\item
We reason as follows: 
$$
\begin{array}{r@{\ }l@{\quad}l}
n\mentval\everyoneAll\phi
\liff&
\Forall{O{\in}\opensne}\Forall{p{\in}O}n,p,O\mentval\phi
&\text{Figure~\ref{fig.3.phi.f}}
\\
\liff&
\Forall{O{\in}\opensne}O\subseteq\{p\in\Pnt \mid n,p,O\mentval\phi\}
&\text{Fact}
\end{array}
$$
\item
We reason as follows:
$$
\begin{array}{r@{\ }l@{\quad}l}
n\mentval\QuorumBox\phi
\liff&
\Exists{O{\in}\opensne}\Forall{p{\in}O}n,p,O\mentval\phi
&\text{Figure~\ref{fig.3.phi.f}}
\\
\liff&
\Exists{O{\in}\opensne}O\subseteq\{p\in\Pnt \mid n,p,O\mentval\phi\}
&\text{Fact}
\end{array}
$$
\item
We reason as follows:
$$
\begin{array}{r@{\ }l@{\quad}l}
n\mentval\CoquorumDiamond\phi
\liff&
\Forall{O{\in}\opensne}\Exists{p{\in}O}n,p,O\mentval\phi
&\text{Figure~\ref{fig.3.phi.f}}
\\
\liff&
\Forall{O{\in}\opensne}O\between\{p\in\Pnt \mid n,p,O\mentval\phi\}
&\text{Fact}
\end{array}
$$
\qedhere\end{enumerate}
\end{proof}



Proposition~\ref{prop.someone.implies.someoneAll} packages up some simple corollaries of Lemma~\ref{lemm.dense.char}, in forms that will be useful for our later proofs: 
\begin{prop}
\label{prop.someone.implies.someoneAll}
Suppose $\phi$ is a closed predicate and $\acontext$ is a context.
Then:
\begin{enumerate*}
\item\label{item.someone.implies.someoneAll.4}
$\acontext\ment\QuorumBox\phi\tand\CoquorumDiamond\psi\limp\acontext\ment\someoneAll(\phi\tand\psi)$.
\item\label{item.someone.implies.someoneAll.4b}
$\acontext\ment\QuorumBox\phi\tand\everyoneAll\psi\limp\acontext\ment\QuorumBox(\phi\tand\psi)$.
\item\label{item.someone.implies.someoneAll.5}
$\acontext\ment\modT(\QuorumBox\phi\tand\CoquorumDiamond\psi)\limp\acontext\ment\modT\someoneAll(\phi\tand\psi)$,\ and
\\
$\acontext\ment(\QuorumBox\phi\tand\CoquorumDiamond\psi)\timpc \someoneAll(\phi\tand\psi)$.
\item\label{item.someone.implies.someoneAll.2}
$\acontext\ment\modT\QuorumBox\phi\limp\acontext\ment\modT\someoneAll\phi$.
\item\label{item.someone.implies.someoneAll.2b}
$\acontext\ment\modT\CoquorumDiamond\phi\limp\acontext\ment\modT\someoneAll\phi$.
\end{enumerate*}
\end{prop}
\begin{proof}
We consider each part in turn:
\begin{enumerate}
\item
Routine from Lemma~\ref{lemm.dense.char}(\ref{item.dense.char.QuorumBox}\&\ref{item.dense.char.CoquorumDiamond}).
\item
Routine from Lemma~\ref{lemm.dense.char}(\ref{item.dense.char.QuorumBox}\&\ref{item.dense.char.everyoneAll}).
\item
From part~\ref{item.someone.implies.someoneAll.4} of this result taking $\modT\phi$ and $\modT\psi$ and the commutation Lemmas~\ref{lemm.someone.commutation} and~\ref{lemm.commuting.connectives}.
Then we use \strongmodusponens.
\item
From part~\ref{item.someone.implies.someoneAll.5} of this result, taking $\psi=\tvT$.
\item
From part~\ref{item.someone.implies.someoneAll.5} of this result, taking $\phi=\tvT$.
\qedhere\end{enumerate}
\end{proof}



 







\subsection{Some comments on time}

\begin{rmrk}
This Subsection notes Lemmas~\ref{lemm.tomorrow.forever.recent} and~\ref{lemm.persist.char}.
They both have to do with time, and underneath all their complicated-looking symbols they are both straightforward:
\begin{itemize*}
\item
Lemma~\ref{lemm.tomorrow.forever.recent} reflects a fact that if something holds now, then in the future it will have held.
\item
Lemma~\ref{lemm.persist.char} reflects the fact that something repeats forever if and only if it happens infinitely often.
\end{itemize*}
\end{rmrk}


\begin{lemm}
\label{lemm.tomorrow.forever.recent}
Suppose $\acontext=(n,p,O,\avaluation)$ is a context and $\phi$ is a closed predicate.
Then:
\begin{enumerate*}
\item\label{item.tomorrow.forever.recent.1}
$n,p,O\mentval\modT\phi$ implies $\Forall{n'{\gneq}n}\,n',p,O\mentval\modT\recent\phi$.
\item\label{item.rur.r}
$n,p,O\mentval\modT\urecent\phi$ implies $\Forall{n'{\geq} n}\,n',p,O\mentval\modT\urecent\phi$.
\item
As corollaries, 
$$
\acontext\ment\phi\timpc \forever\recent\phi
\quad\text{and}\quad
\acontext\ment\urecent\phi\timpc\forever\urecent\phi.
$$
\end{enumerate*}
\end{lemm}
\begin{proof}
This is just from the clauses for $\recent$ and $\urecent$ in Figure~\ref{fig.3.derived}. 
The corollary follows using \strongmodusponens.
\end{proof}



$\infinitely=\forever\sometime$ has a characterisation as `infinitely often', with the convenient feature that its denotation is independent of the stage at which it is evaluated (called $n'$ in the statement of the Lemma below):
\begin{lemm}
\label{lemm.persist.char}
Suppose $\phi$ is a closed predicate and $p\in\Pnt$ and $O\in\opensne$ and $\acontext$ is a context.
Suppose further that $n'\in\Time$ and $\f{tvs}\in\{\threeTrue,\threeValid\}$.
Then we have
$$
\modellabel{\infinitely\phi}{n',p,O,\avaluation}\in\f{tvs}
\quad\liff\quad
\{n\in\Time \mid \modellabel{\phi}{n,p,O,\avaluation}\in\f{tvs}\}\text{ is infinite} . 
$$
\end{lemm}
\begin{proof}
By routine arguments from the clause for $\infinitely$ in Figure~\ref{fig.3.derived}.
\end{proof}






\subsection{Unique existence and affine existence in three-valued logic}
\label{subsect.unique.affine.existence}

\subsubsection{The definition, and main result}

\begin{rmrk}
Distributed algorithms often require an implementation to make (at most) one choice from (zero or more) possibilities.
In our axioms, we model this using unique existence and affine existence:
\begin{itemize*}
\item
\emph{Unique existence} means `there exists precisely one'; we write this $\exists!$, as standard.
\item
\emph{Affine existence} means `there exist zero or one'; we write this $\exists_{01}$.\footnote{We are not aware of a standard symbol for affine existence.}
\end{itemize*}
See for example \rulefont{PaxPropose01} and \rulefont{PaxWrite01} in Figure~\ref{fig.logical.paxos} (a discussion of which is in Remark~\ref{rmrk.paxos.axioms.discussion}(\ref{item.01.discussion})).

So we need to think about what $\exists!$ and $\exists_{01}$ even mean in a three-valued setting.
This turns out to be not completely straightforward, because we have a third truth-value that is neither true nor false.

For concreteness, suppose $\tf P$ is a unary predicate and $\acontext$ is a context, and suppose that 
$\modellabel{\phi[a\ssm v]}{\acontext}=\tvB$ for all $v\in\Val$.
Then what truth-values should $\modellabel{\texiaffine a.\phi}{\acontext}$ and $\modellabel{\texiunique a.\phi}{\acontext}$ have?
 
Now suppose $v,v',v''\in\Val$, and suppose that $\modellabel{\phi[a\ssm v''']}{\acontext}=\tvF$ for all $v'''\in\Val\setminus\{v,v',v''\}$ and
$$
\modellabel{\phi[a\ssm v]}{\acontext}=\tvT
\quad\land\quad
\modellabel{\phi[a\ssm v']}{\acontext}=\tvB
\quad\land\quad
\modellabel{\phi[a\ssm v'']}{\acontext}=\tvB .
$$
\noindent Then what truth-values should $\modellabel{\texiaffine a.\phi}{\acontext}$ and $\modellabel{\texiunique a.\phi}{\acontext}$ have?

Now suppose that $\modellabel{\phi[a\ssm v''']}{\acontext}=\tvF$ for all $v'''\in\Val\setminus\{v\}$ and
$$
\modellabel{\phi[a\ssm v]}{\acontext}=\tvB .
$$
\noindent Then what truth-values should $\modellabel{\texiaffine a.\phi}{\acontext}$ and $\modellabel{\texiunique a.\phi}{\acontext}$ have?

Our answers to all questions above will be: $\tvB$ for both.
This is not automatic: it is a real design decision that we have to make.
\end{rmrk}

\begin{defn}
\label{defn.unique.exist}
Suppose $\phi$ is a predicate.
Define \deffont{affine existence} $\texiaffine a.\phi$ and \deffont{unique existence} $\texiunique  a.\phi$ by:
$$
\begin{array}{r@{\ }l}
\texiaffine a.\phi =& \tall a,a'.(\phi\tand\phi[a\ssm a']) \tnotor a\teq a'
\\
\texiunique  a.\phi =& (\texi a.\phi) \tand (\texiaffine a.\phi)
\end{array}
$$
\end{defn}


\begin{rmrk}
\label{rmrk.unpack.texiunique}
Suppose $\texi a.\phi$ is a closed predicate (i.e. $a\in\tf{VarSymb}$ and $\phi$ is a predicate with free variable symbols at most $\{a\}$), and suppose $\acontext$ is a context.
We spell unpack the denotations of $\texiunique$ and $\texiaffine$ from Definition~\ref{defn.unique.exist}: 
\begin{enumerate*}
\item
If $\exists v,v'\in\Val.v\neq v'\land\modellabel{\phi[a\ssm v]}{\acontext}=\modellabel{\phi[a\ssm v']}{\acontext}=\tvT$, then 
$$
\modellabel{\texiunique a.\phi}{\acontext}=\tvF=\modellabel{\texiaffine a.\phi}{\acontext}.
$$
\item
If $\exists!v\in\Val.\modellabel{\phi[a\ssm v]}{\acontext}=\tvT$ and $\Forall{v\in\Val}\modellabel{\phi[a\ssm v]}{\acontext}\in\{\tvT,\tvF\}$, then 
$$
\modellabel{\texiunique a.\phi}{\acontext}=\tvT=\modellabel{\texiaffine a.\phi}{\acontext}.
$$
\item
If $\exists!v\in\Val.\modellabel{\phi[a\ssm v]}{\acontext}=\tvT$ and $\Exists{v\in\Val}\modellabel{\phi[a\ssm v]}{\acontext}=\tvB$, then 
$$
\modellabel{\texiunique a.\phi}{\acontext}=\tvB=\modellabel{\texiaffine a.\phi}{\acontext}.
$$
\item\label{item.explain.existsunique.4}
If $\neg\Exists{v\in\Val}\modellabel{\phi[a\ssm v]}{\acontext}=\tvT$ and $\exists! v'\in\Val.\modellabel{\phi[a\ssm v]}{\acontext}=\tvB$, then 
$$
\modellabel{\texiunique a.\phi}{\acontext}=\tvB
= \modellabel{\texiaffine a.\phi}{\acontext} . %
$$
(We discuss this in Remark~\ref{rmrk.discuss.clause.4} below.)
\item\label{item.two.b}
If $\neg\Exists{v\in\Val}\modellabel{\phi[a\ssm v]}{\acontext}=\tvT$ and $\exists v,v'\in\Val.v\neq v'\land \modellabel{\phi[a\ssm v]}{\acontext}=\modellabel{\phi[a\ssm v']}{\acontext}=\tvB$, then 
$$
\modellabel{\texiunique a.\phi}{\acontext}
=\tvB=
\modellabel{\texiaffine a.\phi}{\acontext} .
$$
\item
If $\Forall{v\in\Val}\modellabel{\phi[a\ssm v]}{\acontext}=\tvF$, then
$$
\modellabel{\texiunique a.\phi}{\acontext}=\tvF
\quad\text{and}\quad
\modellabel{\texiaffine a.\phi}{\acontext}=\tvT.
$$
\end{enumerate*}
\end{rmrk}

\begin{rmrk}
\label{rmrk.discuss.texiunique}
We make two elementary observations on Remark~\ref{rmrk.unpack.texiunique}:
\begin{enumerate} 
\item
The equality $\teq$ in our logic is two-valued: $\model{v\teq v'}$ returns $\tvT$ or $\tvF$, and never $\tvB$.\footnote{In the terminology of Definition~\ref{defn.correct}(\ref{item.correct.phi}) below, equality predicates are always \emph{correct}.}
\item
Two notions of `unique existence' are in play in this discussion:
\begin{itemize*}
\item
`$\exists!v\in\Val$' refers to the actual unique existence quantifier in the real world (whatever that is).
This computes actual truth-values `true' and `false' and returns true when there actually exists a unique value.
\item
`$\texiunique a$' is part of the (derived) syntax of \QLogic, and is defined in Definition~\ref{defn.unique.exist}.
Its denotation computes elements of $\THREE=\{\tvF,\tvB,\tvT\}$.  
As is standard (though confusingly) we \emph{also} call elements of $\THREE$ truth-values; but they are truth-values of \QLogic, rather than actual truth-values in the real world (whatever that is).
\end{itemize*}
So: $\exists!$ and $\texiunique$ are both unique existence operators, but they live at different levels and they are not the same thing.\footnote{This kind of thing is far from unusual and appears also outside of papers with fancy logics.  For example in programming, a programming language may have multiple versions of the number $1$ (8-bit, floating point, 64-bit, and infinite precision) all of which are called `one', but all of which are different from one other, and also different from the actual $1\in\mathbb N$ (whatever that is).}
\end{enumerate}
\end{rmrk}


We will use unique/affine existence in our proofs via Lemma~\ref{lemm.unique.affine.existence}: 
\begin{lemm}
\label{lemm.unique.affine.existence}
Suppose $\acontext$ is a context and $\texi a.\phi$ is a closed predicate and $v,v'\in\Val$ are values.
Then:
\begin{enumerate*}
\item\label{item.unique.affine.existence.01}
$\acontext\ment\texiaffine a.\phi$ if and only if 
$\exists_{01} v\in\Val.\modellabel{\phi[a\ssm v]}{\acontext}=\tvT$.
\item\label{item.unique.affine.existence.11}
$\acontext\ment\texiunique a.\phi$ if and only if 
$\exists v\in\Val.\modellabel{\phi[a\ssm v]}{\acontext}\in\threeValid \land \exists_{01} v\in\Val.\modellabel{\phi[a\ssm v]}{\acontext}=\tvT$.
\item\label{item.unique.affine.existence.01.implies}
If $\acontext\ment\texiaffine a.\phi$ and $\acontext\ment\modT\phi[a\ssm v]$ and $\acontext\ment\modT\phi[a\ssm v']$, then $v=v'$.
\item\label{item.unique.affine.existence.11.implies}
If $\acontext\ment\texiunique a.\phi$ and $\acontext\ment\modT\phi[a\ssm v]$ and $\acontext\ment\modT\phi[a\ssm v']$, then $v=v'$.
\end{enumerate*}
\end{lemm}
\begin{proof}
We unpack Definition~\ref{defn.validity.judgement}(\ref{item.ctx.ment.phi}) %
and check the possibilities in Remark~\ref{rmrk.unpack.texiunique}.
\end{proof}


\subsubsection{Comments on an edge case}

Lemma~\ref{lemm.texiunique.all.b} notes a (perhaps surprising, at first glance) fact about our denotation, that 
if a unary predicate symbol $\tf P$ takes truth-value $\tvB$ when applied to all values (as would model a participant that has entirely crashed), then $\texiunique\tf P$ and $\texiaffine\tf P$ are valid --- even though $\tf P(v)$ is valid on every value.
This reflects the common side-condition in the literature `if the participant is correct, then \emph{stuff}', which is trivially true if the participant is \emph{not} correct: 
\begin{lemm}
\label{lemm.texiunique.all.b}
Suppose $\acontext$ is a context and $\tf P$ is a unary predicate symbol:\footnote{There is nothing mathematically special about using a \emph{unary} predicate symbol here, but in practice this is the case that we will use.}
Then
$$
\acontext\ment \tall\modB\tf P \limp (\acontext\ment \texiunique \tf P \land \acontext\ment \texiaffine \tf P).  
$$
The above statement uses Notation~\ref{nttn.hos}.
Unpacking this, we obtain:
$$
\acontext\ment \tall a.\modB\tf P(a) \limp (\acontext\ment \texiunique a.\tf P(a)\land \acontext\ment \texiaffine a.\tf P(a)) . 
$$  
\end{lemm}
\begin{proof}
Suppose $\acontext\ment \tall\modB\tf P$.
Unpacking Definition~\ref{defn.validity.judgement}(\ref{item.ctx.ment.phi}) this means $\modellabel{\tall\modB\tf P}{\acontext}\in\threeValid$, and by Figure~\ref{fig.3.derived} it follows that $\Forall{v\in\Val}\modellabel{\tf P(v)}{\acontext}=\tvB$.

By Remark~\ref{rmrk.unpack.texiunique}(\ref{item.two.b}) it follows that $\modellabel{\texiunique\tf P}{\acontext}=\tvB$.
The truth-value $\tvB$ is not true, but it is \emph{valid} by Definition~\ref{defn.validity.judgement}(\ref{item.ctx.ment.phi}), and thus $\acontext\ment\texiaffine\tf P$ and $\acontext\ment\texiunique\tf P$ as required. 
\end{proof}



\begin{rmrk}
\label{rmrk.discuss.clause.4}
Arguably, we would expect in clause~\ref{item.explain.existsunique.4} above 
(where $\neg\Exists{v\in\Val}\modellabel{\phi[a\ssm v]}{\acontext}=\tvT$ and $\exists! v'\in\Val.\modellabel{\phi[a\ssm v]}{\acontext}=\tvB$)
that $\modellabel{\texiaffine a.\phi}{\acontext}$ should equal $\tvT$, not $\tvB$. 
This is not hard to arrange, e.g. as
$$
\texiaffine a.\phi \tor \texiaffine a.\modT\phi .
$$
Either choice works with our later proofs; we never encounter an edge case where it matters.
\end{rmrk}



\subsection{Pointwise predicates (a well-behavedness property)}

For this subsection we fix a signature $\Sigma$ and a model $\mathcal M=(\Pnt,\opens,\Val)$.

\subsubsection{The definition, and basic properties}

\begin{rmrk}
\emph{Being pointwise} is a well-behavedness property on predicates, that their denotation in a context $\acontext=(n,p,O,\avaluation)$ %
does not depend on the open set parameter $O$ of the context.

Many important syntactic classes of predicates are naturally pointwise, including any predicate of the form $\tf P(v_1,\dots,v_n)$ (Lemma~\ref{lemm.atomic.O}).

Being pointwise is desirable for us because it simplifies denotation to holding `at a point';\footnote{Actually: at a stage, point, and valuation, but intuitively we tend to view the stage and valuation as fixed, and the point $p$ and open set $O$ as varying.} and also because several nice properties depend on this (see for example Proposition~\ref{prop.pointwise.dense.char}); and finally also because there is an easy syntactic condition (given in Lemma~\ref{lemm.atomic.O}(\ref{item.atomic.pointwise.2})) that guarantees that a predicate is pointwise.
\end{rmrk}
 

\begin{defn}[Pointwise predicates]
\label{defn.pointwise.meaning}
Suppose $\phi$ is a closed predicate.
\begin{enumerate*}
\item
If $n\in\Time$ and $\avaluation\in\Valuation(\Sigma,\Val)$ then call $\phi$ \deffont{pointwise in $(n,\avaluation)$} when
$$
\Forall{p\in\Pnt}\Forall{O,O'{\in}\opensne}\modellabel{\phi}{n,p,O,\avaluation}=\modellabel{\phi}{n,p,O',\avaluation} .
$$ 
In words, $\phi$ is pointwise in $(n,\avaluation)$ when at every point $p\in\Pnt$, the open set parameter $O\in\opensne$ does not affect the value of the denotation $\modellabel{\phi}{n,p,O,\avaluation}$. 
\item
Call $\phi$ \deffont{absolutely pointwise} when it is pointwise for every $n\in\Time$ and $\avaluation\in\Valuation(\Sigma,\Val)$. 
\end{enumerate*}
\end{defn}

\begin{defn}
\label{defn.pointwise.pred.nttn}
Suppose $\phi$ is closed, and $n\in\Time$ and $p\in\Pnt$ and $\avaluation$ is a valuation.
Then define $\modellabel{\phi}{n,p,\avaluation}$ by
$$
\modellabel{\phi}{n,p,\avaluation}=\bigwedge_{O\in\opensne}\modellabel{\phi}{n,p,O,\avaluation}.
$$
\end{defn}

In a manner of speaking we have already seen $\modellabel{\phi}{n,p,\avaluation}$ from Definition~\ref{defn.pointwise.pred.nttn} in our logic:
\begin{lemm}
\label{lemm.Coquorum.denote}
Suppose $n\in\Time$ and $p\in\Pnt$ and $\avaluation$ is a valuation.
Then
$$
\modellabel{\phi}{n,p,\avaluation}=\modellabel{\Coquorum\phi}{n,p,O,\avaluation}.
$$
\end{lemm}
\begin{proof}
Direct from the clause for $\Coquorum$ from Figure~\ref{fig.3.derived}.
\end{proof}

\begin{lemm}
\label{lemm.pointwise.pred}
Suppose $\phi$ is a closed predicate and $n\in\Time$ and $\avaluation$ is a valuation.
Then $\phi$ is pointwise in $(n,\avaluation)$ in the sense of Definition~\ref{defn.pointwise.meaning} 
if and only if
$$
\Forall{p\in\Pnt}\Forall{O\in\opensne}\modellabel{\phi}{n,p,O,\avaluation}=\modellabel{\phi}{n,p,\avaluation}.
$$
\end{lemm}
\begin{proof}
By a routine calculation from Definitions~\ref{defn.pointwise.meaning} and~\ref{defn.pointwise.pred.nttn}.
\end{proof}

\begin{corr}
\label{corr.pointwise.pred}
Suppose $\phi$ is a closed predicate and $n\in\Time$ and $p\in\Pnt$ and $\avaluation$ is a valuation.
Suppose $\phi$ is pointwise in $(n,\avaluation)$.
Then
$$
\Forall{O\in\opensne}n,p,O\mentval\phi\liff n,p\mentval\phi. 
$$
\end{corr}
\begin{proof}
Routine from Definition~\ref{defn.validity.judgement} and Lemma~\ref{lemm.pointwise.pred}. %
\end{proof}


Being pointwise is particularly nice because of two simple syntactic criteria that guarantee it:
\begin{lemm}[Conditions for being pointwise]
\label{lemm.atomic.O}
Suppose $\phi$ is a predicate.
Then:
\begin{enumerate*}
\item\label{item.atomic.pointwise}
If $\phi$ is atomic (Notation~\ref{nttn.logic.terminology}(\ref{item.atomic.predicate}); having the form $\tf P(v_1,\dots,v_n)$ or $\tneg\tf P(v_1,\dots,v_n)$) then it is absolutely pointwise (Definition~\ref{defn.pointwise.meaning}).
\item\label{item.atomic.pointwise.2}
If $\phi$ does not mention $\someone$ or $\everyone$ except within the scope of a $\Quorum$ or $\Coquorum$ modality, then $\phi$ is absolutely pointwise.\footnote{In particular this holds if $\Quorum$ and $\Coquorum$ and only ever mentioned in combination with $\someone$ and $\everyone$, as in the compound modalities $\someoneAll$, $\everyoneAll$, $\QuorumBox$, and $\CoquorumDiamond$.}
\end{enumerate*} 
\end{lemm}
\begin{proof}
For part~1, we just unpack the definition of $\modellabel{\phi}{}$ in Figure~\ref{fig.3.phi.f} and note that the `$\avaluation(\Pnt)_{n,p}$' used in the clause for a predicate symbol depends just on the values of $n$, $p$, and $\avaluation$, but not on the open sets $O$ and $O'$. 

Part~2 follows by a routine induction on predicate syntax, using part~1 as the base case.
\end{proof}


\subsubsection{Pointwise, as a modality in our logic}

In fact, we can express the property of being pointwise within our logic:
\begin{defn}
\label{defn.pointwise.pred}
Suppose $\phi$ is a predicate.
Then define a predicate $\pointwise{\phi}$ by
$$
\pointwise{\phi} \ =\ \everyoneAll(\Quorum\phi \equiv \Coquorum\phi) .
$$
\end{defn}


\begin{lemm}
\label{lemm.pointwise.heart}
Suppose $\phi$ is a closed predicate and $n\in\Time$ and $p\in\Pnt$ and $\avaluation$ is a valuation.
Then 
$$
n,p\mentval\Quorum\phi \equiv \Coquorum\phi
\quad\text{if and only if}\quad
\bigvee_{O'\in\opensne}\modellabel{\phi}{n,p,O',\avaluation}=\bigwedge_{O'\in\opensne}\modellabel{\phi}{n,p,O',\avaluation}.
$$
(The predicate $\Quorum\phi \equiv \Coquorum\phi$ itself satisfies the conditions for being pointwise from Lemma~\ref{lemm.atomic.O}(\ref{item.atomic.pointwise.2}), so we can sensibly write $n,p\mentval \Quorum\phi \equiv \Coquorum\phi$ above.)
\end{lemm}
\begin{proof}
From Lemma~\ref{lemm.equiv}(\ref{item.equiv.1}) and the clauses for $\Quorum$ in Figure~\ref{fig.3.phi.f}, and $\Coquorum$ from Figure~\ref{fig.3.derived} (or Lemma~\ref{lemm.Coquorum.denote}). 
\end{proof}

\begin{lemm}
Suppose $\phi$ is a closed predicate and $n\in\Time$ and $p\in\Pnt$ and $\avaluation$ is a valuation.
Then $\phi$ is pointwise in $(n,\avaluation)$ if and only if $n,p\mentval\pointwise{\phi}$. 
\end{lemm}
\begin{proof}
By routine calculations using Lemma~\ref{lemm.pointwise.heart}, and the clause for $\everyoneAll$ in Figure~\ref{fig.3.derived}, noting the elementary fact that for every $O\in\opensne$,
$$
\bigwedge_{O'\in\opensne}\modellabel{\phi}{n,p,O',\avaluation}
\leq
\modellabel{\phi}{n,p,O,\avaluation}
\leq
\bigvee_{O'\in\opensne}\modellabel{\phi}{n,p,O',\avaluation} .
$$
\end{proof}



\subsubsection{An interlude: dense sets, and sets with a nonempty open interior}


We recall a pair of standard topological definitions which will turn up repeatedly later, most notably in the context of pointwise predicates as Proposition~\ref{prop.pointwise.dense.char} immediately below, and again later in Lemma~\ref{lemm.is}:
\begin{defn}[Dense subset]
\label{defn.dense.subset}
Suppose $(\Pnt,\opens)$ is a semitopology and suppose $P\subseteq\Pnt$ is a (not necessarily open) set of points.
\begin{enumerate*}
\item\label{item.dense.subset}
Say that $P$ is \deffont{dense} when $\Forall{O{\in}\opensne}O\intersectswith P$.
\item\label{item.noi.subset}
Say that $P$ \deffont{has a nonempty open interior} when $\Exists{O{\in}\opensne}O\subseteq P$.
\item\label{item.dense.noi}
Write $\dense(P)$ when $P$ is dense, and $\noi(P)$ when $P$ has a nonempty open interior. 
\end{enumerate*} 
\end{defn}



\begin{lemm}
\label{lemm.dense.up-closed}
Suppose $(\Pnt,\opens)$ is a semitopology and $P,P'\subseteq\Pnt$.
Then:
\begin{enumerate*}
\item\label{item.dense.up-closed} 
If $P\subseteq P'$ then $\dense(P)\limp \dense(P')$ (being dense is \emph{up-closed}).
\item\label{item.noi.up-closed} 
If $P\subseteq P'$ then $\noi(P)\limp \noi(P')$ (having a non-empty open intersection is \emph{up-closed}).
\item\label{item.noi.noi} 
If $P\in\opensne$ then $\noi(P)$ (a nonempty open set has a nonempty open interior).
\end{enumerate*}
\end{lemm}
\begin{proof}
By routine sets calculations from Definition~\ref{defn.dense.subset}.
\end{proof}

\begin{lemm}
Suppose $(\Pnt,\opens)$ is a semitopology and $P\subseteq\Pnt$.
Then
$$
\dense(P) \liff \neg\noi(\Pnt\setminus P) 
\quad\text{and}\quad 
\noi(P) \liff \neg\dense(\Pnt\setminus P) . 
$$
In this sense, $\noi$ and $\dense$ form a \emph{de Morgan dual pair} of operators --- just as $\forall=\neg\exists\neg$ (as noted in Figure~\ref{fig.easy.equivalences}) and $\exists=\tneg\forall\tneg$.
\end{lemm}
\begin{proof}
We sketch the (easy) argument:
\begin{itemize}
\item
Suppose $P$ is not dense; by Definition~\ref{defn.dense.subset} this means that there exists $O\in\opensne$ such that $P\cap O=\varnothing$, thus $O\subseteq\Pnt\setminus P$, thus $\Pnt\setminus P$ has a nonempty open interior.
\item
Suppose $\Pnt\setminus P$ has a nonempty open interior $\varnothing\neq O\subseteq\Pnt\setminus P$.
Then clearly $P$ is not dense.
\qedhere\end{itemize}
\end{proof}

More on how `being dense' works in semitopologies is in~\cite[Chapter~11]{gabbay:semdca}.
The notion behaves somewhat differently than it does in topologies, but not in ways that will affect us here.

\subsubsection{Application to pointwise predicates}


Technically, Proposition~\ref{prop.pointwise.dense.char} is a corollary of Lemma~\ref{lemm.dense.char} and Corollary~\ref{corr.pointwise.pred}.
The result links the four logical modalities $\someoneAll$, $\everyoneAll$, $\QuorumBox$, and $\CoquorumDiamond$ to elementary topological properties of the sets of points at which these modalities are valid --- being nonempty, being full, having a nonempty open interior, and being dense.
We will use Proposition~\ref{prop.pointwise.dense.char}, both directly and via other lemmas that depend on it.

The most important use-case of Proposition~\ref{prop.pointwise.dense.char} will be when $\phi$ is atomic (Notation~\ref{nttn.logic.terminology}(\ref{item.atomic.predicate}) --- having the form 
$$
\tf P(v_1,\dots,v_n)
\quad\text{or}\quad
\tneg\tf P(v_1,\dots,v_n)
$$
--- since then by Lemma~\ref{lemm.atomic.O} $\phi$ is automatically pointwise:
\begin{prop}
\label{prop.pointwise.dense.char}
Suppose $\phi$ is a closed pointwise (Definition~\ref{defn.pointwise.meaning}) predicate and $n\in\Time$ and $\avaluation$ is a valuation.
Then: 
\begin{enumerate*}
\item\label{item.dense.char.pointwise.someoneAll}
$n\mentval\someoneAll\phi$ if and only if $\varnothing\neq\{p\in\Pnt \mid n,p\mentval\phi\}$. 
\item\label{item.dense.char.pointwise.everyoneAll}
$n\mentval\everyoneAll\phi$ if and only if $\Pnt=\{p\in\Pnt \mid n,p\mentval\phi\}$. 
\item\label{item.dense.char.pointwise.QuorumBox}
$n\mentval\QuorumBox\phi$ if and only if $\noi(\{p\in\Pnt \mid n,p\mentval\phi\})$.
\item\label{item.dense.char.pointwise.CoquorumDiamond}
$n\mentval\CoquorumDiamond\phi$ if and only if $\dense(\{p\in\Pnt \mid n,p\mentval\phi\})$.
\end{enumerate*} 
Recall that:
\begin{itemize*}
\item
$\dense$ and $\noi$ are from Definition~\ref{defn.dense.subset}.
\item
The notations $n,p\mentval\phi$ and $n\mentval\phi$ are from Definition~\ref{defn.validity.judgement}.
\end{itemize*}
\end{prop}
\begin{proof}
If $p$ is pointwise then by Corollary~\ref{corr.pointwise.pred} $n,p,O\mentval\phi$ if and only if $n,p\mentval\phi$, for any $O\in\opensne$, and we just simplify parts~\ref{item.dense.char.someoneAll} to~\ref{item.dense.char.CoquorumDiamond} of Lemma~\ref{lemm.dense.char}:
\begin{enumerate} 
\item
We reason as follows:
$$
\begin{array}{r@{\ }l@{\quad}l}
n\mentval\someoneAll\phi
\liff&
\Exists{O\in\opensne}O\between\{p\in\Pnt \mid n,p\mentval\phi\}
&\text{Lemma~\ref{lemm.dense.char}(\ref{item.dense.char.someoneAll})} 
\\
\liff&
\Pnt\between \{p\in\Pnt \mid n,p\mentval\phi\}
&\text{Fact}
\\
\liff&
\{p\in\Pnt \mid n,p\mentval\phi\}\neq\varnothing
&\text{Fact}
\end{array}
$$
\item
We reason as follows:
$$
\begin{array}{r@{\ }l@{\quad}l}
n\mentval\everyoneAll\phi
\liff&
\Forall{O\in\opensne}O\subseteq\{p\in\Pnt \mid n,p\mentval\phi\}
&\text{Lemma~\ref{lemm.dense.char}(\ref{item.dense.char.everyoneAll})} 
\\
\liff&
\Pnt\subseteq\{p\in\Pnt \mid n,p\mentval\phi\}
&\text{Fact}
\\
\liff&
\{p\in\Pnt \mid n,p\mentval\phi\}=\Pnt
&\text{Fact}
\end{array}
$$
\item
We reason as follows:
$$
\begin{array}{r@{\ }l@{\quad}l}
n\mentval\QuorumBox\phi
\liff&
\Exists{O\in\opensne}O\subseteq\{p\in\Pnt \mid n,p\mentval\phi\}
&\text{Lemma~\ref{lemm.dense.char}(\ref{item.dense.char.QuorumBox})} 
\\
\liff&
\noi(\{p\in\Pnt \mid n,p\mentval\phi\})
&\text{Definition~\ref{defn.dense.subset}}
\end{array}
$$
\item
We reason as follows:
$$
\begin{array}{r@{\ }l@{\quad}l}
n\mentval\CoquorumDiamond\phi
\liff&
\Forall{O\in\opensne}O\between\{p\in\Pnt \mid n,p\mentval\phi\}
&\text{Lemma~\ref{lemm.dense.char}(\ref{item.dense.char.CoquorumDiamond})} 
\\
\liff&
\dense(\{p\in\Pnt \mid n,p\mentval\phi\})
&\text{Definition~\ref{defn.dense.subset}}
\end{array}
$$
\qedhere\end{enumerate}
\end{proof}
See also Lemma~\ref{lemm.allbut.concretely}, which instantiates Proposition~\ref{prop.pointwise.dense.char}(\ref{item.dense.char.pointwise.QuorumBox}\&\ref{item.dense.char.pointwise.CoquorumDiamond}) to a special case
familiar from the literature on distributed systems.


\subsection{Correct predicates (another well-behavedness property)}

\begin{rmrk}
We defined $\threeCorrect=\{\tvT,\tvF\}$ in Definition~\ref{defn.tv.ment}(\ref{item.tv.ment.correct}) and a corresponding modality $\modTF$ in Figure~\ref{fig.3} (which returns $\tvT$ when its argument is in $\threeCorrect$ and $\tvF$ otherwise), and we studied some of their properties, most notably in Lemma~\ref{lemm.TF.correct}.
 
We now extend this to notions of \emph{correctness} on predicates $\phi$ and predicate symbols $\tf P$.

Correctness turns out to be highly relevant to how we will use our logic.
Partly this is because when a predicate is correct, we are returned to the familiar two-valued world of Boolean logic --- but more subtly, correctness is interesting because of how this concept interacts with the predicate and modal parts of our logic; in certain circumstances we can use a combination of correctness \emph{somewhere} and validity \emph{everywhere} to deduce information about truth (see Remark~\ref{rmrk.why.exciting}) and this turns out to be a good logical model for how reasoning on distributed algorithms is often done in practice.

We start with the definition and some basic lemmas, and develop from there:
\end{rmrk}


\begin{defn}[Correctness]
\label{defn.correct}
Suppose $\phi$ is a predicate and $a\in\tf{VarSymb}$ is a variable symbol. 
Then:
\begin{enumerate*}
\item\label{item.correct.phi}
Consistent with the definition of $\threeCorrect=\{\tvT,\tvF\}$ from Definition~\ref{defn.tv.ment}(\ref{item.tv.ment.correct}),
we define $\correct{\phi}$, read \deffont{$\phi$ is correct}, by
$$
\correct{\phi}=\modTF\phi .
$$
Note that $\phi$ being correct does not mean that $\phi$ returns $\tvT$; it means that $\phi$ returns $\tvT$ or $\tvF$ and in particular it does \emph{not} return $\tvB$.
\item\label{item.correct}\label{item.crashed}
Suppose $\tf P\in\PredSymb$ is a unary predicate symbol.
Define a predicate $\correct{\tf P}$ by 
$$
\correct{\tf P}\ =\ \tall a.\correct{\tf P(a)}\ \stackrel{N\ref{nttn.hos}}{=}\ \tall\modTF\tf P
 .
$$
\item
If $\tf P_1,\dots,\tf P_n$ are unary predicate symbols then we may write 
$$
\correct{\tf P_1,\dots,\tf P_n}
\quad\text{for}\quad 
\correct{\tf P_1}\tand\dots\tand\correct{\tf P_n}.
$$
(See for example axiom \rulefont{LdrCorrect} in Figure~\ref{fig.logical.paxos}.) 
\item\label{item.P-correct}
Suppose $n\in\Time$ and $p\in\Pnt$ and $\avaluation$ is a valuation, and suppose $\tf P$ is a predicate symbol.
Then when $n,p\mentval\correct{\tf P}$ we may call $p$ \deffont{$\tf P$-correct}. 
\end{enumerate*}
\end{defn}

\begin{lemm}
\label{lemm.correct.valid}
Suppose $\acontext$ is a context and $\phi$ is a closed predicate.
Then 
$$
\acontext\ment\phi
\quad\text{if and only if}\quad
\modellabel{\phi}{\acontext}\in\threeCorrect.
$$
\end{lemm}
\begin{proof}
This just re-states Lemma~\ref{lemm.tv.ment.TF}(\ref{item.tv.ment.TF.2}).
\end{proof}


\begin{rmrk}
\label{rmrk.tf.synonyms}
Definition~\ref{defn.correct}(\ref{item.correct.phi}) gives us two synonyms for the modality expressing that the truth-value returned by a predicate $\phi$ is either $\tvT$ or $\tvF$ (not $\tvB$): $\modTF$ and $\tf{correct}$.

Although these are synonymous, there is a difference in emphasis:
\begin{itemize*}
\item
$\modTF$ is a logical modality and just means `true or false'.
\item
The literature on distributed systems uses a term `correctness', which indicates that a process succeeds and returns a correct value. 
In our logical model this corresponds to returning $\tvT$ or $\tvF$ --- and thus not returning $\tvB$. 
So when we write $\tf{correct}$, the emphasis is on `not $\tvB$' with emphasis on a process `not failing'.
\end{itemize*}
We use whichever notation seems most natural in context. 
For example, in Figure~\ref{fig.logical.paxos}:
\begin{enumerate*}
\item
In \rulefont{LdrCorrect} we write $\correct{\tf{propose},\tf{write},\tf{decide}}$ because intuitively we are asserting that at a particular time the leader \emph{does not crash} when proposing, writing, or deciding.
\item
In \rulefont{LdrExist} we write $\modTF\tf{leader}$, because $\tf{leader}$ is a predicate that identifies whether a participant is the leader at a particular time, and this is either true or false.\footnote{In practice determining who the leader is require a participant to compute a function, but the computation is not distributed, in the sense that prior agreement amongst participants is assumed on what the function is and what seed data it uses, and each participant computes it locally to get (the same) definitive yes/no answer.}
\end{enumerate*} 
If the reader prefers, they can safely read $\modTF$ every time they see $\tf{correct}$, and they will be precisely correct (pun intended). 
\end{rmrk}



Many predicates of interest are correct, including correctness itself, our truth-modalities, and equality and equivalence:
\begin{lemm}
\label{lemm.equality.correct}
Suppose $\acontext$ is a context and $\phi$ and $\phi'$ are closed predicates and $v,v'\in\Val$.
Then:
\begin{enumerate*}
\item
$\modellabel{\correct{\correct{\phi}}}{\acontext}=\tvT$.
\quad We can (perhaps cheekily) read this as `correctness is correct'.
\item
$\modellabel{\correct{\tf Q\phi}}{\acontext}=\tvT$ for $\tf Q\in\{\modT,\modF,\modB,\modTB,\modTF,\modFB\}$. 
\item\label{item.equality.correct}
$\modellabel{\correct{v\teq v'}}{\acontext}=\tvT$.
In words: equality is correct.
\item\label{item.equivalence.correct}
$\modellabel{\correct{\phi\equiv \phi'}}{\acontext}=\tvT$.
In words: equivalence is correct.
\end{enumerate*}
\end{lemm}
\begin{proof}
By routine checking from Figures~\ref{fig.3}, \ref{fig.3.phi.f}, \ref{fig.3.derived}, and (for $\equiv$)~\ref{fig.3.iff}.
\end{proof}

Lemma~\ref{lemm.correct.P.Pv} and Corollary~\ref{corr.modTF.texi} are straightforward, but we spell it out:
\begin{lemm}
\label{lemm.correct.P.Pv}
Suppose $\tf P\in\PredSymb$ is a unary predicate symbol and $v\in\Val$ and $\acontext$ is a context.
Then 
$$
\modellabel{\correct{\tf P}}{\acontext} = \bigwedge_{v{\in}\Val}\modellabel{\correct{\tf P(v)}}{\acontext}
\quad\text{and}\quad 
\modellabel{\correct{\tf P}}{\acontext}\leq \modellabel{\correct{\tf P(v)}}{\acontext} .
$$
\end{lemm}
\begin{proof}
By Definition~\ref{defn.correct}(\ref{item.correct}) and Figure~\ref{fig.3.derived} 
$$
\modellabel{\correct{\tf P}}{\acontext} 
= 
\modellabel{\tall a.\correct{\tf P(a)}}{\acontext} 
= 
\bigwedge_{v{\in}\Val} \modellabel{\correct{\tf P(v)}}{\acontext} 
\leq 
\modellabel{\correct{\tf P(v)}}{\acontext} .
\qedhere$$
\end{proof}


\begin{corr}
\label{corr.modTF.texi}
Suppose $\acontext$ is a context and $\tf P\in\PredSymb$ is a unary predicate symbol and $a\in\tf{VarSymb}$ is a variable symbol. 
Then 
$$
\acontext\ment\correct{\tf P}
\quad\text{implies}\quad
\acontext\ment\correct{\texi\tf P}\ \land\ \acontext\ment\correct{\tall\tf P}.
$$
\end{corr}
\begin{proof}
Suppose $\acontext\ment\correct{\tf P}$.
By Lemma~\ref{lemm.correct.P.Pv} we see that $\Forall{v\in\Val}\acontext\ment\correct{\tf P(v)}$.
By Lemma~\ref{lemm.correct.valid} $\Forall{v\in\Val}\modellabel{\tf P(v)}{\acontext}\in\threeCorrect$.
Now the clauses for $\texi$ in Figure~\ref{fig.3.phi.f} and $\tall$ in Figure~\ref{fig.3.derived} take a supremum and infimum respectively of $\{\modellabel{\tf P(v)}{\acontext} \mid v\in\Val\}\subseteq\threeCorrect=\{\tvT,\tvF\}$, and this supremum or infimum is also clearly in $\threeCorrect=\{\tvT,\tvF\}$.
We just use Lemma~\ref{lemm.correct.valid}. 
\end{proof}


Proposition~\ref{prop.tv.ment.TF.model}(\ref{item.tv.ment.TF.phi.lr}) is a wrapper for Lemma~\ref{lemm.TF.correct}.
This deceptively simple equivalence turns out to be a key way that correctness can be usefully applied:
\begin{prop}
\label{prop.tv.ment.TF.model}
Suppose $\phi$ is a closed predicate and $\acontext$ is a context and $v\in\Val$.
Then:
\begin{enumerate*}
\item\label{item.tv.ment.TF.phi.lr}
$\acontext\ment\correct{\phi}$ if and only if $\acontext\ment\phi\equiv\modT\phi$. 

In words: correct $\phi$ are characterised precisely by validity being equivalent to truth.
\item\label{item.tv.ment.TF.P}
$\acontext\ment\correct{\tf P}$ if and only if $\Forall{v\in\Val}(\acontext\ment\tf P(v)\equiv\modT\tf P(v))$. 
\item\label{item.modTF.texi.P.valid}
$\acontext\ment\correct{\tf P}$ implies $\acontext\ment\texi\tf P\equiv \modT\texi\tf P$ and $\acontext\ment\tall\tf P\equiv\modT\tall\tf P$. 
\end{enumerate*}
\end{prop}
\begin{proof}
We consider each part in turn:
\begin{enumerate}
\item
Suppose $\acontext\ment\correct{\phi}$.
By Lemma~\ref{lemm.TF.correct} $\modT\modellabel{\phi}{\acontext}=\modellabel{\phi}{\acontext}$.
By Lemma~\ref{lemm.equiv}(\ref{item.equiv.1}) and the clause for $\modT$ in Figure~\ref{fig.3.phi.f}, $\acontext\ment \phi\equiv\modT\phi$.
The reverse implication just reverses the same reasoning.
\item
We combine Lemma~\ref{lemm.correct.P.Pv} with part~\ref{item.tv.ment.TF.phi.lr} of this result for $\phi=\tf P(v)$ for every $v\in\Val$.
\item
From Corollary~\ref{corr.modTF.texi} and Proposition~\ref{prop.tv.ment.TF.model}(\ref{item.tv.ment.TF.phi.lr}).
\qedhere\end{enumerate}
\end{proof}




Proposition~\ref{prop.correct.cobox.boxT} explores how correctness interacts with our compound modalities from Subsection~\ref{subsect.compound.modalities}.
\begin{prop}
\label{prop.correct.cobox.boxT}
Suppose $\phi$ is a closed predicate and $\tf P\in\PredSymb$ is a unary predicate symbol. 
Suppose $n\in\Time$ and $\avaluation$ is a valuation.
Then:
\begin{enumerate*}
\item\label{item.correct.cobox.boxT.1}
$n\mentval\QuorumBox\correct{\phi}\tand\everyoneAll\phi$ implies $n\mentval \modT\QuorumBox\phi$.\footnote{Alert: note that we write $\everyoneAll\phi$ in the right-hand conjunct, not $\CoquorumDiamond\phi$.}
\item\label{item.correct.cobox.boxT.2}
$n\mentval\QuorumBox\correct{\tf P}\tand\everyoneAll\texi\tf P$ implies 
$n\mentval\modT\QuorumBox\texi\tf P$.
\end{enumerate*}
\end{prop}
\begin{proof}
We consider each part in turn:
\begin{enumerate}
\item
Suppose $n\mentval\QuorumBox\correct{\phi}\tand\everyoneAll\phi$.
By Proposition~\ref{prop.someone.implies.someoneAll}(\ref{item.someone.implies.someoneAll.4b})
$n\mentval\QuorumBox(\correct{\phi}\tand\phi)$.
By Proposition~\ref{prop.tv.ment.TF.model}(\ref{item.tv.ment.TF.phi.lr}) and Lemma~\ref{lemm.someone.commutation} (for $\QuorumBox$)
we have $n\mentval\modT\QuorumBox\phi$ as required.
\item
We combine part~\ref{item.correct.cobox.boxT.1} of this result with Corollary~\ref{corr.modTF.texi}.
\qedhere\end{enumerate}
\end{proof}


\begin{rmrk}[Why this is interesting]
\label{rmrk.why.exciting}
The statement of Proposition~\ref{prop.correct.cobox.boxT} is quite heavy on the symbols.
For clarity, let us re-read it, but let us replace $\ment\QuorumBox\tf{correct}$ with the phrase `mostly correct', and $\ment\everyoneAll$ with the phrase `everywhere valid', and $\texi$ with the phrase `for some value'.
Then we get:
\begin{enumerate*}
\item
If $\phi$ is \emph{mostly correct} and \emph{everywhere valid}, then it is \emph{mostly true}.
\item
If $\tf P$ is \emph{mostly correct} and \emph{everywhere valid} for some value, then it is \emph{mostly true} for some value.
\end{enumerate*}
So what we see is that we can promote weak statements about validity everywhere, with stronger but more limited statements about correctness, to obtain limited statements about truth.

The reader familiar with how distributed algorithms work may recognise this style of reasoning.
So: here it is beginning to appear in our logic.
\end{rmrk}

\subsection{$\mrup{a}{\phi}{v}$: a most recent $v$ that makes $\phi[a\ssm v]$ true}
\label{subsect.mrup}

\begin{rmrk}
\label{rmrk.explain.mrup}
We see from Figure~\ref{fig.3.derived} and Remark~\ref{rmrk.comments.on.semantics}(\ref{item.mru.program}) that $\mrup{a}{\phi}{v}$ looks back into the past, searching for a most recent time when $\model{\texi a.\phi}=\tvT$, and checks if at that most recent time, $\model{\phi[a\ssm v]}=\tvT$.
$\mrup{a}{\phi}{v}$ returns $\tvT$ if this search succeeds, and $\tvF$ if this fails.
It never returns $\tvB$.

In words: $\mrup{a}{\phi}{v}$ checks whether $v$ is a value that has most recently made $\phi$ true.

We need $\mrup{a}{\phi}{v}$ for our axiomatisation of Paxos in Figure~\ref{fig.logical.paxos}, where 
\begin{itemize*}
\item
in \rulefont{PaxSend?} we write $\mru{\tf{accept}}{v}$\ (short for $(\mrup{a}{\tf{accept}(a))}{v}$, as per Figure~\ref{fig.hos}), and 
\item
in \rulefont{PaxWrite?} we write $\mru{(\someone\tf{accept})}{v}$\ (short for $\mrup{a}{(\someone\tf{accept}(a))}{v}$).
\end{itemize*} 
We read $\mrup{a}{\phi}{v}$ as 
\begin{quote}
\deffont{$v$ is a value that most recently makes $\phi[a\ssm v]$ true} .
\end{quote}
In the common case that $\phi=\tf P(a)$ for some unary predicate symbol $\tf P$, then we might just say
\begin{quote}
\deffont{$v$ is a value that most recently makes $\tf P$ true}.
\end{quote} 
In this Subsection we study the properties of this construct in more detail.
\end{rmrk}



\begin{lemm}
\label{lemm.mru.unambivalent}
Suppose $\texi a.\phi$ is a closed predicate and $v\in\Val$ and $\acontext$ is a context.
Then:
\begin{enumerate*}
\item
$\acontext\ment\correct{\mrup{a}{\phi}{v}}$.
\item
$\acontext\ment\mrup{a}{\phi}{v}\equiv\modT\mrup{a}{\phi}{v}$.
\end{enumerate*}
\end{lemm}
\begin{proof}
We consider each part in turn:
\begin{enumerate}
\item
By Definition~\ref{defn.correct}(\ref{item.correct.phi}) $\correct{\mrup{a}{\phi}{v}}=\modTF\mrup{a}{\phi}{v}$, and this means that $\mrup{a}{\phi}{v}$ always returns $\tvT$ or $\tvF$ (but not $\tvB$).
We just check the relevant clause in Figure~\ref{fig.3.derived} and see that this is so.
\item
We just use part~1 of this result and Proposition~\ref{prop.tv.ment.TF.model}(\ref{item.tv.ment.TF.phi.lr}).
\qedhere\end{enumerate}
\end{proof}


\begin{rmrk}
\label{rmrk.explain.mru}
\leavevmode
\begin{enumerate}
\item\label{item.explain.mru.T}
Lemma~\ref{lemm.mru.unambivalent}(2) %
notes that $\mrup{a}{\phi}{v}$ and $\modT(\mrup{a}{\phi}{v})$ are extensionally equivalent (Definition~\ref{defn.phi.equivalent}).

We may write one or the other, depending on what seems most clear in a proof.
In particular, if we are using \strongmodusponens then we might retain a $\modT$ modality for clarity.
\item 
We try to head off an easy error: note that $\mrup{a}{\phi}{v}$ is not equivalent to $\modT\recent\phi[a\ssm v]$, nor the (by Lemma~\ref{lemm.someone.commutation}) equivalent form $\recent\modT\phi[a\ssm v]$.
\begin{itemize*}
\item
$\modT\recent\phi[a\ssm v]$ returns $\tvT$ when $\phi[a\ssm v]$ has returned $\tvT$ at some point in the past.
\item
$\mrup{a}{\phi}{v}$ (`$v$ is a value that most recently makes $\phi[a\ssm v]$ true') returns $\tvT$ when $\phi[a\ssm v]$ has returned $\tvT$ in \emph{the most recent stage} such that $\texi a.\phi$ returned $\tvT$ (i.e. such that $\phi[a\ssm v']$ returned $\tvT$ for \emph{any} $v'$).
\end{itemize*}
These are not the same thing.
\end{enumerate}
\end{rmrk}

Remark~\ref{rmrk.explain.mru} notes that $\mrup{a}{\phi}{v}$ is different from $\modT\recent\phi[a\ssm v]$.
That said, there is a close link between them, which will turn out to be important for understanding our axiomatisation of Paxos (see Lemma~\ref{lemm.evaporates.simplify} and surrounding discussion): 
\begin{lemm}
\label{lemm.evaporates}
Suppose $\texi a.\phi$ is a closed predicate and $\acontext=(n,p,O,\avaluation)$ is a context. 
Then
$$
\acontext\ment(\texi v.\mrup{a}{\phi}{v})\equiv \modT\recent\texi a.\phi.
$$
\end{lemm}
\begin{proof}
First, note from Lemma~\ref{lemm.someone.commutation} that $\ment\modT\recent\texi a.\phi\equiv \recent\texi a.\modT\phi$.
And, note that $\texi a.\phi$ and $\texi v.\phi[a\ssm v]$ denote the same syntax, as per Remark~\ref{rmrk.pedant.point}(\ref{item.pedant.texi.v}).

We unpack the clause for $\modellabel{\mrup{a}{\phi}{v}}{\acontext}$ in Figure~\ref{fig.3.derived} and note that it returns $\tvT$ when there exist $v\in\Val$ and $0\leq n'<n$ such that $n',p,O\mentval \modT\phi[a\ssm v]$, and $\tvF$ otherwise.
But this is also what $\modellabel{\recent\texi v.\modT(\phi[a\ssm v])}{\acontext}$ does, so by Lemma~\ref{lemm.equiv}(\ref{item.equiv.1}) we are done.
\end{proof}

Lemma~\ref{lemm.mru.today.yesterday} connects $\mrup{a}{\phi}{v}$ to the fact that our notion of time is discrete and so has a well-defined notion of a `tomorrow' (i.e. a `next timestep'):
\begin{lemm}
\label{lemm.mru.today.yesterday}
Suppose $\texi a.\phi$ is a closed predicate and $v\in\Val$ and $\acontext$ is a context. 
.
Then
$$
\acontext\ment \tomorrow(\mrup{a}{\phi}{v}) \equiv (\modT\phi[a\ssm v] \tor ((\tneg\modT\texi a.\phi) \tand \mrup{a}{\phi}{v})).
$$
\end{lemm}
\begin{proof}
It is a fact of Figure~\ref{fig.3.derived} that $\acontext\ment\tomorrow\yesterday\psi\equiv\psi$ for any closed $\psi$.
The result follows from this observation and from the clause for $\mrup{a}{\phi}{v}$ in Figure~\ref{fig.3.derived}. 
\end{proof}



\begin{corr}
\label{corr.mru.today.to.yesterday}
Suppose $\texi a.\phi$ is a closed predicate and $v\in\Val$ and $\acontext$ is a context. 
Then:
\begin{enumerate*}
\item\label{item.mru.today.implies.tomorrow}
$\acontext\ment \phi[a\ssm v] \timpc \tomorrow(\mrup{a}{\phi}{v})$.
\item\label{item.mru.tomorrow.options}
$\acontext\ment \tomorrow(\mrup{a}{\phi}{v}) \timpc (\phi[a\ssm v]\tor \mrup{a}{\phi}{v})$. 
\end{enumerate*}
\end{corr}
\begin{proof}
Elementary from Lemma~\ref{lemm.mru.today.yesterday} using \strongmodusponens.
\end{proof}


\begin{lemm}
\label{lemm.mru.to.recent}
Suppose $\texi a.\phi$ is a closed predicate %
and $v,v'\in\Val$ and $\acontext$ is a context. 
Then:
\begin{enumerate*}
\item\label{item.mru.to.recent.2}
$\acontext\ment(\mrup{a}{\phi}{v}\tand\mrup{a}{\phi}{v'}) \timpc \recent(\phi[a\ssm v]\tand\phi[a\ssm v'])$.
\item\label{item.mru.to.recent.1}
$\acontext\ment(\mrup{a}{\phi}{v}) \timpc \recent(\phi[a\ssm v])$.

If $\phi=\modality{\tf Q}\tf P(a)$ for a unary predicate symbol $\tf P\in\PredSymb$ and some prefix of modalities $\modality{\tf Q}$, then we can elide the variable symbol $a$ as per Notation~\ref{nttn.hos} and Figure~\ref{fig.hos}, and this becomes
$$
\acontext\ment\mru{(\modality{\tf Q}\tf P)}{v} \timpc \recent\modality{\tf Q}\tf P(v).
$$
\end{enumerate*}
\end{lemm}
\begin{proof}
We consider each part in turn:
\begin{enumerate}
\item
Unpacking the clause for $\mruarrow$ in Figure~\ref{fig.3.derived}, if $v$ and $v'$ are values that most recently make $\phi[a\ssm\text{-}]$ be $\tvT$, then certainly $\recent(\phi[a\ssm v]\tand\phi[a\ssm v'])$.
\item
Follows from part~\ref{item.mru.to.recent.2} as the special case that $v'=v$.
\qedhere\end{enumerate}
\end{proof}


\section{\lowercase{n}-twined semitopologies}
\label{sect.n-twined}



\subsection{Antiseparation properties}

We are particularly interested in models over semitopologies whose open sets have certain \emph{antiseparation properties}, by which we mean properties that involve the existence of nonempty intersections between open sets.\footnote{For comparison, topology is often concerned with \emph{separation} properties, like Hausdorff separation, which have to do with existence of \emph{empty} intersections between open sets.  In this sense, the interest in this Section is the precise dual to where the emphasis traditionally lies in topology.}

The design space of antiseparation properties is rich, as the first author explored in~\cite{gabbay:semdca} --- but for our needs in this paper we will just one canonical property which we will call being \emph{$n$-twined}.
We define the $n$-twined property in Definition~\ref{defn.twined}, we explain why it matters in Remark~\ref{rmrk.motivate.n-twined}, and then we study its properties.


\begin{defn}
\label{defn.twined}
\leavevmode
Suppose $(\Pnt,\opens)$ is a semitopology and $n\in\Ngeqz$.
Then call $(\Pnt,\opens)$ \deffont{$n$-twined} when every collection of $n$ nonempty open sets has a nonempty open intersection.
In symbols:
$$
\Forall{O_1,\dots,O_n\in\opensne}O_1\cap\dots\cap O_n\neq\varnothing.
$$
\end{defn}

\begin{lemm}
\label{lemm.1-twined}
We note two edge cases:
\begin{enumerate*}
\item
A semitopology is $0$-twined if and only if it is nonempty (Notation~\ref{nttn.empty.semitopology}).
\item
Every semitopology is $1$-twined.
\end{enumerate*}
\end{lemm}
\begin{proof}
We consider each part in turn:
\begin{enumerate}
\item
If $n=0$ then we need to show that $\Pnt$ (the intersection of zero nonempty sets) is not equal to $\varnothing$.
by Notation~\ref{nttn.empty.semitopology} this is true precisely when the semitopology is nonempty. 
\item
By Definition~\ref{defn.twined}, we must show that $O\neq\varnothing$ for every $O\in\opensne$.
This is true.
\qedhere\end{enumerate}
\end{proof}

\begin{rmrk}[Discussion of the special cases $n=2$ and $n=3$]
\label{rmrk.motivate.n-twined}
Lemma~\ref{lemm.1-twined} shows that being $0$-twined and $1$-twined are straightforward properties.
In contrast, the $2$- and $3$-twined semitopologies will be very interesting to us in this paper.
We give a high-level discussion of why this is:
\begin{itemize}
\item
Suppose we are in a distributed system running on a $2$-twined semitopology.
We assume that participants are honest.
Suppose further that we have 
\begin{itemize*}
\item
a nonempty open set of participants $O\in\opensne$ who agree that some predicate $\phi$ is true (respectively: valid), and 
\item
a nonempty open set of participants $O'\in\opensne$ who agree that some predicate $\phi'$ is true (respectively: valid).
\end{itemize*}
Then because the system is $2$-twined, there exists some $p\in O\cap O'$ who agrees that both $\phi$ and $\phi'$ are true (respectively: valid).\footnote{We will render this in our logic later on, in Lemma~\ref{lemm.supertwined.qcbt}(\ref{item.intertwined.qcbt.2}).}

This simple observation underlies guarantees of consistent outputs in many distributed algorithms, because it guarantees that $\phi$ and $\phi'$ are not logically conflicting --- e.g. that we do not have something like $\phi'=\modF\phi$ --- because if they were, then an honest $p$ could not believe both.

In the special case that $\phi$ is `the output of our consensus algorithm is $v$' and $\phi'$ is `the output of our consensus algorithm is $v'$', the $2$-twined property makes the strong guarantee that where consensus algorithms succeed (i.e. where participants do not crash), the output is consistent. 
\item
Suppose we are in a distributed system running on a $3$-twined semitopology.
Some participants may be dishonest, in that they may report inconsistent predicates. 

Suppose further that we have nonempty open sets $O,O'\in\opensne$ of participants who agree on the truth (respectively: validity) of predicates $\phi$ and $\phi'$ respectively.
Suppose further that we have a nonempty open set $O''\in\opensne$ of participants who are honest (meaning that we bound dishonest participants to at most some \emph{closed} set). 
Because the system is $3$-twined, there exists some $p\in O\cap O'\cap O''$ --- this $p$ is a participant who agrees on the truth (respectively: validity) of $\phi$ and $\phi'$ and is also honest.\footnote{We will render this in our logic twice: in Lemma~\ref{lemm.supertwined.qcbt}(\ref{item.supertwined.qcbt.2}) as a general result, and then as a useful special case in Corollary~\ref{corr.two.quorums.intersect.correct}(\ref{item.two.quorums.intersect.correct.supertwined}).}

In this case, $\phi\tand\phi'$ must be consistent, because if it were not then the honest participant $p$ could not believe both $\phi$ and $\phi'$.
Just as for the $2$-twined case, this simple argument underlies many consistency arguments in distributed systems in the presence of dishonest participants.
\end{itemize}
There is no barrier to considering $n$-twined semitopologies for $n\geq 4$, and this would just give us stronger guarantees of intersection.

A brief survey of intersection properties in the literature is in Remark~\ref{rmrk.ntwined.and.the.literature}
\end{rmrk}


\subsection{Being $n$-twined and being dense / having nonempty open interior}

The property of being $n$-twined is closely related to the topological notion of being a \emph{dense set} and having a \emph{nonempty open interior} (Definition~\ref{defn.dense.subset}).
We explore this in Lemma~\ref{lemm.is} and Proposition~\ref{prop.noi.dense}; then in Proposition~\ref{prop.allbut.supertwined} we specialise the discussion to the important special case of $\AllBut{N}{f}$.

\begin{lemm}
\label{lemm.is}
Suppose $(\Pnt,\opens)$ is a semitopology and $n\in\mathbb N_{\geq 1}$.
Then:
\begin{enumerate*}
\item\label{item.is.supertwined}
$(\Pnt,\opens)$ is $n$-twined if and only if every collection of $n\minus 1$ nonempty open sets has a dense intersection: 
$$
\Forall{O_1,\dots O_{n\minus 1}\in\opensne}\dense(O_1\cap\dots\cap O_{n\minus 1}).
$$
\item\label{item.is.intertwined}
As a corollary, $(\Pnt,\opens)$ is \xxtwined if and only if every nonempty open set is dense: 
$$
\Forall{O\in\opensne}\dense(O).
$$
\end{enumerate*}
\end{lemm}
\begin{proof}
Both parts follow from Definitions~\ref{defn.dense.subset}(1) and~\ref{defn.twined}, just by pushing around definitions.

One corner case merits brief mention: if $n=1$ then we must show that every collection of $0$ sets (i.e. $\varnothing$) is dense.
So, we need to show that $\bigcap\varnothing$ is dense.
But by convention, the intersection of \emph{no} open sets is the whole of $\Pnt$, and this is indeed dense, so we are done. 
\end{proof}


\begin{prop}
\label{prop.noi.dense}
Suppose $(\Pnt,\opens)$ is a semitopology and $n\in\mathbb N_{\geq 1}$.
Then the following are equivalent:
\begin{enumerate*}
\item
$(\Pnt,\opens)$ is $n$-twined. 
\item
$\Forall{P_1,\dots,P_{n\minus 1}\subseteq\Pnt}(\bigwedge_{1{\leq}i{\leq}n\minus 1}\noi(P_i)) \limp \dense(\bigcap_{1{\leq}i{\leq}n\minus 1}P)$.
\end{enumerate*}
\end{prop}
\begin{proof}
Part~2 trivially implies part~1 because we can take $P_1,\dots,P_n\in\opensne$, so that by Definition~\ref{defn.dense.subset} $\noi(P_i)$ for $1\leq i\leq n\minus 1$ holds by construction by Lemma~\ref{lemm.dense.up-closed}(\ref{item.noi.noi}), and we use Lemma~\ref{lemm.is}(\ref{item.is.supertwined}).

Showing that part~1 implies part~2 is only slightly less trivial.
Suppose $(\Pnt,\opens)$ is $n$-twined and $P_1,\dots,P_{n\minus 1}\subseteq\Pnt$ and $\noi(P_i)$ for each $1\leq i\leq n\minus 1$, meaning by Definition~\ref{defn.dense.subset} that for each such $i$ there exists an $O_i\in\opensne$ such that $O_i\subseteq P_i$.
By Lemma~\ref{lemm.is}(\ref{item.is.supertwined}) $O_1\cap\dots\cap O_{n\minus 1}$ is dense, and by Lemma~\ref{lemm.dense.up-closed}(\ref{item.dense.up-closed}) so is $P_1\cap\dots\cap P_n$.
\end{proof}



Proposition~\ref{prop.allbut.supertwined} provides the canonical example of $n$-twined semitopologies.
It mildly generalises known properties in the literature; the case $n=2$ in Proposition~\ref{prop.allbut.supertwined}(\ref{item.allbut.intertwined}) corresponds to the $Q^2$ property from~\cite[Theorem~1]{DBLP:conf/asiacrypt/DamgardDFN07}, and the case $n=3$ corresponds to the $Q^3$ property of~\cite[Definition 2]{DBLP:journals/dc/AlposCTZ24}.
The topology-flavoured presentations in Proposition~\ref{prop.allbut.supertwined}(\ref{item.allbut.NOI.dense}\&\ref{item.allbut.NOI.dense.3}) are new.
\begin{prop}
\label{prop.allbut.supertwined}
Suppose $N,f\in\Ngeqz$ and $P\subseteq\Nset{N}$ and $n\in\mathbb N_{\geq 1}$.
Then:
\begin{enumerate*}
\item\label{item.allbut.intertwined}\label{item.allbut.supertwined}
$N>n*f$ if and only if $\AllBut{N}{f}$ is $n$-twined.
\item\label{item.allbut.NOI.dense}
$\noi(P)$ (and equivalently $P\in\opensne$) if and only if $\mycard P\geq N\minus f$.\footnote{In general semitopologies, $\noi(P)$ ($P$ having a nonempty open interior) is not the same as $P\in\opensne$ ($P$ being a nonempty open set); but in the case of $\AllBut{N}{f}$, these two properties happen to be equal.}
\item\label{item.allbut.NOI.dense.3}
$\dense(P)$ if and only if $\mycard P\geq f\plus 1$.
\end{enumerate*}
\end{prop}
\begin{proof}
We consider each part in turn:
\begin{enumerate}
\item
Suppose $N>n*f$ and suppose $O_1,\dots,O_n$ are open in $\AllBut{N}{f}$.
From Definition~\ref{defn.allbut} this means that they each contain all but $f$ elements of $\Nset{N}$.
It follows that $O_1\cap\dots \cap O_n$ contains (at least) all but $n*f$ elements of $\Nset{N}$.
Since $N\gneq n*f$, this intersection is nonempty. 

Conversely, if $N\leq n*f$ then it is routine to partition $\Nset{N}$ into $n$ sets, each of which has at most $f$ elements.
Then their complements are a set of $n$ open sets whose intersection is empty.
\item
By routine counting arguments.
\item
By routine counting arguments.
\qedhere\end{enumerate}
\end{proof}



\begin{rmrk}
\label{rmrk.ntwined.and.the.literature}
Being \xxtwined corresponds to the \emph{quorum intersection property} of~\cite[Section~2.1, Condition~B2]{DBLP:journals/tocs/Lamport98}, which requires that any two quorums intersect.
Being \xxxtwined also corresponds to a quorum intersection property with a stronger requirement that (assuming there is a quorum of honest participants) any two quorum must intersect in a participant that is honest~\cite[Definition~5.1]{DBLP:journals/dc/MalkhiR98}. 
  
These intersection properties are closely tied to necessary conditions for some consensus protocols (and other distributed algorithms) to function correctly.
Specifically for the case of the semitopology $\AllBut{N}{f}$ from Definition~\ref{defn.allbut}(\ref{item.allbut}), requiring non-empty quorum intersection is equivalent to the condition \( N > 2f \), where \( f \) denotes the number of processes that may crash during the protocol's execution.
This condition is often referred to as \( Q^2 \)~\cite[Theorem~1]{DBLP:conf/asiacrypt/DamgardDFN07}.
Similarly, being \xxxtwined corresponds to the condition \( N > 3f \), commonly referred to as \( Q^3 \)~\cite[Definition 2]{DBLP:journals/dc/AlposCTZ24}.
A more generalised version of \( Q^3 \) exists, known as \( B^3 \)~\cite[Theorem 2]{DBLP:conf/asiacrypt/DamgardDFN07}\cite[Definition 5]{DBLP:journals/dc/AlposCTZ24}, which accounts for greater heterogeneity. 
Specifically, \( Q^3 \) applies to distributed systems where failures are homogeneous, meaning they are globally known and agreed upon by all participants. 
For example, all participants understand that up to \( f \) processes may fail. 
  
In contrast, in a heterogeneous setting, participants may have differing beliefs about \emph{who} might fail, leading to a more nuanced description of trust assumptions. 
In this context, \( B^3 \) generalises \( Q^3 \) by incorporating these variations in trust and failure assumptions.
Readers interested in exploring this topic further can consult~\cite{DBLP:conf/asiacrypt/DamgardDFN07,DBLP:journals/dc/AlposCTZ24}.
\end{rmrk}
  
  
  
  

\subsection{The canonical properties}

For this subsection we fix a signature $\Sigma$ and a model $\mathcal M=(\Pnt,\opens,\Val)$.

The results in this Subsection express how the underlying semitopology of this model $(\Pnt,\opens)$ being \xxtwined or \xxxtwined manifests itself in the logic, in ways that will be relevant in our later proofs:
\begin{lemm}
\label{lemm.supertwined.qcbt}
Suppose $(\Pnt,\opens)$ is a semitopology and $n\in\Time$ and $\avaluation$ is a valuation and suppose $\phi$, $\psi$, and $\chi$ are pointwise in $(n,\avaluation)$ (Definition~\ref{defn.pointwise.meaning}). 
Then: 
\begin{enumerate*}
\item\label{item.intertwined.qcbt}
If $(\Pnt,\opens)$ is \xxtwined then we have:
\begin{enumerate*}
\item\label{item.intertwined.3}
If $n\mentval \QuorumBox\phi$ then $n\mentval\CoquorumDiamond\phi$. 
\item\label{item.intertwined.qcbt.2}
If $n\mentval \QuorumBox\phi \tand \QuorumBox\psi$ then $n\mentval\someoneAll(\phi\tand\psi)$. 
\item\label{item.intertwined.qcbt.3}
If $n\mentval \QuorumBox\correct{\phi} \tand \QuorumBox\phi$ then $n\mentval\modT\someoneAll\phi$.
\end{enumerate*}
\item\label{item.supertwined.qcbt}
If $(\Pnt,\opens)$ is \xxxtwined then we have: 
\begin{enumerate*}
\item\label{item.supertwined.qcbt.3}
If $n\mentval \QuorumBox\phi \tand \QuorumBox\psi$ then $n\mentval \CoquorumDiamond(\phi\tand\psi)$.
\item\label{item.supertwined.qcbt.2}
If $n\mentval \QuorumBox\phi \tand \QuorumBox\psi \tand \QuorumBox\chi$ then $n\mentval\someoneAll(\phi\tand\psi\tand\chi)$. 
\item\label{item.supertwined.qcbt.4}
If $n\mentval \QuorumBox\correct{\phi} \tand \QuorumBox\phi$ then $n\mentval \modT\CoquorumDiamond\phi$.
\end{enumerate*}
\end{enumerate*}
Recall that $\correct{\phi}$ is from Definition~\ref{defn.correct}(\ref{item.correct.phi}) and just means $\modTF\phi$.
\end{lemm}
\begin{proof}
We reason as follows:
\begin{enumerate}
\item
Suppose $(\Pnt,\opens)$ is \xxtwined. 
\begin{enumerate*}
\item
We reason as follows:
$$
\begin{array}{r@{\ }l@{\qquad}l}
n\mentval\QuorumBox\phi
\liff&
\noi(\{p\in\Pnt \mid n,p,\avaluation\ment\phi\})
&\text{Proposition~\ref{prop.pointwise.dense.char}(\ref{item.dense.char.pointwise.QuorumBox})}
\\
\limp&
\dense(\{p\in\Pnt \mid n,p,\avaluation\ment\phi\}) 
&\text{Proposition~\ref{prop.noi.dense}, $(\Pnt,\opens)$ $2$-twined} %
\\
\liff&
n\mentval\CoquorumDiamond\phi
&\text{Proposition~\ref{prop.pointwise.dense.char}(\ref{item.dense.char.pointwise.CoquorumDiamond})}
\end{array}
$$
\item
Suppose $n\mentval \QuorumBox\phi\tand\QuorumBox\psi$.
By part~\ref{item.intertwined.3} of this result, $n\mentval\CoquorumDiamond\phi$.
By Proposition~\ref{prop.someone.implies.someoneAll}(\ref{item.someone.implies.someoneAll.4}) $n\mentval \someoneAll(\phi\tand\psi)$ as required.
\item
Suppose $n\mentval \QuorumBox\correct{\phi} \tand \QuorumBox\phi$.
By part~\ref{item.intertwined.qcbt.2} of this result, $n\mentval\someoneAll(\correct{\phi}\tand\phi)$.
Using %
Proposition~\ref{prop.tv.ment.TF.model}(\ref{item.tv.ment.TF.phi.lr}) $n\mentval\someoneAll\modT\phi$.
By Lemma~\ref{lemm.someone.commutation} $n\mentval\modT\someoneAll\phi$ as required. 
\end{enumerate*}
\item
The reasoning for the \xxxtwined case is just as for the \xxtwined case, except we need intersection of three open sets instead of two.
\qedhere\end{enumerate} 
\end{proof}




Corollary~\ref{corr.two.quorums.intersect.correct} specialises Lemma~\ref{lemm.supertwined.qcbt} to the useful special case of a correct unary predicate symbol, such that $\phi=\correct{\tf P}$ (Definition~\ref{defn.correct}(\ref{item.correct})) and $\psi=\tf P(v)$ and $\chi=\tf P(v')$.
In our later examples, $\tf P$ will take values such as $\tf{accept}$, $\tf{send}$, $\tf{write}$, and $\tf{decide}$, representing basic actions in a consensus protocol:
\begin{corr}
\label{corr.two.quorums.intersect.correct}
Suppose $(\Pnt,\opens)$ is a semitopology and suppose $\tf P$ is a unary predicate symbol and $v,v'\in\Val$ and $n\in\Time$ and $\avaluation$ is a valuation.
Then:
\begin{enumerate*}
\item
If $(\Pnt,\opens)$ is \xxxtwined then:
\begin{enumerate*}
\item\label{item.two.quorums.intersect.correct.supertwined}
If $n\mentval\QuorumBox\correct{\tf P}\tand\QuorumBox\tf P(v)\tand\QuorumBox\tf P(v')$ then $n\mentval\modT\someoneAll(\tf P(v)\tand\tf P(v'))$.
\item\label{item.QuorumBox.correct.supertwined}
If $n\mentval\QuorumBox\correct{\tf P}\tand\QuorumBox\tf P(v)$ then $n\mentval\modT\CoquorumDiamond\tf P(v)$.
\end{enumerate*}
\item\label{item.two.quorums.intersect.correct.intertwined}
If $(\Pnt,\opens)$ is \xxtwined then
$n\mentval\QuorumBox\correct{\tf P}\tand\QuorumBox\tf P(v)$ implies $n\mentval\modT\someoneAll(\tf P(v))$.
\end{enumerate*}
\end{corr}
\begin{proof}
We consider each part in turn:
\begin{enumerate}
\item[(\ref{item.two.quorums.intersect.correct.supertwined})]
Suppose $(\Pnt,\opens)$ is \xxxtwined and $n\mentval\QuorumBox\correct{\tf P}\tand\QuorumBox\tf P(v)\tand\QuorumBox\tf P(v')$.
By Lemma~\ref{lemm.supertwined.qcbt}(\ref{item.supertwined.qcbt.2}) (since $\tf P(v)$ and $\correct{\tf P}$ are pointwise by Lemma~\ref{lemm.atomic.O}(\ref{item.atomic.pointwise.2})) $n\mentval\someoneAll(\correct{\tf P}\tand\tf P(v)\tand\tf P(v'))$.
By two applications of Proposition~\ref{prop.tv.ment.TF.model}(\ref{item.tv.ment.TF.P}) and Lemma~\ref{lemm.commuting.connectives} (for $\tand$)
$n\mentval\someoneAll\modT(\tf P(v)\tand\tf P(v'))$.
By Lemma~\ref{lemm.someone.commutation} $n\mentval\modT\someoneAll(\tf P(v)\tand\tf P(v'))$ as required.
\item[(\ref{item.QuorumBox.correct.supertwined})] 
Suppose $(\Pnt,\opens)$ is \xxxtwined and $n\mentval\QuorumBox\correct{\tf P}\tand\QuorumBox\tf P(v)$.
Using Lemma~\ref{lemm.correct.P.Pv} it follows that $n\mentval\QuorumBox\correct{\tf P(v)}\tand\QuorumBox\tf P(v)$.
We just apply Lemma~\ref{lemm.supertwined.qcbt}(\ref{item.supertwined.qcbt.4}). 
\item[(\ref{item.two.quorums.intersect.correct.intertwined})]
This follows from Lemma~\ref{lemm.supertwined.qcbt}(\ref{item.intertwined.qcbt.3}) just as part~\ref{item.QuorumBox.correct.supertwined} does from Lemma~\ref{lemm.supertwined.qcbt}(\ref{item.supertwined.qcbt.2}).
\qedhere\end{enumerate}
\end{proof}


Lemma~\ref{lemm.allbut.concretely} unpacks concretely what $\QuorumBox\phi$ and $\CoquorumDiamond\phi$ --- which are the subject of many results including of Lemmas~\ref{lemm.supertwined.qcbt} and~\ref{lemm.intertwined.char} --- mean in the semitopology $\AllBut{N}{f}$ from Definition~\ref{defn.allbut}(\ref{item.allbut}), which has $N$ points, and has open sets being either the empty set, or any set with at least $N\minus f$ elements: 
\begin{enumerate*}
\item
$\QuorumBox\phi$ means `for at least $N\minus f$ elements'; and 
\item
$\CoquorumDiamond\phi$ means `for at least $f\plus 1$ elements.
\end{enumerate*}
This special case is particularly relevant because (as the reader may know) $\AllBut{N}{f}$ is commonly used in the literature on distributed algorithms. 
\begin{lemm}
\label{lemm.allbut.concretely}
Suppose that
\begin{itemize*}
\item
$N,f\in\Ngeqz$, 
\item
$(\Pnt,\opens)=\AllBut{N}{f}$, 
\item
$n\in\Time$ and $\avaluation$ is a valuation, and
\item
$\phi$ is a closed predicate that is pointwise in $(n,\avaluation)$.
\end{itemize*}
Then:
\begin{enumerate*}
\item
$n\mentval\QuorumBox\phi$ if and only if 
$\mycard\{p\in\Pnt \mid n,p\mentval\phi\}\geq N\minus f$. 
\item
$n\mentval\CoquorumDiamond\phi$ if and only if 
$\mycard\{p\in\Pnt \mid n,p\mentval\phi\}\geq f\plus 1$. 
\item
$n\mentval\modT\QuorumBox\phi$ if and only if 
$\mycard\{p\in\Pnt \mid n,p\mentval\modT\phi\}\geq N\minus f$. 
\item
$n\mentval\modT\CoquorumDiamond\phi$ if and only if 
$\mycard\{p\in\Pnt \mid n,p\mentval\modT\phi\}\geq f\plus 1$. 
\end{enumerate*}
\end{lemm}
\begin{proof}
We consider each part in turn:
\begin{enumerate}
\item
The proof is elementary but we spell it out.
We reason as follows:
$$
\begin{array}{r@{\ }l@{\quad}l}
n\mentval\QuorumBox\phi
\liff&
\noi(\{p\in\Pnt \mid n,p\mentval \modT\phi\})
&\text{Proposition~\ref{prop.pointwise.dense.char}(\ref{item.dense.char.pointwise.QuorumBox})}
\\
\liff&
\mycard\{p\in\Pnt \mid n,p\mentval \phi\}\geq N\minus f
&\text{Proposition~\ref{prop.allbut.supertwined}(\ref{item.allbut.NOI.dense})}
\end{array}
$$
\item
Much as the previous part, but using Proposition~\ref{prop.pointwise.dense.char}(\ref{item.dense.char.pointwise.CoquorumDiamond}).
\item
By Lemma~\ref{lemm.someone.commutation}, $n\mentval\modT\QuorumBox\phi$ if and only if $n\mentval\QuorumBox\modT\phi$.
We just use part~1 of this result.
\item
Much as the previous part, using part~2 of this result.
\qedhere\end{enumerate}
\end{proof}

\subsection{Further observations on intersection properties}

For this Subsection we fix a signature $\Sigma$ and a model $\mathcal M=(\Pnt,\opens,\Val)$.
\begin{rmrk}
We collect some further results as follows:
\begin{enumerate*}
\item
Lemma~\ref{lemm.intertwined.char} proves converses to the implications in Lemma~\ref{lemm.supertwined.qcbt}, provided that the signature is nonempty.

This is interesting because it tells us that, provided the signature is nondegenerate, Lemma~\ref{lemm.supertwined.qcbt} accurately captures and reflects the intersection properties that interest us. 
\item
Proposition~\ref{prop.reformulate} is a semitopological distillation of parts~\ref{item.intertwined.qcbt.2} and~\ref{item.supertwined.qcbt.2} of Lemma~\ref{lemm.supertwined.qcbt}.
Proposition~\ref{prop.reformulate} \emph{is} Lemma~\ref{lemm.supertwined.qcbt} and Corollary~\ref{corr.two.quorums.intersect.correct}, in semitopological rather than logical form.
\end{enumerate*}
\end{rmrk}

In Lemma~\ref{lemm.intertwined.char} we could replace every instance of the expression `$\modTF\tf P$' with the (by Definition~\ref{defn.correct}(\ref{item.correct.phi}) and Remark~\ref{rmrk.tf.synonyms}) synonymous expression `$\correct{\tf P}$'.
However, we have not set up a corresponding synonym for $\modTB\tf P$ (it would be `$\tf{valid}[\tf P]$'; we just never wrote that definition) so it seems natural to stay with $\modTF$ here.
\begin{lemm}
\label{lemm.intertwined.char}
Suppose $\tf P\in\ns{PredSymb}$ is a predicate constant symbol (Notation~\ref{nttn.logic.terminology}(\ref{item.nullary.ps}); a nullary predicate symbol).\footnote{The proof works just as well for an $n$-ary predicate symbol; we just pick \emph{any} $\arity(\tf P)$-tuple of values $\vec v\in\Pnt^{\arity(\tf P)}$ --- which is possible because we assumed $\Val$ is nonempty in Definition~\ref{defn.model}(\ref{item.model.val}) --- and we use $\tf P(\vec v)$ in the proof instead.}
Then:
\begin{enumerate*}
\item
If\ \  $\ment(\QuorumBox\modTF\tf P\tand\QuorumBox\modTB\tf P)\timpc \someoneAll\tf P$ then $(\Pnt,\opens)$ is \xxtwined. 
\\
(This is a form of converse to Lemma~\ref{lemm.supertwined.qcbt}(\ref{item.intertwined.qcbt.3}).)
\item
If\ \  $\ment(\QuorumBox\modTF\tf P\tand\QuorumBox\modTB\tf P)\timpc \CoquorumDiamond\tf P$ then $(\Pnt,\opensne)$ is \xxxtwined. 
\\
(This is a form of converse to Lemma~\ref{lemm.supertwined.qcbt}(\ref{item.supertwined.qcbt.4}).) 
\item\label{item.intertwined.char.reverse}
As a corollary, provided that the signature $\Sigma$ is has at least one predicate symbol --- i.e. provided $\PredSymb$ in
it is nonempty (Definition~\ref{defn.signature}(\ref{item.predicate.syntax.empty.signature})) --- then reverse implications to those in Lemma~\ref{lemm.supertwined.qcbt} hold.\footnote{No need to also insist that $\Val$ is nonempty, because we assume this in Definition~\ref{defn.model}(\ref{item.model.val}).}
\end{enumerate*}
Recall from Definition~\ref{defn.validity.judgement}(\ref{item.ment.phi}) that $\ment\phi$ means $\Forall{\acontext\in\tf{Ctx}(\Sigma,\mathcal M)}(\acontext\ment\phi)$.
\end{lemm}
\begin{proof}
We consider each part in turn:
\begin{enumerate}
\item
We prove a contrapositive of \strongmodusponens: we suppose $\opens$ is \emph{not} \xxtwined, so there exist $O,O'\in\opensne$ with $O\cap O'=\varnothing$, and we exhibit a valuation such that $\neg(\mentval\someone\tf P)$ (equivalently: $\mentval\modF\someone\tf P$).

We choose a valuation $\avaluation$ such that for any $n\in\Time$, 
\begin{itemize*}
\item
$\avaluation(\tf P)_{n,p}(\aval)=\tvF$ for $p\in O$, and 
\item
$\avaluation(\tf P)_{n,p}(\aval)=\tvB$ for $p\in \Pnt\setminus O$.
\end{itemize*}
The reader can check that $O\mentval \everyone\modTF\tf P$ and $O'\mentval \everyone\modTB P(\aval)$, but $\mentval\modF\someone\tf P$.
\item
We prove a contrapositive: suppose $\opens$ is \emph{not} \xxxtwined, so there exist $O,O',O''\in\opensne$ with $O\cap O'\cap O''=\varnothing$.

We choose a valuation $\avaluation$ such that for any $n\in\Time$, 
\begin{itemize*}
\item
$\avaluation(\tf P)_{n,p}(\aval)=\tvT$ for $p\in O$, and 
\item
$\avaluation(\tf P)_{n,p}(\aval)=\tvB$ for $p\in O'\setminus O$, and 
\item
$\avaluation(\tf P)_{n,p}(\aval)=\tvF$ for $p\in \Pnt\setminus (O\cup O')$.
\end{itemize*}
The reader can check that $O\cup O'\mentval \everyone\modTF\tf P$ and $O\cup O''\mentval \everyone\modTB P(\aval)$, but $O'\cup O''\mentval\modF\everyone\tf P$.
\item
Just by the arguments above.
\qedhere\end{enumerate}
\end{proof}






\begin{nttn}
\label{nttn.cover}
Suppose $\ns X$ is a set.
Then:
\begin{enumerate*}
\item
A \deffont{cover} of $\ns X$ is a set of sets $\mathcal X\subseteq\powerset(\ns X)$ such that $\bigcup\mathcal X=\ns X$.
\item
A \deffont{partition} of $\ns X$ is a cover of $\ns X$ all of whose elements are pairwise disjoint (so $\Forall{X,X'\in\mathcal X}X\cap X'=\varnothing\lor X=X'$).
\end{enumerate*}
\end{nttn}

The reader might like to try giving a proof of Proposition~\ref{prop.reformulate} using direct calculations (without building on the results we have so far).
This is possible, but it is messy.
The use of logic endows the proofs with a structure and cleanness that the direct proof cannot match: 
\begin{prop}
\label{prop.reformulate}
Suppose $(\Pnt,\opens)$ is a semitopology and $n\in\Time$ and $\avaluation$ is a valuation, and suppose $\phi$ is a closed predicate that is pointwise in $(n,\avaluation)$.
Then:
\begin{enumerate*}
\item\label{item.reformulate.1}
If $(\Pnt,\opens)$ is \xxtwined then $n\mentval\CoquorumDiamond\modTB\phi \tor \CoquorumDiamond\modF\phi$. 
\item\label{item.reformulate.1b}
As a corollary, if $(\Pnt,\opens)$ is \xxtwined then any two-element partition of $\Pnt$ (Notation~\ref{nttn.cover}) contains a dense element (Definition~\ref{defn.dense.subset}(\ref{item.dense.subset})).
\item\label{item.reformulate.2}
If $(\Pnt,\opens)$ is \xxxtwined then 
$n\mentval\CoquorumDiamond\modT\phi \tor \CoquorumDiamond\modB\phi \tor \CoquorumDiamond\modF\phi$. 
\item\label{item.reformulate.2b}
As a corollary, if $(\Pnt,\opens)$ is \xxxtwined then any three-element partition of $\Pnt$ contains a dense element.
\end{enumerate*} 
\end{prop}
\begin{proof}
We consider each part in turn:
\begin{enumerate}
\item
By Lemma~\ref{lemm.supertwined.qcbt}(\ref{item.intertwined.qcbt.2})
$n\mentval\QuorumBox\modF\phi\tand\QuorumBox\modTB\phi$ implies $n\mentval\someoneAll(\modF\phi\tand\modTB\phi)$.
Unpacking what $n\mentval\someoneAll(\modF\phi\tand\modTB\phi)$ means using Definition~\ref{defn.validity.judgement} and Figures~\ref{fig.3.derived} and~\ref{fig.3.phi.f}, we see that this is impossible --- for no $p$ and $O$ can we satisfy both $n,p,O\mentval\modF\phi$ and $n,p,O\mentval\modTB\phi$ together --- so that 
$n\nmentval\QuorumBox\modF\phi\tand\QuorumBox\modTB\phi$.
By Lemma~\ref{lemm.para}(\ref{item.para.em})
$n\mentval\tneg(\QuorumBox\modF\phi\tand\QuorumBox\modTB\phi)$, and simplifying further using de Morgan equivalences as per Figure~\ref{fig.easy.equivalences}, we obtain
$n\mentval\CoquorumDiamond\modTB\phi\tor\CoquorumDiamond\modF\phi$ as required.
\item
Suppose $\{P_1,P_2\}$ partitions $\Pnt$.
Assume a $0$-ary predicate symbol $\tf P$, and let $\varsigma$ be the valuation that sets $\varsigma(\tf P)_{n,p}=\tvT$ if $p\in P_1$ and $\varsigma(\tf P)_{n,p}=\tvF$ if $p\in P_2$.\footnote{We can assume $\tf P$ and $\varsigma$ because this result is about semitopologies, so we are free to impose any model structure on it that we like.  It seems particularly nice here how the mere \emph{existence} of a model for our logic gives us a clean proof of a purely semitopological property which, on the face of it, has nothing to do with our three-valued modal logic.}
By part~\ref{item.reformulate.1} and Lemma~\ref{lemm.tand.tor} either $n\mentval \Coquorum\modTB\tf P$ or $n\mentval\Coquorum\modF\tf P$.
By Proposition~\ref{prop.pointwise.dense.char}(\ref{item.dense.char.pointwise.CoquorumDiamond}) either $P_1$ or $P_2$ is dense.
\item
We reason as for the proof of part~\ref{item.reformulate.1}, using Lemma~\ref{lemm.supertwined.qcbt}(\ref{item.supertwined.qcbt.2}) and starting from $\modFB\phi \tand \modTF\phi \tand \modTB\phi$. 
\item
We reason as for the proof of part~\ref{item.reformulate.2}, for a partition $\{P_1,P_2,P_3\}$ and setting $\varsigma(\tf P)_{n,p}=\tvT$ if $p\in P_1$ and $\varsigma(\tf P)_{n,p}=\tvB$ if $p\in P_2$ and $\varsigma(\tf P)_{n,p}=\tvF$ if $p\in P_3$.
\qedhere\end{enumerate}
\end{proof}




 





\section{Declarative Paxos}
\label{sect.declarative.paxos}

\subsection{Case study: a simple protocol}
\label{subsect.simple}

\begin{rmrk}
\label{rmrk.simple.protocol}
Before Paxos, we will warm up by sketching a simple protocol \theory{Simple}, and then we will give four \QLogic axiomatisations of \theory{Simple}, reflecting four different failure assumptions (messages can/cannot get lost; participants can/cannot crash).
\theory{Simple} is as follows:
\begin{enumerate*}
\item
Any participant $p$ may broadcast a $\tf{propose}$ message to all participants.
\item
If a participant receives a $\tf{propose}$ message, then it responds with $\tf{accept}$. 
\item
If the participant $p$, having broadcast its $\tf{propose}$ message, receives a $\tf{accept}$ message, then it declares $\tf{decide}$.
\end{enumerate*}
\end{rmrk}


\begin{rmrk}
\theory{Simple} is very simple!
Only one value can be proposed, so we do not even bother to state it; we just `propose'.
Any participant accepts immediately that it hears someone propose. 
Any participant that proposes, decides immediately that it hears that someone accepted.

For comparison, in Paxos \emph{multiple values} are available, so we need to make sure that participants all decide on the \emph{same} value if they decide on anything at all. 
This makes things much harder: so Paxos has multiple rounds; there is a designated leader participant at each round; the leader must choose a value to propose; and elaborate communication is in place to ensure that (subject to suitable precisely-stated assumptions) a single choice of agreed value will eventually be decided.

But that is for later.  
Even the simple protocol \theory{Simple} has plenty of subtlety, which will allow us to exhibit (in `toy' form) many of the techniques that we will then apply to Paxos:
\end{rmrk}


\begin{rmrk}
We will axiomatise \theory{Simple} subject to two parameters, yielding four possible axiomatisations:
\begin{enumerate*}
\item
Do all messages arrive, or might there be a finite initial period during which the network may be unreliable so that during that time, messages might not arrive?
\item
Can participants crash (and if so, how many)?
\end{enumerate*}
\end{rmrk}

\subsubsection*{All messages arrive; no participants crash}

First, we assume an ideal situation that all messages arrive, and no participants crash.
We can axiomatise this as follows:
$$
\begin{array}{l@{\qquad}l}
\rulefont{SimpAccept?}&
\tf{accept} \timpc \someoneAll\tf{propose} 
\\
\rulefont{SimpDecide?}&
\tf{decide} \timpc (\tf{propose}\tand\someoneAll\tf{accept}) 
\\
\rulefont{SimpAccept!}&
(\someoneAll\tf{propose}) \timpc \tf{accept}
\\
\rulefont{SimpDecide!}&
(\tf{propose}\tand\someoneAll\tf{accept}) \timpc \tf{decide}
\\
\rulefont{Correct}&
\correct{\tf{propose},\tf{accept},\tf{decide}}
\end{array}
$$
\noindent Recall from Definition~\ref{defn.three.imp} and Proposition~\ref{prop.mp.for.tnotor} that $\timpc$ is a \emph{strong implication} and means `if the LHS is $\tvT$ then so is the RHS', and 
recall from Proposition~\ref{prop.pointwise.dense.char}(\ref{item.dense.char.pointwise.someoneAll}) that $\someoneAll\tf P$ means `$\tf P$ somewhere', where $\tf P\in\{\tf{propose},\tf{accept},\tf{decide}\}$.

Then we explain the axioms (briefly) as follows:
\begin{itemize*}
\item
We call \rulefont{SimpAccept?} and \rulefont{SimpDecide?} \emph{backward rules}; they express backward versions of the rules in Remark~\ref{rmrk.simple.protocol}, namely that:
\begin{itemize*}
\item
if a participant asserts $\tf{accept}$ then some participant must have asserted $\tf{propose}$; and 
\item
if a participant asserts $\tf{decide}$ then it must have asserted $\tf{propose}$, and some participant must have asserted $\tf{accept}$.
\end{itemize*}
\item
We call \rulefont{SimpAccept!} and \rulefont{SimpDecide!} \emph{forward rules}.
They express forward versions of the rules in Remark~\ref{rmrk.simple.protocol}.
Because our failure assumptions are that nothing fails --- all messages arrive; no participants crash --- the forward rules are in this case just the backward rules, reversed.
\item
Axiom \rulefont{Correct} expresses that all predicates return $\tvT$ and $\tvF$ (not $\tvB$), as per Definition~\ref{defn.correct}(\ref{item.correct}). 
This reflects our assumption that no participants crash, because we model a computation crashing %
as the predicate describing it returning $\tvB$.
\end{itemize*}
Note that $\tf{propose}$, $\tf{accept}$, and $\tf{decide}$ above are not messages; they are propositions!
There is no message-passing.
$\tf{accept}$ being true at a participant represents that the participant accepts; this is implemented in the algorithm by the participant broadcasting an accept message, and represented in the logic by $\tf{accept}$ taking truth-value $\tvT$ at that participant.

Our example protocol is simple, but we hope it is already starting to illustrate how the various parts of our logic come together to express the protocol axiomatically.
We continue:

\subsubsection*{All messages arrive; participants can crash}

We assume that all messages arrive, and participants can crash --- but we assume there is a dense set of uncrashed participants.
Recall from Definition~\ref{defn.dense.subset}(\ref{item.dense.subset}) that a \emph{dense set} is one that intersects every open set; in particular, this means that every open set of participants contains a participant that is not crashed.
We model being crashed as returning the third truth-value $\tvB$.
We can axiomatise this as follows:
$$
\begin{array}{l@{\qquad}l}
\rulefont{SimpAccept?}&
\tf{accept} \timpc \someoneAll\tf{propose} 
\\
\rulefont{SimpDecide?}&
\tf{decide} \timpc (\tf{propose}\tand\someoneAll\tf{accept}) 
\\
\rulefont{SimpAccept!}&
(\someoneAll\tf{propose}) \tnotor \tf{accept}
\\
\rulefont{SimpDecide!}&
(\tf{propose}\tand\someoneAll\tf{accept}) \tnotor \tf{decide}
\\
\rulefont{Uncrashed}&
\QuorumBox\modTB\tf P \tnotor \someoneAll\modT\tf P
\quad \tf P\in\{\tf{propose},\tf{accept},\tf{decide}\}
\end{array}
$$ 
\noindent 
The backward rules are unchanged: if a participant accepts then from the form of the rules, it \emph{must} be the case that somebody proposed; and if a participant decides then from the form of the rules, it \emph{must} be the case that the participant proposed and some participant accepted.
 
The forward rules \rulefont{SimpAccept!} and \rulefont{SimpDecide!} have changed, and they now use the weak implication $\tnotor$.
By Definition~\ref{defn.three.imp} and Proposition~\ref{prop.mp.for.tnotor}, $\tnotor$ means `if the LHS is $\tvT$ then the RHS is $\tvT$ or $\tvB$'.
The $\tvB$ truth-value represents `crashed', so that (for example) \rulefont{SimpAccept!} expresses that for each participant $p$, if some (possibly other) participant $q$ asserts $\tf{propose}$, then $p$ asserts $\tf{accept}$ \emph{or} $p$ has crashed. 

Axiom \rulefont{Uncrashed} reflects that any a quorum (open set) of participants contains an uncrashed participant; e.g. in topological language, that the set of uncrashed participants is dense.
More precisely: \rulefont{Uncrashed} asserts that if an open set of participants returns $\tvT$ or $\tvB$ for $\tf P$, then at least one of them (is not crashed and so) returns $\tvT$.

Note from Lemma~\ref{lemm.someone.commutation} that $\Quorum\modTB$ is equivalent to $\modTB\Quorum$, and $\someoneAll\modT$ is equivalent to $\modT\someoneAll$, so that an equivalent form of \rulefont{Uncrashed} is $\modTB\QuorumBox\tf P \tnotor \modT\someoneAll\tf P$.
In our axiomatisations it is usually helpful to pull $\modTB$ and $\modT$ up to the topmost level possible in the predicate, so we may favour that form, but this is just a matter of convenience; meaning is unchanged. 


\subsubsection*{Messages might not arrive; no participants crash}

We assume that there is an initial period during which the network might not deliver messages, but the network eventually stabilises (in the literature this is called \emph{eventual synchrony}).
No participants crash.
We can axiomatise this as follows:
$$
\begin{array}{l@{\qquad}l}
\rulefont{SimpAccept?}&
\tf{accept} \timpc \someoneAll\tf{propose} 
\\
\rulefont{SimpDecide?}&
\tf{decide} \timpc (\tf{propose}\tand\QuorumBox\tf{accept}) 
\\
\rulefont{SimpAccept!}&
\final((\someoneAll\tf{propose}) \timpc \tf{accept})
\\
\rulefont{SimpDecide!}&
\final((\tf{propose}\tand\someoneAll\tf{accept}) \timpc \tf{decide})
\\
\rulefont{Correct}&
\correct{\tf{propose},\tf{accept},\tf{decide}}
\end{array}
$$ 
\noindent The backward rules are unchanged, but the forward rules \rulefont{SimpAccept!} and \rulefont{SimpDecide!} now have a $\final$ modality.
This is from Figure~\ref{fig.3.derived} and it means `eventually, this holds forever', reflecting an assumption that eventually the network stabilises.

Note that \QLogic does not explicitly represent message-passing, but this does not matter as much as one might suppose!
The $\final$ modality captures the essence of what we need, namely, that eventually the forward rules become valid. 

\subsubsection*{Messages might not arrive; participants can crash}

We assume that messages might not arrive, and participants can crash but there is a dense set of uncrashed participants (so that every quorum of participants contains at least one uncrashed participant).
We can axiomatise this just by combining the ideas above, as follows:
$$
\begin{array}{l@{\qquad}l}
\rulefont{SimpAccept?}&
\tf{accept} \timpc \someoneAll\tf{propose} 
\\
\rulefont{SimpDecide?}&
\tf{decide} \timpc (\tf{propose}\tand\someoneAll\tf{accept}) 
\\
\rulefont{SimpAccept!}&
\final((\someoneAll\tf{propose}) \tnotor \tf{accept})
\\
\rulefont{SimpDecide!}&
\final((\tf{propose}\tand\someoneAll\tf{accept}) \tnotor \tf{decide})
\\
\rulefont{Uncrashed}&
\modTB\QuorumBox\tf P \tnotor \modT\someoneAll\tf P
\quad \tf P\in\{\tf{propose},\tf{accept},\tf{decide}\}
\end{array}
$$ 
\noindent The axiomatisations of Declarative Paxos and Simpler Declarative Paxos in Figures~\ref{fig.ThyPaxOne} and~\ref{fig.ThySPax} are more elaborate than the axioms above, because Paxos is a more complicated protocol, but the basic ideas are as above; just scaled up. 


\subsection{A high-level view of the Paxos algorithm}



\begin{rmrk}[High-level view of Paxos]
\label{rmrk.high-level.paxos}
Paxos allows participants to crash, and messages to fail to arrive at their destination, but it assumes all participants are honest (i.e. that if they have not crashed then they follow the rules of the algorithm).
The reader can find detailed treatments of Paxos in \cite{DBLP:journals/tocs/Lamport98,lamport2001paxos,cachin2011paxos} (a Wikipedia article~\cite{wiki:Paxos_(computer_science)} is also accessible and clear).
Here is a high-level view of the algorithm:
\begin{enumerate*}
\item\label{item.l.proposes}
\deffont{Leader proposes:}\quad

A designated leader participant (if it does not crash) sends messages to all participants, \emph{proposing} some unique value $v$.

Message-passing is unreliable, so messages may fail to arrive.
\item\label{item.p.sends}
\deffont{Participant sends:}\quad

If a participant $p$ receives a propose $v$ message from the leader --- it is assumed that every participant knows who the leader is, and other participants cannot forge the leader's signature --- then $p$ sends a \emph{send} message back to the leader, which carries the value $v_p$ that $p$ most recently accepted, along with the timestamp of when $v_p$ was accepted.\footnote{For reasons of compact presentation we do not explicitly represent the timestamps in our axiomatisation; i.e. $\tf{send}$ is a unary predicate in Definition~\ref{defn.declarative.paxos}(\ref{item.ThyPaxOne.sig}), not a binary predicate.
If this bothers the reader then we can just say that, actually, $v_p$ contains a timestamp of the stage $n\in\Time$ it was accepted; i.e. that $v_p$ is actually implemented as a pair $(v_p',n_p)$ of a value and a time.} 
If $p$ has never accepted any value, so no most recently accepted value exists, then $p$ sends back a special $\udfn$ value.
\item\label{item.l.writes}
\deffont{Leader writes:}\quad

The leader participant sends a message to all participants, \emph{writing} the most recently-timestamped value $v_p$ mentioned in any \emph{send} message it received from within the first quorum of participants that it hears back from; if the leader does not hear back from any quorum, then it does not write anything.
\item\label{item.p.accepts}
\deffont{Participant accepts:}\quad

If a participant receives a write message from the leader for some value, then it sends a message back \emph{accepting} that value.
\item\label{item.l.decides}
\deffont{Leader decides:}\quad

If the leader receives accept messages for some value from a quorum of participants, it decides on that value, and sends out a \emph{decide} message to all participants with that value.
\item\label{item.p.decides}
\deffont{Participant (non-leader) decides:}\quad

If a participant receives a decide message from the leader for some value, then it decides on that value and stops.
\end{enumerate*}
\end{rmrk}

\begin{rmrk}[High-level view of our axiomatisation]
\label{rmrk.elisions}
The theory $\ThyPax$ in Definition~\ref{defn.ThyPaxOne} and Figure~\ref{fig.ThyPaxOne} is a declarative abstraction of Paxos-the-algorithm.
It stands in relation to implementations of Paxos in roughly the same way that (say) the mathematical specification of the Fibonacci function stands in relation to its concrete implementations, in the sense that we elide implementational detail to preserve (and reveal) a mathematical ideal of what the implementation does.

\ThyPax elides two features of the concrete implementation:
\begin{enumerate*}
\item
There is no explicit message-passing; there are only predicates returning truth-values.

This may seem counterintuitive --- isn't Paxos \emph{supposed} to be a message-passing algorithm? --- but we shall see that a lot of the character of Paxos survives.
Indeed, arguably (and certainly for the first author) it is much easier to see what is really going on once we have done this.
\item\label{item.elide.algorithmic.time}
There is no abstract machine, and therefore no \deffont{algorithmic time}, by which we mean no sequence of state-changes of the abstract machine. 

The only notion of time is the \emph{stage} $n\in\Time$ from the context as defined in Definition~\ref{defn.contexts}(\ref{item.context.stage}).
We may call this \deffont{logical time} and (when we are modelling Paxos) it represents the sequence of what is in the distributed consensus algorithms literature often called \emph{stages}, or \emph{ballot numbers}.\footnote{In a distributed algorithm, at any point in algorithmic time, different participants in the algorithm may be at different points in logical time.  A practical example: consider several participants racing across 100 metres. In logical time, they are both just counting to 100.  In algorithmic time, one of them will get to 100 \emph{first}, and win.  In this paper, we only care about the logical aspects of the process; there may still be a winner participant, but this would implemented just as a unary predicate $\tf{winner}$ such that $100\ment\texiaffine\tf{winner}$.} 
\end{enumerate*}
We will prove standard Paxos correctness properties for \ThyPax in Section~\ref{sect.paxos.correctness.properties}. 
\end{rmrk}


\subsection{Definition of Declarative Paxos}

\begin{figure}
$$
\begin{array}{l@{\quad}r@{\ }l}
\figunderline{Backward rules, for backward inference} 
\figskip
\rulefont{PaxPropose?}&
& \tall v.\ \tf{propose}(v)\timpc(\tf{leader}\tand v\tneq\udfn)
\figskip
\rulefont{PaxSend?}&
& \tall v.\ \tf{send}(v) \timpc 
\begin{array}[t]{l}
\someoneAll\texi\tf{propose}\ \tand 
\\
\bigl(
(\mru{\tf{accept}}{v})\ \tor  
(\tneg\modT\recent\texi\tf{accept}\tand v\teq\udfn)\bigr)
\end{array}
\figskip
\rulefont{PaxWrite?}&
& \tall v.\ \tf{write}(v) \timpc 
\begin{array}[t]{l}
\tf{leader}\tand 
\\
\Quorum(\Box\texi\tf{send} \tand 
\\
\qquad\bigl((\mru{(\someone\tf{accept})}{v}) \tor
(\tf{propose}(v)\tand\tneg\modT\recent\texi\someone\tf{accept})\bigr)) 
\end{array}
\figskip
\rulefont{PaxAccept?}&
& \tall v.\ \tf{accept}(v) \timpc \someoneAll\tf{write}(v)
\figskip
\rulefont{PaxDecideL?}&
& \tf{leader} \timpc \tf{decide}(v) \timpc \QuorumBox\tf{accept}(v)
\figskip
\rulefont{PaxDecide\tneg L?}&
& \tall v.\ \tneg\tf{leader} \timpc \tf{decide}(v) \timpc \someoneAll(\tf{leader}\tand\tf{decide}(v))
\figskip\figskip
\figunderline{Forward rules, for making progress} 
\figskip
\rulefont{PaxPropose!}&
& \tf{leader}\tnotor\texi\tf{propose}
\figskip
\rulefont{PaxSend!}&
& \final(\someoneAll\texi\tf{propose} \tnotor \texi\tf{send}) 
\figskip
\rulefont{PaxWrite!}&
& \final\bigl((\tf{leader}\tand (\QuorumBox\texi\tf{send}) \tand \texi\tf{propose}) \tnotor \texi\tf{write}\bigr)
\figskip
\rulefont{PaxAccept!}&
& \final(\texi\someoneAll\tf{write}\tnotor \texi\tf{accept}) 
\figskip
\rulefont{PaxDecideL!}&
& \final((\tf{leader} \tand\texi\QuorumBox\tf{accept}) \tnotor \texi\tf{decide}) 
\figskip
\rulefont{PaxDecide\tneg L!}&
& \final((\tneg\tf{leader} \tand \someoneAll(\tf{leader}\tand\texi\tf{decide}))\tnotor \texi\tf{decide}) 
\figskip
\figskip
\figunderline{Other rules, for correctness \& liveness}
\figskip
\rulefont{LdrExist}&
&\modTF\tf{leader}\tand\modT\someoneAll\tf{leader}
\figskip
\rulefont{LdrExt}_\phi&
& \tf{leader}\timpc\someoneAll(\tf{leader}\tand\phi)\timpc\phi
\quad\text{every closed predicate $\phi$}
\figskip
\rulefont{LdrCorrect}&
& \infinitely(\tf{leader}\timpc\correct{\tf{propose},\tf{write},\tf{decide}}) 
\figskip
\rulefont{PaxPCorrect}&
& 
\begin{array}[t]{l}
\QuorumBox\correct{\tf{propose}} \tand \QuorumBox\correct{\tf{send}} \tand \QuorumBox\correct{\tf{accept}} \tand 
\\
\qquad\QuorumBox\correct{\tf{write}} \tand \QuorumBox\correct{\tf{decide}}
\end{array}
\figskip
\rulefont{PaxPropose01}&
& \texiaffine \tf{propose}
\figskip
\rulefont{PaxWrite01}&
& \texiaffine\tf{write}
\figskip
\rulefont{Pax2Twined}&
& \tall v.\ \QuorumBox\tf{accept}(v) \timpc \CoquorumDiamond\tf{accept}(v)
\end{array}
$$
\caption{$\Theta_\ThyPax$: Axioms of Declarative Paxos (Definition~\ref{defn.declarative.paxos})}
\label{fig.logical.paxos}
\label{fig.ThyPaxOne}
\end{figure}


Recall from Definition~\ref{defn.theory.axiom}(\ref{item.theory}) that a \emph{theory} is a pair of a \emph{signature} and a set of \emph{axioms} over that signature.
Then:
\begin{defn}
\label{defn.declarative.paxos}
\label{defn.ThyPaxOne}
We define \deffont{Declarative Paxos} to be the logical theory 
$$
\ThyPax=(\Sigma_{\ThyPax},\Theta_{\ThyPax}),
$$ 
where:
\begin{enumerate*}
\item\label{item.ThyPaxOne.sig}
the signature $\Sigma_{\ThyPax}$ is 
$$
\Sigma_{\ThyPax}=[\tf{leader}:0;\ \tf{propose},\tf{send},\tf{write},\tf{accept},\tf{decide}:1] ,
$$ 
as per the notation in Definition~\ref{defn.signature}(\ref{item.predicate.syntax.signature}), and
\item\label{item.ThyPaxOne.axioms}
the axioms $\Theta_{\ThyPax}$ are as written in Figure~\ref{fig.logical.paxos}.
\end{enumerate*}
In axioms \rulefont{PaxPropose?} and \rulefont{PaxSend?} we assume an element $\udfn\in\Val$, to model the `undefined' value from Paxos (it does not matter what $\udfn\in\Val$ is; we just have to choose some element and use it consistently).
\end{defn}

\begin{rmrk}
\label{rmrk.functional.notation.in.paxos}
For convenient reference, we give pointers to help unpack the notation used in Figure~\ref{fig.logical.paxos}.
Note that this is just an overview focussing on items of notation used in the axioms, to help with reading them; in-depth discussions of the axioms will follow: 
\begin{enumerate*}
\item
We use the functional notation from Notation~\ref{nttn.hos} to elide quantified variables where convenient.
For example, axioms \rulefont{PaxPropose!} and \rulefont{PaxAccept!} written in full are 
$$
\tf{leader}\tnotor\texi\thea.\tf{propose}(\thea)
\quad\text{and}\quad
\final((\texi\thea.\someoneAll\tf{write}(\thea))\tnotor\texi\thea.\tf{accept}(\thea)) .
$$
\item
$\tnotor$ and $\timpc$ are the \emph{weak} and \emph{strong} implications respectively.
Truth-tables are in Figure~\ref{fig.3} and a discussion is in Subsection~\ref{subsect.3.imp}.
In particular, by Proposition~\ref{prop.mp.for.tnotor}, 
\begin{itemize*}
\item
$\phi\tnotor\phi'$ means `if $\phi$ is true then $\phi'$ is valid' (where `valid' means `$\tvT$ or $\tvB$') and 
\item
$\phi\timpc\phi'$ means `if $\phi$ is true then $\phi'$ is true'. 
\end{itemize*}
\item
The compound modalities --- $\someoneAll$, $\everyoneAll$, $\QuorumBox$, and $\CoquorumDiamond$ --- are discussed in Subsection~\ref{subsect.compound.modalities}.
Their meaning in full generality is in Lemma~\ref{lemm.dense.char}, but because the axioms only apply them to \emph{pointwise} predicates, their meanings simplify to `somewhere' ($\someoneAll$), `everywhere' ($\everyoneAll$), `for a set with a nonempty open interior' ($\QuorumBox$), and `for a dense set' ($\CoquorumDiamond$) respectively, as per Proposition~\ref{prop.pointwise.dense.char}. 
\item\label{item.discuss.final}
$\final$ means $\sometime\forever$ and is from Figure~\ref{fig.3.derived}.
It means `eventually, this holds for all future times'. 

In our axiomatisation this is used to reflect a standard assumption in distributed systems of \deffont{eventual network synchrony}, or synonymously \deffont{global stabilisation (of the network)}, meaning that the communications network may be unstable and fail to pass messages at first, but then after a finite but unbounded amount of time, it becomes stable and messages arrive.
Our logic does not explicitly represent messages, but the use of $\final$ in our forward rules reflects this assumption about the underlying implementation. 
\item
$\texiaffine$ in \rulefont{PaxPropose01} and \rulefont{PaxWrite01} means `for zero or one elements', as per Subsection~\ref{subsect.unique.affine.existence}, so these two axioms intuitively state that each participant proposes or writes at most one value.

Note that our logic is three-valued, so what these axioms mean \emph{precisely} is something a little richer than is the case in two-valued logic.
In particular, if $\Forall{v}\modellabel{\tf{write}(v)}{\acontext}=\tvB$ --- corresponding to the case where a participant is crashed --- then $\acontext\ment\texiaffine\tf{write}$, \emph{even though} $\acontext\ment\tf{write}(v)$ holds for every $v$. 
See the discussion and results in Subsection~\ref{subsect.unique.affine.existence}.
\item
$\correct{\tf{propose}}$ means that $\tf{propose}$ returns $\tvT$ or $\tvF$ (not $\tvB$) on every value; see Definition~\ref{defn.correct}(\ref{item.correct}).
Similarly for the other unary predicate symbols in \rulefont{PaxPCorrect}.
\item
$\mru{\tf{accept}}{v}$ means `a most recently accepted value is $v$'.
See the discussion and definition in Subsection~\ref{subsect.mrup}.
\item
Finally, note that we intend the predicates in Figure~\ref{fig.logical.paxos} to be read as axioms.

What it means for an axiom to be valid in a model --- i.e. how we judge whether a model satisfies $\Theta_\ThyPax$, or how we judge whether a predicate is valid in all models of $\Theta_\ThyPax$ --- is unpacked in Figure~\ref{fig.validity} and Definition~\ref{defn.theory.axiom}(\ref{item.axioms.valid}).
We discuss the quantifier scoping in those definitions in Remark~\ref{rmrk.thy.quantifier.scope}.
\end{enumerate*}
\end{rmrk}




\subsection{High-level discussion of the axioms}

\begin{rmrk}
The predicates in Figure~\ref{fig.logical.paxos} are intended as a \emph{theory}, in the sense of Definition~\ref{defn.theory.axiom}, which is why we call them \emph{axioms} and not just `some predicates'.

That is: we fix a signature $\Sigma$ that has (at least) the predicate symbols required to express the axioms in Figure~\ref{fig.logical.paxos}, and we fix a model $\mathcal M$, and we derive properties that hold in all the valuations $\avaluation$ to $\mathcal M$ such that the theory $\Theta_\ThyPax$ from Figure~\ref{fig.validity} is valid.
 
At least as much design effort went into crafting the axioms $\Theta_\ThyPax$ of Declarative Paxos in Figure~\ref{fig.logical.paxos}, as went into designing the three-valued modal fixedpoint logic in which the axioms are expressed.\footnote{Of course there was design feedback in both directions; we put into the logic what we needed to express the axioms, and we refined the axioms to elegantly work within the logic.} 
Therefore, we will take our time to discuss the axioms in detail.
The reader who just wants to see proofs can skip the rest of this Section (for now) and look to Section~\ref{sect.paxos.correctness.properties} onwards.
\end{rmrk}


\begin{rmrk}
\label{rmrk.axioms.paxos.three.groups}
We have organised the axioms in Figure~\ref{fig.logical.paxos} into three groups, which we now discuss:
\begin{enumerate}
\item
\emph{Backward (inference) rules.}

Axioms with names ending with $?$ represent backwards reasoning and so we call them \deffont{backward rules}.
They use the strong implication $\timpc$ (Definition~\ref{defn.three.imp}(\ref{item.three.elementary.strong.implication})), which means `if the LHS is true ($\tvT$) then the RHS is true ($\tvT$)'.

The backward rules represent a strong inference that \emph{if} some $Y$ is observed to be $\tvT$ then it \emph{must} be because some other property $X$ is $\tvT$.
For instance, \rulefont{PaxAccept?} says that if $\tf{accept}(v)$ is true, then $\someoneAll\tf{write}(v)$ must also be true.

We discuss the backward rules individually in Remark~\ref{rmrk.discuss.backward.rules} below.
\item
\emph{Forward (progress) rules.}

Axioms with names ending with $!$ represent forwards reasoning about making progress, so we call them \deffont{forward rules}.
They use $\final$ and the weak implication $\tnotor$ (Definition~\ref{defn.three.imp}(\ref{item.three.elementary.weak.implication})), which means `if the LHS is true ($\tvT$) then the RHS is valid ($\tvT$ or $\tvB$)'.

We derive the forward rules from the backward rules using a systematic recipe which we discuss in Remark~\ref{rmrk.recipe}.

The forward rules represent a weak inference that \emph{if} some precondition $X$ is $\tvT$ then --- provided the network is stable and provided that the participant has not crashed (where `crashed' is interpreted as returning $\tvB$) --- then the postcondition $Y$ is $\tvT$. 

For instance, \rulefont{PaxAccept!} expresses that, after network synchrony (Remark~\ref{rmrk.functional.notation.in.paxos}(\ref{item.discuss.final})), if $\texi\someoneAll\tf{write}$ is true (by Notation~\ref{nttn.hos} this is shorthand for $\texi a.\someoneAll\tf{write}(a)$) then if the participant has not crashed then $\texi\tf{accept}$ (i.e. $\texi a.\tf{accept}(a)$) is true. 

We discuss the forward rules in Remark~\ref{rmrk.discuss.forward.rules} below.
\item
\emph{Other rules, for correctness and liveness.}

Other rules reflect particular conditions of the algorithms.
We discuss them in Remark~\ref{rmrk.paxos.axioms.discussion}.
\end{enumerate}
\end{rmrk}

\begin{rmrk}[Recipe for deriving forward rules from backward rules]
\label{rmrk.recipe}
A design feature of Figure~\ref{fig.logical.paxos} is that each forward rule is obtained from the corresponding backward rule by a \deffont{recipe}.
Namely: for a backward rule of the form 
$$
\tall v.\ \chi \timpc \phi[a\ssm v] \timpc \psi[a\ssm v]
$$
where 
\begin{itemize*}
\item
$\chi$ expresses some possible \emph{side-conditions} (e.g. $\chi=\tf{leader}$ in \rulefont{PaxDecideL?}) and 
\item
$\phi$ expresses some \emph{preconditions}, and 
\item
$\psi$ expresses some \emph{postconditions}, 
\end{itemize*}
our recipe constructs a forward rule of the form
$$
(\texi v.\,\chi \tand \phi)\tnotor \texi v.\psi .
$$
In most of the forward rules this result also get wrapped in a $\final$ modality, representing network synchrony  (Remark~\ref{rmrk.functional.notation.in.paxos}(\ref{item.discuss.final})).

This recipe holds for every pair of rules:
\begin{itemize*}
\item
We see it cleanly in the pairs \rulefont{PaxAccept?} / \rulefont{PaxAccept!}, \rulefont{PaxDecideL?} / \rulefont{PaxDecideL!}, and \rulefont{PaxDecide\tneg L?} / \rulefont{PaxDecide\tneg L!}.
These rules follow our recipe symbol-for-symbol.
\item
We see this also with \rulefont{PaxPropose?} / \rulefont{PaxPropose!}.
Our recipe generates the following version of \rulefont{PaxPropose!}: 
$$
(\texi v.\,\tf{leader}\tand v\tneq\tf{undef})\timpc\texi v.\tf{propose}(v).
$$
The reader can check that this is logically equivalent to what is written in Figure~\ref{fig.logical.paxos}.
\item 
It is less clear how our recipe works for \rulefont{PaxSend?} and \rulefont{PaxSend!} and \rulefont{PaxWrite?} and \rulefont{PaxWrite!}; e.g. \rulefont{PaxSend!} looks different from \rulefont{PaxSend?} --- it is clearly simpler!
Remarkably, these pairs too are consistent with our recipe; see Remark~\ref{rmrk.why.evaporates} and the subsequent discussion and results.
\end{itemize*}
\end{rmrk}

\begin{rmrk}
We make some further elementary comments on our recipe from Remark~\ref{rmrk.recipe}:
\begin{enumerate*}
\item
We uncurry the forward rules --- meaning that we write $(\texi v.\,\chi \tand \phi)\tnotor \texi v.\psi$ instead of $(\texi v.\,\chi) \tnotor \phi\tnotor \texi v.\psi$ --- because this is convenient for some proofs. 
It makes no other difference and either form would be correct.
\item
A forward rule of the form $\texi v.((\chi \tand \phi)\tnotor \psi)$ would be \emph{wrong}.
This is just a different predicate and is not what is intended.
\item
A forward rule of the form $\tall v.((\chi \tand \phi)\tnotor \psi)$ would not necessarily be wrong, but it is stronger than we need for our proofs.
This reflects a general design principle in logical theories that \emph{weaker axioms are better}.
If we can prove the same result from weaker axioms then generally speaking we should, because that is a stronger result.\footnote{In practice, working from weaker axioms tends to force us to think harder and so design a better-engineered proof.}
\end{enumerate*}
\end{rmrk}


\subsection{More detailed discussion of the axioms}


\begin{rmrk}
\label{rmrk.paxos.axioms.discussion}
We discuss the \emph{axioms for correctness and liveness}:
\begin{enumerate}
\item\label{item.discuss.LdrExist}
\rulefont{LdrExist}:

In implementations of Paxos, at each stage there is a unique \emph{leader}, whose identity is computed by applying (what amounts to) a random number generator to an agreed seed.
As the level of abstraction in this paper, we just assume that a leader exists --- we do not care how it is chosen.

We use a predicate constant symbol $\tf{leader}$ to identify the leader.
Because our logic is modal --- across time $n\in\Time$ and participants $p\in\Pnt$ --- we can just assume that a constant symbol $\tf{leader}$ is assigned truth-value $\tvT$ where a participant $p$ that is a leader at a stage $n$, and is assigned truth-value $\tvF$ where a participant $p$ that is not a leader at a stage $n$.

In more technical language: $\avaluation$ is such that $n,p\mentval \modTF\tf{leader}$ for every time $n\in\Time$ and participant $p\in\Pnt$, %
and $p$ is a leader %
precisely when $n,p\mentval\tf{leader}$.
\rulefont{LdrExist} just asserts that a participant either is a leader or it is not, and some leader exists.
\item\label{item.discuss.LdrExistUniq}
\rulefont{LdrExt}:

\rulefont{LdrExt} is a rule-scheme (one for each closed $\phi$) and asserts that leaders are extensionally equal; %
a formal statement and proof are in Lemma~\ref{lemm.LeaderExtUniq}.

In an implementation this is enforced by having a unique leader, but in the logic this is unnecessary --- we don't care how many participants claim to be leaders, so long as they are all saying the same thing! --- and also it is impossible because (although we have modalities $\everyone$ and $\someone$) we do not have an explicit quantification over elements of $\Pnt$ in our logic.
This is a feature, not a bug: we want our logic to be no more expressive than it needs to be.

In fact, the axiom-scheme \rulefont{LdrExt} is stronger than we need for our proofs;
more on this in Subsection~\ref{subsect.ldrexistuniq'}.


Henceforth we may say `the leader' instead of `a leader', because up to extensional equivalence there is just one.
\item
\rulefont{LdrCorrect}:

\rulefont{LdrCorrect} asserts that for infinitely many times $n\in\Time$, the leader $l$ (which exists by the right-hand conjunct in \rulefont{LdrExist}) is $\tf{propose}$-correct, $\tf{write}$-correct, and $\tf{decide}$-correct (terminology from Definition~\ref{defn.correct}(\ref{item.P-correct}).

Unpacking Definition~\ref{defn.correct}(\ref{item.P-correct}), this means that $n,l\mentval \modTF\tf{propose}(v)$ for every $v\in\Val$, and similarly for $\tf{write}$ and $\tf{decide}$.
Intuitively, this is a \emph{liveness assumption} that there is always a future time when the leader is not crashed in the sense that it can correctly $\tvT\tvF$-propose, $\tvT\tvF$-write, and $\tvT\tvF$-decide a value. 
\item
\rulefont{PaxPCorrect}:

Given a valuation $\avaluation$, we unpack Definition~\ref{defn.validity.judgement}(\ref{item.mentval}) to see what $\mentval \rulefont{PaxPCorrect}$ actually means, as per Definition~\ref{defn.theory.axiom}:
\begin{multline*}
\Forall{n\in\Time}\Exists{O\in\opens}\Forall{p\in O}\Forall{v\in\Val}n,p,O\mentval\correct{\tf P(v)}
\\
\text{for each}\ \tf P\in\{\tf{propose},\tf{send},\tf{accept},\tf{write},\tf{decide}\} .
\end{multline*}
This corresponds to a standard correctness assumption in Paxos that at each stage there is a quorum of uncrashed participants.

However, the axiom \rulefont{PaxPCorrect} is subtly more general than the standard correctness assumption, because there is no restriction that the \emph{same} open set $O$ be used for each unary predicate symbol.\footnote{If we wanted the same $O$, then we would scope the $\QuorumBox$ modality further out and write $\QuorumBox(\correct{\tf{propose}} \tand \correct{\tf{send}} \tand \correct{\tf{accept}} \tand 
\correct{\tf{write}} \tand \correct{\tf{decide}})$.  This turns out to be unnecessary.}
It turns out that the proofs require only this assumption, and no more. 
\item\label{item.01.discussion}
\rulefont{PaxPropose01} and \rulefont{PaxWrite01}:

We justify these axioms with respect to the Paxos algorithm as summarised in Remark~\ref{rmrk.high-level.paxos}:
\begin{itemize*}
\item
Step~\ref{item.l.proposes} of Remark~\ref{rmrk.high-level.paxos} states that leader is crashed or proposes one value.
The intuitive correspondence to \rulefont{PaxPropose01} is clear.
\item
Step~\ref{item.l.writes} of Remark~\ref{rmrk.high-level.paxos} states that the leader writes a value derived from \emph{the first} quorum of participants that it hears back from (or none, if it hears back from no such quorum). 
One corollary of this, is that it writes at most one value.
Our logical theory has no state machine and notion of algorithmic time (as per Remark~\ref{rmrk.elisions}(\ref{item.elide.algorithmic.time})) --- so there is no sense in which information from one open set `arrives' at the leader `before' information from another --- so from the point of view of the logic, of all the open sets available, the leader picks at most one to act on.
This is axiom \rulefont{PaxWrite01}. 
\end{itemize*}
It is a fact that \rulefont{PaxPropose01}+\rulefont{PaxPropose!} and \rulefont{PaxWrite01}+\rulefont{PaxWrite!} give us
$$
\tf{leader}\tnotor\texiunique\tf{propose}
\quad\text{and}\quad
\final((\tf{leader}\tand\QuorumBox\texi\tf{send})\tnotor \texiunique\tf{write}) ,
$$ 
and
\rulefont{PaxPropose01}+\rulefont{PaxPropose?} and \rulefont{PaxWrite01}+\rulefont{PaxWrite?} give us
$$
\tneg\tf{leader}\tnotor\tneg\texi\tf{propose}
\quad\text{and}\quad
\tneg\tf{leader}\tnotor\tneg\texi\tf{write} .
$$
So it follows from our axioms that a leader participant proposes and writes one value (assuming it is not crashed, and if messages arrive), and a non-leader participant proposes and writes zero values. 
\item
\rulefont{Pax2Twined}:

This axiom collects a set of instances of Lemma~\ref{lemm.supertwined.qcbt}(\ref{item.intertwined.3}), namely taking $\phi=\tf{accept}(v)$ in that Lemma for every $v\in\Val$.

There is a design choice here: instead of an axiom \rulefont{Pax2Twined}, we could insist specifically that the semitopology in our underlying model be $2$-twined.
Then \rulefont{Pax2Twined} ceases to be an axiom, but we still have it as a fact of the model by directly invoking Lemma~\ref{lemm.supertwined.qcbt}(\ref{item.intertwined.3}).
Or we could be even more specific: we could insist that the underlying semitopology be $\AllBut{N}{f}$ (Definition~\ref{defn.allbut}) for $N> 2*f$; by Proposition~\ref{prop.allbut.supertwined}(\ref{item.allbut.intertwined}) this is $2$-twined, and then we can use Lemma~\ref{lemm.supertwined.qcbt}(\ref{item.intertwined.3}) to obtain \rulefont{Pax2Twined}.

We choose here to express as an axiom precisely the property that our proofs require.
This is concise, general, and yields clean proofs.
Insisting on a $2$-twined semitopology, or specifically on the semitopology $\AllBut{N}{f}$ for $N>2*f$, would also make sense, but in the context of our logical approach, using an axiom seems nice. %
\end{enumerate}
\end{rmrk}

\begin{rmrk}
\label{rmrk.discuss.backward.rules}
We discuss the \emph{axioms for the backward rules}:
\begin{enumerate}
\item
\rulefont{PaxPropose?}:

If we propose some value $v$ then 
\begin{enumerate}
\item
we are the leader and 
\item
$v$ is not the undefined value $\udfn$.
\end{enumerate}
\item
\rulefont{PaxSend?}:

If we send some value $v$ then 
\begin{enumerate}
\item
somebody must have proposed some value and 
\item
\begin{enumerate*}
\item
\emph{either} $v$ is the value we have most recently accepted, 
\item
\emph{or} we never accepted any value and $v$ is the undefined value.
\end{enumerate*}
\end{enumerate}
\item
\rulefont{PaxWrite?}:

If we write some value $v$ then 
\begin{enumerate}
\item
we are a leader, and
\item
there exists an open set of points (= participants) that have all sent us values, and
\item
\begin{enumerate*}
\item
\emph{either} some point in that open set has most recently accepted $v$, 
\item
\emph{or} 
$v$ is the value that we proposed, and no point in that open set has ever accepted any value. 
\end{enumerate*}
\end{enumerate}
\item
\rulefont{PaxAccept?}:

If we accept $v$ then somebody somewhere must have written $v$.
\item
\rulefont{PaxDecideL?}:

If a leader decides $v$ then an open set of points must have all accepted $v$.
\item
\rulefont{PaxDecide{\tneg}L?}:

If we are not a leader and we decide $v$, then some leader somewhere must have decided $v$. 
\end{enumerate}
\end{rmrk}

\begin{rmrk}
\label{rmrk.discuss.forward.rules}
We discuss the \emph{axioms for the forward rules}.

We noted in Remark~\ref{rmrk.recipe} our design methodology of deriving the forward rules from the backward rules; forward rules are rules for making progress, so they are existentially quantified over the variable $v$ and they use the weak implication $\tnotor$ to emulate possible crashes.

The other subtlety is the $\final$ modality wrapping all of the forward rules, except for \rulefont{PaxPropose!}.
\begin{itemize*}
\item
\emph{What $\final$ does for us in most of the forward rules.}

It models possible network failure, and eventual network synchrony.
As touched on in Remark~\ref{rmrk.functional.notation.in.paxos}(\ref{item.discuss.final})), the $\final$ modality reflects the possibility that the network might be down for a while, but eventually --- `finally' --- the network stabilises and messages arrive.

Predicates in our logic are \emph{not} messages, but we get the effect of eventual network synchrony in our logic (even though messages are not explicitly present) by wrapping the implications in our forward rules in the $\final$ modality.

More on this when we introduce the notion of \emph{global stabilisation of logical time} in Subsection~\ref{subsect.gslt}.
\item
\emph{Why $\final$ is absent from \rulefont{PaxPropose!}.}

Because this rule models a property of the implementation that involves no network communication.
The leader does not need to communicate with any other participant to pick a value to propose.

It is true that the leader's propose message might not arrive if the network is down --- but this is reflected by the use of $\final$ in the \rulefont{PaxSend!} rule.
\end{itemize*}
\end{rmrk}


\subsection{Some asides}

\subsubsection{An aside on \rulefont{PaxWrite?}}

\rulefont{PaxWrite?} contains this predicate:
$$
\mru{(\someone\tf{accept})}{v} .
$$
Unpacking what this means (cf. Remark~\ref{rmrk.explain.mrup}), this looks into the past to check whether $v$ is a value that was most recently accepted by some participant.

Note that this slightly differently-scoped predicate would be \emph{wrong}:
$$
\someone(\mru{\tf{accept}}{v}) .
$$
This finds some participant such that $v$ is a value that was most recently accepted by \emph{that} participant (but there might be another participant who more recently accepted some other value).

The reader with a background in logic will know all about how important scoping can be (e.g. $\forall x.\exists y.\phi$ vs. $\exists x.\forall y.\phi$, and $\exists x.(\phi\land \phi')$ vs. $\exists x.\phi\land\exists x.\phi'$), but perhaps the point bears repeating.


\subsubsection{An aside on \rulefont{LdrExt}}

In a Paxos implementation, there is a unique leader at each stage. 
We mentioned in Remark~\ref{rmrk.paxos.axioms.discussion}(\ref{item.discuss.LdrExist}\&\ref{item.discuss.LdrExistUniq}) that \rulefont{LdrExist} asserts that a leader exists and that the \rulefont{LdrExt} axiom-scheme asserts that all leader participants are extensionally equal.
We now make this formal:\footnote{Later on in Subsection~\ref{subsect.ldrexistuniq'} we will note weaker, but more abstract, forms of the axiom.}
\begin{lemm}[Leaders are extensionally equivalent]
\label{lemm.LeaderExtUniq}
Suppose $n\in\Time$ and $\avaluation$ is a valuation, and suppose $n\mentval\rulefont{LdrExt}_\phi$ holds for every closed predicate $\phi$.\footnote{By Definition~\ref{defn.validity.judgement}(\ref{item.ment.n}) this means $\Forall{p\in\Pnt}\Forall{O\in\opensne}(\modellabel{\rulefont{LdrExt}_\phi}{n,p,O,\avaluation}\in\{\tvT,\tvB\})$.} 
Suppose that $l,l'\in\Pnt$ and $n,l\mentval\modT\tf{leader}$ and $n,l'\mentval\modT\tf{leader}$.

Then $l$ and $l'$ are \deffont{extensionally equivalent} at $(n,\avaluation)$, by which we mean that for every closed predicate $\phi$ and every $O\in\opensne$,
$$
\modellabel{\phi}{n,l,O,\avaluation} = \modellabel{\phi}{n,l',O,\avaluation} .
$$
\end{lemm}
\begin{proof}
Suppose $\phi$ is a closed predicate and $O\in\opensne$, and suppose for the sake of argument that $\modellabel{\phi}{n,l,O,\avaluation}=\tvB$.

It is routine to check from Figures~\ref{fig.3.derived} and~\ref{fig.3} and from Lemma~\ref{lemm.dense.char}(\ref{item.dense.char.someoneAll}) that $\modellabel{\someoneAll\modB\phi}{n,l,O,\avaluation}=\tvT$. 
By assumption $n,l',O\mentval \rulefont{LdrExt}_{\modB\phi}$, and by two applications of \strongmodusponens (once for $\modT\tf{leader}$, and another for $\modT\someoneAll(\tf{leader}\tand\modB\phi)$) we conclude that $\modellabel{\modB\phi}{n,l',O,\avaluation}=\tvT$, and therefore $\modellabel{\phi}{n,l',O,\avaluation}=\tvB$ as required. 
 
The cases $\modellabel{\phi}{n,l,O,\avaluation}=\tvT$ and $\modellabel{\phi}{n,l,O,\avaluation}=\tvF$ are be precisely similar, and this suffices to prove the result.
\end{proof}


\subsubsection{An aside on \rulefont{PaxSend?} and \rulefont{PaxSend!}}

\begin{rmrk}
\label{rmrk.why.evaporates}
In Remark~\ref{rmrk.recipe} we proposed a recipe for deriving a forward rule from a backward one: a backward rule of the form
$$
\rulefont{Rule?} 
\qquad
\tall v.\ \chi \timpc \phi[a\ssm v] \timpc \psi[a\ssm v]
$$
generates a forward rule (possibly wrapped in a $\final$ modality) of the form
$$
\rulefont{Rule!}
\qquad
(\texi v.\, \chi \tand \phi)\tnotor \texi v.\psi .
$$
We noted in Remark~\ref{rmrk.recipe} that the recipe works cleanly for some rules --- but it was less clear how it can be applied to derive \rulefont{PaxSend!} from \rulefont{PaxSend?}, and \rulefont{PaxWrite!} from \rulefont{PaxWrite?}.
In these cases, the left-hand sides of \rulefont{PaxSend!} and \rulefont{PaxWrite!} seem much simpler than the right-hand sides of \rulefont{PaxSend?} and \rulefont{PaxWrite?} respectively. 

We now show why, in fact, \rulefont{PaxSend!} and \rulefont{PaxWrite!} are correct and make sense in the context of our recipe, culminating with Corollary~\ref{corr.back.ok}.
\end{rmrk}


\begin{lemm}
\label{lemm.evaporates.simplify}
We have:
\begin{enumerate*}
\item
$\ment \texi v.\bigl((\mru{\tf{accept}}{v})\ \tor (\tneg\modT\recent\texi\tf{accept}\tand v\teq\udfn)\bigr) \ \equiv\ \tvT$.
\item
$\ment \texi v.\bigl(\ (\mru{(\someone\tf{accept})}{v})\ \tor\ (\tneg\modT\recent\texi\someone\tf{accept}\tand \tf{propose}(v))\ \bigr)\ \equiv\ \texi\tf{propose}$.
\end{enumerate*}
Note that:
\begin{itemize*}
\item
The notation $\texi\tf{accept}$ and $\texi\tf{propose}$ and $\texi\someone\tf{accept}$ is from Notation~\ref{nttn.hos}.
\item
We do not restrict to contexts and valuations that satisfy \ThyPax; the meaning of the $\ment$ symbol above is as in Definition~\ref{defn.validity.judgement}(\ref{item.ment.phi}).
\end{itemize*}
\end{lemm}
\begin{proof}
We consider each part in turn:
\begin{enumerate}
\item
We distribute the $\texi v$ through the $\tor$, and then through the $\tand$ on the right-hand side (we can do this because $v$ is not free in the left-hand conjunct) and we see that it suffices to show 
$$
\ment \texi v.(\mru{\tf{accept}}{v})\ \tor\ (\tneg\modT\recent\texi\tf{accept}\tand \texi v.v\teq\udfn)\ \equiv\ \tvT
$$
Using Lemma~\ref{lemm.evaporates} this simplifies to just
$$
\ment \modT\recent\texi\tf{accept}\ \tor\ (\tneg\modT\recent\texi\tf{accept}\tand \texi v.v\teq\udfn) \ \equiv\ \tvT.
$$
This clearly holds, just taking $v\teq\udfn$.
\item
We distribute $\texi v$ into the predicate as we did for the proof of part~1.
We see that it suffices to show 
$$
\ment \bigl(\ \texi v.(\mru{(\someone\tf{accept})}{v})\ \tor\ (\tneg\modT\recent\texi\someone\tf{accept}\tand \texi\tf{propose})\ \bigr)\ \equiv\ \texi\tf{propose}.
$$
We simplify using Lemma~\ref{lemm.evaporates} to
$\ment \texi\tf{propose}\ \equiv\ \texi\tf{propose}$,
which is true.\footnote{Yes, we have a lemma for that: Lemma~\ref{lemm.equiv}(\ref{item.equiv.2}) and reflexivity of equality.}
\qedhere
\end{enumerate}
\end{proof}

We can now justify the left-hand sides of \rulefont{PaxSend!} and \rulefont{PaxWrite!} as being consistent with our recipe as discussed above in Remark~\ref{rmrk.why.evaporates}:
\begin{corr}
\label{corr.back.ok}
We have:
$$
\begin{array}{r@{\ }l}
\ment\texi v.&\bigl(\someoneAll\texi\tf{propose}\ \tand \bigl((\mru{\tf{accept}}{v})\ \tor (\tneg\modT\recent\texi\tf{accept}\tand v\teq\udfn)\bigr)\bigr) 
\\
&\equiv \someoneAll\texi\tf{propose}
\\[2ex]
\ment\texi v.&\bigl(\tf{leader}\tand 
\Quorum(\Box\texi\tf{send} \tand 
\bigl((\mru{(\someone\tf{accept})}{v}) \tor
(\tneg\modT\recent\texi\someone\tf{accept}\tand \tf{propose}(v))\bigr)) 
\bigr)
\\
&\equiv
(\tf{leader}\tand \QuorumBox\texi\tf{send} \tand \texi\tf{propose})
\end{array}
$$
\end{corr}
\begin{proof}
By elementary predicate manipulations, distributing the $\texi v$ quantifier down into the predicate and using Lemma~\ref{lemm.evaporates.simplify}.
\end{proof}

 




\section{Correctness properties}
\label{sect.paxos.correctness.properties}

\subsection{Summary of the properties}

\begin{rmrk}
\label{rmrk.standard.paxos.correctness}
We will derive the following standard correctness properties for Paxos (as laid out in~\cite[Subsection~2.2.1]{cachin2011paxos}) from the theory \ThyPax in Figure~\ref{fig.logical.paxos}:
\begin{enumerate*}
\item
\textbf{Validity:} \quad If a participant decides some value $v$, then $v$ was proposed by some participant.

A formal statement and proof is in Theorem~\ref{thrm.validity}.
\item
\textbf{Agreement:}\quad 
If two participants decide values, then those values are equal.

A formal statement and proof is in Theorem~\ref{thrm.full.agreement}.
\item
\textbf{Termination:}\quad 
Every correct participant eventually decides some value.

A formal statement and proof is in Theorem~\ref{thrm.termination}.
\item
\textbf{Integrity:}\quad 
Every correct (non-crashing) participant decides at most once.

Our logic abstracts away from a direct notion of execution, by design, but this property is represented indirectly in how Theorem~\ref{thrm.termination} is expressed.
See Remark~\ref{rmrk.paxos.integrity} below. 
\end{enumerate*}
\end{rmrk}

\begin{rmrk}[On integrity]
\label{rmrk.paxos.integrity}
In implementations of Paxos, once a participant has decided, it stops.
But in our logic, `stop' is not even a concept.
There is no abstract machine to stop; there are just predicates, with denotations.

Nevertheless, the integrity property is echoed in our maths, as the \emph{infinite repetition} in Theorem~\ref{thrm.termination}.
Once a predicate $\tf{decide}(v)$ returns $\tvT$, it does so \emph{forever} --- and because of Agreement, it never decides any other value.
A predicate eventually always returning the same truth-value is a model of an abstract machine that has computed its result, and then just sits there forever, displaying that same result, unchanged, every time we query its display.
\end{rmrk}

\subsection{Validity}


\begin{rmrk}
\label{rmrk.validity.in.english}
\emph{Validity} is expressed in English as:
\begin{quote}
\emph{If a participant decides some value $v$, then $v$ was proposed by some participant.}
\end{quote}
We render this in our logic in Theorem~\ref{thrm.full.agreement} as
$$
\ThyPax\ment\tall v.(\tf{decide}(v) \timpc \urecent\someoneAll(\tf{leader}\tand\tf{propose}(v))).
$$

In words we can read this as: 
\begin{quote}
If someone decides $v$ now, then at some stage now or in the past a leader of that stage proposed $v$. 
\end{quote}
We will now set about proving this property from our axioms, culminating with Theorem~\ref{thrm.validity}.
\end{rmrk}


\begin{lemm}
\label{lemm.validity.helper}
Suppose $\avaluation$ is a valuation such that $\mentval\ThyPax$ ($\mentval\ThyPax$ is from Figure~\ref{fig.validity}).
Suppose $n\in\Time$ and $v\in\Val$.
Then:
\begin{enumerate*}
\item\label{item.validity.helper.3}
$n\mentval \someoneAll\tf{write}(v) \timpc \urecent\someoneAll\tf{propose}(v)$. 
\item\label{item.validity.helper.2}
$n\mentval \someoneAll\tf{accept}(v) \timpc \urecent\someoneAll\tf{propose}(v)$. 
\end{enumerate*}
\end{lemm}
\begin{proof}
We prove both parts by a simultaneous induction on $n$.
We freely use \strongmodusponens:
\begin{enumerate}
\item
Suppose $n\mentval\modT\someoneAll\tf{write}(v)$.
Then using \rulefont{PaxWrite?} %
$$
n\mentval \modT\Quorum(\mru{\someone\tf{accept}}{v}) 
\quad\text{or}\quad
n\mentval \modT\someoneAll\tf{propose}(v) .
$$
We consider the two cases:
\begin{itemize*}
\item
\emph{If $n\mentval \modT\someoneAll\tf{propose}(v)$,}\ then 
$n\mentval\modT\urecent\someoneAll\tf{propose}(v)$ follows immediately from the clause for %
$\urecent$ in Figure~\ref{fig.3.derived}.
\item
\emph{If $n\mentval \modT\Quorum(\mru{\someone\tf{accept}}{v})$,}\ 
then by Lemma~\ref{lemm.mru.to.recent}(\ref{item.mru.to.recent.1}) $n\mentval\modT\recent\someoneAll\tf{accept}(v)$,
so that by Figure~\ref{fig.3.derived} $n'\mentval\modT\someoneAll\tf{accept}(v)$ for some $n'<n$.
By part~\ref{item.validity.helper.2} of this result for $n'<n$, we have that $n'\mentval\modT\urecent\someoneAll\tf{propose}(v)$, and thus using Lemma~\ref{lemm.tomorrow.forever.recent}(\ref{item.rur.r}) $n\mentval\modT\urecent\someoneAll\tf{propose}(v)$. 
\end{itemize*}
In either case, we are done.
\item
Suppose $n\mentval\modT\someoneAll\tf{accept}(v)$.
By \rulefont{PaxAccept?} $n\mentval\modT\someoneAll\tf{write}(v)$.
We use part~\ref{item.validity.helper.3} of the inductive hypothesis for $n$. 
\qedhere\end{enumerate}
\end{proof}



\begin{thrm}[Validity]
\label{thrm.validity}
We have:
$$
\ThyPax\ment\tall v.(\tf{decide}(v) \timpc \urecent\someoneAll(\tf{leader}\tand\tf{propose}(v))) .
$$
\end{thrm}
\begin{proof}
Suppose $\avaluation$ is a valuation such that $\mentval\ThyPax$. 
We reason using \strongmodusponens. 
Suppose $v\in\Val$ and $n\in\Time$ and $p\in\Pnt$, and suppose $n,p\mentval\modT\tf{decide}(v)$.
By \rulefont{PaxDecideL?} (and \rulefont{PaxDecide\tneg L?} if $p\neq\f{leader}(n)$), $n\mentval\modT\QuorumBox\tf{accept}(v)$.
It follows by Proposition~\ref{prop.someone.implies.someoneAll}(\ref{item.someone.implies.someoneAll.2}) that $n\mentval\modT\someoneAll\tf{accept}(v)$.
By Lemma~\ref{lemm.validity.helper}(\ref{item.validity.helper.2}) $n\mentval\modT\urecent\someoneAll\tf{propose}(v)$, and using \rulefont{PaxPropose?} we conclude that $n\mentval\modT\urecent\someoneAll(\tf{leader}\tand\tf{propose}(v))$ as required. 
\end{proof}




\subsection{Agreement}

\begin{rmrk}
\emph{Agreement} is expressed in English as:
\begin{quote}
\emph{If two participants decide values, then those values are equal.}
\end{quote}
We render this in our logic in Theorem~\ref{thrm.full.agreement} as
$$
\ThyPax\ment\tall v,v'.(\urecent\someoneAll\tf{decide}(v) \timpc \someoneAll\tf{decide}(v') \timpc v'\teq v) .
$$
In words we can read this as: 
\begin{quote}
If someone decides $v$ now or in the past, and someone decides $v'$ now, then $v'=v$.
\end{quote}
It is convenient to prove the `now' and the `in the past' parts of this as two sub-cases: 
\begin{enumerate*}
\item
Proposition~\ref{prop.decide.n.agree} handles the case of $\someoneAll\tf{decide}(v)\tand\someoneAll\tf{decide}(v')$, i.e. when the two decisions are made at the same time.
\item
Theorem~\ref{thrm.full.agreement} then extends this with the case of $\recent\someoneAll\tf{decide}(v)\tand\someoneAll\tf{decide}(v')$, i.e. when one decision precedes the other.
\end{enumerate*}
\end{rmrk}


\subsubsection{Agreement at time $n$}



\begin{lemm}
\label{lemm.various.unique}
Suppose $\avaluation$ is a valuation such that $\mentval\ThyPax$. 
Suppose $n\in\Time$ and $v,v'\in\Val$. 
Then:
\begin{enumerate*}
\item\label{item.various.unique.write}
$n\mentval (\someoneAll\tf{write}(v)\tand \someoneAll\tf{write}(v')) \timpc v\teq v'$. 
\item\label{item.various.unique.accept}
$n\mentval (\someoneAll\tf{accept}(v)\tand \someoneAll\tf{accept}(v')) \timpc v\teq v'$. 
\item\label{item.various.unique.recent.accept}
$n\mentval \bigl(\mru{(\someoneAll\tf{accept})}{v}\tand \mru{(\someoneAll\tf{accept})}{v'}\bigr) \timpc v\teq v'$.
\end{enumerate*}
\end{lemm}
\begin{proof}
We consider each part in turn, using \strongmodusponens (we may also silently commute the $\modT$ modality as convenient):
\begin{enumerate}
\item
Suppose $n\mentval\modT\someoneAll\tf{write}(v)\tand\modT\someoneAll\tf{write}(v')$.
Using \rulefont{PaxWrite?} 
$n\mentval\modT\someoneAll(\tf{write}(v)\tand\tf{leader})\tand\modT\someoneAll(\tf{write}(v')\tand\tf{leader})$. 
Using \rulefont{LdrExt} therefore $n\mentval\modT\someoneAll(\tf{write}(v)\tand\tf{write}(v')\tand\tf{leader})$.
By \rulefont{PaxWrite01} and Lemma~\ref{lemm.unique.affine.existence}(\ref{item.unique.affine.existence.01.implies}) we have $v=v'$ as required.
\item
Suppose $n\mentval \modT\someoneAll\tf{accept}(v)\tand\modT\someoneAll\tf{accept}(v')$.
By \rulefont{PaxAccept?} $n\mentval\modT\someoneAll\tf{write}(v)\tand\modT\someoneAll\tf{write}(v')$.
We use part~\ref{item.various.unique.write} of this result.
\item
Suppose $n\mentval\modT(\mru{(\someoneAll\tf{accept})}{v})\tand\modT(\mru{(\someoneAll\tf{accept})}{v'}\bigr)$.\footnote{By Lemma~\ref{lemm.mru.unambivalent}(2) we do not need to write the $\modT$ here but we do anyway; see Remark~\ref{rmrk.explain.mru}(\ref{item.explain.mru.T}).}
Using Lemma~\ref{lemm.mru.to.recent}(\ref{item.mru.to.recent.2}) (taking $a.\phi=\someoneAll\tf{accept}$ in that Lemma) there exists $n'<n$ such that 
$n'\mentval\modT\someoneAll\tf{accept}(v)\tand\modT\someoneAll\tf{accept}(v')$.
We use part~\ref{item.various.unique.accept} of this result.
\qedhere\end{enumerate}
\end{proof}



\begin{prop}
\label{prop.decide.n.agree}
Suppose $\avaluation$ is a valuation such that $\mentval\ThyPax$. 
Suppose $n\in\Time$ and $v,v'\in\Val$.
Then
$$
n\mentval (\someoneAll\tf{decide}(v) \tand \someoneAll\tf{decide}(v')) \timpc v\teq v' . 
$$
\end{prop}
\begin{proof}
We use \strongmodusponens.
Suppose 
$n\mentval \modT(\someoneAll\tf{decide}(v)\tand\someoneAll\tf{decide}(v'))$.
Using \rulefont{PaxDecide\tneg L?} if needed, and then \rulefont{PaxDecideL?} twice, we have
$n\mentval \modT(\QuorumBox\tf{accept}(v)\tand\QuorumBox\tf{accept}(v'))$ and thus by Proposition~\ref{prop.someone.implies.someoneAll}(\ref{item.someone.implies.someoneAll.2}) and Lemma~\ref{lemm.commuting.connectives}
also $n\mentval \modT(\someoneAll\tf{accept}(v)\tand\someoneAll\tf{accept}(v'))$.
We can now use Lemma~\ref{lemm.various.unique}(\ref{item.various.unique.accept}).
\end{proof}



\begin{rmrk}
Proposition~\ref{prop.decide.n.agree} proves part of the agreement property; that no two correct participants decide differently \emph{at a given time}.
To prove the result across different times will require a bit more work, culminating with Theorem~\ref{thrm.full.agreement}.
\end{rmrk}


\subsubsection{Some fine structure of the rules}
\label{subsect.ldrexistuniq'}

We continue a discussion from Remark~\ref{rmrk.paxos.axioms.discussion}(\ref{item.discuss.LdrExistUniq}) and also Lemma~\ref{lemm.LeaderExtUniq}:

Note that we use the axiom-scheme \rulefont{LdrExt} from Figure~\ref{fig.logical.paxos} just once: in the proof of Lemma~\ref{lemm.various.unique}(\ref{item.various.unique.write}).

Examining how we use this axiom-scheme in the proof above, we see that for our specific purposes it would suffice to assume a single axiom as follows:
$$
\hspace{-6em}\rulefont{LdrExt'}
\qquad\qquad
\tf{leader}\timpc \someoneAll(\tf{leader}\tand\tf{write}(v))\timpc \tf{write}(v) .
$$ 
The axiom scheme \rulefont{LdrExt} asserts that there is a unique leader at each stage, up to extensional equivalence, which is a fair reflection of how the algorithms work where there \emph{is} a unique leader at each stage.

\rulefont{LdrExt'} asserts that there is a unique leader at each stage, up to extensional equivalence \emph{for writing values}.
In the proofs, this turns out to be all we need.

In fact, we can simplify this just to 
$$
\hspace{-6em}\rulefont{SPaxWrite01}
\qquad\qquad
\texiaffine\someoneAll\tf{write} .
$$
In full, this is $\texiaffine v.\someoneAll\tf{write}(v)$, and in words it expresses: 
\begin{quote}
`at each stage, at most one value is written'.
\end{quote}
In the presence of \rulefont{PaxWrite?}, \rulefont{SPaxWrite01} is implied by \rulefont{LdrExt'} --- but it does not imply \rulefont{LdrExt'}, because it permits the possibility that one leader participant writes a value and another leader participant crashes.
The reader can check that \rulefont{SPaxWrite01} is sufficient to complete the proof of Lemma~\ref{lemm.various.unique}(\ref{item.various.unique.write}).

Note also that that \rulefont{SPaxWrite01} implies \rulefont{PaxWrite01}, so that --- for the proofs we care about --- \rulefont{SPaxWrite01} can replace both \rulefont{LdrExt} and \rulefont{PaxWrite01}. 
We will do exactly this, later: the job of $\Theta_\ThyPax$ is to be `declarative Paxos', that is, its purpose is to be declarative \emph{and} to be recognisably Paxos. 
Later on in Figure~\ref{fig.ThySPax} we will define a \emph{simpler} declarative Paxos in (whence the name `SPaxWrite01').
This simplified theory is derived by looking at the proofs here, and seeing what we actually require.










 



\subsubsection{Agreement at possibly different times}



\begin{prop}
\label{prop.integrity.helper}
Suppose $\avaluation$ is a valuation such that $\mentval\ThyPax$, and suppose $v,v'\in\Val$ and $\avaluation$. 
Then:
\begin{enumerate*}
\item\label{item.integrity.helper.accept}
$\mentval\recent\QuorumBox\tf{accept}(v) \timpc \someoneAll\tf{accept}(v')\timpc v'\teq v$.
\item\label{item.integrity.helper.write}
$\mentval\recent\QuorumBox\tf{accept}(v) \timpc \someoneAll\tf{write}(v')\timpc v'\teq v$.
\end{enumerate*}
\end{prop}
\begin{proof}
Suppose $n\in\Time$.
We will reason using \strongmodusponens, so suppose 
\begin{equation}
\label{eq.n.1}
n\mentval\modT\QuorumBox\tf{accept}(v). 
\end{equation}
We will prove by induction on $n'>n$ that for every $v'\in\Val$,
$$
\begin{array}{r@{\ }l}
n'\mentval&\modT(\mru{(\someoneAll\tf{accept})}{v'}\timpc v'\teq v)
\quad\text{and}\quad
\\
n'\mentval&\modT(\someoneAll\tf{write}(v')\timpc v'\teq v)
\end{array}
$$
Clearly, by \strongmodusponens the $\tf{write}$ part of the inductive hypothesis is enough to prove part~\ref{item.integrity.helper.write} of this result.
For the $\tf{accept}$ part of the inductive hypothesis, we see that this implies part~\ref{item.integrity.helper.accept} of this result either by a routine argument using \rulefont{PaxAccept?} and part~\ref{item.integrity.helper.write} of this result --- or using Corollary~\ref{corr.mru.today.to.yesterday}(\ref{item.mru.today.implies.tomorrow}) (for $a.\phi=a.\someoneAll\tf{accept}(a)$ in that result).

We consider cases:
\begin{itemize}
\item
\emph{The base case for $\tf{accept}$, where $n'=n+1$.}

Suppose $n'\mentval\modT(\mru{\someoneAll\tf{accept}}{v'})$.

We assumed in equation~\eqref{eq.n.1} that $n\mentval\modT\QuorumBox\tf{accept}(v)$, so by Proposition~\ref{prop.someone.implies.someoneAll}(\ref{item.someone.implies.someoneAll.2}) also $n\mentval\modT\someoneAll\tf{accept}(v)$, and by Corollary~\ref{corr.mru.today.to.yesterday}(\ref{item.mru.today.implies.tomorrow}) $n'\mentval\modT(\mru{(\someoneAll\tf{accept})}{v})$. 
By Lemma~\ref{lemm.various.unique}(\ref{item.various.unique.recent.accept}) $v'=v$.

\item
\emph{The inductive step for $\tf{accept}$, where $n'>n+1$.}

Suppose $n'\mentval\modT(\mru{\someoneAll\tf{accept}}{v'})$.
By Corollary~\ref{corr.mru.today.to.yesterday}(\ref{item.mru.tomorrow.options}) either $n'\minus 1\mentval\modT\someoneAll\tf{accept}(v')$ or $n'\minus 1\mentval\modT(\mru{\someoneAll\tf{accept}}{v'})$.

In the former case, using \rulefont{PaxAccept?} $n'\minus 1\mentval\someoneAll\tf{write}(v')$, and by the $\tf{write}$ part of the inductive hypothesis, $v'=v$.
In the latter case, we just use the $\tf{accept}$ part of the inductive hypothesis. 
\item
\emph{The case of $\tf{write}$.}

Suppose $p\in\Pnt$ and $n',p\mentval\modT\tf{write}(v')$.
Using \rulefont{PaxWrite?}, 
there exists an $O\in\opensne$ such that 
$$
n',O\mentval \modT\bigl( 
(\mru{\someone\tf{accept}}{v})\ \tor
(\tneg\modT\recent\texi\someone\tf{accept}\tand \tf{propose}(v)) \bigr),
$$
from which it follows using Lemmas~\ref{lemm.commuting.connectives} and~\ref{lemm.tand.tor} and some basic reasoning on truth-values that 
$$
\text{either}\quad
n',O\mentval \modT\mru{\someone\tf{accept}}{v}
\quad\text{or}\quad
n',O\mentval \tneg\modT\recent\texi\someone\tf{accept} .
$$
Now $n<n'$ and we assumed in equation~\eqref{eq.n.1} that $n\mentval\modT\QuorumBox\tf{accept}(v)$, so using 
\rulefont{Pax2Twined} and \strongmodusponens $n\mentval\modT\CoquorumDiamond\tf{accept}(v)$, and in particular (by routine reasoning on the clauses for $\CoquorumDiamond$ in Figure~\ref{fig.3.derived} and for $\someone$ in Figure~\ref{fig.3.phi.f})
$n,O\mentval\modT\someone\tf{accept}(v)$.
Thus $n',O\mentval\modT\texi\someone\tf{accept}$ and by Lemma~\ref{lemm.tomorrow.forever.recent}(\ref{item.tomorrow.forever.recent.1}) (taking $\phi=\texi\someone\tf{accept}$ in that Lemma) 
$n',O\mentval\modT\recent\texi\someone\tf{accept}$.
It follows that we are not in the right-hand (`or') branch above, so we must be in the left-hand (`either') branch, and thus $n',O\mentval\modT\mru{\someone\tf{accept}}{v}$.

By the $\tf{accept}$ part of the inductive hypothesis, $v'=v$.
\qedhere\end{itemize}

\end{proof}


\begin{thrm}[Agreement]
\label{thrm.full.agreement}
We have:
$$
\ThyPax\ment\tall v,v'.(\urecent\someoneAll\tf{decide}(v) \timpc \someoneAll\tf{decide}(v') \timpc v'\teq v).
$$
\end{thrm}
\begin{proof}
Suppose $v,v'\in\Val$ and $n\in\Time$ and $\avaluation$ is a valuation.
We use \strongmodusponens, so suppose
$$
n\mentval\modT\urecent\someoneAll\tf{decide}(v)
\quad\text{and}\quad
n\mentval\modT\someoneAll\tf{decide}(v') .
$$
If $n\mentval\modT\someoneAll\tf{decide}(v)$ then we use Proposition~\ref{prop.decide.n.agree}.

So suppose $n\mentval\modT\recent\someoneAll\tf{decide}(v)$.
By \rulefont{PaxDecideL?} (and \rulefont{PaxDecide\tneg L?} if needed), $n\mentval\modT\recent\QuorumBox\tf{accept}(v)$ and $n\mentval\modT\QuorumBox\tf{accept}(v')$, so that by Proposition~\ref{prop.someone.implies.someoneAll}(\ref{item.someone.implies.someoneAll.2}) also %
$n\mentval\modT\someoneAll\tf{accept}(v')$.
By Proposition~\ref{prop.integrity.helper}(\ref{item.integrity.helper.accept}) $v'=v$ as required.
\end{proof}

\subsection{Termination}
\label{subsect.pax.termination}

\begin{rmrk}
\emph{Termination} is expressed in English as:
\begin{quote}
\emph{Every correct participant eventually decides some value.}
\end{quote}
We render this in our logic in Theorem~\ref{thrm.termination} as 
$$
\ThyPax\ment\infinitely\everyoneAll\texi v.\tf{decide}(v).
$$
In words: 
\begin{quote}
There is always some future time when all participants decide.
\end{quote}
Note that there is no need to explicitly specify `all correct/uncrashed participants' above.
This is taken care of by our logic's in-built notion of validity, since incorrect participants return $\tvB$, which is valid as per Remark~\ref{rmrk.motivate.validity}(\ref{item.tb.valid}) (but not true or correct; cf. Remark~\ref{rmrk.motivate.validity}).
\end{rmrk}


\subsubsection{Global stabilisation of logical time (GSLT)}
\label{subsect.gslt}

Our forward rules are wrapped in a $\final$ modality, to reflect an assumption that while there may be some initial network instability, if we wait long enough the network will stabilise (see the discussion in Remark~\ref{rmrk.discuss.forward.rules}).
This leads us to the notion of \emph{global stabilisation of logical time}, or \emph{GSLT} for short: 
\begin{defn}[GSLT]
\label{defn.gslt}
Suppose $\phi$ is a closed predicate and $\avaluation$ is a valuation.
Then:
\begin{enumerate*}
\item\label{item.gslt.synchronous}
Say that \ $\final\phi$\ (i.e. $\phi$ under a $\final$ modality) is \deffont{synchronous at $n$ in $\avaluation$} when $n\mentval\forever\phi$.\footnote{%
By Definition~\ref{defn.validity.judgement}, $n\ment\final\phi$ means that $n,p,O\mentval\final\phi$ for every $p\in\Pnt$ and $O\in\opens$ and valuation $\varsigma$.} %
That is:
$$
\final\phi\ \ \text{ is synchronous at $n$ in $\avaluation$}
\quad\text{when}\quad
n\mentval\forever\phi .
$$
In words, $\final\phi$ is synchronous when the `$\final$' becomes a `$\forever$'.
\item\label{item.gslt.gslt}
Define $\gslt(\avaluation)\in\Time$ to be the least stage $n$ such that all of the rules in Figure~\ref{fig.3.phi.f} that start with $\final$ --- namely, the forward rules \rulefont{PaxSend!}, \rulefont{PaxAccept!}, \rulefont{PaxDecideL!}, and \rulefont{PaxDecide\tneg L!} --- are synchronous at $n$ in the valuation $\avaluation$.

We call $\gslt(\avaluation)$ the \deffont{Global Stabilisation of Logical Time} in $\avaluation$. 
\end{enumerate*}
\end{defn}

\begin{rmrk}
We make some observations about Definition~\ref{defn.gslt}:
\begin{enumerate}
\item
Unpacking Definition~\ref{defn.gslt}(\ref{item.gslt.synchronous}) using the clause for $\forever$ in Figure~\ref{fig.3.derived}, 
$$
\final\phi \text{ is synchronous at $n$ in $\avaluation$}
\quad\liff\quad
n\mentval\forever\phi 
\quad\liff\quad
\Forall{n'{\geq}n}(n'\mentval\phi),
$$
\noindent
and $\gslt(\avaluation)$ is least such that $\gslt(\avaluation)\mentval\forever(\someoneAll\texi\tf{propose} \tnotor \texi\tf{send})$, and $\gslt(\avaluation)\mentval\forever\bigl((\tf{leader}\tand \QuorumBox\texi\tf{send}) \tnotor \texi\tf{write}\bigr)$, and so on.
\item  
$\gslt$ is well-defined in Definition~\ref{defn.gslt}(\ref{item.gslt.gslt}) because there are finitely many axioms, mentioning finitely many $\final$ modalities, and a finite set of natural numbers has a finite least upper bound.
This statement has real content, e.g. if we had written \rulefont{PaxAccept!} as $\tall v.\final(\someoneAll\tf{write}(v) \tnotor \tf{accept}(v))$ --- which is a subtly weaker assertion than the form used in Figure~\ref{fig.logical.paxos} --- then $\gslt$ would \emph{not} be defined. 
\end{enumerate}
\end{rmrk}




\begin{lemm}
\label{lemm.someone.writes.after.gslt.accept}
Suppose $\avaluation$ is a valuation such that $\mentval\ThyPax$. 
Suppose $n\in\Time$ and $n\geq\gslt(\avaluation)$ and $v\in\Val$. 
Then 
\begin{multline*}
n\mentval\modT\someoneAll\tf{write}(v) 
\liff
n\mentval\everyoneAll\tf{accept}(v) 
\liff
\\
n\mentval\modT\QuorumBox\tf{accept}(v) 
\liff
n\mentval\modT\someoneAll\tf{accept}(v)  .
\end{multline*}
\end{lemm}
\begin{proof}
We prove a cycle of implications, freely using weak and strong modus ponens (Proposition~\ref{prop.mp.for.tnotor}(\ref{item.mp.for.tnotor}\&\ref{item.mp.for.timpc})).

Suppose $n\mentval\modT\someoneAll\tf{write}(v)$. 
By \rulefont{PaxAccept!} (since $n\geq\gslt(\avaluation)$), $\Forall{p\in\Pnt}n\mentval\tf{accept}(v)$.
By Lemma~\ref{lemm.atomic.O}(\ref{item.atomic.pointwise}) $\tf{accept}(v)$ is pointwise, so by Proposition~\ref{prop.pointwise.dense.char}(\ref{item.dense.char.pointwise.everyoneAll}) $n\mentval\everyoneAll\texi\tf{accept}(v)$.


By Proposition~\ref{prop.correct.cobox.boxT}(\ref{item.correct.cobox.boxT.1}) and \rulefont{PaxPCorrect} (for $\tf{accept}$),
$n\mentval\modT\QuorumBox\tf{accept}(v)$, and by Proposition~\ref{prop.someone.implies.someoneAll}(\ref{item.someone.implies.someoneAll.2}) $n\mentval\modT\someoneAll\tf{accept}(v)$.

Now suppose $n\mentval\modT\someoneAll\tf{accept}(v)$.
By \rulefont{PaxAccept?} %
$n\mentval\modT\someoneAll\tf{write}(v)$.  
\end{proof}


\begin{lemm}
\label{lemm.leaderlive.quorumsend}
Suppose $\avaluation$ is a valuation such that $\mentval\ThyPax$. 
Suppose $n\in\Time$ and $n\geq\gslt(\avaluation)$. 
Then:
\begin{enumerate*}
\item\label{item.leaderlive.quorumsend.1}
$n\mentval\someoneAll\texi\tf{propose}\tnotor\everyoneAll\texi\tf{send}$.
\item\label{item.leaderlive.quorumsend.2}
$n\mentval\someoneAll\texi\tf{propose}\timpc\QuorumBox\texi\tf{send}$.
\end{enumerate*}
\end{lemm}
\begin{proof}
We consider each part in turn:
\begin{enumerate*}
\item
We use \weakmodusponens.
Suppose $n\mentval\modT\someoneAll\texi\tf{propose}$.
By \rulefont{PaxSend!} (since $n\geq\gslt(\avaluation)$) 
$$
\Forall{p\in\Pnt}n,p\mentval\texi\tf{send}.
$$
By Lemma~\ref{lemm.atomic.O}(\ref{item.atomic.pointwise.2}) $\texi\tf{send}$ is pointwise, so by Proposition~\ref{prop.pointwise.dense.char}(\ref{item.dense.char.pointwise.everyoneAll}) $n\mentval\everyoneAll\texi\tf{send}$ as required.
\item
We use \strongmodusponens.
Suppose $n\mentval\modT\someoneAll\texi\tf{propose}$.
By part~\ref{item.leaderlive.quorumsend.1} of this result and \weakmodusponens, $n\mentval\everyoneAll\texi\tf{send}$.
By \rulefont{PaxPCorrect} (for $\tf{send}$) $\acontext\ment\QuorumBox\correct{\tf{send}}$, so by 
Proposition~\ref{prop.correct.cobox.boxT}(\ref{item.correct.cobox.boxT.2})
$\acontext\ment\modT\QuorumBox\texi\tf{send}$ as required. 
\qedhere\end{enumerate*}
\end{proof}




\subsubsection{GSLT implies termination}


Recall from Definition~\ref{defn.correct}(\ref{item.correct}) that $\correct{\tf{decide}}=\tall v.\modTF\tf{decide}(v)$. 
\begin{lemm}
\label{lemm.someone.propose.someone.write}
Suppose $\avaluation$ is a valuation such that $\mentval\ThyPax$. 
Suppose $n\in\Time$ and $n\geq\gslt(\avaluation)$, and $p\in\Pnt$.
Then:
\begin{enumerate*}
\item\label{item.someone.propose.someone.write}
If $n,p\mentval\modT\texi\tf{propose}$ and $n,p\mentval\correct{\tf{write}}$ then $n,p\mentval\modT(\tf{leader}\tand\texi\tf{write})$.
\item\label{item.someone.propose.someone.write.live}
$\mentval\infinitely(\correct{\tf{decide}}\tand\modT\someoneAll\texi\tf{write})$.  
\item\label{item.someone.propose.someone.accept}
$\mentval\infinitely(\someoneAll(\tf{leader}\tand\correct{\tf{decide}})\tand\modT\texi\QuorumBox\tf{accept})$.
\end{enumerate*}
\end{lemm}
\begin{proof}
We consider each part in turn:
\begin{enumerate}
\item
Suppose $n,p\mentval\modT\texi\tf{propose}$.
By \rulefont{PaxPropose?} and \strongmodusponens $n,p\mentval\modT\tf{leader}$.
By Lemma~\ref{lemm.leaderlive.quorumsend}(\ref{item.leaderlive.quorumsend.2}) and strong modus ponens $n\mentval\modT\QuorumBox\texi\tf{send}$.
Using \weakmodusponens and \rulefont{PaxWrite!} (since $n\geq\gslt(\avaluation)$) $n,p\mentval\texi\tf{write}$.
By Proposition~\ref{prop.tv.ment.TF.model}(\ref{item.modTF.texi.P.valid}) (since $n,p\mentval\correct{\tf{write}}$) $n,p\mentval\modT\texi\tf{write}$. 
\item
By \rulefont{LdrExist}, and \rulefont{LdrCorrect} with Lemma~\ref{lemm.persist.char} (for $\f{tvs}=\{\tvT\}$ in that Lemma), and \strongmodusponens, 
for infinitely many $n\in\Time$ there exists $l_n\in\Pnt$ such that 
$$
n,l_n\ment\modT\tf{leader}\tand\correct{\tf{propose},\tf{write},\tf{decide}}. 
$$
So for each $n$ and $l_n$ we have $n,l_n\mentval\correct{\tf{decide}}$.
We will now check that $n,l_n\mentval\modT\texi\tf{write}$.
We reason as follows:
\begin{itemize*}
\item
By \rulefont{PaxPropose!} and \weakmodusponens we have that $n,l_n\ment\texi\tf{propose}$.
\item
By Proposition~\ref{prop.tv.ment.TF.model}(\ref{item.modTF.texi.P.valid}) (since $n,l_n\mentval\correct{\tf{propose}}$ and $n,l_n\mentval\texi\tf{propose}$) $n,l_n\mentval\modT\texi\tf{propose}$. 
\item
By part~\ref{item.someone.propose.someone.write} of this result (since $n,l_n\mentval\correct{\tf{write}}$) $n,l_n\mentval\modT\texi\tf{write}$.
\end{itemize*}
\item
We just combine part~\ref{item.someone.propose.someone.write.live} of this result with Lemma~\ref{lemm.someone.writes.after.gslt.accept}.
\qedhere\end{enumerate}
\end{proof}


Recall that Theorem~\ref{thrm.termination} expresses termination (as the name suggests) but also an abstracted form of the Integrity correctness property by virtue of the use of the $\infinitely$ modality, as per Remark~\ref{rmrk.paxos.integrity}:
\begin{thrm}[Termination]
\label{thrm.termination}
We have:
$$
\ThyPax\ment\infinitely\everyoneAll\texi\tf{decide} .
$$ 
\end{thrm}
\begin{proof}
Suppose $\avaluation$ is a valuation and $\mentval\ThyPax$.
By Lemmas~\ref{lemm.someone.propose.someone.write}(\ref{item.someone.propose.someone.accept}) and~\ref{lemm.persist.char} (for $\f{tvs}=\{\tvT\}$ in that Lemma), for infinitely many $n\in\Time$ there exists an $l_n\in\Pnt$ such that 
\begin{enumerate*}
\item
$n,l_n\mentval\tf{leader}$,
\item
$n,l_n\mentval\correct{\tf{decide}}$ (Definition~\ref{defn.correct}(\ref{item.correct})), and
\item
$n\mentval\modT\texi\QuorumBox\tf{accept}$.
\end{enumerate*}
By \rulefont{PaxDecideL!} (since $n\geq\gslt(\avaluation)$) 
and \weakmodusponens $n,l_n\mentval\texi\tf{decide}$.
By Proposition~\ref{prop.tv.ment.TF.model}(\ref{item.modTF.texi.P.valid}) (since $n,l_n\mentval\correct{\tf{decide}}$ and $n,l_n\mentval\texi\tf{decide}$) $n,l_n\mentval\modT\texi\tf{decide}$. 
Then by \rulefont{PaxDecide\tneg L!} and \weakmodusponens also $n,p\mentval\texi\tf{decide}$ for every $p\in\Pnt$, and so by (Lemma~\ref{lemm.atomic.O}(\ref{item.atomic.pointwise.2}) and) Proposition~\ref{prop.pointwise.dense.char}(\ref{item.dense.char.pointwise.everyoneAll}) $n\mentval\everyoneAll\texi\tf{decide}$. 

The above holds for infinitely many $n\in\Time$, so by Lemma~\ref{lemm.persist.char} (for $\f{tvs}=\{\tvT,\tvB\}$ in that Lemma) we have that $\mentval\infinitely\everyoneAll\texi\tf{decide}$ as required. 
\end{proof}


\begin{rmrk}
Theorem~\ref{thrm.termination} proves that $\ThyPax\ment\infinitely\everyoneAll\texi\tf{decide}$ holds.
Note that $\tvB$ (the truth-value associated to a crashed process) is a valid truth-value, so $n\ment\everyoneAll\texi\tf{decide}$ can hold in the trivial case that everyone is crashed.
Thus, what Theorem~\ref{thrm.termination} asserts in English is that infinitely often, everyone who is not crashed, decides (but it might be that everyone is crashed) --- from which it obviously follows that if a participant never crashes then eventually it decides a value, as per Remark~\ref{rmrk.standard.paxos.correctness}.

But in fact, we can prove a bit more:
$$
\ThyPax\ment\infinitely\QuorumBox\modT\texi\tf{decide}
$$
In words: infinitely often, there exists a \emph{quorum of uncrashed participants}, who decide.

This follows from Theorem~\ref{thrm.termination} by easy reasoning using \rulefont{PaxPCorrect}, \rulefont{PaxDecide\tneg L!}, and Proposition~\ref{prop.decide.n.agree}; we leave this to the reader to check. 
\end{rmrk}


\section{Simpler Declarative Paxos}
\label{sect.simpler.paxos}

Logic is useful for simplification and abstraction in many ways.
In particular, we can axiomatise something, prove some properties --- and then look at what we actually needed from the axioms to make the proofs work.
This can feed back to suggest refinements and generalisations of the original axioms.
(For example: many properties of natural numbers are actually properties of fields, or rings, or integral domains, and so on.) 

We can do something similar here, and look for the most elementary set of axioms that are actually required for our correctness proofs as written: 

\subsection{What we used in the correctness proofs for \ThyPax}
\label{subsect.reflection.on.correctness}
 
\begin{rmrk}
\label{rmrk.spax.explanation}
We note a curious thing about our correctness proofs in Section~\ref{sect.paxos.correctness.properties}, that we never used axioms \rulefont{PaxSend?} or \rulefont{PaxPropose01}.

How can this be?
Is this an error?

No.
This is an opportunity to further simplify the Declarative Paxos from Figure~\ref{fig.ThyPaxOne} to the \emph{Simpler Declarative Paxos} in Figure~\ref{fig.ThySPax}.
To understand how and why this works, we need to highlight some details about how the Paxos algorithm works.
\end{rmrk}

\begin{rmrk}[Why we did not use \rulefont{PaxSend?}]
Recall from Remark~\ref{rmrk.high-level.paxos} our high-level view of the Paxos algorithm, and recall that \ThyPax in Figure~\ref{fig.ThyPaxOne} is a declarative version of that algorithm (obtained by abstracting away explicit message-passing and algorithmic time, as per Remark~\ref{rmrk.elisions}).

Note of the algorithm that 
\begin{itemize*}
\item
in step~\ref{item.l.proposes} of Remark~\ref{rmrk.high-level.paxos}, the leader broadcasts a `propose' message; and 
\item
in step~\ref{item.p.sends}, a participant $p$ that receives that `propose' message, replies with a `send' message carrying as its payload a value $v_p$, which is $p$'s most recently accepted value (if any); and
\item
in step~\ref{item.l.writes}, the leader gathers the $v_p$ values from these responses and, if it has a quorum (in topological language: an open set) of responses, it selects a most recent value to use in its write step.
\end{itemize*}
In the algorithm, the only way that the leader $l$ can know of a participant $p$'s value of $v_p$, is if $l$ receives a send message from $p$ with payload $v_p$.
If $l$ receives no such message, then it does not know the value of $v_p$.

In contrast, predicates of \QLogic are not assigned truth-values by participants in an algorithm; they are assigned truth-values as if by an omniscient observer who is standing outside of the model and who has full (omniscient) access to it.
This is reflected in \rulefont{PaxWrite?}, where we write $\mru{(\someone\tf{accept})}{v}$.
This directly examines the model to find the most recently accepted value.
No reference here is required to the payload of the quorum of $\tf{send}$ messages.
\rulefont{PaxWrite?} still insists that such a quorum exists, as per the $\everyone\texi\tf{send}$ that appears in \rulefont{PaxWrite?}, so in this case the value $\mru{(\someone\tf{accept})}{v}$ is computable by the leader, and we could specify a more complex predicate that checks the set of $v_p$ payloads in the send messages of that quorum to select a most recently-accepted value.
In the implementation, this is necessary because the leader $l$ is not omniscient.
But in the logic we have no reason to go to this trouble, because we can just read the information directly from the model using $\mru{(\someone\tf{accept})}{v}$.\footnote{For example: if we want to compute $100!$ in the real world, we have to build a computer to do this for us, which involves many practical details including the construction of a semiconductor industry, compiler design and so on.  In logic, however, we can just say that it is $\prod_{i=1}^{100}i$, because that is what $100!$ \emph{is}.}

So our axioms are just doing their job, which is to be a declarative abstraction of the implementation. 
But since we do not use \rulefont{PaxSend?}, this is evidence that we may treat the $\tf{send}$ step as an implementational detail having to do with how information moves around the network, and elide it.
\end{rmrk}

\begin{rmrk}[Why we did not use \rulefont{PaxPropose01}]
This is due to another particularity of how the Paxos algorithm works.
The leader chooses one value in the propose step (or no value, if the leader has crashed), but it only uses that value in the case that it finds a quorum of participants who have never previously accepted any value (this is the right-hand disjunct in \rulefont{PaxWrite?}).
Just in this case, it writes the value that it chose in the propose step.

So what really matters is that the leader \emph{only writes one value}.
If the reader had proposed two values, this would be fine --- provided that it only writes one of them.

Leaders in the Paxos algorithm do not bother with multiple values in the propose step, because this would be a waste of time; they know they will only ever need at most one.
Thus, \rulefont{PaxPropose01} is right to reflect a true fact of the algorithm that it does not propose any more than the value than it will need, but this an algorithmic efficiency not a logical necessity. 
Likewise, the logic is right to point this out to us by \emph{not} requiring \rulefont{PaxPropose01} in the proofs: so long as only at most one value gets \emph{written} --- which is handled by \rulefont{PaxWrite01} --- none of our correctness properties are affected.
\end{rmrk}

\begin{rmrk}
By examining the proofs further, some further possible simplifications are revealed:
\begin{enumerate*}
\item
We can combine rules \rulefont{SPaxWrite01} and \rulefont{PaxWrite01} and \rulefont{LdrExt}, as per the discussion in Subsection~\ref{subsect.ldrexistuniq'}, to obtain an elegantly concise rule \rulefont{SPaxWrite01}.
\item
We used \rulefont{PaxPCorrect} in our proofs to obtain $\QuorumBox\correct{\tf P}$ only twice: 
\begin{itemize*}
\item
in Lemma~\ref{lemm.someone.writes.after.gslt.accept} for $\tf P=\tf{accept}$, and 
\item
in Lemma~\ref{lemm.leaderlive.quorumsend} for $\tf P=\tf{send}$.
\end{itemize*}
We have elided the \emph{send} step in $\Theta_\ThySPax$, so we only need to consider $\tf{accept}$.
Accordingly we take take \rulefont{SPaxPCorrect} to be just $\QuorumBox\correct{\tf{accept}}$.
\end{enumerate*}
\end{rmrk}

\begin{rmrk}
In conclusion, we see that the proofs in Section~\ref{sect.paxos.correctness.properties} are not only correctness proofs, but can also be read as an active commentary on the axioms.
Our logic has done what logic does, and has separated what is logically necessary from what is implementationally required or implementationally convenient.
This shows us how to simplify the axiomatisation (if we we wish) further away from the implementation, and closer towards what is logically essential: 
\end{rmrk}


\subsection{The definition}

\begin{figure}
$$
\begin{array}{l@{\quad}r@{\ }l}
\figunderline{Backward rules, for backward inference} 
\figskip
\rulefont{SPaxPropose?}&
& \tall v.\ \tf{propose}(v)\timpc\tf{leader}
\figskip
\rulefont{SPaxWrite?}&
& \tall v.\ \tf{write}(v) \timpc 
\begin{array}[t]{l}
\tf{leader}\tand\Quorum( 
\\
\quad\bigl((\mru{(\someone\tf{accept})}{v}) \tor
(\tf{propose}(v)\tand\tneg\modT\recent\texi\someone\tf{accept})\bigr)) 
\end{array}
\figskip
\rulefont{SPaxAccept?}&
& \tall v.\ \tf{accept}(v) \timpc \someoneAll\tf{write}(v)
\figskip
\rulefont{SPaxDecide?}&
& \tall v.\ \tf{decide}(v) \timpc \QuorumBox\tf{accept}(v)
\figskip\figskip
\figunderline{Forward rules, for making progress} 
\figskip
\rulefont{SPaxWrite!}&
& \final(\tf{leader}\tnotor\texi\tf{write})
\figskip
\rulefont{SPaxAccept!}&
& \final(\texi\someoneAll\tf{write}\tnotor \texi\tf{accept})
\figskip
\rulefont{SPaxDecide!}&
& \final(\texi\QuorumBox\tf{accept} \tnotor \texi\tf{decide})
\figskip
\figskip
\figunderline{Other rules, for correctness \& liveness}
\figskip
\rulefont{SLdrExist}&
&\modTF\tf{leader}\tand \modT\someoneAll\tf{leader}
\figskip
\rulefont{SLdrCorrect}&
& \infinitely\someoneAll(\tf{leader}\tand \correct{\tf{propose},\tf{write},\tf{decide}}) 
\figskip
\rulefont{SPaxPCorrect}&
& \QuorumBox\correct{\tf{accept}}
\figskip
\rulefont{SPaxWrite01}&
& \texiaffine\someoneAll\tf{write}
\end{array}
$$
\caption{Simpler Declarative Paxos (Definition~\ref{defn.thyspax})}
\label{fig.logical.paxos'}
\label{fig.ThySPax}
\end{figure}


\begin{defn}
\label{defn.thyspax}
We define \deffont{Simpler Declarative Paxos} to be the logical theory 
$$
\ThySPax=(\Sigma_{\ThySPax},\Theta_{\ThySPax}),
$$ 
where:
\begin{enumerate*}
\item
the signature $\Sigma_{\ThySPax}$ is 
$$
\Sigma_{\ThySPax}=[\tf{leader}:0;\ \tf{propose},\tf{write},\tf{accept},\tf{decide}:1] ,
$$ 
as per the notation in Definition~\ref{defn.signature}(\ref{item.predicate.syntax.signature}), and
\item
the axioms $\Theta_{\ThySPax}$ are as written in Figure~\ref{fig.ThySPax}.
\end{enumerate*}
\end{defn}

The same validity properties hold of \ThySPax from Figure~\ref{fig.ThySPax} as hold for \ThyPax from Figure~\ref{fig.ThyPaxOne}, namely: Validity (Theorem~\ref{thrm.validity}), Agreement (Theorem~\ref{thrm.full.agreement}), and Termination and Integrity (Theorem~\ref{thrm.termination}):
\begin{thrm}
\label{thrm.simpler.dp}
Suppose $\Sigma$ is a signature that includes $\Sigma_\ThySPax$, and suppose $\mathcal M$ is a model.
Then we have:
\begin{enumerate*}
\item
$\ThyPax\ment\tall v.(\tf{decide}(v) \timpc \urecent\someoneAll(\tf{leader}\tand\tf{propose}(v)))$.
\item
$\ThyPax\ment\tall v,v'.(\urecent\someoneAll\tf{decide}(v) \timpc \someoneAll\tf{decide}(v') \timpc v'\teq v)$.
\item
$\ThyPax\ment\infinitely\everyoneAll\texi\tf{decide}$.
\end{enumerate*}
\end{thrm}
\begin{proof}
This is by construction: from the discussion in Subsection~\ref{subsect.reflection.on.correctness}, the theory \rulefont{ThySPax} consists of just those parts of \rulefont{ThyPax} that are actually required in the correctness proofs.
\end{proof}

Thus, Theorem~\ref{thrm.simpler.dp} just notes that Simpler Declarative Paxos extracts from the proofs in Section~\ref{sect.paxos.correctness.properties} the logical essence of Figure~\ref{fig.ThyPaxOne}, as applied to the standard correctness properties of Paxos.
There is no right or wrong about whether we prefer $\ThyPax$ or $\ThySPax$; it depends on what level of abstraction we want to work with at any particular moment.
But, implicit in Theorem~\ref{thrm.simpler.dp} is an observation that \ThySPax is \emph{minimal}, because we put into it only that which we know we need to make some specific correctness proofs work.


\section{Conclusions and future work}
\label{sect.conclusions}

\subsection{Conclusions}







We have presented \QLogic, a three-valued modal fixedpoint logic for declarative specification of consensus algorithms, and we have shown how to use this to axiomatise Paxos, a canonical and nontrivial (though still relatively straightforward) consensus algorithm that can handle benign failures of participants.

We do not presented more complex consensus algorithms in this paper but we have looked into these in preliminary investigations, and they do seem to work. 

This paper is aimed at 
\begin{itemize*}
\item
logicians who might be interested in exploring a logic with a new and unusual combination of features inspired by strong practical motivations, and 
\item
researchers in distributed systems, who might consider (and benefit from) adopting such a framework.
\end{itemize*}
\QLogic offers benefits that could make it a useful addition to the toolkit of anyone working in distributed systems:
\begin{enumerate}
\item 
Much of what is shown in this paper just reflects existing practices, formalised into a new logic.
Therefore, readers may already be reasoning in ways that align with our logic, without explicitly realising it. 
For example, the third truth-value ($\tvB$) in \QLogic corresponds to typical side conditions like ``for every correct participant,'' and the backward axioms mirror the reverse inferences that practitioners often use. 

By formalising these familiar patterns, \QLogic provides a structured and precise way to represent and reason about concepts that are already informally expressed in the literature.
\item 
Logic is a concise, powerful, and unambiguous language. 
It distills paragraphs of English prose into a single line of precise symbols, enabling clear and rigorous reasoning. 

While explanatory English is helpful for accessibility and context, formal logic is the superior tool for working through the details of complex algorithms. 
It ensures precision and eliminates the ambiguities that often arise in informal descriptions.
\item 
\QLogic operates at a new (and higher) level of abstraction, which complements existing logic-based techniques. 
TLA+~\cite{DBLP:journals/toplas/Lamport94} provides a state-based view of distributed algorithms, whereas \QLogic abstracts lower-level implementation details while retaining the essential behaviours and properties of the algorithms.
This abstraction does not obscure the logic of the system; on the contrary we have found that it clarifies it, making the underlying mechanisms and invariants more apparent. 
\item 
One strength of logic is its ability to abstract away what is irrelevant while maintaining precision about what is important. 
In \QLogic, constructs like strong and weak implications ($\timpc$ and $\tnotor$) exemplify this principle, capturing subtle distinctions in meaning with precise formal symbols. 
These constructs distill complex ideas into simple yet powerful tools that can be manipulated symbolically and reasoned about systematically. 

While mastering these constructs requires some initial effort, the benefits of precision and clarity justify the investment.
\item 
Learning \QLogic requires some effort, but for researchers already familiar with (for example) TLA+, the transition to \QLogic should be eminently manageable. 

While axiomatisations in \QLogic are certainly subtle, this is not gratuitous: they reflect the inherent nature of the distributed systems they are designed to model. 
By formalising these complexities, \QLogic captures their structure and reveals and addresses the challenges that are \emph{already present} in the underlying algorithms. 

This transparency makes it a powerful tool for understanding and reasoning about distributed protocols.\footnote{Certainly, the first author (who has a background in logic) has found this study invaluable for gaining a better understanding of consensus algorithms.  It seems likely that other readers might benefit similarly from using \QLogic.}
\item 
Distributed algorithms, and consensus protocols in particular, are foundational to the functionality of blockchain systems and other critical infrastructures. 

These algorithms are recognised to be notoriously challenging to design and implement correctly. 
Ambiguities and informal specifications can lead to vulnerabilities and failures. 

\QLogic bridges the gap between declarative specifications of invariants, and practical verification; enabling researchers to clarify, reason about, and verify the properties of distributed protocols rigorously. 
This is especially valuable in blockchain systems, where clear and robust specifications are essential for trust and reliability.
\end{enumerate}

In conclusion: by providing a formal framework that combines the abstraction of declarative reasoning with the rigour of logical precision, \QLogic offers researchers in distributed systems a powerful tool for modeling, understanding, and verifying complex algorithms. 

\subsection{Related and future work}

In this paper, our focus has been to lay the foundational theory and framework of \QLogic and to demonstrate its application to one of the most well-known consensus protocols, Paxos~\cite{DBLP:journals/tocs/Lamport98}.

Paxos is a sensible starting point, being a cornerstone of distributed systems, but (certainly by modern standards!) it is relatively straightforward.
For example, Paxos only addresses crash failures, limiting its applicability to systems that assume benign (rather than hostile) faults.

A natural direction for future work involves extending our framework to handle more complex and realistic fault models, such as Byzantine faults, where participants can deviate arbitrarily from the algorithm.
Consensus protocols that address Byzantine faults, such as Practical Byzantine Fault Tolerance (PBFT)~\cite{DBLP:conf/osdi/CastroL99} and HotStuff~\cite{DBLP:conf/podc/YinMRGA19}, provide robust mechanisms for fault tolerance and serve as the foundation for many modern blockchain systems, including Tendermint (\url{https://tendermint.com})~\cite{DBLP:journals/corr/abs-1807-04938}. 

Applying \QLogic to these protocols would see how it handles the complexity introduced by arbitrary faults, which is necessary to model and verify more sophisticated algorithms.
Preliminary exploration by the authors in this direction indicates that \QLogic works very well in these more complex situations --- it can indeed axiomatise algorithms that take account of arbitrary, including hostile, faults.
Unpacking this in full is for future papers.

Recent research has explored DAG-based consensus protocols, which deviate from the traditional linear blockchain structure. 
In these protocols, transactions are organized into a directed acyclic graph (DAG) rather than a strict chain, enabling greater parallelism and scalability. 
Protocols such as DAG-rider~\cite{DBLP:conf/podc/KeidarKNS21}, Narwhal and Tusk~\cite{DBLP:conf/eurosys/DanezisKSS22}, and Bullshark~\cite{DBLP:conf/ccs/SpiegelmanGSK22} represent cutting-edge advancements in this domain. 
Extending \QLogic to analyse DAG-based protocols would be another advance, as these systems often involve more intricate interactions between components and require reasoning about concurrent and asynchronous behaviour at a much higher level of complexity.

Another promising avenue for future research involves leveraging the modularity and expressiveness of \QLogic to study hybrid protocols~\cite{DBLP:conf/sp/NeuTT21} that combine elements of classical consensus with newer techniques, such as Gasper~\cite{DBLP:journals/corr/abs-2003-03052} and others~\cite{DBLP:conf/esorics/DAmatoZ23,DBLP:journals/corr/abs-2411-00558}

It would be helpful in future work to develop automated tools for \QLogic: the tools developed for TLA+~\cite{DBLP:books/aw/Lamport2002} could show the way.
TLA+ is a formal specification language and framework designed to model, specify and verify properties of complex systems; particularly distributed algorithms. 
Its TLC model checker~\cite{DBLP:conf/charme/YuML99} provides exhaustive state exploration to verify properties such as safety and liveness, and the PlusCal translator simplifies specification by offering a high-level pseudocode-like interface.
Additionally, the TLA+ Toolbox~\cite{DBLP:journals/corr/abs-1912-10633} integrates these capabilities into a development environment, supporting workflows that are valuable in both research and industrial applications.

Adapting similar tools to \QLogic could enable automated verification of correctness properties for consensus protocols, including invariants and liveness guarantees. 
For example: a higher-level language akin to PlusCal could simplify the specification of protocols using the axiomatic foundations of \QLogic.
Furthermore, leveraging techniques such as the symbolic model checking enhancements presented by Otoni et al.~\cite{DBLP:conf/tacas/OtoniKKES23} or integrating modular proof systems like TLAPS~\cite{DBLP:conf/fm/CousineauDLMRV12} could further enhance the usability and robustness of the framework.

Such tools would streamline the application of \QLogic to real-world systems, and also make it accessible to a broader audience of researchers and practitioners. 
This direction represents an exciting convergence of theoretical advancements and practical tool development, providing a robust framework for analysing both classical and emerging consensus protocols.

Adopting \QLogic requires some investment in time and effort, but the benefits in clarity, correctness, and reliability may be worthwhile.
Consensus algorithms are known to be particularly challenging, so our relatively high-level declarative approach could help to manage (and debug) the complexity of these algorithms.

We hope this work will inspire readers %
to explore and adopt this new perspective, contributing to more robust and trustworthy distributed protocols in the future.














\documentclass[acmjacm]{acmsmall}

\usepackage[shortlabels]{enumitem}
 
\usepackage{latexsym,xspace,calc,amsthm}
\usepackage{amsmath,amssymb} %
\usepackage{centernot}

\setcounter{tocdepth}{4}
\setcounter{secnumdepth}{3}


\usepackage{lmodern}
\usepackage{tgtermes}
\usepackage{fix-cm}


\usepackage{cancel}

\usepackage{savesym}
\savesymbol{subcaption}
\usepackage{subcaption}

\usepackage{graphbox} %

\makeatletter
\def\@biblabel#1{[#1]} %
\def\thebibliography#1{%
    \footnotesize
    \refsection*{{\refname}
        \@mkboth{\uppercase{\refname}}{\uppercase{\refname}}%
    }
    \list{\@biblabel{\@arabic\c@enumiv}}%
       {\settowidth\labelwidth{\@biblabel{#1}}%
        \leftmargin\labelwidth
        \advance\leftmargin\bibindent
        \itemindent-\bibindent
        \itemsep2pt
        \parsep \z@
        \usecounter{enumiv}%
        \let\p@enumiv\@empty
        \renewcommand\theenumiv{\@arabic\c@enumiv}%
    }%
    \let\newblock\@empty
    \sloppy
    \sfcode`\.=1000\relax
}
\makeatother

\let\odlthebibliography=\thebibliography
 \let\endodlthebibliography=\endthebibliography
 \renewenvironment{thebibliography}[1]{%
   \begin{odlthebibliography}{#1}%
     \setlength{\parskip}{0ex}%
     \setlength{\itemsep}{3pt}%
     \fontsize{9.5}{9.5} %
     \selectfont
}%
 {%
   \end{odlthebibliography}%
 }


\usepackage{tikz}
\usepackage{tikz-cd}
\usetikzlibrary{cd,decorations.markings}
\tikzset{
    dharrow/.style={
        <->,
        postaction={decorate,-},
        }
}
\tikzset{
    dhdashedarrow/.style={
        <->,
        dashed,
        postaction={decorate,-},
        }
    }
\tikzset{
    lrharpoonarrow/.style={
        <[harpoon]->[harpoon],
        postaction={decorate,-},
        }
}
\tikzset{
    lrharpoondashedarrow/.style={
        <[harpoon]->[harpoon],
        dashed, %
        postaction={decorate,-},
        }
}
\usetikzlibrary{arrows}
\usetikzlibrary {arrows.meta} 
\usetikzlibrary{calc}
\usetikzlibrary{positioning}
\usetikzlibrary{snakes,automata,chains}
\usetikzlibrary{graphs}

\usepackage{binarytree}

\usepackage{amssymb,stmaryrd,amsmath}
\usepackage{mdwlist} %
\usepackage{float}   %
\usepackage{centernot} %


\usepackage[colorlinks]{hyperref}
\usepackage{breakurl}  %


\usepackage{bm} %





\newtheoremstyle{jamiestyle}%
  {4pt}%
  {0pt}%
  {\it}%
  {0pt}%
  {\sc}%
  {.}%
  { }%
  {}%
\theoremstyle{jamiestyle}
\newtheorem{thrm}{Theorem}[subsection]
\newtheorem{prop}[thrm]{Proposition}
\newtheorem{lemm}[thrm]{Lemma}
\newtheorem{corr}[thrm]{Corollary}
\newtheorem{conj}[thrm]{Conjecture}
\newtheoremstyle{jamienfstyle}%
  {4pt}%
  {0pt}%
  {\normalfont}%
  {0pt}%
  {\sc}%
  {.}%
  { }%
  {}%
\theoremstyle{jamienfstyle}
\newtheorem{nttn}[thrm]{Notation}
\newtheorem{defn}[thrm]{Definition}
\newtheorem{xmpl}[thrm]{Example}
\newtheorem{rmrk}[thrm]{Remark}


\usepackage{color}
\definecolor{mygreen}{rgb}{0,0.6,0}
\definecolor{mygray}{rgb}{0.5,0.5,0.5}
\definecolor{mymauve}{rgb}{0.58,0,0.82}

\usepackage{listings}
 
\definecolor{gray}{RGB}{128, 128, 128}
\definecolor{lightgray}{RGB}{200, 200, 200}
\definecolor{cyan}{RGB}{0, 255, 255}
\definecolor{blue}{RGB}{0, 0, 255}
\definecolor{red}{RGB}{255, 0, 0}
\definecolor{pink}{RGB}{255, 128, 128}
\definecolor{green}{RGB}{0, 128, 0}
\definecolor{lightyellow}{RGB}{255, 255, 200}
\definecolor{purple}{RGB}{128, 0, 128}

\lstdefinestyle{all}
    {basicstyle=\ttfamily\scriptsize,
     keywordstyle=\color{blue}\ttfamily\scriptsize,
     commentstyle=\color{green}\ttfamily\scriptsize,
     stringstyle=\color{red}\ttfamily\scriptsize}

\lstdefinelanguage{hask}{%
    frame=none,
    xleftmargin=2pt,
    belowcaptionskip=\bigskipamount,
    captionpos=b,
    tabsize=2,
    numbers=left,
    numberstyle=\tiny\color{gray},
    emphstyle={\bf},
	morecomment=[s][\color{green}]{\{-}{-\}},
    stringstyle=\mdseries\rmfamily,
    commentstyle=\color{green},
    keywords={},
    keywords=[1]{case, of, data, if, then, else, where, let, in, do},
    keywords=[2]{Chip, Config, CurrencySymbol, TokenName, PubKeyHash, Integer, Value, State, Action, Text, Maybe, Void, TxConstraints,  Contract},
    keywords=[3]{HasNative},
    keywords=[4]{=>},
    keywords=[5]{Just, Nothing, MkChip, MkConfig, SetPrice, Buy},
    keywordstyle=[1]\mdseries\sffamily\color{red},
    keywordstyle=[2]\mdseries\sffamily\color{blue},
    keywordstyle=[3]\mdseries\sffamily\color{green},
    keywordstyle=[4]\mdseries\sffamily,
    keywordstyle=[5]\mdseries\sffamily\color{purple},
    columns=flexible,
    basicstyle=\small\sffamily,
    showstringspaces=false,
    breaklines=false,
    showspaces=false,
    escapeinside={--}{\^^M},escapebegin={\color{green}--},escapeend={},
    literate= {+}{{$+$}}1 {/}{{$/$}}1 {*}{{$*$}}1 {=}{{$=$}}1
              {>}{{$>$}}1 {<}{{$<$}}1 {\\}{{$\lambda$}}1
              {\\\\}{{\char`\\\char`\\}}1
              {->}{{$\rightarrow$}}2 {>=}{{$\geq$}}2 {<-}{{$\leftarrow$}}2
              {<=}{{$\leq$}}2 {=>}{{$\Rightarrow$}}2
              {\ .}{{$\circ$}}2 {\ .\ }{{$\circ$}}2
              {>>}{{>>}}2 {>>=}{{>>=}}2
              {|}{{$\mid$}}1
              {\_}{{\underline{\hspace{2mm}}}}2
}

\lstdefinelanguage{solidity}{%
    frame=none,
    xleftmargin=2pt,
    belowcaptionskip=\bigskipamount,
    captionpos=b,
    tabsize=2,
    numbers=left,
    numberstyle=\tiny\color{gray},
    emphstyle={\bf},
	morecomment=[s][\color{green}]{\{-}{-\}},
    stringstyle=\mdseries\rmfamily,
    commentstyle=\color{green},
    keywords={},
    keywords=[1]{pragma, solidity, contract, event, constructor, require, function, return, emit},
    keywords=[2]{address, uint, mapping},
    keywords=[3]{public, payable, external, view, returns},
    keywords=[4]{=>, +=, -=, =, <=, ==},
    keywords=[5]{msg, sender, transfer, value},
    keywordstyle=[1]\mdseries\sffamily\color{red},
    keywordstyle=[2]\mdseries\sffamily\color{blue},
    keywordstyle=[3]\mdseries\sffamily\color{green},
    keywordstyle=[4]\mdseries\sffamily,
    keywordstyle=[5]\mdseries\sffamily\color{purple},
    columns=flexible,
    basicstyle=\small\sffamily,
    showstringspaces=false,
    breaklines=false,
    showspaces=false,
    escapeinside={--}{\^^M},escapebegin={\color{green}--},escapeend={},
    literate= {+}{{$+$}}1 {/}{{$/$}}1 {*}{{$*$}}1 {=}{{$=$}}1
              {>}{{$>$}}1 {<}{{$<$}}1 {\\}{{$\lambda$}}1
              {\\\\}{{\char`\\\char`\\}}1
              {->}{{$\rightarrow$}}2 {>=}{{$\geq$}}2 {<-}{{$\leftarrow$}}2
              {<=}{{$\leq$}}2 {=>}{{$\Rightarrow$}}2
              {\ .}{{$\circ$}}2 {\ .\ }{{$\circ$}}2
              {>>}{{>>}}2 {>>=}{{>>=}}2
              {|}{{$\mid$}}1
              {\_}{{\underline{\hspace{2mm}}}}2
}




\newcommand\flanks{\ltimes}
\newcommand\setminusskip{\hskip 2pt plus 1pt minus 2pt}
\newcommand\mysetminus{\setminusskip{{\setminus}}\setminusskip}

\newcommand\declaresoundness[2]{The case of \rulefont{#1} with $#2$.}
\newcommand\declaresoundnessshort[1]{The case of \rulefont{#1}.}


\makeatletter
\newcommand\hpn[2][]{%
  \ext@arrow 9999{\hpnfill@}{#1}{#2}}
\newcommand\hpnfill@{%
  \arrowfill@\leftharpoonup\relbar\rightharpoondown}
\makeatother

\newcommand\arrowlength{0.5}
\newcommand\arrowr{{\tikz \draw [semithick,->] (0,0) -- (\arrowlength*0.8,0);}}
\newcommand\Arrowr{{\tikz \draw [semithick,->,decorate,decoration={zigzag,segment length=3,amplitude=0.5}] (0,0) -- (\arrowlength*0.8,0);}}
\newcommand\Arrowrr{{\tikz \draw [semithick,->>,decorate,decoration={zigzag,segment length=3,amplitude=0.5}] (0,0) -- (\arrowlength*0.8,0);}}
\newcommand\arrowlr{{\tikz \draw [semithick,<->] (0,0) -- (\arrowlength*0.8,0);}}
\newcommand\arrowllrr{{\tikz \draw [semithick,<<->>] (0,0) -- (\arrowlength,0);}}
\newcommand\arrowrr{{\tikz \draw [semithick,->>] (0,0) -- (\arrowlength*0.8,0);}}
\newcommand\arrowrrr{{\tikz \draw [thick,-Stealth] (0,0) -- (\arrowlength*0.8,0);}}
\newcommand\arrowll{{\tikz \draw [semithick,<<-] (0,0) -- (\arrowlength*0.8,0);}}
\newcommand\arrowerr{{\tikz \draw [semithick,-.>>] (0,0) -- (\arrowlength*0.8,0);}}
\newcommand\arrowlrr{{\tikz \draw [semithick,<->>] (0,0) -- (\arrowlength,0);}}
\newcommand\arrowllr{{\tikz \draw [semithick,<<->] (0,0) -- (\arrowlength,0);}}
\newcommand\arrowelrr{{\tikz \draw [semithick,<.->>] (0,0) -- (\arrowlength,0);}}
\newcommand\arrowelr{{\tikz \draw [semithick,<.->] (0,0) -- (\arrowlength,0);}}
\newcommand\arrower{{\tikz \draw [semithick,-.>] (0,0) -- (\arrowlength*0.8,0);}}


\newcommand\everyone{\modality{\Box}}
\newcommand\someone{\modality{\Diamond}}
\newcommand\someoneAll{{\mathring{\exists}_{\hspace{-0pt}\someone}}}
\newcommand\everyoneAll{{\mathring{\forall}_{\hspace{-1pt}\everyone}}}

\newcommand\Quorum{\modality{\mathring{\exists}}}
\newcommand\QuorumBox{\modality{\mathring{\exists}_{\Box}}} %
\newcommand\Coquorum{\modality{\mathring{\forall}}}
\newcommand\CoquorumDiamond{\modality{\mathring{\forall}_{\hspace{-1pt}\Diamond}}} %


\newcommand{\ite}[2]{[#1,\ \text{\it o\hspace{-1pt}'wise}\ #2]}
\newcommand\mruarrow{\mathbin{\rlap{\scalebox{0.73}{$\recentsimple$}}\hspace{-2pt}\rightarrow}}
\newcommand{\mru}[2]{#1\mruarrow#2}
\newcommand{\mrup}[3]{#2[#1\mruarrow #3]}
\newcommand{\mruX}[2]{{#1\mathbin{\recent}#2}}
\newcommand{\established}[2]{#1 \mathbin{\Downarrow_{\QuorumBox}} #2}
\newcommand{\lrharpoon}{\mathbin{\hpn{}}}
\newcommand{\nlrharpoon}{\mathbin{\not\hpn{}}}
\newcommand\id{\f{id}}
\newcommand\st{{:}}

\usepackage{tikz}
\newcommand*{\myhash}{%
  \begin{tikzpicture}
    \pgfmathsetlengthmacro\myWidth{.8*width("=")}%
    \pgfmathsetlengthmacro\myHeight{height("H")}%
    \pgfmathsetlengthmacro\mySepY{.3333*\myWidth}%
    \pgfmathsetlengthmacro\mySideBearing{.1*\myWidth}%
    \def\myAngle{70}%
    \pgfmathsetlengthmacro\mySepX{\mySepY/sin(\myAngle)}%
    \pgfmathsetlengthmacro\mySlantX{\myHeight/tan(\myAngle)}%
    \draw[line cap=round]
      (0, {(\myHeight - \mySepY)/2}) -- ++(\myWidth, 0)
      (0, {(\myHeight + \mySepY)/2}) -- ++(\myWidth, 0)
      ({(\myWidth - \mySepX - \mySlantX)/2}, 0)
      -- ({(\myWidth - \mySepX + \mySlantX)/2}, \myHeight)
      ({(\myWidth + \mySepX - \mySlantX)/2}, 0)
      -- ({(\myWidth + \mySepX + \mySlantX)/2}, \myHeight)
    ;%
    \useasboundingbox
      (-\mySideBearing, 0)
      (\myWidth + \mySideBearing, \myHeight)
    ;%
  \end{tikzpicture}%
}

\newcommand\THREE{{\mathbf 3}}
\newcommand\threeValid{\THREE_{\f{valid}}}
\newcommand\threeCorrect{\THREE_{\f{correct}}}
\newcommand\threeTrue{\THREE_{\f{true}}}
\newcommand\threeFalse{\THREE_{\f{false}}}
\newcommand\binaryconnectives{\ensuremath{\{\tand,\tor,\tnotor,\timpc\}}} %
\newcommand\unaryconnectives{\ensuremath{\{\tneg,\modT,\modB,\modF,\modTB,\modTF\}}}
\newcommand\modalityconnectives{\ensuremath{\{\someone,\everyone,\Quorum,\Coquorum,\sometime,\forever,\tomorrow\}}} %
\newcommand\doublemodalities{\ensuremath{\{\QuorumBox,\CoquorumDiamond,\someoneAll,\everyoneAll\}}}
\newcommand\singlemodalities{\ensuremath{\{\someone,\everyone,\Quorum,\Coquorum\}}}
\newcommand\Pnt{\ns{Point}}
\newcommand\Val{\ns{Value}}
\newcommand\Valuation{\ns{Valuation}}
\newcommand\PredSymb{\ns{PredSymb}}

\newcommand\FOUR{{\mathbf 4}}
\newcommand\tvT{{\mathbf t}}
\newcommand\tvF{{\mathbf f}}
\newcommand\tvB{{\mathbf b}}
\newcommand\tvN{{\mathbf n}}
\newcommand\tvsTB{{\mathit{tb}}}
\newcommand\tvsFB{{\mathit{fb}}}
\newcommand\tvsTT{{\mathit{tt}}}
\newcommand\tvsFF{{\mathit{ff}}}
\newcommand*\Rcirc{\ensuremath{\smash{\overset{\circ}{\smash{\mathsf{R}}\rule{0pt}{1.2ex}}\vphantom{\mathbb{R}}}}}
\newcommand\recentsimple{{\reflectbox{$\mathsf R$}}} 
\newcommand\recent{\recentsimple} %
\newcommand\urecent{{\ubar\recentsimple}} %
\newcommand\mrecent[2]{#1\recentsimple #2}

\newcommand\figunderline[1]{\text{\rlap{\underline{\emph{#1}}}}}




\usepackage{stackengine}
\newcommand\ubar[1]{\stackunder[0.6pt]{$#1$}{\rule{1ex}{.18ex}\hspace{0.5pt}}}


\newcommand\figskip{\\[1ex]}
\newcommand\from[1]{{\uparrow}_{#1}}
\newcommand\fs[1]{\mathtt{#1}}
\newcommand\ps[1]{\mathtt{#1}}
\newcommand\thea{a}
\newcommand\theb{b}
\newcommand\thec{c}
\newcommand\aval{v}
\newcommand\xor{\mathbin{\mathsf{\small xor}}}
\newcommand\modality[1]{#1\hspace{1.5pt}}
\newcommand\modT{\modality{\tf T}}
\newcommand\modTB{\modality{\tf T\hspace{-2.5pt}\tf B}}
\newcommand\modFB{\modality{\tf F\hspace{-1.5pt}\tf B}}
\newcommand\modB{\modality{\tf B}}
\newcommand\modTF{\modality{\tf T\hspace{-2.5pt}\tf F}}
\newcommand\modF{\modality{\tf F}}
\newcommand\witno{\prec}
\newcommand\atopen{T}
\newcommand\afilter{F}
\newcommand\apoint{P}
\newcommand\acontext{\f{ctx}}
\newcommand\avaluation{\varsigma} %
\newcommand\indicator[1]{\delta_{#1}} 
\newcommand\charop[1]{\f{char}(#1)}
\newcommand\resilient{R}
\newcommand\rulefont[1]{\ensuremath{{\mathsf{(#1)}}}}
\newcommand\leftopeninterval[1]{(#1]}
\newcommand\rightopeninterval[1]{[#1)}
\newcommand\openinterval[1]{(#1)}
\newcommand\closedinterval[1]{[#1]}
\newcommand\opens{\ns{Open}}
\newcommand\opensne{\opens_{\hspace{-1pt}{\neq}\varnothing}}
\newcommand\regularOpens{\tf{Open}_{\f{reg}}}
\newcommand\topens{\tf{Topen}}
\newcommand\closed{\tf{Closed}}
\newcommand\regularClosed{\tf{Closed}_{\f{reg}}}
\newcommand\thenext[0]{\f{ad}}
\newcommand\closure[1]{|#1|}
\newcommand\dotclosure[1]{\closure{#1}^{\bigcdot}}
\newcommand{\dotarrow}{%
   \mathrel{\ooalign{\hss\raise.85ex\hbox{\scalebox{1.25}{\normalfont .}}%
   \kern0.35ex\hss\cr$\rightarrow$}}}
\newcommand{\dotleftarrow}{\mathrel{\reflectbox{$\dotarrow$}}}
\newcommand\mycard{\myhash}
\newcommand\narrowmath[1]{\scalebox{0.8}[1.0]{$\mathsf{#1}$}}
\newcommand\sometime{\modality{\narrowmath{sometime}}}
\newcommand\forever{\modality{\narrowmath{forever}}}
\newcommand\tomorrow{\modality{\narrowmath{tomorrow}}}
\newcommand\yesterday{\modality{\narrowmath{yesterday}}}
\newcommand\infinitely{\modality{\narrowmath{infinitely}}}
\newcommand\final{\modality{\narrowmath{finally}}}
\newcommand\arity{\f{arity}}
\newcommand\correct[1]{\tf{correct}[#1]}
\newcommand\correctterm{correct\xspace}
\newcommand\pointwise[1]{\tf{pointwise}[#1]}
\newcommand\pointwiseterm{pointwise\xspace}
\newcommand\incorrect[1]{\tf{incorrect}[#1]}
\newcommand\uniform[1]{\tf{uniform}[#1]}
\newcommand\udfn{\tf{undef}}
\newcommand\gslt{\f{GSLT}}
\newcommand\dense{\f{dense}}
\newcommand\noi{\f{noi}}

\newcommand{\upperset}[2]{
  \underset{
            \text{\raisebox{0.9ex}{\smash{\scalebox{0.8}{$#1$}}\ \ }}%
           }
           {
            \text{\raisebox{0ex}{\smash{$#2$}}}
           }
                       }
\newcommand{\convergesto}[1]{\mathrel{\upperset{#1}{\dotarrow}}}
\newcommand\betwixt{\asymp}
\newcommand\notbetween{\mathbin{\cancel{\between}}}
\newcommand\notintertwinedwith{\mathrel{\notbetween}}
\newcommand\intersectswith{\between}
\newcommand\notintersectswith{\notbetween}
\newcommand\stronglybetween{\mathbin{\rlap{$\between$}\hspace{-1.5pt}\between}}
\newcommand\xxtwined{$2$-twined\xspace}
\newcommand\xxxtwined{$3$-twined\xspace}
\newcommand\intertwined[1]{#1_{\between}}
\newcommand\intertwinedwith{\mathrel{\between}}
\newcommand\stronglyintertwinedwith{\mathrel{\stronglybetween}}
\newcommand\leqk{\leq_{\hspace{-.7pt}\intertwinedwith}}
\newcommand\geqk{\geq_{\hspace{-.7pt}\intertwinedwith}}
\newcommand\cw{\leftrightarrow}  %
\newcommand\ctr{\rightarrow}  %
\newcommand\notctr{{\not\rightarrow}}  %
\newcommand\ctl{\leftarrow}  %
\newcommand\cti{\leq}  %
\newcommand\ctieq{{\stackrel{=}{\hookrightarrow}}}  %
\newcommand\notcti{{\not\hookrightarrow}}  %
\newcommand\nbhd[0]{\f{nbhd}}
\newcommand\interior[0]{\f{interior}}
\newcommand\kiss[0]{\f{kiss}}
\newcommand\community[0]{\f{K}}
\newcommand\framecommunity[0]{\f{k}}
\newcommand\kernel[0]{\f{ker}}
\newcommand\Kernel[0]{\f{Ker}}
\newcommand\witness[0]{\f{W}}
\newcommand\Kmod[1]{\tf K_{\hspace{-2pt}#1}}
\newcommand\cast[1]{#1^{\ast c}}
\newcommand\cclo[1]{#1^c}
\newcommand\mbot{\mathsf{e}}
\newcommand\mtop{\mathsf{f}}
\newcommand\isdef[2]{#1\bullet #2}
\newcommand\deffont[1]{{\bfseries #1}}
\newcommand\powerset{\f{pow}}
\newcommand\comppowerset{\f{pow}_{\hspace{-2pt}c}}
\newcommand\finpow{\f{fin}}
\newcommand\finto{\stackrel{\f{fin}}{\rightharpoonup}}
\newcommand\f[1]{\mathit{#1}}
\newcommand\tf[1]{\mathsf{#1}}
\DeclareMathAlphabet{\mathsb}{OT1}{cmss}{sbc}{n}
\newcommand\ns[1]{{\mathsb{#1}}}
\usepackage{slantsc}
\newcommand\theory[1]{\ensuremath{\text{\itshape\scshape #1}}\xspace}
\newcommand\ThyPax{\theory{TyPax}}
\newcommand\ThySPax{\theory{TySPax}}
\newcommand\ThyBB{\theory{TyBB}}
\newcommand\finsubseteq{\mathbin{\subseteq_{\text{\it fin}}}}
\newcommand\at{\text{@}}
\newcommand\tx{\f{tx}}
\newcommand\ty{\f{ty}}
\newcommand\ctx{\f{ctx}}
\newcommand\txs{\f{txs}}
\newcommand\utxi{\f{utxi}}
\newcommand\utxo{\f{utxo}}
\newcommand\stx{\f{stx}}
\newcommand\utxio{\f{utxio}}
\newcommand\UTxIO{\f{UTxIO}}
\newcommand\inlinehask[1]{\lstinline[language=hask]{#1}}
\newcommand\inlinesolidity[1]{\lstinline[language=solidity]{#1}}
\newcommand\valid{\f{valid}}
\newcommand\liff{\Longleftrightarrow}
\newcommand\limp{\Longrightarrow}
\newcommand\nlimp{\centernot{\Longrightarrow}}
\newcommand\aeq{\mathrel{=_{\alpha}}}
\newcommand\ssm{{{:}\text{=}}}
\DeclareMathSymbol{\shortminus}{\mathbin}{AMSa}{"39}
\newcommand\minus{{\shortminus}}
\newcommand\plus{{+}}
\newcommand\Forall[1]{\forall #1.}
\newcommand\Exists[1]{\exists #1.}
\newcommand\sm{\mapsto}
\newcommand\rewritesto{\longrightarrow}
\newcommand\cent{\vdash}
\newcommand\ncent{\not\vdash}
\newcommand\ment{\vDash}
\newcommand\nment{\not\vDash}
\newcommand\mentval{\ment_{\avaluation}}
\newcommand\nmentval{\nment_{\avaluation}}
\newcommand\supp{\f{supp}}
\newcommand\boundary{\f{boundary}}
\newcommand\lmodel{[\hspace{-0.2em}[}
\newcommand\rmodel{]\hspace{-0.2em}]}
\newcommand\model[1]{{\lmodel #1 \rmodel}}
\newcommand\modellabel[2]{{\lmodel #1 \rmodel}_{#2}}
\newcommand\synmodel[1]{{\langle #1 \rangle}}
\newcommand\mone{{\text{-}1}}
\newcommand\fv{\f{fv}}
\newcommand\Ngeqz{\mathbb N_{\scriptscriptstyle\geq 0}}
\newcommand\Time{\Ngeqz}
\newcommand\weakmodusponens{weak modus ponens (Proposition~\ref{prop.mp.for.tnotor}(\ref{item.mp.for.tnotor}))\xspace}
\newcommand\strongmodusponens{strong modus ponens (Proposition~\ref{prop.mp.for.tnotor}(\ref{item.mp.for.timpc}))\xspace}
\newcommand\weakmodusponensnoref{weak modus ponens\xspace}
\newcommand\strongmodusponensnoref{strong modus ponens\xspace}

\makeatletter
\DeclareRobustCommand{\barcent}{\mathbin{\mathpalette\barcent@@\relax}}
\newcommand{\barcent@@}[2]{%
  \vbox{\offinterlineskip
    \sbox\z@{$\m@th#1\cent$}%
    \ialign{%
      \hfil##\hfil\cr
      $\m@th#1{}_{\minus}\kern-\scriptspace$\cr
      \noalign{\kern-.3\ht\z@}
      \box\z@\cr
    }%
  }%
}
\makeatother

\makeatletter
\def\pmb@#1#2{\setbox8\hbox{$\m@th#1{#2}$}%
  \setboxz@h{$\m@th#1\mkern-.1mu$}\pmbraise@\wdz@
  \binrel@{#2}%
  \dimen@-\wd8 %
  \binrel@@{%
    \mkern-.1mu\copy8 %
    \kern\dimen@\mkern-.2mu\copy8 %
    \kern\dimen@\mkern-.3mu\copy8 %
    \kern\dimen@\mkern-.4mu\copy8 %
    \kern\dimen@\mkern.1mu\copy8 %
    \kern\dimen@\mkern.2mu\copy8 %
    \kern\dimen@\mkern.3mu\copy8 %
    \kern\dimen@\mkern.0mu\raise\pmbraise@\copy8 %
    \kern\dimen@\mkern.4mu\box8 %
           }%
}
\makeatother


\newcommand\overbracebig[2]{\overbrace{#2}^{\mbox{$#1$}}}
\newcommand\compressthislight[1]{{\hspace{.8pt}\raisebox{.5pt}{\scalebox{.85}{$#1$}}\hspace{.2pt}}}
\newcommand\compressthis[1]{\pmb{\compressthislight{#1}}}
\newcommand\tempty{\compressthis{\varnothing}}
\newcommand\tcap{{\compressthis{\cap}}}
\newcommand\tcup{{\compressthis{\cup}}}
\newcommand\tsubseteq{{\compressthis\subseteq}}
\newcommand\tneg{{\pmb\neg}}
\newcommand\tnegnf{{\pmb{\sim}}}
\newcommand\tnegneg{{\pmb{{\sim}{\sim}}}}
\newcommand\ttop{{\pmb\top}}
\newcommand\tbot{{\pmb\bot}}
\newcommand\teq{{\pmb{\text{=}}}}
\newcommand\tneq{{\pmb{\neq}}}
\newcommand\tand{\mathrel{\pmb\wedge}}
\newcommand\tor{\mathbin{\pmb\vee}}
\newcommand\tgg{{\compressthis\gg}}
\newcommand\tleq{\compressthis{\leq}}
\newcommand\tlesssim{\compressthis{\lesssim}}
\newcommand\tapprox{\compressthis{\approx}}
\newcommand\timp{{\compressthis{\Rightarrow}}}
\newcommand\ntimp{{\pmb\not\Rightarrow}}
\newcommand\tnotor{\mathrel{\arrowr}}
\newcommand\tnotore{\mathrel{\arrower}}
\newcommand\tlatticeiff{\mathrel{\arrowlr}}
\newcommand\tlatticeiffec{\mathbin{\arrowelr}}
\newcommand\timpc{\mathrel{\arrowrr}}
\newcommand\timpcc{\mathrel{\arrowrrr}}
\newcommand\timpcleft{\mathrel{\arrowll}}
\newcommand\timpec{\mathrel{\arrowerr}}
\newcommand\tlatticeiffc{{\pmb{\bar\equiv}}}
\newcommand\tiff{\compressthis{\Leftrightarrow}} %
\newcommand\tiffc{\mathbin{\arrowlrr}}
\newcommand\tiffcc{\mathbin{\arrowllrr}}
\newcommand\tiffec{\mathbin{\arrowelrr}}
\newcommand\tlam{{\pmb\lambda}}
\newcommand\tsim{\compressthis{{\approx}}}
\newcommand\tin{{\pmb{\in}}}
\newcommand\ntin{{\pmb{\not\in}}}
\newcommand\qall{{\compressthis{\mathsf Q}}}
\newcommand\tall{{\compressthis{\forall}}}
\newcommand\texi{{\compressthis{\exists}}}
\newcommand\texiunique{{\compressthis{\exists\hspace{0.5pt}!}}}
\newcommand\texiaffine{{\compressthis{\exists_{01}}}}
\newcommand\texiR{{\compressthis{\exists_{\mathsf R}}}}
\newcommand\tlot{\textbf{?}}
\newcommand\AllBut[2]{\tf{AllBut}_{#2}(#1)}
\newcommand\Nset[1]{\{0\dots #1\minus 1\}}

\newcommand\oimplies{\mathbin{\rlap{$+$}\hspace{.7pt}{\supset}}}

\newcommand\QLogic{Coalition Logic\xspace}

\makeatletter
\newcommand{\circlearrow}{}%
\DeclareRobustCommand{\circlearrow}{%
  \mathrel{\vphantom{\shortrightarrow}\mathpalette\circle@arrow\relax}%
}
\newcommand{\circle@arrow}[2]{%
  \m@th
  \ooalign{%
    \hidewidth$#1\circ\mkern1mu$\hidewidth\cr
    $#1\longrightarrow$\cr}%
}
\makeatother
\newcommand\minusoplus{{\pmb{\circlearrow}}} %

\makeatletter
\newcommand*\bigcdot{\mathpalette\bigcdot@{.5}}
\newcommand*\bigcdot@[2]{\mathbin{\vcenter{\hbox{\scalebox{#2}{$\m@th#1\bullet$}}}}}
\makeatother

\newcommand\act{{\cdot}}
\newcommand\pact{{\cdot}}
\newcommand\mact{\mathbin{\bigcdot}}

\usepackage{datetime}
\yyyymmdddate

\begin{document}
\title{A declarative approach to specifying distributed algorithms using three-valued modal logic}
\newcommand\titlerunning{\emph{Declarative distributed algorithms via modal logic}} %
\newcommand\authorrunning{\emph{Murdoch J. Gabbay and Luca Zanolini}}
\author{Murdoch J. Gabbay \affil{Heriot-Watt University, Edinburgh, UK}
Luca Zanolini \affil{Ethereum Foundation, London, UK}}

\begin{abstract}
  We present \emph{\QLogic}, a three-valued modal fixed-point logic designed for declaratively specifying and reasoning about distributed algorithms, such as the Paxos consensus algorithm.

  Our methodology represents a distributed algorithm as a logical theory, enabling correctness properties to be derived directly within the framework—or revealing logical errors in the algorithm's design when they exist.
  
  \QLogic adopts a declarative approach, specifying the overall logic of computation without prescribing control flow. Notably, message-passing is not explicitly modeled, distinguishing our framework from approaches like TLA+. This abstraction emphasises the logical essence of distributed algorithms, offering a novel perspective on their specification and reasoning.
  
  We define the syntax and semantics of \QLogic, explore its theoretical properties, and demonstrate its applicability through a detailed treatment of the Paxos consensus algorithm. By presenting Paxos as a logical theory and deriving its standard correctness properties, we showcase the framework’s capacity to handle non-trivial distributed systems.
  
  We envision \QLogic as a versatile tool for specifying and reasoning about distributed algorithms. The Paxos example highlights the framework’s ability to capture intricate details, offering a new lens through which distributed algorithms can be specified, studied, and checked.


\keywords{Modal logic, distributed consensus} 
\end{abstract}


\maketitle
\thispagestyle{empty}

\tableofcontents

\section{Introduction}
\label{sect.intro}

\subsection{Initial overview}

\subsubsection*{Introducing \QLogic}
\emph{\QLogic}is a logical framework for studying distributed algorithms in a \emph{declarative} manner.
In our framework an algorithm is presented as a logical theory (a signature and a set of axioms), and then its correctness properties are derived from those axioms.
We apply our framework to present and study an axiomatisation of Paxos, a canonical distributed consensus algorithm. 

Logic \emph{per se} is widely-applied in distributed algorithms: for example the logic TLA+ (based on Linear Temporal Logic) is an industry standard to model and verify distributed protocols~\cite{DBLP:journals/toplas/Lamport94,DBLP:books/aw/Lamport2002,DBLP:conf/charme/YuML99}. 
However, TLA+ specifications describe state changes in a sequential, step-by-step manner.
A \QLogic axiomatisation feels more abstract: we dispense with step-by-step descriptions of the algorithm and with explicit message-passing, to capture an abstract but still entirely identifiable essence of the algorithm.

Abstracting an algorithm to a logical theory \emph{per se} is also not new: \emph{declarative programming} is based on this.
However, turning a distributed algorithm (like Paxos) into an axiomatic theory, is new, and this paper shows that it can be done.
Examples follow below and in the body of the text.\footnote{Or, if the reader wants to skip the preliminaries and jump to a toy example that illustrates our method, then see Subsection~\ref{subsect.simple}.}

For reasons of length, we focus on the example of Paxos.
This is a canonical example with sufficient complexity to be interesting, but we have checked that \QLogic can be adapted for other algorithms; it is intended as a general logical tool that can be tailored to particular case studies.
In the long run, we hope this will help do all the things that declarative techniques are supposed to: help find errors, decrease development times, increase productivity, and lead to new mathematics.
 
\subsubsection*{What is declarative programming}
By `declarative' we mean that we specify the logic of a computation without specifying its control flow.

Thus our axiomatisation of Paxos (in Figure~\ref{fig.logical.paxos}) relates to a concrete Paxos implementation and/or its TLA+ specification, much as this axiomatic, algebraic specification of the factorial function\footnote{By `axiomatic' we mean that the function is specified just by asserting properties that are true of it.  By `algebraic' we mean that these assertions are purely equational, i.e. we just assert some equalities between terms.}
$$
0!=1
\qquad 
(n+1)!=(n+1)*n!
$$
relates to this concrete implementation of factorial in BASH
\begin{quote}
{\tt\small\hspace{-2em} value=1; while [ \$number -gt 1 ]; do value=\$((value * number)); number=\$((number - 1)); done}
\end{quote}
The reader may be familiar with programming languages based on declarative principles, such as Haskell, Erlang, or Prolog.
Their design philosophy is that we should assert things that we want to be true, and leave it to the compiler to figure out how best to make this happen.
Contrast with C, which is a canonical imperative language.\footnote{We are not interested here in taking sides.  With a variety of tools to choose from, we pick the one that is best for the job.  What matters is that the widest possible \emph{variety} of levels of abstraction is made available.}

This paper %
explores a methodology for specifying distributed algorithms which sits decidedly at the `Haskell/Erlang/Prolog' side of the spectrum, compared to TLA+ which (while still quite high-level) sits more at the C side of the spectrum. 

Consider a tiny concrete example of a simple voting system: we will briefly exhibit this in \QLogic, just to give a flavour of how it works.
Consider a voting system such that 
\begin{quote}
`if a quorum of voters vote for candidate $c$, then $c$ wins the election'.
\end{quote}
We write this axiomatically simply as
$$
(\QuorumBox\tf{vote}(c)) \timpc \tf{wins}(c) .
$$ 
The $\QuorumBox$ is a logical modality that is true when there exists a quorum of participants at which $\tf{vote}(v)$ holds.\footnote{Details of this modality in our framework are in Proposition~\ref{prop.pointwise.dense.char}(\ref{item.dense.char.pointwise.QuorumBox}) and Lemma~\ref{lemm.allbut.concretely}.}

Even without going into the details of what these symbols mean, it is clear that details of message-passing, of how votes are communicated, and of how they are counted up, are absent.
In fact in this case the predicate is just a direct rendering of the English statement above it, without specifying how it is concretely implemented.
Thus, in keeping with a declarative philosophy, we specify what we want the algorithm to do, but not how it should do it: it is the point of the axiom to express formally and precisely our idea --- and to add as little as possible to that, beyond what is necessary to make it logically precise.


\subsubsection*{Who should read this document}
We have in mind two communities: 
\begin{enumerate*}
\item
\emph{logicians}, who might be interested in exploring a logic with an interesting combination of features (three-valued modal fixedpoint logic) with specific practical motivations in programming distributed systems, and 
\item
\emph{researchers in distributed systems}, who might find \QLogic useful for analysing and communicating their designs.
\end{enumerate*}
We will be particularly interested in consensus algorithms as applied to blockchain, because this is an important and safety-critical use case (the second author is part of the core team developing consensus algorithms for Ethereum).

\subsubsection*{References to key points in the text}

The reader who might want to jump to where the action is and read backwards and forwards from there, is welcome to proceed as follows:
\begin{itemize*}
\item
The axioms of Declarative Paxos are in Figure~\ref{fig.logical.paxos}.
\item
The syntax and semantics of \QLogic are in Figures~\ref{fig.predicate.syntax}, \ref{fig.3.phi.f}, and~\ref{fig.3.derived}.
\item
Truth-tables for three-valued logic and its connectives are in Figure~\ref{fig.3}.
\item
An outline of the Paxos algorithm (for the reader who might like to be reminded what it is) is in Remark~\ref{rmrk.high-level.paxos}.
\item
See also toy axiomatisations for an imaginary `simple' protocol in Subsection~\ref{subsect.simple}, designed to be vaguely reminiscent of Paxos, but simplified almost to triviality for the sake of exposition.
\end{itemize*}
 

\subsection{Motivation and background (further details)}

Let us unpack what consensus in distributed systems (like blockchains) is, why it matters, and why it is hard:

\subsubsection*{Distributed systems}
A distributed computing system is is one that distributes its computation over several machines.
Distributed systems are essential to computing: they enable a network of parties to collaborate to achieve common goals~\cite{cachinbook}. 

Being distributed is attractive to software engineers because it can give us desirable properties, including some or all of: scalability, redundancy, lower latency, reliability, and resilience.
For this reason, almost any nontrivial commercial software system that the reader might care to name is implemented in a distributed manner.
Furthermore, blockchain systems (like Bitcoin and Ethereum\footnote{\url{https://bitcoin.org/en/} and \url{https://ethereum.org/en/}}) and peer-to-peer systems are inherently distributed by design; if they were centralised, then they would not even make sense.

\subsubsection*{The challenge of consensus}
Among the many challenges in distributed systems, \emph{consensus} represents a fundamental abstraction.
Consensus captures the problem of 
\begin{quote}
reaching agreement among multiple participants on a common value, in the face of possibly unreliable communication, and in the presence of possibly faulty participants.
\end{quote}
Consensus algorithms are an essential prerequisite for coordination in any distributed system, and they play a particularly critical role in blockchain networks.
In such networks, consensus is the mechanism that ensures all participants agree on the state of the ledger --- the blockchain's source of truth --- thereby solving problems like \emph{double spending} (manipulating the blockchain ledger to spend the same token twice).
Consensus protocols also enable the secure ordering of transactions, and they ensure the integrity and consistency of the blockchain, even in the face of adversarial actors.

Consensus protocols in blockchain systems are particularly challenging to design and implement correctly. 
Indeed, there are numerous examples in which incomplete or insufficiently-specified protocols have led to serious issues, such as network halts or vulnerabilities that adversaries could exploit.
For example:
\begin{enumerate*}
\item
Solana (\url{https://solana.com}) has experienced instances where its consensus mechanism contributed to temporary network halts~\cite{yakovenko2018solana,shoup2022poh}.
\item 
Ripple's (\url{https://ripple.com}) consensus protocol has been analysed for conditions in which safety and liveness may not be guaranteed~\cite{DBLP:conf/opodis/Amores-SesarCM20}. 
\item
Avalanche's (\url{https://www.avax.network}) consensus protocol has faced scrutiny regarding its assumptions and performance in adversarial scenarios~\cite{DBLP:conf/opodis/Amores-SesarCT22}.
\end{enumerate*}
Even mature blockchains like Ethereum are not immune to challenges.
Attacks such as \emph{balancing attacks}~\cite{DBLP:conf/sp/NeuTT21} and \emph{ex-ante reorganisation} (\emph{reorg}) attacks~\cite{DBLP:conf/fc/Schwarz-Schilling22} have highlighted vulnerabilities in Ethereum's protocol.

These examples underscore how subtle flaws or ambiguities in protocol design can have significant consequences, emphasising the need for rigorous analysis and well-defined specifications in this complex and safety-critical field.

Compounding the fact that consensus algorithms are foundational to the security and functionality of blockchain systems, is the fact that designing and implementing these algorithms is notoriously challenging.
Distributed algorithms --- including consensus protocols --- are difficult (at least for human brains) because of the inherent difficulty of coordinating multiple participants, which may be faulty, hostile, drifting in and out of network connectivity~\cite{sleepy}, and/or have different local clocks, states, agendas, and so forth.

In short, this is a notoriously difficult field which is full of `gotchas', and within this difficult field, consensus protocols are amongst the most intricate, demanding, and practically important objects of study. 

\subsubsection*{Introducing a logic-based declarative approach (applied to Paxos)}
In this paper, we develop a logic-based declarative approach to modeling distributed algorithms.
It stands in relation to existing techniques in roughly the same way as declarative programming languages stand to imperative ones. %

Technically speaking, \QLogic is a \emph{three-valued modal fixedpoint logic}. 
This means:
\begin{itemize*}
\item
There are three truth-values: $\tvT$ (true), $\tvB$ (no value returned, i.e. `crashed'), and $\tvF$ (false)~\cite{sep-logic-manyvalued}.
\item
Modalities let us express things like `some participant has property $\phi$' or `a quorum of participants has property $\psi$' or `$\phi$ has been true in the past' or `$\psi$ will be true in the future'.\footnote{In order, these are as follows in our logic: $\someone\phi$, $\QuorumBox\phi$, $\recent\phi$, and $\sometime\phi$.  Precise semantics are in Figures~\ref{fig.3.phi.f} and~\ref{fig.3.derived}.} 
\item
A fixedpoint operator $\mu$ imports the expressive power of iteration~\cite{wiki:Fixedpoint_logic}.
Consensus algorithms are often iterative, so it is natural to include this in our logic.
\end{itemize*} 
We will apply these constructs to give axiomatic models of the Paxos consensus protocol~\cite{DBLP:journals/tocs/Lamport98,lamport2001paxos,cachin2011paxos} and then we will state and derive standard correctness properties from the axioms.
Our example is elementary but nontrivial and strikes (we think) a good balance between being compact but still exhibiting enough richness of structure to illustrate the essential features of how our logic works. 


The first half of this paper will focus on the logic, and the second half of this paper will apply the logic to the examples.
The reader from a logic background might like to start with the first half and refer to examples for motivation; the reader from a distributed systems background might like to do the opposite and start with the examples in the second half, referring back to the logic as required to understand what is being done.
 
\subsubsection*{Related work using logic}
Logic-based approaches have long been used to model and verify distributed protocols, with TLA+~\cite{DBLP:journals/toplas/Lamport94,DBLP:books/aw/Lamport2002,DBLP:conf/charme/YuML99} being one of the most prominent frameworks in the field. 
Introduced by Leslie Lamport in 1999, TLA+ uses Linear Temporal Logic (LTL)~\cite{DBLP:conf/focs/Pnueli77} to describe state changes in a sequential, step-by-step manner, which has become an industry standard for specifying and verifying distributed algorithms.

But TLA+ does things in a very particular way, which is not declarative.
We would argue that it is good for specifying \emph{algorithms}, and so is not particularly \emph{declarative}.
Specifically, TLA+ excels at specifying and verifying system behaviours through state transitions; it is operational and state-based.
This can sometimes distance TLA+ protocol specifications from the invariants that they are intended to uphold.

This separation arises because TLA+ primarily describes how states evolve \emph{over time}, which can make it challenging to directly reason about the essential properties that protocols are designed to maintain. 
So, although TLA+ provides powerful tools for specifying and checking invariants, the process of verifying that these invariants are preserved across all possible state transitions can be complex and unintuitive, particularly when dealing with the intricacies of distributed algorithms.

\subsubsection*{Our approach}
Our novel framework leverages the mathematics of semitopology and three-valued logic. 

Semitopology is from previous work~\cite{gabbay:semdca,gabbay:semtad}.
It provides a natural way to describe quorums --- sets of participants that work together to progress the system state --- by representing them as open sets.
This perspective is particularly suited to distributed systems, where the behaviour of the system emerges from the interactions of coalitions its components, rather than from a sequence of steps by a central controller.

Three-valued logic (a good overview is in \cite[Chapter~5]{bergmann:intmvf}) helps us to reason about the (possibility for) faults that is inherent in distributed systems.
Unlike classical logic with its binary truth-values $\tvT$ and $\tvF$, three-valued logic introduces a third truth value $\tvB$ that captures the indeterminate states often encountered when dealing with faulty parties.
This turns out to give particularly natural and compact descriptions of distributed protocols, distinguishing between normal operation and faulty behaviour. 
It also provides a formal foundation that would be amenable to verification through tools like SAT solvers and theorem provers (building tooling for what we see in this paper would be future work), enhancing both the expressiveness and verifiability of protocol specifications.

Our approach is particularly relevant in the context of blockchain consensus protocols, which are distributed and decentralised systems that require robust and clear specifications to function correctly.
It is not an exaggeration to say that blockchains in particular have put a rocket underneath the development of new consensus protocols~\cite{DBLP:conf/wdag/CachinV17}.
Many of these are not fully specified or formally verified. 
Often, they are even described informally in blog posts and other non-academic sources --- leading to confusion about their properties and behaviours. 
This lack of formalisation creates a fragmented landscape, where the guarantees and limitations of the protocols remain unclear, hindering their adoption and undermining trustworthiness.

We hope that by providing a logic-based declarative language, and the ability to axiomatise consensus protocols, our framework will offer a way to formally specify, reason about, and verify the properties of the many distributed systems under development --- thus lowering costs, increasing reliability, and reducing development time. 
We would further argue that effective formalisation is essential in reducing ambiguity in the design and evaluation of the many new consensus mechanisms being developed in the blockchain space in particular.

Our approach bridges the gap between the declarative specification of invariants and the practical verification of protocol correctness. 

\subsubsection*{We will explain everything, and that's a feature}

Different readers will find some parts of this paper easy, and other parts less easy --- but which parts these are depends on the reader and their background.\footnote{Not a hypothetical: the first author (with a logic background) took a long time to get his head around how Paxos works, whereas the second author (who knows Paxos well and is familiar with logic in general terms, but is not a logician) was not immediately familiar with modal fixedpoint logic and models of axiomatic logical theories.  The community of readers who approach this paper and are already well-versed in both fields may not be the largest.  We hope this paper will help ameliorate this lamentable state of affairs.}

For example:
\begin{itemize*}
\item
three-valued modal logic and its models, and technical results like the Knaster-Taski fixedpoint theorem used in Proposition~\ref{prop.mu.fixedpoint} may be known to the reader with a background in logic (like the first author), but perhaps not to a distributed systems theorist (like the second author); conversely, 
\item
the Paxos algorithm, and standard technical devices like the discussion motivating $n$-twined semitopologies in Remark~\ref{rmrk.motivate.n-twined}, are the bread-and-butter of a distributed systems expert's work --- but this application might be new to someone from a logic background. 
\end{itemize*}
We play safe and explain everything. 
If the reader finds we are going a bit slow in places, then better that than going too fast and skipping details that might help some other reader follow the argument.
 


\subsection{Structure of the paper}

\begin{enumerate*}
\item
Section~\ref{sect.intro} is the Introduction; You are Here.
The rest of the paper splits into two halves: Sections~\ref{sect.three.truth.values} to Section~\ref{sect.n-twined} have to do with logic and semitopologies; and Sections~\ref{sect.declarative.paxos} to~\ref{sect.simpler.paxos} use \QLogic to declaratively specify and study Paxos.
\ \\
\item
Section~\ref{sect.three.truth.values} introduces the basics of three-valued logic, including notable connectives ($\tnotor$ and $\timpc$) and modalities ($\modT$ and $\modTB$).
We also introduce a notion of validity for truth-values.
\item
Section~\ref{sect.quorum.logic} defines semitopologies (our model of the quorums that are ubiquitous in consensus algorithms), then it builds the syntax and semantics of \QLogic.
\item
Section~\ref{sect.predicates} studies \QLogic predicates; notably developing some logical equivalences, and well-behavedness criteria like being \emph{pointwise}.
\item 
Section~\ref{sect.n-twined} returns to semitopologies, and shows how antiseparation properties can be represented in logical form in \QLogic.
One particular antiseparation property corresponds to the \emph{quorum intersection property} which is used in the correctness proofs for Paxos (see Remark~\ref{rmrk.ntwined.and.the.literature}), but we will put this fact in a more general logical/semitopological context.
\ \\
\item
Section~\ref{sect.declarative.paxos} introduces Declarative Paxos; this is representative for the kind of algorithm that we built \QLogic to express.
We take our time to explain how Paxos works, state the axiomatisation, and discuss it in detail.
\item
Section~\ref{sect.paxos.correctness.properties} states the standard correctness properties for Paxos, phrases them in declarative form, and then proves them from the axioms of Declarative Paxos.
What all our work has accomplished is to convert proving correctness properties into a (relatively straightforward) argument of proving facts of an abstract model from its axioms.
\item
Section~\ref{sect.simpler.paxos} revisits the axioms of Declarative Paxos from Section~\ref{sect.declarative.paxos}, in the light of the proofs in Section~\ref{sect.paxos.correctness.properties} --- and then prunes them down to a minimal theory required to prove correctness.

We call this minimal theory \emph{Simpler Declarative Paxos}: it is simpler, and because further details have been elided, it is more abstract.
The reader may prefer Declarative Paxos to Simpler Declarative Paxos because it is more identifiably related to the Paxos algorithm, but if we want to prove correctness properties, then Simpler Declarative Paxos distils the logical essence of what makes those proofs work. 
\ \\
\item
Section~\ref{sect.conclusions} concludes with conclusions and reflections on future work.
\end{enumerate*} 


\section{Three truth-values}
\label{sect.three.truth.values}

\subsection{Basic definition of $\THREE$, and some truth-tables}

\begin{defn}[Three-valued truth-values]
\label{defn.THREE}
Let $\THREE$ be the poset of \deffont{truth-values} with elements $\tvF$ (false), $\tvB$ (both), and $\tvT$ (true) ordered as 
$$
\tvF<\tvB<\tvT.
$$
\end{defn}

\begin{figure} %
$$
\begin{array}{c}
\begin{array}{c@{\qquad}c@{\qquad}c@{\qquad}c@{\qquad}c@{\qquad}c}
\begin{array}{c c c c}
p\tand q & \tvT & \tvB & \tvF
\\
\tvT &  \tvT & \tvB & \tvF
\\
\tvB &  \tvB & \tvB & \tvF 
\\
\tvF &  \tvF & \tvF & \tvF 
\end{array}
&
\begin{array}{c c c c}
p\tor q & \tvT & \tvB & \tvF
\\
\tvT &  \tvT & \tvT & \tvT
\\
\tvB &  \tvT & \tvB & \tvB
\\
\tvF &  \tvT & \tvB & \tvF
\end{array}
&
\begin{array}{c c c c}
p\tnotor q & \tvT & \tvB & \tvF
\\
\tvT &  \tvT & \tvB & \tvF
\\
\tvB &  \tvT & \tvB & \tvB
\\
\tvF &  \tvT & \tvT & \tvT
\end{array}
&
\begin{array}{c c c c}
p \timpc q & \tvT & \tvB & \tvF
\\
\tvT       & \tvT & \tvF & \tvF
\\
\tvB       & \tvT & \tvB & \tvB
\\
\tvF       & \tvT & \tvT & \tvT
\end{array}
\end{array}
\\[7ex]
\begin{array}{c@{\qquad}c@{\qquad}c@{\qquad}c@{\qquad}c@{\qquad}c@{\qquad}c}
\begin{array}{cc}
\tneg p
\\
\tvT & \tvF
\\
\tvB & \tvB
\\
\tvF & \tvT
\end{array}
&
\begin{array}{cc}
\modT p 
\\
\tvT & \tvT
\\
\tvB & \tvF
\\
\tvF & \tvF
\end{array}
&
\begin{array}{cc}
\modF p 
\\
\tvT & \tvF
\\
\tvB & \tvF
\\
\tvF & \tvT
\end{array}
&
\begin{array}{cc}
\modB p
\\
\tvT & \tvF
\\
\tvB & \tvT
\\
\tvF & \tvF
\end{array}
&
\begin{array}{cc}
\modTB p  
\\
\tvT & \tvT
\\
\tvB & \tvT
\\
\tvF & \tvF
\end{array}
&
\begin{array}{cc}
\modTF p  
\\
\tvT & \tvT
\\
\tvB & \tvF
\\
\tvF & \tvT
\end{array}
&
\begin{array}{cc}
\modFB p  
\\
\tvT & \tvF
\\
\tvB & \tvT
\\
\tvF & \tvT
\end{array}
\end{array}
\\[7ex]
\begin{array}{c@{\qquad}c@{\qquad}c@{\qquad}c@{\qquad}c@{\qquad}c}
\begin{array}{c c c c}
p\timp q & \tvT & \tvB & \tvF
\\
\tvT &  \tvT & \tvB & \tvF
\\
\tvB &  \tvT & \tvB & \tvF 
\\
\tvF &  \tvT & \tvT & \tvT 
\end{array}
&
\begin{array}{c c c c}
p\tleq q & \tvT & \tvB & \tvF
\\
\tvT &  \tvT & \tvF & \tvF
\\
\tvB &  \tvT & \tvB & \tvF 
\\
\tvF &  \tvT & \tvT & \tvT 
\end{array}
\end{array}
\end{array}
$$
\emph{Above, the vertical axis of a table indicates values for $p$; the horizontal axis (if nontrivial) denotes values for $q$.}
\caption{Truth-tables for some operations on $\THREE$ (Remark~\ref{rmrk.three.elementary})}
\label{fig.3}
\end{figure}


\begin{figure} %
$$
\begin{array}{c@{\qquad}c@{\qquad}c@{\qquad}c@{\qquad}c@{\qquad}c}
\begin{array}{c c c c}
p\tlatticeiff q & \tvT & \tvB & \tvF
\\
\tvT &  \tvT & \tvB & \tvF
\\
\tvB &  \tvB & \tvB & \tvB
\\
\tvF &  \tvF & \tvB & \tvT
\end{array}
&
\begin{array}{c c c c}
p \tiffcc q & \tvT & \tvB & \tvF
\\
\tvT       & \tvT & \tvF & \tvF
\\
\tvB       & \tvF & \tvB & \tvB
\\
\tvF       & \tvF & \tvB & \tvT
\end{array}
&
\begin{array}{c c c c}
p \equiv q & \tvT & \tvB & \tvF
\\
\tvT       & \tvT & \tvF & \tvF
\\
\tvB       & \tvF & \tvT & \tvF
\\
\tvF       & \tvF & \tvF & \tvT
\end{array}
\end{array}
$$
\caption{Truth-tables for notions of equivalence (Definition~\ref{defn.equivalent})}
\label{fig.3.iff}
\end{figure}


\begin{rmrk}
\label{rmrk.three.elementary}
We start with some elementary observations about $\THREE$.
\begin{enumerate*}
\item\label{item.three.elementary.bounded.lattice}
$\THREE$ %
has bottom element $\tvF$, top element $\tvT$, and meets $\tand$ and joins $\tor$.
This makes it a \emph{finite lattice}.
\item
$\THREE$ has a negation operation $\tneg$.
\item
$\THREE$ has several natural modalities that identify truth-values, including $\modT$, $\modF$, $\modB$, $\modTB$, and $\modTF$.

For example, the reader can examine the truth-table for $\modTF$ in Figure~\ref{fig.3} and see that it identifies elements of $\{\tvT,\tvF\}$.\footnote{One could also call these modalities \emph{characteristic functions} for various subsets of $\THREE$, by which we mean functions which (literally) return $\tvT$ on elements in the set and $\tvF$ on elements not in the set.}
\item
$\THREE$ also has multiple notions of implication and equivalence, some of which are illustrated in Figures~\ref{fig.3} and~\ref{fig.3.iff}.
We will discuss these implications in Subsection~\ref{subsect.3.imp}.
\end{enumerate*}
\end{rmrk}


\begin{rmrk}
We will use the third truth-value $\tvB$ to represent participants in a distributed algorithm \emph{crashing}.
The method is simple: the truth-value of a predicate describing the outcome of a crashed computation, is deemed to be $\tvB$.

This is similar in spirit to how in domain theory an additional value --- usually written $\bot$ --- is added to a denotation to represent nontermination; the function-value of a nonterminating computation is deemed to be $\bot$~\cite{abramsky:domt}.
The nonterminating computation does not return a value, but we represent this absence of a value in the denotation, using $\bot$. 

In our declarative axiomatisation of Paxos, we use $\tvB$ in much the same way, as a truth-value to describe the \emph{absence} of a return value (either $\tvT$ or $\tvF$) from participants that have crashed.
\end{rmrk}



\subsection{Valid and correct truth-values}



\begin{defn}
\label{defn.tv.ment}
\leavevmode
\begin{enumerate*}
\item\label{item.tv.ment.correct}
Define subsets $\threeValid,\threeCorrect,\threeTrue,\threeFalse\subseteq\THREE$ by
$$
\threeValid = \{\tvT,\tvB\}
\quad\text{and}\quad
\threeCorrect = \{\tvT,\tvF\} 
\quad\text{and}\quad
\threeTrue = \{\tvT\} 
\quad\text{and}\quad
\threeFalse = \{\tvF\} .
$$ 
\item
We refer to $\f{tv}\in\THREE$ as being \deffont{valid} / \deffont{invalid},\ \deffont{correct} / \deffont{incorrect},\ and \deffont{true} / \deffont{untrue},\ and \deffont{false} / \deffont{unfalse}\footnote{\dots `unfalse' is `$\tvT$-or-$\tvB$', not `$\tvT$'.  Similarly, `untrue' is $\tvB$-or-$\tvF$.  We cannot resist stating the following gem of three-valued clarity: not-$\tvT$, is not $\tneg\tvT$, and not-$\tvF$ is not $\tneg\tvF$, and not-$\tvB$ \emph{is} not-$\tneg\tvB$, because $\tvB$ is $\tneg\tvB$.} depending on whether it is / is not in $\threeValid$, $\threeCorrect$, $\threeTrue$, and $\threeFalse$ respectively.
\item\label{item.tv.ment}
Define $\ment\f{tv}$ \ (read `\deffont{$\f{tv}$ is valid}') by
$$
\ment\f{tv} 
\quad\text{when}\quad
\f{tv}\in\threeValid ,
\ \ \text{which just means}\ \ 
\f{tv}\in\{\tvT,\tvB\} .
$$
Thus, `$\f{tv}$ is valid' holds precisely when $\f{tv}$ is a valid truth-value.
\end{enumerate*}
\end{defn}

\begin{rmrk}
\label{rmrk.motivate.validity}
We take a moment to discuss Definition~\ref{defn.tv.ment}:
\begin{enumerate}
\item
We can think of $\tvT$ and $\tvF$ as specific `yes/no' results from the successful completion of some computation.
Therefore, we will call $\tvT$ and $\tvF$ \emph{correct}.
\item
We can think of $\tvB$ as representing failure of a computation to complete correctly --- e.g. because a participant has crashed and so can return neither $\tvT$ nor $\tvF$.
\item\label{item.tb.valid}
In Definition~\ref{defn.tv.ment} we set $\threeValid=\{\tvT,\tvB\}$ and define $\ment\f{tv}$ when $\f{tv}\in\threeValid$, and we call $\tvT$ and $\tvB$ \emph{valid}.
\end{enumerate}
It might seem strange at first to let $\tvB$ be valid; if $\tvB$ represents the (non)behaviour of a crashed participant, then surely crashing should \emph{not} be valid?

Not necessarily.
To see why, consider that many of assertions of interest that we will consider will have the form 
\begin{quote}
`if $p$ is not crashed, then $\ment\phi$ at $p$', 
\end{quote}
or equivalently 
\begin{quote}
`either $p$ is crashed, or $\ment\phi$ at $p$'.
\end{quote}
If we let $\tvB$ be valid (i.e. if we define $\ment\tvB$ to hold), then we can simplify the above to the conveniently succinct assertion 
\begin{quote}
`$\ment\phi$ at $p$'.
\end{quote}
The `if $p$ is not crashed' precondition gets swallowed up into the fact that if $p$ is crashed then $\phi$ returns $\tvB$, and $\ment\tvB$ so there is nothing to prove. 
\end{rmrk}


\begin{rmrk}
Later on in Definition~\ref{defn.validity.judgement}(\ref{item.ctx.ment.phi}) we will extend validity to predicates-in-context, just by evaluating a predicate-in-context to a truth-value using a denotation function, and checking that truth-value is valid in the sense we define here. 
\end{rmrk}


We make some elementary observations about $\THREE$ and Definition~\ref{defn.tv.ment}:
\begin{lemm}
\label{lemm.tand.tor}
Suppose $\f{tv}\in\THREE$ is a truth-value.
Then:
\begin{enumerate*}
\item
$\ment\tvT$ and $\ment\tvB$.
\item
$\nment\tvF$.
\item\label{item.tand.tor.tor}
$\ment\f{tv}\tor\f{tv}'$ if and only if ${\ment\f{tv}}\lor{\ment\f{tv}'}$.
\item
$\ment\f{tv}\tand\f{tv}'$ if and only if ${\ment\f{tv}}\land{\ment\f{tv}'}$.
\end{enumerate*}
\end{lemm}
\begin{proof}
Fact of Definition~\ref{defn.tv.ment} and the truth-tables in Figure~\ref{fig.3}.
\end{proof} 

Negation is conspicuous by its absence in Lemma~\ref{lemm.tand.tor}.
We consider this next:
\begin{lemm}
\label{lemm.para}
Suppose $\f{tv}\in\THREE$ is a truth-value.
Then:
\begin{enumerate*}
\item\label{item.para.em}
$\ment\f{tv}\tor\tneg\f{tv}$, and as a corollary $\nment\f{tv}$ implies $\ment\tneg\f{tv}$.
\item\label{item.para.para}
$\ment\f{tv}\tand\tneg\f{tv}$ is possible, and as corollaries: $\ment\f{tv}$ does not necessarily imply $\nment\tneg\f{tv}$; and it is possible that $\ment\f{tv}$ and $\ment\tneg\f{tv}$ both hold.
\end{enumerate*}
\end{lemm}
\begin{proof}
We consider each part in turn:
\begin{enumerate}
\item
Using Lemma~\ref{lemm.tand.tor}(3) it would suffice to show that $\ment\f{tv}$ or $\ment\tneg\f{tv}$.
So suppose $\nment\f{tv}$.
By Definition~\ref{defn.tv.ment}(\ref{item.tv.ment}) $\f{tv}\not\in\threeValid$, thus $\f{tv}=\tvF$ and by Figure~\ref{fig.3} $\tneg\f{tv}=\tvT\in\threeValid$ and so $\ment\tneg\f{tv}$.
\item
We just take $\f{tv}=\tvB$ and check from Figure~\ref{fig.3} that $\tneg\tvB=\tvB$.
We use Lemma~\ref{lemm.tand.tor}(1\&4).
\qedhere\end{enumerate}
\end{proof}

\begin{rmrk}
Some comments on Lemma~\ref{lemm.para}:
\begin{itemize}
\item
Lemma~\ref{lemm.para}(\ref{item.para.para}) notes a typical fact of validity in three-valued logic, that it is naturally \deffont{paraconsistent}, meaning that it is possible for a thing, and the negation of that thing, to both be valid.
This is a feature of $\tvB\in\THREE$, not a bug, and we will make good use of it later when we use $\tvB$ %
as a return value for `I return no correct truth-value'. 
\item
Lemma~\ref{lemm.para}(\ref{item.para.em}) asserts validity of the \emph{excluded middle} (`$\phi$-or-negation-of-$\phi$'), but its corollaries implicitly illustrate another fact of three-valued logic, by what they do \emph{not} say: 
looking at the proof of Lemma~\ref{lemm.para}(\ref{item.para.em}), it is clearly also true that $\ment\f{tv}\tor\modF\f{tv}$, and that $\ment\f{tv}\tor\modT\tneg\f{tv}$, and that $\nment\f{tv}$ implies $\ment\modF\f{tv}$, and also that $\nment\f{tv}$ implies $\ment\modT\tneg\f{tv}$.  

These are stronger assertions, so why did we not state them in the lemma instead?
Because we will not need the result in those forms later.  

In two-valued logic, if something is true then it is probably a useful lemma.
In contrast, in three-valued logic, there is so much more structure with the third truth-value that more things are \emph{true} than are necessarily \emph{useful}.
So, we have to pick and choose our lemmas, and it is not always the case that the strongest form of a result is the most useful.
\end{itemize}
\end{rmrk}


Lemmas~\ref{lemm.tv.ment.TF} and~\ref{lemm.TF.correct} are easy, but they spell out some useful facts about truth-values and validity judgements:
\begin{lemm}
\label{lemm.tv.ment.TF}
Suppose $\f{tv}\in\THREE$ is a truth-value.
Then:
\begin{enumerate*}
\item\label{item.tv.ment.TF.1}
$\threeCorrect\cap\threeValid=\threeTrue$.

In words: a truth-value is true if and only if it is correct and valid.
\item\label{item.tv.ment.TF.2b}
$\ment \modT\f{tv}$ if and only if $\f{tv}\in\threeTrue$. 

In words: `$\modT\f{tv}$ is valid' holds if and only if $\f{tv}$ is true.
\item\label{item.tv.ment.TF.2b.F}
$\ment \modF\f{tv}$ if and only if $\f{tv}\in\threeFalse$. 

In words: `$\modT\f{tv}$ is valid' holds if and only if $\f{tv}$ is true.
\item\label{item.tv.ment.TF.TB}
$\ment \modTB\f{tv}$ if and only if $\f{tv}\in\threeValid$ if and only if $\ment\f{tv}$. 

In words: `$\modTB\f{tv}$ is valid' if and only if $\f{tv}$ is valid.\footnote{Note that $\modTB\f{tv}$ is not necessarily \emph{literally equal} to $\f{tv}$, in the sense of $\modT\f{tv}\equiv\f{tv}$ from Definition~\ref{defn.equivalent}.  For example, $\modTB\tvB=\tvT\neq\tvB$.}
\item\label{item.tv.ment.TF.2}
$\ment \modTF\f{tv}$ if and only if $\f{tv}\in\threeCorrect$. 

In words: `$\modTF\f{tv}$ is valid' if and only if $\f{tv}$ is correct.
\end{enumerate*}
\end{lemm}
\begin{proof}
Part~\ref{item.tv.ment.TF.1} is a fact of sets.
Parts~\ref{item.tv.ment.TF.2b} and~\ref{item.tv.ment.TF.TB} and~\ref{item.tv.ment.TF.2} are just from the truth-table for $\modT$, $\modTB$, and $\modTF$ in Figure~\ref{fig.3} and Definition~\ref{defn.tv.ment}. 
\end{proof}


\begin{lemm}
\label{lemm.TF.correct}
Suppose $\f{tv}\in\THREE$ is a truth-value.
Then the following are all equivalent:
$$
\ment\modTF\f{tv}
\ \liff\ 
\modTB\f{tv}=\modT\f{tv}
\ \liff\ 
\modT\f{tv}=\f{tv} 
\ \liff\ 
\modFB\f{tv}=\modF\f{tv}
\ \liff\ 
\modF\f{tv}=\tneg\f{tv} 
$$ 
Informally --- we do not build the machinery to make this formal, but it may still be a helpful observation --- we can write that the modality $\modTF$ is equivalent to the modal equivalences $\modTB=\modT=\f{id}$ and $\modFB=\modF=\tneg$.
\end{lemm}
\begin{proof}
By Lemma~\ref{lemm.tv.ment.TF}(\ref{item.tv.ment.TF.2}) and Definition~\ref{defn.tv.ment}(\ref{item.tv.ment.correct}),
$\ment\modTF\f{tv}$ when 
$$
\f{tv}=\tvT \ \lor\ \f{tv}=\tvF
\quad\text{and when}\quad
\f{tv}\neq\tvB .
$$
We check the truth-tables in Figure~\ref{fig.3}, and see that all of the other statements above precisely characterise the same condition.
\end{proof}


\subsection{The rich implication and equivalence structure of $\THREE$}
\label{subsect.3.imp}

\subsubsection{Weak and strong implications, and others}


Three-valued logic supports sixteen different kinds of implication~\cite[note~5, page~22]{arieli:idepl}.
We discuss two of them in Definition~\ref{defn.three.imp}, and mention two more in Remark~\ref{rmrk.rich.implication}:
\begin{defn}
\label{defn.three.imp}
\leavevmode
\begin{enumerate*}
\item\label{item.three.elementary.weak.implication}
Define \deffont{weak implication} $\tnotor$, using $\tneg$ and $\tor$, as follows:
$$
\f{tv}\tnotor\f{tv}'=(\tneg\f{tv})\tor\f{tv}'.
$$
\item\label{item.three.elementary.strong.implication}
Define \deffont{strong implication} $\timpc$, using $\tnotor$ and $\modT$, as follows:
$$
\f{tv}\timpc\f{tv}'=\f{tv}\tnotor \modT\f{tv}'. 
$$
\end{enumerate*}
These and other encodings are assembled in Figure~\ref{fig.3.derived}.
\end{defn}

We motivate the weak and strong implications from Definition~\ref{defn.three.imp}(\ref{item.three.elementary.weak.implication}\&\ref{item.three.elementary.strong.implication}):
\begin{rmrk}[How to interpret $\tnotor$ and $\timpc$]
\label{rmrk.interpret.implication}
\leavevmode
For this Remark, let us interpret $\f{tv}\in\THREE$ as the result of a computation process which either succeeds and returns true (truth-value $\tvT$), or it succeeds and returns false (truth-value $\tvF$), or the process does not succeed (e.g. it crashes) and it returns an invalid value / no value (represented by the truth-value $\tvB$).
Recall also from Definition~\ref{defn.tv.ment}(\ref{item.tv.ment}) that $\ment\f{tv}$ means $\f{tv}\in\{\tvT,\tvB\}$.
Then by this reading:
\begin{enumerate*}
\item
$\ment\f{tv}$ means \emph{``the process $\f{tv}$ returns $\tvT$, if valid''}. 
\item
$\ment\f{tv}\tnotor\f{tv}'$ means \emph{``\emph{if} process~$\f{tv}$ returns true \emph{then} process~$\f{tv}'$ returns $\tvT$, if valid''}.

Implication does not imply sequentiality: we are just making a logical assertion about outcomes. 
\item
$\ment\f{tv}\timpc\f{tv}'$ means \emph{``\emph{if} process~$\f{tv}$ returns true \emph{then} process~$\f{tv}'$ is \emph{also} valid, and it returns $\tvT$''}.
\end{enumerate*}
\end{rmrk}


We can sum this up as a \emph{modus ponens} result, which relates the validity $\ment$ from Definition~\ref{defn.tv.ment}, the weak implication $\tnotor$ and strong implication $\timpc$ from Figure~\ref{fig.3}, to real-world implication:
\begin{prop}[Weak and strong modus ponens]
\label{prop.mp.for.tnotor}
Suppose $\f{tv},\f{tv}'\in\THREE$ are truth-values.
Then:
\begin{enumerate*}
\item\label{item.mp.for.tnotor}
$\ment\f{tv}\tnotor\f{tv}'$ if and only if\ \ $\ment\modT\f{tv}$ implies $\ment\f{tv}'$. 

In words: `$\f{tv}$ weakly implies $\f{tv}'$' is valid, if and only if $\f{tv}$ is true implies $\f{tv}'$ is valid (but $\f{tv}'$ need not be true). 

We may call this logical equivalence %
\deffont{weak modus ponens}.
\item\label{item.mp.for.timpc}
$\ment\f{tv}\timpc \f{tv}'$ if and only if $\ment\modT\f{tv}$ implies $\ment\modT\f{tv}'$. 

In words: `$\f{tv}$ strongly implies $\f{tv}'$' is valid if and only if $\f{tv}$ is true implies $\f{tv}'$ is true. 

We may call this logical equivalence \deffont{strong modus ponens}. 
\end{enumerate*}
\end{prop}
\begin{proof}
We consider each part in turn:
\begin{enumerate}
\item
By direct inspection of the truth-table for $\tnotor$ in Figure~\ref{fig.3}, using Definition~\ref{defn.tv.ment}(\ref{item.tv.ment}).
\item
By a similar direct inspection.
\qedhere\end{enumerate}
\end{proof}


\begin{rmrk}[Explicit and implicit implications]
\label{rmrk.rich.implication}
So far, the structure we have seen (the lattice structure, the negation, the modalities, and the two implications $\tnotor$ and $\timpc$) is reflected explicitly in the syntax of our logic in Figures~\ref{fig.predicate.syntax} and~\ref{fig.3.derived}.
For instance, Figure~\ref{fig.predicate.syntax} has a unary negation connective $\tneg$ whose denotation as per Figure~\ref{fig.3.phi.f} is precisely the negation operator $\tneg$ in Figure~\ref{fig.3}.
So far, so explicit.

However, $\THREE$ has further implication structure, which we do not mention explicitly in our logic, but which is still there, is important, and which appears implicitly in the proofs.
Notably: 
\begin{enumerate*}
\item\label{item.rich.implication.timp}
An implication $\timp$ defined by $\f{tv}\timp \f{tv}'=\tneg\f{tv}'\timpc\tneg\f{tv}=(\modTB\f{tv})\tnotor\f{tv}'$ has a truth-table spelled out in Figure~\ref{fig.3} (bottom left).
This can be read as ``if $\f{tv}$ is valid then $\f{tv}'$ is valid''.

$\timp$ is clearly implicit in results like Proposition~\ref{prop.someone.implies.someoneAll} -- see all the parts whose statement uses logical implication $\limp$ --- and in Proposition~\ref{prop.tv.ment.TF.model}.
\item
Another implication is implicit in the use of $\leq$ in results like Lemma~\ref{lemm.positive.monotonicity} and Corollary~\ref{corr.recent.everyone.commute}.
That symbol refers to the lattice order on $\THREE$, but this could easily be internalised in our logic --- e.g. we can set $\f{tv}\tleq \f{tv'}$ to be $(\f{tv}\timpc\f{tv}')\tand(\f{tv}\timp\f{tv}')$, to obtain the connective $\tleq$ also illustrated in Figure~\ref{fig.3} (bottom right).
\end{enumerate*}
So in fact, we use all four of $\timp$ and $\tleq$, and $\tnotor$ and $\timpc$.

The reason that $\tnotor$ and $\timpc$ make it into our logic as (sugared) syntax, and $\timp$ and $\leq$ do not, is that $\tnotor$ and $\timpc$ naturally arise in our axioms (see most of the axioms in Figures~\ref{fig.logical.paxos}) %
whereas we do not need $\timp$ or $\tleq$ in the same way for the axiomatisations in this paper.
Different axiomatisations of different algorithms might make different demands on our logic, and it would be straightforward to include (sugar for) connectives accordingly. 
\end{rmrk}

\subsubsection{Equivalence $\equiv$}

Our three-valued domain of truth-values $\THREE$ has rich notions of equivalence, to match its rich notions of implication.
We mention three of them:
\begin{defn}[Notions of equivalence in $\THREE$]
\label{defn.equivalent}
Define $\phi\tlatticeiff\phi'$ and $\phi\tiffcc\phi'$ and $\phi\equiv\phi'$ by:
$$
\begin{array}{r@{\ }l}
\phi\tlatticeiff\phi' =& (\phi\tnotor\phi')\tand(\phi'\tnotor\phi)
\\
\phi\tiffcc\phi' =& (\phi\timpc\phi')\tand(\phi'\timpc\phi)
\\
\phi\equiv\phi' =& (\modT\phi\tand\modT\phi')\tor(\modB\phi\tand\modB\phi')\tor(\modF\phi\tand\modF\phi')
\end{array}
$$
Truth-tables for $\tlatticeiff$, $\tiffcc$, and $\equiv$ are given in Figure~\ref{fig.3.iff}, and we unpack this further in Remark~\ref{rmrk.equiv}.
\end{defn}


\begin{rmrk}
\label{rmrk.equiv}
We can read 
\begin{itemize*}
\item
$\f{tv}\tlatticeiff \f{tv}'$ as `$\f{tv}$ is valid if and only if $\f{tv}'$ is valid', 
and 
\item
$\f{tv}\tiffcc \f{tv}'$ as `$\f{tv}$ is true if and only if $\f{tv}'$ is true', 
and 
\item
$\f{tv}\equiv \f{tv}'$ as `$\f{tv}$ is literally equal to $\f{tv}'$'.
\end{itemize*}
Of these three notions of equivalence, $\equiv$ will be most useful to us later.
We will consider it in our predicate language in Subsection~\ref{subsect.equivalence.of.predicates}, and (for instance) it will help us to express notions of extensional equivalence in Figure~\ref{fig.easy.equivalences}.

We will call $\equiv$ \deffont{literal equivalence}, and $\equiv$ appears in the sugared syntax of Figure~\ref{fig.3.derived}.
As per the discussion in Remark~\ref{rmrk.rich.implication}, there would be no barrier to including $\tlatticeiff$ or $\tiffcc$ but we do not need to for now.
\end{rmrk}








\section{\QLogic}
\label{sect.quorum.logic}

\subsection{A preliminary: semitopologies}

\begin{rmrk}
A semitopology is like a topology, but generalised to not require that intersections of open sets necessarily be opens.

The intuition of an open set is an \emph{actionable coalition}; a set of participants with the voting power (or other resources) to make updates to their local state.
The reader might be familiar with the notion of \emph{quorum} from literature like~\cite{DBLP:journals/tocs/Lamport98,DBLP:journals/dc/MalkhiR98,naor:loacaq}, or the notion of \emph{winning coalition} in social choice theory (the theory of voting)~\cite[Item~5, page~40]{riker:thepc}.
Semitopologies are similar to this, except that open sets are always closed under arbitrary (possibly empty) unions, whereas 
\begin{itemize*}
\item
quorums may be closed under arbitrary nonempty unions (in which case the correspondence is precise, except that $\varnothing$ is an open set but is not a quorum), or 
\item
quorums are not closed under arbitrary unions (in which case a quorum corresponds to a minimal nonempty open set).
\end{itemize*}
Our presentation of semitopologies will be quite brief --- we just define the parts we need to give semantics to the modalities in \QLogic --- and the reader who wants to know more is referred to previous work~\cite{gabbay:semdca}. 
\end{rmrk} 

\begin{defn}
\label{defn.semitopology}
A \deffont{semitopology} is a pair $(\Pnt,\opens)$ where:
\begin{enumerate*}
\item
$\Pnt$ is a set of \deffont{points}, which following terminology from distributed systems one might also call \deffont{participants}, \deffont{processes}, \deffont{parties}, or \deffont{nodes}.
\item
$\opens\subseteq\powerset(\Pnt)$ is a set of \deffont{open sets} such that:
\begin{enumerate*}
\item
If $\mathcal O\subseteq\powerset(\Pnt)$ then $\bigcup\mathcal O\in\Pnt$.
In words: $\opens$ is closed under arbitrary (possibly empty) unions.
\item
$\Pnt\in\opens$.
\end{enumerate*} 
\end{enumerate*} 
We may write $\opensne$ for the set of nonempty open sets, so
$$
\opensne = \{O\in\opens \mid O\neq\varnothing\} .
$$
\end{defn}


\begin{nttn}
\label{nttn.empty.semitopology}
Suppose $(\Pnt,\opens)$ is a semitopology.
\begin{enumerate*}
\item
We call $(\Pnt,\opens)$ the \deffont{empty semitopology} when its set of points is empty --- i.e. when $\Pnt=\varnothing$ --- so that necessarily $\opens=\{\varnothing\}$.
\item
We call $(\Pnt,\opens)$ a \deffont{nonempty semitopology} when its set of points is nonempty --- i.e. when $\Pnt\neq\varnothing$.
\end{enumerate*}
\end{nttn}


\begin{defn}[All-but-$f$-many semitopology]
\label{defn.allbut}
\leavevmode
\begin{enumerate*}
\item
Suppose $X$ is a finite set and $N\in\Ngeqz$.
Define $\powerset_{\geq N}(X)$ to be the set of subsets of $X$ that contain at least $N$ elements.
In symbols:
$$
\powerset_{\geq N}(X) = \{O\subseteq X \mid \mycard O\geq N\} .
$$
\item\label{item.allbut}
Suppose $N,f\in\Ngeqz$ and $N\geq f$.
Define a semitopology $\AllBut{N}{f}$ by
$$
\AllBut{N}{f} = \bigl(\Nset{N}, \{\varnothing\}\cup\powerset_{\geq N\minus f}(\Nset{N}\bigr)\bigr).
$$
Thus: points are numbers $0\leq n<N$; and open sets are sets of points of cardinality at least $N\minus f$.
\end{enumerate*}
\end{defn}

\begin{lemm}
Suppose $N,f\in\Ngeqz$ and $N\geq f$.
Then:
\begin{enumerate*}
\item
$\AllBut{N}{f}$ is indeed a semitopology.
\item
$\AllBut{N}{f}$ is a topology precisely when $f=N$ or $f=0$.
\end{enumerate*}
\end{lemm}
\begin{proof}
We consider each part in turn:
\begin{enumerate}
\item
It is routine to check that the conditions on open sets in Definition~\ref{defn.semitopology} are satisfied.
The key observation is just the elementary fact that if $\mycard O\geq N\minus f$ and $\mycard O'\geq N\minus f$ then $\mycard(O\cup O')\geq N\minus f$.
\item
For $\AllBut{N}{f}$ to be a topology, open sets need to be closed under intersections.
It is routine to check that this can happen only when $\AllBut{N}{f}$ is the full powerset (so $f=N$) or when open sets are just $\varnothing$ and $\Pnt$ (so $f=0$).
\qedhere\end{enumerate}
\end{proof}

\begin{rmrk}
The original Paxos paper by Lamport~\cite{DBLP:journals/tocs/Lamport98} used semitopologies of the form $\AllBut{N}{f}$ where $N=2n$ and $f=n\minus 1$, for $n\geq 1$ (i.e. open sets are strict majorities of points).
This invites a question: why bother with semitopologies at all?
\begin{enumerate*}
\item
Because we can. 
More specifically, our axiomatisation is parametric over the underlying semitopology.
As good mathematicians, if we can generalise at no cost, then we should.\footnote{\dots and this applies not just in mathematics papers.  Software engineers follow a similar principle and have many names for it: e.g. encapsulation, APIs, modules, classes, typeclasses, and so on.} 
\item
Semitopology not only generalises; it simplifies.
The only properties of $\AllBut{N}{f}$ that we will need, are precisely that it is a semitopology.
If all we need about $\AllBut{N}{f}$ is its semitopological structure, then we should work with that structure because that is what actually matters.
\item
Viewing $\AllBut{N}{f}$ as a semitopology, rather than just as a collection of quorums, adds real meaning and understanding.
In particular, the notions of `dense subset' and `nonempty open interior' will be relevant to understanding the axioms, and these arise naturally in the predicates for declarative axiomatisations of distributed algorithms.
These are clearly native topological notions, so the mathematics is speaking to us here, and what it has to say is inherently topological. 
\end{enumerate*} 
\end{rmrk}


\subsection{Syntax of the logic}

We now develop our three-valued modal fixedpoint logic for reasoning about distributed systems.
We so this in the following order: \emph{signatures} $\Sigma$, \emph{models} $\mathcal M$, \emph{contexts} $\acontext$, and finally \emph{syntax}. 
We develop denotation for this syntax in Subsection~\ref{subsect.denotation}.
For quick reference, the reader can just look at the core syntax and denotation in Figures~\ref{fig.predicate.syntax} and~\ref{fig.3.phi.f}.

 


\begin{defn}[Signatures]
\label{defn.signature}
\leavevmode
\begin{enumerate*}
\item\label{item.signature.varsymb}
Fix $\tf{VarSymb}$ a countably infinite set of \deffont{variable symbols} $\thea,\theb,\thec\in\tf{VarSymb}$.
This set will never change and we may assume it henceforth.
\item
Fix a \deffont{fixedpoint variable symbol} $\tf X$ (we will only need one of these, though there is inherent obstacle to having more).
\item\label{item.predicate.syntax.signature}
A \deffont{signature} $\Sigma$ is a pair 
$$
\Sigma=(\PredSymb,\arity)
$$ 
where:
\begin{enumerate*}
\item\label{item.predsymb}
$\PredSymb$ is a (possibly finite) set of \deffont{predicate symbols} $\tf P,\tf P'\in\PredSymb$.
\item
$\arity$ is an \deffont{arity function} $\arity(\tf P)\in\Ngeqz$ mapping each predicate symbol in $\PredSymb$ to a nonnegative natural number (which is intuitively its number of arguments).
\end{enumerate*}
\item
We may write a signature $\Sigma$ in list form as 
$$
\Sigma=[\tf P,\tf P':i,\ \tf Q:j,\ \dots]
$$ 
meaning that $\tf P,\tf P',\tf Q\in\PredSymb$ and $\arity(\tf P)=\arity(\tf P')=i$ and $\arity(\tf Q)=j$ and so on.
\item\label{item.predicate.syntax.empty.signature}
If the signature in part~\ref{item.predsymb} above contains no predicate symbols then we call it \deffont{empty}; otherwise we call it non-empty.
There is only one empty signature, which we will call \deffont{the empty signature}.\footnote{It is not particularly useful.} 
\end{enumerate*}
\end{defn}


\begin{defn}[Models]
\label{defn.model}
A \deffont{model} $\mathcal M$ is a tuple $\mathcal M=(\Pnt,\opens,\Val)$ where:
\begin{enumerate*}
\item
$(\Pnt,\opens)$ is a semitopology (Definition~\ref{defn.semitopology}), which we may call the \deffont{underlying semitopology} of the model.

We may call $p\in\Pnt$ a \deffont{point} or a \deffont{participant}; also \deffont{party}, \deffont{process}, or \deffont{node} would be consistent with the literature.
For our purposes, these are all synonyms: the terminology \emph{point} is from topology; \emph{participant}, \emph{party}, and \emph{process} are from distributed systems; and \emph{node} is from graph theory (thinking of a network as a graph of communicating processes).
\item\label{item.model.val}
$\Val$ is a nonempty set of \deffont{values}.\footnote{We assume values are nonempty to avoid the degenerate case where a universal quantification over values holds just because there are no values on which to test it.  See in particular the argument in Lemma~\ref{lemm.intertwined.char}(\ref{item.intertwined.char.reverse}).}

For our purposes we will not care what values are, so long as some exist, so we do not put any type system on values or endow them with structure; but we could easily do so if required. 
\end{enumerate*} 
\end{defn}


\begin{defn}[Contexts]
\label{defn.contexts}
Suppose $\Sigma=(\PredSymb,\arity)$ is a signature, and $\mathcal M=(\Pnt,\opens,\Val)$ is a model. 
\begin{enumerate*}
\item\label{item.context.stage}
A \deffont{stage} is a nonnegative number $n\in\Time$.
\item\label{item.contexts.valuation}
A \deffont{valuation} $\avaluation$ inputs a predicate symbol $\tf P\in\PredSymb$ and a stage $n\in\Time$ and a point $p\in\Pnt$ and returns a function $\avaluation(\tf P)_{n,p}:\Val^{\arity(\tf P)}\to\THREE$.

We write the set of valuations $\Valuation(\Sigma,\Val)$ or $\Valuation(\Sigma,\mathcal M)$.
\item\label{item.contexts}
A \deffont{context} $\acontext$ is a tuple $(n,p,O,\avaluation)$ where $n\in\Time$ is a stage, $p\in\Pnt$ is a point, $O\in\opensne$ is a nonempty open set, and $\avaluation\in\tf{valuation}(\Sigma,\Val)$ is a valuation. 

We write the set of contexts $\tf{Ctx}(\Sigma,\mathcal M)$; if $\Sigma$ and $\mathcal M$ are understood, we may write this just $\tf{Ctx}$.
\item
A \deffont{context valuation} is a function $\varkappa\in \THREE^{\tf{Ctx}}$.
\end{enumerate*}
\end{defn}



\begin{figure} %
$$
\begin{array}{r@{\ }l}
t ::=& 
(a\in\tf{VarSymb}) \mid (v\in\Val) %
\\
\phi ::=&  
\tvB 
\mid t\teq t \mid (\tf P(\overbrace{t,\dots,t}^{\arity(\tf P)}) : \tf P\in\PredSymb)
\\
\mid &\tneg\phi \mid \modT\phi \mid \someone\phi \mid \Quorum\phi \mid \texi a.\phi
\\
\mid & \phi\tor\phi'  %
\\
\mid & \yesterday \phi \mid \tomorrow \phi 
\\
\mid & \tf X \mid \mu\tf X.\phi \mid (\varkappa \in \THREE^{\tf{Ctx}})
\end{array}
$$ 
\caption{Term and predicate syntax of \QLogic (Definition~\ref{defn.predicate.syntax})}
\label{fig.predicate.syntax}
\end{figure}


\begin{defn}[Syntax]
\label{defn.predicate.syntax}
Suppose $\Sigma$ is a signature (Definition~\ref{defn.signature}) and $\mathcal M$ is a model (Definition~\ref{defn.model}).
\begin{enumerate*}
\item 
Define the syntax of \deffont{\QLogic} over $\Sigma$ and $\mathcal M$ as in Figure~\ref{fig.predicate.syntax}.

In this syntax, $\texi\thea.\phi$ binds $\thea$ in $\phi$, and $\mu\tf X.\phi$ binds $\tf X$ in $\phi$.
We call a term or predicate with no free variable symbols or fixedpoint variable symbols \deffont{closed}.
\item
One special case will be of particular interest: we will write 
\begin{quote}
`suppose $\texi a.\phi$ is a closed predicate',
\end{quote}
and often we do not really care about $\texi a.\phi$ \emph{per se}, except that when this is a closed predicate it means that $a\in\tf{VarSymb}$ and $\phi$ is a predicate whose set of free variables is at most $\{a\}$.
See for example Remark~\ref{rmrk.unpack.texiunique} and Lemma~\ref{lemm.unique.affine.existence}.
\item
We define substitution of variable symbols / fixedpoint variable symbols for values / context valuations as usual, and write these $\phi[a\ssm v]$ and $\phi[\tf X\ssm \varkappa]$ respectively. %
\end{enumerate*}
\end{defn}


\begin{rmrk}
We call our logic \emph{\QLogic} to reflect that it is a logic for reasoning about how coalitions of participants can work together to carry out algorithms to achieve their goals.
In this document we will apply it to Paxos, but we see this logic as having independent interest and a broader relevance for specifying and reasoning about the kinds of algorithms that get written for distributed systems.
Their hallmark is a requirement for distributed collaborative action --- action by quorums, or actionable coalitions, or winning coalitions; these things have many names --- to progress an algorithm to its next step. 
\end{rmrk}

\begin{xmpl}
We give some example predicates and suggest intuitive meanings for them.
We do not yet have the denotation (Figure~\ref{fig.3.phi.f}) nor the rich derived syntax (Figure~\ref{fig.3.derived}) that we build later, so these examples are simple and informal, but we hope they will be helpful.

Assume our signature has a unary predicate symbol $\tf{output}$ (representing the output of some process), and assume $v,v'\in\Val$ are values.
\begin{enumerate}
\item
$\tf{output}(v)$ means `where we are, today, if the process succeeds, then it outputs $v$' and $\tneg\tf{output}(v')$ means `where we are, if the process succeeds, then it does not output $v'$'.
\item
$\modT\tf{output}(v)$ means `where we are, today, the process succeeds and outputs $v$'. 
\item
$\someone\tf{output}(v)$ means `there is a place where, today, if the process succeeds, then it outputs $v$'.
\item
$\tomorrow\modT\tf{output}(v)$ means `where we are, tomorrow, the process succeeds and outputs $v$'.
\item
$\yesterday\modT\tf{output}(v)$ means `where we are, yesterday, the process succeeded and output $v$'.
This assertion is subtly stronger than it might seem; our notion of time starts at day~0, at which by convention all assertions about yesterday are false.
See Remark~\ref{rmrk.further.comments.denotation}(\ref{item.interpret.yesterday}).
\item
$\tvB$ means `where we are, today, the process (which we have not named with a predicate symbol) crashes'.
\end{enumerate}
More sophisticated examples will appear later; notably in the derived syntax in Figure~\ref{fig.3.derived}, in Remark~\ref{rmrk.comments.on.semantics}, and of course when we apply our logic to axiomatise distributed algorithms.
\end{xmpl}


\begin{rmrk}[Three pedant points]
\label{rmrk.pedant.point}
The thing about being pedantic is that it does not matter, until it does.
We now discuss three pedant points, which may not matter \dots until perhaps a moment arrives then they do:
\begin{enumerate}
\item
We may treat $\Sigma$ and $\mathcal M$ as fixed parameters, and elide them.  
In practice this means that we may write `suppose $\phi$ is a closed predicate' instead of `suppose $\phi$ is a closed predicate in the syntax of \QLogic over $\Sigma$ and $\mathcal M$'.
But note next point.
\item\label{item.model.matters}
Our syntax from Definition~\ref{defn.predicate.syntax} and Figure~\ref{fig.predicate.syntax} contains denotational elements; $v\in\Val$ and $\varkappa\in\THREE^{\tf{Ctx}(\Sigma,\mathcal M)}$.

So when we write `suppose $\phi$ is a closed predicate' then this statement is perfectly correct and well-defined but it \emph{really does matter} what the signature and model are (even if the \emph{precise} choice of model will may not matter to the theorems we prove), because closed predicates $\phi$ can include elements from the denotation such as $\varkappa\in\THREE^{\tf{Ctx}(\Sigma,\mathcal M)}$ (this function-space is determined by both the signature $\Sigma$ and the model $\mathcal M$) and $v\teq v'$ for some $v,v'\in\Val$ ($\Val$ is determined by the choice of model $\mathcal M$).

Adding denotational elements to syntax avoids us having to put this information in the context (which is already a $4$-tuple, as per Definition~\ref{defn.contexts}(\ref{item.contexts})) and simplifies the presentation.
This is a technical device with a long pedigree.\footnote{Adding elements from a model to syntax, and conversely building a model out of syntax, is a standard technical device in logic, with a pedigree that goes back (at least) to the proofs of the upward and downward L\"owenheim-Skolem theorems from the 1920s.  A clear survey of these methods is in Chapter~2 (``Models constructed from constants'') of~\cite{chang:modt}.}

Note that (by design!) the choice of model makes a difference to validity.
For example: in Definition~\ref{defn.phi.equivalent} we define equivalence of predicates (and we explicitly mention $\Sigma$ and $\mathcal M$);
if $\mathcal M$ has only a single value $v$, then $\tall a,b.a\teq b$ and $\tall a.a\teq v$ are equivalent to $\tvT$; and if $\mathcal M$ has more than one value then $\tall a,b.a\teq b$ is equivalent to $\tvF$ --- $\tall a.\phi$ is shorthand for $\tneg\texi a.\tneg\phi$; see Figure~\ref{fig.3.derived}.
In contrast: in Lemma~\ref{lemm.equiv} we elide $\Sigma$ and $\mathcal M$ for brevity, and just write `suppose $\phi$ and $\phi'$ are closed predicates and $\acontext$ is a context', without specifying over what signature or model, because this is clear in context.

The reader can mentally add $\Sigma$ and $\mathcal M$ everywhere if they like.
We will always be precise and explicit where it counts.
\item\label{item.pedant.texi.v}
We may quantify over $v\in\Val$ in our formal syntax of predicates, writing things like $\texi v.\texi v'.v\teq v'$ as predicates.
Technically this is wrong and is not syntax! 
Figure~\ref{fig.predicate.syntax} permits e.g. $\texi a.\texi a'.a\teq a'$ where $a,a'\in\tf{VarSymb}$ (Definition~\ref{defn.signature}(\ref{item.signature.varsymb})); it does not permit $\texi v$ in the formal syntax of predicates where $v\in\Val$, so writing e.g. $\texi v.\texi v'.v\teq v'$ as if it were formal syntax of \QLogic is an abuse of notation.

However, $a$ and $a'$ get instantiated to values $v$ and $v'$ in our denotation (see the line for $\texi$ in Figure~\ref{fig.3.phi.f}).
It seems not worth the disruption of maintaining a distinction between $\texi a.\phi$ (which is syntax, assuming $\phi$ is) and $\texi v.\phi[a\ssm v]$ (which looks like syntax but technically is not).
So in practice this seems to us one example of an abuse of formal notation that is positively useful and clarifies meaning.

If this bothers the reader, then every time they see $\texi v.\phi$ (and $\tall v.\phi$, and similar) in our text, they can either declare that $v$ here is \emph{actually} a variable symbol in $\tf{VarSymb}$ which we happen to have called $v$ (and not a value in $\Val$ after all), or they can pick a fresh variable symbol $a$ and replace every unbound $v$ in $\phi$ with $a$.
Either way they will get well-formed syntax and everything will be as it should be.

Notation exists to serve clarity, and not the other way around, so we will simply ignore this henceforth and carry on writing $\texi v.\phi$, just whenever we feel like it, and this will denote just what it says: `there exists a $v$ such that $\phi$'.
Meaning should always be clear. 

Note that an unbound value $v$ in a predicate is legal syntax as per Figure~\ref{fig.predicate.syntax}; 
e.g. in `$\texi a.a\teq v$', $v$ is a value, and this is fine; the (constructive and useful) abuse of notation comes when we write $\texi v$. 
\end{enumerate}
\end{rmrk}

\subsection{Denotation (semantics) of the syntax}
\label{subsect.denotation}


\begin{figure} %
$$
\begin{array}{r@{\ }l@{\quad}l}
\modellabel{\tvB}{n,p,O,\avaluation} =&\tvB
\\
\modellabel{v\teq v'}{n,p,O,\avaluation} =&
\begin{cases}
\tvT & v=v'
\\
\tvF & v\neq v'
\end{cases}
\\
\modellabel{\tf P(v_1,\dots,v_{\arity(\tf P))}}{n,p,O,\avaluation} =&\avaluation(\tf P)_{n,p}(v_1,\dots,v_{\arity(\tf P)}) %
\\[2ex]
\modellabel{\phi\tor\phi'}{n,p,O,\avaluation} =& \modellabel{\phi}{n,p,O,\avaluation}\tor\modellabel{\phi'}{n,p,O,\avaluation}
\\
\modellabel{\tneg\phi}{n,p,O,\avaluation} =& \tneg\,\modellabel{\phi}{n,p,O,\avaluation}  
\\
\modellabel{\modT\phi}{n,p,O,\avaluation} =& \modT\modellabel{\phi}{n,p,O,\avaluation}  
\\
\modellabel{\texi\thea.\phi}{n,p,O,\avaluation} =& \bigvee_{v{\in}\Val} \modellabel{\phi[\thea\ssm v]}{n,p,O,\avaluation}
\\
\modellabel{\someone\phi}{n,p,O,\avaluation} =& \bigvee_{p'\in O} \modellabel{\phi}{n,p',O,\avaluation} 
\\
\modellabel{\Quorum\phi}{n,p,O,\avaluation} =& \bigvee_{O'\in\opensne} \modellabel{\phi}{n,p,O',\avaluation} 
\\[2ex]
\modellabel{\yesterday\phi}{n,p,O,\avaluation} =&
\bigvee_{\f{max}(0,n\minus 1)\leq n'<n}\modellabel{\phi}{n',p,O,\avaluation}
\\
\modellabel{\tomorrow\phi}{n,p,O,\avaluation} =&
\modellabel{\phi}{n\plus 1,p,O,\avaluation}
\\[2ex]
\modellabel{\mu\tf X.\phi}{n,p,O,\avaluation}=&
\bigl(\bigwedge\{ \varkappa\in\THREE^{\tf{Ctx}}\mid \model{\phi[\tf X\ssm\varkappa]}\leq\varkappa\}\bigr)(n,p,O,\avaluation)
\figskip
\modellabel{\varkappa}{n,p,O,\avaluation}=&
\varkappa(n,p,O,\avaluation)
\end{array}
$$ 
\caption{Denotational semantics of \QLogic (Definition~\ref{defn.denotation})}
\label{fig.3.phi.f}
\end{figure}


\begin{figure} %
$$
\begin{array}{r@{\ }l@{\ \ }r@{\ }l}
\tvF =& \modT\tvB
&
\modellabel{\tvF}{\acontext} =&\tvF
\figskip
\tvT =& \tneg\tvF
&
\modellabel{\tvT}{\acontext} =&\tvT
\\[2ex]
\modF\phi =& \modT\tneg\phi
&
\modellabel{\modF\phi}{\acontext} =&\modF\modellabel{\phi}{\acontext}
\figskip
\modB\phi =& \tneg\modT\phi\tand\tneg\modF\phi
&
\modellabel{\modB\phi}{\acontext} =&\modB\modellabel{\phi}{\acontext}
\figskip
\modTB\phi =& \tneg\modT\tneg\phi
&
\modellabel{\modTB\phi}{\acontext} =&\modTB\modellabel{\phi}{\acontext}
\figskip
\modTF\phi =& \modT\phi\tor\modF\phi
&
\modellabel{\modTF\phi}{\acontext} =&\modTF\modellabel{\phi}{\acontext}
\figskip
\modFB\phi =& \modF\phi\tor\modB\phi
&
\modellabel{\modFB\phi}{\acontext} =&\modFB\modellabel{\phi}{\acontext}
\\[2ex]
\phi\tand\phi' =& \tneg(\tneg\phi\tor \tneg\phi')
&
\modellabel{\phi\tand\phi'}{\acontext} =&\modellabel{\phi}{\acontext}\tand\modellabel{\phi'}{\acontext}
\figskip
\phi\tnotor\phi' =& \tneg\phi\tor\phi'
&
\modellabel{\phi\tnotor\phi'}{\acontext} =&\modellabel{\phi}{\acontext}\tnotor\modellabel{\phi'}{\acontext}
\figskip
\phi\timpc\phi' =& \phi\tnotor\modT\phi'
&
\modellabel{\phi\timpc\phi'}{\acontext} =&\modellabel{\phi}{\acontext}\timpc\modellabel{\phi'}{\acontext}
\figskip
\phi\equiv\phi' =& (\modT\phi\tand\modT\phi')\tor(\modB\phi\tand\modB\phi') \tor(\modF\phi\tand\modF\phi') \hspace{-5em}
\\
&&
\modellabel{\phi\equiv\phi'}{\acontext} =&\modellabel{\phi}{\acontext}\equiv\modellabel{\phi'}{\acontext}
\\[2ex]
\tall a.\phi =& \tneg\texi a.\tneg\phi
&
\modellabel{\tall a.\phi}{\acontext} =&\bigwedge_{v{\in}\Val}\modellabel{\phi[a\ssm v]}{\acontext}
\figskip
\everyone \phi =& \tneg\someone\tneg\phi
&
\modellabel{\everyone\phi}{n,p,O,\avaluation} =&\bigwedge_{p'{\in}\Pnt}\modellabel{\phi}{n,p',O,\avaluation}
\figskip
\Coquorum\phi =& \tneg\Quorum\tneg\phi
&
\modellabel{\Coquorum\phi}{n,p,O,\avaluation} =&\bigwedge_{O'{\in}\opensne}\modellabel{\phi}{n,p,O',\avaluation}
\\[2ex]
\QuorumBox\phi =& \Quorum\everyone\phi
&
\modellabel{\QuorumBox\phi}{n,p,O,\avaluation} =&\bigvee_{O'{\in}\opensne}\bigwedge_{p'\in O'}\modellabel{\phi}{n,p',O',\avaluation}
\figskip
\CoquorumDiamond\phi =& \Coquorum\someone\phi
&
\modellabel{\CoquorumDiamond\phi}{n,p,O,\avaluation} =&\bigwedge_{O'{\in}\opensne}\bigvee_{p'\in O'}\modellabel{\phi}{n,p',O',\avaluation}
\figskip
\someoneAll\phi =& \Quorum\someone\phi
&
\modellabel{\someoneAll\phi}{n,p,O,\avaluation} =&\bigvee_{O'{\in}\opensne}\bigvee_{p'\in O'}\modellabel{\phi}{n,p',O',\avaluation}
\figskip
\everyoneAll\phi =& \Coquorum\everyone\phi
&
\modellabel{\everyoneAll\phi}{n,p,O,\avaluation} =&\bigwedge_{O'{\in}\opensne}\bigwedge_{p'\in O'}\modellabel{\phi}{n,p',O',\avaluation}
\\[2ex]
\forever\phi =& \mu\tf X.\tomorrow(\phi\tand \tf X)
&
\modellabel{\forever\phi}{n,p,O,\avaluation} =&\bigwedge_{n'> n}\modellabel{\phi}{n',p,O,\avaluation}
\figskip
\sometime\phi =& \mu\tf X.\tomorrow(\phi\tor \tf X)
&
\modellabel{\sometime\phi}{n,p,O,\avaluation} =&\bigvee_{n'> n}\modellabel{\phi}{n',p,O,\avaluation}
\figskip
\infinitely\phi =& \forever\sometime\phi
&
\modellabel{\infinitely\phi}{n,p,O,\avaluation} =&\bigwedge_{n'> n}\bigvee_{n''>n'}\modellabel{\phi}{n'',p,O,\avaluation}
\figskip
\final\phi =& \sometime\forever\phi
&
\modellabel{\final\phi}{n,p,O,\avaluation} =&\bigvee_{n'> n}\bigwedge_{n''>n'}\modellabel{\phi}{n'',p,O,\avaluation}
\figskip
\recent\phi =& \mu\tf X.\yesterday(\phi\tor\tf X)
&
\modellabel{\recent\phi}{n,p,O,\avaluation} =&\bigvee_{0\leq n'<n}\modellabel{\phi}{n',p,O,\avaluation}
\figskip
\urecent\phi =& \mu\tf X.(\phi\tor\yesterday\tf X)
&
\modellabel{\urecent\phi}{n,p,O,\avaluation} =&\bigvee_{0\leq n'\leq n}\modellabel{\phi}{n',p,O,\avaluation}
\end{array}
$$
$$
\begin{array}{r@{\ }l}
\mrup{a}{\phi}{v} =&
\mu\tf X.(\yesterday(\modT\phi[a\ssm v] \tor ((\tneg\modT\texi a.\phi) \tand \tf X)))
\figskip
\modellabel{\mrup{a}{\phi}{v}}{n,p,O,\avaluation} =&
\begin{cases}
\tvF 
&\Forall{0\leq n'<n} \modellabel{\texi a.\phi}{n',p,O,\avaluation}\neq\tvT
\\
\modT\modellabel{\phi[\thea\ssm v]}{\f{max},p,O,\avaluation}
&\Exists{0\leq n'<n} \modellabel{\texi a.\phi}{n',p,O,\avaluation}=\tvT
\end{cases}
\\
&\qquad\text{where}\quad
\f{max}=\bigvee\{0\leq n'<n \mid \modellabel{\texi a.\phi}{n',p,O,\avaluation}=\tvT\}
\end{array}
$$ 
\caption{Derived syntactic sugar expressions (Definition~\ref{defn.3.derived})}
\label{fig.3.derived}
\end{figure}



Recall from Definition~\ref{defn.signature}(\ref{item.predicate.syntax.signature}) the notion of a \emph{signature} and from Definition~\ref{defn.model} the notion of a \emph{model}:
\begin{defn}[Denotation]
\label{defn.denotation}
Suppose $\Sigma=(\PredSymb,\arity)$ is a signature, and $\mathcal M$ is a model. 
Suppose $\phi$ is a closed predicate (Definition~\ref{defn.predicate.syntax}(1)).
Define the \deffont{denotation} of $\phi$ 
$$
\model{\phi}\in\THREE^{\tf{Ctx}(\Sigma,\mathcal M)}
\quad\text{such that}\quad 
\acontext \longmapsto \modellabel{\phi}{\acontext}
$$
by the rules in Figure~\ref{fig.3.phi.f}.

We will usually treat $\Sigma$ and $\mathcal M$ as fixed parameters, so we may write $\tf{Ctx}(\Sigma,\mathcal M)$ just as $\tf{Ctx}$.
\end{defn}




\begin{defn}[Derived expressions]
\label{defn.3.derived}
The syntax of \QLogic in Figure~\ref{fig.predicate.syntax} and Definition~\ref{defn.predicate.syntax} is very expressive.
In the left-hand column of Figure~\ref{fig.3.derived} we define some derived expressions --- i.e. \emph{syntactic sugar}.
For clarity and convenience we unpack their denotation in the right-hand column.

In principle we could fully unpack our sugar down to the primitive syntax from Figure~\ref{fig.3.phi.f}, but in practice we work directly with Figure~\ref{fig.3.derived} because this is more convenient.
We will discuss syntax further below, notably in Remarks~\ref{rmrk.initial.comments.denotation}, \ref{rmrk.further.comments.denotation}, and~\ref{rmrk.comments.on.semantics}.
\end{defn}


\begin{rmrk}
Usually when we talk about denotation as defined in Definition~\ref{defn.denotation} above, the signature $\Sigma$ and model $\mathcal M$ are either understood or do not matter, so we may not mention them explicitly.
Meaning will always be clear.
\end{rmrk}

\begin{rmrk}[Initial comments on the denotation]
\label{rmrk.initial.comments.denotation}
We make some initial comments on Figure~\ref{fig.3.phi.f}:
\begin{enumerate*}
\item\label{item.first.mu}
$\mu\tf X.\phi$ is a \emph{fixedpoint operator}.
This is a standard method to include the power of recursion / iteration in a logic; an instance of $\tf X$ inside $\mu\tf X.\phi$ can be thought of as a recursive call back to $\phi$.
$\mu\tf X.\phi$ is what would be written $\mathrm{LFP}_{\tf X}\phi$ in~\cite[Chapter~4]{immerman:desc} (a Wikipedia article is also accessible and clear~\cite{wiki:Fixedpoint_logic}).
We will discuss this in Subsection~\ref{subsection.muX}.
\item
In the clause $\model{\tvB}=\tvB$ in Figure~\ref{fig.3.phi.f}, the $\tvB$ on the left is the syntax symbol from Figure~\ref{fig.predicate.syntax}, and the $\tvB$ on the right is $\tvB$-the-truth-value from Definition~\ref{defn.THREE}.
\item
In the clause for $\model{v\teq v'}$ in Figure~\ref{fig.3.phi.f}, the $\teq$ on the left is the syntax symbol from Figure~\ref{fig.predicate.syntax}, and the $\tvT$, $\tvF$, and $=$ on the right are truth-values and (actual) equality.
\item
Similarly for the clauses for $\model{\tneg\phi}$ and $\model{\phi\tor\phi'}$ and $\model{\modT\phi}$; the symbols $\tneg$, $\tor$, and $\modT$ on the left are syntax, and the symbols $\tneg$, $\tor$, and $\modT$ on the right refer to functions on truth-values as per Figure~\ref{fig.3}.
\end{enumerate*} 
\end{rmrk}


\begin{rmrk}[Further comments] 
\label{rmrk.further.comments.denotation}
\leavevmode
\begin{enumerate}
\item\label{item.interpret.yesterday}
The clause 
$\modellabel{\yesterday\phi}{n,p,O,\avaluation} = \bigvee_{\f{max}(0,n\minus 1)\leq n'<n}\modellabel{\phi}{n',p,O,\avaluation}$ is just a one-line encoding of
$$
\modellabel{\yesterday\phi}{n,p,O,\avaluation}=
\begin{cases}
\tvF & n=0
\\
\modellabel{\phi}{n\minus 1,p,O,\avaluation} & n>0 . 
\end{cases}
$$
\item
We see from the right-hand column in Figure~\ref{fig.3.derived} that $\recent\phi$ returns the greatest truth-value that $\phi$ has taken in the past, and $\urecent\phi$ returns the greatest truth-value that $\phi$ has taken today or in the past.

The notation $\urecent$ (`$\recent$', underlined) is a pun\footnote{A bad one perhaps, and the reader may or may not find this pun-acceptable.} on $<$ and $\leq$ (`$<$', underlined).
\item
Informally, we see that the only naturally \emph{non-monotone} operator in the (primitive or derived) syntax is $\tneg$, which is anti-monotone.
This design feature is what powers Lemma~\ref{lemm.positive.monotonicity}.
\end{enumerate}
\end{rmrk}


\begin{nttn}
\label{nttn.logic.terminology}
We take a moment to recall some standard conventions and terminology:
\begin{enumerate*}
\item
Implications $\tnotor$ and $\timpc$ associate to the right, so that $\phi\timpc\psi\timpc\chi$ means $\phi\timpc(\psi\timpc\chi)$.
\item
Unary connectives (like $\tneg$, $\modT$, and $\someone$) bind more strongly than binary connectives (like $\tand$ and $\tnotor$), so that $\QuorumBox\tneg\phi\tand\psi$ means $(\QuorumBox\tneg\phi)\tand\psi$.
\item
The scope of a quantification $\texi a$ and $\tall a$ extends as far to the right as possible. 
Thus $\tall\thea.\phi\tand\psi$ means $\tall\thea.(\phi\tand\psi)$.
Note however that $\QuorumBox\texi\thea.\phi\tand\psi$ still means $(\QuorumBox(\texi\thea.\phi))\tand\psi$, because the $\texi$ quantifier is trapped in the scope of the $\QuorumBox$. 
\item
All the above said, we will (as is usual) use spacing and some common sense to convey scope.
For example 
$$
\texi\thea.\phi\mathbin{\tnotor}\psi
\quad\text{means}\quad
\texi\thea.(\phi\tnotor\psi), 
$$
and
$$
\texi\thea.\phi \ \tnotor \psi
\quad\text{means}\quad
(\texi\thea.\phi)\tnotor\psi. 
$$
In practice, meaning will always be clear.
\item
\begin{enumerate*}
\item\label{item.nullary.ps}
If $\arity(\tf P)=0$ then we may call $\tf P$ \deffont{nullary} or a \deffont{predicate constant symbol}.\footnote{$\tvT$, $\tvB$, and $\tvF$ are predicate constant symbols, but they are special so we give them a distinct identity in the syntax.}
\item\label{item.unary.ps}
If $\arity(\tf P)=1$ then we may call $\tf P$ a \deffont{unary} predicate symbol.
\item\label{item.atomic.predicate}
We may call a predicate of the form $\tf P(v_1,\dots,v_n)$ or $\tneg\tf P(v_1,\dots,v_n)$ an \deffont{atomic predicate}. 
\end{enumerate*}
\end{enumerate*} 
\end{nttn}



\subsection{The fixedpoint operator $\mu \tf X$}
\label{subsection.muX}

\begin{rmrk}
\label{rmrk.comments.on.semantics}
As per Remark~\ref{rmrk.initial.comments.denotation}(\ref{item.first.mu}), our fixedpoint operator $\mu\tf X.\phi$ is a method to include the power of recursion / iteration in the logic.

A mention of $\tf X$ inside $\phi$ in $\mu\tf X.\phi$ can be thought of as a recursive call back to $\phi$, and we can think of $\mu\tf X.\phi$ as encoding an iterative loop that repeatedly executes a computation encoded in $\phi$, and loops if it encounters $\tf X$ in $\phi$.

With this intuition, we can look at some of the derived syntax in Figure~\ref{fig.3.derived} that uses $\mu$:
\begin{enumerate}
\item
$\recent\phi=\mu\tf X.\yesterday(\phi\tor\tf X)$ can be imagined as the following loop:
\begin{enumerate*}
\item
Set an accumulator to $\tvF$.\footnote{We mean `accumulator' in the sense of an `accumulator variable' inside a machine that (we imagine) is running this pseudocode.  This is not a technical term from logic.}
\item\label{recent.step.0}
If we are at day~0 then terminate and return the value of the accumulator; otherwise go to yesterday.
\item
Compute the truth-value of $\phi$, and $\tor$ this truth-value with the accumulator.
\item
Loop back to step~\ref{recent.step.0}.
\end{enumerate*} 
\item
$\forever\phi=\mu\tf X.\tomorrow(\phi\tand\tf X)$ can be viewed as the following loop:
\begin{enumerate*}
\item
Set an accumulator to $\tvT$.
\item\label{forever.step.0}
Go to tomorrow.
\item
Compute the truth-value of $\phi$, and $\tand$ this truth-value with the accumulator.
\item
Loop back to step~\ref{forever.step.0}.
\item
After infinite time (see Remark~\ref{rmrk.oo} about why that makes sense) return the limit value of the accumulator.
\end{enumerate*} 
\item
$\sometime\phi=\mu\tf X.\tomorrow(\phi\tor\tf X)$ can be viewed as executing the following loop:
\begin{enumerate*}
\item
Set an accumulator to $\tvF$.
\item\label{sometime.step.0}
Go to tomorrow.
\item
Compute the truth-value of $\phi$ and $\tor$ with the accumulator.
\item
Loop back to step~\ref{sometime.step.0}.
\item
After infinite time, return the limit value of the accumulator.
\end{enumerate*}
\item\label{item.mru.program}
The clause 
$\mrup{a}{\phi}{v} = \mu\tf X.(\yesterday(\modT\phi[a\ssm v] \tor ((\tneg\modT\texi a.\phi) \tand \tf X)))$
looks more imposing, but the idea is simply to 
\begin{itemize*}
\item
iterate back in time looking for when $\texi a.\phi$ has truth-value $\tvT$, i.e. when \emph{some} value makes $\phi(a)$ be true, and 
\item
if we find such a time then we check if \emph{at that time} $\phi[a\ssm v]$ has truth-value $\tvT$ and 
\item
only if both steps above succeed, do we return $\tvT$.
\end{itemize*}
Thus, we return $\tvT$ when $v$ is one of the values that \emph{most recently} makes $\phi(a)$ be true.

In detail:
\begin{enumerate*}
\item\label{mu.step.0}
If we are at day~0, terminate with value $\tvF$; otherwise go to yesterday.
\item
If $\phi[a\ssm v]$ has truth-value $\tvT$ then terminate with value $\tvT$; otherwise continue to the next step.
\item
If $\texi a.\phi$ has truth-value $\tvT$ then terminate with value $\tvF$; otherwise go to step~\ref{mu.step.0}.
\end{enumerate*}
\end{enumerate}
\end{rmrk}

\begin{rmrk}[Comment on the omniscient observer]
\label{rmrk.oo}
The programs in Remark~\ref{rmrk.comments.on.semantics} above are executed by an omniscient observer who stands outside of the model's notion of time and space.
This is why (for example) $\forever$ can run an infinite loop --- the omniscient observer can observe it because the observer is looking at the model from outside.\footnote{Let's make this concrete.  Consider the natural numbers $\Ngeqz$.  We can easily observe this and note that $n\plus 1\gneq n$ for every $n\in\Ngeqz$; whereas a local observer trapped on the number line and constrained to perceive this line as the passage of time can make similar observation, but from the point of view of that observer, the calculation involved is neverending.  

Note that even if we could speak with this local observer, it would be no good to appeal to it to `just realise that it is on the number line', because in order to check this to a mathematical certainty the observer would have to \dots traverse all of its local notion of time to \emph{check}, amongst other things, that $n\plus 1\gneq n$.  After all, the observer might be on $\mathbb Z/p$ for some astronomically enormous $p$, i.e. their time might be on a large loop.  

The reader interested in a branch of mathematics explicitly devoted to studying these questions, and who is not afraid by the possibility of ending up seriously doubting the nature of their own reality, can start by reading up on nonstandard models of arithmetic.} 
\end{rmrk}

\begin{rmrk}
The rest of this Subsection will show how our denotation for $\mu$ from Figure~\ref{fig.3.phi.f} does indeed yield a fixedpoint operator.
The culminating result is Proposition~\ref{prop.mu.fixedpoint} (which just says that $\mu$ is indeed a fixedpoint operator) but the key technical lemma is Lemma~\ref{lemm.positive.monotonicity} (monotonicity of positive predicates).

The reader is welcome to read the rest of this Subsection, but this is all just a standard use of the Knaster-Tarski fixedpoint theorem~\cite[Theorem~2.35, page~50]{priestley:intlo}.

In fact, the real design effort has already happened: we designed the derived syntax of Figure~\ref{fig.3.derived} so that it would respect the positivity condition of Definition~\ref{defn.positive}, without having to fully expand it out.
For instance, we defined $\modTB\phi$ to be $\tneg\modT\tneg\phi$ (which clearly puts $\phi$ in a positive position) instead of $\modT\phi\tor\modB\phi$ (which is logically equivalent, but which does not).
\end{rmrk} 



\begin{defn}
\label{defn.three.ctx.lattice}
Recall from Remark~\ref{rmrk.three.elementary}(\ref{item.three.elementary.bounded.lattice}) that $\THREE$ with ordering $\tvF<\tvB<\tvT$ is a finite lattice.
As is standard, the function space $\THREE^{\tf{Ctx}}$ (functions from $\tf{Ctx}$ to $\THREE$) inherits this ordering and lattice structure \deffont{pointwise}, as indicated in Figure~\ref{fig.chi.3.lattice}.

This makes $\THREE^{\tf{Ctx}}$ into a \deffont{bounded complete lattice}, meaning that it has a least and greatest element, and arbitrary joins (least upper bounds) and meets (greatest lower bounds).\footnote{Having a least and greatest element follows from having arbitrary joins and meets; the least element is the join of the empty set, and the greatest element is the meet of the empty set.  So in fact a complete lattice is already a bounded complete lattice.} 
\end{defn}


\begin{figure}
$$
\begin{array}{r@{\ }l}
\tbot =& \lambda\acontext.\tvF
\\
\ttop =& \lambda\acontext.\tvT
\\
\bigvee\mathcal X =& \lambda\acontext.\bigvee\{ \varkappa\,\acontext \mid \varkappa\in\mathcal X\}
\\
\bigwedge\mathcal X =& \lambda\acontext.\bigwedge\{ \varkappa\,\acontext \mid \varkappa\in\mathcal X\}
\\[2ex]
\varkappa\leq\varkappa' \liff& \Forall{\acontext} \varkappa\,\acontext\leq\varkappa'\,\acontext
\end{array}
$$
Above, $\mathcal X\subseteq \THREE^{\tf{Ctx}}$ and $\varkappa,\varkappa'\in\THREE^{\tf{Ctx}}$. 
\caption{Bounded complete lattice structure on the function space $\THREE^{\tf{Ctx}}$}
\label{fig.chi.3.lattice}
\end{figure}

\begin{rmrk}
\label{rmrk.leq.phi}
Our denotation $\model{\text{-}}$ from Definition~\ref{defn.denotation} and Figure~\ref{fig.3.phi.f} maps predicates to functions in $\THREE^{\tf{Ctx}}$.
One particular consequence of Definition~\ref{defn.three.ctx.lattice} is that (denotations of) predicates are naturally ordered, and we can write 
$$
\model{\phi}\leq\model{\psi}
\quad\text{which just means}\quad
\Forall{\acontext\in\tf{Ctx}}\modellabel{\phi}{\acontext}\leq\modellabel{\psi}{\acontext}.
$$
We use this fact for example in Lemma~\ref{lemm.positive.monotonicity}.
We can also use this fact, and the complete lattice structure on $\THREE^{\tf{Ctx}}$ from Definition~\ref{defn.three.ctx.lattice}, to simplify the presentation of the clause for $\mu\tf X.\phi$ in Figure~\ref{fig.3.phi.f} as follows:
$$
\model{\mu\tf X.\phi}=
\bigwedge\{ \varkappa\in\THREE^{\tf{Ctx}}\mid \model{\phi[\tf X\ssm\varkappa]}\leq\varkappa\} .
$$
\end{rmrk}

\begin{defn}
\label{defn.positive}
Recall the predicate syntax from Figure~\ref{fig.predicate.syntax} %
and suppose $\phi$ is a predicate.
Then:
\begin{enumerate*}
\item\label{item.positive.occurrence}
Call a free occurrence of $\tf X$ in $\phi$ 
\begin{itemize*}
\item
\deffont{positive} when the occurrence appears under an even number (which may be $0$) of negations $\tneg$, and 
\item
\deffont{negative} when the occurrence appears under an odd number of negations.
\end{itemize*}
For example: all free occurrences of $\tf X$ are positive in $\tf X$ and $\mu\tf X.\tneg\tf X$ and $\yesterday\tf X$; and the free occurrence of $\tf X$ in $\tneg\tf X$ is negative. 
\item
Call a predicate $\phi$ \deffont{positive} when every free occurence of $\tf X$ in $\phi$ is positive.
\end{enumerate*}
\end{defn}


\begin{lemm}
\label{lemm.positive.monotonicity}
Suppose $\mu\tf X.\phi$ is a closed predicate.
Then:
\begin{enumerate*}
\item
If $\phi$ is positive then $\model{\phi}$ is monotone as a function from $\THREE^{\tf{Ctx}}$ with the partial order $\leq$ from Figure~\ref{fig.chi.3.lattice}, to itself.
In symbols, for every $\varkappa,\varkappa'\in\THREE^{\tf{Ctx}}$:
$$
\varkappa\leq\varkappa'
\quad\text{implies}\quad
\model{\phi[\tf X\ssm\varkappa]}\leq\model{\phi[\tf X\ssm\varkappa']} . 
$$
\item
If $\tneg\phi$ is positive then $\model{\phi}$ is antimonotone as a function from $\THREE^{\tf{Ctx}}$ with the partial order $\leq$ from Figure~\ref{fig.chi.3.lattice}, to itself.
In symbols, for every $\varkappa,\varkappa'\in\THREE^{\tf{Ctx}}$:
$$
\varkappa\leq\varkappa'
\quad\text{implies}\quad
\model{\phi[\tf X\ssm\varkappa']}\leq\model{\phi[\tf X\ssm\varkappa]} . 
$$
\end{enumerate*}
(The meaning of $\model{\text{-}}\leq\model{\text{-}}$ is unpacked in Remark~\ref{rmrk.leq.phi}.)

\end{lemm}
\begin{proof}
Both parts follow by a routine simultaneous induction on the syntax of $\phi$, using the rules in Figure~\ref{fig.3.phi.f}.
\end{proof}

\begin{prop}
\label{prop.mu.fixedpoint}
Suppose $\mu\tf X.\phi$ is a closed predicate and $\phi$ is positive.
Then:
\begin{enumerate*}
\item
$\model{\mu\tf X.\phi}$ is a least fixedpoint of the function $\varphi:\THREE^{\tf{Ctx}}\to\THREE^{\tf{Ctx}}$ given by
$$
\varphi(\varkappa)=\lambda\acontext.\modellabel{\phi[\tf X\ssm\varkappa]}{\acontext}. 
$$
\item
If $\acontext$ is a context then
$$
\modellabel{\mu\tf X.\phi}{\acontext}=\modellabel{\phi[\tf X\ssm \model{\mu\tf X.\phi}]}{\acontext}.
$$
\end{enumerate*}
\end{prop}
\begin{proof}
By the Knaster-Tarski fixedpoint theorem for $\THREE^{\tf{Ctx}}$ as a bounded lattice as per Definition~\ref{defn.three.ctx.lattice}.
A concise but exceptionally clear presentation is in~\cite[Theorem~2.35, page~50]{priestley:intlo} (a Wikipedia article~\cite{wiki:Knaster-Tarski_theorem} is also accessible and clear). 
\end{proof}



\subsection{Validity}



We defined in Definition~\ref{defn.tv.ment} what it means for a truth-value to be \emph{valid}; now we extend this notion to predicates-in-context: 
\begin{defn}[Validity]
\label{defn.validity.judgement}
Suppose that:
\begin{itemize*}
\item
$\Sigma$ is a signature (Definition~\ref{defn.signature}(\ref{item.predicate.syntax.signature})).
We take $\Sigma$ to be fixed for the rest of this Definition, which means that we do not explicitly include it (e.g. as a subscript) in the notation.
\item
$\mathcal M$ is a model (Definition~\ref{defn.model}).
Again, we fix $\mathcal M$ as a parameter for the rest of this definition.
\end{itemize*}
Suppose further that:
\begin{itemize*}
\item
$\acontext=(n,p,O,\avaluation)$ is a context in $\tf{Ctx}(\Sigma,\mathcal M)$, so that by Definition~\ref{defn.contexts}(\ref{item.contexts}) $n\in\Time$ and $p\in\Pnt$ and $O\in\opensne$ and $\avaluation$ is a valuation to $\mathcal M$.
\item
$\phi$ is a closed predicate in $\Sigma$ and $\mathcal M$ (cf. Remark~\ref{rmrk.pedant.point}(\ref{item.model.matters})).
\end{itemize*}
Then we define: %
\begin{enumerate*}
\item\label{item.ctx.ment.phi}
$\acontext\ment\phi$ means $\ment\modellabel{\phi}{\acontext}$, which by Definition~\ref{defn.tv.ment}(\ref{item.tv.ment}) means $\modellabel{\phi}{\acontext}\in\threeValid$, which by Definition~\ref{defn.tv.ment}(\ref{item.tv.ment.correct}) means $\modellabel{\phi}{\acontext}=\tvT\lor\modellabel{\phi}{\acontext}=\tvB$.

Because this is a core definition, we sum this up for the reader's clear reference:
$$
\acontext\ment\phi
\quad\liff\quad
{\ment\modellabel{\phi}{\acontext}}
\quad\liff\quad
\modellabel{\phi}{\acontext}\in\threeValid
\quad\liff\quad
\modellabel{\phi}{\acontext}=\tvT\lor\modellabel{\phi}{\acontext}=\tvB.
$$
We read this as \deffont{$\phi$ is valid in the context $\acontext$.}
\item
$n,p\mentval\phi$ means $\Forall{O\in\opensne}(n,p,O\mentval\phi)$.
\item\label{item.ment.n}
$n\mentval\phi$ means $\Forall{p\in\Pnt}\Forall{O\in\opensne}(n,p,O\mentval\phi)$.
\item\label{item.mentval}
$\mentval\phi$ means $\Forall{n\in\Time}\Forall{p\in\Pnt}\Forall{O\in\opensne}(n,p,O\mentval\phi)$.
\item\label{item.ment.phi}
$\ment\phi$ means $\Forall{\acontext\in\tf{Ctx}(\Sigma,\mathcal M)}(\acontext\ment\phi)$.
\end{enumerate*}
\end{defn}

\begin{figure}
$$
\begin{array}{l@{\quad}l@{\quad}l}
\mentval\Theta & \text{means} & \Forall{\theta\in\Theta}({\mentval\theta})
\figskip
\Theta\ment\phi & \text{means} &
\Forall{\avaluation\in\Valuation(\Sigma,\mathcal M)}({\mentval\Theta})\limp ({\mentval\phi}) 
\end{array}
$$
\caption{Validity of $\phi$ in a theory $\Theta$ (Definition~\ref{defn.theory.axiom})}
\label{fig.validity}
\end{figure}


Later on, when we start to build declarative axiomatisations of consensus algorithms, we will express these as \emph{theories}.
We include the relevant definition here, since it follows on naturally from Definition~\ref{defn.validity.judgement}: 
\begin{defn}[Validity in a theory]
\label{defn.theory.axiom}
Suppose that:
\begin{itemize*}
\item
$\Sigma$ is a signature and $\mathcal M=(\Pnt,\opens,\Val)$ is a model.
\item
$\phi$ is a closed predicate in $\Sigma$ and $\mathcal M$.
\item
$\Theta=\{\theta_1,\dots,\theta_n\}$ is a set of closed predicates.\footnote{It would be unusual for $\Theta$, which is intended to be a collection of axioms, to mention specific elements $v\in\Val$ or $\varkappa\in\THREE^{\tf{Ctx}(\Sigma,\mathcal M)}$ from a model (cf. Remark~\ref{rmrk.pedant.point}), but the maths does not seem to mind, so we do not bother to exclude this possibility.}
\end{itemize*}
Then:
\begin{enumerate*}
\item\label{item.theory}
A \deffont{theory} is a pair $(\Sigma,\Theta)$ of 
\begin{itemize*}
\item
a signature $\Sigma$, and 
\item
a set $\Theta$ of \deffont{axioms}, which are closed predicates over $\Sigma$ and $\mathcal M$.
\end{itemize*}
\item\label{item.axioms.valid}
Given a valuation $\avaluation\in\Valuation(\Sigma,\mathcal M)$ (Definition~\ref{defn.contexts}(\ref{item.contexts.valuation})), we define 
$$
\mentval\Theta
\quad\text{and}\quad
\Theta\ment\phi
$$ 
as per Figure~\ref{fig.validity}. 
When $\Theta\ment\phi$ holds, we say that $\phi$ is \deffont{valid in the theory $\Theta$} (in $\Sigma$ and $\mathcal M$). 
\end{enumerate*}
Examples of theories are in Figures~\ref{fig.ThyPaxOne} and~\ref{fig.ThySPax}.
\end{defn}


\begin{rmrk}
\label{rmrk.thy.quantifier.scope}
Note the scope of the quantifiers in Definition~\ref{defn.theory.axiom}:
\begin{enumerate*}
\item
The assertion $\Theta\ment\phi$ does \emph{not} mean 
$$
\Forall{\avaluation}(\mentval(\theta_1\tand\dots\tand\theta_n) \timp \phi),
$$ 
which is equivalent to 
$$
\Forall{\acontext}(\acontext\ment(\theta_1\tand\dots\tand\theta_n)\limp \acontext\mentval\phi) .
$$ 
Above, the short double arrow `$\timp$' means `if the LHS is valid then so is the RHS'; see Remark~\ref{rmrk.rich.implication}(\ref{item.rich.implication.timp}).
The long double arrow `$\limp$' is just logical implication. 
\item
The assertion $\Theta\ment\phi$ means 
$$
\Forall{\avaluation}(\mentval\theta_1\tand\dots\tand\theta_n) \limp (\mentval\phi) .
$$
In words: if $\Theta$ holds everywhere in the model under the valuation $\avaluation$, then $\phi$ holds everywhere in the model under the valuation $\avaluation$.
\end{enumerate*}
We spell this out further.
The assertion $\Theta\ment\phi$ means:
$$
\Forall{\avaluation}(\Forall{n,p,O}n,p,O\mentval\theta_1\tand\dots\tand\theta_n) \limp (\Forall{n,p,O}n,P,O\mentval\phi),
$$
where above $n\in\Time$, $p\in\Pnt$, and $O\in\opensne$.

This pattern of quantifiers is what makes $\Theta$ a theory, rather than just a set of assumptions: as a theory, $\Theta$ is a restriction on the valuation $\avaluation$.
\end{rmrk}


\begin{rmrk}
One final pedant point: 
in Definition~\ref{defn.theory.axiom}(\ref{item.theory}) we defined a theory to be a pair $(\Sigma,\Theta)$ of a signature and a set of closed predicates over a signature $\Sigma$ \emph{and} over a model $\mathcal M$.
This technical convenience lets us write axioms that directly reference specific values in the model (instead of introducing explicit constant symbols to our signature, and interpreting them as \dots values in the model).

We will exploit to assume a distinguished `undefined' value $\udfn$ in axiom \rulefont{PaxPropose?} in Figure~\ref{fig.ThyPaxOne}, as per Definition~\ref{defn.ThyPaxOne}(\ref{item.ThyPaxOne.axioms}).
\end{rmrk}

\subsection{Comments on notation: eliding parameters, and functional notation}

\begin{nttn}[Eliding parameters]
\label{nttn.validity.irrelevant.context}
It will be quite common that we make validity judgements where it is clear that part of the context is irrelevant.
For example the reader can check the denotational semantics in Figure~\ref{fig.3.phi.f} and the definition of validity in Definition~\ref{defn.validity.judgement}(\ref{item.ctx.ment.phi}), to see that:
\begin{enumerate*}
\item
$\acontext \ment \tvB$ is valid; we do not need to know the values in $\acontext$. 
\item
The validity of $n,p,O\mentval \tf P(v)$ depends on $\avaluation$, $n$, and $p$ --- but not on $O$.
\item
The validity of $n,p,O\mentval\QuorumBox\phi$ depends on $n$ and $\avaluation$ (and of course on $\phi$), but it does not depend on the values of $p$ and $O$.
\end{enumerate*}
In such cases, we may omit parameters in the context if they do not matter --- we may even omit the entire context if appropriate.
Thus, we may write $\ment\tvT$, $\ment \thea\teq\thea$, and $n\mentval\QuorumBox\phi$.

This notation remains within the scope of Definition~\ref{defn.validity.judgement} --- which universally quantifies over any omitted parameter --- so omitting parameters in this way is not an abuse of notation.
However, if the reader sees us omit a value from a context, there may be two reasons for it:
\begin{enumerate*}
\item
It may indicate a universal quantification over all possible values.
For example, this is clearly what we intend when we write $\mentval\theta$ in Figure~\ref{fig.validity} (top-right).
\item
It may indicate that we have elided irrelevant parameters.
For example, this is clearly what we intend in the example of $\ment\tvT$ above.
\end{enumerate*}
In practice, it will always be clear which is intended and meaning will always be clear.\footnote{Actually, meaning will be \emph{clearer}.  

In early drafts of this paper, we dutifully included all parameters, whether relevant or not.  
But this turned out to be just too confusing: the notational overhead of carrying around values that we did not care about was burdensome; but furthermore, it was actually easier to understand the sequents once we have dropped irrelevant parameters, because this helped us to keep track of which parameters \emph{did} matter.}
\end{nttn}

\begin{figure}
$$
\begin{array}{l@{\quad}l@{\quad}l@{\quad}l@{\quad}l}
\texiunique  a.\tf{write}(a)
&
\QuorumBox\texi a.\tf{send}(a)
&
\texi a.\everyoneAll\tf{decide}(a)
&
\tall a.\modTF \tf{propose}(a)
&
\mrup{a}{(\someone\tf{accept}(a))}{v}
\\[1ex]
& \text{may be written as}
\\[1ex]
\texiunique \tf{write} 
&
\QuorumBox\texi\tf{send} 
&
\texi\everyoneAll\tf{decide}
&
\tall\modTF\tf{propose} 
&
\mru{(\someone\tf{accept})}{v}
\end{array}
$$
\caption{Abbreviations inspired by higher-order logic (Notation~\ref{nttn.hos})}
\label{fig.hos}
\end{figure}


\begin{nttn}[Functional notation]
\label{nttn.hos}
We will frequently apply quantifiers (like $\texi$, $\texiunique$, or $\tall$) to unary predicate symbols, possibly with modalities.
For example:
$$
\tall a.\QuorumBox\tf{decide}(a)
\quad\text{is}\quad\text{$\tall$ quantifying $a$ in\ $(\QuorumBox\tf{decide})$\ applied to $a$.}
$$ 
For clarity, we may use a \deffont{functional notation} inspired by higher-order logic (HOL) to elide the variable, as illustrated in Figure~\ref{fig.hos}.
Meaning will always be clear.

See also Remark~\ref{rmrk.functional.notation.in.paxos} for examples of this notation applied in an axiomatisation.
\end{nttn}






\section{Properties of predicates}
\label{sect.predicates}

This Section collects elementary technical results which we will need later.
Some of these results, and their proofs, are more technical than others, but they are not necessarily surprising if we look at what they mean in the model.
Mostly the statements are simple and it is obvious by common sense that they should work (simple is good, because it means greater productivity and fewer errors).

A typical example is Lemma~\ref{lemm.tomorrow.forever.recent}(\ref{item.tomorrow.forever.recent.1}): once we look at what it actually says about the model, this is that ``if $\phi$ is true today, then in the future it will have been true in the past''.
For saying this a mathematician earns a living? --- we might think.

What is less evident is \emph{which} elementary properties will be important later, nor is it obvious what definitions and structures make these properties all fit together nicely.
So the reader should note that the simplicity of this Section is a design feature, which was only attained in retrospect and after some work.

Broadly speaking, we start this Section with the most elementary results, and build complexity from there.
The reader is welcome to read this Section now, or to skip it and refer back to its results as they are used in later proofs.


\subsection{Equivalence of predicates, and some (de Morgan) equivalences}
\label{subsect.equivalence.of.predicates}

For the rest of this Section we make a fixed but arbitrary choice of signature $\Sigma$ (Definition~\ref{defn.signature}(\ref{item.predicate.syntax.signature})) and model $\mathcal M$ (Definition~\ref{defn.model}).
All predicates considered will be over $\Sigma$ and $\mathcal M$, unless otherwise stated.

We generalise equivalence $\equiv$ (Definition~\ref{defn.equivalent}) from truth-values, to predicates viewed as functions from contexts to truth-values:
\begin{defn}[Equivalence of predicates]
\label{defn.phi.equivalent}
Suppose $\phi$ and $\phi'$ are closed predicates.

Call $\phi$ and $\phi'$ \deffont{extensionally equivalent} when 
$$
\model{\phi}=\model{\phi'} .
$$
Unpacking Definition~\ref{defn.denotation}, this means 
$$
\Forall{\acontext\in\tf{Ctx}(\Sigma,\mathcal M)}(\modellabel{\phi}{\acontext}=\modellabel{\phi'}{\acontext}) .
$$
For brevity, we may write `extensionally equivalent' just as `equivalent' henceforth.
\end{defn}



\begin{lemm}
\label{lemm.equiv}
Suppose $\phi$ and $\phi'$ are closed predicates and $\acontext$ is a context.
Then: 
\begin{enumerate*}
\item\label{item.equiv.1}
$\acontext\ment\phi\equiv\phi'$ if and only if $\modellabel{\phi}{\acontext}=\modellabel{\phi'}{\acontext}$.
\item\label{item.equiv.2}
$\ment\phi\equiv\phi'$ if and only if $\model{\phi}=\model{\phi'}$ (i.e. $\phi$ and $\phi'$ are equivalent as per Definition~\ref{defn.phi.equivalent}).
\end{enumerate*}
\end{lemm}
\begin{proof}
We consider each part in turn:
\begin{enumerate}
\item
By Definition~\ref{defn.validity.judgement}(\ref{item.ctx.ment.phi}) $\acontext\ment\phi\equiv\phi'$ means $\modellabel{\phi\equiv\phi'}{\acontext}\in\threeValid$, and by the clause for $\equiv$ in Figure~\ref{fig.3.derived} this is if and only if $\modellabel{\phi}{\acontext}=\modellabel{\phi'}{\acontext}$.
\item
By Definition~\ref{defn.validity.judgement}(\ref{item.ment.phi}), $\ment\phi\equiv\phi'$ means precisely $\Forall{\acontext}\acontext\ment\phi\equiv\phi'$.
We use part~\ref{item.equiv.1} of this result.
\qedhere\end{enumerate}
\end{proof}


We will see many equivalences in what follows.
We start by using Definition~\ref{defn.phi.equivalent} and Lemma~\ref{lemm.equiv} to express some standard equivalences in logic, having to do with implication being expressible using negation and disjunction, and with \emph{de Morgan} dualities between modalities and connectives:
\begin{lemm}
\label{lemm.minimal.syntax}
For each of the equivalences $\f{LHS}\equiv\f{RHS}$ listed in Figure~\ref{fig.easy.equivalences}, we have that $\f{LHS}$ is extensionally equal to $\f{RHS}$ in the sense of Definition~\ref{defn.phi.equivalent}.
\end{lemm}
\begin{proof}
By routine calculations using Figures~\ref{fig.3.phi.f} and~\ref{fig.3.derived} and facts of the lattice $\THREE=(\tvF<\tvB<\tvT)$.
\end{proof}


\begin{figure}
$$
\begin{array}{r@{\ }l@{\quad}r@{\ }l@{\quad}r@{\ }l@{\quad}r@{\ }l}
\phi\tand\phi' \equiv& \tneg(\tneg\phi\tor\tneg\phi')
&
\phi\tor\phi' \equiv& \tneg(\tneg\phi\tand\tneg\phi')
&
\tneg\phi\equiv& \phi\timpc\tvF \equiv\phi\tnotor\tvF %
\\
\phi\timpc\phi' \equiv& \phi\tnotor\modT\phi'
&
\phi\tnotor\phi' \equiv& \tneg\phi\tor\phi'
&
\phi\timpc\phi' \equiv& \phi\tnotor\modT\phi'
\\
\everyone\phi\equiv&\tneg\someone\tneg\phi
&
\someone\phi\equiv&\tneg\everyone\tneg\phi
\\
\Coquorum\phi\equiv&\tneg\Quorum\tneg\phi
&
\Quorum\phi\equiv&\tneg\Coquorum\tneg\phi
\\
\forever\phi\equiv&\tneg\sometime\tneg\phi
&
\sometime\phi\equiv&\tneg\forever\tneg\phi
\\
\tall a.\phi\equiv&\tneg\texi a.\tneg\phi
&
\texi a.\phi\equiv&\tneg\tall a.\tneg\phi
&
\modTF\phi\equiv&\modT\phi\tor\modT\tneg\phi
\\
\modB\phi\equiv&\tneg\modTF\phi
&
\modT\phi\equiv&\tneg\modFB\phi
&
\modF\phi\equiv&\tneg\modTB\phi
\end{array}
$$
\caption{Some easy equivalences}
\label{fig.easy.equivalences}
\end{figure}


\subsection{Two commutation lemmas}
\label{subsect.two.commutation.results}

We note some simple commutation lemmas.
These results are all easy facts of our model, but it only takes a few lines to spell them out:

\begin{lemm}[Commuting connectives]
\label{lemm.commuting.connectives}
Suppose $\phi$ and $\phi'$ are closed predicates and suppose
$$
\bullet\in\{\tor,\tand\} 
\ \ \text{and}\ \ 
\tf M\in\{\modT,\modTB\} . 
$$
Then
$$
\ment\tf M(\phi\bullet\phi')\equiv (\tf M\phi\bullet\tf M\phi') .
$$
\end{lemm}
\begin{proof}
Routine from Lemma~\ref{lemm.equiv}(\ref{item.equiv.2}) and Figures~\ref{fig.3.phi.f} and~\ref{fig.3.derived}, and facts of $\THREE$ as a lattice.
\end{proof}


\begin{lemm}[Commuting modalities]
\label{lemm.someone.commutation}
Suppose $\phi$ is a closed predicate and suppose 
$$
\mathsf Q\in\{\texi,\tall,\someone,\everyone,\Quorum,\Coquorum,\recent,\urecent,\yesterday,\tomorrow,\forever,\sometime,\infinitely,\final\}.
$$
Then 
$$
\ment\mathsf Q\modT\phi \equiv \modT\mathsf Q\phi
\quad\text{and}\quad
\ment\mathsf Q\modTB\phi \equiv \modTB\mathsf Q\phi.
$$
\end{lemm}
\begin{proof}
Routine from Lemma~\ref{lemm.equiv}(\ref{item.equiv.2}) and Figures~\ref{fig.3.phi.f} and~\ref{fig.3.derived}, and facts of $\THREE$ as a lattice.
\end{proof}


\subsection{Lattice inequalities}

Corollary~\ref{corr.recent.everyone.commute} is in the same spirit as Lemmas~\ref{lemm.commuting.connectives} and~\ref{lemm.someone.commutation}. 
In practice we usually just use it a direct fact of the model, but we spell it out anyway.
It relies on this simple fact about lattices: 
\begin{lemm} 
\label{lemm.quantifier.swap}
Suppose $(\mathcal L,\tand,\tor)$ is a complete lattice and $X$ and $Y$ are sets and $f:(X\times Y)\to\mathcal L$ is any function. 
Then:
\begin{enumerate*}
\item\label{item.quantifier.swap.1}
$\bigvee_{x}\bigwedge_{y}f(x,y)\leq \bigwedge_{y}\bigvee_{x}f(x,y)$, 
\item\label{item.quantifier.swap.1b}
It is not in general the case that 
$\bigwedge_{y}\bigvee_{x}f(x,y) \leq \bigvee_{x}\bigwedge_{y}f(x,y)$.
\item\label{item.quantifier.swap.2}
$\bigvee_{x}\bigvee_{y}f(x,y) = \bigvee_{y}\bigvee_{x}f(x,y)$.
\end{enumerate*}
\end{lemm}
\begin{proof}
By easy lattice calculations. 
\end{proof}


\begin{corr}
\label{corr.recent.everyone.commute}
Suppose $\mathsf Q\phi$ is a closed predicate, where 
$$
\mathsf Q\in\{\texi a,\tall a \mid a{\in}\tf{VarSymb}\}\cup\{\someone,\everyone,\Quorum,\Coquorum\},
$$
and suppose $\acontext$ is a context.
Then:
\begin{enumerate*}
\item\label{item.recent.everyone.commute.1}
$\modellabel{\recent\mathsf Q\phi}{\acontext}\leq \modellabel{\mathsf Q\recent\phi}{\acontext}$. 
\item
$\modellabel{\mathsf Q\recent\phi}{\acontext}\leq \modellabel{\recent\mathsf Q\phi}{\acontext}$ for 
$\mathsf Q\in\{\texi a\mid a{\in}\tf{VarSymb}\}\cup\{\someone,\Quorum\}$.
\item
It need not be the case that $\modellabel{\mathsf Q\recent\phi}{\acontext}\leq \modellabel{\recent\mathsf Q}{\acontext}$ for 
$\mathsf Q\in\{\tall a\mid a{\in}\tf{VarSymb}\}\cup\{\everyone,\Coquorum\}$.
\end{enumerate*}
\end{corr}
\begin{proof}
Suppose $n\in\Time$ and $p\in\Pnt$ and $O\in\opensne$ and $\avaluation$ is a valuation.
We freely use \strongmodusponens.
\begin{enumerate}
\item
We consider $\mathsf Q=\Coquorum$.
We use the clause for $\recent$ in Figure~\ref{fig.3.derived}, the clause for $\Coquorum$ in Figure~\ref{fig.3.phi.f}, and then Lemma~\ref{lemm.quantifier.swap}(\ref{item.quantifier.swap.1}): 
$$
\bigvee_{0{\leq}n'{<}n}\bigwedge_{O'{\in}\opensne}\modellabel{\phi}{n',p,O',\avaluation}
\stackrel{L\ref{lemm.quantifier.swap}(\ref{item.quantifier.swap.1})}{\leq}
\bigwedge_{O'{\in}\opensne}\bigvee_{0{\leq}n'{<}n}\modellabel{\phi}{n',p,O',\avaluation} .
$$
Now we consider $\mathsf Q=\tall a$:
$$
\bigvee_{0{\leq}n'{<}n}\bigwedge_{v{\in}\Val}\modellabel{\phi[a\ssm v]}{n',p,O,\avaluation}
\stackrel{L\ref{lemm.quantifier.swap}(\ref{item.quantifier.swap.1})}{\leq}
\bigwedge_{v{\in}\Val}\bigvee_{0{\leq}n'{<}n}\modellabel{\phi[a\ssm v]}{n',p,O,\avaluation} .
$$
The other cases are precisely similar.
\item
We consider just $\mathsf Q=\someone$; the other cases are precisely similar.
We just use Lemma~\ref{lemm.quantifier.swap}(\ref{item.quantifier.swap.2}): 
$$
\bigvee_{p'{\in}\Pnt}\bigvee_{0{\leq}n'{<}n}\modellabel{\phi}{n',p',O,\avaluation}
\stackrel{L\ref{lemm.quantifier.swap}(\ref{item.quantifier.swap.2})}{=}
\bigvee_{0{\leq}n'{<}n}\bigvee_{p'{\in}\Pnt}\modellabel{\phi}{n',p',O,\avaluation} .
$$
\item
We consider just $\mathsf Q=\tall$; the other cases are precisely similar.
We note by Lemma~\ref{lemm.quantifier.swap}(\ref{item.quantifier.swap.1b}) that 
$$
\bigwedge_{v{\in}\Val}\bigvee_{0{\leq}n'_v{<}n}\modellabel{\phi[a\ssm v]}{n'_v,p,O,\avaluation}
\leq
\bigvee_{0{\leq}n'{<}n}\bigwedge_{v\in\Val}\modellabel{\phi[a\ssm v]}{n',p,O,\avaluation} 
\ \ \ \text{does not hold in general}.
$$
\qedhere\end{enumerate}
\end{proof}





\subsection{The compound modalities $\someoneAll$, $\everyoneAll$, $\QuorumBox$, and $\CoquorumDiamond$}
\label{subsect.compound.modalities}

\begin{rmrk}
\label{rmrk.compound.modalities}
We will make much use of the modalities $\someoneAll$, $\everyoneAll$, $\QuorumBox$, and $\CoquorumDiamond$.

We call these four combinations \deffont{compound modalities}, because they are built out of $\Quorum$ and $\Coquorum$ (which quantify over open sets) and $\someone$ and $\everyone$ (which quantify over points in an open set). 
In this Subsection we study some of their properties.
\end{rmrk}



Lemma~\ref{lemm.dense.char} just unpacks definitions:
\begin{lemm}
\label{lemm.dense.char}
Suppose $\phi$ is a closed predicate and $n\in\Time$ and $\avaluation$ is a valuation.
Then: 
\begin{enumerate*}
\item\label{item.dense.char.someoneAll}
$n\mentval\someoneAll\phi$ if and only if $\Exists{O\in\opensne}O\between\{p\in\Pnt \mid n,p,O\mentval\phi\}$.
\item\label{item.dense.char.everyoneAll}
$n\mentval\everyoneAll\phi$ if and only if $\Forall{O\in\opensne}O\subseteq\{p\in\Pnt \mid n,p,O\mentval\phi\}$.
\item\label{item.dense.char.QuorumBox}
$n\mentval\QuorumBox\phi$ if and only if $\Exists{O\in\opensne}O\subseteq\{p\in\Pnt \mid n,p,O\mentval\phi\}$.
\item\label{item.dense.char.CoquorumDiamond}
$n\mentval\CoquorumDiamond\phi$ if and only if $\Forall{O\in\opensne}O\between\{p\in\Pnt \mid n,p,O\mentval\phi\}$.
\end{enumerate*} 
\end{lemm}
\begin{proof}
We consider each part in turn:
\begin{enumerate}
\item
We reason as follows: 
$$
\begin{array}{r@{\ }l@{\quad}l}
n\mentval\someoneAll\phi
\liff&
\Exists{O{\in}\opensne}\Exists{p{\in}O}n,p,O\mentval\phi
&\text{Figure~\ref{fig.3.phi.f}}
\\
\liff&
\Exists{O{\in}\opensne}O\between \{p\in\Pnt \mid n,p,O\mentval\phi\}
&\text{Fact}
\end{array}
$$
\item
We reason as follows: 
$$
\begin{array}{r@{\ }l@{\quad}l}
n\mentval\everyoneAll\phi
\liff&
\Forall{O{\in}\opensne}\Forall{p{\in}O}n,p,O\mentval\phi
&\text{Figure~\ref{fig.3.phi.f}}
\\
\liff&
\Forall{O{\in}\opensne}O\subseteq\{p\in\Pnt \mid n,p,O\mentval\phi\}
&\text{Fact}
\end{array}
$$
\item
We reason as follows:
$$
\begin{array}{r@{\ }l@{\quad}l}
n\mentval\QuorumBox\phi
\liff&
\Exists{O{\in}\opensne}\Forall{p{\in}O}n,p,O\mentval\phi
&\text{Figure~\ref{fig.3.phi.f}}
\\
\liff&
\Exists{O{\in}\opensne}O\subseteq\{p\in\Pnt \mid n,p,O\mentval\phi\}
&\text{Fact}
\end{array}
$$
\item
We reason as follows:
$$
\begin{array}{r@{\ }l@{\quad}l}
n\mentval\CoquorumDiamond\phi
\liff&
\Forall{O{\in}\opensne}\Exists{p{\in}O}n,p,O\mentval\phi
&\text{Figure~\ref{fig.3.phi.f}}
\\
\liff&
\Forall{O{\in}\opensne}O\between\{p\in\Pnt \mid n,p,O\mentval\phi\}
&\text{Fact}
\end{array}
$$
\qedhere\end{enumerate}
\end{proof}



Proposition~\ref{prop.someone.implies.someoneAll} packages up some simple corollaries of Lemma~\ref{lemm.dense.char}, in forms that will be useful for our later proofs: 
\begin{prop}
\label{prop.someone.implies.someoneAll}
Suppose $\phi$ is a closed predicate and $\acontext$ is a context.
Then:
\begin{enumerate*}
\item\label{item.someone.implies.someoneAll.4}
$\acontext\ment\QuorumBox\phi\tand\CoquorumDiamond\psi\limp\acontext\ment\someoneAll(\phi\tand\psi)$.
\item\label{item.someone.implies.someoneAll.4b}
$\acontext\ment\QuorumBox\phi\tand\everyoneAll\psi\limp\acontext\ment\QuorumBox(\phi\tand\psi)$.
\item\label{item.someone.implies.someoneAll.5}
$\acontext\ment\modT(\QuorumBox\phi\tand\CoquorumDiamond\psi)\limp\acontext\ment\modT\someoneAll(\phi\tand\psi)$,\ and
\\
$\acontext\ment(\QuorumBox\phi\tand\CoquorumDiamond\psi)\timpc \someoneAll(\phi\tand\psi)$.
\item\label{item.someone.implies.someoneAll.2}
$\acontext\ment\modT\QuorumBox\phi\limp\acontext\ment\modT\someoneAll\phi$.
\item\label{item.someone.implies.someoneAll.2b}
$\acontext\ment\modT\CoquorumDiamond\phi\limp\acontext\ment\modT\someoneAll\phi$.
\end{enumerate*}
\end{prop}
\begin{proof}
We consider each part in turn:
\begin{enumerate}
\item
Routine from Lemma~\ref{lemm.dense.char}(\ref{item.dense.char.QuorumBox}\&\ref{item.dense.char.CoquorumDiamond}).
\item
Routine from Lemma~\ref{lemm.dense.char}(\ref{item.dense.char.QuorumBox}\&\ref{item.dense.char.everyoneAll}).
\item
From part~\ref{item.someone.implies.someoneAll.4} of this result taking $\modT\phi$ and $\modT\psi$ and the commutation Lemmas~\ref{lemm.someone.commutation} and~\ref{lemm.commuting.connectives}.
Then we use \strongmodusponens.
\item
From part~\ref{item.someone.implies.someoneAll.5} of this result, taking $\psi=\tvT$.
\item
From part~\ref{item.someone.implies.someoneAll.5} of this result, taking $\phi=\tvT$.
\qedhere\end{enumerate}
\end{proof}



 







\subsection{Some comments on time}

\begin{rmrk}
This Subsection notes Lemmas~\ref{lemm.tomorrow.forever.recent} and~\ref{lemm.persist.char}.
They both have to do with time, and underneath all their complicated-looking symbols they are both straightforward:
\begin{itemize*}
\item
Lemma~\ref{lemm.tomorrow.forever.recent} reflects a fact that if something holds now, then in the future it will have held.
\item
Lemma~\ref{lemm.persist.char} reflects the fact that something repeats forever if and only if it happens infinitely often.
\end{itemize*}
\end{rmrk}


\begin{lemm}
\label{lemm.tomorrow.forever.recent}
Suppose $\acontext=(n,p,O,\avaluation)$ is a context and $\phi$ is a closed predicate.
Then:
\begin{enumerate*}
\item\label{item.tomorrow.forever.recent.1}
$n,p,O\mentval\modT\phi$ implies $\Forall{n'{\gneq}n}\,n',p,O\mentval\modT\recent\phi$.
\item\label{item.rur.r}
$n,p,O\mentval\modT\urecent\phi$ implies $\Forall{n'{\geq} n}\,n',p,O\mentval\modT\urecent\phi$.
\item
As corollaries, 
$$
\acontext\ment\phi\timpc \forever\recent\phi
\quad\text{and}\quad
\acontext\ment\urecent\phi\timpc\forever\urecent\phi.
$$
\end{enumerate*}
\end{lemm}
\begin{proof}
This is just from the clauses for $\recent$ and $\urecent$ in Figure~\ref{fig.3.derived}. 
The corollary follows using \strongmodusponens.
\end{proof}



$\infinitely=\forever\sometime$ has a characterisation as `infinitely often', with the convenient feature that its denotation is independent of the stage at which it is evaluated (called $n'$ in the statement of the Lemma below):
\begin{lemm}
\label{lemm.persist.char}
Suppose $\phi$ is a closed predicate and $p\in\Pnt$ and $O\in\opensne$ and $\acontext$ is a context.
Suppose further that $n'\in\Time$ and $\f{tvs}\in\{\threeTrue,\threeValid\}$.
Then we have
$$
\modellabel{\infinitely\phi}{n',p,O,\avaluation}\in\f{tvs}
\quad\liff\quad
\{n\in\Time \mid \modellabel{\phi}{n,p,O,\avaluation}\in\f{tvs}\}\text{ is infinite} . 
$$
\end{lemm}
\begin{proof}
By routine arguments from the clause for $\infinitely$ in Figure~\ref{fig.3.derived}.
\end{proof}






\subsection{Unique existence and affine existence in three-valued logic}
\label{subsect.unique.affine.existence}

\subsubsection{The definition, and main result}

\begin{rmrk}
Distributed algorithms often require an implementation to make (at most) one choice from (zero or more) possibilities.
In our axioms, we model this using unique existence and affine existence:
\begin{itemize*}
\item
\emph{Unique existence} means `there exists precisely one'; we write this $\exists!$, as standard.
\item
\emph{Affine existence} means `there exist zero or one'; we write this $\exists_{01}$.\footnote{We are not aware of a standard symbol for affine existence.}
\end{itemize*}
See for example \rulefont{PaxPropose01} and \rulefont{PaxWrite01} in Figure~\ref{fig.logical.paxos} (a discussion of which is in Remark~\ref{rmrk.paxos.axioms.discussion}(\ref{item.01.discussion})).

So we need to think about what $\exists!$ and $\exists_{01}$ even mean in a three-valued setting.
This turns out to be not completely straightforward, because we have a third truth-value that is neither true nor false.

For concreteness, suppose $\tf P$ is a unary predicate and $\acontext$ is a context, and suppose that 
$\modellabel{\phi[a\ssm v]}{\acontext}=\tvB$ for all $v\in\Val$.
Then what truth-values should $\modellabel{\texiaffine a.\phi}{\acontext}$ and $\modellabel{\texiunique a.\phi}{\acontext}$ have?
 
Now suppose $v,v',v''\in\Val$, and suppose that $\modellabel{\phi[a\ssm v''']}{\acontext}=\tvF$ for all $v'''\in\Val\setminus\{v,v',v''\}$ and
$$
\modellabel{\phi[a\ssm v]}{\acontext}=\tvT
\quad\land\quad
\modellabel{\phi[a\ssm v']}{\acontext}=\tvB
\quad\land\quad
\modellabel{\phi[a\ssm v'']}{\acontext}=\tvB .
$$
\noindent Then what truth-values should $\modellabel{\texiaffine a.\phi}{\acontext}$ and $\modellabel{\texiunique a.\phi}{\acontext}$ have?

Now suppose that $\modellabel{\phi[a\ssm v''']}{\acontext}=\tvF$ for all $v'''\in\Val\setminus\{v\}$ and
$$
\modellabel{\phi[a\ssm v]}{\acontext}=\tvB .
$$
\noindent Then what truth-values should $\modellabel{\texiaffine a.\phi}{\acontext}$ and $\modellabel{\texiunique a.\phi}{\acontext}$ have?

Our answers to all questions above will be: $\tvB$ for both.
This is not automatic: it is a real design decision that we have to make.
\end{rmrk}

\begin{defn}
\label{defn.unique.exist}
Suppose $\phi$ is a predicate.
Define \deffont{affine existence} $\texiaffine a.\phi$ and \deffont{unique existence} $\texiunique  a.\phi$ by:
$$
\begin{array}{r@{\ }l}
\texiaffine a.\phi =& \tall a,a'.(\phi\tand\phi[a\ssm a']) \tnotor a\teq a'
\\
\texiunique  a.\phi =& (\texi a.\phi) \tand (\texiaffine a.\phi)
\end{array}
$$
\end{defn}


\begin{rmrk}
\label{rmrk.unpack.texiunique}
Suppose $\texi a.\phi$ is a closed predicate (i.e. $a\in\tf{VarSymb}$ and $\phi$ is a predicate with free variable symbols at most $\{a\}$), and suppose $\acontext$ is a context.
We spell unpack the denotations of $\texiunique$ and $\texiaffine$ from Definition~\ref{defn.unique.exist}: 
\begin{enumerate*}
\item
If $\exists v,v'\in\Val.v\neq v'\land\modellabel{\phi[a\ssm v]}{\acontext}=\modellabel{\phi[a\ssm v']}{\acontext}=\tvT$, then 
$$
\modellabel{\texiunique a.\phi}{\acontext}=\tvF=\modellabel{\texiaffine a.\phi}{\acontext}.
$$
\item
If $\exists!v\in\Val.\modellabel{\phi[a\ssm v]}{\acontext}=\tvT$ and $\Forall{v\in\Val}\modellabel{\phi[a\ssm v]}{\acontext}\in\{\tvT,\tvF\}$, then 
$$
\modellabel{\texiunique a.\phi}{\acontext}=\tvT=\modellabel{\texiaffine a.\phi}{\acontext}.
$$
\item
If $\exists!v\in\Val.\modellabel{\phi[a\ssm v]}{\acontext}=\tvT$ and $\Exists{v\in\Val}\modellabel{\phi[a\ssm v]}{\acontext}=\tvB$, then 
$$
\modellabel{\texiunique a.\phi}{\acontext}=\tvB=\modellabel{\texiaffine a.\phi}{\acontext}.
$$
\item\label{item.explain.existsunique.4}
If $\neg\Exists{v\in\Val}\modellabel{\phi[a\ssm v]}{\acontext}=\tvT$ and $\exists! v'\in\Val.\modellabel{\phi[a\ssm v]}{\acontext}=\tvB$, then 
$$
\modellabel{\texiunique a.\phi}{\acontext}=\tvB
= \modellabel{\texiaffine a.\phi}{\acontext} . %
$$
(We discuss this in Remark~\ref{rmrk.discuss.clause.4} below.)
\item\label{item.two.b}
If $\neg\Exists{v\in\Val}\modellabel{\phi[a\ssm v]}{\acontext}=\tvT$ and $\exists v,v'\in\Val.v\neq v'\land \modellabel{\phi[a\ssm v]}{\acontext}=\modellabel{\phi[a\ssm v']}{\acontext}=\tvB$, then 
$$
\modellabel{\texiunique a.\phi}{\acontext}
=\tvB=
\modellabel{\texiaffine a.\phi}{\acontext} .
$$
\item
If $\Forall{v\in\Val}\modellabel{\phi[a\ssm v]}{\acontext}=\tvF$, then
$$
\modellabel{\texiunique a.\phi}{\acontext}=\tvF
\quad\text{and}\quad
\modellabel{\texiaffine a.\phi}{\acontext}=\tvT.
$$
\end{enumerate*}
\end{rmrk}

\begin{rmrk}
\label{rmrk.discuss.texiunique}
We make two elementary observations on Remark~\ref{rmrk.unpack.texiunique}:
\begin{enumerate} 
\item
The equality $\teq$ in our logic is two-valued: $\model{v\teq v'}$ returns $\tvT$ or $\tvF$, and never $\tvB$.\footnote{In the terminology of Definition~\ref{defn.correct}(\ref{item.correct.phi}) below, equality predicates are always \emph{correct}.}
\item
Two notions of `unique existence' are in play in this discussion:
\begin{itemize*}
\item
`$\exists!v\in\Val$' refers to the actual unique existence quantifier in the real world (whatever that is).
This computes actual truth-values `true' and `false' and returns true when there actually exists a unique value.
\item
`$\texiunique a$' is part of the (derived) syntax of \QLogic, and is defined in Definition~\ref{defn.unique.exist}.
Its denotation computes elements of $\THREE=\{\tvF,\tvB,\tvT\}$.  
As is standard (though confusingly) we \emph{also} call elements of $\THREE$ truth-values; but they are truth-values of \QLogic, rather than actual truth-values in the real world (whatever that is).
\end{itemize*}
So: $\exists!$ and $\texiunique$ are both unique existence operators, but they live at different levels and they are not the same thing.\footnote{This kind of thing is far from unusual and appears also outside of papers with fancy logics.  For example in programming, a programming language may have multiple versions of the number $1$ (8-bit, floating point, 64-bit, and infinite precision) all of which are called `one', but all of which are different from one other, and also different from the actual $1\in\mathbb N$ (whatever that is).}
\end{enumerate}
\end{rmrk}


We will use unique/affine existence in our proofs via Lemma~\ref{lemm.unique.affine.existence}: 
\begin{lemm}
\label{lemm.unique.affine.existence}
Suppose $\acontext$ is a context and $\texi a.\phi$ is a closed predicate and $v,v'\in\Val$ are values.
Then:
\begin{enumerate*}
\item\label{item.unique.affine.existence.01}
$\acontext\ment\texiaffine a.\phi$ if and only if 
$\exists_{01} v\in\Val.\modellabel{\phi[a\ssm v]}{\acontext}=\tvT$.
\item\label{item.unique.affine.existence.11}
$\acontext\ment\texiunique a.\phi$ if and only if 
$\exists v\in\Val.\modellabel{\phi[a\ssm v]}{\acontext}\in\threeValid \land \exists_{01} v\in\Val.\modellabel{\phi[a\ssm v]}{\acontext}=\tvT$.
\item\label{item.unique.affine.existence.01.implies}
If $\acontext\ment\texiaffine a.\phi$ and $\acontext\ment\modT\phi[a\ssm v]$ and $\acontext\ment\modT\phi[a\ssm v']$, then $v=v'$.
\item\label{item.unique.affine.existence.11.implies}
If $\acontext\ment\texiunique a.\phi$ and $\acontext\ment\modT\phi[a\ssm v]$ and $\acontext\ment\modT\phi[a\ssm v']$, then $v=v'$.
\end{enumerate*}
\end{lemm}
\begin{proof}
We unpack Definition~\ref{defn.validity.judgement}(\ref{item.ctx.ment.phi}) %
and check the possibilities in Remark~\ref{rmrk.unpack.texiunique}.
\end{proof}


\subsubsection{Comments on an edge case}

Lemma~\ref{lemm.texiunique.all.b} notes a (perhaps surprising, at first glance) fact about our denotation, that 
if a unary predicate symbol $\tf P$ takes truth-value $\tvB$ when applied to all values (as would model a participant that has entirely crashed), then $\texiunique\tf P$ and $\texiaffine\tf P$ are valid --- even though $\tf P(v)$ is valid on every value.
This reflects the common side-condition in the literature `if the participant is correct, then \emph{stuff}', which is trivially true if the participant is \emph{not} correct: 
\begin{lemm}
\label{lemm.texiunique.all.b}
Suppose $\acontext$ is a context and $\tf P$ is a unary predicate symbol:\footnote{There is nothing mathematically special about using a \emph{unary} predicate symbol here, but in practice this is the case that we will use.}
Then
$$
\acontext\ment \tall\modB\tf P \limp (\acontext\ment \texiunique \tf P \land \acontext\ment \texiaffine \tf P).  
$$
The above statement uses Notation~\ref{nttn.hos}.
Unpacking this, we obtain:
$$
\acontext\ment \tall a.\modB\tf P(a) \limp (\acontext\ment \texiunique a.\tf P(a)\land \acontext\ment \texiaffine a.\tf P(a)) . 
$$  
\end{lemm}
\begin{proof}
Suppose $\acontext\ment \tall\modB\tf P$.
Unpacking Definition~\ref{defn.validity.judgement}(\ref{item.ctx.ment.phi}) this means $\modellabel{\tall\modB\tf P}{\acontext}\in\threeValid$, and by Figure~\ref{fig.3.derived} it follows that $\Forall{v\in\Val}\modellabel{\tf P(v)}{\acontext}=\tvB$.

By Remark~\ref{rmrk.unpack.texiunique}(\ref{item.two.b}) it follows that $\modellabel{\texiunique\tf P}{\acontext}=\tvB$.
The truth-value $\tvB$ is not true, but it is \emph{valid} by Definition~\ref{defn.validity.judgement}(\ref{item.ctx.ment.phi}), and thus $\acontext\ment\texiaffine\tf P$ and $\acontext\ment\texiunique\tf P$ as required. 
\end{proof}



\begin{rmrk}
\label{rmrk.discuss.clause.4}
Arguably, we would expect in clause~\ref{item.explain.existsunique.4} above 
(where $\neg\Exists{v\in\Val}\modellabel{\phi[a\ssm v]}{\acontext}=\tvT$ and $\exists! v'\in\Val.\modellabel{\phi[a\ssm v]}{\acontext}=\tvB$)
that $\modellabel{\texiaffine a.\phi}{\acontext}$ should equal $\tvT$, not $\tvB$. 
This is not hard to arrange, e.g. as
$$
\texiaffine a.\phi \tor \texiaffine a.\modT\phi .
$$
Either choice works with our later proofs; we never encounter an edge case where it matters.
\end{rmrk}



\subsection{Pointwise predicates (a well-behavedness property)}

For this subsection we fix a signature $\Sigma$ and a model $\mathcal M=(\Pnt,\opens,\Val)$.

\subsubsection{The definition, and basic properties}

\begin{rmrk}
\emph{Being pointwise} is a well-behavedness property on predicates, that their denotation in a context $\acontext=(n,p,O,\avaluation)$ %
does not depend on the open set parameter $O$ of the context.

Many important syntactic classes of predicates are naturally pointwise, including any predicate of the form $\tf P(v_1,\dots,v_n)$ (Lemma~\ref{lemm.atomic.O}).

Being pointwise is desirable for us because it simplifies denotation to holding `at a point';\footnote{Actually: at a stage, point, and valuation, but intuitively we tend to view the stage and valuation as fixed, and the point $p$ and open set $O$ as varying.} and also because several nice properties depend on this (see for example Proposition~\ref{prop.pointwise.dense.char}); and finally also because there is an easy syntactic condition (given in Lemma~\ref{lemm.atomic.O}(\ref{item.atomic.pointwise.2})) that guarantees that a predicate is pointwise.
\end{rmrk}
 

\begin{defn}[Pointwise predicates]
\label{defn.pointwise.meaning}
Suppose $\phi$ is a closed predicate.
\begin{enumerate*}
\item
If $n\in\Time$ and $\avaluation\in\Valuation(\Sigma,\Val)$ then call $\phi$ \deffont{pointwise in $(n,\avaluation)$} when
$$
\Forall{p\in\Pnt}\Forall{O,O'{\in}\opensne}\modellabel{\phi}{n,p,O,\avaluation}=\modellabel{\phi}{n,p,O',\avaluation} .
$$ 
In words, $\phi$ is pointwise in $(n,\avaluation)$ when at every point $p\in\Pnt$, the open set parameter $O\in\opensne$ does not affect the value of the denotation $\modellabel{\phi}{n,p,O,\avaluation}$. 
\item
Call $\phi$ \deffont{absolutely pointwise} when it is pointwise for every $n\in\Time$ and $\avaluation\in\Valuation(\Sigma,\Val)$. 
\end{enumerate*}
\end{defn}

\begin{defn}
\label{defn.pointwise.pred.nttn}
Suppose $\phi$ is closed, and $n\in\Time$ and $p\in\Pnt$ and $\avaluation$ is a valuation.
Then define $\modellabel{\phi}{n,p,\avaluation}$ by
$$
\modellabel{\phi}{n,p,\avaluation}=\bigwedge_{O\in\opensne}\modellabel{\phi}{n,p,O,\avaluation}.
$$
\end{defn}

In a manner of speaking we have already seen $\modellabel{\phi}{n,p,\avaluation}$ from Definition~\ref{defn.pointwise.pred.nttn} in our logic:
\begin{lemm}
\label{lemm.Coquorum.denote}
Suppose $n\in\Time$ and $p\in\Pnt$ and $\avaluation$ is a valuation.
Then
$$
\modellabel{\phi}{n,p,\avaluation}=\modellabel{\Coquorum\phi}{n,p,O,\avaluation}.
$$
\end{lemm}
\begin{proof}
Direct from the clause for $\Coquorum$ from Figure~\ref{fig.3.derived}.
\end{proof}

\begin{lemm}
\label{lemm.pointwise.pred}
Suppose $\phi$ is a closed predicate and $n\in\Time$ and $\avaluation$ is a valuation.
Then $\phi$ is pointwise in $(n,\avaluation)$ in the sense of Definition~\ref{defn.pointwise.meaning} 
if and only if
$$
\Forall{p\in\Pnt}\Forall{O\in\opensne}\modellabel{\phi}{n,p,O,\avaluation}=\modellabel{\phi}{n,p,\avaluation}.
$$
\end{lemm}
\begin{proof}
By a routine calculation from Definitions~\ref{defn.pointwise.meaning} and~\ref{defn.pointwise.pred.nttn}.
\end{proof}

\begin{corr}
\label{corr.pointwise.pred}
Suppose $\phi$ is a closed predicate and $n\in\Time$ and $p\in\Pnt$ and $\avaluation$ is a valuation.
Suppose $\phi$ is pointwise in $(n,\avaluation)$.
Then
$$
\Forall{O\in\opensne}n,p,O\mentval\phi\liff n,p\mentval\phi. 
$$
\end{corr}
\begin{proof}
Routine from Definition~\ref{defn.validity.judgement} and Lemma~\ref{lemm.pointwise.pred}. %
\end{proof}


Being pointwise is particularly nice because of two simple syntactic criteria that guarantee it:
\begin{lemm}[Conditions for being pointwise]
\label{lemm.atomic.O}
Suppose $\phi$ is a predicate.
Then:
\begin{enumerate*}
\item\label{item.atomic.pointwise}
If $\phi$ is atomic (Notation~\ref{nttn.logic.terminology}(\ref{item.atomic.predicate}); having the form $\tf P(v_1,\dots,v_n)$ or $\tneg\tf P(v_1,\dots,v_n)$) then it is absolutely pointwise (Definition~\ref{defn.pointwise.meaning}).
\item\label{item.atomic.pointwise.2}
If $\phi$ does not mention $\someone$ or $\everyone$ except within the scope of a $\Quorum$ or $\Coquorum$ modality, then $\phi$ is absolutely pointwise.\footnote{In particular this holds if $\Quorum$ and $\Coquorum$ and only ever mentioned in combination with $\someone$ and $\everyone$, as in the compound modalities $\someoneAll$, $\everyoneAll$, $\QuorumBox$, and $\CoquorumDiamond$.}
\end{enumerate*} 
\end{lemm}
\begin{proof}
For part~1, we just unpack the definition of $\modellabel{\phi}{}$ in Figure~\ref{fig.3.phi.f} and note that the `$\avaluation(\Pnt)_{n,p}$' used in the clause for a predicate symbol depends just on the values of $n$, $p$, and $\avaluation$, but not on the open sets $O$ and $O'$. 

Part~2 follows by a routine induction on predicate syntax, using part~1 as the base case.
\end{proof}


\subsubsection{Pointwise, as a modality in our logic}

In fact, we can express the property of being pointwise within our logic:
\begin{defn}
\label{defn.pointwise.pred}
Suppose $\phi$ is a predicate.
Then define a predicate $\pointwise{\phi}$ by
$$
\pointwise{\phi} \ =\ \everyoneAll(\Quorum\phi \equiv \Coquorum\phi) .
$$
\end{defn}


\begin{lemm}
\label{lemm.pointwise.heart}
Suppose $\phi$ is a closed predicate and $n\in\Time$ and $p\in\Pnt$ and $\avaluation$ is a valuation.
Then 
$$
n,p\mentval\Quorum\phi \equiv \Coquorum\phi
\quad\text{if and only if}\quad
\bigvee_{O'\in\opensne}\modellabel{\phi}{n,p,O',\avaluation}=\bigwedge_{O'\in\opensne}\modellabel{\phi}{n,p,O',\avaluation}.
$$
(The predicate $\Quorum\phi \equiv \Coquorum\phi$ itself satisfies the conditions for being pointwise from Lemma~\ref{lemm.atomic.O}(\ref{item.atomic.pointwise.2}), so we can sensibly write $n,p\mentval \Quorum\phi \equiv \Coquorum\phi$ above.)
\end{lemm}
\begin{proof}
From Lemma~\ref{lemm.equiv}(\ref{item.equiv.1}) and the clauses for $\Quorum$ in Figure~\ref{fig.3.phi.f}, and $\Coquorum$ from Figure~\ref{fig.3.derived} (or Lemma~\ref{lemm.Coquorum.denote}). 
\end{proof}

\begin{lemm}
Suppose $\phi$ is a closed predicate and $n\in\Time$ and $p\in\Pnt$ and $\avaluation$ is a valuation.
Then $\phi$ is pointwise in $(n,\avaluation)$ if and only if $n,p\mentval\pointwise{\phi}$. 
\end{lemm}
\begin{proof}
By routine calculations using Lemma~\ref{lemm.pointwise.heart}, and the clause for $\everyoneAll$ in Figure~\ref{fig.3.derived}, noting the elementary fact that for every $O\in\opensne$,
$$
\bigwedge_{O'\in\opensne}\modellabel{\phi}{n,p,O',\avaluation}
\leq
\modellabel{\phi}{n,p,O,\avaluation}
\leq
\bigvee_{O'\in\opensne}\modellabel{\phi}{n,p,O',\avaluation} .
$$
\end{proof}



\subsubsection{An interlude: dense sets, and sets with a nonempty open interior}


We recall a pair of standard topological definitions which will turn up repeatedly later, most notably in the context of pointwise predicates as Proposition~\ref{prop.pointwise.dense.char} immediately below, and again later in Lemma~\ref{lemm.is}:
\begin{defn}[Dense subset]
\label{defn.dense.subset}
Suppose $(\Pnt,\opens)$ is a semitopology and suppose $P\subseteq\Pnt$ is a (not necessarily open) set of points.
\begin{enumerate*}
\item\label{item.dense.subset}
Say that $P$ is \deffont{dense} when $\Forall{O{\in}\opensne}O\intersectswith P$.
\item\label{item.noi.subset}
Say that $P$ \deffont{has a nonempty open interior} when $\Exists{O{\in}\opensne}O\subseteq P$.
\item\label{item.dense.noi}
Write $\dense(P)$ when $P$ is dense, and $\noi(P)$ when $P$ has a nonempty open interior. 
\end{enumerate*} 
\end{defn}



\begin{lemm}
\label{lemm.dense.up-closed}
Suppose $(\Pnt,\opens)$ is a semitopology and $P,P'\subseteq\Pnt$.
Then:
\begin{enumerate*}
\item\label{item.dense.up-closed} 
If $P\subseteq P'$ then $\dense(P)\limp \dense(P')$ (being dense is \emph{up-closed}).
\item\label{item.noi.up-closed} 
If $P\subseteq P'$ then $\noi(P)\limp \noi(P')$ (having a non-empty open intersection is \emph{up-closed}).
\item\label{item.noi.noi} 
If $P\in\opensne$ then $\noi(P)$ (a nonempty open set has a nonempty open interior).
\end{enumerate*}
\end{lemm}
\begin{proof}
By routine sets calculations from Definition~\ref{defn.dense.subset}.
\end{proof}

\begin{lemm}
Suppose $(\Pnt,\opens)$ is a semitopology and $P\subseteq\Pnt$.
Then
$$
\dense(P) \liff \neg\noi(\Pnt\setminus P) 
\quad\text{and}\quad 
\noi(P) \liff \neg\dense(\Pnt\setminus P) . 
$$
In this sense, $\noi$ and $\dense$ form a \emph{de Morgan dual pair} of operators --- just as $\forall=\neg\exists\neg$ (as noted in Figure~\ref{fig.easy.equivalences}) and $\exists=\tneg\forall\tneg$.
\end{lemm}
\begin{proof}
We sketch the (easy) argument:
\begin{itemize}
\item
Suppose $P$ is not dense; by Definition~\ref{defn.dense.subset} this means that there exists $O\in\opensne$ such that $P\cap O=\varnothing$, thus $O\subseteq\Pnt\setminus P$, thus $\Pnt\setminus P$ has a nonempty open interior.
\item
Suppose $\Pnt\setminus P$ has a nonempty open interior $\varnothing\neq O\subseteq\Pnt\setminus P$.
Then clearly $P$ is not dense.
\qedhere\end{itemize}
\end{proof}

More on how `being dense' works in semitopologies is in~\cite[Chapter~11]{gabbay:semdca}.
The notion behaves somewhat differently than it does in topologies, but not in ways that will affect us here.

\subsubsection{Application to pointwise predicates}


Technically, Proposition~\ref{prop.pointwise.dense.char} is a corollary of Lemma~\ref{lemm.dense.char} and Corollary~\ref{corr.pointwise.pred}.
The result links the four logical modalities $\someoneAll$, $\everyoneAll$, $\QuorumBox$, and $\CoquorumDiamond$ to elementary topological properties of the sets of points at which these modalities are valid --- being nonempty, being full, having a nonempty open interior, and being dense.
We will use Proposition~\ref{prop.pointwise.dense.char}, both directly and via other lemmas that depend on it.

The most important use-case of Proposition~\ref{prop.pointwise.dense.char} will be when $\phi$ is atomic (Notation~\ref{nttn.logic.terminology}(\ref{item.atomic.predicate}) --- having the form 
$$
\tf P(v_1,\dots,v_n)
\quad\text{or}\quad
\tneg\tf P(v_1,\dots,v_n)
$$
--- since then by Lemma~\ref{lemm.atomic.O} $\phi$ is automatically pointwise:
\begin{prop}
\label{prop.pointwise.dense.char}
Suppose $\phi$ is a closed pointwise (Definition~\ref{defn.pointwise.meaning}) predicate and $n\in\Time$ and $\avaluation$ is a valuation.
Then: 
\begin{enumerate*}
\item\label{item.dense.char.pointwise.someoneAll}
$n\mentval\someoneAll\phi$ if and only if $\varnothing\neq\{p\in\Pnt \mid n,p\mentval\phi\}$. 
\item\label{item.dense.char.pointwise.everyoneAll}
$n\mentval\everyoneAll\phi$ if and only if $\Pnt=\{p\in\Pnt \mid n,p\mentval\phi\}$. 
\item\label{item.dense.char.pointwise.QuorumBox}
$n\mentval\QuorumBox\phi$ if and only if $\noi(\{p\in\Pnt \mid n,p\mentval\phi\})$.
\item\label{item.dense.char.pointwise.CoquorumDiamond}
$n\mentval\CoquorumDiamond\phi$ if and only if $\dense(\{p\in\Pnt \mid n,p\mentval\phi\})$.
\end{enumerate*} 
Recall that:
\begin{itemize*}
\item
$\dense$ and $\noi$ are from Definition~\ref{defn.dense.subset}.
\item
The notations $n,p\mentval\phi$ and $n\mentval\phi$ are from Definition~\ref{defn.validity.judgement}.
\end{itemize*}
\end{prop}
\begin{proof}
If $p$ is pointwise then by Corollary~\ref{corr.pointwise.pred} $n,p,O\mentval\phi$ if and only if $n,p\mentval\phi$, for any $O\in\opensne$, and we just simplify parts~\ref{item.dense.char.someoneAll} to~\ref{item.dense.char.CoquorumDiamond} of Lemma~\ref{lemm.dense.char}:
\begin{enumerate} 
\item
We reason as follows:
$$
\begin{array}{r@{\ }l@{\quad}l}
n\mentval\someoneAll\phi
\liff&
\Exists{O\in\opensne}O\between\{p\in\Pnt \mid n,p\mentval\phi\}
&\text{Lemma~\ref{lemm.dense.char}(\ref{item.dense.char.someoneAll})} 
\\
\liff&
\Pnt\between \{p\in\Pnt \mid n,p\mentval\phi\}
&\text{Fact}
\\
\liff&
\{p\in\Pnt \mid n,p\mentval\phi\}\neq\varnothing
&\text{Fact}
\end{array}
$$
\item
We reason as follows:
$$
\begin{array}{r@{\ }l@{\quad}l}
n\mentval\everyoneAll\phi
\liff&
\Forall{O\in\opensne}O\subseteq\{p\in\Pnt \mid n,p\mentval\phi\}
&\text{Lemma~\ref{lemm.dense.char}(\ref{item.dense.char.everyoneAll})} 
\\
\liff&
\Pnt\subseteq\{p\in\Pnt \mid n,p\mentval\phi\}
&\text{Fact}
\\
\liff&
\{p\in\Pnt \mid n,p\mentval\phi\}=\Pnt
&\text{Fact}
\end{array}
$$
\item
We reason as follows:
$$
\begin{array}{r@{\ }l@{\quad}l}
n\mentval\QuorumBox\phi
\liff&
\Exists{O\in\opensne}O\subseteq\{p\in\Pnt \mid n,p\mentval\phi\}
&\text{Lemma~\ref{lemm.dense.char}(\ref{item.dense.char.QuorumBox})} 
\\
\liff&
\noi(\{p\in\Pnt \mid n,p\mentval\phi\})
&\text{Definition~\ref{defn.dense.subset}}
\end{array}
$$
\item
We reason as follows:
$$
\begin{array}{r@{\ }l@{\quad}l}
n\mentval\CoquorumDiamond\phi
\liff&
\Forall{O\in\opensne}O\between\{p\in\Pnt \mid n,p\mentval\phi\}
&\text{Lemma~\ref{lemm.dense.char}(\ref{item.dense.char.CoquorumDiamond})} 
\\
\liff&
\dense(\{p\in\Pnt \mid n,p\mentval\phi\})
&\text{Definition~\ref{defn.dense.subset}}
\end{array}
$$
\qedhere\end{enumerate}
\end{proof}
See also Lemma~\ref{lemm.allbut.concretely}, which instantiates Proposition~\ref{prop.pointwise.dense.char}(\ref{item.dense.char.pointwise.QuorumBox}\&\ref{item.dense.char.pointwise.CoquorumDiamond}) to a special case
familiar from the literature on distributed systems.


\subsection{Correct predicates (another well-behavedness property)}

\begin{rmrk}
We defined $\threeCorrect=\{\tvT,\tvF\}$ in Definition~\ref{defn.tv.ment}(\ref{item.tv.ment.correct}) and a corresponding modality $\modTF$ in Figure~\ref{fig.3} (which returns $\tvT$ when its argument is in $\threeCorrect$ and $\tvF$ otherwise), and we studied some of their properties, most notably in Lemma~\ref{lemm.TF.correct}.
 
We now extend this to notions of \emph{correctness} on predicates $\phi$ and predicate symbols $\tf P$.

Correctness turns out to be highly relevant to how we will use our logic.
Partly this is because when a predicate is correct, we are returned to the familiar two-valued world of Boolean logic --- but more subtly, correctness is interesting because of how this concept interacts with the predicate and modal parts of our logic; in certain circumstances we can use a combination of correctness \emph{somewhere} and validity \emph{everywhere} to deduce information about truth (see Remark~\ref{rmrk.why.exciting}) and this turns out to be a good logical model for how reasoning on distributed algorithms is often done in practice.

We start with the definition and some basic lemmas, and develop from there:
\end{rmrk}


\begin{defn}[Correctness]
\label{defn.correct}
Suppose $\phi$ is a predicate and $a\in\tf{VarSymb}$ is a variable symbol. 
Then:
\begin{enumerate*}
\item\label{item.correct.phi}
Consistent with the definition of $\threeCorrect=\{\tvT,\tvF\}$ from Definition~\ref{defn.tv.ment}(\ref{item.tv.ment.correct}),
we define $\correct{\phi}$, read \deffont{$\phi$ is correct}, by
$$
\correct{\phi}=\modTF\phi .
$$
Note that $\phi$ being correct does not mean that $\phi$ returns $\tvT$; it means that $\phi$ returns $\tvT$ or $\tvF$ and in particular it does \emph{not} return $\tvB$.
\item\label{item.correct}\label{item.crashed}
Suppose $\tf P\in\PredSymb$ is a unary predicate symbol.
Define a predicate $\correct{\tf P}$ by 
$$
\correct{\tf P}\ =\ \tall a.\correct{\tf P(a)}\ \stackrel{N\ref{nttn.hos}}{=}\ \tall\modTF\tf P
 .
$$
\item
If $\tf P_1,\dots,\tf P_n$ are unary predicate symbols then we may write 
$$
\correct{\tf P_1,\dots,\tf P_n}
\quad\text{for}\quad 
\correct{\tf P_1}\tand\dots\tand\correct{\tf P_n}.
$$
(See for example axiom \rulefont{LdrCorrect} in Figure~\ref{fig.logical.paxos}.) 
\item\label{item.P-correct}
Suppose $n\in\Time$ and $p\in\Pnt$ and $\avaluation$ is a valuation, and suppose $\tf P$ is a predicate symbol.
Then when $n,p\mentval\correct{\tf P}$ we may call $p$ \deffont{$\tf P$-correct}. 
\end{enumerate*}
\end{defn}

\begin{lemm}
\label{lemm.correct.valid}
Suppose $\acontext$ is a context and $\phi$ is a closed predicate.
Then 
$$
\acontext\ment\phi
\quad\text{if and only if}\quad
\modellabel{\phi}{\acontext}\in\threeCorrect.
$$
\end{lemm}
\begin{proof}
This just re-states Lemma~\ref{lemm.tv.ment.TF}(\ref{item.tv.ment.TF.2}).
\end{proof}


\begin{rmrk}
\label{rmrk.tf.synonyms}
Definition~\ref{defn.correct}(\ref{item.correct.phi}) gives us two synonyms for the modality expressing that the truth-value returned by a predicate $\phi$ is either $\tvT$ or $\tvF$ (not $\tvB$): $\modTF$ and $\tf{correct}$.

Although these are synonymous, there is a difference in emphasis:
\begin{itemize*}
\item
$\modTF$ is a logical modality and just means `true or false'.
\item
The literature on distributed systems uses a term `correctness', which indicates that a process succeeds and returns a correct value. 
In our logical model this corresponds to returning $\tvT$ or $\tvF$ --- and thus not returning $\tvB$. 
So when we write $\tf{correct}$, the emphasis is on `not $\tvB$' with emphasis on a process `not failing'.
\end{itemize*}
We use whichever notation seems most natural in context. 
For example, in Figure~\ref{fig.logical.paxos}:
\begin{enumerate*}
\item
In \rulefont{LdrCorrect} we write $\correct{\tf{propose},\tf{write},\tf{decide}}$ because intuitively we are asserting that at a particular time the leader \emph{does not crash} when proposing, writing, or deciding.
\item
In \rulefont{LdrExist} we write $\modTF\tf{leader}$, because $\tf{leader}$ is a predicate that identifies whether a participant is the leader at a particular time, and this is either true or false.\footnote{In practice determining who the leader is require a participant to compute a function, but the computation is not distributed, in the sense that prior agreement amongst participants is assumed on what the function is and what seed data it uses, and each participant computes it locally to get (the same) definitive yes/no answer.}
\end{enumerate*} 
If the reader prefers, they can safely read $\modTF$ every time they see $\tf{correct}$, and they will be precisely correct (pun intended). 
\end{rmrk}



Many predicates of interest are correct, including correctness itself, our truth-modalities, and equality and equivalence:
\begin{lemm}
\label{lemm.equality.correct}
Suppose $\acontext$ is a context and $\phi$ and $\phi'$ are closed predicates and $v,v'\in\Val$.
Then:
\begin{enumerate*}
\item
$\modellabel{\correct{\correct{\phi}}}{\acontext}=\tvT$.
\quad We can (perhaps cheekily) read this as `correctness is correct'.
\item
$\modellabel{\correct{\tf Q\phi}}{\acontext}=\tvT$ for $\tf Q\in\{\modT,\modF,\modB,\modTB,\modTF,\modFB\}$. 
\item\label{item.equality.correct}
$\modellabel{\correct{v\teq v'}}{\acontext}=\tvT$.
In words: equality is correct.
\item\label{item.equivalence.correct}
$\modellabel{\correct{\phi\equiv \phi'}}{\acontext}=\tvT$.
In words: equivalence is correct.
\end{enumerate*}
\end{lemm}
\begin{proof}
By routine checking from Figures~\ref{fig.3}, \ref{fig.3.phi.f}, \ref{fig.3.derived}, and (for $\equiv$)~\ref{fig.3.iff}.
\end{proof}

Lemma~\ref{lemm.correct.P.Pv} and Corollary~\ref{corr.modTF.texi} are straightforward, but we spell it out:
\begin{lemm}
\label{lemm.correct.P.Pv}
Suppose $\tf P\in\PredSymb$ is a unary predicate symbol and $v\in\Val$ and $\acontext$ is a context.
Then 
$$
\modellabel{\correct{\tf P}}{\acontext} = \bigwedge_{v{\in}\Val}\modellabel{\correct{\tf P(v)}}{\acontext}
\quad\text{and}\quad 
\modellabel{\correct{\tf P}}{\acontext}\leq \modellabel{\correct{\tf P(v)}}{\acontext} .
$$
\end{lemm}
\begin{proof}
By Definition~\ref{defn.correct}(\ref{item.correct}) and Figure~\ref{fig.3.derived} 
$$
\modellabel{\correct{\tf P}}{\acontext} 
= 
\modellabel{\tall a.\correct{\tf P(a)}}{\acontext} 
= 
\bigwedge_{v{\in}\Val} \modellabel{\correct{\tf P(v)}}{\acontext} 
\leq 
\modellabel{\correct{\tf P(v)}}{\acontext} .
\qedhere$$
\end{proof}


\begin{corr}
\label{corr.modTF.texi}
Suppose $\acontext$ is a context and $\tf P\in\PredSymb$ is a unary predicate symbol and $a\in\tf{VarSymb}$ is a variable symbol. 
Then 
$$
\acontext\ment\correct{\tf P}
\quad\text{implies}\quad
\acontext\ment\correct{\texi\tf P}\ \land\ \acontext\ment\correct{\tall\tf P}.
$$
\end{corr}
\begin{proof}
Suppose $\acontext\ment\correct{\tf P}$.
By Lemma~\ref{lemm.correct.P.Pv} we see that $\Forall{v\in\Val}\acontext\ment\correct{\tf P(v)}$.
By Lemma~\ref{lemm.correct.valid} $\Forall{v\in\Val}\modellabel{\tf P(v)}{\acontext}\in\threeCorrect$.
Now the clauses for $\texi$ in Figure~\ref{fig.3.phi.f} and $\tall$ in Figure~\ref{fig.3.derived} take a supremum and infimum respectively of $\{\modellabel{\tf P(v)}{\acontext} \mid v\in\Val\}\subseteq\threeCorrect=\{\tvT,\tvF\}$, and this supremum or infimum is also clearly in $\threeCorrect=\{\tvT,\tvF\}$.
We just use Lemma~\ref{lemm.correct.valid}. 
\end{proof}


Proposition~\ref{prop.tv.ment.TF.model}(\ref{item.tv.ment.TF.phi.lr}) is a wrapper for Lemma~\ref{lemm.TF.correct}.
This deceptively simple equivalence turns out to be a key way that correctness can be usefully applied:
\begin{prop}
\label{prop.tv.ment.TF.model}
Suppose $\phi$ is a closed predicate and $\acontext$ is a context and $v\in\Val$.
Then:
\begin{enumerate*}
\item\label{item.tv.ment.TF.phi.lr}
$\acontext\ment\correct{\phi}$ if and only if $\acontext\ment\phi\equiv\modT\phi$. 

In words: correct $\phi$ are characterised precisely by validity being equivalent to truth.
\item\label{item.tv.ment.TF.P}
$\acontext\ment\correct{\tf P}$ if and only if $\Forall{v\in\Val}(\acontext\ment\tf P(v)\equiv\modT\tf P(v))$. 
\item\label{item.modTF.texi.P.valid}
$\acontext\ment\correct{\tf P}$ implies $\acontext\ment\texi\tf P\equiv \modT\texi\tf P$ and $\acontext\ment\tall\tf P\equiv\modT\tall\tf P$. 
\end{enumerate*}
\end{prop}
\begin{proof}
We consider each part in turn:
\begin{enumerate}
\item
Suppose $\acontext\ment\correct{\phi}$.
By Lemma~\ref{lemm.TF.correct} $\modT\modellabel{\phi}{\acontext}=\modellabel{\phi}{\acontext}$.
By Lemma~\ref{lemm.equiv}(\ref{item.equiv.1}) and the clause for $\modT$ in Figure~\ref{fig.3.phi.f}, $\acontext\ment \phi\equiv\modT\phi$.
The reverse implication just reverses the same reasoning.
\item
We combine Lemma~\ref{lemm.correct.P.Pv} with part~\ref{item.tv.ment.TF.phi.lr} of this result for $\phi=\tf P(v)$ for every $v\in\Val$.
\item
From Corollary~\ref{corr.modTF.texi} and Proposition~\ref{prop.tv.ment.TF.model}(\ref{item.tv.ment.TF.phi.lr}).
\qedhere\end{enumerate}
\end{proof}




Proposition~\ref{prop.correct.cobox.boxT} explores how correctness interacts with our compound modalities from Subsection~\ref{subsect.compound.modalities}.
\begin{prop}
\label{prop.correct.cobox.boxT}
Suppose $\phi$ is a closed predicate and $\tf P\in\PredSymb$ is a unary predicate symbol. 
Suppose $n\in\Time$ and $\avaluation$ is a valuation.
Then:
\begin{enumerate*}
\item\label{item.correct.cobox.boxT.1}
$n\mentval\QuorumBox\correct{\phi}\tand\everyoneAll\phi$ implies $n\mentval \modT\QuorumBox\phi$.\footnote{Alert: note that we write $\everyoneAll\phi$ in the right-hand conjunct, not $\CoquorumDiamond\phi$.}
\item\label{item.correct.cobox.boxT.2}
$n\mentval\QuorumBox\correct{\tf P}\tand\everyoneAll\texi\tf P$ implies 
$n\mentval\modT\QuorumBox\texi\tf P$.
\end{enumerate*}
\end{prop}
\begin{proof}
We consider each part in turn:
\begin{enumerate}
\item
Suppose $n\mentval\QuorumBox\correct{\phi}\tand\everyoneAll\phi$.
By Proposition~\ref{prop.someone.implies.someoneAll}(\ref{item.someone.implies.someoneAll.4b})
$n\mentval\QuorumBox(\correct{\phi}\tand\phi)$.
By Proposition~\ref{prop.tv.ment.TF.model}(\ref{item.tv.ment.TF.phi.lr}) and Lemma~\ref{lemm.someone.commutation} (for $\QuorumBox$)
we have $n\mentval\modT\QuorumBox\phi$ as required.
\item
We combine part~\ref{item.correct.cobox.boxT.1} of this result with Corollary~\ref{corr.modTF.texi}.
\qedhere\end{enumerate}
\end{proof}


\begin{rmrk}[Why this is interesting]
\label{rmrk.why.exciting}
The statement of Proposition~\ref{prop.correct.cobox.boxT} is quite heavy on the symbols.
For clarity, let us re-read it, but let us replace $\ment\QuorumBox\tf{correct}$ with the phrase `mostly correct', and $\ment\everyoneAll$ with the phrase `everywhere valid', and $\texi$ with the phrase `for some value'.
Then we get:
\begin{enumerate*}
\item
If $\phi$ is \emph{mostly correct} and \emph{everywhere valid}, then it is \emph{mostly true}.
\item
If $\tf P$ is \emph{mostly correct} and \emph{everywhere valid} for some value, then it is \emph{mostly true} for some value.
\end{enumerate*}
So what we see is that we can promote weak statements about validity everywhere, with stronger but more limited statements about correctness, to obtain limited statements about truth.

The reader familiar with how distributed algorithms work may recognise this style of reasoning.
So: here it is beginning to appear in our logic.
\end{rmrk}

\subsection{$\mrup{a}{\phi}{v}$: a most recent $v$ that makes $\phi[a\ssm v]$ true}
\label{subsect.mrup}

\begin{rmrk}
\label{rmrk.explain.mrup}
We see from Figure~\ref{fig.3.derived} and Remark~\ref{rmrk.comments.on.semantics}(\ref{item.mru.program}) that $\mrup{a}{\phi}{v}$ looks back into the past, searching for a most recent time when $\model{\texi a.\phi}=\tvT$, and checks if at that most recent time, $\model{\phi[a\ssm v]}=\tvT$.
$\mrup{a}{\phi}{v}$ returns $\tvT$ if this search succeeds, and $\tvF$ if this fails.
It never returns $\tvB$.

In words: $\mrup{a}{\phi}{v}$ checks whether $v$ is a value that has most recently made $\phi$ true.

We need $\mrup{a}{\phi}{v}$ for our axiomatisation of Paxos in Figure~\ref{fig.logical.paxos}, where 
\begin{itemize*}
\item
in \rulefont{PaxSend?} we write $\mru{\tf{accept}}{v}$\ (short for $(\mrup{a}{\tf{accept}(a))}{v}$, as per Figure~\ref{fig.hos}), and 
\item
in \rulefont{PaxWrite?} we write $\mru{(\someone\tf{accept})}{v}$\ (short for $\mrup{a}{(\someone\tf{accept}(a))}{v}$).
\end{itemize*} 
We read $\mrup{a}{\phi}{v}$ as 
\begin{quote}
\deffont{$v$ is a value that most recently makes $\phi[a\ssm v]$ true} .
\end{quote}
In the common case that $\phi=\tf P(a)$ for some unary predicate symbol $\tf P$, then we might just say
\begin{quote}
\deffont{$v$ is a value that most recently makes $\tf P$ true}.
\end{quote} 
In this Subsection we study the properties of this construct in more detail.
\end{rmrk}



\begin{lemm}
\label{lemm.mru.unambivalent}
Suppose $\texi a.\phi$ is a closed predicate and $v\in\Val$ and $\acontext$ is a context.
Then:
\begin{enumerate*}
\item
$\acontext\ment\correct{\mrup{a}{\phi}{v}}$.
\item
$\acontext\ment\mrup{a}{\phi}{v}\equiv\modT\mrup{a}{\phi}{v}$.
\end{enumerate*}
\end{lemm}
\begin{proof}
We consider each part in turn:
\begin{enumerate}
\item
By Definition~\ref{defn.correct}(\ref{item.correct.phi}) $\correct{\mrup{a}{\phi}{v}}=\modTF\mrup{a}{\phi}{v}$, and this means that $\mrup{a}{\phi}{v}$ always returns $\tvT$ or $\tvF$ (but not $\tvB$).
We just check the relevant clause in Figure~\ref{fig.3.derived} and see that this is so.
\item
We just use part~1 of this result and Proposition~\ref{prop.tv.ment.TF.model}(\ref{item.tv.ment.TF.phi.lr}).
\qedhere\end{enumerate}
\end{proof}


\begin{rmrk}
\label{rmrk.explain.mru}
\leavevmode
\begin{enumerate}
\item\label{item.explain.mru.T}
Lemma~\ref{lemm.mru.unambivalent}(2) %
notes that $\mrup{a}{\phi}{v}$ and $\modT(\mrup{a}{\phi}{v})$ are extensionally equivalent (Definition~\ref{defn.phi.equivalent}).

We may write one or the other, depending on what seems most clear in a proof.
In particular, if we are using \strongmodusponens then we might retain a $\modT$ modality for clarity.
\item 
We try to head off an easy error: note that $\mrup{a}{\phi}{v}$ is not equivalent to $\modT\recent\phi[a\ssm v]$, nor the (by Lemma~\ref{lemm.someone.commutation}) equivalent form $\recent\modT\phi[a\ssm v]$.
\begin{itemize*}
\item
$\modT\recent\phi[a\ssm v]$ returns $\tvT$ when $\phi[a\ssm v]$ has returned $\tvT$ at some point in the past.
\item
$\mrup{a}{\phi}{v}$ (`$v$ is a value that most recently makes $\phi[a\ssm v]$ true') returns $\tvT$ when $\phi[a\ssm v]$ has returned $\tvT$ in \emph{the most recent stage} such that $\texi a.\phi$ returned $\tvT$ (i.e. such that $\phi[a\ssm v']$ returned $\tvT$ for \emph{any} $v'$).
\end{itemize*}
These are not the same thing.
\end{enumerate}
\end{rmrk}

Remark~\ref{rmrk.explain.mru} notes that $\mrup{a}{\phi}{v}$ is different from $\modT\recent\phi[a\ssm v]$.
That said, there is a close link between them, which will turn out to be important for understanding our axiomatisation of Paxos (see Lemma~\ref{lemm.evaporates.simplify} and surrounding discussion): 
\begin{lemm}
\label{lemm.evaporates}
Suppose $\texi a.\phi$ is a closed predicate and $\acontext=(n,p,O,\avaluation)$ is a context. 
Then
$$
\acontext\ment(\texi v.\mrup{a}{\phi}{v})\equiv \modT\recent\texi a.\phi.
$$
\end{lemm}
\begin{proof}
First, note from Lemma~\ref{lemm.someone.commutation} that $\ment\modT\recent\texi a.\phi\equiv \recent\texi a.\modT\phi$.
And, note that $\texi a.\phi$ and $\texi v.\phi[a\ssm v]$ denote the same syntax, as per Remark~\ref{rmrk.pedant.point}(\ref{item.pedant.texi.v}).

We unpack the clause for $\modellabel{\mrup{a}{\phi}{v}}{\acontext}$ in Figure~\ref{fig.3.derived} and note that it returns $\tvT$ when there exist $v\in\Val$ and $0\leq n'<n$ such that $n',p,O\mentval \modT\phi[a\ssm v]$, and $\tvF$ otherwise.
But this is also what $\modellabel{\recent\texi v.\modT(\phi[a\ssm v])}{\acontext}$ does, so by Lemma~\ref{lemm.equiv}(\ref{item.equiv.1}) we are done.
\end{proof}

Lemma~\ref{lemm.mru.today.yesterday} connects $\mrup{a}{\phi}{v}$ to the fact that our notion of time is discrete and so has a well-defined notion of a `tomorrow' (i.e. a `next timestep'):
\begin{lemm}
\label{lemm.mru.today.yesterday}
Suppose $\texi a.\phi$ is a closed predicate and $v\in\Val$ and $\acontext$ is a context. 
.
Then
$$
\acontext\ment \tomorrow(\mrup{a}{\phi}{v}) \equiv (\modT\phi[a\ssm v] \tor ((\tneg\modT\texi a.\phi) \tand \mrup{a}{\phi}{v})).
$$
\end{lemm}
\begin{proof}
It is a fact of Figure~\ref{fig.3.derived} that $\acontext\ment\tomorrow\yesterday\psi\equiv\psi$ for any closed $\psi$.
The result follows from this observation and from the clause for $\mrup{a}{\phi}{v}$ in Figure~\ref{fig.3.derived}. 
\end{proof}



\begin{corr}
\label{corr.mru.today.to.yesterday}
Suppose $\texi a.\phi$ is a closed predicate and $v\in\Val$ and $\acontext$ is a context. 
Then:
\begin{enumerate*}
\item\label{item.mru.today.implies.tomorrow}
$\acontext\ment \phi[a\ssm v] \timpc \tomorrow(\mrup{a}{\phi}{v})$.
\item\label{item.mru.tomorrow.options}
$\acontext\ment \tomorrow(\mrup{a}{\phi}{v}) \timpc (\phi[a\ssm v]\tor \mrup{a}{\phi}{v})$. 
\end{enumerate*}
\end{corr}
\begin{proof}
Elementary from Lemma~\ref{lemm.mru.today.yesterday} using \strongmodusponens.
\end{proof}


\begin{lemm}
\label{lemm.mru.to.recent}
Suppose $\texi a.\phi$ is a closed predicate %
and $v,v'\in\Val$ and $\acontext$ is a context. 
Then:
\begin{enumerate*}
\item\label{item.mru.to.recent.2}
$\acontext\ment(\mrup{a}{\phi}{v}\tand\mrup{a}{\phi}{v'}) \timpc \recent(\phi[a\ssm v]\tand\phi[a\ssm v'])$.
\item\label{item.mru.to.recent.1}
$\acontext\ment(\mrup{a}{\phi}{v}) \timpc \recent(\phi[a\ssm v])$.

If $\phi=\modality{\tf Q}\tf P(a)$ for a unary predicate symbol $\tf P\in\PredSymb$ and some prefix of modalities $\modality{\tf Q}$, then we can elide the variable symbol $a$ as per Notation~\ref{nttn.hos} and Figure~\ref{fig.hos}, and this becomes
$$
\acontext\ment\mru{(\modality{\tf Q}\tf P)}{v} \timpc \recent\modality{\tf Q}\tf P(v).
$$
\end{enumerate*}
\end{lemm}
\begin{proof}
We consider each part in turn:
\begin{enumerate}
\item
Unpacking the clause for $\mruarrow$ in Figure~\ref{fig.3.derived}, if $v$ and $v'$ are values that most recently make $\phi[a\ssm\text{-}]$ be $\tvT$, then certainly $\recent(\phi[a\ssm v]\tand\phi[a\ssm v'])$.
\item
Follows from part~\ref{item.mru.to.recent.2} as the special case that $v'=v$.
\qedhere\end{enumerate}
\end{proof}


\section{\lowercase{n}-twined semitopologies}
\label{sect.n-twined}



\subsection{Antiseparation properties}

We are particularly interested in models over semitopologies whose open sets have certain \emph{antiseparation properties}, by which we mean properties that involve the existence of nonempty intersections between open sets.\footnote{For comparison, topology is often concerned with \emph{separation} properties, like Hausdorff separation, which have to do with existence of \emph{empty} intersections between open sets.  In this sense, the interest in this Section is the precise dual to where the emphasis traditionally lies in topology.}

The design space of antiseparation properties is rich, as the first author explored in~\cite{gabbay:semdca} --- but for our needs in this paper we will just one canonical property which we will call being \emph{$n$-twined}.
We define the $n$-twined property in Definition~\ref{defn.twined}, we explain why it matters in Remark~\ref{rmrk.motivate.n-twined}, and then we study its properties.


\begin{defn}
\label{defn.twined}
\leavevmode
Suppose $(\Pnt,\opens)$ is a semitopology and $n\in\Ngeqz$.
Then call $(\Pnt,\opens)$ \deffont{$n$-twined} when every collection of $n$ nonempty open sets has a nonempty open intersection.
In symbols:
$$
\Forall{O_1,\dots,O_n\in\opensne}O_1\cap\dots\cap O_n\neq\varnothing.
$$
\end{defn}

\begin{lemm}
\label{lemm.1-twined}
We note two edge cases:
\begin{enumerate*}
\item
A semitopology is $0$-twined if and only if it is nonempty (Notation~\ref{nttn.empty.semitopology}).
\item
Every semitopology is $1$-twined.
\end{enumerate*}
\end{lemm}
\begin{proof}
We consider each part in turn:
\begin{enumerate}
\item
If $n=0$ then we need to show that $\Pnt$ (the intersection of zero nonempty sets) is not equal to $\varnothing$.
by Notation~\ref{nttn.empty.semitopology} this is true precisely when the semitopology is nonempty. 
\item
By Definition~\ref{defn.twined}, we must show that $O\neq\varnothing$ for every $O\in\opensne$.
This is true.
\qedhere\end{enumerate}
\end{proof}

\begin{rmrk}[Discussion of the special cases $n=2$ and $n=3$]
\label{rmrk.motivate.n-twined}
Lemma~\ref{lemm.1-twined} shows that being $0$-twined and $1$-twined are straightforward properties.
In contrast, the $2$- and $3$-twined semitopologies will be very interesting to us in this paper.
We give a high-level discussion of why this is:
\begin{itemize}
\item
Suppose we are in a distributed system running on a $2$-twined semitopology.
We assume that participants are honest.
Suppose further that we have 
\begin{itemize*}
\item
a nonempty open set of participants $O\in\opensne$ who agree that some predicate $\phi$ is true (respectively: valid), and 
\item
a nonempty open set of participants $O'\in\opensne$ who agree that some predicate $\phi'$ is true (respectively: valid).
\end{itemize*}
Then because the system is $2$-twined, there exists some $p\in O\cap O'$ who agrees that both $\phi$ and $\phi'$ are true (respectively: valid).\footnote{We will render this in our logic later on, in Lemma~\ref{lemm.supertwined.qcbt}(\ref{item.intertwined.qcbt.2}).}

This simple observation underlies guarantees of consistent outputs in many distributed algorithms, because it guarantees that $\phi$ and $\phi'$ are not logically conflicting --- e.g. that we do not have something like $\phi'=\modF\phi$ --- because if they were, then an honest $p$ could not believe both.

In the special case that $\phi$ is `the output of our consensus algorithm is $v$' and $\phi'$ is `the output of our consensus algorithm is $v'$', the $2$-twined property makes the strong guarantee that where consensus algorithms succeed (i.e. where participants do not crash), the output is consistent. 
\item
Suppose we are in a distributed system running on a $3$-twined semitopology.
Some participants may be dishonest, in that they may report inconsistent predicates. 

Suppose further that we have nonempty open sets $O,O'\in\opensne$ of participants who agree on the truth (respectively: validity) of predicates $\phi$ and $\phi'$ respectively.
Suppose further that we have a nonempty open set $O''\in\opensne$ of participants who are honest (meaning that we bound dishonest participants to at most some \emph{closed} set). 
Because the system is $3$-twined, there exists some $p\in O\cap O'\cap O''$ --- this $p$ is a participant who agrees on the truth (respectively: validity) of $\phi$ and $\phi'$ and is also honest.\footnote{We will render this in our logic twice: in Lemma~\ref{lemm.supertwined.qcbt}(\ref{item.supertwined.qcbt.2}) as a general result, and then as a useful special case in Corollary~\ref{corr.two.quorums.intersect.correct}(\ref{item.two.quorums.intersect.correct.supertwined}).}

In this case, $\phi\tand\phi'$ must be consistent, because if it were not then the honest participant $p$ could not believe both $\phi$ and $\phi'$.
Just as for the $2$-twined case, this simple argument underlies many consistency arguments in distributed systems in the presence of dishonest participants.
\end{itemize}
There is no barrier to considering $n$-twined semitopologies for $n\geq 4$, and this would just give us stronger guarantees of intersection.

A brief survey of intersection properties in the literature is in Remark~\ref{rmrk.ntwined.and.the.literature}
\end{rmrk}


\subsection{Being $n$-twined and being dense / having nonempty open interior}

The property of being $n$-twined is closely related to the topological notion of being a \emph{dense set} and having a \emph{nonempty open interior} (Definition~\ref{defn.dense.subset}).
We explore this in Lemma~\ref{lemm.is} and Proposition~\ref{prop.noi.dense}; then in Proposition~\ref{prop.allbut.supertwined} we specialise the discussion to the important special case of $\AllBut{N}{f}$.

\begin{lemm}
\label{lemm.is}
Suppose $(\Pnt,\opens)$ is a semitopology and $n\in\mathbb N_{\geq 1}$.
Then:
\begin{enumerate*}
\item\label{item.is.supertwined}
$(\Pnt,\opens)$ is $n$-twined if and only if every collection of $n\minus 1$ nonempty open sets has a dense intersection: 
$$
\Forall{O_1,\dots O_{n\minus 1}\in\opensne}\dense(O_1\cap\dots\cap O_{n\minus 1}).
$$
\item\label{item.is.intertwined}
As a corollary, $(\Pnt,\opens)$ is \xxtwined if and only if every nonempty open set is dense: 
$$
\Forall{O\in\opensne}\dense(O).
$$
\end{enumerate*}
\end{lemm}
\begin{proof}
Both parts follow from Definitions~\ref{defn.dense.subset}(1) and~\ref{defn.twined}, just by pushing around definitions.

One corner case merits brief mention: if $n=1$ then we must show that every collection of $0$ sets (i.e. $\varnothing$) is dense.
So, we need to show that $\bigcap\varnothing$ is dense.
But by convention, the intersection of \emph{no} open sets is the whole of $\Pnt$, and this is indeed dense, so we are done. 
\end{proof}


\begin{prop}
\label{prop.noi.dense}
Suppose $(\Pnt,\opens)$ is a semitopology and $n\in\mathbb N_{\geq 1}$.
Then the following are equivalent:
\begin{enumerate*}
\item
$(\Pnt,\opens)$ is $n$-twined. 
\item
$\Forall{P_1,\dots,P_{n\minus 1}\subseteq\Pnt}(\bigwedge_{1{\leq}i{\leq}n\minus 1}\noi(P_i)) \limp \dense(\bigcap_{1{\leq}i{\leq}n\minus 1}P)$.
\end{enumerate*}
\end{prop}
\begin{proof}
Part~2 trivially implies part~1 because we can take $P_1,\dots,P_n\in\opensne$, so that by Definition~\ref{defn.dense.subset} $\noi(P_i)$ for $1\leq i\leq n\minus 1$ holds by construction by Lemma~\ref{lemm.dense.up-closed}(\ref{item.noi.noi}), and we use Lemma~\ref{lemm.is}(\ref{item.is.supertwined}).

Showing that part~1 implies part~2 is only slightly less trivial.
Suppose $(\Pnt,\opens)$ is $n$-twined and $P_1,\dots,P_{n\minus 1}\subseteq\Pnt$ and $\noi(P_i)$ for each $1\leq i\leq n\minus 1$, meaning by Definition~\ref{defn.dense.subset} that for each such $i$ there exists an $O_i\in\opensne$ such that $O_i\subseteq P_i$.
By Lemma~\ref{lemm.is}(\ref{item.is.supertwined}) $O_1\cap\dots\cap O_{n\minus 1}$ is dense, and by Lemma~\ref{lemm.dense.up-closed}(\ref{item.dense.up-closed}) so is $P_1\cap\dots\cap P_n$.
\end{proof}



Proposition~\ref{prop.allbut.supertwined} provides the canonical example of $n$-twined semitopologies.
It mildly generalises known properties in the literature; the case $n=2$ in Proposition~\ref{prop.allbut.supertwined}(\ref{item.allbut.intertwined}) corresponds to the $Q^2$ property from~\cite[Theorem~1]{DBLP:conf/asiacrypt/DamgardDFN07}, and the case $n=3$ corresponds to the $Q^3$ property of~\cite[Definition 2]{DBLP:journals/dc/AlposCTZ24}.
The topology-flavoured presentations in Proposition~\ref{prop.allbut.supertwined}(\ref{item.allbut.NOI.dense}\&\ref{item.allbut.NOI.dense.3}) are new.
\begin{prop}
\label{prop.allbut.supertwined}
Suppose $N,f\in\Ngeqz$ and $P\subseteq\Nset{N}$ and $n\in\mathbb N_{\geq 1}$.
Then:
\begin{enumerate*}
\item\label{item.allbut.intertwined}\label{item.allbut.supertwined}
$N>n*f$ if and only if $\AllBut{N}{f}$ is $n$-twined.
\item\label{item.allbut.NOI.dense}
$\noi(P)$ (and equivalently $P\in\opensne$) if and only if $\mycard P\geq N\minus f$.\footnote{In general semitopologies, $\noi(P)$ ($P$ having a nonempty open interior) is not the same as $P\in\opensne$ ($P$ being a nonempty open set); but in the case of $\AllBut{N}{f}$, these two properties happen to be equal.}
\item\label{item.allbut.NOI.dense.3}
$\dense(P)$ if and only if $\mycard P\geq f\plus 1$.
\end{enumerate*}
\end{prop}
\begin{proof}
We consider each part in turn:
\begin{enumerate}
\item
Suppose $N>n*f$ and suppose $O_1,\dots,O_n$ are open in $\AllBut{N}{f}$.
From Definition~\ref{defn.allbut} this means that they each contain all but $f$ elements of $\Nset{N}$.
It follows that $O_1\cap\dots \cap O_n$ contains (at least) all but $n*f$ elements of $\Nset{N}$.
Since $N\gneq n*f$, this intersection is nonempty. 

Conversely, if $N\leq n*f$ then it is routine to partition $\Nset{N}$ into $n$ sets, each of which has at most $f$ elements.
Then their complements are a set of $n$ open sets whose intersection is empty.
\item
By routine counting arguments.
\item
By routine counting arguments.
\qedhere\end{enumerate}
\end{proof}



\begin{rmrk}
\label{rmrk.ntwined.and.the.literature}
Being \xxtwined corresponds to the \emph{quorum intersection property} of~\cite[Section~2.1, Condition~B2]{DBLP:journals/tocs/Lamport98}, which requires that any two quorums intersect.
Being \xxxtwined also corresponds to a quorum intersection property with a stronger requirement that (assuming there is a quorum of honest participants) any two quorum must intersect in a participant that is honest~\cite[Definition~5.1]{DBLP:journals/dc/MalkhiR98}. 
  
These intersection properties are closely tied to necessary conditions for some consensus protocols (and other distributed algorithms) to function correctly.
Specifically for the case of the semitopology $\AllBut{N}{f}$ from Definition~\ref{defn.allbut}(\ref{item.allbut}), requiring non-empty quorum intersection is equivalent to the condition \( N > 2f \), where \( f \) denotes the number of processes that may crash during the protocol's execution.
This condition is often referred to as \( Q^2 \)~\cite[Theorem~1]{DBLP:conf/asiacrypt/DamgardDFN07}.
Similarly, being \xxxtwined corresponds to the condition \( N > 3f \), commonly referred to as \( Q^3 \)~\cite[Definition 2]{DBLP:journals/dc/AlposCTZ24}.
A more generalised version of \( Q^3 \) exists, known as \( B^3 \)~\cite[Theorem 2]{DBLP:conf/asiacrypt/DamgardDFN07}\cite[Definition 5]{DBLP:journals/dc/AlposCTZ24}, which accounts for greater heterogeneity. 
Specifically, \( Q^3 \) applies to distributed systems where failures are homogeneous, meaning they are globally known and agreed upon by all participants. 
For example, all participants understand that up to \( f \) processes may fail. 
  
In contrast, in a heterogeneous setting, participants may have differing beliefs about \emph{who} might fail, leading to a more nuanced description of trust assumptions. 
In this context, \( B^3 \) generalises \( Q^3 \) by incorporating these variations in trust and failure assumptions.
Readers interested in exploring this topic further can consult~\cite{DBLP:conf/asiacrypt/DamgardDFN07,DBLP:journals/dc/AlposCTZ24}.
\end{rmrk}
  
  
  
  

\subsection{The canonical properties}

For this subsection we fix a signature $\Sigma$ and a model $\mathcal M=(\Pnt,\opens,\Val)$.

The results in this Subsection express how the underlying semitopology of this model $(\Pnt,\opens)$ being \xxtwined or \xxxtwined manifests itself in the logic, in ways that will be relevant in our later proofs:
\begin{lemm}
\label{lemm.supertwined.qcbt}
Suppose $(\Pnt,\opens)$ is a semitopology and $n\in\Time$ and $\avaluation$ is a valuation and suppose $\phi$, $\psi$, and $\chi$ are pointwise in $(n,\avaluation)$ (Definition~\ref{defn.pointwise.meaning}). 
Then: 
\begin{enumerate*}
\item\label{item.intertwined.qcbt}
If $(\Pnt,\opens)$ is \xxtwined then we have:
\begin{enumerate*}
\item\label{item.intertwined.3}
If $n\mentval \QuorumBox\phi$ then $n\mentval\CoquorumDiamond\phi$. 
\item\label{item.intertwined.qcbt.2}
If $n\mentval \QuorumBox\phi \tand \QuorumBox\psi$ then $n\mentval\someoneAll(\phi\tand\psi)$. 
\item\label{item.intertwined.qcbt.3}
If $n\mentval \QuorumBox\correct{\phi} \tand \QuorumBox\phi$ then $n\mentval\modT\someoneAll\phi$.
\end{enumerate*}
\item\label{item.supertwined.qcbt}
If $(\Pnt,\opens)$ is \xxxtwined then we have: 
\begin{enumerate*}
\item\label{item.supertwined.qcbt.3}
If $n\mentval \QuorumBox\phi \tand \QuorumBox\psi$ then $n\mentval \CoquorumDiamond(\phi\tand\psi)$.
\item\label{item.supertwined.qcbt.2}
If $n\mentval \QuorumBox\phi \tand \QuorumBox\psi \tand \QuorumBox\chi$ then $n\mentval\someoneAll(\phi\tand\psi\tand\chi)$. 
\item\label{item.supertwined.qcbt.4}
If $n\mentval \QuorumBox\correct{\phi} \tand \QuorumBox\phi$ then $n\mentval \modT\CoquorumDiamond\phi$.
\end{enumerate*}
\end{enumerate*}
Recall that $\correct{\phi}$ is from Definition~\ref{defn.correct}(\ref{item.correct.phi}) and just means $\modTF\phi$.
\end{lemm}
\begin{proof}
We reason as follows:
\begin{enumerate}
\item
Suppose $(\Pnt,\opens)$ is \xxtwined. 
\begin{enumerate*}
\item
We reason as follows:
$$
\begin{array}{r@{\ }l@{\qquad}l}
n\mentval\QuorumBox\phi
\liff&
\noi(\{p\in\Pnt \mid n,p,\avaluation\ment\phi\})
&\text{Proposition~\ref{prop.pointwise.dense.char}(\ref{item.dense.char.pointwise.QuorumBox})}
\\
\limp&
\dense(\{p\in\Pnt \mid n,p,\avaluation\ment\phi\}) 
&\text{Proposition~\ref{prop.noi.dense}, $(\Pnt,\opens)$ $2$-twined} %
\\
\liff&
n\mentval\CoquorumDiamond\phi
&\text{Proposition~\ref{prop.pointwise.dense.char}(\ref{item.dense.char.pointwise.CoquorumDiamond})}
\end{array}
$$
\item
Suppose $n\mentval \QuorumBox\phi\tand\QuorumBox\psi$.
By part~\ref{item.intertwined.3} of this result, $n\mentval\CoquorumDiamond\phi$.
By Proposition~\ref{prop.someone.implies.someoneAll}(\ref{item.someone.implies.someoneAll.4}) $n\mentval \someoneAll(\phi\tand\psi)$ as required.
\item
Suppose $n\mentval \QuorumBox\correct{\phi} \tand \QuorumBox\phi$.
By part~\ref{item.intertwined.qcbt.2} of this result, $n\mentval\someoneAll(\correct{\phi}\tand\phi)$.
Using %
Proposition~\ref{prop.tv.ment.TF.model}(\ref{item.tv.ment.TF.phi.lr}) $n\mentval\someoneAll\modT\phi$.
By Lemma~\ref{lemm.someone.commutation} $n\mentval\modT\someoneAll\phi$ as required. 
\end{enumerate*}
\item
The reasoning for the \xxxtwined case is just as for the \xxtwined case, except we need intersection of three open sets instead of two.
\qedhere\end{enumerate} 
\end{proof}




Corollary~\ref{corr.two.quorums.intersect.correct} specialises Lemma~\ref{lemm.supertwined.qcbt} to the useful special case of a correct unary predicate symbol, such that $\phi=\correct{\tf P}$ (Definition~\ref{defn.correct}(\ref{item.correct})) and $\psi=\tf P(v)$ and $\chi=\tf P(v')$.
In our later examples, $\tf P$ will take values such as $\tf{accept}$, $\tf{send}$, $\tf{write}$, and $\tf{decide}$, representing basic actions in a consensus protocol:
\begin{corr}
\label{corr.two.quorums.intersect.correct}
Suppose $(\Pnt,\opens)$ is a semitopology and suppose $\tf P$ is a unary predicate symbol and $v,v'\in\Val$ and $n\in\Time$ and $\avaluation$ is a valuation.
Then:
\begin{enumerate*}
\item
If $(\Pnt,\opens)$ is \xxxtwined then:
\begin{enumerate*}
\item\label{item.two.quorums.intersect.correct.supertwined}
If $n\mentval\QuorumBox\correct{\tf P}\tand\QuorumBox\tf P(v)\tand\QuorumBox\tf P(v')$ then $n\mentval\modT\someoneAll(\tf P(v)\tand\tf P(v'))$.
\item\label{item.QuorumBox.correct.supertwined}
If $n\mentval\QuorumBox\correct{\tf P}\tand\QuorumBox\tf P(v)$ then $n\mentval\modT\CoquorumDiamond\tf P(v)$.
\end{enumerate*}
\item\label{item.two.quorums.intersect.correct.intertwined}
If $(\Pnt,\opens)$ is \xxtwined then
$n\mentval\QuorumBox\correct{\tf P}\tand\QuorumBox\tf P(v)$ implies $n\mentval\modT\someoneAll(\tf P(v))$.
\end{enumerate*}
\end{corr}
\begin{proof}
We consider each part in turn:
\begin{enumerate}
\item[(\ref{item.two.quorums.intersect.correct.supertwined})]
Suppose $(\Pnt,\opens)$ is \xxxtwined and $n\mentval\QuorumBox\correct{\tf P}\tand\QuorumBox\tf P(v)\tand\QuorumBox\tf P(v')$.
By Lemma~\ref{lemm.supertwined.qcbt}(\ref{item.supertwined.qcbt.2}) (since $\tf P(v)$ and $\correct{\tf P}$ are pointwise by Lemma~\ref{lemm.atomic.O}(\ref{item.atomic.pointwise.2})) $n\mentval\someoneAll(\correct{\tf P}\tand\tf P(v)\tand\tf P(v'))$.
By two applications of Proposition~\ref{prop.tv.ment.TF.model}(\ref{item.tv.ment.TF.P}) and Lemma~\ref{lemm.commuting.connectives} (for $\tand$)
$n\mentval\someoneAll\modT(\tf P(v)\tand\tf P(v'))$.
By Lemma~\ref{lemm.someone.commutation} $n\mentval\modT\someoneAll(\tf P(v)\tand\tf P(v'))$ as required.
\item[(\ref{item.QuorumBox.correct.supertwined})] 
Suppose $(\Pnt,\opens)$ is \xxxtwined and $n\mentval\QuorumBox\correct{\tf P}\tand\QuorumBox\tf P(v)$.
Using Lemma~\ref{lemm.correct.P.Pv} it follows that $n\mentval\QuorumBox\correct{\tf P(v)}\tand\QuorumBox\tf P(v)$.
We just apply Lemma~\ref{lemm.supertwined.qcbt}(\ref{item.supertwined.qcbt.4}). 
\item[(\ref{item.two.quorums.intersect.correct.intertwined})]
This follows from Lemma~\ref{lemm.supertwined.qcbt}(\ref{item.intertwined.qcbt.3}) just as part~\ref{item.QuorumBox.correct.supertwined} does from Lemma~\ref{lemm.supertwined.qcbt}(\ref{item.supertwined.qcbt.2}).
\qedhere\end{enumerate}
\end{proof}


Lemma~\ref{lemm.allbut.concretely} unpacks concretely what $\QuorumBox\phi$ and $\CoquorumDiamond\phi$ --- which are the subject of many results including of Lemmas~\ref{lemm.supertwined.qcbt} and~\ref{lemm.intertwined.char} --- mean in the semitopology $\AllBut{N}{f}$ from Definition~\ref{defn.allbut}(\ref{item.allbut}), which has $N$ points, and has open sets being either the empty set, or any set with at least $N\minus f$ elements: 
\begin{enumerate*}
\item
$\QuorumBox\phi$ means `for at least $N\minus f$ elements'; and 
\item
$\CoquorumDiamond\phi$ means `for at least $f\plus 1$ elements.
\end{enumerate*}
This special case is particularly relevant because (as the reader may know) $\AllBut{N}{f}$ is commonly used in the literature on distributed algorithms. 
\begin{lemm}
\label{lemm.allbut.concretely}
Suppose that
\begin{itemize*}
\item
$N,f\in\Ngeqz$, 
\item
$(\Pnt,\opens)=\AllBut{N}{f}$, 
\item
$n\in\Time$ and $\avaluation$ is a valuation, and
\item
$\phi$ is a closed predicate that is pointwise in $(n,\avaluation)$.
\end{itemize*}
Then:
\begin{enumerate*}
\item
$n\mentval\QuorumBox\phi$ if and only if 
$\mycard\{p\in\Pnt \mid n,p\mentval\phi\}\geq N\minus f$. 
\item
$n\mentval\CoquorumDiamond\phi$ if and only if 
$\mycard\{p\in\Pnt \mid n,p\mentval\phi\}\geq f\plus 1$. 
\item
$n\mentval\modT\QuorumBox\phi$ if and only if 
$\mycard\{p\in\Pnt \mid n,p\mentval\modT\phi\}\geq N\minus f$. 
\item
$n\mentval\modT\CoquorumDiamond\phi$ if and only if 
$\mycard\{p\in\Pnt \mid n,p\mentval\modT\phi\}\geq f\plus 1$. 
\end{enumerate*}
\end{lemm}
\begin{proof}
We consider each part in turn:
\begin{enumerate}
\item
The proof is elementary but we spell it out.
We reason as follows:
$$
\begin{array}{r@{\ }l@{\quad}l}
n\mentval\QuorumBox\phi
\liff&
\noi(\{p\in\Pnt \mid n,p\mentval \modT\phi\})
&\text{Proposition~\ref{prop.pointwise.dense.char}(\ref{item.dense.char.pointwise.QuorumBox})}
\\
\liff&
\mycard\{p\in\Pnt \mid n,p\mentval \phi\}\geq N\minus f
&\text{Proposition~\ref{prop.allbut.supertwined}(\ref{item.allbut.NOI.dense})}
\end{array}
$$
\item
Much as the previous part, but using Proposition~\ref{prop.pointwise.dense.char}(\ref{item.dense.char.pointwise.CoquorumDiamond}).
\item
By Lemma~\ref{lemm.someone.commutation}, $n\mentval\modT\QuorumBox\phi$ if and only if $n\mentval\QuorumBox\modT\phi$.
We just use part~1 of this result.
\item
Much as the previous part, using part~2 of this result.
\qedhere\end{enumerate}
\end{proof}

\subsection{Further observations on intersection properties}

For this Subsection we fix a signature $\Sigma$ and a model $\mathcal M=(\Pnt,\opens,\Val)$.
\begin{rmrk}
We collect some further results as follows:
\begin{enumerate*}
\item
Lemma~\ref{lemm.intertwined.char} proves converses to the implications in Lemma~\ref{lemm.supertwined.qcbt}, provided that the signature is nonempty.

This is interesting because it tells us that, provided the signature is nondegenerate, Lemma~\ref{lemm.supertwined.qcbt} accurately captures and reflects the intersection properties that interest us. 
\item
Proposition~\ref{prop.reformulate} is a semitopological distillation of parts~\ref{item.intertwined.qcbt.2} and~\ref{item.supertwined.qcbt.2} of Lemma~\ref{lemm.supertwined.qcbt}.
Proposition~\ref{prop.reformulate} \emph{is} Lemma~\ref{lemm.supertwined.qcbt} and Corollary~\ref{corr.two.quorums.intersect.correct}, in semitopological rather than logical form.
\end{enumerate*}
\end{rmrk}

In Lemma~\ref{lemm.intertwined.char} we could replace every instance of the expression `$\modTF\tf P$' with the (by Definition~\ref{defn.correct}(\ref{item.correct.phi}) and Remark~\ref{rmrk.tf.synonyms}) synonymous expression `$\correct{\tf P}$'.
However, we have not set up a corresponding synonym for $\modTB\tf P$ (it would be `$\tf{valid}[\tf P]$'; we just never wrote that definition) so it seems natural to stay with $\modTF$ here.
\begin{lemm}
\label{lemm.intertwined.char}
Suppose $\tf P\in\ns{PredSymb}$ is a predicate constant symbol (Notation~\ref{nttn.logic.terminology}(\ref{item.nullary.ps}); a nullary predicate symbol).\footnote{The proof works just as well for an $n$-ary predicate symbol; we just pick \emph{any} $\arity(\tf P)$-tuple of values $\vec v\in\Pnt^{\arity(\tf P)}$ --- which is possible because we assumed $\Val$ is nonempty in Definition~\ref{defn.model}(\ref{item.model.val}) --- and we use $\tf P(\vec v)$ in the proof instead.}
Then:
\begin{enumerate*}
\item
If\ \  $\ment(\QuorumBox\modTF\tf P\tand\QuorumBox\modTB\tf P)\timpc \someoneAll\tf P$ then $(\Pnt,\opens)$ is \xxtwined. 
\\
(This is a form of converse to Lemma~\ref{lemm.supertwined.qcbt}(\ref{item.intertwined.qcbt.3}).)
\item
If\ \  $\ment(\QuorumBox\modTF\tf P\tand\QuorumBox\modTB\tf P)\timpc \CoquorumDiamond\tf P$ then $(\Pnt,\opensne)$ is \xxxtwined. 
\\
(This is a form of converse to Lemma~\ref{lemm.supertwined.qcbt}(\ref{item.supertwined.qcbt.4}).) 
\item\label{item.intertwined.char.reverse}
As a corollary, provided that the signature $\Sigma$ is has at least one predicate symbol --- i.e. provided $\PredSymb$ in
it is nonempty (Definition~\ref{defn.signature}(\ref{item.predicate.syntax.empty.signature})) --- then reverse implications to those in Lemma~\ref{lemm.supertwined.qcbt} hold.\footnote{No need to also insist that $\Val$ is nonempty, because we assume this in Definition~\ref{defn.model}(\ref{item.model.val}).}
\end{enumerate*}
Recall from Definition~\ref{defn.validity.judgement}(\ref{item.ment.phi}) that $\ment\phi$ means $\Forall{\acontext\in\tf{Ctx}(\Sigma,\mathcal M)}(\acontext\ment\phi)$.
\end{lemm}
\begin{proof}
We consider each part in turn:
\begin{enumerate}
\item
We prove a contrapositive of \strongmodusponens: we suppose $\opens$ is \emph{not} \xxtwined, so there exist $O,O'\in\opensne$ with $O\cap O'=\varnothing$, and we exhibit a valuation such that $\neg(\mentval\someone\tf P)$ (equivalently: $\mentval\modF\someone\tf P$).

We choose a valuation $\avaluation$ such that for any $n\in\Time$, 
\begin{itemize*}
\item
$\avaluation(\tf P)_{n,p}(\aval)=\tvF$ for $p\in O$, and 
\item
$\avaluation(\tf P)_{n,p}(\aval)=\tvB$ for $p\in \Pnt\setminus O$.
\end{itemize*}
The reader can check that $O\mentval \everyone\modTF\tf P$ and $O'\mentval \everyone\modTB P(\aval)$, but $\mentval\modF\someone\tf P$.
\item
We prove a contrapositive: suppose $\opens$ is \emph{not} \xxxtwined, so there exist $O,O',O''\in\opensne$ with $O\cap O'\cap O''=\varnothing$.

We choose a valuation $\avaluation$ such that for any $n\in\Time$, 
\begin{itemize*}
\item
$\avaluation(\tf P)_{n,p}(\aval)=\tvT$ for $p\in O$, and 
\item
$\avaluation(\tf P)_{n,p}(\aval)=\tvB$ for $p\in O'\setminus O$, and 
\item
$\avaluation(\tf P)_{n,p}(\aval)=\tvF$ for $p\in \Pnt\setminus (O\cup O')$.
\end{itemize*}
The reader can check that $O\cup O'\mentval \everyone\modTF\tf P$ and $O\cup O''\mentval \everyone\modTB P(\aval)$, but $O'\cup O''\mentval\modF\everyone\tf P$.
\item
Just by the arguments above.
\qedhere\end{enumerate}
\end{proof}






\begin{nttn}
\label{nttn.cover}
Suppose $\ns X$ is a set.
Then:
\begin{enumerate*}
\item
A \deffont{cover} of $\ns X$ is a set of sets $\mathcal X\subseteq\powerset(\ns X)$ such that $\bigcup\mathcal X=\ns X$.
\item
A \deffont{partition} of $\ns X$ is a cover of $\ns X$ all of whose elements are pairwise disjoint (so $\Forall{X,X'\in\mathcal X}X\cap X'=\varnothing\lor X=X'$).
\end{enumerate*}
\end{nttn}

The reader might like to try giving a proof of Proposition~\ref{prop.reformulate} using direct calculations (without building on the results we have so far).
This is possible, but it is messy.
The use of logic endows the proofs with a structure and cleanness that the direct proof cannot match: 
\begin{prop}
\label{prop.reformulate}
Suppose $(\Pnt,\opens)$ is a semitopology and $n\in\Time$ and $\avaluation$ is a valuation, and suppose $\phi$ is a closed predicate that is pointwise in $(n,\avaluation)$.
Then:
\begin{enumerate*}
\item\label{item.reformulate.1}
If $(\Pnt,\opens)$ is \xxtwined then $n\mentval\CoquorumDiamond\modTB\phi \tor \CoquorumDiamond\modF\phi$. 
\item\label{item.reformulate.1b}
As a corollary, if $(\Pnt,\opens)$ is \xxtwined then any two-element partition of $\Pnt$ (Notation~\ref{nttn.cover}) contains a dense element (Definition~\ref{defn.dense.subset}(\ref{item.dense.subset})).
\item\label{item.reformulate.2}
If $(\Pnt,\opens)$ is \xxxtwined then 
$n\mentval\CoquorumDiamond\modT\phi \tor \CoquorumDiamond\modB\phi \tor \CoquorumDiamond\modF\phi$. 
\item\label{item.reformulate.2b}
As a corollary, if $(\Pnt,\opens)$ is \xxxtwined then any three-element partition of $\Pnt$ contains a dense element.
\end{enumerate*} 
\end{prop}
\begin{proof}
We consider each part in turn:
\begin{enumerate}
\item
By Lemma~\ref{lemm.supertwined.qcbt}(\ref{item.intertwined.qcbt.2})
$n\mentval\QuorumBox\modF\phi\tand\QuorumBox\modTB\phi$ implies $n\mentval\someoneAll(\modF\phi\tand\modTB\phi)$.
Unpacking what $n\mentval\someoneAll(\modF\phi\tand\modTB\phi)$ means using Definition~\ref{defn.validity.judgement} and Figures~\ref{fig.3.derived} and~\ref{fig.3.phi.f}, we see that this is impossible --- for no $p$ and $O$ can we satisfy both $n,p,O\mentval\modF\phi$ and $n,p,O\mentval\modTB\phi$ together --- so that 
$n\nmentval\QuorumBox\modF\phi\tand\QuorumBox\modTB\phi$.
By Lemma~\ref{lemm.para}(\ref{item.para.em})
$n\mentval\tneg(\QuorumBox\modF\phi\tand\QuorumBox\modTB\phi)$, and simplifying further using de Morgan equivalences as per Figure~\ref{fig.easy.equivalences}, we obtain
$n\mentval\CoquorumDiamond\modTB\phi\tor\CoquorumDiamond\modF\phi$ as required.
\item
Suppose $\{P_1,P_2\}$ partitions $\Pnt$.
Assume a $0$-ary predicate symbol $\tf P$, and let $\varsigma$ be the valuation that sets $\varsigma(\tf P)_{n,p}=\tvT$ if $p\in P_1$ and $\varsigma(\tf P)_{n,p}=\tvF$ if $p\in P_2$.\footnote{We can assume $\tf P$ and $\varsigma$ because this result is about semitopologies, so we are free to impose any model structure on it that we like.  It seems particularly nice here how the mere \emph{existence} of a model for our logic gives us a clean proof of a purely semitopological property which, on the face of it, has nothing to do with our three-valued modal logic.}
By part~\ref{item.reformulate.1} and Lemma~\ref{lemm.tand.tor} either $n\mentval \Coquorum\modTB\tf P$ or $n\mentval\Coquorum\modF\tf P$.
By Proposition~\ref{prop.pointwise.dense.char}(\ref{item.dense.char.pointwise.CoquorumDiamond}) either $P_1$ or $P_2$ is dense.
\item
We reason as for the proof of part~\ref{item.reformulate.1}, using Lemma~\ref{lemm.supertwined.qcbt}(\ref{item.supertwined.qcbt.2}) and starting from $\modFB\phi \tand \modTF\phi \tand \modTB\phi$. 
\item
We reason as for the proof of part~\ref{item.reformulate.2}, for a partition $\{P_1,P_2,P_3\}$ and setting $\varsigma(\tf P)_{n,p}=\tvT$ if $p\in P_1$ and $\varsigma(\tf P)_{n,p}=\tvB$ if $p\in P_2$ and $\varsigma(\tf P)_{n,p}=\tvF$ if $p\in P_3$.
\qedhere\end{enumerate}
\end{proof}




 





\section{Declarative Paxos}
\label{sect.declarative.paxos}

\subsection{Case study: a simple protocol}
\label{subsect.simple}

\begin{rmrk}
\label{rmrk.simple.protocol}
Before Paxos, we will warm up by sketching a simple protocol \theory{Simple}, and then we will give four \QLogic axiomatisations of \theory{Simple}, reflecting four different failure assumptions (messages can/cannot get lost; participants can/cannot crash).
\theory{Simple} is as follows:
\begin{enumerate*}
\item
Any participant $p$ may broadcast a $\tf{propose}$ message to all participants.
\item
If a participant receives a $\tf{propose}$ message, then it responds with $\tf{accept}$. 
\item
If the participant $p$, having broadcast its $\tf{propose}$ message, receives a $\tf{accept}$ message, then it declares $\tf{decide}$.
\end{enumerate*}
\end{rmrk}


\begin{rmrk}
\theory{Simple} is very simple!
Only one value can be proposed, so we do not even bother to state it; we just `propose'.
Any participant accepts immediately that it hears someone propose. 
Any participant that proposes, decides immediately that it hears that someone accepted.

For comparison, in Paxos \emph{multiple values} are available, so we need to make sure that participants all decide on the \emph{same} value if they decide on anything at all. 
This makes things much harder: so Paxos has multiple rounds; there is a designated leader participant at each round; the leader must choose a value to propose; and elaborate communication is in place to ensure that (subject to suitable precisely-stated assumptions) a single choice of agreed value will eventually be decided.

But that is for later.  
Even the simple protocol \theory{Simple} has plenty of subtlety, which will allow us to exhibit (in `toy' form) many of the techniques that we will then apply to Paxos:
\end{rmrk}


\begin{rmrk}
We will axiomatise \theory{Simple} subject to two parameters, yielding four possible axiomatisations:
\begin{enumerate*}
\item
Do all messages arrive, or might there be a finite initial period during which the network may be unreliable so that during that time, messages might not arrive?
\item
Can participants crash (and if so, how many)?
\end{enumerate*}
\end{rmrk}

\subsubsection*{All messages arrive; no participants crash}

First, we assume an ideal situation that all messages arrive, and no participants crash.
We can axiomatise this as follows:
$$
\begin{array}{l@{\qquad}l}
\rulefont{SimpAccept?}&
\tf{accept} \timpc \someoneAll\tf{propose} 
\\
\rulefont{SimpDecide?}&
\tf{decide} \timpc (\tf{propose}\tand\someoneAll\tf{accept}) 
\\
\rulefont{SimpAccept!}&
(\someoneAll\tf{propose}) \timpc \tf{accept}
\\
\rulefont{SimpDecide!}&
(\tf{propose}\tand\someoneAll\tf{accept}) \timpc \tf{decide}
\\
\rulefont{Correct}&
\correct{\tf{propose},\tf{accept},\tf{decide}}
\end{array}
$$
\noindent Recall from Definition~\ref{defn.three.imp} and Proposition~\ref{prop.mp.for.tnotor} that $\timpc$ is a \emph{strong implication} and means `if the LHS is $\tvT$ then so is the RHS', and 
recall from Proposition~\ref{prop.pointwise.dense.char}(\ref{item.dense.char.pointwise.someoneAll}) that $\someoneAll\tf P$ means `$\tf P$ somewhere', where $\tf P\in\{\tf{propose},\tf{accept},\tf{decide}\}$.

Then we explain the axioms (briefly) as follows:
\begin{itemize*}
\item
We call \rulefont{SimpAccept?} and \rulefont{SimpDecide?} \emph{backward rules}; they express backward versions of the rules in Remark~\ref{rmrk.simple.protocol}, namely that:
\begin{itemize*}
\item
if a participant asserts $\tf{accept}$ then some participant must have asserted $\tf{propose}$; and 
\item
if a participant asserts $\tf{decide}$ then it must have asserted $\tf{propose}$, and some participant must have asserted $\tf{accept}$.
\end{itemize*}
\item
We call \rulefont{SimpAccept!} and \rulefont{SimpDecide!} \emph{forward rules}.
They express forward versions of the rules in Remark~\ref{rmrk.simple.protocol}.
Because our failure assumptions are that nothing fails --- all messages arrive; no participants crash --- the forward rules are in this case just the backward rules, reversed.
\item
Axiom \rulefont{Correct} expresses that all predicates return $\tvT$ and $\tvF$ (not $\tvB$), as per Definition~\ref{defn.correct}(\ref{item.correct}). 
This reflects our assumption that no participants crash, because we model a computation crashing %
as the predicate describing it returning $\tvB$.
\end{itemize*}
Note that $\tf{propose}$, $\tf{accept}$, and $\tf{decide}$ above are not messages; they are propositions!
There is no message-passing.
$\tf{accept}$ being true at a participant represents that the participant accepts; this is implemented in the algorithm by the participant broadcasting an accept message, and represented in the logic by $\tf{accept}$ taking truth-value $\tvT$ at that participant.

Our example protocol is simple, but we hope it is already starting to illustrate how the various parts of our logic come together to express the protocol axiomatically.
We continue:

\subsubsection*{All messages arrive; participants can crash}

We assume that all messages arrive, and participants can crash --- but we assume there is a dense set of uncrashed participants.
Recall from Definition~\ref{defn.dense.subset}(\ref{item.dense.subset}) that a \emph{dense set} is one that intersects every open set; in particular, this means that every open set of participants contains a participant that is not crashed.
We model being crashed as returning the third truth-value $\tvB$.
We can axiomatise this as follows:
$$
\begin{array}{l@{\qquad}l}
\rulefont{SimpAccept?}&
\tf{accept} \timpc \someoneAll\tf{propose} 
\\
\rulefont{SimpDecide?}&
\tf{decide} \timpc (\tf{propose}\tand\someoneAll\tf{accept}) 
\\
\rulefont{SimpAccept!}&
(\someoneAll\tf{propose}) \tnotor \tf{accept}
\\
\rulefont{SimpDecide!}&
(\tf{propose}\tand\someoneAll\tf{accept}) \tnotor \tf{decide}
\\
\rulefont{Uncrashed}&
\QuorumBox\modTB\tf P \tnotor \someoneAll\modT\tf P
\quad \tf P\in\{\tf{propose},\tf{accept},\tf{decide}\}
\end{array}
$$ 
\noindent 
The backward rules are unchanged: if a participant accepts then from the form of the rules, it \emph{must} be the case that somebody proposed; and if a participant decides then from the form of the rules, it \emph{must} be the case that the participant proposed and some participant accepted.
 
The forward rules \rulefont{SimpAccept!} and \rulefont{SimpDecide!} have changed, and they now use the weak implication $\tnotor$.
By Definition~\ref{defn.three.imp} and Proposition~\ref{prop.mp.for.tnotor}, $\tnotor$ means `if the LHS is $\tvT$ then the RHS is $\tvT$ or $\tvB$'.
The $\tvB$ truth-value represents `crashed', so that (for example) \rulefont{SimpAccept!} expresses that for each participant $p$, if some (possibly other) participant $q$ asserts $\tf{propose}$, then $p$ asserts $\tf{accept}$ \emph{or} $p$ has crashed. 

Axiom \rulefont{Uncrashed} reflects that any a quorum (open set) of participants contains an uncrashed participant; e.g. in topological language, that the set of uncrashed participants is dense.
More precisely: \rulefont{Uncrashed} asserts that if an open set of participants returns $\tvT$ or $\tvB$ for $\tf P$, then at least one of them (is not crashed and so) returns $\tvT$.

Note from Lemma~\ref{lemm.someone.commutation} that $\Quorum\modTB$ is equivalent to $\modTB\Quorum$, and $\someoneAll\modT$ is equivalent to $\modT\someoneAll$, so that an equivalent form of \rulefont{Uncrashed} is $\modTB\QuorumBox\tf P \tnotor \modT\someoneAll\tf P$.
In our axiomatisations it is usually helpful to pull $\modTB$ and $\modT$ up to the topmost level possible in the predicate, so we may favour that form, but this is just a matter of convenience; meaning is unchanged. 


\subsubsection*{Messages might not arrive; no participants crash}

We assume that there is an initial period during which the network might not deliver messages, but the network eventually stabilises (in the literature this is called \emph{eventual synchrony}).
No participants crash.
We can axiomatise this as follows:
$$
\begin{array}{l@{\qquad}l}
\rulefont{SimpAccept?}&
\tf{accept} \timpc \someoneAll\tf{propose} 
\\
\rulefont{SimpDecide?}&
\tf{decide} \timpc (\tf{propose}\tand\QuorumBox\tf{accept}) 
\\
\rulefont{SimpAccept!}&
\final((\someoneAll\tf{propose}) \timpc \tf{accept})
\\
\rulefont{SimpDecide!}&
\final((\tf{propose}\tand\someoneAll\tf{accept}) \timpc \tf{decide})
\\
\rulefont{Correct}&
\correct{\tf{propose},\tf{accept},\tf{decide}}
\end{array}
$$ 
\noindent The backward rules are unchanged, but the forward rules \rulefont{SimpAccept!} and \rulefont{SimpDecide!} now have a $\final$ modality.
This is from Figure~\ref{fig.3.derived} and it means `eventually, this holds forever', reflecting an assumption that eventually the network stabilises.

Note that \QLogic does not explicitly represent message-passing, but this does not matter as much as one might suppose!
The $\final$ modality captures the essence of what we need, namely, that eventually the forward rules become valid. 

\subsubsection*{Messages might not arrive; participants can crash}

We assume that messages might not arrive, and participants can crash but there is a dense set of uncrashed participants (so that every quorum of participants contains at least one uncrashed participant).
We can axiomatise this just by combining the ideas above, as follows:
$$
\begin{array}{l@{\qquad}l}
\rulefont{SimpAccept?}&
\tf{accept} \timpc \someoneAll\tf{propose} 
\\
\rulefont{SimpDecide?}&
\tf{decide} \timpc (\tf{propose}\tand\someoneAll\tf{accept}) 
\\
\rulefont{SimpAccept!}&
\final((\someoneAll\tf{propose}) \tnotor \tf{accept})
\\
\rulefont{SimpDecide!}&
\final((\tf{propose}\tand\someoneAll\tf{accept}) \tnotor \tf{decide})
\\
\rulefont{Uncrashed}&
\modTB\QuorumBox\tf P \tnotor \modT\someoneAll\tf P
\quad \tf P\in\{\tf{propose},\tf{accept},\tf{decide}\}
\end{array}
$$ 
\noindent The axiomatisations of Declarative Paxos and Simpler Declarative Paxos in Figures~\ref{fig.ThyPaxOne} and~\ref{fig.ThySPax} are more elaborate than the axioms above, because Paxos is a more complicated protocol, but the basic ideas are as above; just scaled up. 


\subsection{A high-level view of the Paxos algorithm}



\begin{rmrk}[High-level view of Paxos]
\label{rmrk.high-level.paxos}
Paxos allows participants to crash, and messages to fail to arrive at their destination, but it assumes all participants are honest (i.e. that if they have not crashed then they follow the rules of the algorithm).
The reader can find detailed treatments of Paxos in \cite{DBLP:journals/tocs/Lamport98,lamport2001paxos,cachin2011paxos} (a Wikipedia article~\cite{wiki:Paxos_(computer_science)} is also accessible and clear).
Here is a high-level view of the algorithm:
\begin{enumerate*}
\item\label{item.l.proposes}
\deffont{Leader proposes:}\quad

A designated leader participant (if it does not crash) sends messages to all participants, \emph{proposing} some unique value $v$.

Message-passing is unreliable, so messages may fail to arrive.
\item\label{item.p.sends}
\deffont{Participant sends:}\quad

If a participant $p$ receives a propose $v$ message from the leader --- it is assumed that every participant knows who the leader is, and other participants cannot forge the leader's signature --- then $p$ sends a \emph{send} message back to the leader, which carries the value $v_p$ that $p$ most recently accepted, along with the timestamp of when $v_p$ was accepted.\footnote{For reasons of compact presentation we do not explicitly represent the timestamps in our axiomatisation; i.e. $\tf{send}$ is a unary predicate in Definition~\ref{defn.declarative.paxos}(\ref{item.ThyPaxOne.sig}), not a binary predicate.
If this bothers the reader then we can just say that, actually, $v_p$ contains a timestamp of the stage $n\in\Time$ it was accepted; i.e. that $v_p$ is actually implemented as a pair $(v_p',n_p)$ of a value and a time.} 
If $p$ has never accepted any value, so no most recently accepted value exists, then $p$ sends back a special $\udfn$ value.
\item\label{item.l.writes}
\deffont{Leader writes:}\quad

The leader participant sends a message to all participants, \emph{writing} the most recently-timestamped value $v_p$ mentioned in any \emph{send} message it received from within the first quorum of participants that it hears back from; if the leader does not hear back from any quorum, then it does not write anything.
\item\label{item.p.accepts}
\deffont{Participant accepts:}\quad

If a participant receives a write message from the leader for some value, then it sends a message back \emph{accepting} that value.
\item\label{item.l.decides}
\deffont{Leader decides:}\quad

If the leader receives accept messages for some value from a quorum of participants, it decides on that value, and sends out a \emph{decide} message to all participants with that value.
\item\label{item.p.decides}
\deffont{Participant (non-leader) decides:}\quad

If a participant receives a decide message from the leader for some value, then it decides on that value and stops.
\end{enumerate*}
\end{rmrk}

\begin{rmrk}[High-level view of our axiomatisation]
\label{rmrk.elisions}
The theory $\ThyPax$ in Definition~\ref{defn.ThyPaxOne} and Figure~\ref{fig.ThyPaxOne} is a declarative abstraction of Paxos-the-algorithm.
It stands in relation to implementations of Paxos in roughly the same way that (say) the mathematical specification of the Fibonacci function stands in relation to its concrete implementations, in the sense that we elide implementational detail to preserve (and reveal) a mathematical ideal of what the implementation does.

\ThyPax elides two features of the concrete implementation:
\begin{enumerate*}
\item
There is no explicit message-passing; there are only predicates returning truth-values.

This may seem counterintuitive --- isn't Paxos \emph{supposed} to be a message-passing algorithm? --- but we shall see that a lot of the character of Paxos survives.
Indeed, arguably (and certainly for the first author) it is much easier to see what is really going on once we have done this.
\item\label{item.elide.algorithmic.time}
There is no abstract machine, and therefore no \deffont{algorithmic time}, by which we mean no sequence of state-changes of the abstract machine. 

The only notion of time is the \emph{stage} $n\in\Time$ from the context as defined in Definition~\ref{defn.contexts}(\ref{item.context.stage}).
We may call this \deffont{logical time} and (when we are modelling Paxos) it represents the sequence of what is in the distributed consensus algorithms literature often called \emph{stages}, or \emph{ballot numbers}.\footnote{In a distributed algorithm, at any point in algorithmic time, different participants in the algorithm may be at different points in logical time.  A practical example: consider several participants racing across 100 metres. In logical time, they are both just counting to 100.  In algorithmic time, one of them will get to 100 \emph{first}, and win.  In this paper, we only care about the logical aspects of the process; there may still be a winner participant, but this would implemented just as a unary predicate $\tf{winner}$ such that $100\ment\texiaffine\tf{winner}$.} 
\end{enumerate*}
We will prove standard Paxos correctness properties for \ThyPax in Section~\ref{sect.paxos.correctness.properties}. 
\end{rmrk}


\subsection{Definition of Declarative Paxos}

\begin{figure}
$$
\begin{array}{l@{\quad}r@{\ }l}
\figunderline{Backward rules, for backward inference} 
\figskip
\rulefont{PaxPropose?}&
& \tall v.\ \tf{propose}(v)\timpc(\tf{leader}\tand v\tneq\udfn)
\figskip
\rulefont{PaxSend?}&
& \tall v.\ \tf{send}(v) \timpc 
\begin{array}[t]{l}
\someoneAll\texi\tf{propose}\ \tand 
\\
\bigl(
(\mru{\tf{accept}}{v})\ \tor  
(\tneg\modT\recent\texi\tf{accept}\tand v\teq\udfn)\bigr)
\end{array}
\figskip
\rulefont{PaxWrite?}&
& \tall v.\ \tf{write}(v) \timpc 
\begin{array}[t]{l}
\tf{leader}\tand 
\\
\Quorum(\Box\texi\tf{send} \tand 
\\
\qquad\bigl((\mru{(\someone\tf{accept})}{v}) \tor
(\tf{propose}(v)\tand\tneg\modT\recent\texi\someone\tf{accept})\bigr)) 
\end{array}
\figskip
\rulefont{PaxAccept?}&
& \tall v.\ \tf{accept}(v) \timpc \someoneAll\tf{write}(v)
\figskip
\rulefont{PaxDecideL?}&
& \tf{leader} \timpc \tf{decide}(v) \timpc \QuorumBox\tf{accept}(v)
\figskip
\rulefont{PaxDecide\tneg L?}&
& \tall v.\ \tneg\tf{leader} \timpc \tf{decide}(v) \timpc \someoneAll(\tf{leader}\tand\tf{decide}(v))
\figskip\figskip
\figunderline{Forward rules, for making progress} 
\figskip
\rulefont{PaxPropose!}&
& \tf{leader}\tnotor\texi\tf{propose}
\figskip
\rulefont{PaxSend!}&
& \final(\someoneAll\texi\tf{propose} \tnotor \texi\tf{send}) 
\figskip
\rulefont{PaxWrite!}&
& \final\bigl((\tf{leader}\tand (\QuorumBox\texi\tf{send}) \tand \texi\tf{propose}) \tnotor \texi\tf{write}\bigr)
\figskip
\rulefont{PaxAccept!}&
& \final(\texi\someoneAll\tf{write}\tnotor \texi\tf{accept}) 
\figskip
\rulefont{PaxDecideL!}&
& \final((\tf{leader} \tand\texi\QuorumBox\tf{accept}) \tnotor \texi\tf{decide}) 
\figskip
\rulefont{PaxDecide\tneg L!}&
& \final((\tneg\tf{leader} \tand \someoneAll(\tf{leader}\tand\texi\tf{decide}))\tnotor \texi\tf{decide}) 
\figskip
\figskip
\figunderline{Other rules, for correctness \& liveness}
\figskip
\rulefont{LdrExist}&
&\modTF\tf{leader}\tand\modT\someoneAll\tf{leader}
\figskip
\rulefont{LdrExt}_\phi&
& \tf{leader}\timpc\someoneAll(\tf{leader}\tand\phi)\timpc\phi
\quad\text{every closed predicate $\phi$}
\figskip
\rulefont{LdrCorrect}&
& \infinitely(\tf{leader}\timpc\correct{\tf{propose},\tf{write},\tf{decide}}) 
\figskip
\rulefont{PaxPCorrect}&
& 
\begin{array}[t]{l}
\QuorumBox\correct{\tf{propose}} \tand \QuorumBox\correct{\tf{send}} \tand \QuorumBox\correct{\tf{accept}} \tand 
\\
\qquad\QuorumBox\correct{\tf{write}} \tand \QuorumBox\correct{\tf{decide}}
\end{array}
\figskip
\rulefont{PaxPropose01}&
& \texiaffine \tf{propose}
\figskip
\rulefont{PaxWrite01}&
& \texiaffine\tf{write}
\figskip
\rulefont{Pax2Twined}&
& \tall v.\ \QuorumBox\tf{accept}(v) \timpc \CoquorumDiamond\tf{accept}(v)
\end{array}
$$
\caption{$\Theta_\ThyPax$: Axioms of Declarative Paxos (Definition~\ref{defn.declarative.paxos})}
\label{fig.logical.paxos}
\label{fig.ThyPaxOne}
\end{figure}


Recall from Definition~\ref{defn.theory.axiom}(\ref{item.theory}) that a \emph{theory} is a pair of a \emph{signature} and a set of \emph{axioms} over that signature.
Then:
\begin{defn}
\label{defn.declarative.paxos}
\label{defn.ThyPaxOne}
We define \deffont{Declarative Paxos} to be the logical theory 
$$
\ThyPax=(\Sigma_{\ThyPax},\Theta_{\ThyPax}),
$$ 
where:
\begin{enumerate*}
\item\label{item.ThyPaxOne.sig}
the signature $\Sigma_{\ThyPax}$ is 
$$
\Sigma_{\ThyPax}=[\tf{leader}:0;\ \tf{propose},\tf{send},\tf{write},\tf{accept},\tf{decide}:1] ,
$$ 
as per the notation in Definition~\ref{defn.signature}(\ref{item.predicate.syntax.signature}), and
\item\label{item.ThyPaxOne.axioms}
the axioms $\Theta_{\ThyPax}$ are as written in Figure~\ref{fig.logical.paxos}.
\end{enumerate*}
In axioms \rulefont{PaxPropose?} and \rulefont{PaxSend?} we assume an element $\udfn\in\Val$, to model the `undefined' value from Paxos (it does not matter what $\udfn\in\Val$ is; we just have to choose some element and use it consistently).
\end{defn}

\begin{rmrk}
\label{rmrk.functional.notation.in.paxos}
For convenient reference, we give pointers to help unpack the notation used in Figure~\ref{fig.logical.paxos}.
Note that this is just an overview focussing on items of notation used in the axioms, to help with reading them; in-depth discussions of the axioms will follow: 
\begin{enumerate*}
\item
We use the functional notation from Notation~\ref{nttn.hos} to elide quantified variables where convenient.
For example, axioms \rulefont{PaxPropose!} and \rulefont{PaxAccept!} written in full are 
$$
\tf{leader}\tnotor\texi\thea.\tf{propose}(\thea)
\quad\text{and}\quad
\final((\texi\thea.\someoneAll\tf{write}(\thea))\tnotor\texi\thea.\tf{accept}(\thea)) .
$$
\item
$\tnotor$ and $\timpc$ are the \emph{weak} and \emph{strong} implications respectively.
Truth-tables are in Figure~\ref{fig.3} and a discussion is in Subsection~\ref{subsect.3.imp}.
In particular, by Proposition~\ref{prop.mp.for.tnotor}, 
\begin{itemize*}
\item
$\phi\tnotor\phi'$ means `if $\phi$ is true then $\phi'$ is valid' (where `valid' means `$\tvT$ or $\tvB$') and 
\item
$\phi\timpc\phi'$ means `if $\phi$ is true then $\phi'$ is true'. 
\end{itemize*}
\item
The compound modalities --- $\someoneAll$, $\everyoneAll$, $\QuorumBox$, and $\CoquorumDiamond$ --- are discussed in Subsection~\ref{subsect.compound.modalities}.
Their meaning in full generality is in Lemma~\ref{lemm.dense.char}, but because the axioms only apply them to \emph{pointwise} predicates, their meanings simplify to `somewhere' ($\someoneAll$), `everywhere' ($\everyoneAll$), `for a set with a nonempty open interior' ($\QuorumBox$), and `for a dense set' ($\CoquorumDiamond$) respectively, as per Proposition~\ref{prop.pointwise.dense.char}. 
\item\label{item.discuss.final}
$\final$ means $\sometime\forever$ and is from Figure~\ref{fig.3.derived}.
It means `eventually, this holds for all future times'. 

In our axiomatisation this is used to reflect a standard assumption in distributed systems of \deffont{eventual network synchrony}, or synonymously \deffont{global stabilisation (of the network)}, meaning that the communications network may be unstable and fail to pass messages at first, but then after a finite but unbounded amount of time, it becomes stable and messages arrive.
Our logic does not explicitly represent messages, but the use of $\final$ in our forward rules reflects this assumption about the underlying implementation. 
\item
$\texiaffine$ in \rulefont{PaxPropose01} and \rulefont{PaxWrite01} means `for zero or one elements', as per Subsection~\ref{subsect.unique.affine.existence}, so these two axioms intuitively state that each participant proposes or writes at most one value.

Note that our logic is three-valued, so what these axioms mean \emph{precisely} is something a little richer than is the case in two-valued logic.
In particular, if $\Forall{v}\modellabel{\tf{write}(v)}{\acontext}=\tvB$ --- corresponding to the case where a participant is crashed --- then $\acontext\ment\texiaffine\tf{write}$, \emph{even though} $\acontext\ment\tf{write}(v)$ holds for every $v$. 
See the discussion and results in Subsection~\ref{subsect.unique.affine.existence}.
\item
$\correct{\tf{propose}}$ means that $\tf{propose}$ returns $\tvT$ or $\tvF$ (not $\tvB$) on every value; see Definition~\ref{defn.correct}(\ref{item.correct}).
Similarly for the other unary predicate symbols in \rulefont{PaxPCorrect}.
\item
$\mru{\tf{accept}}{v}$ means `a most recently accepted value is $v$'.
See the discussion and definition in Subsection~\ref{subsect.mrup}.
\item
Finally, note that we intend the predicates in Figure~\ref{fig.logical.paxos} to be read as axioms.

What it means for an axiom to be valid in a model --- i.e. how we judge whether a model satisfies $\Theta_\ThyPax$, or how we judge whether a predicate is valid in all models of $\Theta_\ThyPax$ --- is unpacked in Figure~\ref{fig.validity} and Definition~\ref{defn.theory.axiom}(\ref{item.axioms.valid}).
We discuss the quantifier scoping in those definitions in Remark~\ref{rmrk.thy.quantifier.scope}.
\end{enumerate*}
\end{rmrk}




\subsection{High-level discussion of the axioms}

\begin{rmrk}
The predicates in Figure~\ref{fig.logical.paxos} are intended as a \emph{theory}, in the sense of Definition~\ref{defn.theory.axiom}, which is why we call them \emph{axioms} and not just `some predicates'.

That is: we fix a signature $\Sigma$ that has (at least) the predicate symbols required to express the axioms in Figure~\ref{fig.logical.paxos}, and we fix a model $\mathcal M$, and we derive properties that hold in all the valuations $\avaluation$ to $\mathcal M$ such that the theory $\Theta_\ThyPax$ from Figure~\ref{fig.validity} is valid.
 
At least as much design effort went into crafting the axioms $\Theta_\ThyPax$ of Declarative Paxos in Figure~\ref{fig.logical.paxos}, as went into designing the three-valued modal fixedpoint logic in which the axioms are expressed.\footnote{Of course there was design feedback in both directions; we put into the logic what we needed to express the axioms, and we refined the axioms to elegantly work within the logic.} 
Therefore, we will take our time to discuss the axioms in detail.
The reader who just wants to see proofs can skip the rest of this Section (for now) and look to Section~\ref{sect.paxos.correctness.properties} onwards.
\end{rmrk}


\begin{rmrk}
\label{rmrk.axioms.paxos.three.groups}
We have organised the axioms in Figure~\ref{fig.logical.paxos} into three groups, which we now discuss:
\begin{enumerate}
\item
\emph{Backward (inference) rules.}

Axioms with names ending with $?$ represent backwards reasoning and so we call them \deffont{backward rules}.
They use the strong implication $\timpc$ (Definition~\ref{defn.three.imp}(\ref{item.three.elementary.strong.implication})), which means `if the LHS is true ($\tvT$) then the RHS is true ($\tvT$)'.

The backward rules represent a strong inference that \emph{if} some $Y$ is observed to be $\tvT$ then it \emph{must} be because some other property $X$ is $\tvT$.
For instance, \rulefont{PaxAccept?} says that if $\tf{accept}(v)$ is true, then $\someoneAll\tf{write}(v)$ must also be true.

We discuss the backward rules individually in Remark~\ref{rmrk.discuss.backward.rules} below.
\item
\emph{Forward (progress) rules.}

Axioms with names ending with $!$ represent forwards reasoning about making progress, so we call them \deffont{forward rules}.
They use $\final$ and the weak implication $\tnotor$ (Definition~\ref{defn.three.imp}(\ref{item.three.elementary.weak.implication})), which means `if the LHS is true ($\tvT$) then the RHS is valid ($\tvT$ or $\tvB$)'.

We derive the forward rules from the backward rules using a systematic recipe which we discuss in Remark~\ref{rmrk.recipe}.

The forward rules represent a weak inference that \emph{if} some precondition $X$ is $\tvT$ then --- provided the network is stable and provided that the participant has not crashed (where `crashed' is interpreted as returning $\tvB$) --- then the postcondition $Y$ is $\tvT$. 

For instance, \rulefont{PaxAccept!} expresses that, after network synchrony (Remark~\ref{rmrk.functional.notation.in.paxos}(\ref{item.discuss.final})), if $\texi\someoneAll\tf{write}$ is true (by Notation~\ref{nttn.hos} this is shorthand for $\texi a.\someoneAll\tf{write}(a)$) then if the participant has not crashed then $\texi\tf{accept}$ (i.e. $\texi a.\tf{accept}(a)$) is true. 

We discuss the forward rules in Remark~\ref{rmrk.discuss.forward.rules} below.
\item
\emph{Other rules, for correctness and liveness.}

Other rules reflect particular conditions of the algorithms.
We discuss them in Remark~\ref{rmrk.paxos.axioms.discussion}.
\end{enumerate}
\end{rmrk}

\begin{rmrk}[Recipe for deriving forward rules from backward rules]
\label{rmrk.recipe}
A design feature of Figure~\ref{fig.logical.paxos} is that each forward rule is obtained from the corresponding backward rule by a \deffont{recipe}.
Namely: for a backward rule of the form 
$$
\tall v.\ \chi \timpc \phi[a\ssm v] \timpc \psi[a\ssm v]
$$
where 
\begin{itemize*}
\item
$\chi$ expresses some possible \emph{side-conditions} (e.g. $\chi=\tf{leader}$ in \rulefont{PaxDecideL?}) and 
\item
$\phi$ expresses some \emph{preconditions}, and 
\item
$\psi$ expresses some \emph{postconditions}, 
\end{itemize*}
our recipe constructs a forward rule of the form
$$
(\texi v.\,\chi \tand \phi)\tnotor \texi v.\psi .
$$
In most of the forward rules this result also get wrapped in a $\final$ modality, representing network synchrony  (Remark~\ref{rmrk.functional.notation.in.paxos}(\ref{item.discuss.final})).

This recipe holds for every pair of rules:
\begin{itemize*}
\item
We see it cleanly in the pairs \rulefont{PaxAccept?} / \rulefont{PaxAccept!}, \rulefont{PaxDecideL?} / \rulefont{PaxDecideL!}, and \rulefont{PaxDecide\tneg L?} / \rulefont{PaxDecide\tneg L!}.
These rules follow our recipe symbol-for-symbol.
\item
We see this also with \rulefont{PaxPropose?} / \rulefont{PaxPropose!}.
Our recipe generates the following version of \rulefont{PaxPropose!}: 
$$
(\texi v.\,\tf{leader}\tand v\tneq\tf{undef})\timpc\texi v.\tf{propose}(v).
$$
The reader can check that this is logically equivalent to what is written in Figure~\ref{fig.logical.paxos}.
\item 
It is less clear how our recipe works for \rulefont{PaxSend?} and \rulefont{PaxSend!} and \rulefont{PaxWrite?} and \rulefont{PaxWrite!}; e.g. \rulefont{PaxSend!} looks different from \rulefont{PaxSend?} --- it is clearly simpler!
Remarkably, these pairs too are consistent with our recipe; see Remark~\ref{rmrk.why.evaporates} and the subsequent discussion and results.
\end{itemize*}
\end{rmrk}

\begin{rmrk}
We make some further elementary comments on our recipe from Remark~\ref{rmrk.recipe}:
\begin{enumerate*}
\item
We uncurry the forward rules --- meaning that we write $(\texi v.\,\chi \tand \phi)\tnotor \texi v.\psi$ instead of $(\texi v.\,\chi) \tnotor \phi\tnotor \texi v.\psi$ --- because this is convenient for some proofs. 
It makes no other difference and either form would be correct.
\item
A forward rule of the form $\texi v.((\chi \tand \phi)\tnotor \psi)$ would be \emph{wrong}.
This is just a different predicate and is not what is intended.
\item
A forward rule of the form $\tall v.((\chi \tand \phi)\tnotor \psi)$ would not necessarily be wrong, but it is stronger than we need for our proofs.
This reflects a general design principle in logical theories that \emph{weaker axioms are better}.
If we can prove the same result from weaker axioms then generally speaking we should, because that is a stronger result.\footnote{In practice, working from weaker axioms tends to force us to think harder and so design a better-engineered proof.}
\end{enumerate*}
\end{rmrk}


\subsection{More detailed discussion of the axioms}


\begin{rmrk}
\label{rmrk.paxos.axioms.discussion}
We discuss the \emph{axioms for correctness and liveness}:
\begin{enumerate}
\item\label{item.discuss.LdrExist}
\rulefont{LdrExist}:

In implementations of Paxos, at each stage there is a unique \emph{leader}, whose identity is computed by applying (what amounts to) a random number generator to an agreed seed.
As the level of abstraction in this paper, we just assume that a leader exists --- we do not care how it is chosen.

We use a predicate constant symbol $\tf{leader}$ to identify the leader.
Because our logic is modal --- across time $n\in\Time$ and participants $p\in\Pnt$ --- we can just assume that a constant symbol $\tf{leader}$ is assigned truth-value $\tvT$ where a participant $p$ that is a leader at a stage $n$, and is assigned truth-value $\tvF$ where a participant $p$ that is not a leader at a stage $n$.

In more technical language: $\avaluation$ is such that $n,p\mentval \modTF\tf{leader}$ for every time $n\in\Time$ and participant $p\in\Pnt$, %
and $p$ is a leader %
precisely when $n,p\mentval\tf{leader}$.
\rulefont{LdrExist} just asserts that a participant either is a leader or it is not, and some leader exists.
\item\label{item.discuss.LdrExistUniq}
\rulefont{LdrExt}:

\rulefont{LdrExt} is a rule-scheme (one for each closed $\phi$) and asserts that leaders are extensionally equal; %
a formal statement and proof are in Lemma~\ref{lemm.LeaderExtUniq}.

In an implementation this is enforced by having a unique leader, but in the logic this is unnecessary --- we don't care how many participants claim to be leaders, so long as they are all saying the same thing! --- and also it is impossible because (although we have modalities $\everyone$ and $\someone$) we do not have an explicit quantification over elements of $\Pnt$ in our logic.
This is a feature, not a bug: we want our logic to be no more expressive than it needs to be.

In fact, the axiom-scheme \rulefont{LdrExt} is stronger than we need for our proofs;
more on this in Subsection~\ref{subsect.ldrexistuniq'}.


Henceforth we may say `the leader' instead of `a leader', because up to extensional equivalence there is just one.
\item
\rulefont{LdrCorrect}:

\rulefont{LdrCorrect} asserts that for infinitely many times $n\in\Time$, the leader $l$ (which exists by the right-hand conjunct in \rulefont{LdrExist}) is $\tf{propose}$-correct, $\tf{write}$-correct, and $\tf{decide}$-correct (terminology from Definition~\ref{defn.correct}(\ref{item.P-correct}).

Unpacking Definition~\ref{defn.correct}(\ref{item.P-correct}), this means that $n,l\mentval \modTF\tf{propose}(v)$ for every $v\in\Val$, and similarly for $\tf{write}$ and $\tf{decide}$.
Intuitively, this is a \emph{liveness assumption} that there is always a future time when the leader is not crashed in the sense that it can correctly $\tvT\tvF$-propose, $\tvT\tvF$-write, and $\tvT\tvF$-decide a value. 
\item
\rulefont{PaxPCorrect}:

Given a valuation $\avaluation$, we unpack Definition~\ref{defn.validity.judgement}(\ref{item.mentval}) to see what $\mentval \rulefont{PaxPCorrect}$ actually means, as per Definition~\ref{defn.theory.axiom}:
\begin{multline*}
\Forall{n\in\Time}\Exists{O\in\opens}\Forall{p\in O}\Forall{v\in\Val}n,p,O\mentval\correct{\tf P(v)}
\\
\text{for each}\ \tf P\in\{\tf{propose},\tf{send},\tf{accept},\tf{write},\tf{decide}\} .
\end{multline*}
This corresponds to a standard correctness assumption in Paxos that at each stage there is a quorum of uncrashed participants.

However, the axiom \rulefont{PaxPCorrect} is subtly more general than the standard correctness assumption, because there is no restriction that the \emph{same} open set $O$ be used for each unary predicate symbol.\footnote{If we wanted the same $O$, then we would scope the $\QuorumBox$ modality further out and write $\QuorumBox(\correct{\tf{propose}} \tand \correct{\tf{send}} \tand \correct{\tf{accept}} \tand 
\correct{\tf{write}} \tand \correct{\tf{decide}})$.  This turns out to be unnecessary.}
It turns out that the proofs require only this assumption, and no more. 
\item\label{item.01.discussion}
\rulefont{PaxPropose01} and \rulefont{PaxWrite01}:

We justify these axioms with respect to the Paxos algorithm as summarised in Remark~\ref{rmrk.high-level.paxos}:
\begin{itemize*}
\item
Step~\ref{item.l.proposes} of Remark~\ref{rmrk.high-level.paxos} states that leader is crashed or proposes one value.
The intuitive correspondence to \rulefont{PaxPropose01} is clear.
\item
Step~\ref{item.l.writes} of Remark~\ref{rmrk.high-level.paxos} states that the leader writes a value derived from \emph{the first} quorum of participants that it hears back from (or none, if it hears back from no such quorum). 
One corollary of this, is that it writes at most one value.
Our logical theory has no state machine and notion of algorithmic time (as per Remark~\ref{rmrk.elisions}(\ref{item.elide.algorithmic.time})) --- so there is no sense in which information from one open set `arrives' at the leader `before' information from another --- so from the point of view of the logic, of all the open sets available, the leader picks at most one to act on.
This is axiom \rulefont{PaxWrite01}. 
\end{itemize*}
It is a fact that \rulefont{PaxPropose01}+\rulefont{PaxPropose!} and \rulefont{PaxWrite01}+\rulefont{PaxWrite!} give us
$$
\tf{leader}\tnotor\texiunique\tf{propose}
\quad\text{and}\quad
\final((\tf{leader}\tand\QuorumBox\texi\tf{send})\tnotor \texiunique\tf{write}) ,
$$ 
and
\rulefont{PaxPropose01}+\rulefont{PaxPropose?} and \rulefont{PaxWrite01}+\rulefont{PaxWrite?} give us
$$
\tneg\tf{leader}\tnotor\tneg\texi\tf{propose}
\quad\text{and}\quad
\tneg\tf{leader}\tnotor\tneg\texi\tf{write} .
$$
So it follows from our axioms that a leader participant proposes and writes one value (assuming it is not crashed, and if messages arrive), and a non-leader participant proposes and writes zero values. 
\item
\rulefont{Pax2Twined}:

This axiom collects a set of instances of Lemma~\ref{lemm.supertwined.qcbt}(\ref{item.intertwined.3}), namely taking $\phi=\tf{accept}(v)$ in that Lemma for every $v\in\Val$.

There is a design choice here: instead of an axiom \rulefont{Pax2Twined}, we could insist specifically that the semitopology in our underlying model be $2$-twined.
Then \rulefont{Pax2Twined} ceases to be an axiom, but we still have it as a fact of the model by directly invoking Lemma~\ref{lemm.supertwined.qcbt}(\ref{item.intertwined.3}).
Or we could be even more specific: we could insist that the underlying semitopology be $\AllBut{N}{f}$ (Definition~\ref{defn.allbut}) for $N> 2*f$; by Proposition~\ref{prop.allbut.supertwined}(\ref{item.allbut.intertwined}) this is $2$-twined, and then we can use Lemma~\ref{lemm.supertwined.qcbt}(\ref{item.intertwined.3}) to obtain \rulefont{Pax2Twined}.

We choose here to express as an axiom precisely the property that our proofs require.
This is concise, general, and yields clean proofs.
Insisting on a $2$-twined semitopology, or specifically on the semitopology $\AllBut{N}{f}$ for $N>2*f$, would also make sense, but in the context of our logical approach, using an axiom seems nice. %
\end{enumerate}
\end{rmrk}

\begin{rmrk}
\label{rmrk.discuss.backward.rules}
We discuss the \emph{axioms for the backward rules}:
\begin{enumerate}
\item
\rulefont{PaxPropose?}:

If we propose some value $v$ then 
\begin{enumerate}
\item
we are the leader and 
\item
$v$ is not the undefined value $\udfn$.
\end{enumerate}
\item
\rulefont{PaxSend?}:

If we send some value $v$ then 
\begin{enumerate}
\item
somebody must have proposed some value and 
\item
\begin{enumerate*}
\item
\emph{either} $v$ is the value we have most recently accepted, 
\item
\emph{or} we never accepted any value and $v$ is the undefined value.
\end{enumerate*}
\end{enumerate}
\item
\rulefont{PaxWrite?}:

If we write some value $v$ then 
\begin{enumerate}
\item
we are a leader, and
\item
there exists an open set of points (= participants) that have all sent us values, and
\item
\begin{enumerate*}
\item
\emph{either} some point in that open set has most recently accepted $v$, 
\item
\emph{or} 
$v$ is the value that we proposed, and no point in that open set has ever accepted any value. 
\end{enumerate*}
\end{enumerate}
\item
\rulefont{PaxAccept?}:

If we accept $v$ then somebody somewhere must have written $v$.
\item
\rulefont{PaxDecideL?}:

If a leader decides $v$ then an open set of points must have all accepted $v$.
\item
\rulefont{PaxDecide{\tneg}L?}:

If we are not a leader and we decide $v$, then some leader somewhere must have decided $v$. 
\end{enumerate}
\end{rmrk}

\begin{rmrk}
\label{rmrk.discuss.forward.rules}
We discuss the \emph{axioms for the forward rules}.

We noted in Remark~\ref{rmrk.recipe} our design methodology of deriving the forward rules from the backward rules; forward rules are rules for making progress, so they are existentially quantified over the variable $v$ and they use the weak implication $\tnotor$ to emulate possible crashes.

The other subtlety is the $\final$ modality wrapping all of the forward rules, except for \rulefont{PaxPropose!}.
\begin{itemize*}
\item
\emph{What $\final$ does for us in most of the forward rules.}

It models possible network failure, and eventual network synchrony.
As touched on in Remark~\ref{rmrk.functional.notation.in.paxos}(\ref{item.discuss.final})), the $\final$ modality reflects the possibility that the network might be down for a while, but eventually --- `finally' --- the network stabilises and messages arrive.

Predicates in our logic are \emph{not} messages, but we get the effect of eventual network synchrony in our logic (even though messages are not explicitly present) by wrapping the implications in our forward rules in the $\final$ modality.

More on this when we introduce the notion of \emph{global stabilisation of logical time} in Subsection~\ref{subsect.gslt}.
\item
\emph{Why $\final$ is absent from \rulefont{PaxPropose!}.}

Because this rule models a property of the implementation that involves no network communication.
The leader does not need to communicate with any other participant to pick a value to propose.

It is true that the leader's propose message might not arrive if the network is down --- but this is reflected by the use of $\final$ in the \rulefont{PaxSend!} rule.
\end{itemize*}
\end{rmrk}


\subsection{Some asides}

\subsubsection{An aside on \rulefont{PaxWrite?}}

\rulefont{PaxWrite?} contains this predicate:
$$
\mru{(\someone\tf{accept})}{v} .
$$
Unpacking what this means (cf. Remark~\ref{rmrk.explain.mrup}), this looks into the past to check whether $v$ is a value that was most recently accepted by some participant.

Note that this slightly differently-scoped predicate would be \emph{wrong}:
$$
\someone(\mru{\tf{accept}}{v}) .
$$
This finds some participant such that $v$ is a value that was most recently accepted by \emph{that} participant (but there might be another participant who more recently accepted some other value).

The reader with a background in logic will know all about how important scoping can be (e.g. $\forall x.\exists y.\phi$ vs. $\exists x.\forall y.\phi$, and $\exists x.(\phi\land \phi')$ vs. $\exists x.\phi\land\exists x.\phi'$), but perhaps the point bears repeating.


\subsubsection{An aside on \rulefont{LdrExt}}

In a Paxos implementation, there is a unique leader at each stage. 
We mentioned in Remark~\ref{rmrk.paxos.axioms.discussion}(\ref{item.discuss.LdrExist}\&\ref{item.discuss.LdrExistUniq}) that \rulefont{LdrExist} asserts that a leader exists and that the \rulefont{LdrExt} axiom-scheme asserts that all leader participants are extensionally equal.
We now make this formal:\footnote{Later on in Subsection~\ref{subsect.ldrexistuniq'} we will note weaker, but more abstract, forms of the axiom.}
\begin{lemm}[Leaders are extensionally equivalent]
\label{lemm.LeaderExtUniq}
Suppose $n\in\Time$ and $\avaluation$ is a valuation, and suppose $n\mentval\rulefont{LdrExt}_\phi$ holds for every closed predicate $\phi$.\footnote{By Definition~\ref{defn.validity.judgement}(\ref{item.ment.n}) this means $\Forall{p\in\Pnt}\Forall{O\in\opensne}(\modellabel{\rulefont{LdrExt}_\phi}{n,p,O,\avaluation}\in\{\tvT,\tvB\})$.} 
Suppose that $l,l'\in\Pnt$ and $n,l\mentval\modT\tf{leader}$ and $n,l'\mentval\modT\tf{leader}$.

Then $l$ and $l'$ are \deffont{extensionally equivalent} at $(n,\avaluation)$, by which we mean that for every closed predicate $\phi$ and every $O\in\opensne$,
$$
\modellabel{\phi}{n,l,O,\avaluation} = \modellabel{\phi}{n,l',O,\avaluation} .
$$
\end{lemm}
\begin{proof}
Suppose $\phi$ is a closed predicate and $O\in\opensne$, and suppose for the sake of argument that $\modellabel{\phi}{n,l,O,\avaluation}=\tvB$.

It is routine to check from Figures~\ref{fig.3.derived} and~\ref{fig.3} and from Lemma~\ref{lemm.dense.char}(\ref{item.dense.char.someoneAll}) that $\modellabel{\someoneAll\modB\phi}{n,l,O,\avaluation}=\tvT$. 
By assumption $n,l',O\mentval \rulefont{LdrExt}_{\modB\phi}$, and by two applications of \strongmodusponens (once for $\modT\tf{leader}$, and another for $\modT\someoneAll(\tf{leader}\tand\modB\phi)$) we conclude that $\modellabel{\modB\phi}{n,l',O,\avaluation}=\tvT$, and therefore $\modellabel{\phi}{n,l',O,\avaluation}=\tvB$ as required. 
 
The cases $\modellabel{\phi}{n,l,O,\avaluation}=\tvT$ and $\modellabel{\phi}{n,l,O,\avaluation}=\tvF$ are be precisely similar, and this suffices to prove the result.
\end{proof}


\subsubsection{An aside on \rulefont{PaxSend?} and \rulefont{PaxSend!}}

\begin{rmrk}
\label{rmrk.why.evaporates}
In Remark~\ref{rmrk.recipe} we proposed a recipe for deriving a forward rule from a backward one: a backward rule of the form
$$
\rulefont{Rule?} 
\qquad
\tall v.\ \chi \timpc \phi[a\ssm v] \timpc \psi[a\ssm v]
$$
generates a forward rule (possibly wrapped in a $\final$ modality) of the form
$$
\rulefont{Rule!}
\qquad
(\texi v.\, \chi \tand \phi)\tnotor \texi v.\psi .
$$
We noted in Remark~\ref{rmrk.recipe} that the recipe works cleanly for some rules --- but it was less clear how it can be applied to derive \rulefont{PaxSend!} from \rulefont{PaxSend?}, and \rulefont{PaxWrite!} from \rulefont{PaxWrite?}.
In these cases, the left-hand sides of \rulefont{PaxSend!} and \rulefont{PaxWrite!} seem much simpler than the right-hand sides of \rulefont{PaxSend?} and \rulefont{PaxWrite?} respectively. 

We now show why, in fact, \rulefont{PaxSend!} and \rulefont{PaxWrite!} are correct and make sense in the context of our recipe, culminating with Corollary~\ref{corr.back.ok}.
\end{rmrk}


\begin{lemm}
\label{lemm.evaporates.simplify}
We have:
\begin{enumerate*}
\item
$\ment \texi v.\bigl((\mru{\tf{accept}}{v})\ \tor (\tneg\modT\recent\texi\tf{accept}\tand v\teq\udfn)\bigr) \ \equiv\ \tvT$.
\item
$\ment \texi v.\bigl(\ (\mru{(\someone\tf{accept})}{v})\ \tor\ (\tneg\modT\recent\texi\someone\tf{accept}\tand \tf{propose}(v))\ \bigr)\ \equiv\ \texi\tf{propose}$.
\end{enumerate*}
Note that:
\begin{itemize*}
\item
The notation $\texi\tf{accept}$ and $\texi\tf{propose}$ and $\texi\someone\tf{accept}$ is from Notation~\ref{nttn.hos}.
\item
We do not restrict to contexts and valuations that satisfy \ThyPax; the meaning of the $\ment$ symbol above is as in Definition~\ref{defn.validity.judgement}(\ref{item.ment.phi}).
\end{itemize*}
\end{lemm}
\begin{proof}
We consider each part in turn:
\begin{enumerate}
\item
We distribute the $\texi v$ through the $\tor$, and then through the $\tand$ on the right-hand side (we can do this because $v$ is not free in the left-hand conjunct) and we see that it suffices to show 
$$
\ment \texi v.(\mru{\tf{accept}}{v})\ \tor\ (\tneg\modT\recent\texi\tf{accept}\tand \texi v.v\teq\udfn)\ \equiv\ \tvT
$$
Using Lemma~\ref{lemm.evaporates} this simplifies to just
$$
\ment \modT\recent\texi\tf{accept}\ \tor\ (\tneg\modT\recent\texi\tf{accept}\tand \texi v.v\teq\udfn) \ \equiv\ \tvT.
$$
This clearly holds, just taking $v\teq\udfn$.
\item
We distribute $\texi v$ into the predicate as we did for the proof of part~1.
We see that it suffices to show 
$$
\ment \bigl(\ \texi v.(\mru{(\someone\tf{accept})}{v})\ \tor\ (\tneg\modT\recent\texi\someone\tf{accept}\tand \texi\tf{propose})\ \bigr)\ \equiv\ \texi\tf{propose}.
$$
We simplify using Lemma~\ref{lemm.evaporates} to
$\ment \texi\tf{propose}\ \equiv\ \texi\tf{propose}$,
which is true.\footnote{Yes, we have a lemma for that: Lemma~\ref{lemm.equiv}(\ref{item.equiv.2}) and reflexivity of equality.}
\qedhere
\end{enumerate}
\end{proof}

We can now justify the left-hand sides of \rulefont{PaxSend!} and \rulefont{PaxWrite!} as being consistent with our recipe as discussed above in Remark~\ref{rmrk.why.evaporates}:
\begin{corr}
\label{corr.back.ok}
We have:
$$
\begin{array}{r@{\ }l}
\ment\texi v.&\bigl(\someoneAll\texi\tf{propose}\ \tand \bigl((\mru{\tf{accept}}{v})\ \tor (\tneg\modT\recent\texi\tf{accept}\tand v\teq\udfn)\bigr)\bigr) 
\\
&\equiv \someoneAll\texi\tf{propose}
\\[2ex]
\ment\texi v.&\bigl(\tf{leader}\tand 
\Quorum(\Box\texi\tf{send} \tand 
\bigl((\mru{(\someone\tf{accept})}{v}) \tor
(\tneg\modT\recent\texi\someone\tf{accept}\tand \tf{propose}(v))\bigr)) 
\bigr)
\\
&\equiv
(\tf{leader}\tand \QuorumBox\texi\tf{send} \tand \texi\tf{propose})
\end{array}
$$
\end{corr}
\begin{proof}
By elementary predicate manipulations, distributing the $\texi v$ quantifier down into the predicate and using Lemma~\ref{lemm.evaporates.simplify}.
\end{proof}

 




\section{Correctness properties}
\label{sect.paxos.correctness.properties}

\subsection{Summary of the properties}

\begin{rmrk}
\label{rmrk.standard.paxos.correctness}
We will derive the following standard correctness properties for Paxos (as laid out in~\cite[Subsection~2.2.1]{cachin2011paxos}) from the theory \ThyPax in Figure~\ref{fig.logical.paxos}:
\begin{enumerate*}
\item
\textbf{Validity:} \quad If a participant decides some value $v$, then $v$ was proposed by some participant.

A formal statement and proof is in Theorem~\ref{thrm.validity}.
\item
\textbf{Agreement:}\quad 
If two participants decide values, then those values are equal.

A formal statement and proof is in Theorem~\ref{thrm.full.agreement}.
\item
\textbf{Termination:}\quad 
Every correct participant eventually decides some value.

A formal statement and proof is in Theorem~\ref{thrm.termination}.
\item
\textbf{Integrity:}\quad 
Every correct (non-crashing) participant decides at most once.

Our logic abstracts away from a direct notion of execution, by design, but this property is represented indirectly in how Theorem~\ref{thrm.termination} is expressed.
See Remark~\ref{rmrk.paxos.integrity} below. 
\end{enumerate*}
\end{rmrk}

\begin{rmrk}[On integrity]
\label{rmrk.paxos.integrity}
In implementations of Paxos, once a participant has decided, it stops.
But in our logic, `stop' is not even a concept.
There is no abstract machine to stop; there are just predicates, with denotations.

Nevertheless, the integrity property is echoed in our maths, as the \emph{infinite repetition} in Theorem~\ref{thrm.termination}.
Once a predicate $\tf{decide}(v)$ returns $\tvT$, it does so \emph{forever} --- and because of Agreement, it never decides any other value.
A predicate eventually always returning the same truth-value is a model of an abstract machine that has computed its result, and then just sits there forever, displaying that same result, unchanged, every time we query its display.
\end{rmrk}

\subsection{Validity}


\begin{rmrk}
\label{rmrk.validity.in.english}
\emph{Validity} is expressed in English as:
\begin{quote}
\emph{If a participant decides some value $v$, then $v$ was proposed by some participant.}
\end{quote}
We render this in our logic in Theorem~\ref{thrm.full.agreement} as
$$
\ThyPax\ment\tall v.(\tf{decide}(v) \timpc \urecent\someoneAll(\tf{leader}\tand\tf{propose}(v))).
$$

In words we can read this as: 
\begin{quote}
If someone decides $v$ now, then at some stage now or in the past a leader of that stage proposed $v$. 
\end{quote}
We will now set about proving this property from our axioms, culminating with Theorem~\ref{thrm.validity}.
\end{rmrk}


\begin{lemm}
\label{lemm.validity.helper}
Suppose $\avaluation$ is a valuation such that $\mentval\ThyPax$ ($\mentval\ThyPax$ is from Figure~\ref{fig.validity}).
Suppose $n\in\Time$ and $v\in\Val$.
Then:
\begin{enumerate*}
\item\label{item.validity.helper.3}
$n\mentval \someoneAll\tf{write}(v) \timpc \urecent\someoneAll\tf{propose}(v)$. 
\item\label{item.validity.helper.2}
$n\mentval \someoneAll\tf{accept}(v) \timpc \urecent\someoneAll\tf{propose}(v)$. 
\end{enumerate*}
\end{lemm}
\begin{proof}
We prove both parts by a simultaneous induction on $n$.
We freely use \strongmodusponens:
\begin{enumerate}
\item
Suppose $n\mentval\modT\someoneAll\tf{write}(v)$.
Then using \rulefont{PaxWrite?} %
$$
n\mentval \modT\Quorum(\mru{\someone\tf{accept}}{v}) 
\quad\text{or}\quad
n\mentval \modT\someoneAll\tf{propose}(v) .
$$
We consider the two cases:
\begin{itemize*}
\item
\emph{If $n\mentval \modT\someoneAll\tf{propose}(v)$,}\ then 
$n\mentval\modT\urecent\someoneAll\tf{propose}(v)$ follows immediately from the clause for %
$\urecent$ in Figure~\ref{fig.3.derived}.
\item
\emph{If $n\mentval \modT\Quorum(\mru{\someone\tf{accept}}{v})$,}\ 
then by Lemma~\ref{lemm.mru.to.recent}(\ref{item.mru.to.recent.1}) $n\mentval\modT\recent\someoneAll\tf{accept}(v)$,
so that by Figure~\ref{fig.3.derived} $n'\mentval\modT\someoneAll\tf{accept}(v)$ for some $n'<n$.
By part~\ref{item.validity.helper.2} of this result for $n'<n$, we have that $n'\mentval\modT\urecent\someoneAll\tf{propose}(v)$, and thus using Lemma~\ref{lemm.tomorrow.forever.recent}(\ref{item.rur.r}) $n\mentval\modT\urecent\someoneAll\tf{propose}(v)$. 
\end{itemize*}
In either case, we are done.
\item
Suppose $n\mentval\modT\someoneAll\tf{accept}(v)$.
By \rulefont{PaxAccept?} $n\mentval\modT\someoneAll\tf{write}(v)$.
We use part~\ref{item.validity.helper.3} of the inductive hypothesis for $n$. 
\qedhere\end{enumerate}
\end{proof}



\begin{thrm}[Validity]
\label{thrm.validity}
We have:
$$
\ThyPax\ment\tall v.(\tf{decide}(v) \timpc \urecent\someoneAll(\tf{leader}\tand\tf{propose}(v))) .
$$
\end{thrm}
\begin{proof}
Suppose $\avaluation$ is a valuation such that $\mentval\ThyPax$. 
We reason using \strongmodusponens. 
Suppose $v\in\Val$ and $n\in\Time$ and $p\in\Pnt$, and suppose $n,p\mentval\modT\tf{decide}(v)$.
By \rulefont{PaxDecideL?} (and \rulefont{PaxDecide\tneg L?} if $p\neq\f{leader}(n)$), $n\mentval\modT\QuorumBox\tf{accept}(v)$.
It follows by Proposition~\ref{prop.someone.implies.someoneAll}(\ref{item.someone.implies.someoneAll.2}) that $n\mentval\modT\someoneAll\tf{accept}(v)$.
By Lemma~\ref{lemm.validity.helper}(\ref{item.validity.helper.2}) $n\mentval\modT\urecent\someoneAll\tf{propose}(v)$, and using \rulefont{PaxPropose?} we conclude that $n\mentval\modT\urecent\someoneAll(\tf{leader}\tand\tf{propose}(v))$ as required. 
\end{proof}




\subsection{Agreement}

\begin{rmrk}
\emph{Agreement} is expressed in English as:
\begin{quote}
\emph{If two participants decide values, then those values are equal.}
\end{quote}
We render this in our logic in Theorem~\ref{thrm.full.agreement} as
$$
\ThyPax\ment\tall v,v'.(\urecent\someoneAll\tf{decide}(v) \timpc \someoneAll\tf{decide}(v') \timpc v'\teq v) .
$$
In words we can read this as: 
\begin{quote}
If someone decides $v$ now or in the past, and someone decides $v'$ now, then $v'=v$.
\end{quote}
It is convenient to prove the `now' and the `in the past' parts of this as two sub-cases: 
\begin{enumerate*}
\item
Proposition~\ref{prop.decide.n.agree} handles the case of $\someoneAll\tf{decide}(v)\tand\someoneAll\tf{decide}(v')$, i.e. when the two decisions are made at the same time.
\item
Theorem~\ref{thrm.full.agreement} then extends this with the case of $\recent\someoneAll\tf{decide}(v)\tand\someoneAll\tf{decide}(v')$, i.e. when one decision precedes the other.
\end{enumerate*}
\end{rmrk}


\subsubsection{Agreement at time $n$}



\begin{lemm}
\label{lemm.various.unique}
Suppose $\avaluation$ is a valuation such that $\mentval\ThyPax$. 
Suppose $n\in\Time$ and $v,v'\in\Val$. 
Then:
\begin{enumerate*}
\item\label{item.various.unique.write}
$n\mentval (\someoneAll\tf{write}(v)\tand \someoneAll\tf{write}(v')) \timpc v\teq v'$. 
\item\label{item.various.unique.accept}
$n\mentval (\someoneAll\tf{accept}(v)\tand \someoneAll\tf{accept}(v')) \timpc v\teq v'$. 
\item\label{item.various.unique.recent.accept}
$n\mentval \bigl(\mru{(\someoneAll\tf{accept})}{v}\tand \mru{(\someoneAll\tf{accept})}{v'}\bigr) \timpc v\teq v'$.
\end{enumerate*}
\end{lemm}
\begin{proof}
We consider each part in turn, using \strongmodusponens (we may also silently commute the $\modT$ modality as convenient):
\begin{enumerate}
\item
Suppose $n\mentval\modT\someoneAll\tf{write}(v)\tand\modT\someoneAll\tf{write}(v')$.
Using \rulefont{PaxWrite?} 
$n\mentval\modT\someoneAll(\tf{write}(v)\tand\tf{leader})\tand\modT\someoneAll(\tf{write}(v')\tand\tf{leader})$. 
Using \rulefont{LdrExt} therefore $n\mentval\modT\someoneAll(\tf{write}(v)\tand\tf{write}(v')\tand\tf{leader})$.
By \rulefont{PaxWrite01} and Lemma~\ref{lemm.unique.affine.existence}(\ref{item.unique.affine.existence.01.implies}) we have $v=v'$ as required.
\item
Suppose $n\mentval \modT\someoneAll\tf{accept}(v)\tand\modT\someoneAll\tf{accept}(v')$.
By \rulefont{PaxAccept?} $n\mentval\modT\someoneAll\tf{write}(v)\tand\modT\someoneAll\tf{write}(v')$.
We use part~\ref{item.various.unique.write} of this result.
\item
Suppose $n\mentval\modT(\mru{(\someoneAll\tf{accept})}{v})\tand\modT(\mru{(\someoneAll\tf{accept})}{v'}\bigr)$.\footnote{By Lemma~\ref{lemm.mru.unambivalent}(2) we do not need to write the $\modT$ here but we do anyway; see Remark~\ref{rmrk.explain.mru}(\ref{item.explain.mru.T}).}
Using Lemma~\ref{lemm.mru.to.recent}(\ref{item.mru.to.recent.2}) (taking $a.\phi=\someoneAll\tf{accept}$ in that Lemma) there exists $n'<n$ such that 
$n'\mentval\modT\someoneAll\tf{accept}(v)\tand\modT\someoneAll\tf{accept}(v')$.
We use part~\ref{item.various.unique.accept} of this result.
\qedhere\end{enumerate}
\end{proof}



\begin{prop}
\label{prop.decide.n.agree}
Suppose $\avaluation$ is a valuation such that $\mentval\ThyPax$. 
Suppose $n\in\Time$ and $v,v'\in\Val$.
Then
$$
n\mentval (\someoneAll\tf{decide}(v) \tand \someoneAll\tf{decide}(v')) \timpc v\teq v' . 
$$
\end{prop}
\begin{proof}
We use \strongmodusponens.
Suppose 
$n\mentval \modT(\someoneAll\tf{decide}(v)\tand\someoneAll\tf{decide}(v'))$.
Using \rulefont{PaxDecide\tneg L?} if needed, and then \rulefont{PaxDecideL?} twice, we have
$n\mentval \modT(\QuorumBox\tf{accept}(v)\tand\QuorumBox\tf{accept}(v'))$ and thus by Proposition~\ref{prop.someone.implies.someoneAll}(\ref{item.someone.implies.someoneAll.2}) and Lemma~\ref{lemm.commuting.connectives}
also $n\mentval \modT(\someoneAll\tf{accept}(v)\tand\someoneAll\tf{accept}(v'))$.
We can now use Lemma~\ref{lemm.various.unique}(\ref{item.various.unique.accept}).
\end{proof}



\begin{rmrk}
Proposition~\ref{prop.decide.n.agree} proves part of the agreement property; that no two correct participants decide differently \emph{at a given time}.
To prove the result across different times will require a bit more work, culminating with Theorem~\ref{thrm.full.agreement}.
\end{rmrk}


\subsubsection{Some fine structure of the rules}
\label{subsect.ldrexistuniq'}

We continue a discussion from Remark~\ref{rmrk.paxos.axioms.discussion}(\ref{item.discuss.LdrExistUniq}) and also Lemma~\ref{lemm.LeaderExtUniq}:

Note that we use the axiom-scheme \rulefont{LdrExt} from Figure~\ref{fig.logical.paxos} just once: in the proof of Lemma~\ref{lemm.various.unique}(\ref{item.various.unique.write}).

Examining how we use this axiom-scheme in the proof above, we see that for our specific purposes it would suffice to assume a single axiom as follows:
$$
\hspace{-6em}\rulefont{LdrExt'}
\qquad\qquad
\tf{leader}\timpc \someoneAll(\tf{leader}\tand\tf{write}(v))\timpc \tf{write}(v) .
$$ 
The axiom scheme \rulefont{LdrExt} asserts that there is a unique leader at each stage, up to extensional equivalence, which is a fair reflection of how the algorithms work where there \emph{is} a unique leader at each stage.

\rulefont{LdrExt'} asserts that there is a unique leader at each stage, up to extensional equivalence \emph{for writing values}.
In the proofs, this turns out to be all we need.

In fact, we can simplify this just to 
$$
\hspace{-6em}\rulefont{SPaxWrite01}
\qquad\qquad
\texiaffine\someoneAll\tf{write} .
$$
In full, this is $\texiaffine v.\someoneAll\tf{write}(v)$, and in words it expresses: 
\begin{quote}
`at each stage, at most one value is written'.
\end{quote}
In the presence of \rulefont{PaxWrite?}, \rulefont{SPaxWrite01} is implied by \rulefont{LdrExt'} --- but it does not imply \rulefont{LdrExt'}, because it permits the possibility that one leader participant writes a value and another leader participant crashes.
The reader can check that \rulefont{SPaxWrite01} is sufficient to complete the proof of Lemma~\ref{lemm.various.unique}(\ref{item.various.unique.write}).

Note also that that \rulefont{SPaxWrite01} implies \rulefont{PaxWrite01}, so that --- for the proofs we care about --- \rulefont{SPaxWrite01} can replace both \rulefont{LdrExt} and \rulefont{PaxWrite01}. 
We will do exactly this, later: the job of $\Theta_\ThyPax$ is to be `declarative Paxos', that is, its purpose is to be declarative \emph{and} to be recognisably Paxos. 
Later on in Figure~\ref{fig.ThySPax} we will define a \emph{simpler} declarative Paxos in (whence the name `SPaxWrite01').
This simplified theory is derived by looking at the proofs here, and seeing what we actually require.










 



\subsubsection{Agreement at possibly different times}



\begin{prop}
\label{prop.integrity.helper}
Suppose $\avaluation$ is a valuation such that $\mentval\ThyPax$, and suppose $v,v'\in\Val$ and $\avaluation$. 
Then:
\begin{enumerate*}
\item\label{item.integrity.helper.accept}
$\mentval\recent\QuorumBox\tf{accept}(v) \timpc \someoneAll\tf{accept}(v')\timpc v'\teq v$.
\item\label{item.integrity.helper.write}
$\mentval\recent\QuorumBox\tf{accept}(v) \timpc \someoneAll\tf{write}(v')\timpc v'\teq v$.
\end{enumerate*}
\end{prop}
\begin{proof}
Suppose $n\in\Time$.
We will reason using \strongmodusponens, so suppose 
\begin{equation}
\label{eq.n.1}
n\mentval\modT\QuorumBox\tf{accept}(v). 
\end{equation}
We will prove by induction on $n'>n$ that for every $v'\in\Val$,
$$
\begin{array}{r@{\ }l}
n'\mentval&\modT(\mru{(\someoneAll\tf{accept})}{v'}\timpc v'\teq v)
\quad\text{and}\quad
\\
n'\mentval&\modT(\someoneAll\tf{write}(v')\timpc v'\teq v)
\end{array}
$$
Clearly, by \strongmodusponens the $\tf{write}$ part of the inductive hypothesis is enough to prove part~\ref{item.integrity.helper.write} of this result.
For the $\tf{accept}$ part of the inductive hypothesis, we see that this implies part~\ref{item.integrity.helper.accept} of this result either by a routine argument using \rulefont{PaxAccept?} and part~\ref{item.integrity.helper.write} of this result --- or using Corollary~\ref{corr.mru.today.to.yesterday}(\ref{item.mru.today.implies.tomorrow}) (for $a.\phi=a.\someoneAll\tf{accept}(a)$ in that result).

We consider cases:
\begin{itemize}
\item
\emph{The base case for $\tf{accept}$, where $n'=n+1$.}

Suppose $n'\mentval\modT(\mru{\someoneAll\tf{accept}}{v'})$.

We assumed in equation~\eqref{eq.n.1} that $n\mentval\modT\QuorumBox\tf{accept}(v)$, so by Proposition~\ref{prop.someone.implies.someoneAll}(\ref{item.someone.implies.someoneAll.2}) also $n\mentval\modT\someoneAll\tf{accept}(v)$, and by Corollary~\ref{corr.mru.today.to.yesterday}(\ref{item.mru.today.implies.tomorrow}) $n'\mentval\modT(\mru{(\someoneAll\tf{accept})}{v})$. 
By Lemma~\ref{lemm.various.unique}(\ref{item.various.unique.recent.accept}) $v'=v$.

\item
\emph{The inductive step for $\tf{accept}$, where $n'>n+1$.}

Suppose $n'\mentval\modT(\mru{\someoneAll\tf{accept}}{v'})$.
By Corollary~\ref{corr.mru.today.to.yesterday}(\ref{item.mru.tomorrow.options}) either $n'\minus 1\mentval\modT\someoneAll\tf{accept}(v')$ or $n'\minus 1\mentval\modT(\mru{\someoneAll\tf{accept}}{v'})$.

In the former case, using \rulefont{PaxAccept?} $n'\minus 1\mentval\someoneAll\tf{write}(v')$, and by the $\tf{write}$ part of the inductive hypothesis, $v'=v$.
In the latter case, we just use the $\tf{accept}$ part of the inductive hypothesis. 
\item
\emph{The case of $\tf{write}$.}

Suppose $p\in\Pnt$ and $n',p\mentval\modT\tf{write}(v')$.
Using \rulefont{PaxWrite?}, 
there exists an $O\in\opensne$ such that 
$$
n',O\mentval \modT\bigl( 
(\mru{\someone\tf{accept}}{v})\ \tor
(\tneg\modT\recent\texi\someone\tf{accept}\tand \tf{propose}(v)) \bigr),
$$
from which it follows using Lemmas~\ref{lemm.commuting.connectives} and~\ref{lemm.tand.tor} and some basic reasoning on truth-values that 
$$
\text{either}\quad
n',O\mentval \modT\mru{\someone\tf{accept}}{v}
\quad\text{or}\quad
n',O\mentval \tneg\modT\recent\texi\someone\tf{accept} .
$$
Now $n<n'$ and we assumed in equation~\eqref{eq.n.1} that $n\mentval\modT\QuorumBox\tf{accept}(v)$, so using 
\rulefont{Pax2Twined} and \strongmodusponens $n\mentval\modT\CoquorumDiamond\tf{accept}(v)$, and in particular (by routine reasoning on the clauses for $\CoquorumDiamond$ in Figure~\ref{fig.3.derived} and for $\someone$ in Figure~\ref{fig.3.phi.f})
$n,O\mentval\modT\someone\tf{accept}(v)$.
Thus $n',O\mentval\modT\texi\someone\tf{accept}$ and by Lemma~\ref{lemm.tomorrow.forever.recent}(\ref{item.tomorrow.forever.recent.1}) (taking $\phi=\texi\someone\tf{accept}$ in that Lemma) 
$n',O\mentval\modT\recent\texi\someone\tf{accept}$.
It follows that we are not in the right-hand (`or') branch above, so we must be in the left-hand (`either') branch, and thus $n',O\mentval\modT\mru{\someone\tf{accept}}{v}$.

By the $\tf{accept}$ part of the inductive hypothesis, $v'=v$.
\qedhere\end{itemize}

\end{proof}


\begin{thrm}[Agreement]
\label{thrm.full.agreement}
We have:
$$
\ThyPax\ment\tall v,v'.(\urecent\someoneAll\tf{decide}(v) \timpc \someoneAll\tf{decide}(v') \timpc v'\teq v).
$$
\end{thrm}
\begin{proof}
Suppose $v,v'\in\Val$ and $n\in\Time$ and $\avaluation$ is a valuation.
We use \strongmodusponens, so suppose
$$
n\mentval\modT\urecent\someoneAll\tf{decide}(v)
\quad\text{and}\quad
n\mentval\modT\someoneAll\tf{decide}(v') .
$$
If $n\mentval\modT\someoneAll\tf{decide}(v)$ then we use Proposition~\ref{prop.decide.n.agree}.

So suppose $n\mentval\modT\recent\someoneAll\tf{decide}(v)$.
By \rulefont{PaxDecideL?} (and \rulefont{PaxDecide\tneg L?} if needed), $n\mentval\modT\recent\QuorumBox\tf{accept}(v)$ and $n\mentval\modT\QuorumBox\tf{accept}(v')$, so that by Proposition~\ref{prop.someone.implies.someoneAll}(\ref{item.someone.implies.someoneAll.2}) also %
$n\mentval\modT\someoneAll\tf{accept}(v')$.
By Proposition~\ref{prop.integrity.helper}(\ref{item.integrity.helper.accept}) $v'=v$ as required.
\end{proof}

\subsection{Termination}
\label{subsect.pax.termination}

\begin{rmrk}
\emph{Termination} is expressed in English as:
\begin{quote}
\emph{Every correct participant eventually decides some value.}
\end{quote}
We render this in our logic in Theorem~\ref{thrm.termination} as 
$$
\ThyPax\ment\infinitely\everyoneAll\texi v.\tf{decide}(v).
$$
In words: 
\begin{quote}
There is always some future time when all participants decide.
\end{quote}
Note that there is no need to explicitly specify `all correct/uncrashed participants' above.
This is taken care of by our logic's in-built notion of validity, since incorrect participants return $\tvB$, which is valid as per Remark~\ref{rmrk.motivate.validity}(\ref{item.tb.valid}) (but not true or correct; cf. Remark~\ref{rmrk.motivate.validity}).
\end{rmrk}


\subsubsection{Global stabilisation of logical time (GSLT)}
\label{subsect.gslt}

Our forward rules are wrapped in a $\final$ modality, to reflect an assumption that while there may be some initial network instability, if we wait long enough the network will stabilise (see the discussion in Remark~\ref{rmrk.discuss.forward.rules}).
This leads us to the notion of \emph{global stabilisation of logical time}, or \emph{GSLT} for short: 
\begin{defn}[GSLT]
\label{defn.gslt}
Suppose $\phi$ is a closed predicate and $\avaluation$ is a valuation.
Then:
\begin{enumerate*}
\item\label{item.gslt.synchronous}
Say that \ $\final\phi$\ (i.e. $\phi$ under a $\final$ modality) is \deffont{synchronous at $n$ in $\avaluation$} when $n\mentval\forever\phi$.\footnote{%
By Definition~\ref{defn.validity.judgement}, $n\ment\final\phi$ means that $n,p,O\mentval\final\phi$ for every $p\in\Pnt$ and $O\in\opens$ and valuation $\varsigma$.} %
That is:
$$
\final\phi\ \ \text{ is synchronous at $n$ in $\avaluation$}
\quad\text{when}\quad
n\mentval\forever\phi .
$$
In words, $\final\phi$ is synchronous when the `$\final$' becomes a `$\forever$'.
\item\label{item.gslt.gslt}
Define $\gslt(\avaluation)\in\Time$ to be the least stage $n$ such that all of the rules in Figure~\ref{fig.3.phi.f} that start with $\final$ --- namely, the forward rules \rulefont{PaxSend!}, \rulefont{PaxAccept!}, \rulefont{PaxDecideL!}, and \rulefont{PaxDecide\tneg L!} --- are synchronous at $n$ in the valuation $\avaluation$.

We call $\gslt(\avaluation)$ the \deffont{Global Stabilisation of Logical Time} in $\avaluation$. 
\end{enumerate*}
\end{defn}

\begin{rmrk}
We make some observations about Definition~\ref{defn.gslt}:
\begin{enumerate}
\item
Unpacking Definition~\ref{defn.gslt}(\ref{item.gslt.synchronous}) using the clause for $\forever$ in Figure~\ref{fig.3.derived}, 
$$
\final\phi \text{ is synchronous at $n$ in $\avaluation$}
\quad\liff\quad
n\mentval\forever\phi 
\quad\liff\quad
\Forall{n'{\geq}n}(n'\mentval\phi),
$$
\noindent
and $\gslt(\avaluation)$ is least such that $\gslt(\avaluation)\mentval\forever(\someoneAll\texi\tf{propose} \tnotor \texi\tf{send})$, and $\gslt(\avaluation)\mentval\forever\bigl((\tf{leader}\tand \QuorumBox\texi\tf{send}) \tnotor \texi\tf{write}\bigr)$, and so on.
\item  
$\gslt$ is well-defined in Definition~\ref{defn.gslt}(\ref{item.gslt.gslt}) because there are finitely many axioms, mentioning finitely many $\final$ modalities, and a finite set of natural numbers has a finite least upper bound.
This statement has real content, e.g. if we had written \rulefont{PaxAccept!} as $\tall v.\final(\someoneAll\tf{write}(v) \tnotor \tf{accept}(v))$ --- which is a subtly weaker assertion than the form used in Figure~\ref{fig.logical.paxos} --- then $\gslt$ would \emph{not} be defined. 
\end{enumerate}
\end{rmrk}




\begin{lemm}
\label{lemm.someone.writes.after.gslt.accept}
Suppose $\avaluation$ is a valuation such that $\mentval\ThyPax$. 
Suppose $n\in\Time$ and $n\geq\gslt(\avaluation)$ and $v\in\Val$. 
Then 
\begin{multline*}
n\mentval\modT\someoneAll\tf{write}(v) 
\liff
n\mentval\everyoneAll\tf{accept}(v) 
\liff
\\
n\mentval\modT\QuorumBox\tf{accept}(v) 
\liff
n\mentval\modT\someoneAll\tf{accept}(v)  .
\end{multline*}
\end{lemm}
\begin{proof}
We prove a cycle of implications, freely using weak and strong modus ponens (Proposition~\ref{prop.mp.for.tnotor}(\ref{item.mp.for.tnotor}\&\ref{item.mp.for.timpc})).

Suppose $n\mentval\modT\someoneAll\tf{write}(v)$. 
By \rulefont{PaxAccept!} (since $n\geq\gslt(\avaluation)$), $\Forall{p\in\Pnt}n\mentval\tf{accept}(v)$.
By Lemma~\ref{lemm.atomic.O}(\ref{item.atomic.pointwise}) $\tf{accept}(v)$ is pointwise, so by Proposition~\ref{prop.pointwise.dense.char}(\ref{item.dense.char.pointwise.everyoneAll}) $n\mentval\everyoneAll\texi\tf{accept}(v)$.


By Proposition~\ref{prop.correct.cobox.boxT}(\ref{item.correct.cobox.boxT.1}) and \rulefont{PaxPCorrect} (for $\tf{accept}$),
$n\mentval\modT\QuorumBox\tf{accept}(v)$, and by Proposition~\ref{prop.someone.implies.someoneAll}(\ref{item.someone.implies.someoneAll.2}) $n\mentval\modT\someoneAll\tf{accept}(v)$.

Now suppose $n\mentval\modT\someoneAll\tf{accept}(v)$.
By \rulefont{PaxAccept?} %
$n\mentval\modT\someoneAll\tf{write}(v)$.  
\end{proof}


\begin{lemm}
\label{lemm.leaderlive.quorumsend}
Suppose $\avaluation$ is a valuation such that $\mentval\ThyPax$. 
Suppose $n\in\Time$ and $n\geq\gslt(\avaluation)$. 
Then:
\begin{enumerate*}
\item\label{item.leaderlive.quorumsend.1}
$n\mentval\someoneAll\texi\tf{propose}\tnotor\everyoneAll\texi\tf{send}$.
\item\label{item.leaderlive.quorumsend.2}
$n\mentval\someoneAll\texi\tf{propose}\timpc\QuorumBox\texi\tf{send}$.
\end{enumerate*}
\end{lemm}
\begin{proof}
We consider each part in turn:
\begin{enumerate*}
\item
We use \weakmodusponens.
Suppose $n\mentval\modT\someoneAll\texi\tf{propose}$.
By \rulefont{PaxSend!} (since $n\geq\gslt(\avaluation)$) 
$$
\Forall{p\in\Pnt}n,p\mentval\texi\tf{send}.
$$
By Lemma~\ref{lemm.atomic.O}(\ref{item.atomic.pointwise.2}) $\texi\tf{send}$ is pointwise, so by Proposition~\ref{prop.pointwise.dense.char}(\ref{item.dense.char.pointwise.everyoneAll}) $n\mentval\everyoneAll\texi\tf{send}$ as required.
\item
We use \strongmodusponens.
Suppose $n\mentval\modT\someoneAll\texi\tf{propose}$.
By part~\ref{item.leaderlive.quorumsend.1} of this result and \weakmodusponens, $n\mentval\everyoneAll\texi\tf{send}$.
By \rulefont{PaxPCorrect} (for $\tf{send}$) $\acontext\ment\QuorumBox\correct{\tf{send}}$, so by 
Proposition~\ref{prop.correct.cobox.boxT}(\ref{item.correct.cobox.boxT.2})
$\acontext\ment\modT\QuorumBox\texi\tf{send}$ as required. 
\qedhere\end{enumerate*}
\end{proof}




\subsubsection{GSLT implies termination}


Recall from Definition~\ref{defn.correct}(\ref{item.correct}) that $\correct{\tf{decide}}=\tall v.\modTF\tf{decide}(v)$. 
\begin{lemm}
\label{lemm.someone.propose.someone.write}
Suppose $\avaluation$ is a valuation such that $\mentval\ThyPax$. 
Suppose $n\in\Time$ and $n\geq\gslt(\avaluation)$, and $p\in\Pnt$.
Then:
\begin{enumerate*}
\item\label{item.someone.propose.someone.write}
If $n,p\mentval\modT\texi\tf{propose}$ and $n,p\mentval\correct{\tf{write}}$ then $n,p\mentval\modT(\tf{leader}\tand\texi\tf{write})$.
\item\label{item.someone.propose.someone.write.live}
$\mentval\infinitely(\correct{\tf{decide}}\tand\modT\someoneAll\texi\tf{write})$.  
\item\label{item.someone.propose.someone.accept}
$\mentval\infinitely(\someoneAll(\tf{leader}\tand\correct{\tf{decide}})\tand\modT\texi\QuorumBox\tf{accept})$.
\end{enumerate*}
\end{lemm}
\begin{proof}
We consider each part in turn:
\begin{enumerate}
\item
Suppose $n,p\mentval\modT\texi\tf{propose}$.
By \rulefont{PaxPropose?} and \strongmodusponens $n,p\mentval\modT\tf{leader}$.
By Lemma~\ref{lemm.leaderlive.quorumsend}(\ref{item.leaderlive.quorumsend.2}) and strong modus ponens $n\mentval\modT\QuorumBox\texi\tf{send}$.
Using \weakmodusponens and \rulefont{PaxWrite!} (since $n\geq\gslt(\avaluation)$) $n,p\mentval\texi\tf{write}$.
By Proposition~\ref{prop.tv.ment.TF.model}(\ref{item.modTF.texi.P.valid}) (since $n,p\mentval\correct{\tf{write}}$) $n,p\mentval\modT\texi\tf{write}$. 
\item
By \rulefont{LdrExist}, and \rulefont{LdrCorrect} with Lemma~\ref{lemm.persist.char} (for $\f{tvs}=\{\tvT\}$ in that Lemma), and \strongmodusponens, 
for infinitely many $n\in\Time$ there exists $l_n\in\Pnt$ such that 
$$
n,l_n\ment\modT\tf{leader}\tand\correct{\tf{propose},\tf{write},\tf{decide}}. 
$$
So for each $n$ and $l_n$ we have $n,l_n\mentval\correct{\tf{decide}}$.
We will now check that $n,l_n\mentval\modT\texi\tf{write}$.
We reason as follows:
\begin{itemize*}
\item
By \rulefont{PaxPropose!} and \weakmodusponens we have that $n,l_n\ment\texi\tf{propose}$.
\item
By Proposition~\ref{prop.tv.ment.TF.model}(\ref{item.modTF.texi.P.valid}) (since $n,l_n\mentval\correct{\tf{propose}}$ and $n,l_n\mentval\texi\tf{propose}$) $n,l_n\mentval\modT\texi\tf{propose}$. 
\item
By part~\ref{item.someone.propose.someone.write} of this result (since $n,l_n\mentval\correct{\tf{write}}$) $n,l_n\mentval\modT\texi\tf{write}$.
\end{itemize*}
\item
We just combine part~\ref{item.someone.propose.someone.write.live} of this result with Lemma~\ref{lemm.someone.writes.after.gslt.accept}.
\qedhere\end{enumerate}
\end{proof}


Recall that Theorem~\ref{thrm.termination} expresses termination (as the name suggests) but also an abstracted form of the Integrity correctness property by virtue of the use of the $\infinitely$ modality, as per Remark~\ref{rmrk.paxos.integrity}:
\begin{thrm}[Termination]
\label{thrm.termination}
We have:
$$
\ThyPax\ment\infinitely\everyoneAll\texi\tf{decide} .
$$ 
\end{thrm}
\begin{proof}
Suppose $\avaluation$ is a valuation and $\mentval\ThyPax$.
By Lemmas~\ref{lemm.someone.propose.someone.write}(\ref{item.someone.propose.someone.accept}) and~\ref{lemm.persist.char} (for $\f{tvs}=\{\tvT\}$ in that Lemma), for infinitely many $n\in\Time$ there exists an $l_n\in\Pnt$ such that 
\begin{enumerate*}
\item
$n,l_n\mentval\tf{leader}$,
\item
$n,l_n\mentval\correct{\tf{decide}}$ (Definition~\ref{defn.correct}(\ref{item.correct})), and
\item
$n\mentval\modT\texi\QuorumBox\tf{accept}$.
\end{enumerate*}
By \rulefont{PaxDecideL!} (since $n\geq\gslt(\avaluation)$) 
and \weakmodusponens $n,l_n\mentval\texi\tf{decide}$.
By Proposition~\ref{prop.tv.ment.TF.model}(\ref{item.modTF.texi.P.valid}) (since $n,l_n\mentval\correct{\tf{decide}}$ and $n,l_n\mentval\texi\tf{decide}$) $n,l_n\mentval\modT\texi\tf{decide}$. 
Then by \rulefont{PaxDecide\tneg L!} and \weakmodusponens also $n,p\mentval\texi\tf{decide}$ for every $p\in\Pnt$, and so by (Lemma~\ref{lemm.atomic.O}(\ref{item.atomic.pointwise.2}) and) Proposition~\ref{prop.pointwise.dense.char}(\ref{item.dense.char.pointwise.everyoneAll}) $n\mentval\everyoneAll\texi\tf{decide}$. 

The above holds for infinitely many $n\in\Time$, so by Lemma~\ref{lemm.persist.char} (for $\f{tvs}=\{\tvT,\tvB\}$ in that Lemma) we have that $\mentval\infinitely\everyoneAll\texi\tf{decide}$ as required. 
\end{proof}


\begin{rmrk}
Theorem~\ref{thrm.termination} proves that $\ThyPax\ment\infinitely\everyoneAll\texi\tf{decide}$ holds.
Note that $\tvB$ (the truth-value associated to a crashed process) is a valid truth-value, so $n\ment\everyoneAll\texi\tf{decide}$ can hold in the trivial case that everyone is crashed.
Thus, what Theorem~\ref{thrm.termination} asserts in English is that infinitely often, everyone who is not crashed, decides (but it might be that everyone is crashed) --- from which it obviously follows that if a participant never crashes then eventually it decides a value, as per Remark~\ref{rmrk.standard.paxos.correctness}.

But in fact, we can prove a bit more:
$$
\ThyPax\ment\infinitely\QuorumBox\modT\texi\tf{decide}
$$
In words: infinitely often, there exists a \emph{quorum of uncrashed participants}, who decide.

This follows from Theorem~\ref{thrm.termination} by easy reasoning using \rulefont{PaxPCorrect}, \rulefont{PaxDecide\tneg L!}, and Proposition~\ref{prop.decide.n.agree}; we leave this to the reader to check. 
\end{rmrk}


\section{Simpler Declarative Paxos}
\label{sect.simpler.paxos}

Logic is useful for simplification and abstraction in many ways.
In particular, we can axiomatise something, prove some properties --- and then look at what we actually needed from the axioms to make the proofs work.
This can feed back to suggest refinements and generalisations of the original axioms.
(For example: many properties of natural numbers are actually properties of fields, or rings, or integral domains, and so on.) 

We can do something similar here, and look for the most elementary set of axioms that are actually required for our correctness proofs as written: 

\subsection{What we used in the correctness proofs for \ThyPax}
\label{subsect.reflection.on.correctness}
 
\begin{rmrk}
\label{rmrk.spax.explanation}
We note a curious thing about our correctness proofs in Section~\ref{sect.paxos.correctness.properties}, that we never used axioms \rulefont{PaxSend?} or \rulefont{PaxPropose01}.

How can this be?
Is this an error?

No.
This is an opportunity to further simplify the Declarative Paxos from Figure~\ref{fig.ThyPaxOne} to the \emph{Simpler Declarative Paxos} in Figure~\ref{fig.ThySPax}.
To understand how and why this works, we need to highlight some details about how the Paxos algorithm works.
\end{rmrk}

\begin{rmrk}[Why we did not use \rulefont{PaxSend?}]
Recall from Remark~\ref{rmrk.high-level.paxos} our high-level view of the Paxos algorithm, and recall that \ThyPax in Figure~\ref{fig.ThyPaxOne} is a declarative version of that algorithm (obtained by abstracting away explicit message-passing and algorithmic time, as per Remark~\ref{rmrk.elisions}).

Note of the algorithm that 
\begin{itemize*}
\item
in step~\ref{item.l.proposes} of Remark~\ref{rmrk.high-level.paxos}, the leader broadcasts a `propose' message; and 
\item
in step~\ref{item.p.sends}, a participant $p$ that receives that `propose' message, replies with a `send' message carrying as its payload a value $v_p$, which is $p$'s most recently accepted value (if any); and
\item
in step~\ref{item.l.writes}, the leader gathers the $v_p$ values from these responses and, if it has a quorum (in topological language: an open set) of responses, it selects a most recent value to use in its write step.
\end{itemize*}
In the algorithm, the only way that the leader $l$ can know of a participant $p$'s value of $v_p$, is if $l$ receives a send message from $p$ with payload $v_p$.
If $l$ receives no such message, then it does not know the value of $v_p$.

In contrast, predicates of \QLogic are not assigned truth-values by participants in an algorithm; they are assigned truth-values as if by an omniscient observer who is standing outside of the model and who has full (omniscient) access to it.
This is reflected in \rulefont{PaxWrite?}, where we write $\mru{(\someone\tf{accept})}{v}$.
This directly examines the model to find the most recently accepted value.
No reference here is required to the payload of the quorum of $\tf{send}$ messages.
\rulefont{PaxWrite?} still insists that such a quorum exists, as per the $\everyone\texi\tf{send}$ that appears in \rulefont{PaxWrite?}, so in this case the value $\mru{(\someone\tf{accept})}{v}$ is computable by the leader, and we could specify a more complex predicate that checks the set of $v_p$ payloads in the send messages of that quorum to select a most recently-accepted value.
In the implementation, this is necessary because the leader $l$ is not omniscient.
But in the logic we have no reason to go to this trouble, because we can just read the information directly from the model using $\mru{(\someone\tf{accept})}{v}$.\footnote{For example: if we want to compute $100!$ in the real world, we have to build a computer to do this for us, which involves many practical details including the construction of a semiconductor industry, compiler design and so on.  In logic, however, we can just say that it is $\prod_{i=1}^{100}i$, because that is what $100!$ \emph{is}.}

So our axioms are just doing their job, which is to be a declarative abstraction of the implementation. 
But since we do not use \rulefont{PaxSend?}, this is evidence that we may treat the $\tf{send}$ step as an implementational detail having to do with how information moves around the network, and elide it.
\end{rmrk}

\begin{rmrk}[Why we did not use \rulefont{PaxPropose01}]
This is due to another particularity of how the Paxos algorithm works.
The leader chooses one value in the propose step (or no value, if the leader has crashed), but it only uses that value in the case that it finds a quorum of participants who have never previously accepted any value (this is the right-hand disjunct in \rulefont{PaxWrite?}).
Just in this case, it writes the value that it chose in the propose step.

So what really matters is that the leader \emph{only writes one value}.
If the reader had proposed two values, this would be fine --- provided that it only writes one of them.

Leaders in the Paxos algorithm do not bother with multiple values in the propose step, because this would be a waste of time; they know they will only ever need at most one.
Thus, \rulefont{PaxPropose01} is right to reflect a true fact of the algorithm that it does not propose any more than the value than it will need, but this an algorithmic efficiency not a logical necessity. 
Likewise, the logic is right to point this out to us by \emph{not} requiring \rulefont{PaxPropose01} in the proofs: so long as only at most one value gets \emph{written} --- which is handled by \rulefont{PaxWrite01} --- none of our correctness properties are affected.
\end{rmrk}

\begin{rmrk}
By examining the proofs further, some further possible simplifications are revealed:
\begin{enumerate*}
\item
We can combine rules \rulefont{SPaxWrite01} and \rulefont{PaxWrite01} and \rulefont{LdrExt}, as per the discussion in Subsection~\ref{subsect.ldrexistuniq'}, to obtain an elegantly concise rule \rulefont{SPaxWrite01}.
\item
We used \rulefont{PaxPCorrect} in our proofs to obtain $\QuorumBox\correct{\tf P}$ only twice: 
\begin{itemize*}
\item
in Lemma~\ref{lemm.someone.writes.after.gslt.accept} for $\tf P=\tf{accept}$, and 
\item
in Lemma~\ref{lemm.leaderlive.quorumsend} for $\tf P=\tf{send}$.
\end{itemize*}
We have elided the \emph{send} step in $\Theta_\ThySPax$, so we only need to consider $\tf{accept}$.
Accordingly we take take \rulefont{SPaxPCorrect} to be just $\QuorumBox\correct{\tf{accept}}$.
\end{enumerate*}
\end{rmrk}

\begin{rmrk}
In conclusion, we see that the proofs in Section~\ref{sect.paxos.correctness.properties} are not only correctness proofs, but can also be read as an active commentary on the axioms.
Our logic has done what logic does, and has separated what is logically necessary from what is implementationally required or implementationally convenient.
This shows us how to simplify the axiomatisation (if we we wish) further away from the implementation, and closer towards what is logically essential: 
\end{rmrk}


\subsection{The definition}

\begin{figure}
$$
\begin{array}{l@{\quad}r@{\ }l}
\figunderline{Backward rules, for backward inference} 
\figskip
\rulefont{SPaxPropose?}&
& \tall v.\ \tf{propose}(v)\timpc\tf{leader}
\figskip
\rulefont{SPaxWrite?}&
& \tall v.\ \tf{write}(v) \timpc 
\begin{array}[t]{l}
\tf{leader}\tand\Quorum( 
\\
\quad\bigl((\mru{(\someone\tf{accept})}{v}) \tor
(\tf{propose}(v)\tand\tneg\modT\recent\texi\someone\tf{accept})\bigr)) 
\end{array}
\figskip
\rulefont{SPaxAccept?}&
& \tall v.\ \tf{accept}(v) \timpc \someoneAll\tf{write}(v)
\figskip
\rulefont{SPaxDecide?}&
& \tall v.\ \tf{decide}(v) \timpc \QuorumBox\tf{accept}(v)
\figskip\figskip
\figunderline{Forward rules, for making progress} 
\figskip
\rulefont{SPaxWrite!}&
& \final(\tf{leader}\tnotor\texi\tf{write})
\figskip
\rulefont{SPaxAccept!}&
& \final(\texi\someoneAll\tf{write}\tnotor \texi\tf{accept})
\figskip
\rulefont{SPaxDecide!}&
& \final(\texi\QuorumBox\tf{accept} \tnotor \texi\tf{decide})
\figskip
\figskip
\figunderline{Other rules, for correctness \& liveness}
\figskip
\rulefont{SLdrExist}&
&\modTF\tf{leader}\tand \modT\someoneAll\tf{leader}
\figskip
\rulefont{SLdrCorrect}&
& \infinitely\someoneAll(\tf{leader}\tand \correct{\tf{propose},\tf{write},\tf{decide}}) 
\figskip
\rulefont{SPaxPCorrect}&
& \QuorumBox\correct{\tf{accept}}
\figskip
\rulefont{SPaxWrite01}&
& \texiaffine\someoneAll\tf{write}
\end{array}
$$
\caption{Simpler Declarative Paxos (Definition~\ref{defn.thyspax})}
\label{fig.logical.paxos'}
\label{fig.ThySPax}
\end{figure}


\begin{defn}
\label{defn.thyspax}
We define \deffont{Simpler Declarative Paxos} to be the logical theory 
$$
\ThySPax=(\Sigma_{\ThySPax},\Theta_{\ThySPax}),
$$ 
where:
\begin{enumerate*}
\item
the signature $\Sigma_{\ThySPax}$ is 
$$
\Sigma_{\ThySPax}=[\tf{leader}:0;\ \tf{propose},\tf{write},\tf{accept},\tf{decide}:1] ,
$$ 
as per the notation in Definition~\ref{defn.signature}(\ref{item.predicate.syntax.signature}), and
\item
the axioms $\Theta_{\ThySPax}$ are as written in Figure~\ref{fig.ThySPax}.
\end{enumerate*}
\end{defn}

The same validity properties hold of \ThySPax from Figure~\ref{fig.ThySPax} as hold for \ThyPax from Figure~\ref{fig.ThyPaxOne}, namely: Validity (Theorem~\ref{thrm.validity}), Agreement (Theorem~\ref{thrm.full.agreement}), and Termination and Integrity (Theorem~\ref{thrm.termination}):
\begin{thrm}
\label{thrm.simpler.dp}
Suppose $\Sigma$ is a signature that includes $\Sigma_\ThySPax$, and suppose $\mathcal M$ is a model.
Then we have:
\begin{enumerate*}
\item
$\ThyPax\ment\tall v.(\tf{decide}(v) \timpc \urecent\someoneAll(\tf{leader}\tand\tf{propose}(v)))$.
\item
$\ThyPax\ment\tall v,v'.(\urecent\someoneAll\tf{decide}(v) \timpc \someoneAll\tf{decide}(v') \timpc v'\teq v)$.
\item
$\ThyPax\ment\infinitely\everyoneAll\texi\tf{decide}$.
\end{enumerate*}
\end{thrm}
\begin{proof}
This is by construction: from the discussion in Subsection~\ref{subsect.reflection.on.correctness}, the theory \rulefont{ThySPax} consists of just those parts of \rulefont{ThyPax} that are actually required in the correctness proofs.
\end{proof}

Thus, Theorem~\ref{thrm.simpler.dp} just notes that Simpler Declarative Paxos extracts from the proofs in Section~\ref{sect.paxos.correctness.properties} the logical essence of Figure~\ref{fig.ThyPaxOne}, as applied to the standard correctness properties of Paxos.
There is no right or wrong about whether we prefer $\ThyPax$ or $\ThySPax$; it depends on what level of abstraction we want to work with at any particular moment.
But, implicit in Theorem~\ref{thrm.simpler.dp} is an observation that \ThySPax is \emph{minimal}, because we put into it only that which we know we need to make some specific correctness proofs work.


\section{Conclusions and future work}
\label{sect.conclusions}

\subsection{Conclusions}







We have presented \QLogic, a three-valued modal fixedpoint logic for declarative specification of consensus algorithms, and we have shown how to use this to axiomatise Paxos, a canonical and nontrivial (though still relatively straightforward) consensus algorithm that can handle benign failures of participants.

We do not presented more complex consensus algorithms in this paper but we have looked into these in preliminary investigations, and they do seem to work. 

This paper is aimed at 
\begin{itemize*}
\item
logicians who might be interested in exploring a logic with a new and unusual combination of features inspired by strong practical motivations, and 
\item
researchers in distributed systems, who might consider (and benefit from) adopting such a framework.
\end{itemize*}
\QLogic offers benefits that could make it a useful addition to the toolkit of anyone working in distributed systems:
\begin{enumerate}
\item 
Much of what is shown in this paper just reflects existing practices, formalised into a new logic.
Therefore, readers may already be reasoning in ways that align with our logic, without explicitly realising it. 
For example, the third truth-value ($\tvB$) in \QLogic corresponds to typical side conditions like ``for every correct participant,'' and the backward axioms mirror the reverse inferences that practitioners often use. 

By formalising these familiar patterns, \QLogic provides a structured and precise way to represent and reason about concepts that are already informally expressed in the literature.
\item 
Logic is a concise, powerful, and unambiguous language. 
It distills paragraphs of English prose into a single line of precise symbols, enabling clear and rigorous reasoning. 

While explanatory English is helpful for accessibility and context, formal logic is the superior tool for working through the details of complex algorithms. 
It ensures precision and eliminates the ambiguities that often arise in informal descriptions.
\item 
\QLogic operates at a new (and higher) level of abstraction, which complements existing logic-based techniques. 
TLA+~\cite{DBLP:journals/toplas/Lamport94} provides a state-based view of distributed algorithms, whereas \QLogic abstracts lower-level implementation details while retaining the essential behaviours and properties of the algorithms.
This abstraction does not obscure the logic of the system; on the contrary we have found that it clarifies it, making the underlying mechanisms and invariants more apparent. 
\item 
One strength of logic is its ability to abstract away what is irrelevant while maintaining precision about what is important. 
In \QLogic, constructs like strong and weak implications ($\timpc$ and $\tnotor$) exemplify this principle, capturing subtle distinctions in meaning with precise formal symbols. 
These constructs distill complex ideas into simple yet powerful tools that can be manipulated symbolically and reasoned about systematically. 

While mastering these constructs requires some initial effort, the benefits of precision and clarity justify the investment.
\item 
Learning \QLogic requires some effort, but for researchers already familiar with (for example) TLA+, the transition to \QLogic should be eminently manageable. 

While axiomatisations in \QLogic are certainly subtle, this is not gratuitous: they reflect the inherent nature of the distributed systems they are designed to model. 
By formalising these complexities, \QLogic captures their structure and reveals and addresses the challenges that are \emph{already present} in the underlying algorithms. 

This transparency makes it a powerful tool for understanding and reasoning about distributed protocols.\footnote{Certainly, the first author (who has a background in logic) has found this study invaluable for gaining a better understanding of consensus algorithms.  It seems likely that other readers might benefit similarly from using \QLogic.}
\item 
Distributed algorithms, and consensus protocols in particular, are foundational to the functionality of blockchain systems and other critical infrastructures. 

These algorithms are recognised to be notoriously challenging to design and implement correctly. 
Ambiguities and informal specifications can lead to vulnerabilities and failures. 

\QLogic bridges the gap between declarative specifications of invariants, and practical verification; enabling researchers to clarify, reason about, and verify the properties of distributed protocols rigorously. 
This is especially valuable in blockchain systems, where clear and robust specifications are essential for trust and reliability.
\end{enumerate}

In conclusion: by providing a formal framework that combines the abstraction of declarative reasoning with the rigour of logical precision, \QLogic offers researchers in distributed systems a powerful tool for modeling, understanding, and verifying complex algorithms. 

\subsection{Related and future work}

In this paper, our focus has been to lay the foundational theory and framework of \QLogic and to demonstrate its application to one of the most well-known consensus protocols, Paxos~\cite{DBLP:journals/tocs/Lamport98}.

Paxos is a sensible starting point, being a cornerstone of distributed systems, but (certainly by modern standards!) it is relatively straightforward.
For example, Paxos only addresses crash failures, limiting its applicability to systems that assume benign (rather than hostile) faults.

A natural direction for future work involves extending our framework to handle more complex and realistic fault models, such as Byzantine faults, where participants can deviate arbitrarily from the algorithm.
Consensus protocols that address Byzantine faults, such as Practical Byzantine Fault Tolerance (PBFT)~\cite{DBLP:conf/osdi/CastroL99} and HotStuff~\cite{DBLP:conf/podc/YinMRGA19}, provide robust mechanisms for fault tolerance and serve as the foundation for many modern blockchain systems, including Tendermint (\url{https://tendermint.com})~\cite{DBLP:journals/corr/abs-1807-04938}. 

Applying \QLogic to these protocols would see how it handles the complexity introduced by arbitrary faults, which is necessary to model and verify more sophisticated algorithms.
Preliminary exploration by the authors in this direction indicates that \QLogic works very well in these more complex situations --- it can indeed axiomatise algorithms that take account of arbitrary, including hostile, faults.
Unpacking this in full is for future papers.

Recent research has explored DAG-based consensus protocols, which deviate from the traditional linear blockchain structure. 
In these protocols, transactions are organized into a directed acyclic graph (DAG) rather than a strict chain, enabling greater parallelism and scalability. 
Protocols such as DAG-rider~\cite{DBLP:conf/podc/KeidarKNS21}, Narwhal and Tusk~\cite{DBLP:conf/eurosys/DanezisKSS22}, and Bullshark~\cite{DBLP:conf/ccs/SpiegelmanGSK22} represent cutting-edge advancements in this domain. 
Extending \QLogic to analyse DAG-based protocols would be another advance, as these systems often involve more intricate interactions between components and require reasoning about concurrent and asynchronous behaviour at a much higher level of complexity.

Another promising avenue for future research involves leveraging the modularity and expressiveness of \QLogic to study hybrid protocols~\cite{DBLP:conf/sp/NeuTT21} that combine elements of classical consensus with newer techniques, such as Gasper~\cite{DBLP:journals/corr/abs-2003-03052} and others~\cite{DBLP:conf/esorics/DAmatoZ23,DBLP:journals/corr/abs-2411-00558}

It would be helpful in future work to develop automated tools for \QLogic: the tools developed for TLA+~\cite{DBLP:books/aw/Lamport2002} could show the way.
TLA+ is a formal specification language and framework designed to model, specify and verify properties of complex systems; particularly distributed algorithms. 
Its TLC model checker~\cite{DBLP:conf/charme/YuML99} provides exhaustive state exploration to verify properties such as safety and liveness, and the PlusCal translator simplifies specification by offering a high-level pseudocode-like interface.
Additionally, the TLA+ Toolbox~\cite{DBLP:journals/corr/abs-1912-10633} integrates these capabilities into a development environment, supporting workflows that are valuable in both research and industrial applications.

Adapting similar tools to \QLogic could enable automated verification of correctness properties for consensus protocols, including invariants and liveness guarantees. 
For example: a higher-level language akin to PlusCal could simplify the specification of protocols using the axiomatic foundations of \QLogic.
Furthermore, leveraging techniques such as the symbolic model checking enhancements presented by Otoni et al.~\cite{DBLP:conf/tacas/OtoniKKES23} or integrating modular proof systems like TLAPS~\cite{DBLP:conf/fm/CousineauDLMRV12} could further enhance the usability and robustness of the framework.

Such tools would streamline the application of \QLogic to real-world systems, and also make it accessible to a broader audience of researchers and practitioners. 
This direction represents an exciting convergence of theoretical advancements and practical tool development, providing a robust framework for analysing both classical and emerging consensus protocols.

Adopting \QLogic requires some investment in time and effort, but the benefits in clarity, correctness, and reliability may be worthwhile.
Consensus algorithms are known to be particularly challenging, so our relatively high-level declarative approach could help to manage (and debug) the complexity of these algorithms.

We hope this work will inspire readers %
to explore and adopt this new perspective, contributing to more robust and trustworthy distributed protocols in the future.



\input{paxos.bbl}

\end{document}


\end{document}


\end{document}


\end{document}
